\documentclass[12pt]{article}

% ===== 切換模式：只保留一行未註解 =====
% \def\docmode{chinese}
% \def\docmode{german}
\def\docmode{bilingual}

% ===== 雙語對照橫式紙張：只保留一行未註解 =====
\def\bilingualpaper{a4}
% \def\bilingualpaper{a3}

% bverfge_layout.tex
% 共用排版基礎設施，供 Schwangerschaftsabbruch I 與 II 共享。
% 由各文件以 \documentclass 後 % bverfge_layout.tex
% 共用排版基礎設施，供 Schwangerschaftsabbruch I 與 II 共享。
% 由各文件以 \documentclass 後 % bverfge_layout.tex
% 共用排版基礎設施，供 Schwangerschaftsabbruch I 與 II 共享。
% 由各文件以 \documentclass 後 \input{../../共用LaTeX/bverfge_layout} 載入。
% 不含：\documentclass、\documentversion、\ggversionde/zh、\tocde/\toczh。
% 日期與首頁版本列由本檔自動產生；各文件只需手動指定 \documentversion。

\usepackage{xeCJK}
\usepackage{fontspec}
\usepackage{geometry}
\usepackage{setspace}
\usepackage{hyperref}
\usepackage{xurl}
\usepackage{bookmark}
\usepackage{enumitem}
\usepackage{xcolor}
\usepackage{titlesec}
\usepackage{array}
\usepackage{longtable}
\usepackage{tcolorbox}
\usepackage{environ}
\usepackage{pdflscape}
\usepackage{etoolbox}
\usepackage{ragged2e}
\usepackage{paracol}
\usepackage{multicol}
\usepackage{needspace}
\usepackage{fancyhdr}
\usepackage{tikz}
\usepackage{zref-abspage}
\tcbuselibrary{breakable,skins}
\usetikzlibrary{calc}
\usepackage{pgfcalendar}

\hypersetup{
  unicode=true,
  bookmarksopen=true,
  bookmarksnumbered=false,
  hidelinks
}

% ===== 自動推導，不要手動改下面 =====
% （\docmode 與 \bilingualpaper 由各文件在 \input{bverfge_layout} 之前設定；
%   預設 chinese / a4）
\ifdefined\docmode\else\def\docmode{chinese}\fi
\ifdefined\bilingualpaper\else\def\bilingualpaper{a4}\fi
\newif\ifshoworiginal
\newif\ifshowtranslation
\newif\ifbilinguallandscape
\newif\ifbilingualcolumns
\newif\ifintranslation
\newif\iffirstbilingualsection

\showoriginalfalse
\showtranslationfalse
\bilinguallandscapefalse
\bilingualcolumnsfalse
\intranslationfalse
\firstbilingualsectionfalse

\def\modechinese{chinese}
\def\modegerman{german}
\def\modebilingual{bilingual}
\def\bilingualpaperaThree{a3}
\def\bilingualpaperaFour{a4}

\ifx\docmode\modechinese
  \showtranslationtrue
\fi

\ifx\docmode\modegerman
  \showoriginaltrue
\fi

\ifx\docmode\modebilingual
  \showoriginaltrue
  \showtranslationtrue
  \bilinguallandscapetrue
  \bilingualcolumnstrue
\fi

% 編排模式說明：
% chinese  → \showoriginalfalse  \showtranslationtrue  （A4 直式，純中文）
% german   → \showoriginaltrue   \showtranslationfalse （A4 直式，純德文）
% bilingual→ \showoriginaltrue   \showtranslationtrue  \bilingualcolumnstrue（A3/A4 橫式對照）

\ifbilingualcolumns
  \ifx\bilingualpaper\bilingualpaperaFour
    \geometry{
      a4paper,
      landscape,
      left=1.25cm,
      right=2.25cm,
      top=1.25cm,
      bottom=1.75cm,
      footskip=8.5mm,
      marginparwidth=0.75cm,
      marginparsep=0.25cm
    }
  \else
    \geometry{
      a3paper,
      landscape,
      left=1.55cm,
      right=2.75cm,
      top=1.55cm,
      bottom=2.05cm,
      footskip=9.5mm,
      marginparwidth=0.9cm,
      marginparsep=0.35cm
    }
  \fi
\else
\geometry{
  a4paper,
  portrait,
  left=2.2cm,
  right=2.6cm,
  top=2.0cm,
  bottom=2.0cm,
  marginparwidth=0.8cm,
  marginparsep=0.3cm
}
\fi
% ===== 時間戳基礎設施 =====
\newcount\stamptimeval
\newcount\stampminuteval
\newcommand{\stamphh}{%
  \stamptimeval=\time
  \divide\stamptimeval by 60
  \ifnum\stamptimeval<10 0\fi
  \the\stamptimeval}
\newcommand{\stampmm}{%
  \stamptimeval=\time
  \stampminuteval=\time
  \divide\stampminuteval by 60
  \multiply\stampminuteval by 60
  \advance\stamptimeval by -\stampminuteval
  \ifnum\stamptimeval<10 0\fi
  \the\stamptimeval}
\newcommand{\stampwkde}[1]{%
  \ifcase#1Montag\or Dienstag\or Mittwoch\or Donnerstag\or Freitag\or Samstag\or Sonntag\fi}
\newcommand{\stampwkzh}[1]{%
  \ifcase#1 一\or 二\or 三\or 四\or 五\or 六\or 日\fi}
\newcommand{\stampmonthde}[1]{%
  \ifcase#1\or Januar\or Februar\or März\or April\or Mai\or Juni\or
  Juli\or August\or September\or Oktober\or November\or Dezember\fi}
\newcount\stampjdval
\newcount\stampwdval
\pgfcalendardatetojulian{\the\year-\the\month-\the\day}{\stampjdval}
\pgfcalendarjuliantoweekday{\the\stampjdval}{\stampwdval}
\newcommand{\timestampzh}{%
  \songfont 最後修改：\the\year 年\the\month 月\the\day 日%
  % （星期\stampwkzh{\the\stampwdval}）%
  % \space\stamphh:\stampmm
  }

\newcommand{\documentdatezh}{\the\year 年 \the\month 月 \the\day 日}
\newcommand{\documentdatede}{\the\day. \stampmonthde{\the\month} \the\year}
\newcommand{\bverfgetranslationnote}{%
  {\songfont 翻譯說明：初譯使用 Claude Code 與 GPT 協助迭代；仍請以德文原文為準。}%
}
\newcommand{\bverfgecourseinfo}{%
  {\songfont 課程：114 學年度第 2 學期；憲法專題研究（四）；授課老師：湯德宗教授；導讀日期：115 年 6 月 1 日。}%
}
\newcommand{\bverfgeeditorinfo}{%
  {\songfont 編者：王逸帆（R10A21126），國立台灣大學法律系研究所碩士班。}%
}
\newcommand{\bverfgetitleversionline}{%
  \ifbilingualcolumns
    Fassung \documentversion\enspace\textperiodcentered\enspace Stand: \documentdatede
    \quad
    {\songfont 版本 \documentversion\enspace\textperiodcentered\enspace 修訂日期：\documentdatezh\par}%
  \else
    \ifshowtranslation
      {\songfont 版本 \documentversion\enspace\textperiodcentered\enspace 修訂日期：\documentdatezh\par}%
    \else
      Fassung \documentversion\enspace\textperiodcentered\enspace Stand: \documentdatede\par
    \fi
  \fi
}
\newcommand{\bverfgetitlecornerinfo}{%
  \ifshowtranslation
    \vspace{0.72em}%
    \noindent\hfill
    \begin{minipage}{0.66\textwidth}
      \raggedleft\tiny\color{pendingink}\songfont
      \bverfgecourseinfo\par
      \bverfgeeditorinfo
    \end{minipage}%
  \fi
}
\newcommand{\bverfgetitlemeta}{%
  \vfill
  \begin{center}%
    \footnotesize\color{pendingink}%
    \bverfgetitleversionline
    \ifshowtranslation
      \vspace{0.28em}{\scriptsize\bverfgetranslationnote\par}%
    \fi
  \end{center}%
  \bverfgetitlecornerinfo
}

\pagestyle{fancy}
\fancyhf{}
\renewcommand{\headrulewidth}{0pt}
\renewcommand{\footrulewidth}{0pt}
\fancyfoot[C]{\thepage}
\ifbilingualcolumns
  \fancyfoot[R]{\scriptsize\color{pendingink}\timestampzh}%
\else
  \ifshoworiginal
  \else
    \fancyfoot[R]{\scriptsize\color{pendingink}\timestampzh\hspace{0.45cm}}%
  \fi
\fi
\setmainfont{Times New Roman}
\setsansfont{Helvetica Neue}
\setmonofont{Menlo}
\IfFileExists{/Users/iw/Library/Fonts/kaiu.ttf}{
  \setCJKmainfont[Path=/Users/iw/Library/Fonts/,AutoFakeBold=4]{kaiu.ttf}
}{
  \IfFontExistsTF{標楷體}{
    \setCJKmainfont[AutoFakeBold=4]{標楷體}
  }{
    \setCJKmainfont{Songti TC}
  }
}
\IfFontExistsTF{Songti TC}{
  \setCJKfamilyfont{song}{Songti TC}
  \setCJKmonofont{Songti TC}
}{
  \setCJKfamilyfont{song}[Path=/Users/iw/Library/Fonts/,AutoFakeBold=4]{kaiu.ttf}
  \setCJKmonofont[Path=/Users/iw/Library/Fonts/,AutoFakeBold=4]{kaiu.ttf}
}
\newcommand{\songfont}{\CJKfamily{song}}

% ===== 封面／目錄專用打字機字體（僅限封面、目錄頁使用，不影響正文）=====
\IfFontExistsTF{Radio Newsman}{%
  \newfontfamily\bverfgetitlede{Radio Newsman}%
}{%
  \newfontfamily\bverfgetitlede{Courier New}%
}
\IfFontExistsTF{Erikas Farbband}{%
  \newfontfamily\bverfgebodyde{Erikas Farbband}%
}{%
  \let\bverfgebodyde\bverfgetitlede
}
\IfFontExistsTF{Maoken Heavy Labourer}{%
  \setCJKfamilyfont{bverfgetitlecjk}{Maoken Heavy Labourer}%
}{%
  \setCJKfamilyfont{bverfgetitlecjk}{Songti TC}%
}
\IfFontExistsTF{Huiwen-mincho}{%
  \setCJKfamilyfont{bverfgebodycjk}{Huiwen-mincho}%
}{%
  \setCJKfamilyfont{bverfgebodycjk}{Songti TC}%
}
\newcommand{\bverfgetitlezh}{\bverfgetitlede\CJKfamily{bverfgetitlecjk}}
\newcommand{\bverfgebodyzh}{\bverfgebodyde\CJKfamily{bverfgebodycjk}}

\setstretch{1.35}
\setlist{itemsep=0.2em, topsep=0.4em}
\setlength{\parindent}{0pt}
\setlength{\parskip}{0.45em}
\raggedbottom
\setcounter{secnumdepth}{0}
\setcounter{tocdepth}{2}

\newcommand{\lawboxsourcefont}{\footnotesize}
\newcommand{\lawboxtransfont}{\footnotesize}
\newcommand{\lawboxsourcestretch}{1.02}
\newcommand{\lawboxtransstretch}{0.92}

\titleformat{\section}{\centering\Large\bfseries\songfont}{\thesection}{1em}{}
\titleformat{\subsection}{\large\bfseries\songfont}{\thesubsection}{1em}{}
\titleformat{\subsubsection}{\normalsize\bfseries\songfont}{\thesubsubsection}{1em}{}
\titleformat{\paragraph}{\normalsize\bfseries\songfont}{\theparagraph}{1em}{}
\titlespacing*{\section}{0pt}{1.8em}{1.0em}
\titlespacing*{\subsection}{0pt}{1.2em}{0.6em}
\titlespacing*{\subsubsection}{0pt}{1.0em}{0.45em}
\titlespacing*{\paragraph}{0pt}{0.6em}{0.4em}

\definecolor{noteblue}{HTML}{A9C9F5}
\definecolor{noteteal}{HTML}{AEE1D6}
\definecolor{notegray}{HTML}{D7DADF}
\definecolor{caseink}{HTML}{000000}
\definecolor{casepaper}{HTML}{FCFEFD}
\definecolor{sourcepaper}{HTML}{FAFBF8}
\definecolor{sourceline}{HTML}{D7DED8}
\definecolor{sourceink}{HTML}{000000}
\definecolor{pendingink}{HTML}{000000}

\tcbset{
  caseboxbase/.style={
    enhanced,
    breakable,
    width=0.985\linewidth,
    colback=casepaper,
    coltitle=caseink,
    fonttitle=\bfseries\songfont,
    boxrule=0.28pt,
    arc=1.2mm,
    left=2.4mm,
    right=2.4mm,
    top=1.5mm,
    bottom=1.5mm,
    boxsep=1.5mm,
    before skip=0.65em,
    after skip=0.65em,
    before upper={\setlength{\parindent}{0pt}\setlength{\parskip}{0.35em}},
  }
}

\newcommand{\bilingualstart}{}
\newcommand{\bilingualstop}{}
\ifbilingualcolumns
  \renewcommand{\bilingualstart}{\small\setlength{\columnsep}{1.25cm}\setlength{\columnseprule}{0.12pt}\columnratio{0.5,0.5}\multicolumnfootnotes\flushbottom\sloppy\interlinepenalty=80\clubpenalty=800\widowpenalty=800\displaywidowpenalty=800\begin{paracol}{2}\global\intranslationfalse\global\firstbilingualsectiontrue{\footnotesize\color{sourceink}Deutsch\par}\switchcolumn{\footnotesize\songfont\color{sourceink}中文譯文\par}\switchcolumn*}
  \renewcommand{\bilingualstop}{\ifintranslation\switchcolumn*\global\intranslationfalse\fi\end{paracol}\clearpage\raggedbottom}
\fi
\newif\ifrandnumbers
\randnumbersfalse
\newcounter{randnumber}
\newcounter{randnumberlocal}
\newcounter{lastsourcerandstart}
\newif\iflastsourcehasrandnumbers
\lastsourcehasrandnumbersfalse
\newif\ifinlawbox
\inlawboxfalse
\newif\iflawboxhaspendingrn
\lawboxhaspendingrnfalse
\newcommand{\lawboxpendingrn}{}
\newcommand{\startrandnumbers}{\global\randnumberstrue\setcounter{randnumber}{1}\global\lastsourcehasrandnumbersfalse}
\newcommand{\stoprandnumbers}{\global\randnumbersfalse\global\lastsourcehasrandnumbersfalse}
\newlength{\randnumberwidth}
\setlength{\randnumberwidth}{0.95cm}
\newlength{\randnumberrightoffset}
\setlength{\randnumberrightoffset}{0.85cm}
\newlength{\randnumberlawboxrightoffset}
\setlength{\randnumberlawboxrightoffset}{1.40cm}
\newcommand{\randnumberbox}[1]{%
  \leavevmode
  \makebox[0pt][l]{%
    \smash{%
      \tikz[remember picture,overlay,baseline=(rnmark.base)]{%
        \coordinate (rnmark) at (0,0);%
        \ifinlawbox
          \path let \p1=(current page.east), \p2=(rnmark.base) in
            node[
              anchor=base east,
              inner sep=0pt,
              outer sep=0pt,
              text width=\randnumberwidth,
              align=right,
              text=pendingink,
              font=\normalfont\mdseries\scriptsize\ttfamily
            ] at (\x1-\randnumberlawboxrightoffset,\y2) {{\normalfont\mdseries\scriptsize\ttfamily #1}};%
        \else
          \path let \p1=(current page.east), \p2=(rnmark.base) in
            node[
              anchor=base east,
              inner sep=0pt,
              outer sep=0pt,
              text width=\randnumberwidth,
              align=right,
              text=pendingink,
              font=\normalfont\mdseries\scriptsize\ttfamily
            ] at (\x1-\randnumberrightoffset,\y2) {{\normalfont\mdseries\scriptsize\ttfamily #1}};%
        \fi
      }%
    }%
  }%
  \ignorespaces
}
\newcommand{\lawboxstorerandnumber}[1]{%
  \gdef\lawboxpendingrn{#1}%
  \global\lawboxhaspendingrntrue
  \ignorespaces
}
\newcommand{\lawboxuserandnumber}{%
  \iflawboxhaspendingrn
    \randnumberbox{\lawboxpendingrn}%
    \global\lawboxhaspendingrnfalse
  \fi
}
\newcommand{\sourcern}[1]{%
  \ifrandnumbers
    \ifshoworiginal
      \ifshowtranslation\else
        \ifinlawbox\lawboxstorerandnumber{#1}\else\randnumberbox{#1}\fi
      \fi
    \fi
  \fi
}
\newcommand{\transrn}[1]{%
  \ifrandnumbers
    \ifshowtranslation
      \ifinlawbox\lawboxstorerandnumber{#1}\else\randnumberbox{#1}\fi
    \fi
  \fi
}
% ===== 課堂模式：德文原文縮寫註解 =====
% 日常切換只改各判決檔的一行：
%   \classroommodeon        開啟課堂版；註解掉這行即回到一般版。
% 模板預設：
%   1. 課堂模式關閉。
%   2. 課堂模式開啟時，縮寫註解包含「德文全稱＋中文譯名」。
%   3. 課堂模式開啟時，GG/條文編者註仍會自動顯示。
% 進階微調才需要在單案檔覆寫：
%   \classroomabbrzhoff     縮寫註解僅顯示德文全稱。
%   \showeditornotestrue    一般版也顯示編者註。
% 註解策略：同一實際 PDF 頁面中，同一縮寫只在第一次出現處加註。
% 這裡用 zref-abspage 的絕對頁序，不用 \thepage，避免文件內頁碼重置時誤判。
\newif\ifclassroommode
\newif\ifclassroomabbrzh
\classroommodefalse
\classroomabbrzhtrue
\newcommand{\classroommodeon}{\classroommodetrue}
\newcommand{\classroommodeoff}{\classroommodefalse}
\newcommand{\classroomabbrzhon}{\classroomabbrzhtrue}
\newcommand{\classroomabbrzhoff}{\classroomabbrzhfalse}
\newcommand{\abbrnotelabel}{\emph{Abk.}: }
\newcommand{\abbrnote}[3]{%
  \ifclassroommode
    \ifshoworiginal
      \ifintranslation\else
        \footnote{\normalfont\footnotesize\abbrnotelabel #2\ifclassroomabbrzh；{\songfont #3}\fi}%
      \fi
    \fi
  \fi
}
\newcommand{\abbrnoteonce}[4]{%
  #2%
  \ifclassroommode
    \ifshoworiginal
      \ifintranslation\else
        \begingroup
          \edef\abbrcurrentabspage{\number\numexpr\value{abspage}+1\relax}%
          \ifcsname abbrnoteseen@#1@\abbrcurrentabspage\endcsname\else
            \global\expandafter\let\csname abbrnoteseen@#1@\abbrcurrentabspage\endcsname\empty
            \abbrnote{#1}{#3}{#4}%
          \fi
        \endgroup
      \fi
    \fi
  \fi
}
\newcommand{\abbrAbs}{\abbrnoteonce{abs}{Abs.}{Absatz}{項}}
\newcommand{\abbrAAO}{\abbrnoteonce{aao}{a.a.O.}{am angegebenen Ort}{前引處}}
\newcommand{\abbrAE}{\abbrnoteonce{ae}{AE}{Alternativ-Entwurf}{替代草案}}
\newcommand{\abbrArt}{\abbrnoteonce{art}{Art.}{Artikel}{條}}
\newcommand{\abbrAufl}{\abbrnoteonce{aufl}{Aufl.}{Auflage}{版}}
\newcommand{\abbrBd}{\abbrnoteonce{bd}{Bd.}{Band}{卷}}
\newcommand{\abbrBGBl}{\abbrnoteonce{bgbl}{BGBl.}{Bundesgesetzblatt}{聯邦法律公報}}
\newcommand{\abbrBGHSt}{\abbrnoteonce{bghst}{BGHSt}{Entscheidungen des Bundesgerichtshofes in Strafsachen}{聯邦普通法院刑事判決彙編}}
\newcommand{\abbrBRDrucks}{\abbrnoteonce{brdrucks}{BRDrucks.}{Bundesrats-Drucksache}{聯邦參議院文件}}
\newcommand{\abbrBTDrucks}{\abbrnoteonce{btdrucks}{BTDrucks.}{Bundestags-Drucksache}{聯邦議院文件}}
\newcommand{\abbrBundesgesetzbl}{\abbrnoteonce{bundesgesetzbl}{Bundesgesetzbl.}{Bundesgesetzblatt}{聯邦法律公報}}
\newcommand{\abbrBvF}{\abbrnoteonce{bvf}{BvF}{Bundesverfassungsgerichtsverfahren}{聯邦憲法法院程序案號}}
\newcommand{\abbrBvL}{\abbrnoteonce{bvl}{BvL}{Bundesverfassungsgerichtsverfahren}{聯邦憲法法院程序案號}}
\newcommand{\abbrBvR}{\abbrnoteonce{bvr}{BvR}{Bundesverfassungsgerichtsverfahren}{聯邦憲法法院程序案號}}
\newcommand{\abbrBVerfG}{\abbrnoteonce{bverfg}{BVerfG}{Bundesverfassungsgericht}{聯邦憲法法院}}
\newcommand{\abbrBVerfGE}{\abbrnoteonce{bverfge}{BVerfGE}{Entscheidungen des Bundesverfassungsgerichts}{聯邦憲法法院判決集}}
\newcommand{\abbrBVerfGG}{\abbrnoteonce{bverfgg}{BVerfGG}{Bundesverfassungsgerichtsgesetz}{聯邦憲法法院法}}
\newcommand{\abbrDDR}{\abbrnoteonce{ddr}{DDR}{Deutsche Demokratische Republik}{德意志民主共和國；東德}}
\newcommand{\abbrDJT}{\abbrnoteonce{djt}{DJT}{Deutscher Juristentag}{德國法學家大會}}
\newcommand{\abbrDr}{\abbrnoteonce{dr}{Dr.}{Doktor}{Dr.}}
\newcommand{\abbrEuGRZ}{\abbrnoteonce{eugrz}{EuGRZ}{Europäische Grundrechte-Zeitschrift}{歐洲基本權期刊}}
\newcommand{\abbrEV}{\abbrnoteonce{ev}{EV}{Einigungsvertrag}{統一條約}}
\newcommand{\abbrF}{\abbrnoteonce{f}{f.}{folgende}{以下一頁；下一段}}
\newcommand{\abbrFF}{\abbrnoteonce{ff}{ff.}{fortfolgende}{以下數頁；以下各段}}
\newcommand{\abbrGBlDDR}{\abbrnoteonce{gblddr}{GBl. DDR}{Gesetzblatt der Deutschen Demokratischen Republik}{德意志民主共和國法律公報}}
\newcommand{\abbrGG}{\abbrnoteonce{gg}{GG}{Grundgesetz}{基本法}}
\newcommand{\abbrGVBl}{\abbrnoteonce{gvbl}{GVBl.}{Gesetz- und Verordnungsblatt}{法律及命令公報}}
\newcommand{\abbrJR}{\abbrnoteonce{jr}{JR}{Juristische Rundschau}{法學評論}}
\newcommand{\abbrJZ}{\abbrnoteonce{jz}{JZ}{JuristenZeitung}{法學家報}}
\newcommand{\abbrIVm}{\abbrnoteonce{ivm}{i.V.m.}{in Verbindung mit}{結合；連同}}
\newcommand{\abbrMdB}{\abbrnoteonce{mdb}{MdB}{Mitglied des Bundestages}{聯邦議院議員}}
\newcommand{\abbrMwN}{\abbrnoteonce{mwn}{m.w.N.}{mit weiteren Nachweisen}{附進一步文獻／判例出處}}
\newcommand{\abbrNF}{\abbrnoteonce{nf}{n.F.}{neue Fassung}{新版本}}
\newcommand{\abbrAF}{\abbrnoteonce{af}{a.F.}{alte Fassung}{舊版本}}
\newcommand{\abbrNr}{\abbrnoteonce{nr}{Nr.}{Nummer}{款、目或號，依語境}}
\newcommand{\abbrProf}{\abbrnoteonce{prof}{Prof.}{Professor}{教授}}
\newcommand{\abbrRdnr}{\abbrnoteonce{rdnr}{Rdnr.}{Randnummer}{邊碼；段碼}}
\newcommand{\abbrRdnrn}{\abbrnoteonce{rdnrn}{Rdnrn.}{Randnummern}{邊碼；段碼}}
\newcommand{\abbrRGBl}{\abbrnoteonce{rgbl}{RGBl.}{Reichsgesetzblatt}{帝國法律公報}}
\newcommand{\abbrRGSt}{\abbrnoteonce{rgst}{RGSt}{Entscheidungen des Reichsgerichts in Strafsachen}{帝國法院刑事判決彙編}}
\newcommand{\abbrRspr}{\abbrnoteonce{rspr}{Rspr.}{Rechtsprechung}{判例；司法實務}}
\newcommand{\abbrRVO}{\abbrnoteonce{rvo}{RVO}{Reichsversicherungsordnung}{帝國保險法}}
\newcommand{\abbrS}{\abbrnoteonce{s}{S.}{Seite}{頁}}
\newcommand{\abbrSFHG}{\abbrnoteonce{sfhg}{SFHG}{Schwangeren- und Familienhilfegesetz}{孕婦及家庭扶助法}}
\newcommand{\abbrSGBV}{\abbrnoteonce{sgbv}{SGB V}{Fünftes Buch Sozialgesetzbuch}{社會法典第五編}}
\newcommand{\abbrStAG}{\abbrnoteonce{staeg}{StÄG}{Strafrechtsänderungsgesetz}{刑法修正法}}
\newcommand{\abbrStenBer}{\abbrnoteonce{stenber}{StenBer.}{Stenographischer Bericht}{速記紀錄}}
\newcommand{\abbrStGB}{\abbrnoteonce{stgb}{StGB}{Strafgesetzbuch}{刑法}}
\newcommand{\abbrStREG}{\abbrnoteonce{streg}{StREG}{Strafrechtsreform-Ergänzungsgesetz}{刑法改革補充法}}
\newcommand{\abbrStrRG}{\abbrnoteonce{strrg}{StrRG}{Strafrechtsreformgesetz}{刑法改革法}}
\newcommand{\abbrUS}{\abbrnoteonce{us}{U.S.}{United States Reports}{美國最高法院判決彙編}}
\newcommand{\abbrVgl}{\abbrnoteonce{vgl}{vgl.}{vergleiche}{參見}}
\newcommand{\abbrVVDStRL}{\abbrnoteonce{vvdstrl}{VVDStRL}{Veröffentlichungen der Vereinigung der Deutschen Staatsrechtslehrer}{德國公法教師協會論文集}}
\newcommand{\abbrWP}{\abbrnoteonce{wp}{WP.}{Wahlperiode}{屆期}}
\newcommand{\abbrWp}{\abbrnoteonce{wp}{Wp.}{Wahlperiode}{屆期}}
\newcommand{\abbrZStrW}{\abbrnoteonce{zstrw}{ZStrW}{Zeitschrift für die gesamte Strafrechtswissenschaft}{整體刑法學期刊}}
\newcommand{\recordsourcerandstart}{%
  \ifrandnumbers
    \setcounter{lastsourcerandstart}{\value{randnumber}}%
    \global\lastsourcehasrandnumberstrue
  \else
    \global\lastsourcehasrandnumbersfalse
  \fi
}
\newlength{\biparagraphskip}
\setlength{\biparagraphskip}{0.52em}
\newlength{\lawboxblockbeforeskip}
\setlength{\lawboxblockbeforeskip}{0.72em}
\newlength{\lawboxblockafterskip}
\setlength{\lawboxblockafterskip}{0.92em}
\newif\iflawboxpartendedblock
\lawboxpartendedblockfalse
\newcommand{\sourceparstyle}{\footnotesize\color{sourceink}\setstretch{1.12}\RaggedRight\emergencystretch=3em\setlength{\parskip}{0.42em}\obeylines\selectfont}
\newcommand{\transparstyle}{\songfont\color{caseink}\setstretch{1.12}\setlength{\parindent}{0pt}\setlength{\parskip}{0.42em}}
\newcommand{\bilingualpairskip}{%
  \par\addvspace{\biparagraphskip}%
  \switchcolumn
  \par\addvspace{\biparagraphskip}%
  \switchcolumn*%
  \global\intranslationfalse
}
\newcommand{\bilingualboxpairskip}{%
  \par
  \switchcolumn
  \par
  \switchcolumn*%
  \global\intranslationfalse
}
\newcommand{\bilinguallawboxblockskip}{%
  \switchcolumn*%
  \par\addvspace{\lawboxblockafterskip}%
  \switchcolumn
  \par\addvspace{\lawboxblockafterskip}%
  \switchcolumn*%
  \global\intranslationfalse
}
\newcommand{\bikeepwithnext}[1][4]{%
  \ifbilingualcolumns
    \syncleft
    \Needspace{#1\baselineskip}%
    \switchcolumn
    \Needspace{#1\baselineskip}%
    \switchcolumn*%
    \global\intranslationfalse
  \else
    \Needspace{#1\baselineskip}%
  \fi
}
\NewEnviron{sourcepara}[1][]{
  \recordsourcerandstart
  \ifshoworiginal
    \ifbilingualcolumns
      \ifintranslation\switchcolumn*\global\intranslationfalse\fi
    \else
      \Needspace{5\baselineskip}
    \fi
    \ifbilingualcolumns\else\par\addvspace{0.35em}\fi
    {\sourceparstyle
    \noindent\BODY\par}
    \ifbilingualcolumns\else\addvspace{0.25em}\fi
    \ifbilingualcolumns
      \switchcolumn
      \global\intranslationtrue
    \fi
  \else
    \ifrandnumbers
      \begingroup
        \setbox0=\vbox{\hsize=\linewidth\sourceparstyle\noindent\BODY\par}%
      \endgroup
    \fi
  \fi
}
\NewEnviron{transpara}{
  \ifshowtranslation
    \setcounter{randnumberlocal}{\value{lastsourcerandstart}}%
    \ifbilingualcolumns
      \ifintranslation\else\switchcolumn\global\intranslationtrue\fi
      {\transparstyle\setlength{\leftskip}{0.55em}\setlength{\rightskip}{0.15em plus 1em}\addtolength{\linewidth}{-0.55em}\BODY\par}
      \switchcolumn*
      \bilingualpairskip
    \else
      {\transparstyle\BODY\par}
    \fi
  \else
    \ifbilingualcolumns
      \ifintranslation\switchcolumn*\global\intranslationfalse\fi
    \fi
  \fi
}
\NewEnviron{sourceboxpara}[1][]{
  \global\lawboxpartendedblockfalse
  \recordsourcerandstart
  \ifshoworiginal
    \ifbilingualcolumns
      \ifintranslation\switchcolumn*\global\intranslationfalse\fi
    \else
      \Needspace{3\baselineskip}
    \fi
    {\sourceparstyle
    \BODY\par}
    \ifbilingualcolumns\else
      \iflawboxpartendedblock\addvspace{\lawboxblockafterskip}\fi
    \fi
    \ifbilingualcolumns
      \switchcolumn
      \global\intranslationtrue
    \fi
  \fi
}
\NewEnviron{transboxpara}{
  \global\lawboxpartendedblockfalse
  \ifshowtranslation
    \setcounter{randnumberlocal}{\value{lastsourcerandstart}}%
    \ifbilingualcolumns
      \ifintranslation\else\switchcolumn\global\intranslationtrue\fi
      {\transparstyle\BODY\par}
      \iflawboxpartendedblock
        \bilinguallawboxblockskip
      \else
        \switchcolumn*
        \bilingualboxpairskip
      \fi
    \else
      {\transparstyle\BODY\par}
      \iflawboxpartendedblock\addvspace{\lawboxblockafterskip}\fi
    \fi
  \else
    \ifbilingualcolumns
      \ifintranslation\switchcolumn*\global\intranslationfalse\fi
    \fi
  \fi
}
\newcommand{\leitsatznum}[1]{\noindent\hangindent=3.0em\hangafter=1\makebox[2.25em][r]{\bfseries #1.}\hspace{0.55em}\ignorespaces}
\newcommand{\syncleft}{\ifbilingualcolumns\ifintranslation\switchcolumn*\global\intranslationfalse\fi\fi}
\pretocmd{\section}{\syncleft\clearpage}{}{}
\pretocmd{\subsection}{\syncleft}{}{}
\pretocmd{\subsubsection}{\syncleft}{}{}
\newcommand{\biheadingfont}{\songfont}
\newcommand{\bitocentry}[2]{%
  \ifshoworiginal
    \ifshowtranslation
      \ifstrequal{#1}{#2}{#1}{#1\space #2}%
    \else
      #1%
    \fi
  \else
    #2%
  \fi
}
\pdfstringdefDisableCommands{%
  \def\bitocentry#1#2{#1}%
}
\newcommand{\biheading}[7]{%
  \syncleft#1\bikeepwithnext[#3]\phantomsection\addcontentsline{toc}{#2}{\bitocentry{#6}{#7}}%
  \ifbilingualcolumns
    {#4 #6\par}%
    \switchcolumn
    {#4\biheadingfont #7\par}%
    \switchcolumn*%
    \global\intranslationfalse
    \par\addvspace{#5}%
    \switchcolumn
    \par\addvspace{#5}%
    \switchcolumn*%
  \else
    \ifshoworiginal{#4 #6\par}\fi
    \ifshowtranslation{#4\biheadingfont #7\par}\fi
    \addvspace{#5}%
  \fi
}
\newcommand{\termsectionheading}[1]{%
  \par\Needspace{9\baselineskip}%
  \vspace{0.65em}%
  {\songfont\bfseries #1\par}%
  \vspace{0.2em}%
}
% 章節換頁：
%   雙語 paracol 欄內不要用 \clearpage；它會在同步兩欄時吐出只有欄線/頁碼的空頁。
%   \newpage 足以讓下一個 section 起新頁，又不會強制清空所有待排材料。
%   單語模式不在 paracol 內，仍可用 \clearpage。
\newcommand{\bisectionpagebreak}{\ifbilingualcolumns\newpage\else\clearpage\fi}
\newcommand{\bisection}[2]{%
  \biheading{\iffirstbilingualsection\global\firstbilingualsectionfalse\else\bisectionpagebreak\fi}{section}{6}{\centering\Large\bfseries}{0.8em}{#1}{#2}%
}
\newcommand{\bisubsection}[2]{%
  \biheading{}{subsection}{5}{\large\bfseries}{0.55em}{#1}{#2}%
}
\newcommand{\bisubsubsection}[2]{%
  \biheading{}{subsubsection}{4}{\normalsize\bfseries}{0.45em}{#1}{#2}%
}
\newcommand{\bicompositeheading}[4]{%
  \biheading{}{subsection}{5}{\large\bfseries}{0.16em}{#1}{#2}%
  \biheading{}{subsubsection}{3}{\normalsize\bfseries}{0.45em}{#3}{#4}%
}
\newif\iflawboxfirstitem
\lawboxfirstitemfalse
\newcommand{\lawboxprespace}[1]{%
  \par
  \iflawboxfirstitem
    \global\lawboxfirstitemfalse
  \else
    \addvspace{#1}%
  \fi
}
\newtcolorbox{lawboxinner}{
  caseboxbase,
  colframe=sourceline!85!black,
  colback=white,
  left=1.05em,
  right=1.05em,
  top=0.62em,
  bottom=0.62em,
  before upper={\global\inlawboxtrue\global\lawboxfirstitemtrue\global\lawboxhaspendingrnfalse\ifintranslation\lawboxtransfont\else\lawboxsourcefont\fi\RaggedRight\setlength{\parindent}{0pt}\setlength{\parskip}{0pt}\ifintranslation\setstretch{\lawboxtransstretch}\else\setstretch{\lawboxsourcestretch}\fi},
  after upper={\global\lawboxhaspendingrnfalse\global\inlawboxfalse}
}
\tcbset{
  lawboxpartbase/.style={
    enhanced,
    breakable=false,
    width=0.985\linewidth,
    colback=white,
    frame hidden,
    boxrule=0pt,
    arc=0pt,
    left=1.05em,
    right=1.05em,
    boxsep=1.5mm,
    before skip=0pt,
    after skip=0pt,
    before upper={\global\inlawboxtrue\global\lawboxfirstitemtrue\global\lawboxhaspendingrnfalse\ifintranslation\lawboxtransfont\else\lawboxsourcefont\fi\RaggedRight\setlength{\parindent}{0pt}\setlength{\parskip}{0pt}\ifintranslation\setstretch{\lawboxtransstretch}\else\setstretch{\lawboxsourcestretch}\fi},
    after upper={\global\lawboxhaspendingrnfalse\global\inlawboxfalse},
    borderline west={0.28pt}{0pt}{sourceline!85!black},
    borderline east={0.28pt}{0pt}{sourceline!85!black},
  },
  lawboxpartfirst/.style={
    lawboxpartbase,
    top=0.62em,
    bottom=0.04em,
    before skip=\lawboxblockbeforeskip,
    borderline north={0.28pt}{0pt}{sourceline!85!black},
  },
  lawboxpartmiddle/.style={
    lawboxpartbase,
    top=0.04em,
    bottom=0.04em,
  },
  lawboxpartlast/.style={
    lawboxpartbase,
    top=0.04em,
    bottom=0.62em,
    after skip=0pt,
    borderline south={0.28pt}{0pt}{sourceline!85!black},
  }
}
\newtcolorbox{lawboxpartinner}[1]{#1}
\newtcolorbox{termboxinner}{
  caseboxbase,
  colframe=noteblue!70!black,
  colback=casepaper,
  before upper={\songfont\setlength{\parindent}{0pt}\setlength{\parskip}{0.35em}}
}
\newtcolorbox{deboxinner}{
  caseboxbase,
  colframe=notegray!72!black,
  colback=white,
  left=1.05em,
  right=1.05em,
  top=0.62em,
  bottom=0.62em,
  before upper={\global\inlawboxtrue\global\lawboxfirstitemtrue\global\lawboxhaspendingrnfalse\small\setlength{\parindent}{0pt}\setlength{\parskip}{0.35em}},
  after upper={\global\lawboxhaspendingrnfalse\global\inlawboxfalse}
}
\NewEnviron{lawbox}{
  \begin{lawboxinner}\BODY\end{lawboxinner}
}
\NewEnviron{lawboxpart}[2]{
  \begin{lawboxpartinner}{lawboxpart#1,equal height group=#2}\BODY\end{lawboxpartinner}%
  \ifstrequal{#1}{last}{\global\lawboxpartendedblocktrue}{}%
}
\NewEnviron{termbox}{
  \par\addvspace{0.55em}
  {\color{noteblue!70!black}\rule{\linewidth}{0.28pt}}\par
  \begingroup
    \songfont\small\color{caseink}
    \setlength{\parindent}{0pt}
    \setlength{\parskip}{0.24em}
    \BODY
    \par
  \endgroup
  {\color{noteblue!70!black}\rule{\linewidth}{0.28pt}}\par
  \addvspace{0.55em}
}
\NewEnviron{debox}{
  \begin{deboxinner}\BODY\end{deboxinner}
}
\providecommand{\paragraphmark}[1]{}
\newcommand{\lawboxtitleafter}{\addvspace{0.06em}}
\newcommand{\lawtitle}[1]{\lawboxprespace{0.22em}{\lawboxtransfont\bfseries\lawboxuserandnumber #1}\par\lawboxtitleafter}
\newcommand{\absno}[2]{\lawboxprespace{0.04em}\noindent\lawboxuserandnumber\hangindent=2.6em\hangafter=1\makebox[2.6em][l]{(#1)}#2\par}
\newcommand{\nrno}[2]{\lawboxprespace{0pt}\noindent\lawboxuserandnumber\hspace*{2.6em}\hangindent=4.4em\hangafter=1\makebox[1.8em][l]{#1.}#2\par}
\newcommand{\letterno}[2]{\lawboxprespace{0pt}\noindent\lawboxuserandnumber\hspace*{4.4em}\hangindent=6.0em\hangafter=1\makebox[1.6em][l]{#1)}#2\par}
\newcommand{\sourcelawtitle}[1]{\lawboxprespace{0.22em}{\lawboxsourcefont\bfseries\lawboxuserandnumber #1}\par\lawboxtitleafter}
\newcommand{\sourcelawsubtitle}[1]{\lawboxprespace{0.04em}{\lawboxsourcefont\bfseries\lawboxuserandnumber #1}\par\lawboxtitleafter}
\newcommand{\sourceabsno}[2]{\lawboxprespace{0.04em}\noindent\lawboxuserandnumber\hangindent=2.45em\hangafter=1\makebox[2.45em][l]{(#1)}#2\par}
\newcommand{\sourcenrno}[2]{\lawboxprespace{0pt}\noindent\lawboxuserandnumber\hspace*{2.45em}\hangindent=4.25em\hangafter=1\makebox[1.8em][l]{#1.}#2\par}
\newcommand{\sourceletterno}[2]{\lawboxprespace{0pt}\noindent\lawboxuserandnumber\hspace*{4.25em}\hangindent=5.85em\hangafter=1\makebox[1.6em][l]{#1)}#2\par}
\newcommand{\sourcequotepara}[1]{\lawboxprespace{0.04em}\noindent\lawboxuserandnumber\hangindent=1.2em\hangafter=1 #1\par}
\newcommand{\sourcelawline}[1]{\lawboxprespace{0.04em}\noindent\lawboxuserandnumber #1\par}
\newcommand{\lawline}[1]{\lawboxprespace{0.04em}\noindent\lawboxuserandnumber #1\par}
\newcommand{\pendingtranslation}{\par{\songfont\footnotesize\color{pendingink}（中文譯文待補。）}\par}
\newcommand{\tocpart}[1]{\par\addvspace{0.52em}{\footnotesize\bfseries\color{pendingink}#1\par}\addvspace{0.16em}}
\newcommand{\tocdashline}{\par\addvspace{0.48em}\noindent{\color{notegray}\xleaders\hbox to 0.82em{\hfil\rule{0.42em}{0.28pt}\hfil}\hskip\linewidth}\par\addvspace{0.38em}}
% \tocpanel{font-cmd}{content}：將目錄置中於所在欄寬（雙語欄適用）。
% A3 橫式使用 0.58\linewidth，A4 橫式使用 0.84\linewidth，
% 使兩種紙張下目錄欄內容實際寬度相近（均約 100–105 mm）。
\newcommand{\tocpanel}[2]{%
  \makebox[\linewidth][c]{%
    \ifx\bilingualpaper\bilingualpaperaThree
      \begin{minipage}{0.58\linewidth}%
    \else
      \begin{minipage}{0.84\linewidth}%
    \fi
    \toccolstyle
    #1#2%
    \end{minipage}%
  }%
}
% ===== 首頁簡目錄的兩層 description list 縮排 =====
% enumitem 的 description list 可把每一列拆成：
%   [label]  labelsep  body text
% 這裡的「起點」都以所在 minipage/欄位的左緣為 0。
%
% 核心規則：
%   labelindent = label 欄位從左緣開始的位置。
%   labelwidth  = label 欄位本身寬度；標籤太長（如 Abw. Meinung）就加大它。
%   labelsep    = label 欄位右緣到正文的距離。
%   leftmargin  = 正文 body text 的起點。判斷層級視覺時主要看這個。
%
% 所以：
%   想讓「Begründung / 判決理由」整列正文往右或往左：改 tocitems 的 leftmargin。
%   想讓「A. Verfahren und Gegenstand / A. 程序與審查標的」正文往右或往左：
%     改 tocsubitems 的 leftmargin。
%   想讓 A./B./C. 這些標籤本身往右或往左，但正文不動：改 tocsubitems 的 labelindent。
%   想讓 A./B./C. 與正文間距變大或變小：改 tocsubitems 的 labelsep。
%
% 層級原則：
%   tocsubitems.leftmargin 必須大於 tocitems.leftmargin，
%   否則子層級正文會看起來比「Gründe / 理由」還靠左。
\newlist{tocitems}{description}{1}
\setlist[tocitems]{%
  labelindent=0pt,    % 第一層 label 欄位起點；0pt 表示貼齊目錄欄左緣。
  leftmargin=2.54cm,  % 第一層正文起點；例如 Begründung / 判決理由 的左緣。
  labelwidth=2.16cm,  % 第一層 label 欄位寬度；需容納 Leitsätze、Abw. Meinung、不同意見等。
  labelsep=0.32cm,    % 第一層 label 與正文的水平間距；只改間距，不改正文起點。
  itemsep=0.28em,     % 第一層各項之間的垂直距離。
  topsep=0.20em,      % 第一層 list 上下與前後內容的垂直距離。
  parsep=0pt,         % 同一 item 內段落之間的距離；目錄項通常不需要。
  partopsep=0pt,      % list 若緊接段落開始時的額外距離；關閉以免目錄高度飄動。
  align=left,         % label 欄內左對齊；長短標籤不靠右貼齊。
  font=\normalfont\bfseries % label 的字型；正文不受這行影響。
}
\newlist{tocsubitems}{description}{1}
\setlist[tocsubitems]{%
  labelindent=1cm,    % 第二層 label 起點；只控制 A./B./C. 本身的位置。
  leftmargin=3.00cm,  % 第二層正文起點；必須大於 2.54cm，才會比 Begründung / 判決理由更右。
  labelwidth=0.5cm,   % 第二層 label 欄寬；A./B./C. 很短，可小於第一層。
  labelsep=1.5cm,    % 第二層 label 與正文間距；A. 與 Verfahren / 程序 的距離看這裡。
  itemsep=0.12em,     % 第二層各項較密，避免 A.–E. 佔太高。
  topsep=0.04em,      % 第二層 list 與 Gründe / 理由之間的垂直距離。
  parsep=0pt,         % 同一第二層 item 內段落距離；目錄項通常不需要。
  partopsep=0pt,      % 關閉額外段落起始距離，避免版面不穩。
  align=left,         % A./B./C. label 欄內左對齊。
  font=\normalfont\bfseries, % A./B./C. label 加粗；正文不受這行影響。
  before={}           % 不在第二層 list 前自動插入額外內容。
}
\newlist{termitems}{description}{1}
\ifbilingualcolumns
  \setlist[termitems]{%
    leftmargin=5.2cm,
    labelwidth=4.8cm,
    labelsep=0.35cm,
    itemsep=0.16em,
    topsep=0.2em,
    parsep=0pt,
    partopsep=0pt,
    font=\normalfont\bfseries,
    style=nextline
  }
\else
  \setlist[termitems]{%
    leftmargin=0pt,
    labelwidth=0pt,
    labelsep=0pt,
    itemsep=0.24em,
    topsep=0.2em,
    parsep=0pt,
    partopsep=0pt,
    font=\normalfont\bfseries,
    style=nextline
  }
\fi
\newlist{abbritems}{description}{1}
\setlist[abbritems]{%
  leftmargin=2.38cm,
  labelwidth=2.02cm,
  labelsep=0.28cm,
  itemsep=0.16em,
  topsep=0.2em,
  parsep=0pt,
  partopsep=0pt,
  align=left,
  font=\normalfont\bfseries
}
\ifbilingualcolumns
  \newenvironment{termcolumns}{\begin{multicols}{2}}{\end{multicols}}
\else
  \newenvironment{termcolumns}{}{}
\fi

% ===== 編者註腳開關 =====
% 一般版預設不顯示編者註，避免正文腳註過密。
% 若開啟 \classroommodeon，GG/條文編者註仍會自動顯示，無需另開本開關。
% 若一般版也要顯示編者註，可在單案檔 \input 後加 \showeditornotestrue。
\newif\ifshoweditornotes
\showeditornotesfalse

% ===== 目錄排版輔助 =====
\newcommand{\toccolstyle}{\raggedright\arraybackslash}
\newcommand{\tocsingleindent}{\leftskip=1.45cm\rightskip=1.2cm}

% ===== 行內引文環境（德文原文與中文譯文各一，帶左側灰色邊線）=====
\newcommand{\quoteparbreak}{\par\addvspace{0.48em}}
\newcommand{\sourcequoteinner}[1]{%
  \begin{tcolorbox}[
    enhanced, breakable, width=\dimexpr\linewidth-0.8em\relax, frame hidden,
    colback=white, coltext=caseink, boxrule=0pt,
    borderline west={0.55pt}{0.88em}{notegray},
    left=1.78em, right=0.72em, top=0.18em, bottom=0.18em,
    boxsep=0pt, before skip=0.24em, after skip=0.24em,
    before upper={\normalfont\footnotesize\color{caseink}\setlength{\parindent}{0pt}\setlength{\parskip}{0.34em}\raggedright\setlength{\rightskip}{0pt plus 1em}}
  ]%
  #1\par
  \end{tcolorbox}%
}
\newcommand{\transquoteinner}[1]{%
  \begin{tcolorbox}[
    enhanced, breakable, width=\dimexpr\linewidth-0.8em\relax, frame hidden,
    colback=white, coltext=caseink, boxrule=0pt,
    borderline west={0.55pt}{0.88em}{notegray},
    left=1.78em, right=0.72em, top=0.18em, bottom=0.18em,
    boxsep=0pt, before skip=0.24em, after skip=0.24em,
    before upper={\songfont\small\color{caseink}\setlength{\parindent}{0pt}\setlength{\parskip}{0.34em}\raggedright\setlength{\rightskip}{0pt plus 1em}}
  ]%
  #1\par
  \end{tcolorbox}%
}

% ===== 大綱（Gliederung）Rn 輔助命令 =====
\newcommand{\outrn}[2]{\enspace{\normalfont\footnotesize\color{pendingink}(Rn\,#1–#2)}}
\newcommand{\outrnsingle}[1]{\enspace{\normalfont\footnotesize\color{pendingink}(Rn\,#1)}}
\newcommand{\outrnzh}[2]{\enspace{\normalfont\footnotesize\color{pendingink}（第\,#1–#2\,段）}}
\newcommand{\outrnzhs}[1]{\enspace{\normalfont\footnotesize\color{pendingink}（第\,#1\,段）}}
% ===== 大綱雙語對照表格輔助命令（三欄：段碼 │ 德文 │ 中文）=====
% 欄 1：段碼（由欄定義自動格式化：小字、等寬、右對齊、pendingink）
% 欄 2：德文標籤＋正文
% 欄 3：中文標籤＋正文
%
% 無段碼的列（\omainrow、\osecrow），欄 1 以 & 空置。
% 有段碼的列（\osecrowrn、\osubrow、\ointrow），#1 為預格式化字串
%   例：\osubrow{2–49}{I.}{...}{I.}{...}
%       \ointrow{107}{...}{...}
%
% \opartrow：段組標題（橫跨三欄）
\newcommand{\opartrow}[2]{%
  \noalign{\vspace{0.30em}}%
  \multicolumn{3}{@{}l@{}}{%
    \footnotesize\bfseries\color{pendingink}#1\hfill\songfont#2%
  }\\[0.06em]%
}
% \odashrow：虛線分隔（橫跨三欄）
\newcommand{\odashrow}{%
  \noalign{\vspace{0.28em}}%
  \multicolumn{3}{@{}l@{}}{\color{notegray}%
    \xleaders\hbox to 0.82em{\hfil\rule{0.42em}{0.28pt}\hfil}\hskip 0.88\linewidth}%
  \\[0.24em]%
}
% \odissentpart：不同意見書區段標頭。比一般 \opartrow 更像附錄性轉場。
\newcommand{\odissentpart}[2]{%
  \noalign{\vspace{0.42em}}%
  \multicolumn{3}{@{}p{0.88\textwidth}@{}}{%
    \color{notegray}\rule{\linewidth}{0.18pt}%
  }\\[-0.18em]%
  \multicolumn{3}{@{}p{0.88\textwidth}@{}}{%
    \makebox[\linewidth][s]{%
      \footnotesize\bfseries\color{pendingink}#1%
      \hfill{\songfont #2}%
    }%
  }\\[-0.02em]%
}
% \odissentopinion：不同意見書的署名與一行摘要，右欄保留段碼範圍。
\newcommand{\odissentopinion}[5]{%
  \footnotesize\hangindent=0.52cm\hangafter=1%
    {\bfseries\color{caseink}#2}\enspace{\color{notegray}\textperiodcentered}\enspace\color{caseink}#3%
  &\footnotesize\hangindent=0.52cm\hangafter=1%
    {\songfont\bfseries\color{caseink}#4}\enspace{\color{notegray}\textperiodcentered}\enspace\songfont\color{caseink}#5%
  &#1%
  \\[0.12em]%
}
% \omainrow：頂層項目，無段碼（Leitsätze、Urteil、Gründe …）
\newcommand{\omainrow}[4]{%
  \footnotesize\hangindent=1.92cm\hangafter=1%
    \makebox[1.92cm][l]{\bfseries\color{caseink}#1}\color{caseink}#2%
  &\footnotesize\hangindent=1.92cm\hangafter=1%
    \makebox[1.92cm][l]{\bfseries\color{caseink}#3}\songfont\color{caseink}#4%
  &%
  \\[0.15em]%
}
% \omainrowrn：頂層項目，有段碼（不同意見等標籤寬於 \osecrow 框者）
\newcommand{\omainrowrn}[5]{%
  \footnotesize\hangindent=1.92cm\hangafter=1%
    \makebox[1.92cm][l]{\bfseries\color{caseink}#2}\color{caseink}#3%
  &\footnotesize\hangindent=1.92cm\hangafter=1%
    \makebox[1.92cm][l]{\bfseries\color{caseink}#4}\songfont\color{caseink}#5%
  &#1%
  \\[0.15em]%
}
% \osecrow：一級子項，無段碼（A.、C.、D. 等有子項者）
\newcommand{\osecrow}[4]{%
  \footnotesize\hspace{1.10cm}\hangindent=1.98cm\hangafter=1%
    \makebox[0.86cm][l]{\bfseries\color{caseink}#1}\color{caseink}#2%
  &\footnotesize\hspace{1.10cm}\hangindent=1.98cm\hangafter=1%
    \makebox[0.86cm][l]{\bfseries\color{caseink}#3}\songfont\color{caseink}#4%
  &%
  \\[0.11em]%
}
% \osecrowrn：一級子項，有段碼（B.、E. 等完整段落範圍者）
% #1=段碼字串（如 101–106）, #2=DE標籤, #3=DE正文, #4=ZH標籤, #5=ZH正文
\newcommand{\osecrowrn}[5]{%
  \footnotesize\hspace{1.10cm}\hangindent=1.98cm\hangafter=1%
    \makebox[0.86cm][l]{\bfseries\color{caseink}#2}\color{caseink}#3%
  &\footnotesize\hspace{1.10cm}\hangindent=1.98cm\hangafter=1%
    \makebox[0.86cm][l]{\bfseries\color{caseink}#4}\songfont\color{caseink}#5%
  &#1%
  \\[0.11em]%
}
% \osubrow：二級子項，有段碼範圍（I.、II.、A.I. …）
% #1=段碼字串, #2=DE標籤, #3=DE正文, #4=ZH標籤, #5=ZH正文
\newcommand{\osubrow}[5]{%
  \footnotesize\hspace{1.40cm}\hangindent=2.30cm\hangafter=1%
    \makebox[0.90cm][l]{\bfseries\color{caseink}#2}\color{caseink}#3%
  &\footnotesize\hspace{1.40cm}\hangindent=2.30cm\hangafter=1%
    \makebox[0.90cm][l]{\bfseries\color{caseink}#4}\songfont\color{caseink}#5%
  &#1%
  \\[0.09em]%
}
% \ointrow：引入段落，有單一段碼，無標籤
% #1=段碼字串, #2=DE正文, #3=ZH正文
\newcommand{\ointrow}[3]{%
  \footnotesize\hspace{0.52cm}\hangindent=0.52cm\hangafter=1\color{caseink}#2%
  &\footnotesize\hspace{0.52cm}\hangindent=0.52cm\hangafter=1\songfont\color{caseink}#3%
  &#1%
  \\[0.09em]%
}
% \osubsubrow：三級子項（1.、2.、3.）
% #1=段碼字串, #2=DE標籤, #3=DE正文, #4=ZH標籤, #5=ZH正文
\newcommand{\osubsubrow}[5]{%
  \footnotesize\hspace{1.80cm}\hangindent=2.48cm\hangafter=1%
    \makebox[0.68cm][l]{\bfseries\color{caseink}#2}\color{caseink}#3%
  &\footnotesize\hspace{1.80cm}\hangindent=2.48cm\hangafter=1%
    \makebox[0.68cm][l]{\bfseries\color{caseink}#4}\songfont\color{caseink}#5%
  &#1%
  \\[0.08em]%
}
% \osubsubsubrow：四級子項（a）、b））
% #1=段碼字串, #2=DE標籤, #3=DE正文, #4=ZH標籤, #5=ZH正文
\newcommand{\osubsubsubrow}[5]{%
  \footnotesize\hspace{2.10cm}\hangindent=2.68cm\hangafter=1%
    \makebox[0.58cm][l]{\bfseries\color{caseink}#2}\color{caseink}#3%
  &\footnotesize\hspace{2.10cm}\hangindent=2.68cm\hangafter=1%
    \makebox[0.58cm][l]{\bfseries\color{caseink}#4}\songfont\color{caseink}#5%
  &#1%
  \\[0.07em]%
}
% \osubsubsubsubrow：五級子項（aa）、bb））
% 標籤框 0.62cm：足以容納「aa)」在 footnotesize bold 下（約 0.50cm）並留有間距
% #1=段碼字串, #2=DE標籤, #3=DE正文, #4=ZH標籤, #5=ZH正文
\newcommand{\osubsubsubsubrow}[5]{%
  \footnotesize\hspace{2.38cm}\hangindent=3.06cm\hangafter=1%
    \makebox[0.68cm][l]{\bfseries\color{caseink}#2}\color{caseink}#3%
  &\footnotesize\hspace{2.38cm}\hangindent=3.06cm\hangafter=1%
    \makebox[0.68cm][l]{\bfseries\color{caseink}#4}\songfont\color{caseink}#5%
  &#1%
  \\[0.06em]%
}
% \outlinepage：雙語對照大綱頁包裝環境（longtable 自動跨頁）
% 三欄：德文欄 + 中文欄 + 段碼欄（最右，不折行）
% 段碼欄寬度：1.80cm，足以容納最長段碼如「309–348」
% DE/ZH 欄寬：(0.87\textwidth - 2.08cm) / 2
%   驗算：2×欄 + 0.05\textwidth + 0.28cm gap + 1.80cm = 0.87tw - 2.08cm + 0.05tw + 2.08cm = 0.92tw ✓
\newenvironment{outlinepage}{%
  \vspace*{0.40cm}%
  \setlength{\LTleft}{0.04\textwidth}%
  \setlength{\LTright}{0.04\textwidth}%
  \setlength{\tabcolsep}{0pt}%
  % 大綱列距由 longtable 的 \arraystretch 與各 row macro 結尾的 \\[...em] 共同決定。
  % 這裡控制「每一列本身的表格 strut 高度」；數值越大，A./I./1. 各項之間越像有空行。
  % A3 原本用 0.62，視覺上沒有空行；A4 也沿用同值，避免切到 A4 後項目被撐開。
  % 若日後覺得大綱太擠，只微調這個值（例如 0.68），不要改各 \omainrow/\osubrow 的縮排。
  \renewcommand{\arraystretch}{0.62}%
  \begin{longtable}{%
    >{\vspace{0pt}\raggedright\arraybackslash}p{\dimexpr(0.87\textwidth-2.08cm)/2\relax}%
    @{\hspace{0.025\textwidth}}%
    !{\color{notegray}\vrule width 0.12pt}%
    @{\hspace{0.025\textwidth}}%
    >{\vspace{0pt}\raggedright\arraybackslash}p{\dimexpr(0.87\textwidth-2.08cm)/2\relax}%
    @{\hspace{0.28cm}}%
    >{\vspace{0pt}\footnotesize\normalfont\color{pendingink}\raggedright\arraybackslash}p{1.80cm}%
    @{}%
  }%
  % 首頁欄標籤：不要橫跨置中；分別放在德文欄與中文欄的實際欄位起點。
  \footnotesize\bfseries\color{sourceink}Gliederung%
  &\footnotesize\bfseries\songfont\color{sourceink}大綱%
  &%
  \\[0.36em]%
  \endfirsthead%
  % 續頁欄標籤：同樣分欄定位，以灰色細體區分；無需標注「續」。
  \footnotesize\color{pendingink}Gliederung%
  &\footnotesize\songfont\color{pendingink}大綱%
  &%
  \\[0.24em]%
  \endhead%
}{%
  \end{longtable}%
}
 載入。
% 不含：\documentclass、\documentversion、\ggversionde/zh、\tocde/\toczh。
% 日期與首頁版本列由本檔自動產生；各文件只需手動指定 \documentversion。

\usepackage{xeCJK}
\usepackage{fontspec}
\usepackage{geometry}
\usepackage{setspace}
\usepackage{hyperref}
\usepackage{xurl}
\usepackage{bookmark}
\usepackage{enumitem}
\usepackage{xcolor}
\usepackage{titlesec}
\usepackage{array}
\usepackage{longtable}
\usepackage{tcolorbox}
\usepackage{environ}
\usepackage{pdflscape}
\usepackage{etoolbox}
\usepackage{ragged2e}
\usepackage{paracol}
\usepackage{multicol}
\usepackage{needspace}
\usepackage{fancyhdr}
\usepackage{tikz}
\usepackage{zref-abspage}
\tcbuselibrary{breakable,skins}
\usetikzlibrary{calc}
\usepackage{pgfcalendar}

\hypersetup{
  unicode=true,
  bookmarksopen=true,
  bookmarksnumbered=false,
  hidelinks
}

% ===== 自動推導，不要手動改下面 =====
% （\docmode 與 \bilingualpaper 由各文件在 % bverfge_layout.tex
% 共用排版基礎設施，供 Schwangerschaftsabbruch I 與 II 共享。
% 由各文件以 \documentclass 後 \input{../../共用LaTeX/bverfge_layout} 載入。
% 不含：\documentclass、\documentversion、\ggversionde/zh、\tocde/\toczh。
% 日期與首頁版本列由本檔自動產生；各文件只需手動指定 \documentversion。

\usepackage{xeCJK}
\usepackage{fontspec}
\usepackage{geometry}
\usepackage{setspace}
\usepackage{hyperref}
\usepackage{xurl}
\usepackage{bookmark}
\usepackage{enumitem}
\usepackage{xcolor}
\usepackage{titlesec}
\usepackage{array}
\usepackage{longtable}
\usepackage{tcolorbox}
\usepackage{environ}
\usepackage{pdflscape}
\usepackage{etoolbox}
\usepackage{ragged2e}
\usepackage{paracol}
\usepackage{multicol}
\usepackage{needspace}
\usepackage{fancyhdr}
\usepackage{tikz}
\usepackage{zref-abspage}
\tcbuselibrary{breakable,skins}
\usetikzlibrary{calc}
\usepackage{pgfcalendar}

\hypersetup{
  unicode=true,
  bookmarksopen=true,
  bookmarksnumbered=false,
  hidelinks
}

% ===== 自動推導，不要手動改下面 =====
% （\docmode 與 \bilingualpaper 由各文件在 \input{bverfge_layout} 之前設定；
%   預設 chinese / a4）
\ifdefined\docmode\else\def\docmode{chinese}\fi
\ifdefined\bilingualpaper\else\def\bilingualpaper{a4}\fi
\newif\ifshoworiginal
\newif\ifshowtranslation
\newif\ifbilinguallandscape
\newif\ifbilingualcolumns
\newif\ifintranslation
\newif\iffirstbilingualsection

\showoriginalfalse
\showtranslationfalse
\bilinguallandscapefalse
\bilingualcolumnsfalse
\intranslationfalse
\firstbilingualsectionfalse

\def\modechinese{chinese}
\def\modegerman{german}
\def\modebilingual{bilingual}
\def\bilingualpaperaThree{a3}
\def\bilingualpaperaFour{a4}

\ifx\docmode\modechinese
  \showtranslationtrue
\fi

\ifx\docmode\modegerman
  \showoriginaltrue
\fi

\ifx\docmode\modebilingual
  \showoriginaltrue
  \showtranslationtrue
  \bilinguallandscapetrue
  \bilingualcolumnstrue
\fi

% 編排模式說明：
% chinese  → \showoriginalfalse  \showtranslationtrue  （A4 直式，純中文）
% german   → \showoriginaltrue   \showtranslationfalse （A4 直式，純德文）
% bilingual→ \showoriginaltrue   \showtranslationtrue  \bilingualcolumnstrue（A3/A4 橫式對照）

\ifbilingualcolumns
  \ifx\bilingualpaper\bilingualpaperaFour
    \geometry{
      a4paper,
      landscape,
      left=1.25cm,
      right=2.25cm,
      top=1.25cm,
      bottom=1.75cm,
      footskip=8.5mm,
      marginparwidth=0.75cm,
      marginparsep=0.25cm
    }
  \else
    \geometry{
      a3paper,
      landscape,
      left=1.55cm,
      right=2.75cm,
      top=1.55cm,
      bottom=2.05cm,
      footskip=9.5mm,
      marginparwidth=0.9cm,
      marginparsep=0.35cm
    }
  \fi
\else
\geometry{
  a4paper,
  portrait,
  left=2.2cm,
  right=2.6cm,
  top=2.0cm,
  bottom=2.0cm,
  marginparwidth=0.8cm,
  marginparsep=0.3cm
}
\fi
% ===== 時間戳基礎設施 =====
\newcount\stamptimeval
\newcount\stampminuteval
\newcommand{\stamphh}{%
  \stamptimeval=\time
  \divide\stamptimeval by 60
  \ifnum\stamptimeval<10 0\fi
  \the\stamptimeval}
\newcommand{\stampmm}{%
  \stamptimeval=\time
  \stampminuteval=\time
  \divide\stampminuteval by 60
  \multiply\stampminuteval by 60
  \advance\stamptimeval by -\stampminuteval
  \ifnum\stamptimeval<10 0\fi
  \the\stamptimeval}
\newcommand{\stampwkde}[1]{%
  \ifcase#1Montag\or Dienstag\or Mittwoch\or Donnerstag\or Freitag\or Samstag\or Sonntag\fi}
\newcommand{\stampwkzh}[1]{%
  \ifcase#1 一\or 二\or 三\or 四\or 五\or 六\or 日\fi}
\newcommand{\stampmonthde}[1]{%
  \ifcase#1\or Januar\or Februar\or März\or April\or Mai\or Juni\or
  Juli\or August\or September\or Oktober\or November\or Dezember\fi}
\newcount\stampjdval
\newcount\stampwdval
\pgfcalendardatetojulian{\the\year-\the\month-\the\day}{\stampjdval}
\pgfcalendarjuliantoweekday{\the\stampjdval}{\stampwdval}
\newcommand{\timestampzh}{%
  \songfont 最後修改：\the\year 年\the\month 月\the\day 日%
  % （星期\stampwkzh{\the\stampwdval}）%
  % \space\stamphh:\stampmm
  }

\newcommand{\documentdatezh}{\the\year 年 \the\month 月 \the\day 日}
\newcommand{\documentdatede}{\the\day. \stampmonthde{\the\month} \the\year}
\newcommand{\bverfgetranslationnote}{%
  {\songfont 翻譯說明：初譯使用 Claude Code 與 GPT 協助迭代；仍請以德文原文為準。}%
}
\newcommand{\bverfgecourseinfo}{%
  {\songfont 課程：114 學年度第 2 學期；憲法專題研究（四）；授課老師：湯德宗教授；導讀日期：115 年 6 月 1 日。}%
}
\newcommand{\bverfgeeditorinfo}{%
  {\songfont 編者：王逸帆（R10A21126），國立台灣大學法律系研究所碩士班。}%
}
\newcommand{\bverfgetitleversionline}{%
  \ifbilingualcolumns
    Fassung \documentversion\enspace\textperiodcentered\enspace Stand: \documentdatede
    \quad
    {\songfont 版本 \documentversion\enspace\textperiodcentered\enspace 修訂日期：\documentdatezh\par}%
  \else
    \ifshowtranslation
      {\songfont 版本 \documentversion\enspace\textperiodcentered\enspace 修訂日期：\documentdatezh\par}%
    \else
      Fassung \documentversion\enspace\textperiodcentered\enspace Stand: \documentdatede\par
    \fi
  \fi
}
\newcommand{\bverfgetitlecornerinfo}{%
  \ifshowtranslation
    \vspace{0.72em}%
    \noindent\hfill
    \begin{minipage}{0.66\textwidth}
      \raggedleft\tiny\color{pendingink}\songfont
      \bverfgecourseinfo\par
      \bverfgeeditorinfo
    \end{minipage}%
  \fi
}
\newcommand{\bverfgetitlemeta}{%
  \vfill
  \begin{center}%
    \footnotesize\color{pendingink}%
    \bverfgetitleversionline
    \ifshowtranslation
      \vspace{0.28em}{\scriptsize\bverfgetranslationnote\par}%
    \fi
  \end{center}%
  \bverfgetitlecornerinfo
}

\pagestyle{fancy}
\fancyhf{}
\renewcommand{\headrulewidth}{0pt}
\renewcommand{\footrulewidth}{0pt}
\fancyfoot[C]{\thepage}
\ifbilingualcolumns
  \fancyfoot[R]{\scriptsize\color{pendingink}\timestampzh}%
\else
  \ifshoworiginal
  \else
    \fancyfoot[R]{\scriptsize\color{pendingink}\timestampzh\hspace{0.45cm}}%
  \fi
\fi
\setmainfont{Times New Roman}
\setsansfont{Helvetica Neue}
\setmonofont{Menlo}
\IfFileExists{/Users/iw/Library/Fonts/kaiu.ttf}{
  \setCJKmainfont[Path=/Users/iw/Library/Fonts/,AutoFakeBold=4]{kaiu.ttf}
}{
  \IfFontExistsTF{標楷體}{
    \setCJKmainfont[AutoFakeBold=4]{標楷體}
  }{
    \setCJKmainfont{Songti TC}
  }
}
\IfFontExistsTF{Songti TC}{
  \setCJKfamilyfont{song}{Songti TC}
  \setCJKmonofont{Songti TC}
}{
  \setCJKfamilyfont{song}[Path=/Users/iw/Library/Fonts/,AutoFakeBold=4]{kaiu.ttf}
  \setCJKmonofont[Path=/Users/iw/Library/Fonts/,AutoFakeBold=4]{kaiu.ttf}
}
\newcommand{\songfont}{\CJKfamily{song}}

% ===== 封面／目錄專用打字機字體（僅限封面、目錄頁使用，不影響正文）=====
\IfFontExistsTF{Radio Newsman}{%
  \newfontfamily\bverfgetitlede{Radio Newsman}%
}{%
  \newfontfamily\bverfgetitlede{Courier New}%
}
\IfFontExistsTF{Erikas Farbband}{%
  \newfontfamily\bverfgebodyde{Erikas Farbband}%
}{%
  \let\bverfgebodyde\bverfgetitlede
}
\IfFontExistsTF{Maoken Heavy Labourer}{%
  \setCJKfamilyfont{bverfgetitlecjk}{Maoken Heavy Labourer}%
}{%
  \setCJKfamilyfont{bverfgetitlecjk}{Songti TC}%
}
\IfFontExistsTF{Huiwen-mincho}{%
  \setCJKfamilyfont{bverfgebodycjk}{Huiwen-mincho}%
}{%
  \setCJKfamilyfont{bverfgebodycjk}{Songti TC}%
}
\newcommand{\bverfgetitlezh}{\bverfgetitlede\CJKfamily{bverfgetitlecjk}}
\newcommand{\bverfgebodyzh}{\bverfgebodyde\CJKfamily{bverfgebodycjk}}

\setstretch{1.35}
\setlist{itemsep=0.2em, topsep=0.4em}
\setlength{\parindent}{0pt}
\setlength{\parskip}{0.45em}
\raggedbottom
\setcounter{secnumdepth}{0}
\setcounter{tocdepth}{2}

\newcommand{\lawboxsourcefont}{\footnotesize}
\newcommand{\lawboxtransfont}{\footnotesize}
\newcommand{\lawboxsourcestretch}{1.02}
\newcommand{\lawboxtransstretch}{0.92}

\titleformat{\section}{\centering\Large\bfseries\songfont}{\thesection}{1em}{}
\titleformat{\subsection}{\large\bfseries\songfont}{\thesubsection}{1em}{}
\titleformat{\subsubsection}{\normalsize\bfseries\songfont}{\thesubsubsection}{1em}{}
\titleformat{\paragraph}{\normalsize\bfseries\songfont}{\theparagraph}{1em}{}
\titlespacing*{\section}{0pt}{1.8em}{1.0em}
\titlespacing*{\subsection}{0pt}{1.2em}{0.6em}
\titlespacing*{\subsubsection}{0pt}{1.0em}{0.45em}
\titlespacing*{\paragraph}{0pt}{0.6em}{0.4em}

\definecolor{noteblue}{HTML}{A9C9F5}
\definecolor{noteteal}{HTML}{AEE1D6}
\definecolor{notegray}{HTML}{D7DADF}
\definecolor{caseink}{HTML}{000000}
\definecolor{casepaper}{HTML}{FCFEFD}
\definecolor{sourcepaper}{HTML}{FAFBF8}
\definecolor{sourceline}{HTML}{D7DED8}
\definecolor{sourceink}{HTML}{000000}
\definecolor{pendingink}{HTML}{000000}

\tcbset{
  caseboxbase/.style={
    enhanced,
    breakable,
    width=0.985\linewidth,
    colback=casepaper,
    coltitle=caseink,
    fonttitle=\bfseries\songfont,
    boxrule=0.28pt,
    arc=1.2mm,
    left=2.4mm,
    right=2.4mm,
    top=1.5mm,
    bottom=1.5mm,
    boxsep=1.5mm,
    before skip=0.65em,
    after skip=0.65em,
    before upper={\setlength{\parindent}{0pt}\setlength{\parskip}{0.35em}},
  }
}

\newcommand{\bilingualstart}{}
\newcommand{\bilingualstop}{}
\ifbilingualcolumns
  \renewcommand{\bilingualstart}{\small\setlength{\columnsep}{1.25cm}\setlength{\columnseprule}{0.12pt}\columnratio{0.5,0.5}\multicolumnfootnotes\flushbottom\sloppy\interlinepenalty=80\clubpenalty=800\widowpenalty=800\displaywidowpenalty=800\begin{paracol}{2}\global\intranslationfalse\global\firstbilingualsectiontrue{\footnotesize\color{sourceink}Deutsch\par}\switchcolumn{\footnotesize\songfont\color{sourceink}中文譯文\par}\switchcolumn*}
  \renewcommand{\bilingualstop}{\ifintranslation\switchcolumn*\global\intranslationfalse\fi\end{paracol}\clearpage\raggedbottom}
\fi
\newif\ifrandnumbers
\randnumbersfalse
\newcounter{randnumber}
\newcounter{randnumberlocal}
\newcounter{lastsourcerandstart}
\newif\iflastsourcehasrandnumbers
\lastsourcehasrandnumbersfalse
\newif\ifinlawbox
\inlawboxfalse
\newif\iflawboxhaspendingrn
\lawboxhaspendingrnfalse
\newcommand{\lawboxpendingrn}{}
\newcommand{\startrandnumbers}{\global\randnumberstrue\setcounter{randnumber}{1}\global\lastsourcehasrandnumbersfalse}
\newcommand{\stoprandnumbers}{\global\randnumbersfalse\global\lastsourcehasrandnumbersfalse}
\newlength{\randnumberwidth}
\setlength{\randnumberwidth}{0.95cm}
\newlength{\randnumberrightoffset}
\setlength{\randnumberrightoffset}{0.85cm}
\newlength{\randnumberlawboxrightoffset}
\setlength{\randnumberlawboxrightoffset}{1.40cm}
\newcommand{\randnumberbox}[1]{%
  \leavevmode
  \makebox[0pt][l]{%
    \smash{%
      \tikz[remember picture,overlay,baseline=(rnmark.base)]{%
        \coordinate (rnmark) at (0,0);%
        \ifinlawbox
          \path let \p1=(current page.east), \p2=(rnmark.base) in
            node[
              anchor=base east,
              inner sep=0pt,
              outer sep=0pt,
              text width=\randnumberwidth,
              align=right,
              text=pendingink,
              font=\normalfont\mdseries\scriptsize\ttfamily
            ] at (\x1-\randnumberlawboxrightoffset,\y2) {{\normalfont\mdseries\scriptsize\ttfamily #1}};%
        \else
          \path let \p1=(current page.east), \p2=(rnmark.base) in
            node[
              anchor=base east,
              inner sep=0pt,
              outer sep=0pt,
              text width=\randnumberwidth,
              align=right,
              text=pendingink,
              font=\normalfont\mdseries\scriptsize\ttfamily
            ] at (\x1-\randnumberrightoffset,\y2) {{\normalfont\mdseries\scriptsize\ttfamily #1}};%
        \fi
      }%
    }%
  }%
  \ignorespaces
}
\newcommand{\lawboxstorerandnumber}[1]{%
  \gdef\lawboxpendingrn{#1}%
  \global\lawboxhaspendingrntrue
  \ignorespaces
}
\newcommand{\lawboxuserandnumber}{%
  \iflawboxhaspendingrn
    \randnumberbox{\lawboxpendingrn}%
    \global\lawboxhaspendingrnfalse
  \fi
}
\newcommand{\sourcern}[1]{%
  \ifrandnumbers
    \ifshoworiginal
      \ifshowtranslation\else
        \ifinlawbox\lawboxstorerandnumber{#1}\else\randnumberbox{#1}\fi
      \fi
    \fi
  \fi
}
\newcommand{\transrn}[1]{%
  \ifrandnumbers
    \ifshowtranslation
      \ifinlawbox\lawboxstorerandnumber{#1}\else\randnumberbox{#1}\fi
    \fi
  \fi
}
% ===== 課堂模式：德文原文縮寫註解 =====
% 日常切換只改各判決檔的一行：
%   \classroommodeon        開啟課堂版；註解掉這行即回到一般版。
% 模板預設：
%   1. 課堂模式關閉。
%   2. 課堂模式開啟時，縮寫註解包含「德文全稱＋中文譯名」。
%   3. 課堂模式開啟時，GG/條文編者註仍會自動顯示。
% 進階微調才需要在單案檔覆寫：
%   \classroomabbrzhoff     縮寫註解僅顯示德文全稱。
%   \showeditornotestrue    一般版也顯示編者註。
% 註解策略：同一實際 PDF 頁面中，同一縮寫只在第一次出現處加註。
% 這裡用 zref-abspage 的絕對頁序，不用 \thepage，避免文件內頁碼重置時誤判。
\newif\ifclassroommode
\newif\ifclassroomabbrzh
\classroommodefalse
\classroomabbrzhtrue
\newcommand{\classroommodeon}{\classroommodetrue}
\newcommand{\classroommodeoff}{\classroommodefalse}
\newcommand{\classroomabbrzhon}{\classroomabbrzhtrue}
\newcommand{\classroomabbrzhoff}{\classroomabbrzhfalse}
\newcommand{\abbrnotelabel}{\emph{Abk.}: }
\newcommand{\abbrnote}[3]{%
  \ifclassroommode
    \ifshoworiginal
      \ifintranslation\else
        \footnote{\normalfont\footnotesize\abbrnotelabel #2\ifclassroomabbrzh；{\songfont #3}\fi}%
      \fi
    \fi
  \fi
}
\newcommand{\abbrnoteonce}[4]{%
  #2%
  \ifclassroommode
    \ifshoworiginal
      \ifintranslation\else
        \begingroup
          \edef\abbrcurrentabspage{\number\numexpr\value{abspage}+1\relax}%
          \ifcsname abbrnoteseen@#1@\abbrcurrentabspage\endcsname\else
            \global\expandafter\let\csname abbrnoteseen@#1@\abbrcurrentabspage\endcsname\empty
            \abbrnote{#1}{#3}{#4}%
          \fi
        \endgroup
      \fi
    \fi
  \fi
}
\newcommand{\abbrAbs}{\abbrnoteonce{abs}{Abs.}{Absatz}{項}}
\newcommand{\abbrAAO}{\abbrnoteonce{aao}{a.a.O.}{am angegebenen Ort}{前引處}}
\newcommand{\abbrAE}{\abbrnoteonce{ae}{AE}{Alternativ-Entwurf}{替代草案}}
\newcommand{\abbrArt}{\abbrnoteonce{art}{Art.}{Artikel}{條}}
\newcommand{\abbrAufl}{\abbrnoteonce{aufl}{Aufl.}{Auflage}{版}}
\newcommand{\abbrBd}{\abbrnoteonce{bd}{Bd.}{Band}{卷}}
\newcommand{\abbrBGBl}{\abbrnoteonce{bgbl}{BGBl.}{Bundesgesetzblatt}{聯邦法律公報}}
\newcommand{\abbrBGHSt}{\abbrnoteonce{bghst}{BGHSt}{Entscheidungen des Bundesgerichtshofes in Strafsachen}{聯邦普通法院刑事判決彙編}}
\newcommand{\abbrBRDrucks}{\abbrnoteonce{brdrucks}{BRDrucks.}{Bundesrats-Drucksache}{聯邦參議院文件}}
\newcommand{\abbrBTDrucks}{\abbrnoteonce{btdrucks}{BTDrucks.}{Bundestags-Drucksache}{聯邦議院文件}}
\newcommand{\abbrBundesgesetzbl}{\abbrnoteonce{bundesgesetzbl}{Bundesgesetzbl.}{Bundesgesetzblatt}{聯邦法律公報}}
\newcommand{\abbrBvF}{\abbrnoteonce{bvf}{BvF}{Bundesverfassungsgerichtsverfahren}{聯邦憲法法院程序案號}}
\newcommand{\abbrBvL}{\abbrnoteonce{bvl}{BvL}{Bundesverfassungsgerichtsverfahren}{聯邦憲法法院程序案號}}
\newcommand{\abbrBvR}{\abbrnoteonce{bvr}{BvR}{Bundesverfassungsgerichtsverfahren}{聯邦憲法法院程序案號}}
\newcommand{\abbrBVerfG}{\abbrnoteonce{bverfg}{BVerfG}{Bundesverfassungsgericht}{聯邦憲法法院}}
\newcommand{\abbrBVerfGE}{\abbrnoteonce{bverfge}{BVerfGE}{Entscheidungen des Bundesverfassungsgerichts}{聯邦憲法法院判決集}}
\newcommand{\abbrBVerfGG}{\abbrnoteonce{bverfgg}{BVerfGG}{Bundesverfassungsgerichtsgesetz}{聯邦憲法法院法}}
\newcommand{\abbrDDR}{\abbrnoteonce{ddr}{DDR}{Deutsche Demokratische Republik}{德意志民主共和國；東德}}
\newcommand{\abbrDJT}{\abbrnoteonce{djt}{DJT}{Deutscher Juristentag}{德國法學家大會}}
\newcommand{\abbrDr}{\abbrnoteonce{dr}{Dr.}{Doktor}{Dr.}}
\newcommand{\abbrEuGRZ}{\abbrnoteonce{eugrz}{EuGRZ}{Europäische Grundrechte-Zeitschrift}{歐洲基本權期刊}}
\newcommand{\abbrEV}{\abbrnoteonce{ev}{EV}{Einigungsvertrag}{統一條約}}
\newcommand{\abbrF}{\abbrnoteonce{f}{f.}{folgende}{以下一頁；下一段}}
\newcommand{\abbrFF}{\abbrnoteonce{ff}{ff.}{fortfolgende}{以下數頁；以下各段}}
\newcommand{\abbrGBlDDR}{\abbrnoteonce{gblddr}{GBl. DDR}{Gesetzblatt der Deutschen Demokratischen Republik}{德意志民主共和國法律公報}}
\newcommand{\abbrGG}{\abbrnoteonce{gg}{GG}{Grundgesetz}{基本法}}
\newcommand{\abbrGVBl}{\abbrnoteonce{gvbl}{GVBl.}{Gesetz- und Verordnungsblatt}{法律及命令公報}}
\newcommand{\abbrJR}{\abbrnoteonce{jr}{JR}{Juristische Rundschau}{法學評論}}
\newcommand{\abbrJZ}{\abbrnoteonce{jz}{JZ}{JuristenZeitung}{法學家報}}
\newcommand{\abbrIVm}{\abbrnoteonce{ivm}{i.V.m.}{in Verbindung mit}{結合；連同}}
\newcommand{\abbrMdB}{\abbrnoteonce{mdb}{MdB}{Mitglied des Bundestages}{聯邦議院議員}}
\newcommand{\abbrMwN}{\abbrnoteonce{mwn}{m.w.N.}{mit weiteren Nachweisen}{附進一步文獻／判例出處}}
\newcommand{\abbrNF}{\abbrnoteonce{nf}{n.F.}{neue Fassung}{新版本}}
\newcommand{\abbrAF}{\abbrnoteonce{af}{a.F.}{alte Fassung}{舊版本}}
\newcommand{\abbrNr}{\abbrnoteonce{nr}{Nr.}{Nummer}{款、目或號，依語境}}
\newcommand{\abbrProf}{\abbrnoteonce{prof}{Prof.}{Professor}{教授}}
\newcommand{\abbrRdnr}{\abbrnoteonce{rdnr}{Rdnr.}{Randnummer}{邊碼；段碼}}
\newcommand{\abbrRdnrn}{\abbrnoteonce{rdnrn}{Rdnrn.}{Randnummern}{邊碼；段碼}}
\newcommand{\abbrRGBl}{\abbrnoteonce{rgbl}{RGBl.}{Reichsgesetzblatt}{帝國法律公報}}
\newcommand{\abbrRGSt}{\abbrnoteonce{rgst}{RGSt}{Entscheidungen des Reichsgerichts in Strafsachen}{帝國法院刑事判決彙編}}
\newcommand{\abbrRspr}{\abbrnoteonce{rspr}{Rspr.}{Rechtsprechung}{判例；司法實務}}
\newcommand{\abbrRVO}{\abbrnoteonce{rvo}{RVO}{Reichsversicherungsordnung}{帝國保險法}}
\newcommand{\abbrS}{\abbrnoteonce{s}{S.}{Seite}{頁}}
\newcommand{\abbrSFHG}{\abbrnoteonce{sfhg}{SFHG}{Schwangeren- und Familienhilfegesetz}{孕婦及家庭扶助法}}
\newcommand{\abbrSGBV}{\abbrnoteonce{sgbv}{SGB V}{Fünftes Buch Sozialgesetzbuch}{社會法典第五編}}
\newcommand{\abbrStAG}{\abbrnoteonce{staeg}{StÄG}{Strafrechtsänderungsgesetz}{刑法修正法}}
\newcommand{\abbrStenBer}{\abbrnoteonce{stenber}{StenBer.}{Stenographischer Bericht}{速記紀錄}}
\newcommand{\abbrStGB}{\abbrnoteonce{stgb}{StGB}{Strafgesetzbuch}{刑法}}
\newcommand{\abbrStREG}{\abbrnoteonce{streg}{StREG}{Strafrechtsreform-Ergänzungsgesetz}{刑法改革補充法}}
\newcommand{\abbrStrRG}{\abbrnoteonce{strrg}{StrRG}{Strafrechtsreformgesetz}{刑法改革法}}
\newcommand{\abbrUS}{\abbrnoteonce{us}{U.S.}{United States Reports}{美國最高法院判決彙編}}
\newcommand{\abbrVgl}{\abbrnoteonce{vgl}{vgl.}{vergleiche}{參見}}
\newcommand{\abbrVVDStRL}{\abbrnoteonce{vvdstrl}{VVDStRL}{Veröffentlichungen der Vereinigung der Deutschen Staatsrechtslehrer}{德國公法教師協會論文集}}
\newcommand{\abbrWP}{\abbrnoteonce{wp}{WP.}{Wahlperiode}{屆期}}
\newcommand{\abbrWp}{\abbrnoteonce{wp}{Wp.}{Wahlperiode}{屆期}}
\newcommand{\abbrZStrW}{\abbrnoteonce{zstrw}{ZStrW}{Zeitschrift für die gesamte Strafrechtswissenschaft}{整體刑法學期刊}}
\newcommand{\recordsourcerandstart}{%
  \ifrandnumbers
    \setcounter{lastsourcerandstart}{\value{randnumber}}%
    \global\lastsourcehasrandnumberstrue
  \else
    \global\lastsourcehasrandnumbersfalse
  \fi
}
\newlength{\biparagraphskip}
\setlength{\biparagraphskip}{0.52em}
\newlength{\lawboxblockbeforeskip}
\setlength{\lawboxblockbeforeskip}{0.72em}
\newlength{\lawboxblockafterskip}
\setlength{\lawboxblockafterskip}{0.92em}
\newif\iflawboxpartendedblock
\lawboxpartendedblockfalse
\newcommand{\sourceparstyle}{\footnotesize\color{sourceink}\setstretch{1.12}\RaggedRight\emergencystretch=3em\setlength{\parskip}{0.42em}\obeylines\selectfont}
\newcommand{\transparstyle}{\songfont\color{caseink}\setstretch{1.12}\setlength{\parindent}{0pt}\setlength{\parskip}{0.42em}}
\newcommand{\bilingualpairskip}{%
  \par\addvspace{\biparagraphskip}%
  \switchcolumn
  \par\addvspace{\biparagraphskip}%
  \switchcolumn*%
  \global\intranslationfalse
}
\newcommand{\bilingualboxpairskip}{%
  \par
  \switchcolumn
  \par
  \switchcolumn*%
  \global\intranslationfalse
}
\newcommand{\bilinguallawboxblockskip}{%
  \switchcolumn*%
  \par\addvspace{\lawboxblockafterskip}%
  \switchcolumn
  \par\addvspace{\lawboxblockafterskip}%
  \switchcolumn*%
  \global\intranslationfalse
}
\newcommand{\bikeepwithnext}[1][4]{%
  \ifbilingualcolumns
    \syncleft
    \Needspace{#1\baselineskip}%
    \switchcolumn
    \Needspace{#1\baselineskip}%
    \switchcolumn*%
    \global\intranslationfalse
  \else
    \Needspace{#1\baselineskip}%
  \fi
}
\NewEnviron{sourcepara}[1][]{
  \recordsourcerandstart
  \ifshoworiginal
    \ifbilingualcolumns
      \ifintranslation\switchcolumn*\global\intranslationfalse\fi
    \else
      \Needspace{5\baselineskip}
    \fi
    \ifbilingualcolumns\else\par\addvspace{0.35em}\fi
    {\sourceparstyle
    \noindent\BODY\par}
    \ifbilingualcolumns\else\addvspace{0.25em}\fi
    \ifbilingualcolumns
      \switchcolumn
      \global\intranslationtrue
    \fi
  \else
    \ifrandnumbers
      \begingroup
        \setbox0=\vbox{\hsize=\linewidth\sourceparstyle\noindent\BODY\par}%
      \endgroup
    \fi
  \fi
}
\NewEnviron{transpara}{
  \ifshowtranslation
    \setcounter{randnumberlocal}{\value{lastsourcerandstart}}%
    \ifbilingualcolumns
      \ifintranslation\else\switchcolumn\global\intranslationtrue\fi
      {\transparstyle\setlength{\leftskip}{0.55em}\setlength{\rightskip}{0.15em plus 1em}\addtolength{\linewidth}{-0.55em}\BODY\par}
      \switchcolumn*
      \bilingualpairskip
    \else
      {\transparstyle\BODY\par}
    \fi
  \else
    \ifbilingualcolumns
      \ifintranslation\switchcolumn*\global\intranslationfalse\fi
    \fi
  \fi
}
\NewEnviron{sourceboxpara}[1][]{
  \global\lawboxpartendedblockfalse
  \recordsourcerandstart
  \ifshoworiginal
    \ifbilingualcolumns
      \ifintranslation\switchcolumn*\global\intranslationfalse\fi
    \else
      \Needspace{3\baselineskip}
    \fi
    {\sourceparstyle
    \BODY\par}
    \ifbilingualcolumns\else
      \iflawboxpartendedblock\addvspace{\lawboxblockafterskip}\fi
    \fi
    \ifbilingualcolumns
      \switchcolumn
      \global\intranslationtrue
    \fi
  \fi
}
\NewEnviron{transboxpara}{
  \global\lawboxpartendedblockfalse
  \ifshowtranslation
    \setcounter{randnumberlocal}{\value{lastsourcerandstart}}%
    \ifbilingualcolumns
      \ifintranslation\else\switchcolumn\global\intranslationtrue\fi
      {\transparstyle\BODY\par}
      \iflawboxpartendedblock
        \bilinguallawboxblockskip
      \else
        \switchcolumn*
        \bilingualboxpairskip
      \fi
    \else
      {\transparstyle\BODY\par}
      \iflawboxpartendedblock\addvspace{\lawboxblockafterskip}\fi
    \fi
  \else
    \ifbilingualcolumns
      \ifintranslation\switchcolumn*\global\intranslationfalse\fi
    \fi
  \fi
}
\newcommand{\leitsatznum}[1]{\noindent\hangindent=3.0em\hangafter=1\makebox[2.25em][r]{\bfseries #1.}\hspace{0.55em}\ignorespaces}
\newcommand{\syncleft}{\ifbilingualcolumns\ifintranslation\switchcolumn*\global\intranslationfalse\fi\fi}
\pretocmd{\section}{\syncleft\clearpage}{}{}
\pretocmd{\subsection}{\syncleft}{}{}
\pretocmd{\subsubsection}{\syncleft}{}{}
\newcommand{\biheadingfont}{\songfont}
\newcommand{\bitocentry}[2]{%
  \ifshoworiginal
    \ifshowtranslation
      \ifstrequal{#1}{#2}{#1}{#1\space #2}%
    \else
      #1%
    \fi
  \else
    #2%
  \fi
}
\pdfstringdefDisableCommands{%
  \def\bitocentry#1#2{#1}%
}
\newcommand{\biheading}[7]{%
  \syncleft#1\bikeepwithnext[#3]\phantomsection\addcontentsline{toc}{#2}{\bitocentry{#6}{#7}}%
  \ifbilingualcolumns
    {#4 #6\par}%
    \switchcolumn
    {#4\biheadingfont #7\par}%
    \switchcolumn*%
    \global\intranslationfalse
    \par\addvspace{#5}%
    \switchcolumn
    \par\addvspace{#5}%
    \switchcolumn*%
  \else
    \ifshoworiginal{#4 #6\par}\fi
    \ifshowtranslation{#4\biheadingfont #7\par}\fi
    \addvspace{#5}%
  \fi
}
\newcommand{\termsectionheading}[1]{%
  \par\Needspace{9\baselineskip}%
  \vspace{0.65em}%
  {\songfont\bfseries #1\par}%
  \vspace{0.2em}%
}
% 章節換頁：
%   雙語 paracol 欄內不要用 \clearpage；它會在同步兩欄時吐出只有欄線/頁碼的空頁。
%   \newpage 足以讓下一個 section 起新頁，又不會強制清空所有待排材料。
%   單語模式不在 paracol 內，仍可用 \clearpage。
\newcommand{\bisectionpagebreak}{\ifbilingualcolumns\newpage\else\clearpage\fi}
\newcommand{\bisection}[2]{%
  \biheading{\iffirstbilingualsection\global\firstbilingualsectionfalse\else\bisectionpagebreak\fi}{section}{6}{\centering\Large\bfseries}{0.8em}{#1}{#2}%
}
\newcommand{\bisubsection}[2]{%
  \biheading{}{subsection}{5}{\large\bfseries}{0.55em}{#1}{#2}%
}
\newcommand{\bisubsubsection}[2]{%
  \biheading{}{subsubsection}{4}{\normalsize\bfseries}{0.45em}{#1}{#2}%
}
\newcommand{\bicompositeheading}[4]{%
  \biheading{}{subsection}{5}{\large\bfseries}{0.16em}{#1}{#2}%
  \biheading{}{subsubsection}{3}{\normalsize\bfseries}{0.45em}{#3}{#4}%
}
\newif\iflawboxfirstitem
\lawboxfirstitemfalse
\newcommand{\lawboxprespace}[1]{%
  \par
  \iflawboxfirstitem
    \global\lawboxfirstitemfalse
  \else
    \addvspace{#1}%
  \fi
}
\newtcolorbox{lawboxinner}{
  caseboxbase,
  colframe=sourceline!85!black,
  colback=white,
  left=1.05em,
  right=1.05em,
  top=0.62em,
  bottom=0.62em,
  before upper={\global\inlawboxtrue\global\lawboxfirstitemtrue\global\lawboxhaspendingrnfalse\ifintranslation\lawboxtransfont\else\lawboxsourcefont\fi\RaggedRight\setlength{\parindent}{0pt}\setlength{\parskip}{0pt}\ifintranslation\setstretch{\lawboxtransstretch}\else\setstretch{\lawboxsourcestretch}\fi},
  after upper={\global\lawboxhaspendingrnfalse\global\inlawboxfalse}
}
\tcbset{
  lawboxpartbase/.style={
    enhanced,
    breakable=false,
    width=0.985\linewidth,
    colback=white,
    frame hidden,
    boxrule=0pt,
    arc=0pt,
    left=1.05em,
    right=1.05em,
    boxsep=1.5mm,
    before skip=0pt,
    after skip=0pt,
    before upper={\global\inlawboxtrue\global\lawboxfirstitemtrue\global\lawboxhaspendingrnfalse\ifintranslation\lawboxtransfont\else\lawboxsourcefont\fi\RaggedRight\setlength{\parindent}{0pt}\setlength{\parskip}{0pt}\ifintranslation\setstretch{\lawboxtransstretch}\else\setstretch{\lawboxsourcestretch}\fi},
    after upper={\global\lawboxhaspendingrnfalse\global\inlawboxfalse},
    borderline west={0.28pt}{0pt}{sourceline!85!black},
    borderline east={0.28pt}{0pt}{sourceline!85!black},
  },
  lawboxpartfirst/.style={
    lawboxpartbase,
    top=0.62em,
    bottom=0.04em,
    before skip=\lawboxblockbeforeskip,
    borderline north={0.28pt}{0pt}{sourceline!85!black},
  },
  lawboxpartmiddle/.style={
    lawboxpartbase,
    top=0.04em,
    bottom=0.04em,
  },
  lawboxpartlast/.style={
    lawboxpartbase,
    top=0.04em,
    bottom=0.62em,
    after skip=0pt,
    borderline south={0.28pt}{0pt}{sourceline!85!black},
  }
}
\newtcolorbox{lawboxpartinner}[1]{#1}
\newtcolorbox{termboxinner}{
  caseboxbase,
  colframe=noteblue!70!black,
  colback=casepaper,
  before upper={\songfont\setlength{\parindent}{0pt}\setlength{\parskip}{0.35em}}
}
\newtcolorbox{deboxinner}{
  caseboxbase,
  colframe=notegray!72!black,
  colback=white,
  left=1.05em,
  right=1.05em,
  top=0.62em,
  bottom=0.62em,
  before upper={\global\inlawboxtrue\global\lawboxfirstitemtrue\global\lawboxhaspendingrnfalse\small\setlength{\parindent}{0pt}\setlength{\parskip}{0.35em}},
  after upper={\global\lawboxhaspendingrnfalse\global\inlawboxfalse}
}
\NewEnviron{lawbox}{
  \begin{lawboxinner}\BODY\end{lawboxinner}
}
\NewEnviron{lawboxpart}[2]{
  \begin{lawboxpartinner}{lawboxpart#1,equal height group=#2}\BODY\end{lawboxpartinner}%
  \ifstrequal{#1}{last}{\global\lawboxpartendedblocktrue}{}%
}
\NewEnviron{termbox}{
  \par\addvspace{0.55em}
  {\color{noteblue!70!black}\rule{\linewidth}{0.28pt}}\par
  \begingroup
    \songfont\small\color{caseink}
    \setlength{\parindent}{0pt}
    \setlength{\parskip}{0.24em}
    \BODY
    \par
  \endgroup
  {\color{noteblue!70!black}\rule{\linewidth}{0.28pt}}\par
  \addvspace{0.55em}
}
\NewEnviron{debox}{
  \begin{deboxinner}\BODY\end{deboxinner}
}
\providecommand{\paragraphmark}[1]{}
\newcommand{\lawboxtitleafter}{\addvspace{0.06em}}
\newcommand{\lawtitle}[1]{\lawboxprespace{0.22em}{\lawboxtransfont\bfseries\lawboxuserandnumber #1}\par\lawboxtitleafter}
\newcommand{\absno}[2]{\lawboxprespace{0.04em}\noindent\lawboxuserandnumber\hangindent=2.6em\hangafter=1\makebox[2.6em][l]{(#1)}#2\par}
\newcommand{\nrno}[2]{\lawboxprespace{0pt}\noindent\lawboxuserandnumber\hspace*{2.6em}\hangindent=4.4em\hangafter=1\makebox[1.8em][l]{#1.}#2\par}
\newcommand{\letterno}[2]{\lawboxprespace{0pt}\noindent\lawboxuserandnumber\hspace*{4.4em}\hangindent=6.0em\hangafter=1\makebox[1.6em][l]{#1)}#2\par}
\newcommand{\sourcelawtitle}[1]{\lawboxprespace{0.22em}{\lawboxsourcefont\bfseries\lawboxuserandnumber #1}\par\lawboxtitleafter}
\newcommand{\sourcelawsubtitle}[1]{\lawboxprespace{0.04em}{\lawboxsourcefont\bfseries\lawboxuserandnumber #1}\par\lawboxtitleafter}
\newcommand{\sourceabsno}[2]{\lawboxprespace{0.04em}\noindent\lawboxuserandnumber\hangindent=2.45em\hangafter=1\makebox[2.45em][l]{(#1)}#2\par}
\newcommand{\sourcenrno}[2]{\lawboxprespace{0pt}\noindent\lawboxuserandnumber\hspace*{2.45em}\hangindent=4.25em\hangafter=1\makebox[1.8em][l]{#1.}#2\par}
\newcommand{\sourceletterno}[2]{\lawboxprespace{0pt}\noindent\lawboxuserandnumber\hspace*{4.25em}\hangindent=5.85em\hangafter=1\makebox[1.6em][l]{#1)}#2\par}
\newcommand{\sourcequotepara}[1]{\lawboxprespace{0.04em}\noindent\lawboxuserandnumber\hangindent=1.2em\hangafter=1 #1\par}
\newcommand{\sourcelawline}[1]{\lawboxprespace{0.04em}\noindent\lawboxuserandnumber #1\par}
\newcommand{\lawline}[1]{\lawboxprespace{0.04em}\noindent\lawboxuserandnumber #1\par}
\newcommand{\pendingtranslation}{\par{\songfont\footnotesize\color{pendingink}（中文譯文待補。）}\par}
\newcommand{\tocpart}[1]{\par\addvspace{0.52em}{\footnotesize\bfseries\color{pendingink}#1\par}\addvspace{0.16em}}
\newcommand{\tocdashline}{\par\addvspace{0.48em}\noindent{\color{notegray}\xleaders\hbox to 0.82em{\hfil\rule{0.42em}{0.28pt}\hfil}\hskip\linewidth}\par\addvspace{0.38em}}
% \tocpanel{font-cmd}{content}：將目錄置中於所在欄寬（雙語欄適用）。
% A3 橫式使用 0.58\linewidth，A4 橫式使用 0.84\linewidth，
% 使兩種紙張下目錄欄內容實際寬度相近（均約 100–105 mm）。
\newcommand{\tocpanel}[2]{%
  \makebox[\linewidth][c]{%
    \ifx\bilingualpaper\bilingualpaperaThree
      \begin{minipage}{0.58\linewidth}%
    \else
      \begin{minipage}{0.84\linewidth}%
    \fi
    \toccolstyle
    #1#2%
    \end{minipage}%
  }%
}
% ===== 首頁簡目錄的兩層 description list 縮排 =====
% enumitem 的 description list 可把每一列拆成：
%   [label]  labelsep  body text
% 這裡的「起點」都以所在 minipage/欄位的左緣為 0。
%
% 核心規則：
%   labelindent = label 欄位從左緣開始的位置。
%   labelwidth  = label 欄位本身寬度；標籤太長（如 Abw. Meinung）就加大它。
%   labelsep    = label 欄位右緣到正文的距離。
%   leftmargin  = 正文 body text 的起點。判斷層級視覺時主要看這個。
%
% 所以：
%   想讓「Begründung / 判決理由」整列正文往右或往左：改 tocitems 的 leftmargin。
%   想讓「A. Verfahren und Gegenstand / A. 程序與審查標的」正文往右或往左：
%     改 tocsubitems 的 leftmargin。
%   想讓 A./B./C. 這些標籤本身往右或往左，但正文不動：改 tocsubitems 的 labelindent。
%   想讓 A./B./C. 與正文間距變大或變小：改 tocsubitems 的 labelsep。
%
% 層級原則：
%   tocsubitems.leftmargin 必須大於 tocitems.leftmargin，
%   否則子層級正文會看起來比「Gründe / 理由」還靠左。
\newlist{tocitems}{description}{1}
\setlist[tocitems]{%
  labelindent=0pt,    % 第一層 label 欄位起點；0pt 表示貼齊目錄欄左緣。
  leftmargin=2.54cm,  % 第一層正文起點；例如 Begründung / 判決理由 的左緣。
  labelwidth=2.16cm,  % 第一層 label 欄位寬度；需容納 Leitsätze、Abw. Meinung、不同意見等。
  labelsep=0.32cm,    % 第一層 label 與正文的水平間距；只改間距，不改正文起點。
  itemsep=0.28em,     % 第一層各項之間的垂直距離。
  topsep=0.20em,      % 第一層 list 上下與前後內容的垂直距離。
  parsep=0pt,         % 同一 item 內段落之間的距離；目錄項通常不需要。
  partopsep=0pt,      % list 若緊接段落開始時的額外距離；關閉以免目錄高度飄動。
  align=left,         % label 欄內左對齊；長短標籤不靠右貼齊。
  font=\normalfont\bfseries % label 的字型；正文不受這行影響。
}
\newlist{tocsubitems}{description}{1}
\setlist[tocsubitems]{%
  labelindent=1cm,    % 第二層 label 起點；只控制 A./B./C. 本身的位置。
  leftmargin=3.00cm,  % 第二層正文起點；必須大於 2.54cm，才會比 Begründung / 判決理由更右。
  labelwidth=0.5cm,   % 第二層 label 欄寬；A./B./C. 很短，可小於第一層。
  labelsep=1.5cm,    % 第二層 label 與正文間距；A. 與 Verfahren / 程序 的距離看這裡。
  itemsep=0.12em,     % 第二層各項較密，避免 A.–E. 佔太高。
  topsep=0.04em,      % 第二層 list 與 Gründe / 理由之間的垂直距離。
  parsep=0pt,         % 同一第二層 item 內段落距離；目錄項通常不需要。
  partopsep=0pt,      % 關閉額外段落起始距離，避免版面不穩。
  align=left,         % A./B./C. label 欄內左對齊。
  font=\normalfont\bfseries, % A./B./C. label 加粗；正文不受這行影響。
  before={}           % 不在第二層 list 前自動插入額外內容。
}
\newlist{termitems}{description}{1}
\ifbilingualcolumns
  \setlist[termitems]{%
    leftmargin=5.2cm,
    labelwidth=4.8cm,
    labelsep=0.35cm,
    itemsep=0.16em,
    topsep=0.2em,
    parsep=0pt,
    partopsep=0pt,
    font=\normalfont\bfseries,
    style=nextline
  }
\else
  \setlist[termitems]{%
    leftmargin=0pt,
    labelwidth=0pt,
    labelsep=0pt,
    itemsep=0.24em,
    topsep=0.2em,
    parsep=0pt,
    partopsep=0pt,
    font=\normalfont\bfseries,
    style=nextline
  }
\fi
\newlist{abbritems}{description}{1}
\setlist[abbritems]{%
  leftmargin=2.38cm,
  labelwidth=2.02cm,
  labelsep=0.28cm,
  itemsep=0.16em,
  topsep=0.2em,
  parsep=0pt,
  partopsep=0pt,
  align=left,
  font=\normalfont\bfseries
}
\ifbilingualcolumns
  \newenvironment{termcolumns}{\begin{multicols}{2}}{\end{multicols}}
\else
  \newenvironment{termcolumns}{}{}
\fi

% ===== 編者註腳開關 =====
% 一般版預設不顯示編者註，避免正文腳註過密。
% 若開啟 \classroommodeon，GG/條文編者註仍會自動顯示，無需另開本開關。
% 若一般版也要顯示編者註，可在單案檔 \input 後加 \showeditornotestrue。
\newif\ifshoweditornotes
\showeditornotesfalse

% ===== 目錄排版輔助 =====
\newcommand{\toccolstyle}{\raggedright\arraybackslash}
\newcommand{\tocsingleindent}{\leftskip=1.45cm\rightskip=1.2cm}

% ===== 行內引文環境（德文原文與中文譯文各一，帶左側灰色邊線）=====
\newcommand{\quoteparbreak}{\par\addvspace{0.48em}}
\newcommand{\sourcequoteinner}[1]{%
  \begin{tcolorbox}[
    enhanced, breakable, width=\dimexpr\linewidth-0.8em\relax, frame hidden,
    colback=white, coltext=caseink, boxrule=0pt,
    borderline west={0.55pt}{0.88em}{notegray},
    left=1.78em, right=0.72em, top=0.18em, bottom=0.18em,
    boxsep=0pt, before skip=0.24em, after skip=0.24em,
    before upper={\normalfont\footnotesize\color{caseink}\setlength{\parindent}{0pt}\setlength{\parskip}{0.34em}\raggedright\setlength{\rightskip}{0pt plus 1em}}
  ]%
  #1\par
  \end{tcolorbox}%
}
\newcommand{\transquoteinner}[1]{%
  \begin{tcolorbox}[
    enhanced, breakable, width=\dimexpr\linewidth-0.8em\relax, frame hidden,
    colback=white, coltext=caseink, boxrule=0pt,
    borderline west={0.55pt}{0.88em}{notegray},
    left=1.78em, right=0.72em, top=0.18em, bottom=0.18em,
    boxsep=0pt, before skip=0.24em, after skip=0.24em,
    before upper={\songfont\small\color{caseink}\setlength{\parindent}{0pt}\setlength{\parskip}{0.34em}\raggedright\setlength{\rightskip}{0pt plus 1em}}
  ]%
  #1\par
  \end{tcolorbox}%
}

% ===== 大綱（Gliederung）Rn 輔助命令 =====
\newcommand{\outrn}[2]{\enspace{\normalfont\footnotesize\color{pendingink}(Rn\,#1–#2)}}
\newcommand{\outrnsingle}[1]{\enspace{\normalfont\footnotesize\color{pendingink}(Rn\,#1)}}
\newcommand{\outrnzh}[2]{\enspace{\normalfont\footnotesize\color{pendingink}（第\,#1–#2\,段）}}
\newcommand{\outrnzhs}[1]{\enspace{\normalfont\footnotesize\color{pendingink}（第\,#1\,段）}}
% ===== 大綱雙語對照表格輔助命令（三欄：段碼 │ 德文 │ 中文）=====
% 欄 1：段碼（由欄定義自動格式化：小字、等寬、右對齊、pendingink）
% 欄 2：德文標籤＋正文
% 欄 3：中文標籤＋正文
%
% 無段碼的列（\omainrow、\osecrow），欄 1 以 & 空置。
% 有段碼的列（\osecrowrn、\osubrow、\ointrow），#1 為預格式化字串
%   例：\osubrow{2–49}{I.}{...}{I.}{...}
%       \ointrow{107}{...}{...}
%
% \opartrow：段組標題（橫跨三欄）
\newcommand{\opartrow}[2]{%
  \noalign{\vspace{0.30em}}%
  \multicolumn{3}{@{}l@{}}{%
    \footnotesize\bfseries\color{pendingink}#1\hfill\songfont#2%
  }\\[0.06em]%
}
% \odashrow：虛線分隔（橫跨三欄）
\newcommand{\odashrow}{%
  \noalign{\vspace{0.28em}}%
  \multicolumn{3}{@{}l@{}}{\color{notegray}%
    \xleaders\hbox to 0.82em{\hfil\rule{0.42em}{0.28pt}\hfil}\hskip 0.88\linewidth}%
  \\[0.24em]%
}
% \odissentpart：不同意見書區段標頭。比一般 \opartrow 更像附錄性轉場。
\newcommand{\odissentpart}[2]{%
  \noalign{\vspace{0.42em}}%
  \multicolumn{3}{@{}p{0.88\textwidth}@{}}{%
    \color{notegray}\rule{\linewidth}{0.18pt}%
  }\\[-0.18em]%
  \multicolumn{3}{@{}p{0.88\textwidth}@{}}{%
    \makebox[\linewidth][s]{%
      \footnotesize\bfseries\color{pendingink}#1%
      \hfill{\songfont #2}%
    }%
  }\\[-0.02em]%
}
% \odissentopinion：不同意見書的署名與一行摘要，右欄保留段碼範圍。
\newcommand{\odissentopinion}[5]{%
  \footnotesize\hangindent=0.52cm\hangafter=1%
    {\bfseries\color{caseink}#2}\enspace{\color{notegray}\textperiodcentered}\enspace\color{caseink}#3%
  &\footnotesize\hangindent=0.52cm\hangafter=1%
    {\songfont\bfseries\color{caseink}#4}\enspace{\color{notegray}\textperiodcentered}\enspace\songfont\color{caseink}#5%
  &#1%
  \\[0.12em]%
}
% \omainrow：頂層項目，無段碼（Leitsätze、Urteil、Gründe …）
\newcommand{\omainrow}[4]{%
  \footnotesize\hangindent=1.92cm\hangafter=1%
    \makebox[1.92cm][l]{\bfseries\color{caseink}#1}\color{caseink}#2%
  &\footnotesize\hangindent=1.92cm\hangafter=1%
    \makebox[1.92cm][l]{\bfseries\color{caseink}#3}\songfont\color{caseink}#4%
  &%
  \\[0.15em]%
}
% \omainrowrn：頂層項目，有段碼（不同意見等標籤寬於 \osecrow 框者）
\newcommand{\omainrowrn}[5]{%
  \footnotesize\hangindent=1.92cm\hangafter=1%
    \makebox[1.92cm][l]{\bfseries\color{caseink}#2}\color{caseink}#3%
  &\footnotesize\hangindent=1.92cm\hangafter=1%
    \makebox[1.92cm][l]{\bfseries\color{caseink}#4}\songfont\color{caseink}#5%
  &#1%
  \\[0.15em]%
}
% \osecrow：一級子項，無段碼（A.、C.、D. 等有子項者）
\newcommand{\osecrow}[4]{%
  \footnotesize\hspace{1.10cm}\hangindent=1.98cm\hangafter=1%
    \makebox[0.86cm][l]{\bfseries\color{caseink}#1}\color{caseink}#2%
  &\footnotesize\hspace{1.10cm}\hangindent=1.98cm\hangafter=1%
    \makebox[0.86cm][l]{\bfseries\color{caseink}#3}\songfont\color{caseink}#4%
  &%
  \\[0.11em]%
}
% \osecrowrn：一級子項，有段碼（B.、E. 等完整段落範圍者）
% #1=段碼字串（如 101–106）, #2=DE標籤, #3=DE正文, #4=ZH標籤, #5=ZH正文
\newcommand{\osecrowrn}[5]{%
  \footnotesize\hspace{1.10cm}\hangindent=1.98cm\hangafter=1%
    \makebox[0.86cm][l]{\bfseries\color{caseink}#2}\color{caseink}#3%
  &\footnotesize\hspace{1.10cm}\hangindent=1.98cm\hangafter=1%
    \makebox[0.86cm][l]{\bfseries\color{caseink}#4}\songfont\color{caseink}#5%
  &#1%
  \\[0.11em]%
}
% \osubrow：二級子項，有段碼範圍（I.、II.、A.I. …）
% #1=段碼字串, #2=DE標籤, #3=DE正文, #4=ZH標籤, #5=ZH正文
\newcommand{\osubrow}[5]{%
  \footnotesize\hspace{1.40cm}\hangindent=2.30cm\hangafter=1%
    \makebox[0.90cm][l]{\bfseries\color{caseink}#2}\color{caseink}#3%
  &\footnotesize\hspace{1.40cm}\hangindent=2.30cm\hangafter=1%
    \makebox[0.90cm][l]{\bfseries\color{caseink}#4}\songfont\color{caseink}#5%
  &#1%
  \\[0.09em]%
}
% \ointrow：引入段落，有單一段碼，無標籤
% #1=段碼字串, #2=DE正文, #3=ZH正文
\newcommand{\ointrow}[3]{%
  \footnotesize\hspace{0.52cm}\hangindent=0.52cm\hangafter=1\color{caseink}#2%
  &\footnotesize\hspace{0.52cm}\hangindent=0.52cm\hangafter=1\songfont\color{caseink}#3%
  &#1%
  \\[0.09em]%
}
% \osubsubrow：三級子項（1.、2.、3.）
% #1=段碼字串, #2=DE標籤, #3=DE正文, #4=ZH標籤, #5=ZH正文
\newcommand{\osubsubrow}[5]{%
  \footnotesize\hspace{1.80cm}\hangindent=2.48cm\hangafter=1%
    \makebox[0.68cm][l]{\bfseries\color{caseink}#2}\color{caseink}#3%
  &\footnotesize\hspace{1.80cm}\hangindent=2.48cm\hangafter=1%
    \makebox[0.68cm][l]{\bfseries\color{caseink}#4}\songfont\color{caseink}#5%
  &#1%
  \\[0.08em]%
}
% \osubsubsubrow：四級子項（a）、b））
% #1=段碼字串, #2=DE標籤, #3=DE正文, #4=ZH標籤, #5=ZH正文
\newcommand{\osubsubsubrow}[5]{%
  \footnotesize\hspace{2.10cm}\hangindent=2.68cm\hangafter=1%
    \makebox[0.58cm][l]{\bfseries\color{caseink}#2}\color{caseink}#3%
  &\footnotesize\hspace{2.10cm}\hangindent=2.68cm\hangafter=1%
    \makebox[0.58cm][l]{\bfseries\color{caseink}#4}\songfont\color{caseink}#5%
  &#1%
  \\[0.07em]%
}
% \osubsubsubsubrow：五級子項（aa）、bb））
% 標籤框 0.62cm：足以容納「aa)」在 footnotesize bold 下（約 0.50cm）並留有間距
% #1=段碼字串, #2=DE標籤, #3=DE正文, #4=ZH標籤, #5=ZH正文
\newcommand{\osubsubsubsubrow}[5]{%
  \footnotesize\hspace{2.38cm}\hangindent=3.06cm\hangafter=1%
    \makebox[0.68cm][l]{\bfseries\color{caseink}#2}\color{caseink}#3%
  &\footnotesize\hspace{2.38cm}\hangindent=3.06cm\hangafter=1%
    \makebox[0.68cm][l]{\bfseries\color{caseink}#4}\songfont\color{caseink}#5%
  &#1%
  \\[0.06em]%
}
% \outlinepage：雙語對照大綱頁包裝環境（longtable 自動跨頁）
% 三欄：德文欄 + 中文欄 + 段碼欄（最右，不折行）
% 段碼欄寬度：1.80cm，足以容納最長段碼如「309–348」
% DE/ZH 欄寬：(0.87\textwidth - 2.08cm) / 2
%   驗算：2×欄 + 0.05\textwidth + 0.28cm gap + 1.80cm = 0.87tw - 2.08cm + 0.05tw + 2.08cm = 0.92tw ✓
\newenvironment{outlinepage}{%
  \vspace*{0.40cm}%
  \setlength{\LTleft}{0.04\textwidth}%
  \setlength{\LTright}{0.04\textwidth}%
  \setlength{\tabcolsep}{0pt}%
  % 大綱列距由 longtable 的 \arraystretch 與各 row macro 結尾的 \\[...em] 共同決定。
  % 這裡控制「每一列本身的表格 strut 高度」；數值越大，A./I./1. 各項之間越像有空行。
  % A3 原本用 0.62，視覺上沒有空行；A4 也沿用同值，避免切到 A4 後項目被撐開。
  % 若日後覺得大綱太擠，只微調這個值（例如 0.68），不要改各 \omainrow/\osubrow 的縮排。
  \renewcommand{\arraystretch}{0.62}%
  \begin{longtable}{%
    >{\vspace{0pt}\raggedright\arraybackslash}p{\dimexpr(0.87\textwidth-2.08cm)/2\relax}%
    @{\hspace{0.025\textwidth}}%
    !{\color{notegray}\vrule width 0.12pt}%
    @{\hspace{0.025\textwidth}}%
    >{\vspace{0pt}\raggedright\arraybackslash}p{\dimexpr(0.87\textwidth-2.08cm)/2\relax}%
    @{\hspace{0.28cm}}%
    >{\vspace{0pt}\footnotesize\normalfont\color{pendingink}\raggedright\arraybackslash}p{1.80cm}%
    @{}%
  }%
  % 首頁欄標籤：不要橫跨置中；分別放在德文欄與中文欄的實際欄位起點。
  \footnotesize\bfseries\color{sourceink}Gliederung%
  &\footnotesize\bfseries\songfont\color{sourceink}大綱%
  &%
  \\[0.36em]%
  \endfirsthead%
  % 續頁欄標籤：同樣分欄定位，以灰色細體區分；無需標注「續」。
  \footnotesize\color{pendingink}Gliederung%
  &\footnotesize\songfont\color{pendingink}大綱%
  &%
  \\[0.24em]%
  \endhead%
}{%
  \end{longtable}%
}
 之前設定；
%   預設 chinese / a4）
\ifdefined\docmode\else\def\docmode{chinese}\fi
\ifdefined\bilingualpaper\else\def\bilingualpaper{a4}\fi
\newif\ifshoworiginal
\newif\ifshowtranslation
\newif\ifbilinguallandscape
\newif\ifbilingualcolumns
\newif\ifintranslation
\newif\iffirstbilingualsection

\showoriginalfalse
\showtranslationfalse
\bilinguallandscapefalse
\bilingualcolumnsfalse
\intranslationfalse
\firstbilingualsectionfalse

\def\modechinese{chinese}
\def\modegerman{german}
\def\modebilingual{bilingual}
\def\bilingualpaperaThree{a3}
\def\bilingualpaperaFour{a4}

\ifx\docmode\modechinese
  \showtranslationtrue
\fi

\ifx\docmode\modegerman
  \showoriginaltrue
\fi

\ifx\docmode\modebilingual
  \showoriginaltrue
  \showtranslationtrue
  \bilinguallandscapetrue
  \bilingualcolumnstrue
\fi

% 編排模式說明：
% chinese  → \showoriginalfalse  \showtranslationtrue  （A4 直式，純中文）
% german   → \showoriginaltrue   \showtranslationfalse （A4 直式，純德文）
% bilingual→ \showoriginaltrue   \showtranslationtrue  \bilingualcolumnstrue（A3/A4 橫式對照）

\ifbilingualcolumns
  \ifx\bilingualpaper\bilingualpaperaFour
    \geometry{
      a4paper,
      landscape,
      left=1.25cm,
      right=2.25cm,
      top=1.25cm,
      bottom=1.75cm,
      footskip=8.5mm,
      marginparwidth=0.75cm,
      marginparsep=0.25cm
    }
  \else
    \geometry{
      a3paper,
      landscape,
      left=1.55cm,
      right=2.75cm,
      top=1.55cm,
      bottom=2.05cm,
      footskip=9.5mm,
      marginparwidth=0.9cm,
      marginparsep=0.35cm
    }
  \fi
\else
\geometry{
  a4paper,
  portrait,
  left=2.2cm,
  right=2.6cm,
  top=2.0cm,
  bottom=2.0cm,
  marginparwidth=0.8cm,
  marginparsep=0.3cm
}
\fi
% ===== 時間戳基礎設施 =====
\newcount\stamptimeval
\newcount\stampminuteval
\newcommand{\stamphh}{%
  \stamptimeval=\time
  \divide\stamptimeval by 60
  \ifnum\stamptimeval<10 0\fi
  \the\stamptimeval}
\newcommand{\stampmm}{%
  \stamptimeval=\time
  \stampminuteval=\time
  \divide\stampminuteval by 60
  \multiply\stampminuteval by 60
  \advance\stamptimeval by -\stampminuteval
  \ifnum\stamptimeval<10 0\fi
  \the\stamptimeval}
\newcommand{\stampwkde}[1]{%
  \ifcase#1Montag\or Dienstag\or Mittwoch\or Donnerstag\or Freitag\or Samstag\or Sonntag\fi}
\newcommand{\stampwkzh}[1]{%
  \ifcase#1 一\or 二\or 三\or 四\or 五\or 六\or 日\fi}
\newcommand{\stampmonthde}[1]{%
  \ifcase#1\or Januar\or Februar\or März\or April\or Mai\or Juni\or
  Juli\or August\or September\or Oktober\or November\or Dezember\fi}
\newcount\stampjdval
\newcount\stampwdval
\pgfcalendardatetojulian{\the\year-\the\month-\the\day}{\stampjdval}
\pgfcalendarjuliantoweekday{\the\stampjdval}{\stampwdval}
\newcommand{\timestampzh}{%
  \songfont 最後修改：\the\year 年\the\month 月\the\day 日%
  % （星期\stampwkzh{\the\stampwdval}）%
  % \space\stamphh:\stampmm
  }

\newcommand{\documentdatezh}{\the\year 年 \the\month 月 \the\day 日}
\newcommand{\documentdatede}{\the\day. \stampmonthde{\the\month} \the\year}
\newcommand{\bverfgetranslationnote}{%
  {\songfont 翻譯說明：初譯使用 Claude Code 與 GPT 協助迭代；仍請以德文原文為準。}%
}
\newcommand{\bverfgecourseinfo}{%
  {\songfont 課程：114 學年度第 2 學期；憲法專題研究（四）；授課老師：湯德宗教授；導讀日期：115 年 6 月 1 日。}%
}
\newcommand{\bverfgeeditorinfo}{%
  {\songfont 編者：王逸帆（R10A21126），國立台灣大學法律系研究所碩士班。}%
}
\newcommand{\bverfgetitleversionline}{%
  \ifbilingualcolumns
    Fassung \documentversion\enspace\textperiodcentered\enspace Stand: \documentdatede
    \quad
    {\songfont 版本 \documentversion\enspace\textperiodcentered\enspace 修訂日期：\documentdatezh\par}%
  \else
    \ifshowtranslation
      {\songfont 版本 \documentversion\enspace\textperiodcentered\enspace 修訂日期：\documentdatezh\par}%
    \else
      Fassung \documentversion\enspace\textperiodcentered\enspace Stand: \documentdatede\par
    \fi
  \fi
}
\newcommand{\bverfgetitlecornerinfo}{%
  \ifshowtranslation
    \vspace{0.72em}%
    \noindent\hfill
    \begin{minipage}{0.66\textwidth}
      \raggedleft\tiny\color{pendingink}\songfont
      \bverfgecourseinfo\par
      \bverfgeeditorinfo
    \end{minipage}%
  \fi
}
\newcommand{\bverfgetitlemeta}{%
  \vfill
  \begin{center}%
    \footnotesize\color{pendingink}%
    \bverfgetitleversionline
    \ifshowtranslation
      \vspace{0.28em}{\scriptsize\bverfgetranslationnote\par}%
    \fi
  \end{center}%
  \bverfgetitlecornerinfo
}

\pagestyle{fancy}
\fancyhf{}
\renewcommand{\headrulewidth}{0pt}
\renewcommand{\footrulewidth}{0pt}
\fancyfoot[C]{\thepage}
\ifbilingualcolumns
  \fancyfoot[R]{\scriptsize\color{pendingink}\timestampzh}%
\else
  \ifshoworiginal
  \else
    \fancyfoot[R]{\scriptsize\color{pendingink}\timestampzh\hspace{0.45cm}}%
  \fi
\fi
\setmainfont{Times New Roman}
\setsansfont{Helvetica Neue}
\setmonofont{Menlo}
\IfFileExists{/Users/iw/Library/Fonts/kaiu.ttf}{
  \setCJKmainfont[Path=/Users/iw/Library/Fonts/,AutoFakeBold=4]{kaiu.ttf}
}{
  \IfFontExistsTF{標楷體}{
    \setCJKmainfont[AutoFakeBold=4]{標楷體}
  }{
    \setCJKmainfont{Songti TC}
  }
}
\IfFontExistsTF{Songti TC}{
  \setCJKfamilyfont{song}{Songti TC}
  \setCJKmonofont{Songti TC}
}{
  \setCJKfamilyfont{song}[Path=/Users/iw/Library/Fonts/,AutoFakeBold=4]{kaiu.ttf}
  \setCJKmonofont[Path=/Users/iw/Library/Fonts/,AutoFakeBold=4]{kaiu.ttf}
}
\newcommand{\songfont}{\CJKfamily{song}}

% ===== 封面／目錄專用打字機字體（僅限封面、目錄頁使用，不影響正文）=====
\IfFontExistsTF{Radio Newsman}{%
  \newfontfamily\bverfgetitlede{Radio Newsman}%
}{%
  \newfontfamily\bverfgetitlede{Courier New}%
}
\IfFontExistsTF{Erikas Farbband}{%
  \newfontfamily\bverfgebodyde{Erikas Farbband}%
}{%
  \let\bverfgebodyde\bverfgetitlede
}
\IfFontExistsTF{Maoken Heavy Labourer}{%
  \setCJKfamilyfont{bverfgetitlecjk}{Maoken Heavy Labourer}%
}{%
  \setCJKfamilyfont{bverfgetitlecjk}{Songti TC}%
}
\IfFontExistsTF{Huiwen-mincho}{%
  \setCJKfamilyfont{bverfgebodycjk}{Huiwen-mincho}%
}{%
  \setCJKfamilyfont{bverfgebodycjk}{Songti TC}%
}
\newcommand{\bverfgetitlezh}{\bverfgetitlede\CJKfamily{bverfgetitlecjk}}
\newcommand{\bverfgebodyzh}{\bverfgebodyde\CJKfamily{bverfgebodycjk}}

\setstretch{1.35}
\setlist{itemsep=0.2em, topsep=0.4em}
\setlength{\parindent}{0pt}
\setlength{\parskip}{0.45em}
\raggedbottom
\setcounter{secnumdepth}{0}
\setcounter{tocdepth}{2}

\newcommand{\lawboxsourcefont}{\footnotesize}
\newcommand{\lawboxtransfont}{\footnotesize}
\newcommand{\lawboxsourcestretch}{1.02}
\newcommand{\lawboxtransstretch}{0.92}

\titleformat{\section}{\centering\Large\bfseries\songfont}{\thesection}{1em}{}
\titleformat{\subsection}{\large\bfseries\songfont}{\thesubsection}{1em}{}
\titleformat{\subsubsection}{\normalsize\bfseries\songfont}{\thesubsubsection}{1em}{}
\titleformat{\paragraph}{\normalsize\bfseries\songfont}{\theparagraph}{1em}{}
\titlespacing*{\section}{0pt}{1.8em}{1.0em}
\titlespacing*{\subsection}{0pt}{1.2em}{0.6em}
\titlespacing*{\subsubsection}{0pt}{1.0em}{0.45em}
\titlespacing*{\paragraph}{0pt}{0.6em}{0.4em}

\definecolor{noteblue}{HTML}{A9C9F5}
\definecolor{noteteal}{HTML}{AEE1D6}
\definecolor{notegray}{HTML}{D7DADF}
\definecolor{caseink}{HTML}{000000}
\definecolor{casepaper}{HTML}{FCFEFD}
\definecolor{sourcepaper}{HTML}{FAFBF8}
\definecolor{sourceline}{HTML}{D7DED8}
\definecolor{sourceink}{HTML}{000000}
\definecolor{pendingink}{HTML}{000000}

\tcbset{
  caseboxbase/.style={
    enhanced,
    breakable,
    width=0.985\linewidth,
    colback=casepaper,
    coltitle=caseink,
    fonttitle=\bfseries\songfont,
    boxrule=0.28pt,
    arc=1.2mm,
    left=2.4mm,
    right=2.4mm,
    top=1.5mm,
    bottom=1.5mm,
    boxsep=1.5mm,
    before skip=0.65em,
    after skip=0.65em,
    before upper={\setlength{\parindent}{0pt}\setlength{\parskip}{0.35em}},
  }
}

\newcommand{\bilingualstart}{}
\newcommand{\bilingualstop}{}
\ifbilingualcolumns
  \renewcommand{\bilingualstart}{\small\setlength{\columnsep}{1.25cm}\setlength{\columnseprule}{0.12pt}\columnratio{0.5,0.5}\multicolumnfootnotes\flushbottom\sloppy\interlinepenalty=80\clubpenalty=800\widowpenalty=800\displaywidowpenalty=800\begin{paracol}{2}\global\intranslationfalse\global\firstbilingualsectiontrue{\footnotesize\color{sourceink}Deutsch\par}\switchcolumn{\footnotesize\songfont\color{sourceink}中文譯文\par}\switchcolumn*}
  \renewcommand{\bilingualstop}{\ifintranslation\switchcolumn*\global\intranslationfalse\fi\end{paracol}\clearpage\raggedbottom}
\fi
\newif\ifrandnumbers
\randnumbersfalse
\newcounter{randnumber}
\newcounter{randnumberlocal}
\newcounter{lastsourcerandstart}
\newif\iflastsourcehasrandnumbers
\lastsourcehasrandnumbersfalse
\newif\ifinlawbox
\inlawboxfalse
\newif\iflawboxhaspendingrn
\lawboxhaspendingrnfalse
\newcommand{\lawboxpendingrn}{}
\newcommand{\startrandnumbers}{\global\randnumberstrue\setcounter{randnumber}{1}\global\lastsourcehasrandnumbersfalse}
\newcommand{\stoprandnumbers}{\global\randnumbersfalse\global\lastsourcehasrandnumbersfalse}
\newlength{\randnumberwidth}
\setlength{\randnumberwidth}{0.95cm}
\newlength{\randnumberrightoffset}
\setlength{\randnumberrightoffset}{0.85cm}
\newlength{\randnumberlawboxrightoffset}
\setlength{\randnumberlawboxrightoffset}{1.40cm}
\newcommand{\randnumberbox}[1]{%
  \leavevmode
  \makebox[0pt][l]{%
    \smash{%
      \tikz[remember picture,overlay,baseline=(rnmark.base)]{%
        \coordinate (rnmark) at (0,0);%
        \ifinlawbox
          \path let \p1=(current page.east), \p2=(rnmark.base) in
            node[
              anchor=base east,
              inner sep=0pt,
              outer sep=0pt,
              text width=\randnumberwidth,
              align=right,
              text=pendingink,
              font=\normalfont\mdseries\scriptsize\ttfamily
            ] at (\x1-\randnumberlawboxrightoffset,\y2) {{\normalfont\mdseries\scriptsize\ttfamily #1}};%
        \else
          \path let \p1=(current page.east), \p2=(rnmark.base) in
            node[
              anchor=base east,
              inner sep=0pt,
              outer sep=0pt,
              text width=\randnumberwidth,
              align=right,
              text=pendingink,
              font=\normalfont\mdseries\scriptsize\ttfamily
            ] at (\x1-\randnumberrightoffset,\y2) {{\normalfont\mdseries\scriptsize\ttfamily #1}};%
        \fi
      }%
    }%
  }%
  \ignorespaces
}
\newcommand{\lawboxstorerandnumber}[1]{%
  \gdef\lawboxpendingrn{#1}%
  \global\lawboxhaspendingrntrue
  \ignorespaces
}
\newcommand{\lawboxuserandnumber}{%
  \iflawboxhaspendingrn
    \randnumberbox{\lawboxpendingrn}%
    \global\lawboxhaspendingrnfalse
  \fi
}
\newcommand{\sourcern}[1]{%
  \ifrandnumbers
    \ifshoworiginal
      \ifshowtranslation\else
        \ifinlawbox\lawboxstorerandnumber{#1}\else\randnumberbox{#1}\fi
      \fi
    \fi
  \fi
}
\newcommand{\transrn}[1]{%
  \ifrandnumbers
    \ifshowtranslation
      \ifinlawbox\lawboxstorerandnumber{#1}\else\randnumberbox{#1}\fi
    \fi
  \fi
}
% ===== 課堂模式：德文原文縮寫註解 =====
% 日常切換只改各判決檔的一行：
%   \classroommodeon        開啟課堂版；註解掉這行即回到一般版。
% 模板預設：
%   1. 課堂模式關閉。
%   2. 課堂模式開啟時，縮寫註解包含「德文全稱＋中文譯名」。
%   3. 課堂模式開啟時，GG/條文編者註仍會自動顯示。
% 進階微調才需要在單案檔覆寫：
%   \classroomabbrzhoff     縮寫註解僅顯示德文全稱。
%   \showeditornotestrue    一般版也顯示編者註。
% 註解策略：同一實際 PDF 頁面中，同一縮寫只在第一次出現處加註。
% 這裡用 zref-abspage 的絕對頁序，不用 \thepage，避免文件內頁碼重置時誤判。
\newif\ifclassroommode
\newif\ifclassroomabbrzh
\classroommodefalse
\classroomabbrzhtrue
\newcommand{\classroommodeon}{\classroommodetrue}
\newcommand{\classroommodeoff}{\classroommodefalse}
\newcommand{\classroomabbrzhon}{\classroomabbrzhtrue}
\newcommand{\classroomabbrzhoff}{\classroomabbrzhfalse}
\newcommand{\abbrnotelabel}{\emph{Abk.}: }
\newcommand{\abbrnote}[3]{%
  \ifclassroommode
    \ifshoworiginal
      \ifintranslation\else
        \footnote{\normalfont\footnotesize\abbrnotelabel #2\ifclassroomabbrzh；{\songfont #3}\fi}%
      \fi
    \fi
  \fi
}
\newcommand{\abbrnoteonce}[4]{%
  #2%
  \ifclassroommode
    \ifshoworiginal
      \ifintranslation\else
        \begingroup
          \edef\abbrcurrentabspage{\number\numexpr\value{abspage}+1\relax}%
          \ifcsname abbrnoteseen@#1@\abbrcurrentabspage\endcsname\else
            \global\expandafter\let\csname abbrnoteseen@#1@\abbrcurrentabspage\endcsname\empty
            \abbrnote{#1}{#3}{#4}%
          \fi
        \endgroup
      \fi
    \fi
  \fi
}
\newcommand{\abbrAbs}{\abbrnoteonce{abs}{Abs.}{Absatz}{項}}
\newcommand{\abbrAAO}{\abbrnoteonce{aao}{a.a.O.}{am angegebenen Ort}{前引處}}
\newcommand{\abbrAE}{\abbrnoteonce{ae}{AE}{Alternativ-Entwurf}{替代草案}}
\newcommand{\abbrArt}{\abbrnoteonce{art}{Art.}{Artikel}{條}}
\newcommand{\abbrAufl}{\abbrnoteonce{aufl}{Aufl.}{Auflage}{版}}
\newcommand{\abbrBd}{\abbrnoteonce{bd}{Bd.}{Band}{卷}}
\newcommand{\abbrBGBl}{\abbrnoteonce{bgbl}{BGBl.}{Bundesgesetzblatt}{聯邦法律公報}}
\newcommand{\abbrBGHSt}{\abbrnoteonce{bghst}{BGHSt}{Entscheidungen des Bundesgerichtshofes in Strafsachen}{聯邦普通法院刑事判決彙編}}
\newcommand{\abbrBRDrucks}{\abbrnoteonce{brdrucks}{BRDrucks.}{Bundesrats-Drucksache}{聯邦參議院文件}}
\newcommand{\abbrBTDrucks}{\abbrnoteonce{btdrucks}{BTDrucks.}{Bundestags-Drucksache}{聯邦議院文件}}
\newcommand{\abbrBundesgesetzbl}{\abbrnoteonce{bundesgesetzbl}{Bundesgesetzbl.}{Bundesgesetzblatt}{聯邦法律公報}}
\newcommand{\abbrBvF}{\abbrnoteonce{bvf}{BvF}{Bundesverfassungsgerichtsverfahren}{聯邦憲法法院程序案號}}
\newcommand{\abbrBvL}{\abbrnoteonce{bvl}{BvL}{Bundesverfassungsgerichtsverfahren}{聯邦憲法法院程序案號}}
\newcommand{\abbrBvR}{\abbrnoteonce{bvr}{BvR}{Bundesverfassungsgerichtsverfahren}{聯邦憲法法院程序案號}}
\newcommand{\abbrBVerfG}{\abbrnoteonce{bverfg}{BVerfG}{Bundesverfassungsgericht}{聯邦憲法法院}}
\newcommand{\abbrBVerfGE}{\abbrnoteonce{bverfge}{BVerfGE}{Entscheidungen des Bundesverfassungsgerichts}{聯邦憲法法院判決集}}
\newcommand{\abbrBVerfGG}{\abbrnoteonce{bverfgg}{BVerfGG}{Bundesverfassungsgerichtsgesetz}{聯邦憲法法院法}}
\newcommand{\abbrDDR}{\abbrnoteonce{ddr}{DDR}{Deutsche Demokratische Republik}{德意志民主共和國；東德}}
\newcommand{\abbrDJT}{\abbrnoteonce{djt}{DJT}{Deutscher Juristentag}{德國法學家大會}}
\newcommand{\abbrDr}{\abbrnoteonce{dr}{Dr.}{Doktor}{Dr.}}
\newcommand{\abbrEuGRZ}{\abbrnoteonce{eugrz}{EuGRZ}{Europäische Grundrechte-Zeitschrift}{歐洲基本權期刊}}
\newcommand{\abbrEV}{\abbrnoteonce{ev}{EV}{Einigungsvertrag}{統一條約}}
\newcommand{\abbrF}{\abbrnoteonce{f}{f.}{folgende}{以下一頁；下一段}}
\newcommand{\abbrFF}{\abbrnoteonce{ff}{ff.}{fortfolgende}{以下數頁；以下各段}}
\newcommand{\abbrGBlDDR}{\abbrnoteonce{gblddr}{GBl. DDR}{Gesetzblatt der Deutschen Demokratischen Republik}{德意志民主共和國法律公報}}
\newcommand{\abbrGG}{\abbrnoteonce{gg}{GG}{Grundgesetz}{基本法}}
\newcommand{\abbrGVBl}{\abbrnoteonce{gvbl}{GVBl.}{Gesetz- und Verordnungsblatt}{法律及命令公報}}
\newcommand{\abbrJR}{\abbrnoteonce{jr}{JR}{Juristische Rundschau}{法學評論}}
\newcommand{\abbrJZ}{\abbrnoteonce{jz}{JZ}{JuristenZeitung}{法學家報}}
\newcommand{\abbrIVm}{\abbrnoteonce{ivm}{i.V.m.}{in Verbindung mit}{結合；連同}}
\newcommand{\abbrMdB}{\abbrnoteonce{mdb}{MdB}{Mitglied des Bundestages}{聯邦議院議員}}
\newcommand{\abbrMwN}{\abbrnoteonce{mwn}{m.w.N.}{mit weiteren Nachweisen}{附進一步文獻／判例出處}}
\newcommand{\abbrNF}{\abbrnoteonce{nf}{n.F.}{neue Fassung}{新版本}}
\newcommand{\abbrAF}{\abbrnoteonce{af}{a.F.}{alte Fassung}{舊版本}}
\newcommand{\abbrNr}{\abbrnoteonce{nr}{Nr.}{Nummer}{款、目或號，依語境}}
\newcommand{\abbrProf}{\abbrnoteonce{prof}{Prof.}{Professor}{教授}}
\newcommand{\abbrRdnr}{\abbrnoteonce{rdnr}{Rdnr.}{Randnummer}{邊碼；段碼}}
\newcommand{\abbrRdnrn}{\abbrnoteonce{rdnrn}{Rdnrn.}{Randnummern}{邊碼；段碼}}
\newcommand{\abbrRGBl}{\abbrnoteonce{rgbl}{RGBl.}{Reichsgesetzblatt}{帝國法律公報}}
\newcommand{\abbrRGSt}{\abbrnoteonce{rgst}{RGSt}{Entscheidungen des Reichsgerichts in Strafsachen}{帝國法院刑事判決彙編}}
\newcommand{\abbrRspr}{\abbrnoteonce{rspr}{Rspr.}{Rechtsprechung}{判例；司法實務}}
\newcommand{\abbrRVO}{\abbrnoteonce{rvo}{RVO}{Reichsversicherungsordnung}{帝國保險法}}
\newcommand{\abbrS}{\abbrnoteonce{s}{S.}{Seite}{頁}}
\newcommand{\abbrSFHG}{\abbrnoteonce{sfhg}{SFHG}{Schwangeren- und Familienhilfegesetz}{孕婦及家庭扶助法}}
\newcommand{\abbrSGBV}{\abbrnoteonce{sgbv}{SGB V}{Fünftes Buch Sozialgesetzbuch}{社會法典第五編}}
\newcommand{\abbrStAG}{\abbrnoteonce{staeg}{StÄG}{Strafrechtsänderungsgesetz}{刑法修正法}}
\newcommand{\abbrStenBer}{\abbrnoteonce{stenber}{StenBer.}{Stenographischer Bericht}{速記紀錄}}
\newcommand{\abbrStGB}{\abbrnoteonce{stgb}{StGB}{Strafgesetzbuch}{刑法}}
\newcommand{\abbrStREG}{\abbrnoteonce{streg}{StREG}{Strafrechtsreform-Ergänzungsgesetz}{刑法改革補充法}}
\newcommand{\abbrStrRG}{\abbrnoteonce{strrg}{StrRG}{Strafrechtsreformgesetz}{刑法改革法}}
\newcommand{\abbrUS}{\abbrnoteonce{us}{U.S.}{United States Reports}{美國最高法院判決彙編}}
\newcommand{\abbrVgl}{\abbrnoteonce{vgl}{vgl.}{vergleiche}{參見}}
\newcommand{\abbrVVDStRL}{\abbrnoteonce{vvdstrl}{VVDStRL}{Veröffentlichungen der Vereinigung der Deutschen Staatsrechtslehrer}{德國公法教師協會論文集}}
\newcommand{\abbrWP}{\abbrnoteonce{wp}{WP.}{Wahlperiode}{屆期}}
\newcommand{\abbrWp}{\abbrnoteonce{wp}{Wp.}{Wahlperiode}{屆期}}
\newcommand{\abbrZStrW}{\abbrnoteonce{zstrw}{ZStrW}{Zeitschrift für die gesamte Strafrechtswissenschaft}{整體刑法學期刊}}
\newcommand{\recordsourcerandstart}{%
  \ifrandnumbers
    \setcounter{lastsourcerandstart}{\value{randnumber}}%
    \global\lastsourcehasrandnumberstrue
  \else
    \global\lastsourcehasrandnumbersfalse
  \fi
}
\newlength{\biparagraphskip}
\setlength{\biparagraphskip}{0.52em}
\newlength{\lawboxblockbeforeskip}
\setlength{\lawboxblockbeforeskip}{0.72em}
\newlength{\lawboxblockafterskip}
\setlength{\lawboxblockafterskip}{0.92em}
\newif\iflawboxpartendedblock
\lawboxpartendedblockfalse
\newcommand{\sourceparstyle}{\footnotesize\color{sourceink}\setstretch{1.12}\RaggedRight\emergencystretch=3em\setlength{\parskip}{0.42em}\obeylines\selectfont}
\newcommand{\transparstyle}{\songfont\color{caseink}\setstretch{1.12}\setlength{\parindent}{0pt}\setlength{\parskip}{0.42em}}
\newcommand{\bilingualpairskip}{%
  \par\addvspace{\biparagraphskip}%
  \switchcolumn
  \par\addvspace{\biparagraphskip}%
  \switchcolumn*%
  \global\intranslationfalse
}
\newcommand{\bilingualboxpairskip}{%
  \par
  \switchcolumn
  \par
  \switchcolumn*%
  \global\intranslationfalse
}
\newcommand{\bilinguallawboxblockskip}{%
  \switchcolumn*%
  \par\addvspace{\lawboxblockafterskip}%
  \switchcolumn
  \par\addvspace{\lawboxblockafterskip}%
  \switchcolumn*%
  \global\intranslationfalse
}
\newcommand{\bikeepwithnext}[1][4]{%
  \ifbilingualcolumns
    \syncleft
    \Needspace{#1\baselineskip}%
    \switchcolumn
    \Needspace{#1\baselineskip}%
    \switchcolumn*%
    \global\intranslationfalse
  \else
    \Needspace{#1\baselineskip}%
  \fi
}
\NewEnviron{sourcepara}[1][]{
  \recordsourcerandstart
  \ifshoworiginal
    \ifbilingualcolumns
      \ifintranslation\switchcolumn*\global\intranslationfalse\fi
    \else
      \Needspace{5\baselineskip}
    \fi
    \ifbilingualcolumns\else\par\addvspace{0.35em}\fi
    {\sourceparstyle
    \noindent\BODY\par}
    \ifbilingualcolumns\else\addvspace{0.25em}\fi
    \ifbilingualcolumns
      \switchcolumn
      \global\intranslationtrue
    \fi
  \else
    \ifrandnumbers
      \begingroup
        \setbox0=\vbox{\hsize=\linewidth\sourceparstyle\noindent\BODY\par}%
      \endgroup
    \fi
  \fi
}
\NewEnviron{transpara}{
  \ifshowtranslation
    \setcounter{randnumberlocal}{\value{lastsourcerandstart}}%
    \ifbilingualcolumns
      \ifintranslation\else\switchcolumn\global\intranslationtrue\fi
      {\transparstyle\setlength{\leftskip}{0.55em}\setlength{\rightskip}{0.15em plus 1em}\addtolength{\linewidth}{-0.55em}\BODY\par}
      \switchcolumn*
      \bilingualpairskip
    \else
      {\transparstyle\BODY\par}
    \fi
  \else
    \ifbilingualcolumns
      \ifintranslation\switchcolumn*\global\intranslationfalse\fi
    \fi
  \fi
}
\NewEnviron{sourceboxpara}[1][]{
  \global\lawboxpartendedblockfalse
  \recordsourcerandstart
  \ifshoworiginal
    \ifbilingualcolumns
      \ifintranslation\switchcolumn*\global\intranslationfalse\fi
    \else
      \Needspace{3\baselineskip}
    \fi
    {\sourceparstyle
    \BODY\par}
    \ifbilingualcolumns\else
      \iflawboxpartendedblock\addvspace{\lawboxblockafterskip}\fi
    \fi
    \ifbilingualcolumns
      \switchcolumn
      \global\intranslationtrue
    \fi
  \fi
}
\NewEnviron{transboxpara}{
  \global\lawboxpartendedblockfalse
  \ifshowtranslation
    \setcounter{randnumberlocal}{\value{lastsourcerandstart}}%
    \ifbilingualcolumns
      \ifintranslation\else\switchcolumn\global\intranslationtrue\fi
      {\transparstyle\BODY\par}
      \iflawboxpartendedblock
        \bilinguallawboxblockskip
      \else
        \switchcolumn*
        \bilingualboxpairskip
      \fi
    \else
      {\transparstyle\BODY\par}
      \iflawboxpartendedblock\addvspace{\lawboxblockafterskip}\fi
    \fi
  \else
    \ifbilingualcolumns
      \ifintranslation\switchcolumn*\global\intranslationfalse\fi
    \fi
  \fi
}
\newcommand{\leitsatznum}[1]{\noindent\hangindent=3.0em\hangafter=1\makebox[2.25em][r]{\bfseries #1.}\hspace{0.55em}\ignorespaces}
\newcommand{\syncleft}{\ifbilingualcolumns\ifintranslation\switchcolumn*\global\intranslationfalse\fi\fi}
\pretocmd{\section}{\syncleft\clearpage}{}{}
\pretocmd{\subsection}{\syncleft}{}{}
\pretocmd{\subsubsection}{\syncleft}{}{}
\newcommand{\biheadingfont}{\songfont}
\newcommand{\bitocentry}[2]{%
  \ifshoworiginal
    \ifshowtranslation
      \ifstrequal{#1}{#2}{#1}{#1\space #2}%
    \else
      #1%
    \fi
  \else
    #2%
  \fi
}
\pdfstringdefDisableCommands{%
  \def\bitocentry#1#2{#1}%
}
\newcommand{\biheading}[7]{%
  \syncleft#1\bikeepwithnext[#3]\phantomsection\addcontentsline{toc}{#2}{\bitocentry{#6}{#7}}%
  \ifbilingualcolumns
    {#4 #6\par}%
    \switchcolumn
    {#4\biheadingfont #7\par}%
    \switchcolumn*%
    \global\intranslationfalse
    \par\addvspace{#5}%
    \switchcolumn
    \par\addvspace{#5}%
    \switchcolumn*%
  \else
    \ifshoworiginal{#4 #6\par}\fi
    \ifshowtranslation{#4\biheadingfont #7\par}\fi
    \addvspace{#5}%
  \fi
}
\newcommand{\termsectionheading}[1]{%
  \par\Needspace{9\baselineskip}%
  \vspace{0.65em}%
  {\songfont\bfseries #1\par}%
  \vspace{0.2em}%
}
% 章節換頁：
%   雙語 paracol 欄內不要用 \clearpage；它會在同步兩欄時吐出只有欄線/頁碼的空頁。
%   \newpage 足以讓下一個 section 起新頁，又不會強制清空所有待排材料。
%   單語模式不在 paracol 內，仍可用 \clearpage。
\newcommand{\bisectionpagebreak}{\ifbilingualcolumns\newpage\else\clearpage\fi}
\newcommand{\bisection}[2]{%
  \biheading{\iffirstbilingualsection\global\firstbilingualsectionfalse\else\bisectionpagebreak\fi}{section}{6}{\centering\Large\bfseries}{0.8em}{#1}{#2}%
}
\newcommand{\bisubsection}[2]{%
  \biheading{}{subsection}{5}{\large\bfseries}{0.55em}{#1}{#2}%
}
\newcommand{\bisubsubsection}[2]{%
  \biheading{}{subsubsection}{4}{\normalsize\bfseries}{0.45em}{#1}{#2}%
}
\newcommand{\bicompositeheading}[4]{%
  \biheading{}{subsection}{5}{\large\bfseries}{0.16em}{#1}{#2}%
  \biheading{}{subsubsection}{3}{\normalsize\bfseries}{0.45em}{#3}{#4}%
}
\newif\iflawboxfirstitem
\lawboxfirstitemfalse
\newcommand{\lawboxprespace}[1]{%
  \par
  \iflawboxfirstitem
    \global\lawboxfirstitemfalse
  \else
    \addvspace{#1}%
  \fi
}
\newtcolorbox{lawboxinner}{
  caseboxbase,
  colframe=sourceline!85!black,
  colback=white,
  left=1.05em,
  right=1.05em,
  top=0.62em,
  bottom=0.62em,
  before upper={\global\inlawboxtrue\global\lawboxfirstitemtrue\global\lawboxhaspendingrnfalse\ifintranslation\lawboxtransfont\else\lawboxsourcefont\fi\RaggedRight\setlength{\parindent}{0pt}\setlength{\parskip}{0pt}\ifintranslation\setstretch{\lawboxtransstretch}\else\setstretch{\lawboxsourcestretch}\fi},
  after upper={\global\lawboxhaspendingrnfalse\global\inlawboxfalse}
}
\tcbset{
  lawboxpartbase/.style={
    enhanced,
    breakable=false,
    width=0.985\linewidth,
    colback=white,
    frame hidden,
    boxrule=0pt,
    arc=0pt,
    left=1.05em,
    right=1.05em,
    boxsep=1.5mm,
    before skip=0pt,
    after skip=0pt,
    before upper={\global\inlawboxtrue\global\lawboxfirstitemtrue\global\lawboxhaspendingrnfalse\ifintranslation\lawboxtransfont\else\lawboxsourcefont\fi\RaggedRight\setlength{\parindent}{0pt}\setlength{\parskip}{0pt}\ifintranslation\setstretch{\lawboxtransstretch}\else\setstretch{\lawboxsourcestretch}\fi},
    after upper={\global\lawboxhaspendingrnfalse\global\inlawboxfalse},
    borderline west={0.28pt}{0pt}{sourceline!85!black},
    borderline east={0.28pt}{0pt}{sourceline!85!black},
  },
  lawboxpartfirst/.style={
    lawboxpartbase,
    top=0.62em,
    bottom=0.04em,
    before skip=\lawboxblockbeforeskip,
    borderline north={0.28pt}{0pt}{sourceline!85!black},
  },
  lawboxpartmiddle/.style={
    lawboxpartbase,
    top=0.04em,
    bottom=0.04em,
  },
  lawboxpartlast/.style={
    lawboxpartbase,
    top=0.04em,
    bottom=0.62em,
    after skip=0pt,
    borderline south={0.28pt}{0pt}{sourceline!85!black},
  }
}
\newtcolorbox{lawboxpartinner}[1]{#1}
\newtcolorbox{termboxinner}{
  caseboxbase,
  colframe=noteblue!70!black,
  colback=casepaper,
  before upper={\songfont\setlength{\parindent}{0pt}\setlength{\parskip}{0.35em}}
}
\newtcolorbox{deboxinner}{
  caseboxbase,
  colframe=notegray!72!black,
  colback=white,
  left=1.05em,
  right=1.05em,
  top=0.62em,
  bottom=0.62em,
  before upper={\global\inlawboxtrue\global\lawboxfirstitemtrue\global\lawboxhaspendingrnfalse\small\setlength{\parindent}{0pt}\setlength{\parskip}{0.35em}},
  after upper={\global\lawboxhaspendingrnfalse\global\inlawboxfalse}
}
\NewEnviron{lawbox}{
  \begin{lawboxinner}\BODY\end{lawboxinner}
}
\NewEnviron{lawboxpart}[2]{
  \begin{lawboxpartinner}{lawboxpart#1,equal height group=#2}\BODY\end{lawboxpartinner}%
  \ifstrequal{#1}{last}{\global\lawboxpartendedblocktrue}{}%
}
\NewEnviron{termbox}{
  \par\addvspace{0.55em}
  {\color{noteblue!70!black}\rule{\linewidth}{0.28pt}}\par
  \begingroup
    \songfont\small\color{caseink}
    \setlength{\parindent}{0pt}
    \setlength{\parskip}{0.24em}
    \BODY
    \par
  \endgroup
  {\color{noteblue!70!black}\rule{\linewidth}{0.28pt}}\par
  \addvspace{0.55em}
}
\NewEnviron{debox}{
  \begin{deboxinner}\BODY\end{deboxinner}
}
\providecommand{\paragraphmark}[1]{}
\newcommand{\lawboxtitleafter}{\addvspace{0.06em}}
\newcommand{\lawtitle}[1]{\lawboxprespace{0.22em}{\lawboxtransfont\bfseries\lawboxuserandnumber #1}\par\lawboxtitleafter}
\newcommand{\absno}[2]{\lawboxprespace{0.04em}\noindent\lawboxuserandnumber\hangindent=2.6em\hangafter=1\makebox[2.6em][l]{(#1)}#2\par}
\newcommand{\nrno}[2]{\lawboxprespace{0pt}\noindent\lawboxuserandnumber\hspace*{2.6em}\hangindent=4.4em\hangafter=1\makebox[1.8em][l]{#1.}#2\par}
\newcommand{\letterno}[2]{\lawboxprespace{0pt}\noindent\lawboxuserandnumber\hspace*{4.4em}\hangindent=6.0em\hangafter=1\makebox[1.6em][l]{#1)}#2\par}
\newcommand{\sourcelawtitle}[1]{\lawboxprespace{0.22em}{\lawboxsourcefont\bfseries\lawboxuserandnumber #1}\par\lawboxtitleafter}
\newcommand{\sourcelawsubtitle}[1]{\lawboxprespace{0.04em}{\lawboxsourcefont\bfseries\lawboxuserandnumber #1}\par\lawboxtitleafter}
\newcommand{\sourceabsno}[2]{\lawboxprespace{0.04em}\noindent\lawboxuserandnumber\hangindent=2.45em\hangafter=1\makebox[2.45em][l]{(#1)}#2\par}
\newcommand{\sourcenrno}[2]{\lawboxprespace{0pt}\noindent\lawboxuserandnumber\hspace*{2.45em}\hangindent=4.25em\hangafter=1\makebox[1.8em][l]{#1.}#2\par}
\newcommand{\sourceletterno}[2]{\lawboxprespace{0pt}\noindent\lawboxuserandnumber\hspace*{4.25em}\hangindent=5.85em\hangafter=1\makebox[1.6em][l]{#1)}#2\par}
\newcommand{\sourcequotepara}[1]{\lawboxprespace{0.04em}\noindent\lawboxuserandnumber\hangindent=1.2em\hangafter=1 #1\par}
\newcommand{\sourcelawline}[1]{\lawboxprespace{0.04em}\noindent\lawboxuserandnumber #1\par}
\newcommand{\lawline}[1]{\lawboxprespace{0.04em}\noindent\lawboxuserandnumber #1\par}
\newcommand{\pendingtranslation}{\par{\songfont\footnotesize\color{pendingink}（中文譯文待補。）}\par}
\newcommand{\tocpart}[1]{\par\addvspace{0.52em}{\footnotesize\bfseries\color{pendingink}#1\par}\addvspace{0.16em}}
\newcommand{\tocdashline}{\par\addvspace{0.48em}\noindent{\color{notegray}\xleaders\hbox to 0.82em{\hfil\rule{0.42em}{0.28pt}\hfil}\hskip\linewidth}\par\addvspace{0.38em}}
% \tocpanel{font-cmd}{content}：將目錄置中於所在欄寬（雙語欄適用）。
% A3 橫式使用 0.58\linewidth，A4 橫式使用 0.84\linewidth，
% 使兩種紙張下目錄欄內容實際寬度相近（均約 100–105 mm）。
\newcommand{\tocpanel}[2]{%
  \makebox[\linewidth][c]{%
    \ifx\bilingualpaper\bilingualpaperaThree
      \begin{minipage}{0.58\linewidth}%
    \else
      \begin{minipage}{0.84\linewidth}%
    \fi
    \toccolstyle
    #1#2%
    \end{minipage}%
  }%
}
% ===== 首頁簡目錄的兩層 description list 縮排 =====
% enumitem 的 description list 可把每一列拆成：
%   [label]  labelsep  body text
% 這裡的「起點」都以所在 minipage/欄位的左緣為 0。
%
% 核心規則：
%   labelindent = label 欄位從左緣開始的位置。
%   labelwidth  = label 欄位本身寬度；標籤太長（如 Abw. Meinung）就加大它。
%   labelsep    = label 欄位右緣到正文的距離。
%   leftmargin  = 正文 body text 的起點。判斷層級視覺時主要看這個。
%
% 所以：
%   想讓「Begründung / 判決理由」整列正文往右或往左：改 tocitems 的 leftmargin。
%   想讓「A. Verfahren und Gegenstand / A. 程序與審查標的」正文往右或往左：
%     改 tocsubitems 的 leftmargin。
%   想讓 A./B./C. 這些標籤本身往右或往左，但正文不動：改 tocsubitems 的 labelindent。
%   想讓 A./B./C. 與正文間距變大或變小：改 tocsubitems 的 labelsep。
%
% 層級原則：
%   tocsubitems.leftmargin 必須大於 tocitems.leftmargin，
%   否則子層級正文會看起來比「Gründe / 理由」還靠左。
\newlist{tocitems}{description}{1}
\setlist[tocitems]{%
  labelindent=0pt,    % 第一層 label 欄位起點；0pt 表示貼齊目錄欄左緣。
  leftmargin=2.54cm,  % 第一層正文起點；例如 Begründung / 判決理由 的左緣。
  labelwidth=2.16cm,  % 第一層 label 欄位寬度；需容納 Leitsätze、Abw. Meinung、不同意見等。
  labelsep=0.32cm,    % 第一層 label 與正文的水平間距；只改間距，不改正文起點。
  itemsep=0.28em,     % 第一層各項之間的垂直距離。
  topsep=0.20em,      % 第一層 list 上下與前後內容的垂直距離。
  parsep=0pt,         % 同一 item 內段落之間的距離；目錄項通常不需要。
  partopsep=0pt,      % list 若緊接段落開始時的額外距離；關閉以免目錄高度飄動。
  align=left,         % label 欄內左對齊；長短標籤不靠右貼齊。
  font=\normalfont\bfseries % label 的字型；正文不受這行影響。
}
\newlist{tocsubitems}{description}{1}
\setlist[tocsubitems]{%
  labelindent=1cm,    % 第二層 label 起點；只控制 A./B./C. 本身的位置。
  leftmargin=3.00cm,  % 第二層正文起點；必須大於 2.54cm，才會比 Begründung / 判決理由更右。
  labelwidth=0.5cm,   % 第二層 label 欄寬；A./B./C. 很短，可小於第一層。
  labelsep=1.5cm,    % 第二層 label 與正文間距；A. 與 Verfahren / 程序 的距離看這裡。
  itemsep=0.12em,     % 第二層各項較密，避免 A.–E. 佔太高。
  topsep=0.04em,      % 第二層 list 與 Gründe / 理由之間的垂直距離。
  parsep=0pt,         % 同一第二層 item 內段落距離；目錄項通常不需要。
  partopsep=0pt,      % 關閉額外段落起始距離，避免版面不穩。
  align=left,         % A./B./C. label 欄內左對齊。
  font=\normalfont\bfseries, % A./B./C. label 加粗；正文不受這行影響。
  before={}           % 不在第二層 list 前自動插入額外內容。
}
\newlist{termitems}{description}{1}
\ifbilingualcolumns
  \setlist[termitems]{%
    leftmargin=5.2cm,
    labelwidth=4.8cm,
    labelsep=0.35cm,
    itemsep=0.16em,
    topsep=0.2em,
    parsep=0pt,
    partopsep=0pt,
    font=\normalfont\bfseries,
    style=nextline
  }
\else
  \setlist[termitems]{%
    leftmargin=0pt,
    labelwidth=0pt,
    labelsep=0pt,
    itemsep=0.24em,
    topsep=0.2em,
    parsep=0pt,
    partopsep=0pt,
    font=\normalfont\bfseries,
    style=nextline
  }
\fi
\newlist{abbritems}{description}{1}
\setlist[abbritems]{%
  leftmargin=2.38cm,
  labelwidth=2.02cm,
  labelsep=0.28cm,
  itemsep=0.16em,
  topsep=0.2em,
  parsep=0pt,
  partopsep=0pt,
  align=left,
  font=\normalfont\bfseries
}
\ifbilingualcolumns
  \newenvironment{termcolumns}{\begin{multicols}{2}}{\end{multicols}}
\else
  \newenvironment{termcolumns}{}{}
\fi

% ===== 編者註腳開關 =====
% 一般版預設不顯示編者註，避免正文腳註過密。
% 若開啟 \classroommodeon，GG/條文編者註仍會自動顯示，無需另開本開關。
% 若一般版也要顯示編者註，可在單案檔 \input 後加 \showeditornotestrue。
\newif\ifshoweditornotes
\showeditornotesfalse

% ===== 目錄排版輔助 =====
\newcommand{\toccolstyle}{\raggedright\arraybackslash}
\newcommand{\tocsingleindent}{\leftskip=1.45cm\rightskip=1.2cm}

% ===== 行內引文環境（德文原文與中文譯文各一，帶左側灰色邊線）=====
\newcommand{\quoteparbreak}{\par\addvspace{0.48em}}
\newcommand{\sourcequoteinner}[1]{%
  \begin{tcolorbox}[
    enhanced, breakable, width=\dimexpr\linewidth-0.8em\relax, frame hidden,
    colback=white, coltext=caseink, boxrule=0pt,
    borderline west={0.55pt}{0.88em}{notegray},
    left=1.78em, right=0.72em, top=0.18em, bottom=0.18em,
    boxsep=0pt, before skip=0.24em, after skip=0.24em,
    before upper={\normalfont\footnotesize\color{caseink}\setlength{\parindent}{0pt}\setlength{\parskip}{0.34em}\raggedright\setlength{\rightskip}{0pt plus 1em}}
  ]%
  #1\par
  \end{tcolorbox}%
}
\newcommand{\transquoteinner}[1]{%
  \begin{tcolorbox}[
    enhanced, breakable, width=\dimexpr\linewidth-0.8em\relax, frame hidden,
    colback=white, coltext=caseink, boxrule=0pt,
    borderline west={0.55pt}{0.88em}{notegray},
    left=1.78em, right=0.72em, top=0.18em, bottom=0.18em,
    boxsep=0pt, before skip=0.24em, after skip=0.24em,
    before upper={\songfont\small\color{caseink}\setlength{\parindent}{0pt}\setlength{\parskip}{0.34em}\raggedright\setlength{\rightskip}{0pt plus 1em}}
  ]%
  #1\par
  \end{tcolorbox}%
}

% ===== 大綱（Gliederung）Rn 輔助命令 =====
\newcommand{\outrn}[2]{\enspace{\normalfont\footnotesize\color{pendingink}(Rn\,#1–#2)}}
\newcommand{\outrnsingle}[1]{\enspace{\normalfont\footnotesize\color{pendingink}(Rn\,#1)}}
\newcommand{\outrnzh}[2]{\enspace{\normalfont\footnotesize\color{pendingink}（第\,#1–#2\,段）}}
\newcommand{\outrnzhs}[1]{\enspace{\normalfont\footnotesize\color{pendingink}（第\,#1\,段）}}
% ===== 大綱雙語對照表格輔助命令（三欄：段碼 │ 德文 │ 中文）=====
% 欄 1：段碼（由欄定義自動格式化：小字、等寬、右對齊、pendingink）
% 欄 2：德文標籤＋正文
% 欄 3：中文標籤＋正文
%
% 無段碼的列（\omainrow、\osecrow），欄 1 以 & 空置。
% 有段碼的列（\osecrowrn、\osubrow、\ointrow），#1 為預格式化字串
%   例：\osubrow{2–49}{I.}{...}{I.}{...}
%       \ointrow{107}{...}{...}
%
% \opartrow：段組標題（橫跨三欄）
\newcommand{\opartrow}[2]{%
  \noalign{\vspace{0.30em}}%
  \multicolumn{3}{@{}l@{}}{%
    \footnotesize\bfseries\color{pendingink}#1\hfill\songfont#2%
  }\\[0.06em]%
}
% \odashrow：虛線分隔（橫跨三欄）
\newcommand{\odashrow}{%
  \noalign{\vspace{0.28em}}%
  \multicolumn{3}{@{}l@{}}{\color{notegray}%
    \xleaders\hbox to 0.82em{\hfil\rule{0.42em}{0.28pt}\hfil}\hskip 0.88\linewidth}%
  \\[0.24em]%
}
% \odissentpart：不同意見書區段標頭。比一般 \opartrow 更像附錄性轉場。
\newcommand{\odissentpart}[2]{%
  \noalign{\vspace{0.42em}}%
  \multicolumn{3}{@{}p{0.88\textwidth}@{}}{%
    \color{notegray}\rule{\linewidth}{0.18pt}%
  }\\[-0.18em]%
  \multicolumn{3}{@{}p{0.88\textwidth}@{}}{%
    \makebox[\linewidth][s]{%
      \footnotesize\bfseries\color{pendingink}#1%
      \hfill{\songfont #2}%
    }%
  }\\[-0.02em]%
}
% \odissentopinion：不同意見書的署名與一行摘要，右欄保留段碼範圍。
\newcommand{\odissentopinion}[5]{%
  \footnotesize\hangindent=0.52cm\hangafter=1%
    {\bfseries\color{caseink}#2}\enspace{\color{notegray}\textperiodcentered}\enspace\color{caseink}#3%
  &\footnotesize\hangindent=0.52cm\hangafter=1%
    {\songfont\bfseries\color{caseink}#4}\enspace{\color{notegray}\textperiodcentered}\enspace\songfont\color{caseink}#5%
  &#1%
  \\[0.12em]%
}
% \omainrow：頂層項目，無段碼（Leitsätze、Urteil、Gründe …）
\newcommand{\omainrow}[4]{%
  \footnotesize\hangindent=1.92cm\hangafter=1%
    \makebox[1.92cm][l]{\bfseries\color{caseink}#1}\color{caseink}#2%
  &\footnotesize\hangindent=1.92cm\hangafter=1%
    \makebox[1.92cm][l]{\bfseries\color{caseink}#3}\songfont\color{caseink}#4%
  &%
  \\[0.15em]%
}
% \omainrowrn：頂層項目，有段碼（不同意見等標籤寬於 \osecrow 框者）
\newcommand{\omainrowrn}[5]{%
  \footnotesize\hangindent=1.92cm\hangafter=1%
    \makebox[1.92cm][l]{\bfseries\color{caseink}#2}\color{caseink}#3%
  &\footnotesize\hangindent=1.92cm\hangafter=1%
    \makebox[1.92cm][l]{\bfseries\color{caseink}#4}\songfont\color{caseink}#5%
  &#1%
  \\[0.15em]%
}
% \osecrow：一級子項，無段碼（A.、C.、D. 等有子項者）
\newcommand{\osecrow}[4]{%
  \footnotesize\hspace{1.10cm}\hangindent=1.98cm\hangafter=1%
    \makebox[0.86cm][l]{\bfseries\color{caseink}#1}\color{caseink}#2%
  &\footnotesize\hspace{1.10cm}\hangindent=1.98cm\hangafter=1%
    \makebox[0.86cm][l]{\bfseries\color{caseink}#3}\songfont\color{caseink}#4%
  &%
  \\[0.11em]%
}
% \osecrowrn：一級子項，有段碼（B.、E. 等完整段落範圍者）
% #1=段碼字串（如 101–106）, #2=DE標籤, #3=DE正文, #4=ZH標籤, #5=ZH正文
\newcommand{\osecrowrn}[5]{%
  \footnotesize\hspace{1.10cm}\hangindent=1.98cm\hangafter=1%
    \makebox[0.86cm][l]{\bfseries\color{caseink}#2}\color{caseink}#3%
  &\footnotesize\hspace{1.10cm}\hangindent=1.98cm\hangafter=1%
    \makebox[0.86cm][l]{\bfseries\color{caseink}#4}\songfont\color{caseink}#5%
  &#1%
  \\[0.11em]%
}
% \osubrow：二級子項，有段碼範圍（I.、II.、A.I. …）
% #1=段碼字串, #2=DE標籤, #3=DE正文, #4=ZH標籤, #5=ZH正文
\newcommand{\osubrow}[5]{%
  \footnotesize\hspace{1.40cm}\hangindent=2.30cm\hangafter=1%
    \makebox[0.90cm][l]{\bfseries\color{caseink}#2}\color{caseink}#3%
  &\footnotesize\hspace{1.40cm}\hangindent=2.30cm\hangafter=1%
    \makebox[0.90cm][l]{\bfseries\color{caseink}#4}\songfont\color{caseink}#5%
  &#1%
  \\[0.09em]%
}
% \ointrow：引入段落，有單一段碼，無標籤
% #1=段碼字串, #2=DE正文, #3=ZH正文
\newcommand{\ointrow}[3]{%
  \footnotesize\hspace{0.52cm}\hangindent=0.52cm\hangafter=1\color{caseink}#2%
  &\footnotesize\hspace{0.52cm}\hangindent=0.52cm\hangafter=1\songfont\color{caseink}#3%
  &#1%
  \\[0.09em]%
}
% \osubsubrow：三級子項（1.、2.、3.）
% #1=段碼字串, #2=DE標籤, #3=DE正文, #4=ZH標籤, #5=ZH正文
\newcommand{\osubsubrow}[5]{%
  \footnotesize\hspace{1.80cm}\hangindent=2.48cm\hangafter=1%
    \makebox[0.68cm][l]{\bfseries\color{caseink}#2}\color{caseink}#3%
  &\footnotesize\hspace{1.80cm}\hangindent=2.48cm\hangafter=1%
    \makebox[0.68cm][l]{\bfseries\color{caseink}#4}\songfont\color{caseink}#5%
  &#1%
  \\[0.08em]%
}
% \osubsubsubrow：四級子項（a）、b））
% #1=段碼字串, #2=DE標籤, #3=DE正文, #4=ZH標籤, #5=ZH正文
\newcommand{\osubsubsubrow}[5]{%
  \footnotesize\hspace{2.10cm}\hangindent=2.68cm\hangafter=1%
    \makebox[0.58cm][l]{\bfseries\color{caseink}#2}\color{caseink}#3%
  &\footnotesize\hspace{2.10cm}\hangindent=2.68cm\hangafter=1%
    \makebox[0.58cm][l]{\bfseries\color{caseink}#4}\songfont\color{caseink}#5%
  &#1%
  \\[0.07em]%
}
% \osubsubsubsubrow：五級子項（aa）、bb））
% 標籤框 0.62cm：足以容納「aa)」在 footnotesize bold 下（約 0.50cm）並留有間距
% #1=段碼字串, #2=DE標籤, #3=DE正文, #4=ZH標籤, #5=ZH正文
\newcommand{\osubsubsubsubrow}[5]{%
  \footnotesize\hspace{2.38cm}\hangindent=3.06cm\hangafter=1%
    \makebox[0.68cm][l]{\bfseries\color{caseink}#2}\color{caseink}#3%
  &\footnotesize\hspace{2.38cm}\hangindent=3.06cm\hangafter=1%
    \makebox[0.68cm][l]{\bfseries\color{caseink}#4}\songfont\color{caseink}#5%
  &#1%
  \\[0.06em]%
}
% \outlinepage：雙語對照大綱頁包裝環境（longtable 自動跨頁）
% 三欄：德文欄 + 中文欄 + 段碼欄（最右，不折行）
% 段碼欄寬度：1.80cm，足以容納最長段碼如「309–348」
% DE/ZH 欄寬：(0.87\textwidth - 2.08cm) / 2
%   驗算：2×欄 + 0.05\textwidth + 0.28cm gap + 1.80cm = 0.87tw - 2.08cm + 0.05tw + 2.08cm = 0.92tw ✓
\newenvironment{outlinepage}{%
  \vspace*{0.40cm}%
  \setlength{\LTleft}{0.04\textwidth}%
  \setlength{\LTright}{0.04\textwidth}%
  \setlength{\tabcolsep}{0pt}%
  % 大綱列距由 longtable 的 \arraystretch 與各 row macro 結尾的 \\[...em] 共同決定。
  % 這裡控制「每一列本身的表格 strut 高度」；數值越大，A./I./1. 各項之間越像有空行。
  % A3 原本用 0.62，視覺上沒有空行；A4 也沿用同值，避免切到 A4 後項目被撐開。
  % 若日後覺得大綱太擠，只微調這個值（例如 0.68），不要改各 \omainrow/\osubrow 的縮排。
  \renewcommand{\arraystretch}{0.62}%
  \begin{longtable}{%
    >{\vspace{0pt}\raggedright\arraybackslash}p{\dimexpr(0.87\textwidth-2.08cm)/2\relax}%
    @{\hspace{0.025\textwidth}}%
    !{\color{notegray}\vrule width 0.12pt}%
    @{\hspace{0.025\textwidth}}%
    >{\vspace{0pt}\raggedright\arraybackslash}p{\dimexpr(0.87\textwidth-2.08cm)/2\relax}%
    @{\hspace{0.28cm}}%
    >{\vspace{0pt}\footnotesize\normalfont\color{pendingink}\raggedright\arraybackslash}p{1.80cm}%
    @{}%
  }%
  % 首頁欄標籤：不要橫跨置中；分別放在德文欄與中文欄的實際欄位起點。
  \footnotesize\bfseries\color{sourceink}Gliederung%
  &\footnotesize\bfseries\songfont\color{sourceink}大綱%
  &%
  \\[0.36em]%
  \endfirsthead%
  % 續頁欄標籤：同樣分欄定位，以灰色細體區分；無需標注「續」。
  \footnotesize\color{pendingink}Gliederung%
  &\footnotesize\songfont\color{pendingink}大綱%
  &%
  \\[0.24em]%
  \endhead%
}{%
  \end{longtable}%
}
 載入。
% 不含：\documentclass、\documentversion、\ggversionde/zh、\tocde/\toczh。
% 日期與首頁版本列由本檔自動產生；各文件只需手動指定 \documentversion。

\usepackage{xeCJK}
\usepackage{fontspec}
\usepackage{geometry}
\usepackage{setspace}
\usepackage{hyperref}
\usepackage{xurl}
\usepackage{bookmark}
\usepackage{enumitem}
\usepackage{xcolor}
\usepackage{titlesec}
\usepackage{array}
\usepackage{longtable}
\usepackage{tcolorbox}
\usepackage{environ}
\usepackage{pdflscape}
\usepackage{etoolbox}
\usepackage{ragged2e}
\usepackage{paracol}
\usepackage{multicol}
\usepackage{needspace}
\usepackage{fancyhdr}
\usepackage{tikz}
\usepackage{zref-abspage}
\tcbuselibrary{breakable,skins}
\usetikzlibrary{calc}
\usepackage{pgfcalendar}

\hypersetup{
  unicode=true,
  bookmarksopen=true,
  bookmarksnumbered=false,
  hidelinks
}

% ===== 自動推導，不要手動改下面 =====
% （\docmode 與 \bilingualpaper 由各文件在 % bverfge_layout.tex
% 共用排版基礎設施，供 Schwangerschaftsabbruch I 與 II 共享。
% 由各文件以 \documentclass 後 % bverfge_layout.tex
% 共用排版基礎設施，供 Schwangerschaftsabbruch I 與 II 共享。
% 由各文件以 \documentclass 後 \input{../../共用LaTeX/bverfge_layout} 載入。
% 不含：\documentclass、\documentversion、\ggversionde/zh、\tocde/\toczh。
% 日期與首頁版本列由本檔自動產生；各文件只需手動指定 \documentversion。

\usepackage{xeCJK}
\usepackage{fontspec}
\usepackage{geometry}
\usepackage{setspace}
\usepackage{hyperref}
\usepackage{xurl}
\usepackage{bookmark}
\usepackage{enumitem}
\usepackage{xcolor}
\usepackage{titlesec}
\usepackage{array}
\usepackage{longtable}
\usepackage{tcolorbox}
\usepackage{environ}
\usepackage{pdflscape}
\usepackage{etoolbox}
\usepackage{ragged2e}
\usepackage{paracol}
\usepackage{multicol}
\usepackage{needspace}
\usepackage{fancyhdr}
\usepackage{tikz}
\usepackage{zref-abspage}
\tcbuselibrary{breakable,skins}
\usetikzlibrary{calc}
\usepackage{pgfcalendar}

\hypersetup{
  unicode=true,
  bookmarksopen=true,
  bookmarksnumbered=false,
  hidelinks
}

% ===== 自動推導，不要手動改下面 =====
% （\docmode 與 \bilingualpaper 由各文件在 \input{bverfge_layout} 之前設定；
%   預設 chinese / a4）
\ifdefined\docmode\else\def\docmode{chinese}\fi
\ifdefined\bilingualpaper\else\def\bilingualpaper{a4}\fi
\newif\ifshoworiginal
\newif\ifshowtranslation
\newif\ifbilinguallandscape
\newif\ifbilingualcolumns
\newif\ifintranslation
\newif\iffirstbilingualsection

\showoriginalfalse
\showtranslationfalse
\bilinguallandscapefalse
\bilingualcolumnsfalse
\intranslationfalse
\firstbilingualsectionfalse

\def\modechinese{chinese}
\def\modegerman{german}
\def\modebilingual{bilingual}
\def\bilingualpaperaThree{a3}
\def\bilingualpaperaFour{a4}

\ifx\docmode\modechinese
  \showtranslationtrue
\fi

\ifx\docmode\modegerman
  \showoriginaltrue
\fi

\ifx\docmode\modebilingual
  \showoriginaltrue
  \showtranslationtrue
  \bilinguallandscapetrue
  \bilingualcolumnstrue
\fi

% 編排模式說明：
% chinese  → \showoriginalfalse  \showtranslationtrue  （A4 直式，純中文）
% german   → \showoriginaltrue   \showtranslationfalse （A4 直式，純德文）
% bilingual→ \showoriginaltrue   \showtranslationtrue  \bilingualcolumnstrue（A3/A4 橫式對照）

\ifbilingualcolumns
  \ifx\bilingualpaper\bilingualpaperaFour
    \geometry{
      a4paper,
      landscape,
      left=1.25cm,
      right=2.25cm,
      top=1.25cm,
      bottom=1.75cm,
      footskip=8.5mm,
      marginparwidth=0.75cm,
      marginparsep=0.25cm
    }
  \else
    \geometry{
      a3paper,
      landscape,
      left=1.55cm,
      right=2.75cm,
      top=1.55cm,
      bottom=2.05cm,
      footskip=9.5mm,
      marginparwidth=0.9cm,
      marginparsep=0.35cm
    }
  \fi
\else
\geometry{
  a4paper,
  portrait,
  left=2.2cm,
  right=2.6cm,
  top=2.0cm,
  bottom=2.0cm,
  marginparwidth=0.8cm,
  marginparsep=0.3cm
}
\fi
% ===== 時間戳基礎設施 =====
\newcount\stamptimeval
\newcount\stampminuteval
\newcommand{\stamphh}{%
  \stamptimeval=\time
  \divide\stamptimeval by 60
  \ifnum\stamptimeval<10 0\fi
  \the\stamptimeval}
\newcommand{\stampmm}{%
  \stamptimeval=\time
  \stampminuteval=\time
  \divide\stampminuteval by 60
  \multiply\stampminuteval by 60
  \advance\stamptimeval by -\stampminuteval
  \ifnum\stamptimeval<10 0\fi
  \the\stamptimeval}
\newcommand{\stampwkde}[1]{%
  \ifcase#1Montag\or Dienstag\or Mittwoch\or Donnerstag\or Freitag\or Samstag\or Sonntag\fi}
\newcommand{\stampwkzh}[1]{%
  \ifcase#1 一\or 二\or 三\or 四\or 五\or 六\or 日\fi}
\newcommand{\stampmonthde}[1]{%
  \ifcase#1\or Januar\or Februar\or März\or April\or Mai\or Juni\or
  Juli\or August\or September\or Oktober\or November\or Dezember\fi}
\newcount\stampjdval
\newcount\stampwdval
\pgfcalendardatetojulian{\the\year-\the\month-\the\day}{\stampjdval}
\pgfcalendarjuliantoweekday{\the\stampjdval}{\stampwdval}
\newcommand{\timestampzh}{%
  \songfont 最後修改：\the\year 年\the\month 月\the\day 日%
  % （星期\stampwkzh{\the\stampwdval}）%
  % \space\stamphh:\stampmm
  }

\newcommand{\documentdatezh}{\the\year 年 \the\month 月 \the\day 日}
\newcommand{\documentdatede}{\the\day. \stampmonthde{\the\month} \the\year}
\newcommand{\bverfgetranslationnote}{%
  {\songfont 翻譯說明：初譯使用 Claude Code 與 GPT 協助迭代；仍請以德文原文為準。}%
}
\newcommand{\bverfgecourseinfo}{%
  {\songfont 課程：114 學年度第 2 學期；憲法專題研究（四）；授課老師：湯德宗教授；導讀日期：115 年 6 月 1 日。}%
}
\newcommand{\bverfgeeditorinfo}{%
  {\songfont 編者：王逸帆（R10A21126），國立台灣大學法律系研究所碩士班。}%
}
\newcommand{\bverfgetitleversionline}{%
  \ifbilingualcolumns
    Fassung \documentversion\enspace\textperiodcentered\enspace Stand: \documentdatede
    \quad
    {\songfont 版本 \documentversion\enspace\textperiodcentered\enspace 修訂日期：\documentdatezh\par}%
  \else
    \ifshowtranslation
      {\songfont 版本 \documentversion\enspace\textperiodcentered\enspace 修訂日期：\documentdatezh\par}%
    \else
      Fassung \documentversion\enspace\textperiodcentered\enspace Stand: \documentdatede\par
    \fi
  \fi
}
\newcommand{\bverfgetitlecornerinfo}{%
  \ifshowtranslation
    \vspace{0.72em}%
    \noindent\hfill
    \begin{minipage}{0.66\textwidth}
      \raggedleft\tiny\color{pendingink}\songfont
      \bverfgecourseinfo\par
      \bverfgeeditorinfo
    \end{minipage}%
  \fi
}
\newcommand{\bverfgetitlemeta}{%
  \vfill
  \begin{center}%
    \footnotesize\color{pendingink}%
    \bverfgetitleversionline
    \ifshowtranslation
      \vspace{0.28em}{\scriptsize\bverfgetranslationnote\par}%
    \fi
  \end{center}%
  \bverfgetitlecornerinfo
}

\pagestyle{fancy}
\fancyhf{}
\renewcommand{\headrulewidth}{0pt}
\renewcommand{\footrulewidth}{0pt}
\fancyfoot[C]{\thepage}
\ifbilingualcolumns
  \fancyfoot[R]{\scriptsize\color{pendingink}\timestampzh}%
\else
  \ifshoworiginal
  \else
    \fancyfoot[R]{\scriptsize\color{pendingink}\timestampzh\hspace{0.45cm}}%
  \fi
\fi
\setmainfont{Times New Roman}
\setsansfont{Helvetica Neue}
\setmonofont{Menlo}
\IfFileExists{/Users/iw/Library/Fonts/kaiu.ttf}{
  \setCJKmainfont[Path=/Users/iw/Library/Fonts/,AutoFakeBold=4]{kaiu.ttf}
}{
  \IfFontExistsTF{標楷體}{
    \setCJKmainfont[AutoFakeBold=4]{標楷體}
  }{
    \setCJKmainfont{Songti TC}
  }
}
\IfFontExistsTF{Songti TC}{
  \setCJKfamilyfont{song}{Songti TC}
  \setCJKmonofont{Songti TC}
}{
  \setCJKfamilyfont{song}[Path=/Users/iw/Library/Fonts/,AutoFakeBold=4]{kaiu.ttf}
  \setCJKmonofont[Path=/Users/iw/Library/Fonts/,AutoFakeBold=4]{kaiu.ttf}
}
\newcommand{\songfont}{\CJKfamily{song}}

% ===== 封面／目錄專用打字機字體（僅限封面、目錄頁使用，不影響正文）=====
\IfFontExistsTF{Radio Newsman}{%
  \newfontfamily\bverfgetitlede{Radio Newsman}%
}{%
  \newfontfamily\bverfgetitlede{Courier New}%
}
\IfFontExistsTF{Erikas Farbband}{%
  \newfontfamily\bverfgebodyde{Erikas Farbband}%
}{%
  \let\bverfgebodyde\bverfgetitlede
}
\IfFontExistsTF{Maoken Heavy Labourer}{%
  \setCJKfamilyfont{bverfgetitlecjk}{Maoken Heavy Labourer}%
}{%
  \setCJKfamilyfont{bverfgetitlecjk}{Songti TC}%
}
\IfFontExistsTF{Huiwen-mincho}{%
  \setCJKfamilyfont{bverfgebodycjk}{Huiwen-mincho}%
}{%
  \setCJKfamilyfont{bverfgebodycjk}{Songti TC}%
}
\newcommand{\bverfgetitlezh}{\bverfgetitlede\CJKfamily{bverfgetitlecjk}}
\newcommand{\bverfgebodyzh}{\bverfgebodyde\CJKfamily{bverfgebodycjk}}

\setstretch{1.35}
\setlist{itemsep=0.2em, topsep=0.4em}
\setlength{\parindent}{0pt}
\setlength{\parskip}{0.45em}
\raggedbottom
\setcounter{secnumdepth}{0}
\setcounter{tocdepth}{2}

\newcommand{\lawboxsourcefont}{\footnotesize}
\newcommand{\lawboxtransfont}{\footnotesize}
\newcommand{\lawboxsourcestretch}{1.02}
\newcommand{\lawboxtransstretch}{0.92}

\titleformat{\section}{\centering\Large\bfseries\songfont}{\thesection}{1em}{}
\titleformat{\subsection}{\large\bfseries\songfont}{\thesubsection}{1em}{}
\titleformat{\subsubsection}{\normalsize\bfseries\songfont}{\thesubsubsection}{1em}{}
\titleformat{\paragraph}{\normalsize\bfseries\songfont}{\theparagraph}{1em}{}
\titlespacing*{\section}{0pt}{1.8em}{1.0em}
\titlespacing*{\subsection}{0pt}{1.2em}{0.6em}
\titlespacing*{\subsubsection}{0pt}{1.0em}{0.45em}
\titlespacing*{\paragraph}{0pt}{0.6em}{0.4em}

\definecolor{noteblue}{HTML}{A9C9F5}
\definecolor{noteteal}{HTML}{AEE1D6}
\definecolor{notegray}{HTML}{D7DADF}
\definecolor{caseink}{HTML}{000000}
\definecolor{casepaper}{HTML}{FCFEFD}
\definecolor{sourcepaper}{HTML}{FAFBF8}
\definecolor{sourceline}{HTML}{D7DED8}
\definecolor{sourceink}{HTML}{000000}
\definecolor{pendingink}{HTML}{000000}

\tcbset{
  caseboxbase/.style={
    enhanced,
    breakable,
    width=0.985\linewidth,
    colback=casepaper,
    coltitle=caseink,
    fonttitle=\bfseries\songfont,
    boxrule=0.28pt,
    arc=1.2mm,
    left=2.4mm,
    right=2.4mm,
    top=1.5mm,
    bottom=1.5mm,
    boxsep=1.5mm,
    before skip=0.65em,
    after skip=0.65em,
    before upper={\setlength{\parindent}{0pt}\setlength{\parskip}{0.35em}},
  }
}

\newcommand{\bilingualstart}{}
\newcommand{\bilingualstop}{}
\ifbilingualcolumns
  \renewcommand{\bilingualstart}{\small\setlength{\columnsep}{1.25cm}\setlength{\columnseprule}{0.12pt}\columnratio{0.5,0.5}\multicolumnfootnotes\flushbottom\sloppy\interlinepenalty=80\clubpenalty=800\widowpenalty=800\displaywidowpenalty=800\begin{paracol}{2}\global\intranslationfalse\global\firstbilingualsectiontrue{\footnotesize\color{sourceink}Deutsch\par}\switchcolumn{\footnotesize\songfont\color{sourceink}中文譯文\par}\switchcolumn*}
  \renewcommand{\bilingualstop}{\ifintranslation\switchcolumn*\global\intranslationfalse\fi\end{paracol}\clearpage\raggedbottom}
\fi
\newif\ifrandnumbers
\randnumbersfalse
\newcounter{randnumber}
\newcounter{randnumberlocal}
\newcounter{lastsourcerandstart}
\newif\iflastsourcehasrandnumbers
\lastsourcehasrandnumbersfalse
\newif\ifinlawbox
\inlawboxfalse
\newif\iflawboxhaspendingrn
\lawboxhaspendingrnfalse
\newcommand{\lawboxpendingrn}{}
\newcommand{\startrandnumbers}{\global\randnumberstrue\setcounter{randnumber}{1}\global\lastsourcehasrandnumbersfalse}
\newcommand{\stoprandnumbers}{\global\randnumbersfalse\global\lastsourcehasrandnumbersfalse}
\newlength{\randnumberwidth}
\setlength{\randnumberwidth}{0.95cm}
\newlength{\randnumberrightoffset}
\setlength{\randnumberrightoffset}{0.85cm}
\newlength{\randnumberlawboxrightoffset}
\setlength{\randnumberlawboxrightoffset}{1.40cm}
\newcommand{\randnumberbox}[1]{%
  \leavevmode
  \makebox[0pt][l]{%
    \smash{%
      \tikz[remember picture,overlay,baseline=(rnmark.base)]{%
        \coordinate (rnmark) at (0,0);%
        \ifinlawbox
          \path let \p1=(current page.east), \p2=(rnmark.base) in
            node[
              anchor=base east,
              inner sep=0pt,
              outer sep=0pt,
              text width=\randnumberwidth,
              align=right,
              text=pendingink,
              font=\normalfont\mdseries\scriptsize\ttfamily
            ] at (\x1-\randnumberlawboxrightoffset,\y2) {{\normalfont\mdseries\scriptsize\ttfamily #1}};%
        \else
          \path let \p1=(current page.east), \p2=(rnmark.base) in
            node[
              anchor=base east,
              inner sep=0pt,
              outer sep=0pt,
              text width=\randnumberwidth,
              align=right,
              text=pendingink,
              font=\normalfont\mdseries\scriptsize\ttfamily
            ] at (\x1-\randnumberrightoffset,\y2) {{\normalfont\mdseries\scriptsize\ttfamily #1}};%
        \fi
      }%
    }%
  }%
  \ignorespaces
}
\newcommand{\lawboxstorerandnumber}[1]{%
  \gdef\lawboxpendingrn{#1}%
  \global\lawboxhaspendingrntrue
  \ignorespaces
}
\newcommand{\lawboxuserandnumber}{%
  \iflawboxhaspendingrn
    \randnumberbox{\lawboxpendingrn}%
    \global\lawboxhaspendingrnfalse
  \fi
}
\newcommand{\sourcern}[1]{%
  \ifrandnumbers
    \ifshoworiginal
      \ifshowtranslation\else
        \ifinlawbox\lawboxstorerandnumber{#1}\else\randnumberbox{#1}\fi
      \fi
    \fi
  \fi
}
\newcommand{\transrn}[1]{%
  \ifrandnumbers
    \ifshowtranslation
      \ifinlawbox\lawboxstorerandnumber{#1}\else\randnumberbox{#1}\fi
    \fi
  \fi
}
% ===== 課堂模式：德文原文縮寫註解 =====
% 日常切換只改各判決檔的一行：
%   \classroommodeon        開啟課堂版；註解掉這行即回到一般版。
% 模板預設：
%   1. 課堂模式關閉。
%   2. 課堂模式開啟時，縮寫註解包含「德文全稱＋中文譯名」。
%   3. 課堂模式開啟時，GG/條文編者註仍會自動顯示。
% 進階微調才需要在單案檔覆寫：
%   \classroomabbrzhoff     縮寫註解僅顯示德文全稱。
%   \showeditornotestrue    一般版也顯示編者註。
% 註解策略：同一實際 PDF 頁面中，同一縮寫只在第一次出現處加註。
% 這裡用 zref-abspage 的絕對頁序，不用 \thepage，避免文件內頁碼重置時誤判。
\newif\ifclassroommode
\newif\ifclassroomabbrzh
\classroommodefalse
\classroomabbrzhtrue
\newcommand{\classroommodeon}{\classroommodetrue}
\newcommand{\classroommodeoff}{\classroommodefalse}
\newcommand{\classroomabbrzhon}{\classroomabbrzhtrue}
\newcommand{\classroomabbrzhoff}{\classroomabbrzhfalse}
\newcommand{\abbrnotelabel}{\emph{Abk.}: }
\newcommand{\abbrnote}[3]{%
  \ifclassroommode
    \ifshoworiginal
      \ifintranslation\else
        \footnote{\normalfont\footnotesize\abbrnotelabel #2\ifclassroomabbrzh；{\songfont #3}\fi}%
      \fi
    \fi
  \fi
}
\newcommand{\abbrnoteonce}[4]{%
  #2%
  \ifclassroommode
    \ifshoworiginal
      \ifintranslation\else
        \begingroup
          \edef\abbrcurrentabspage{\number\numexpr\value{abspage}+1\relax}%
          \ifcsname abbrnoteseen@#1@\abbrcurrentabspage\endcsname\else
            \global\expandafter\let\csname abbrnoteseen@#1@\abbrcurrentabspage\endcsname\empty
            \abbrnote{#1}{#3}{#4}%
          \fi
        \endgroup
      \fi
    \fi
  \fi
}
\newcommand{\abbrAbs}{\abbrnoteonce{abs}{Abs.}{Absatz}{項}}
\newcommand{\abbrAAO}{\abbrnoteonce{aao}{a.a.O.}{am angegebenen Ort}{前引處}}
\newcommand{\abbrAE}{\abbrnoteonce{ae}{AE}{Alternativ-Entwurf}{替代草案}}
\newcommand{\abbrArt}{\abbrnoteonce{art}{Art.}{Artikel}{條}}
\newcommand{\abbrAufl}{\abbrnoteonce{aufl}{Aufl.}{Auflage}{版}}
\newcommand{\abbrBd}{\abbrnoteonce{bd}{Bd.}{Band}{卷}}
\newcommand{\abbrBGBl}{\abbrnoteonce{bgbl}{BGBl.}{Bundesgesetzblatt}{聯邦法律公報}}
\newcommand{\abbrBGHSt}{\abbrnoteonce{bghst}{BGHSt}{Entscheidungen des Bundesgerichtshofes in Strafsachen}{聯邦普通法院刑事判決彙編}}
\newcommand{\abbrBRDrucks}{\abbrnoteonce{brdrucks}{BRDrucks.}{Bundesrats-Drucksache}{聯邦參議院文件}}
\newcommand{\abbrBTDrucks}{\abbrnoteonce{btdrucks}{BTDrucks.}{Bundestags-Drucksache}{聯邦議院文件}}
\newcommand{\abbrBundesgesetzbl}{\abbrnoteonce{bundesgesetzbl}{Bundesgesetzbl.}{Bundesgesetzblatt}{聯邦法律公報}}
\newcommand{\abbrBvF}{\abbrnoteonce{bvf}{BvF}{Bundesverfassungsgerichtsverfahren}{聯邦憲法法院程序案號}}
\newcommand{\abbrBvL}{\abbrnoteonce{bvl}{BvL}{Bundesverfassungsgerichtsverfahren}{聯邦憲法法院程序案號}}
\newcommand{\abbrBvR}{\abbrnoteonce{bvr}{BvR}{Bundesverfassungsgerichtsverfahren}{聯邦憲法法院程序案號}}
\newcommand{\abbrBVerfG}{\abbrnoteonce{bverfg}{BVerfG}{Bundesverfassungsgericht}{聯邦憲法法院}}
\newcommand{\abbrBVerfGE}{\abbrnoteonce{bverfge}{BVerfGE}{Entscheidungen des Bundesverfassungsgerichts}{聯邦憲法法院判決集}}
\newcommand{\abbrBVerfGG}{\abbrnoteonce{bverfgg}{BVerfGG}{Bundesverfassungsgerichtsgesetz}{聯邦憲法法院法}}
\newcommand{\abbrDDR}{\abbrnoteonce{ddr}{DDR}{Deutsche Demokratische Republik}{德意志民主共和國；東德}}
\newcommand{\abbrDJT}{\abbrnoteonce{djt}{DJT}{Deutscher Juristentag}{德國法學家大會}}
\newcommand{\abbrDr}{\abbrnoteonce{dr}{Dr.}{Doktor}{Dr.}}
\newcommand{\abbrEuGRZ}{\abbrnoteonce{eugrz}{EuGRZ}{Europäische Grundrechte-Zeitschrift}{歐洲基本權期刊}}
\newcommand{\abbrEV}{\abbrnoteonce{ev}{EV}{Einigungsvertrag}{統一條約}}
\newcommand{\abbrF}{\abbrnoteonce{f}{f.}{folgende}{以下一頁；下一段}}
\newcommand{\abbrFF}{\abbrnoteonce{ff}{ff.}{fortfolgende}{以下數頁；以下各段}}
\newcommand{\abbrGBlDDR}{\abbrnoteonce{gblddr}{GBl. DDR}{Gesetzblatt der Deutschen Demokratischen Republik}{德意志民主共和國法律公報}}
\newcommand{\abbrGG}{\abbrnoteonce{gg}{GG}{Grundgesetz}{基本法}}
\newcommand{\abbrGVBl}{\abbrnoteonce{gvbl}{GVBl.}{Gesetz- und Verordnungsblatt}{法律及命令公報}}
\newcommand{\abbrJR}{\abbrnoteonce{jr}{JR}{Juristische Rundschau}{法學評論}}
\newcommand{\abbrJZ}{\abbrnoteonce{jz}{JZ}{JuristenZeitung}{法學家報}}
\newcommand{\abbrIVm}{\abbrnoteonce{ivm}{i.V.m.}{in Verbindung mit}{結合；連同}}
\newcommand{\abbrMdB}{\abbrnoteonce{mdb}{MdB}{Mitglied des Bundestages}{聯邦議院議員}}
\newcommand{\abbrMwN}{\abbrnoteonce{mwn}{m.w.N.}{mit weiteren Nachweisen}{附進一步文獻／判例出處}}
\newcommand{\abbrNF}{\abbrnoteonce{nf}{n.F.}{neue Fassung}{新版本}}
\newcommand{\abbrAF}{\abbrnoteonce{af}{a.F.}{alte Fassung}{舊版本}}
\newcommand{\abbrNr}{\abbrnoteonce{nr}{Nr.}{Nummer}{款、目或號，依語境}}
\newcommand{\abbrProf}{\abbrnoteonce{prof}{Prof.}{Professor}{教授}}
\newcommand{\abbrRdnr}{\abbrnoteonce{rdnr}{Rdnr.}{Randnummer}{邊碼；段碼}}
\newcommand{\abbrRdnrn}{\abbrnoteonce{rdnrn}{Rdnrn.}{Randnummern}{邊碼；段碼}}
\newcommand{\abbrRGBl}{\abbrnoteonce{rgbl}{RGBl.}{Reichsgesetzblatt}{帝國法律公報}}
\newcommand{\abbrRGSt}{\abbrnoteonce{rgst}{RGSt}{Entscheidungen des Reichsgerichts in Strafsachen}{帝國法院刑事判決彙編}}
\newcommand{\abbrRspr}{\abbrnoteonce{rspr}{Rspr.}{Rechtsprechung}{判例；司法實務}}
\newcommand{\abbrRVO}{\abbrnoteonce{rvo}{RVO}{Reichsversicherungsordnung}{帝國保險法}}
\newcommand{\abbrS}{\abbrnoteonce{s}{S.}{Seite}{頁}}
\newcommand{\abbrSFHG}{\abbrnoteonce{sfhg}{SFHG}{Schwangeren- und Familienhilfegesetz}{孕婦及家庭扶助法}}
\newcommand{\abbrSGBV}{\abbrnoteonce{sgbv}{SGB V}{Fünftes Buch Sozialgesetzbuch}{社會法典第五編}}
\newcommand{\abbrStAG}{\abbrnoteonce{staeg}{StÄG}{Strafrechtsänderungsgesetz}{刑法修正法}}
\newcommand{\abbrStenBer}{\abbrnoteonce{stenber}{StenBer.}{Stenographischer Bericht}{速記紀錄}}
\newcommand{\abbrStGB}{\abbrnoteonce{stgb}{StGB}{Strafgesetzbuch}{刑法}}
\newcommand{\abbrStREG}{\abbrnoteonce{streg}{StREG}{Strafrechtsreform-Ergänzungsgesetz}{刑法改革補充法}}
\newcommand{\abbrStrRG}{\abbrnoteonce{strrg}{StrRG}{Strafrechtsreformgesetz}{刑法改革法}}
\newcommand{\abbrUS}{\abbrnoteonce{us}{U.S.}{United States Reports}{美國最高法院判決彙編}}
\newcommand{\abbrVgl}{\abbrnoteonce{vgl}{vgl.}{vergleiche}{參見}}
\newcommand{\abbrVVDStRL}{\abbrnoteonce{vvdstrl}{VVDStRL}{Veröffentlichungen der Vereinigung der Deutschen Staatsrechtslehrer}{德國公法教師協會論文集}}
\newcommand{\abbrWP}{\abbrnoteonce{wp}{WP.}{Wahlperiode}{屆期}}
\newcommand{\abbrWp}{\abbrnoteonce{wp}{Wp.}{Wahlperiode}{屆期}}
\newcommand{\abbrZStrW}{\abbrnoteonce{zstrw}{ZStrW}{Zeitschrift für die gesamte Strafrechtswissenschaft}{整體刑法學期刊}}
\newcommand{\recordsourcerandstart}{%
  \ifrandnumbers
    \setcounter{lastsourcerandstart}{\value{randnumber}}%
    \global\lastsourcehasrandnumberstrue
  \else
    \global\lastsourcehasrandnumbersfalse
  \fi
}
\newlength{\biparagraphskip}
\setlength{\biparagraphskip}{0.52em}
\newlength{\lawboxblockbeforeskip}
\setlength{\lawboxblockbeforeskip}{0.72em}
\newlength{\lawboxblockafterskip}
\setlength{\lawboxblockafterskip}{0.92em}
\newif\iflawboxpartendedblock
\lawboxpartendedblockfalse
\newcommand{\sourceparstyle}{\footnotesize\color{sourceink}\setstretch{1.12}\RaggedRight\emergencystretch=3em\setlength{\parskip}{0.42em}\obeylines\selectfont}
\newcommand{\transparstyle}{\songfont\color{caseink}\setstretch{1.12}\setlength{\parindent}{0pt}\setlength{\parskip}{0.42em}}
\newcommand{\bilingualpairskip}{%
  \par\addvspace{\biparagraphskip}%
  \switchcolumn
  \par\addvspace{\biparagraphskip}%
  \switchcolumn*%
  \global\intranslationfalse
}
\newcommand{\bilingualboxpairskip}{%
  \par
  \switchcolumn
  \par
  \switchcolumn*%
  \global\intranslationfalse
}
\newcommand{\bilinguallawboxblockskip}{%
  \switchcolumn*%
  \par\addvspace{\lawboxblockafterskip}%
  \switchcolumn
  \par\addvspace{\lawboxblockafterskip}%
  \switchcolumn*%
  \global\intranslationfalse
}
\newcommand{\bikeepwithnext}[1][4]{%
  \ifbilingualcolumns
    \syncleft
    \Needspace{#1\baselineskip}%
    \switchcolumn
    \Needspace{#1\baselineskip}%
    \switchcolumn*%
    \global\intranslationfalse
  \else
    \Needspace{#1\baselineskip}%
  \fi
}
\NewEnviron{sourcepara}[1][]{
  \recordsourcerandstart
  \ifshoworiginal
    \ifbilingualcolumns
      \ifintranslation\switchcolumn*\global\intranslationfalse\fi
    \else
      \Needspace{5\baselineskip}
    \fi
    \ifbilingualcolumns\else\par\addvspace{0.35em}\fi
    {\sourceparstyle
    \noindent\BODY\par}
    \ifbilingualcolumns\else\addvspace{0.25em}\fi
    \ifbilingualcolumns
      \switchcolumn
      \global\intranslationtrue
    \fi
  \else
    \ifrandnumbers
      \begingroup
        \setbox0=\vbox{\hsize=\linewidth\sourceparstyle\noindent\BODY\par}%
      \endgroup
    \fi
  \fi
}
\NewEnviron{transpara}{
  \ifshowtranslation
    \setcounter{randnumberlocal}{\value{lastsourcerandstart}}%
    \ifbilingualcolumns
      \ifintranslation\else\switchcolumn\global\intranslationtrue\fi
      {\transparstyle\setlength{\leftskip}{0.55em}\setlength{\rightskip}{0.15em plus 1em}\addtolength{\linewidth}{-0.55em}\BODY\par}
      \switchcolumn*
      \bilingualpairskip
    \else
      {\transparstyle\BODY\par}
    \fi
  \else
    \ifbilingualcolumns
      \ifintranslation\switchcolumn*\global\intranslationfalse\fi
    \fi
  \fi
}
\NewEnviron{sourceboxpara}[1][]{
  \global\lawboxpartendedblockfalse
  \recordsourcerandstart
  \ifshoworiginal
    \ifbilingualcolumns
      \ifintranslation\switchcolumn*\global\intranslationfalse\fi
    \else
      \Needspace{3\baselineskip}
    \fi
    {\sourceparstyle
    \BODY\par}
    \ifbilingualcolumns\else
      \iflawboxpartendedblock\addvspace{\lawboxblockafterskip}\fi
    \fi
    \ifbilingualcolumns
      \switchcolumn
      \global\intranslationtrue
    \fi
  \fi
}
\NewEnviron{transboxpara}{
  \global\lawboxpartendedblockfalse
  \ifshowtranslation
    \setcounter{randnumberlocal}{\value{lastsourcerandstart}}%
    \ifbilingualcolumns
      \ifintranslation\else\switchcolumn\global\intranslationtrue\fi
      {\transparstyle\BODY\par}
      \iflawboxpartendedblock
        \bilinguallawboxblockskip
      \else
        \switchcolumn*
        \bilingualboxpairskip
      \fi
    \else
      {\transparstyle\BODY\par}
      \iflawboxpartendedblock\addvspace{\lawboxblockafterskip}\fi
    \fi
  \else
    \ifbilingualcolumns
      \ifintranslation\switchcolumn*\global\intranslationfalse\fi
    \fi
  \fi
}
\newcommand{\leitsatznum}[1]{\noindent\hangindent=3.0em\hangafter=1\makebox[2.25em][r]{\bfseries #1.}\hspace{0.55em}\ignorespaces}
\newcommand{\syncleft}{\ifbilingualcolumns\ifintranslation\switchcolumn*\global\intranslationfalse\fi\fi}
\pretocmd{\section}{\syncleft\clearpage}{}{}
\pretocmd{\subsection}{\syncleft}{}{}
\pretocmd{\subsubsection}{\syncleft}{}{}
\newcommand{\biheadingfont}{\songfont}
\newcommand{\bitocentry}[2]{%
  \ifshoworiginal
    \ifshowtranslation
      \ifstrequal{#1}{#2}{#1}{#1\space #2}%
    \else
      #1%
    \fi
  \else
    #2%
  \fi
}
\pdfstringdefDisableCommands{%
  \def\bitocentry#1#2{#1}%
}
\newcommand{\biheading}[7]{%
  \syncleft#1\bikeepwithnext[#3]\phantomsection\addcontentsline{toc}{#2}{\bitocentry{#6}{#7}}%
  \ifbilingualcolumns
    {#4 #6\par}%
    \switchcolumn
    {#4\biheadingfont #7\par}%
    \switchcolumn*%
    \global\intranslationfalse
    \par\addvspace{#5}%
    \switchcolumn
    \par\addvspace{#5}%
    \switchcolumn*%
  \else
    \ifshoworiginal{#4 #6\par}\fi
    \ifshowtranslation{#4\biheadingfont #7\par}\fi
    \addvspace{#5}%
  \fi
}
\newcommand{\termsectionheading}[1]{%
  \par\Needspace{9\baselineskip}%
  \vspace{0.65em}%
  {\songfont\bfseries #1\par}%
  \vspace{0.2em}%
}
% 章節換頁：
%   雙語 paracol 欄內不要用 \clearpage；它會在同步兩欄時吐出只有欄線/頁碼的空頁。
%   \newpage 足以讓下一個 section 起新頁，又不會強制清空所有待排材料。
%   單語模式不在 paracol 內，仍可用 \clearpage。
\newcommand{\bisectionpagebreak}{\ifbilingualcolumns\newpage\else\clearpage\fi}
\newcommand{\bisection}[2]{%
  \biheading{\iffirstbilingualsection\global\firstbilingualsectionfalse\else\bisectionpagebreak\fi}{section}{6}{\centering\Large\bfseries}{0.8em}{#1}{#2}%
}
\newcommand{\bisubsection}[2]{%
  \biheading{}{subsection}{5}{\large\bfseries}{0.55em}{#1}{#2}%
}
\newcommand{\bisubsubsection}[2]{%
  \biheading{}{subsubsection}{4}{\normalsize\bfseries}{0.45em}{#1}{#2}%
}
\newcommand{\bicompositeheading}[4]{%
  \biheading{}{subsection}{5}{\large\bfseries}{0.16em}{#1}{#2}%
  \biheading{}{subsubsection}{3}{\normalsize\bfseries}{0.45em}{#3}{#4}%
}
\newif\iflawboxfirstitem
\lawboxfirstitemfalse
\newcommand{\lawboxprespace}[1]{%
  \par
  \iflawboxfirstitem
    \global\lawboxfirstitemfalse
  \else
    \addvspace{#1}%
  \fi
}
\newtcolorbox{lawboxinner}{
  caseboxbase,
  colframe=sourceline!85!black,
  colback=white,
  left=1.05em,
  right=1.05em,
  top=0.62em,
  bottom=0.62em,
  before upper={\global\inlawboxtrue\global\lawboxfirstitemtrue\global\lawboxhaspendingrnfalse\ifintranslation\lawboxtransfont\else\lawboxsourcefont\fi\RaggedRight\setlength{\parindent}{0pt}\setlength{\parskip}{0pt}\ifintranslation\setstretch{\lawboxtransstretch}\else\setstretch{\lawboxsourcestretch}\fi},
  after upper={\global\lawboxhaspendingrnfalse\global\inlawboxfalse}
}
\tcbset{
  lawboxpartbase/.style={
    enhanced,
    breakable=false,
    width=0.985\linewidth,
    colback=white,
    frame hidden,
    boxrule=0pt,
    arc=0pt,
    left=1.05em,
    right=1.05em,
    boxsep=1.5mm,
    before skip=0pt,
    after skip=0pt,
    before upper={\global\inlawboxtrue\global\lawboxfirstitemtrue\global\lawboxhaspendingrnfalse\ifintranslation\lawboxtransfont\else\lawboxsourcefont\fi\RaggedRight\setlength{\parindent}{0pt}\setlength{\parskip}{0pt}\ifintranslation\setstretch{\lawboxtransstretch}\else\setstretch{\lawboxsourcestretch}\fi},
    after upper={\global\lawboxhaspendingrnfalse\global\inlawboxfalse},
    borderline west={0.28pt}{0pt}{sourceline!85!black},
    borderline east={0.28pt}{0pt}{sourceline!85!black},
  },
  lawboxpartfirst/.style={
    lawboxpartbase,
    top=0.62em,
    bottom=0.04em,
    before skip=\lawboxblockbeforeskip,
    borderline north={0.28pt}{0pt}{sourceline!85!black},
  },
  lawboxpartmiddle/.style={
    lawboxpartbase,
    top=0.04em,
    bottom=0.04em,
  },
  lawboxpartlast/.style={
    lawboxpartbase,
    top=0.04em,
    bottom=0.62em,
    after skip=0pt,
    borderline south={0.28pt}{0pt}{sourceline!85!black},
  }
}
\newtcolorbox{lawboxpartinner}[1]{#1}
\newtcolorbox{termboxinner}{
  caseboxbase,
  colframe=noteblue!70!black,
  colback=casepaper,
  before upper={\songfont\setlength{\parindent}{0pt}\setlength{\parskip}{0.35em}}
}
\newtcolorbox{deboxinner}{
  caseboxbase,
  colframe=notegray!72!black,
  colback=white,
  left=1.05em,
  right=1.05em,
  top=0.62em,
  bottom=0.62em,
  before upper={\global\inlawboxtrue\global\lawboxfirstitemtrue\global\lawboxhaspendingrnfalse\small\setlength{\parindent}{0pt}\setlength{\parskip}{0.35em}},
  after upper={\global\lawboxhaspendingrnfalse\global\inlawboxfalse}
}
\NewEnviron{lawbox}{
  \begin{lawboxinner}\BODY\end{lawboxinner}
}
\NewEnviron{lawboxpart}[2]{
  \begin{lawboxpartinner}{lawboxpart#1,equal height group=#2}\BODY\end{lawboxpartinner}%
  \ifstrequal{#1}{last}{\global\lawboxpartendedblocktrue}{}%
}
\NewEnviron{termbox}{
  \par\addvspace{0.55em}
  {\color{noteblue!70!black}\rule{\linewidth}{0.28pt}}\par
  \begingroup
    \songfont\small\color{caseink}
    \setlength{\parindent}{0pt}
    \setlength{\parskip}{0.24em}
    \BODY
    \par
  \endgroup
  {\color{noteblue!70!black}\rule{\linewidth}{0.28pt}}\par
  \addvspace{0.55em}
}
\NewEnviron{debox}{
  \begin{deboxinner}\BODY\end{deboxinner}
}
\providecommand{\paragraphmark}[1]{}
\newcommand{\lawboxtitleafter}{\addvspace{0.06em}}
\newcommand{\lawtitle}[1]{\lawboxprespace{0.22em}{\lawboxtransfont\bfseries\lawboxuserandnumber #1}\par\lawboxtitleafter}
\newcommand{\absno}[2]{\lawboxprespace{0.04em}\noindent\lawboxuserandnumber\hangindent=2.6em\hangafter=1\makebox[2.6em][l]{(#1)}#2\par}
\newcommand{\nrno}[2]{\lawboxprespace{0pt}\noindent\lawboxuserandnumber\hspace*{2.6em}\hangindent=4.4em\hangafter=1\makebox[1.8em][l]{#1.}#2\par}
\newcommand{\letterno}[2]{\lawboxprespace{0pt}\noindent\lawboxuserandnumber\hspace*{4.4em}\hangindent=6.0em\hangafter=1\makebox[1.6em][l]{#1)}#2\par}
\newcommand{\sourcelawtitle}[1]{\lawboxprespace{0.22em}{\lawboxsourcefont\bfseries\lawboxuserandnumber #1}\par\lawboxtitleafter}
\newcommand{\sourcelawsubtitle}[1]{\lawboxprespace{0.04em}{\lawboxsourcefont\bfseries\lawboxuserandnumber #1}\par\lawboxtitleafter}
\newcommand{\sourceabsno}[2]{\lawboxprespace{0.04em}\noindent\lawboxuserandnumber\hangindent=2.45em\hangafter=1\makebox[2.45em][l]{(#1)}#2\par}
\newcommand{\sourcenrno}[2]{\lawboxprespace{0pt}\noindent\lawboxuserandnumber\hspace*{2.45em}\hangindent=4.25em\hangafter=1\makebox[1.8em][l]{#1.}#2\par}
\newcommand{\sourceletterno}[2]{\lawboxprespace{0pt}\noindent\lawboxuserandnumber\hspace*{4.25em}\hangindent=5.85em\hangafter=1\makebox[1.6em][l]{#1)}#2\par}
\newcommand{\sourcequotepara}[1]{\lawboxprespace{0.04em}\noindent\lawboxuserandnumber\hangindent=1.2em\hangafter=1 #1\par}
\newcommand{\sourcelawline}[1]{\lawboxprespace{0.04em}\noindent\lawboxuserandnumber #1\par}
\newcommand{\lawline}[1]{\lawboxprespace{0.04em}\noindent\lawboxuserandnumber #1\par}
\newcommand{\pendingtranslation}{\par{\songfont\footnotesize\color{pendingink}（中文譯文待補。）}\par}
\newcommand{\tocpart}[1]{\par\addvspace{0.52em}{\footnotesize\bfseries\color{pendingink}#1\par}\addvspace{0.16em}}
\newcommand{\tocdashline}{\par\addvspace{0.48em}\noindent{\color{notegray}\xleaders\hbox to 0.82em{\hfil\rule{0.42em}{0.28pt}\hfil}\hskip\linewidth}\par\addvspace{0.38em}}
% \tocpanel{font-cmd}{content}：將目錄置中於所在欄寬（雙語欄適用）。
% A3 橫式使用 0.58\linewidth，A4 橫式使用 0.84\linewidth，
% 使兩種紙張下目錄欄內容實際寬度相近（均約 100–105 mm）。
\newcommand{\tocpanel}[2]{%
  \makebox[\linewidth][c]{%
    \ifx\bilingualpaper\bilingualpaperaThree
      \begin{minipage}{0.58\linewidth}%
    \else
      \begin{minipage}{0.84\linewidth}%
    \fi
    \toccolstyle
    #1#2%
    \end{minipage}%
  }%
}
% ===== 首頁簡目錄的兩層 description list 縮排 =====
% enumitem 的 description list 可把每一列拆成：
%   [label]  labelsep  body text
% 這裡的「起點」都以所在 minipage/欄位的左緣為 0。
%
% 核心規則：
%   labelindent = label 欄位從左緣開始的位置。
%   labelwidth  = label 欄位本身寬度；標籤太長（如 Abw. Meinung）就加大它。
%   labelsep    = label 欄位右緣到正文的距離。
%   leftmargin  = 正文 body text 的起點。判斷層級視覺時主要看這個。
%
% 所以：
%   想讓「Begründung / 判決理由」整列正文往右或往左：改 tocitems 的 leftmargin。
%   想讓「A. Verfahren und Gegenstand / A. 程序與審查標的」正文往右或往左：
%     改 tocsubitems 的 leftmargin。
%   想讓 A./B./C. 這些標籤本身往右或往左，但正文不動：改 tocsubitems 的 labelindent。
%   想讓 A./B./C. 與正文間距變大或變小：改 tocsubitems 的 labelsep。
%
% 層級原則：
%   tocsubitems.leftmargin 必須大於 tocitems.leftmargin，
%   否則子層級正文會看起來比「Gründe / 理由」還靠左。
\newlist{tocitems}{description}{1}
\setlist[tocitems]{%
  labelindent=0pt,    % 第一層 label 欄位起點；0pt 表示貼齊目錄欄左緣。
  leftmargin=2.54cm,  % 第一層正文起點；例如 Begründung / 判決理由 的左緣。
  labelwidth=2.16cm,  % 第一層 label 欄位寬度；需容納 Leitsätze、Abw. Meinung、不同意見等。
  labelsep=0.32cm,    % 第一層 label 與正文的水平間距；只改間距，不改正文起點。
  itemsep=0.28em,     % 第一層各項之間的垂直距離。
  topsep=0.20em,      % 第一層 list 上下與前後內容的垂直距離。
  parsep=0pt,         % 同一 item 內段落之間的距離；目錄項通常不需要。
  partopsep=0pt,      % list 若緊接段落開始時的額外距離；關閉以免目錄高度飄動。
  align=left,         % label 欄內左對齊；長短標籤不靠右貼齊。
  font=\normalfont\bfseries % label 的字型；正文不受這行影響。
}
\newlist{tocsubitems}{description}{1}
\setlist[tocsubitems]{%
  labelindent=1cm,    % 第二層 label 起點；只控制 A./B./C. 本身的位置。
  leftmargin=3.00cm,  % 第二層正文起點；必須大於 2.54cm，才會比 Begründung / 判決理由更右。
  labelwidth=0.5cm,   % 第二層 label 欄寬；A./B./C. 很短，可小於第一層。
  labelsep=1.5cm,    % 第二層 label 與正文間距；A. 與 Verfahren / 程序 的距離看這裡。
  itemsep=0.12em,     % 第二層各項較密，避免 A.–E. 佔太高。
  topsep=0.04em,      % 第二層 list 與 Gründe / 理由之間的垂直距離。
  parsep=0pt,         % 同一第二層 item 內段落距離；目錄項通常不需要。
  partopsep=0pt,      % 關閉額外段落起始距離，避免版面不穩。
  align=left,         % A./B./C. label 欄內左對齊。
  font=\normalfont\bfseries, % A./B./C. label 加粗；正文不受這行影響。
  before={}           % 不在第二層 list 前自動插入額外內容。
}
\newlist{termitems}{description}{1}
\ifbilingualcolumns
  \setlist[termitems]{%
    leftmargin=5.2cm,
    labelwidth=4.8cm,
    labelsep=0.35cm,
    itemsep=0.16em,
    topsep=0.2em,
    parsep=0pt,
    partopsep=0pt,
    font=\normalfont\bfseries,
    style=nextline
  }
\else
  \setlist[termitems]{%
    leftmargin=0pt,
    labelwidth=0pt,
    labelsep=0pt,
    itemsep=0.24em,
    topsep=0.2em,
    parsep=0pt,
    partopsep=0pt,
    font=\normalfont\bfseries,
    style=nextline
  }
\fi
\newlist{abbritems}{description}{1}
\setlist[abbritems]{%
  leftmargin=2.38cm,
  labelwidth=2.02cm,
  labelsep=0.28cm,
  itemsep=0.16em,
  topsep=0.2em,
  parsep=0pt,
  partopsep=0pt,
  align=left,
  font=\normalfont\bfseries
}
\ifbilingualcolumns
  \newenvironment{termcolumns}{\begin{multicols}{2}}{\end{multicols}}
\else
  \newenvironment{termcolumns}{}{}
\fi

% ===== 編者註腳開關 =====
% 一般版預設不顯示編者註，避免正文腳註過密。
% 若開啟 \classroommodeon，GG/條文編者註仍會自動顯示，無需另開本開關。
% 若一般版也要顯示編者註，可在單案檔 \input 後加 \showeditornotestrue。
\newif\ifshoweditornotes
\showeditornotesfalse

% ===== 目錄排版輔助 =====
\newcommand{\toccolstyle}{\raggedright\arraybackslash}
\newcommand{\tocsingleindent}{\leftskip=1.45cm\rightskip=1.2cm}

% ===== 行內引文環境（德文原文與中文譯文各一，帶左側灰色邊線）=====
\newcommand{\quoteparbreak}{\par\addvspace{0.48em}}
\newcommand{\sourcequoteinner}[1]{%
  \begin{tcolorbox}[
    enhanced, breakable, width=\dimexpr\linewidth-0.8em\relax, frame hidden,
    colback=white, coltext=caseink, boxrule=0pt,
    borderline west={0.55pt}{0.88em}{notegray},
    left=1.78em, right=0.72em, top=0.18em, bottom=0.18em,
    boxsep=0pt, before skip=0.24em, after skip=0.24em,
    before upper={\normalfont\footnotesize\color{caseink}\setlength{\parindent}{0pt}\setlength{\parskip}{0.34em}\raggedright\setlength{\rightskip}{0pt plus 1em}}
  ]%
  #1\par
  \end{tcolorbox}%
}
\newcommand{\transquoteinner}[1]{%
  \begin{tcolorbox}[
    enhanced, breakable, width=\dimexpr\linewidth-0.8em\relax, frame hidden,
    colback=white, coltext=caseink, boxrule=0pt,
    borderline west={0.55pt}{0.88em}{notegray},
    left=1.78em, right=0.72em, top=0.18em, bottom=0.18em,
    boxsep=0pt, before skip=0.24em, after skip=0.24em,
    before upper={\songfont\small\color{caseink}\setlength{\parindent}{0pt}\setlength{\parskip}{0.34em}\raggedright\setlength{\rightskip}{0pt plus 1em}}
  ]%
  #1\par
  \end{tcolorbox}%
}

% ===== 大綱（Gliederung）Rn 輔助命令 =====
\newcommand{\outrn}[2]{\enspace{\normalfont\footnotesize\color{pendingink}(Rn\,#1–#2)}}
\newcommand{\outrnsingle}[1]{\enspace{\normalfont\footnotesize\color{pendingink}(Rn\,#1)}}
\newcommand{\outrnzh}[2]{\enspace{\normalfont\footnotesize\color{pendingink}（第\,#1–#2\,段）}}
\newcommand{\outrnzhs}[1]{\enspace{\normalfont\footnotesize\color{pendingink}（第\,#1\,段）}}
% ===== 大綱雙語對照表格輔助命令（三欄：段碼 │ 德文 │ 中文）=====
% 欄 1：段碼（由欄定義自動格式化：小字、等寬、右對齊、pendingink）
% 欄 2：德文標籤＋正文
% 欄 3：中文標籤＋正文
%
% 無段碼的列（\omainrow、\osecrow），欄 1 以 & 空置。
% 有段碼的列（\osecrowrn、\osubrow、\ointrow），#1 為預格式化字串
%   例：\osubrow{2–49}{I.}{...}{I.}{...}
%       \ointrow{107}{...}{...}
%
% \opartrow：段組標題（橫跨三欄）
\newcommand{\opartrow}[2]{%
  \noalign{\vspace{0.30em}}%
  \multicolumn{3}{@{}l@{}}{%
    \footnotesize\bfseries\color{pendingink}#1\hfill\songfont#2%
  }\\[0.06em]%
}
% \odashrow：虛線分隔（橫跨三欄）
\newcommand{\odashrow}{%
  \noalign{\vspace{0.28em}}%
  \multicolumn{3}{@{}l@{}}{\color{notegray}%
    \xleaders\hbox to 0.82em{\hfil\rule{0.42em}{0.28pt}\hfil}\hskip 0.88\linewidth}%
  \\[0.24em]%
}
% \odissentpart：不同意見書區段標頭。比一般 \opartrow 更像附錄性轉場。
\newcommand{\odissentpart}[2]{%
  \noalign{\vspace{0.42em}}%
  \multicolumn{3}{@{}p{0.88\textwidth}@{}}{%
    \color{notegray}\rule{\linewidth}{0.18pt}%
  }\\[-0.18em]%
  \multicolumn{3}{@{}p{0.88\textwidth}@{}}{%
    \makebox[\linewidth][s]{%
      \footnotesize\bfseries\color{pendingink}#1%
      \hfill{\songfont #2}%
    }%
  }\\[-0.02em]%
}
% \odissentopinion：不同意見書的署名與一行摘要，右欄保留段碼範圍。
\newcommand{\odissentopinion}[5]{%
  \footnotesize\hangindent=0.52cm\hangafter=1%
    {\bfseries\color{caseink}#2}\enspace{\color{notegray}\textperiodcentered}\enspace\color{caseink}#3%
  &\footnotesize\hangindent=0.52cm\hangafter=1%
    {\songfont\bfseries\color{caseink}#4}\enspace{\color{notegray}\textperiodcentered}\enspace\songfont\color{caseink}#5%
  &#1%
  \\[0.12em]%
}
% \omainrow：頂層項目，無段碼（Leitsätze、Urteil、Gründe …）
\newcommand{\omainrow}[4]{%
  \footnotesize\hangindent=1.92cm\hangafter=1%
    \makebox[1.92cm][l]{\bfseries\color{caseink}#1}\color{caseink}#2%
  &\footnotesize\hangindent=1.92cm\hangafter=1%
    \makebox[1.92cm][l]{\bfseries\color{caseink}#3}\songfont\color{caseink}#4%
  &%
  \\[0.15em]%
}
% \omainrowrn：頂層項目，有段碼（不同意見等標籤寬於 \osecrow 框者）
\newcommand{\omainrowrn}[5]{%
  \footnotesize\hangindent=1.92cm\hangafter=1%
    \makebox[1.92cm][l]{\bfseries\color{caseink}#2}\color{caseink}#3%
  &\footnotesize\hangindent=1.92cm\hangafter=1%
    \makebox[1.92cm][l]{\bfseries\color{caseink}#4}\songfont\color{caseink}#5%
  &#1%
  \\[0.15em]%
}
% \osecrow：一級子項，無段碼（A.、C.、D. 等有子項者）
\newcommand{\osecrow}[4]{%
  \footnotesize\hspace{1.10cm}\hangindent=1.98cm\hangafter=1%
    \makebox[0.86cm][l]{\bfseries\color{caseink}#1}\color{caseink}#2%
  &\footnotesize\hspace{1.10cm}\hangindent=1.98cm\hangafter=1%
    \makebox[0.86cm][l]{\bfseries\color{caseink}#3}\songfont\color{caseink}#4%
  &%
  \\[0.11em]%
}
% \osecrowrn：一級子項，有段碼（B.、E. 等完整段落範圍者）
% #1=段碼字串（如 101–106）, #2=DE標籤, #3=DE正文, #4=ZH標籤, #5=ZH正文
\newcommand{\osecrowrn}[5]{%
  \footnotesize\hspace{1.10cm}\hangindent=1.98cm\hangafter=1%
    \makebox[0.86cm][l]{\bfseries\color{caseink}#2}\color{caseink}#3%
  &\footnotesize\hspace{1.10cm}\hangindent=1.98cm\hangafter=1%
    \makebox[0.86cm][l]{\bfseries\color{caseink}#4}\songfont\color{caseink}#5%
  &#1%
  \\[0.11em]%
}
% \osubrow：二級子項，有段碼範圍（I.、II.、A.I. …）
% #1=段碼字串, #2=DE標籤, #3=DE正文, #4=ZH標籤, #5=ZH正文
\newcommand{\osubrow}[5]{%
  \footnotesize\hspace{1.40cm}\hangindent=2.30cm\hangafter=1%
    \makebox[0.90cm][l]{\bfseries\color{caseink}#2}\color{caseink}#3%
  &\footnotesize\hspace{1.40cm}\hangindent=2.30cm\hangafter=1%
    \makebox[0.90cm][l]{\bfseries\color{caseink}#4}\songfont\color{caseink}#5%
  &#1%
  \\[0.09em]%
}
% \ointrow：引入段落，有單一段碼，無標籤
% #1=段碼字串, #2=DE正文, #3=ZH正文
\newcommand{\ointrow}[3]{%
  \footnotesize\hspace{0.52cm}\hangindent=0.52cm\hangafter=1\color{caseink}#2%
  &\footnotesize\hspace{0.52cm}\hangindent=0.52cm\hangafter=1\songfont\color{caseink}#3%
  &#1%
  \\[0.09em]%
}
% \osubsubrow：三級子項（1.、2.、3.）
% #1=段碼字串, #2=DE標籤, #3=DE正文, #4=ZH標籤, #5=ZH正文
\newcommand{\osubsubrow}[5]{%
  \footnotesize\hspace{1.80cm}\hangindent=2.48cm\hangafter=1%
    \makebox[0.68cm][l]{\bfseries\color{caseink}#2}\color{caseink}#3%
  &\footnotesize\hspace{1.80cm}\hangindent=2.48cm\hangafter=1%
    \makebox[0.68cm][l]{\bfseries\color{caseink}#4}\songfont\color{caseink}#5%
  &#1%
  \\[0.08em]%
}
% \osubsubsubrow：四級子項（a）、b））
% #1=段碼字串, #2=DE標籤, #3=DE正文, #4=ZH標籤, #5=ZH正文
\newcommand{\osubsubsubrow}[5]{%
  \footnotesize\hspace{2.10cm}\hangindent=2.68cm\hangafter=1%
    \makebox[0.58cm][l]{\bfseries\color{caseink}#2}\color{caseink}#3%
  &\footnotesize\hspace{2.10cm}\hangindent=2.68cm\hangafter=1%
    \makebox[0.58cm][l]{\bfseries\color{caseink}#4}\songfont\color{caseink}#5%
  &#1%
  \\[0.07em]%
}
% \osubsubsubsubrow：五級子項（aa）、bb））
% 標籤框 0.62cm：足以容納「aa)」在 footnotesize bold 下（約 0.50cm）並留有間距
% #1=段碼字串, #2=DE標籤, #3=DE正文, #4=ZH標籤, #5=ZH正文
\newcommand{\osubsubsubsubrow}[5]{%
  \footnotesize\hspace{2.38cm}\hangindent=3.06cm\hangafter=1%
    \makebox[0.68cm][l]{\bfseries\color{caseink}#2}\color{caseink}#3%
  &\footnotesize\hspace{2.38cm}\hangindent=3.06cm\hangafter=1%
    \makebox[0.68cm][l]{\bfseries\color{caseink}#4}\songfont\color{caseink}#5%
  &#1%
  \\[0.06em]%
}
% \outlinepage：雙語對照大綱頁包裝環境（longtable 自動跨頁）
% 三欄：德文欄 + 中文欄 + 段碼欄（最右，不折行）
% 段碼欄寬度：1.80cm，足以容納最長段碼如「309–348」
% DE/ZH 欄寬：(0.87\textwidth - 2.08cm) / 2
%   驗算：2×欄 + 0.05\textwidth + 0.28cm gap + 1.80cm = 0.87tw - 2.08cm + 0.05tw + 2.08cm = 0.92tw ✓
\newenvironment{outlinepage}{%
  \vspace*{0.40cm}%
  \setlength{\LTleft}{0.04\textwidth}%
  \setlength{\LTright}{0.04\textwidth}%
  \setlength{\tabcolsep}{0pt}%
  % 大綱列距由 longtable 的 \arraystretch 與各 row macro 結尾的 \\[...em] 共同決定。
  % 這裡控制「每一列本身的表格 strut 高度」；數值越大，A./I./1. 各項之間越像有空行。
  % A3 原本用 0.62，視覺上沒有空行；A4 也沿用同值，避免切到 A4 後項目被撐開。
  % 若日後覺得大綱太擠，只微調這個值（例如 0.68），不要改各 \omainrow/\osubrow 的縮排。
  \renewcommand{\arraystretch}{0.62}%
  \begin{longtable}{%
    >{\vspace{0pt}\raggedright\arraybackslash}p{\dimexpr(0.87\textwidth-2.08cm)/2\relax}%
    @{\hspace{0.025\textwidth}}%
    !{\color{notegray}\vrule width 0.12pt}%
    @{\hspace{0.025\textwidth}}%
    >{\vspace{0pt}\raggedright\arraybackslash}p{\dimexpr(0.87\textwidth-2.08cm)/2\relax}%
    @{\hspace{0.28cm}}%
    >{\vspace{0pt}\footnotesize\normalfont\color{pendingink}\raggedright\arraybackslash}p{1.80cm}%
    @{}%
  }%
  % 首頁欄標籤：不要橫跨置中；分別放在德文欄與中文欄的實際欄位起點。
  \footnotesize\bfseries\color{sourceink}Gliederung%
  &\footnotesize\bfseries\songfont\color{sourceink}大綱%
  &%
  \\[0.36em]%
  \endfirsthead%
  % 續頁欄標籤：同樣分欄定位，以灰色細體區分；無需標注「續」。
  \footnotesize\color{pendingink}Gliederung%
  &\footnotesize\songfont\color{pendingink}大綱%
  &%
  \\[0.24em]%
  \endhead%
}{%
  \end{longtable}%
}
 載入。
% 不含：\documentclass、\documentversion、\ggversionde/zh、\tocde/\toczh。
% 日期與首頁版本列由本檔自動產生；各文件只需手動指定 \documentversion。

\usepackage{xeCJK}
\usepackage{fontspec}
\usepackage{geometry}
\usepackage{setspace}
\usepackage{hyperref}
\usepackage{xurl}
\usepackage{bookmark}
\usepackage{enumitem}
\usepackage{xcolor}
\usepackage{titlesec}
\usepackage{array}
\usepackage{longtable}
\usepackage{tcolorbox}
\usepackage{environ}
\usepackage{pdflscape}
\usepackage{etoolbox}
\usepackage{ragged2e}
\usepackage{paracol}
\usepackage{multicol}
\usepackage{needspace}
\usepackage{fancyhdr}
\usepackage{tikz}
\usepackage{zref-abspage}
\tcbuselibrary{breakable,skins}
\usetikzlibrary{calc}
\usepackage{pgfcalendar}

\hypersetup{
  unicode=true,
  bookmarksopen=true,
  bookmarksnumbered=false,
  hidelinks
}

% ===== 自動推導，不要手動改下面 =====
% （\docmode 與 \bilingualpaper 由各文件在 % bverfge_layout.tex
% 共用排版基礎設施，供 Schwangerschaftsabbruch I 與 II 共享。
% 由各文件以 \documentclass 後 \input{../../共用LaTeX/bverfge_layout} 載入。
% 不含：\documentclass、\documentversion、\ggversionde/zh、\tocde/\toczh。
% 日期與首頁版本列由本檔自動產生；各文件只需手動指定 \documentversion。

\usepackage{xeCJK}
\usepackage{fontspec}
\usepackage{geometry}
\usepackage{setspace}
\usepackage{hyperref}
\usepackage{xurl}
\usepackage{bookmark}
\usepackage{enumitem}
\usepackage{xcolor}
\usepackage{titlesec}
\usepackage{array}
\usepackage{longtable}
\usepackage{tcolorbox}
\usepackage{environ}
\usepackage{pdflscape}
\usepackage{etoolbox}
\usepackage{ragged2e}
\usepackage{paracol}
\usepackage{multicol}
\usepackage{needspace}
\usepackage{fancyhdr}
\usepackage{tikz}
\usepackage{zref-abspage}
\tcbuselibrary{breakable,skins}
\usetikzlibrary{calc}
\usepackage{pgfcalendar}

\hypersetup{
  unicode=true,
  bookmarksopen=true,
  bookmarksnumbered=false,
  hidelinks
}

% ===== 自動推導，不要手動改下面 =====
% （\docmode 與 \bilingualpaper 由各文件在 \input{bverfge_layout} 之前設定；
%   預設 chinese / a4）
\ifdefined\docmode\else\def\docmode{chinese}\fi
\ifdefined\bilingualpaper\else\def\bilingualpaper{a4}\fi
\newif\ifshoworiginal
\newif\ifshowtranslation
\newif\ifbilinguallandscape
\newif\ifbilingualcolumns
\newif\ifintranslation
\newif\iffirstbilingualsection

\showoriginalfalse
\showtranslationfalse
\bilinguallandscapefalse
\bilingualcolumnsfalse
\intranslationfalse
\firstbilingualsectionfalse

\def\modechinese{chinese}
\def\modegerman{german}
\def\modebilingual{bilingual}
\def\bilingualpaperaThree{a3}
\def\bilingualpaperaFour{a4}

\ifx\docmode\modechinese
  \showtranslationtrue
\fi

\ifx\docmode\modegerman
  \showoriginaltrue
\fi

\ifx\docmode\modebilingual
  \showoriginaltrue
  \showtranslationtrue
  \bilinguallandscapetrue
  \bilingualcolumnstrue
\fi

% 編排模式說明：
% chinese  → \showoriginalfalse  \showtranslationtrue  （A4 直式，純中文）
% german   → \showoriginaltrue   \showtranslationfalse （A4 直式，純德文）
% bilingual→ \showoriginaltrue   \showtranslationtrue  \bilingualcolumnstrue（A3/A4 橫式對照）

\ifbilingualcolumns
  \ifx\bilingualpaper\bilingualpaperaFour
    \geometry{
      a4paper,
      landscape,
      left=1.25cm,
      right=2.25cm,
      top=1.25cm,
      bottom=1.75cm,
      footskip=8.5mm,
      marginparwidth=0.75cm,
      marginparsep=0.25cm
    }
  \else
    \geometry{
      a3paper,
      landscape,
      left=1.55cm,
      right=2.75cm,
      top=1.55cm,
      bottom=2.05cm,
      footskip=9.5mm,
      marginparwidth=0.9cm,
      marginparsep=0.35cm
    }
  \fi
\else
\geometry{
  a4paper,
  portrait,
  left=2.2cm,
  right=2.6cm,
  top=2.0cm,
  bottom=2.0cm,
  marginparwidth=0.8cm,
  marginparsep=0.3cm
}
\fi
% ===== 時間戳基礎設施 =====
\newcount\stamptimeval
\newcount\stampminuteval
\newcommand{\stamphh}{%
  \stamptimeval=\time
  \divide\stamptimeval by 60
  \ifnum\stamptimeval<10 0\fi
  \the\stamptimeval}
\newcommand{\stampmm}{%
  \stamptimeval=\time
  \stampminuteval=\time
  \divide\stampminuteval by 60
  \multiply\stampminuteval by 60
  \advance\stamptimeval by -\stampminuteval
  \ifnum\stamptimeval<10 0\fi
  \the\stamptimeval}
\newcommand{\stampwkde}[1]{%
  \ifcase#1Montag\or Dienstag\or Mittwoch\or Donnerstag\or Freitag\or Samstag\or Sonntag\fi}
\newcommand{\stampwkzh}[1]{%
  \ifcase#1 一\or 二\or 三\or 四\or 五\or 六\or 日\fi}
\newcommand{\stampmonthde}[1]{%
  \ifcase#1\or Januar\or Februar\or März\or April\or Mai\or Juni\or
  Juli\or August\or September\or Oktober\or November\or Dezember\fi}
\newcount\stampjdval
\newcount\stampwdval
\pgfcalendardatetojulian{\the\year-\the\month-\the\day}{\stampjdval}
\pgfcalendarjuliantoweekday{\the\stampjdval}{\stampwdval}
\newcommand{\timestampzh}{%
  \songfont 最後修改：\the\year 年\the\month 月\the\day 日%
  % （星期\stampwkzh{\the\stampwdval}）%
  % \space\stamphh:\stampmm
  }

\newcommand{\documentdatezh}{\the\year 年 \the\month 月 \the\day 日}
\newcommand{\documentdatede}{\the\day. \stampmonthde{\the\month} \the\year}
\newcommand{\bverfgetranslationnote}{%
  {\songfont 翻譯說明：初譯使用 Claude Code 與 GPT 協助迭代；仍請以德文原文為準。}%
}
\newcommand{\bverfgecourseinfo}{%
  {\songfont 課程：114 學年度第 2 學期；憲法專題研究（四）；授課老師：湯德宗教授；導讀日期：115 年 6 月 1 日。}%
}
\newcommand{\bverfgeeditorinfo}{%
  {\songfont 編者：王逸帆（R10A21126），國立台灣大學法律系研究所碩士班。}%
}
\newcommand{\bverfgetitleversionline}{%
  \ifbilingualcolumns
    Fassung \documentversion\enspace\textperiodcentered\enspace Stand: \documentdatede
    \quad
    {\songfont 版本 \documentversion\enspace\textperiodcentered\enspace 修訂日期：\documentdatezh\par}%
  \else
    \ifshowtranslation
      {\songfont 版本 \documentversion\enspace\textperiodcentered\enspace 修訂日期：\documentdatezh\par}%
    \else
      Fassung \documentversion\enspace\textperiodcentered\enspace Stand: \documentdatede\par
    \fi
  \fi
}
\newcommand{\bverfgetitlecornerinfo}{%
  \ifshowtranslation
    \vspace{0.72em}%
    \noindent\hfill
    \begin{minipage}{0.66\textwidth}
      \raggedleft\tiny\color{pendingink}\songfont
      \bverfgecourseinfo\par
      \bverfgeeditorinfo
    \end{minipage}%
  \fi
}
\newcommand{\bverfgetitlemeta}{%
  \vfill
  \begin{center}%
    \footnotesize\color{pendingink}%
    \bverfgetitleversionline
    \ifshowtranslation
      \vspace{0.28em}{\scriptsize\bverfgetranslationnote\par}%
    \fi
  \end{center}%
  \bverfgetitlecornerinfo
}

\pagestyle{fancy}
\fancyhf{}
\renewcommand{\headrulewidth}{0pt}
\renewcommand{\footrulewidth}{0pt}
\fancyfoot[C]{\thepage}
\ifbilingualcolumns
  \fancyfoot[R]{\scriptsize\color{pendingink}\timestampzh}%
\else
  \ifshoworiginal
  \else
    \fancyfoot[R]{\scriptsize\color{pendingink}\timestampzh\hspace{0.45cm}}%
  \fi
\fi
\setmainfont{Times New Roman}
\setsansfont{Helvetica Neue}
\setmonofont{Menlo}
\IfFileExists{/Users/iw/Library/Fonts/kaiu.ttf}{
  \setCJKmainfont[Path=/Users/iw/Library/Fonts/,AutoFakeBold=4]{kaiu.ttf}
}{
  \IfFontExistsTF{標楷體}{
    \setCJKmainfont[AutoFakeBold=4]{標楷體}
  }{
    \setCJKmainfont{Songti TC}
  }
}
\IfFontExistsTF{Songti TC}{
  \setCJKfamilyfont{song}{Songti TC}
  \setCJKmonofont{Songti TC}
}{
  \setCJKfamilyfont{song}[Path=/Users/iw/Library/Fonts/,AutoFakeBold=4]{kaiu.ttf}
  \setCJKmonofont[Path=/Users/iw/Library/Fonts/,AutoFakeBold=4]{kaiu.ttf}
}
\newcommand{\songfont}{\CJKfamily{song}}

% ===== 封面／目錄專用打字機字體（僅限封面、目錄頁使用，不影響正文）=====
\IfFontExistsTF{Radio Newsman}{%
  \newfontfamily\bverfgetitlede{Radio Newsman}%
}{%
  \newfontfamily\bverfgetitlede{Courier New}%
}
\IfFontExistsTF{Erikas Farbband}{%
  \newfontfamily\bverfgebodyde{Erikas Farbband}%
}{%
  \let\bverfgebodyde\bverfgetitlede
}
\IfFontExistsTF{Maoken Heavy Labourer}{%
  \setCJKfamilyfont{bverfgetitlecjk}{Maoken Heavy Labourer}%
}{%
  \setCJKfamilyfont{bverfgetitlecjk}{Songti TC}%
}
\IfFontExistsTF{Huiwen-mincho}{%
  \setCJKfamilyfont{bverfgebodycjk}{Huiwen-mincho}%
}{%
  \setCJKfamilyfont{bverfgebodycjk}{Songti TC}%
}
\newcommand{\bverfgetitlezh}{\bverfgetitlede\CJKfamily{bverfgetitlecjk}}
\newcommand{\bverfgebodyzh}{\bverfgebodyde\CJKfamily{bverfgebodycjk}}

\setstretch{1.35}
\setlist{itemsep=0.2em, topsep=0.4em}
\setlength{\parindent}{0pt}
\setlength{\parskip}{0.45em}
\raggedbottom
\setcounter{secnumdepth}{0}
\setcounter{tocdepth}{2}

\newcommand{\lawboxsourcefont}{\footnotesize}
\newcommand{\lawboxtransfont}{\footnotesize}
\newcommand{\lawboxsourcestretch}{1.02}
\newcommand{\lawboxtransstretch}{0.92}

\titleformat{\section}{\centering\Large\bfseries\songfont}{\thesection}{1em}{}
\titleformat{\subsection}{\large\bfseries\songfont}{\thesubsection}{1em}{}
\titleformat{\subsubsection}{\normalsize\bfseries\songfont}{\thesubsubsection}{1em}{}
\titleformat{\paragraph}{\normalsize\bfseries\songfont}{\theparagraph}{1em}{}
\titlespacing*{\section}{0pt}{1.8em}{1.0em}
\titlespacing*{\subsection}{0pt}{1.2em}{0.6em}
\titlespacing*{\subsubsection}{0pt}{1.0em}{0.45em}
\titlespacing*{\paragraph}{0pt}{0.6em}{0.4em}

\definecolor{noteblue}{HTML}{A9C9F5}
\definecolor{noteteal}{HTML}{AEE1D6}
\definecolor{notegray}{HTML}{D7DADF}
\definecolor{caseink}{HTML}{000000}
\definecolor{casepaper}{HTML}{FCFEFD}
\definecolor{sourcepaper}{HTML}{FAFBF8}
\definecolor{sourceline}{HTML}{D7DED8}
\definecolor{sourceink}{HTML}{000000}
\definecolor{pendingink}{HTML}{000000}

\tcbset{
  caseboxbase/.style={
    enhanced,
    breakable,
    width=0.985\linewidth,
    colback=casepaper,
    coltitle=caseink,
    fonttitle=\bfseries\songfont,
    boxrule=0.28pt,
    arc=1.2mm,
    left=2.4mm,
    right=2.4mm,
    top=1.5mm,
    bottom=1.5mm,
    boxsep=1.5mm,
    before skip=0.65em,
    after skip=0.65em,
    before upper={\setlength{\parindent}{0pt}\setlength{\parskip}{0.35em}},
  }
}

\newcommand{\bilingualstart}{}
\newcommand{\bilingualstop}{}
\ifbilingualcolumns
  \renewcommand{\bilingualstart}{\small\setlength{\columnsep}{1.25cm}\setlength{\columnseprule}{0.12pt}\columnratio{0.5,0.5}\multicolumnfootnotes\flushbottom\sloppy\interlinepenalty=80\clubpenalty=800\widowpenalty=800\displaywidowpenalty=800\begin{paracol}{2}\global\intranslationfalse\global\firstbilingualsectiontrue{\footnotesize\color{sourceink}Deutsch\par}\switchcolumn{\footnotesize\songfont\color{sourceink}中文譯文\par}\switchcolumn*}
  \renewcommand{\bilingualstop}{\ifintranslation\switchcolumn*\global\intranslationfalse\fi\end{paracol}\clearpage\raggedbottom}
\fi
\newif\ifrandnumbers
\randnumbersfalse
\newcounter{randnumber}
\newcounter{randnumberlocal}
\newcounter{lastsourcerandstart}
\newif\iflastsourcehasrandnumbers
\lastsourcehasrandnumbersfalse
\newif\ifinlawbox
\inlawboxfalse
\newif\iflawboxhaspendingrn
\lawboxhaspendingrnfalse
\newcommand{\lawboxpendingrn}{}
\newcommand{\startrandnumbers}{\global\randnumberstrue\setcounter{randnumber}{1}\global\lastsourcehasrandnumbersfalse}
\newcommand{\stoprandnumbers}{\global\randnumbersfalse\global\lastsourcehasrandnumbersfalse}
\newlength{\randnumberwidth}
\setlength{\randnumberwidth}{0.95cm}
\newlength{\randnumberrightoffset}
\setlength{\randnumberrightoffset}{0.85cm}
\newlength{\randnumberlawboxrightoffset}
\setlength{\randnumberlawboxrightoffset}{1.40cm}
\newcommand{\randnumberbox}[1]{%
  \leavevmode
  \makebox[0pt][l]{%
    \smash{%
      \tikz[remember picture,overlay,baseline=(rnmark.base)]{%
        \coordinate (rnmark) at (0,0);%
        \ifinlawbox
          \path let \p1=(current page.east), \p2=(rnmark.base) in
            node[
              anchor=base east,
              inner sep=0pt,
              outer sep=0pt,
              text width=\randnumberwidth,
              align=right,
              text=pendingink,
              font=\normalfont\mdseries\scriptsize\ttfamily
            ] at (\x1-\randnumberlawboxrightoffset,\y2) {{\normalfont\mdseries\scriptsize\ttfamily #1}};%
        \else
          \path let \p1=(current page.east), \p2=(rnmark.base) in
            node[
              anchor=base east,
              inner sep=0pt,
              outer sep=0pt,
              text width=\randnumberwidth,
              align=right,
              text=pendingink,
              font=\normalfont\mdseries\scriptsize\ttfamily
            ] at (\x1-\randnumberrightoffset,\y2) {{\normalfont\mdseries\scriptsize\ttfamily #1}};%
        \fi
      }%
    }%
  }%
  \ignorespaces
}
\newcommand{\lawboxstorerandnumber}[1]{%
  \gdef\lawboxpendingrn{#1}%
  \global\lawboxhaspendingrntrue
  \ignorespaces
}
\newcommand{\lawboxuserandnumber}{%
  \iflawboxhaspendingrn
    \randnumberbox{\lawboxpendingrn}%
    \global\lawboxhaspendingrnfalse
  \fi
}
\newcommand{\sourcern}[1]{%
  \ifrandnumbers
    \ifshoworiginal
      \ifshowtranslation\else
        \ifinlawbox\lawboxstorerandnumber{#1}\else\randnumberbox{#1}\fi
      \fi
    \fi
  \fi
}
\newcommand{\transrn}[1]{%
  \ifrandnumbers
    \ifshowtranslation
      \ifinlawbox\lawboxstorerandnumber{#1}\else\randnumberbox{#1}\fi
    \fi
  \fi
}
% ===== 課堂模式：德文原文縮寫註解 =====
% 日常切換只改各判決檔的一行：
%   \classroommodeon        開啟課堂版；註解掉這行即回到一般版。
% 模板預設：
%   1. 課堂模式關閉。
%   2. 課堂模式開啟時，縮寫註解包含「德文全稱＋中文譯名」。
%   3. 課堂模式開啟時，GG/條文編者註仍會自動顯示。
% 進階微調才需要在單案檔覆寫：
%   \classroomabbrzhoff     縮寫註解僅顯示德文全稱。
%   \showeditornotestrue    一般版也顯示編者註。
% 註解策略：同一實際 PDF 頁面中，同一縮寫只在第一次出現處加註。
% 這裡用 zref-abspage 的絕對頁序，不用 \thepage，避免文件內頁碼重置時誤判。
\newif\ifclassroommode
\newif\ifclassroomabbrzh
\classroommodefalse
\classroomabbrzhtrue
\newcommand{\classroommodeon}{\classroommodetrue}
\newcommand{\classroommodeoff}{\classroommodefalse}
\newcommand{\classroomabbrzhon}{\classroomabbrzhtrue}
\newcommand{\classroomabbrzhoff}{\classroomabbrzhfalse}
\newcommand{\abbrnotelabel}{\emph{Abk.}: }
\newcommand{\abbrnote}[3]{%
  \ifclassroommode
    \ifshoworiginal
      \ifintranslation\else
        \footnote{\normalfont\footnotesize\abbrnotelabel #2\ifclassroomabbrzh；{\songfont #3}\fi}%
      \fi
    \fi
  \fi
}
\newcommand{\abbrnoteonce}[4]{%
  #2%
  \ifclassroommode
    \ifshoworiginal
      \ifintranslation\else
        \begingroup
          \edef\abbrcurrentabspage{\number\numexpr\value{abspage}+1\relax}%
          \ifcsname abbrnoteseen@#1@\abbrcurrentabspage\endcsname\else
            \global\expandafter\let\csname abbrnoteseen@#1@\abbrcurrentabspage\endcsname\empty
            \abbrnote{#1}{#3}{#4}%
          \fi
        \endgroup
      \fi
    \fi
  \fi
}
\newcommand{\abbrAbs}{\abbrnoteonce{abs}{Abs.}{Absatz}{項}}
\newcommand{\abbrAAO}{\abbrnoteonce{aao}{a.a.O.}{am angegebenen Ort}{前引處}}
\newcommand{\abbrAE}{\abbrnoteonce{ae}{AE}{Alternativ-Entwurf}{替代草案}}
\newcommand{\abbrArt}{\abbrnoteonce{art}{Art.}{Artikel}{條}}
\newcommand{\abbrAufl}{\abbrnoteonce{aufl}{Aufl.}{Auflage}{版}}
\newcommand{\abbrBd}{\abbrnoteonce{bd}{Bd.}{Band}{卷}}
\newcommand{\abbrBGBl}{\abbrnoteonce{bgbl}{BGBl.}{Bundesgesetzblatt}{聯邦法律公報}}
\newcommand{\abbrBGHSt}{\abbrnoteonce{bghst}{BGHSt}{Entscheidungen des Bundesgerichtshofes in Strafsachen}{聯邦普通法院刑事判決彙編}}
\newcommand{\abbrBRDrucks}{\abbrnoteonce{brdrucks}{BRDrucks.}{Bundesrats-Drucksache}{聯邦參議院文件}}
\newcommand{\abbrBTDrucks}{\abbrnoteonce{btdrucks}{BTDrucks.}{Bundestags-Drucksache}{聯邦議院文件}}
\newcommand{\abbrBundesgesetzbl}{\abbrnoteonce{bundesgesetzbl}{Bundesgesetzbl.}{Bundesgesetzblatt}{聯邦法律公報}}
\newcommand{\abbrBvF}{\abbrnoteonce{bvf}{BvF}{Bundesverfassungsgerichtsverfahren}{聯邦憲法法院程序案號}}
\newcommand{\abbrBvL}{\abbrnoteonce{bvl}{BvL}{Bundesverfassungsgerichtsverfahren}{聯邦憲法法院程序案號}}
\newcommand{\abbrBvR}{\abbrnoteonce{bvr}{BvR}{Bundesverfassungsgerichtsverfahren}{聯邦憲法法院程序案號}}
\newcommand{\abbrBVerfG}{\abbrnoteonce{bverfg}{BVerfG}{Bundesverfassungsgericht}{聯邦憲法法院}}
\newcommand{\abbrBVerfGE}{\abbrnoteonce{bverfge}{BVerfGE}{Entscheidungen des Bundesverfassungsgerichts}{聯邦憲法法院判決集}}
\newcommand{\abbrBVerfGG}{\abbrnoteonce{bverfgg}{BVerfGG}{Bundesverfassungsgerichtsgesetz}{聯邦憲法法院法}}
\newcommand{\abbrDDR}{\abbrnoteonce{ddr}{DDR}{Deutsche Demokratische Republik}{德意志民主共和國；東德}}
\newcommand{\abbrDJT}{\abbrnoteonce{djt}{DJT}{Deutscher Juristentag}{德國法學家大會}}
\newcommand{\abbrDr}{\abbrnoteonce{dr}{Dr.}{Doktor}{Dr.}}
\newcommand{\abbrEuGRZ}{\abbrnoteonce{eugrz}{EuGRZ}{Europäische Grundrechte-Zeitschrift}{歐洲基本權期刊}}
\newcommand{\abbrEV}{\abbrnoteonce{ev}{EV}{Einigungsvertrag}{統一條約}}
\newcommand{\abbrF}{\abbrnoteonce{f}{f.}{folgende}{以下一頁；下一段}}
\newcommand{\abbrFF}{\abbrnoteonce{ff}{ff.}{fortfolgende}{以下數頁；以下各段}}
\newcommand{\abbrGBlDDR}{\abbrnoteonce{gblddr}{GBl. DDR}{Gesetzblatt der Deutschen Demokratischen Republik}{德意志民主共和國法律公報}}
\newcommand{\abbrGG}{\abbrnoteonce{gg}{GG}{Grundgesetz}{基本法}}
\newcommand{\abbrGVBl}{\abbrnoteonce{gvbl}{GVBl.}{Gesetz- und Verordnungsblatt}{法律及命令公報}}
\newcommand{\abbrJR}{\abbrnoteonce{jr}{JR}{Juristische Rundschau}{法學評論}}
\newcommand{\abbrJZ}{\abbrnoteonce{jz}{JZ}{JuristenZeitung}{法學家報}}
\newcommand{\abbrIVm}{\abbrnoteonce{ivm}{i.V.m.}{in Verbindung mit}{結合；連同}}
\newcommand{\abbrMdB}{\abbrnoteonce{mdb}{MdB}{Mitglied des Bundestages}{聯邦議院議員}}
\newcommand{\abbrMwN}{\abbrnoteonce{mwn}{m.w.N.}{mit weiteren Nachweisen}{附進一步文獻／判例出處}}
\newcommand{\abbrNF}{\abbrnoteonce{nf}{n.F.}{neue Fassung}{新版本}}
\newcommand{\abbrAF}{\abbrnoteonce{af}{a.F.}{alte Fassung}{舊版本}}
\newcommand{\abbrNr}{\abbrnoteonce{nr}{Nr.}{Nummer}{款、目或號，依語境}}
\newcommand{\abbrProf}{\abbrnoteonce{prof}{Prof.}{Professor}{教授}}
\newcommand{\abbrRdnr}{\abbrnoteonce{rdnr}{Rdnr.}{Randnummer}{邊碼；段碼}}
\newcommand{\abbrRdnrn}{\abbrnoteonce{rdnrn}{Rdnrn.}{Randnummern}{邊碼；段碼}}
\newcommand{\abbrRGBl}{\abbrnoteonce{rgbl}{RGBl.}{Reichsgesetzblatt}{帝國法律公報}}
\newcommand{\abbrRGSt}{\abbrnoteonce{rgst}{RGSt}{Entscheidungen des Reichsgerichts in Strafsachen}{帝國法院刑事判決彙編}}
\newcommand{\abbrRspr}{\abbrnoteonce{rspr}{Rspr.}{Rechtsprechung}{判例；司法實務}}
\newcommand{\abbrRVO}{\abbrnoteonce{rvo}{RVO}{Reichsversicherungsordnung}{帝國保險法}}
\newcommand{\abbrS}{\abbrnoteonce{s}{S.}{Seite}{頁}}
\newcommand{\abbrSFHG}{\abbrnoteonce{sfhg}{SFHG}{Schwangeren- und Familienhilfegesetz}{孕婦及家庭扶助法}}
\newcommand{\abbrSGBV}{\abbrnoteonce{sgbv}{SGB V}{Fünftes Buch Sozialgesetzbuch}{社會法典第五編}}
\newcommand{\abbrStAG}{\abbrnoteonce{staeg}{StÄG}{Strafrechtsänderungsgesetz}{刑法修正法}}
\newcommand{\abbrStenBer}{\abbrnoteonce{stenber}{StenBer.}{Stenographischer Bericht}{速記紀錄}}
\newcommand{\abbrStGB}{\abbrnoteonce{stgb}{StGB}{Strafgesetzbuch}{刑法}}
\newcommand{\abbrStREG}{\abbrnoteonce{streg}{StREG}{Strafrechtsreform-Ergänzungsgesetz}{刑法改革補充法}}
\newcommand{\abbrStrRG}{\abbrnoteonce{strrg}{StrRG}{Strafrechtsreformgesetz}{刑法改革法}}
\newcommand{\abbrUS}{\abbrnoteonce{us}{U.S.}{United States Reports}{美國最高法院判決彙編}}
\newcommand{\abbrVgl}{\abbrnoteonce{vgl}{vgl.}{vergleiche}{參見}}
\newcommand{\abbrVVDStRL}{\abbrnoteonce{vvdstrl}{VVDStRL}{Veröffentlichungen der Vereinigung der Deutschen Staatsrechtslehrer}{德國公法教師協會論文集}}
\newcommand{\abbrWP}{\abbrnoteonce{wp}{WP.}{Wahlperiode}{屆期}}
\newcommand{\abbrWp}{\abbrnoteonce{wp}{Wp.}{Wahlperiode}{屆期}}
\newcommand{\abbrZStrW}{\abbrnoteonce{zstrw}{ZStrW}{Zeitschrift für die gesamte Strafrechtswissenschaft}{整體刑法學期刊}}
\newcommand{\recordsourcerandstart}{%
  \ifrandnumbers
    \setcounter{lastsourcerandstart}{\value{randnumber}}%
    \global\lastsourcehasrandnumberstrue
  \else
    \global\lastsourcehasrandnumbersfalse
  \fi
}
\newlength{\biparagraphskip}
\setlength{\biparagraphskip}{0.52em}
\newlength{\lawboxblockbeforeskip}
\setlength{\lawboxblockbeforeskip}{0.72em}
\newlength{\lawboxblockafterskip}
\setlength{\lawboxblockafterskip}{0.92em}
\newif\iflawboxpartendedblock
\lawboxpartendedblockfalse
\newcommand{\sourceparstyle}{\footnotesize\color{sourceink}\setstretch{1.12}\RaggedRight\emergencystretch=3em\setlength{\parskip}{0.42em}\obeylines\selectfont}
\newcommand{\transparstyle}{\songfont\color{caseink}\setstretch{1.12}\setlength{\parindent}{0pt}\setlength{\parskip}{0.42em}}
\newcommand{\bilingualpairskip}{%
  \par\addvspace{\biparagraphskip}%
  \switchcolumn
  \par\addvspace{\biparagraphskip}%
  \switchcolumn*%
  \global\intranslationfalse
}
\newcommand{\bilingualboxpairskip}{%
  \par
  \switchcolumn
  \par
  \switchcolumn*%
  \global\intranslationfalse
}
\newcommand{\bilinguallawboxblockskip}{%
  \switchcolumn*%
  \par\addvspace{\lawboxblockafterskip}%
  \switchcolumn
  \par\addvspace{\lawboxblockafterskip}%
  \switchcolumn*%
  \global\intranslationfalse
}
\newcommand{\bikeepwithnext}[1][4]{%
  \ifbilingualcolumns
    \syncleft
    \Needspace{#1\baselineskip}%
    \switchcolumn
    \Needspace{#1\baselineskip}%
    \switchcolumn*%
    \global\intranslationfalse
  \else
    \Needspace{#1\baselineskip}%
  \fi
}
\NewEnviron{sourcepara}[1][]{
  \recordsourcerandstart
  \ifshoworiginal
    \ifbilingualcolumns
      \ifintranslation\switchcolumn*\global\intranslationfalse\fi
    \else
      \Needspace{5\baselineskip}
    \fi
    \ifbilingualcolumns\else\par\addvspace{0.35em}\fi
    {\sourceparstyle
    \noindent\BODY\par}
    \ifbilingualcolumns\else\addvspace{0.25em}\fi
    \ifbilingualcolumns
      \switchcolumn
      \global\intranslationtrue
    \fi
  \else
    \ifrandnumbers
      \begingroup
        \setbox0=\vbox{\hsize=\linewidth\sourceparstyle\noindent\BODY\par}%
      \endgroup
    \fi
  \fi
}
\NewEnviron{transpara}{
  \ifshowtranslation
    \setcounter{randnumberlocal}{\value{lastsourcerandstart}}%
    \ifbilingualcolumns
      \ifintranslation\else\switchcolumn\global\intranslationtrue\fi
      {\transparstyle\setlength{\leftskip}{0.55em}\setlength{\rightskip}{0.15em plus 1em}\addtolength{\linewidth}{-0.55em}\BODY\par}
      \switchcolumn*
      \bilingualpairskip
    \else
      {\transparstyle\BODY\par}
    \fi
  \else
    \ifbilingualcolumns
      \ifintranslation\switchcolumn*\global\intranslationfalse\fi
    \fi
  \fi
}
\NewEnviron{sourceboxpara}[1][]{
  \global\lawboxpartendedblockfalse
  \recordsourcerandstart
  \ifshoworiginal
    \ifbilingualcolumns
      \ifintranslation\switchcolumn*\global\intranslationfalse\fi
    \else
      \Needspace{3\baselineskip}
    \fi
    {\sourceparstyle
    \BODY\par}
    \ifbilingualcolumns\else
      \iflawboxpartendedblock\addvspace{\lawboxblockafterskip}\fi
    \fi
    \ifbilingualcolumns
      \switchcolumn
      \global\intranslationtrue
    \fi
  \fi
}
\NewEnviron{transboxpara}{
  \global\lawboxpartendedblockfalse
  \ifshowtranslation
    \setcounter{randnumberlocal}{\value{lastsourcerandstart}}%
    \ifbilingualcolumns
      \ifintranslation\else\switchcolumn\global\intranslationtrue\fi
      {\transparstyle\BODY\par}
      \iflawboxpartendedblock
        \bilinguallawboxblockskip
      \else
        \switchcolumn*
        \bilingualboxpairskip
      \fi
    \else
      {\transparstyle\BODY\par}
      \iflawboxpartendedblock\addvspace{\lawboxblockafterskip}\fi
    \fi
  \else
    \ifbilingualcolumns
      \ifintranslation\switchcolumn*\global\intranslationfalse\fi
    \fi
  \fi
}
\newcommand{\leitsatznum}[1]{\noindent\hangindent=3.0em\hangafter=1\makebox[2.25em][r]{\bfseries #1.}\hspace{0.55em}\ignorespaces}
\newcommand{\syncleft}{\ifbilingualcolumns\ifintranslation\switchcolumn*\global\intranslationfalse\fi\fi}
\pretocmd{\section}{\syncleft\clearpage}{}{}
\pretocmd{\subsection}{\syncleft}{}{}
\pretocmd{\subsubsection}{\syncleft}{}{}
\newcommand{\biheadingfont}{\songfont}
\newcommand{\bitocentry}[2]{%
  \ifshoworiginal
    \ifshowtranslation
      \ifstrequal{#1}{#2}{#1}{#1\space #2}%
    \else
      #1%
    \fi
  \else
    #2%
  \fi
}
\pdfstringdefDisableCommands{%
  \def\bitocentry#1#2{#1}%
}
\newcommand{\biheading}[7]{%
  \syncleft#1\bikeepwithnext[#3]\phantomsection\addcontentsline{toc}{#2}{\bitocentry{#6}{#7}}%
  \ifbilingualcolumns
    {#4 #6\par}%
    \switchcolumn
    {#4\biheadingfont #7\par}%
    \switchcolumn*%
    \global\intranslationfalse
    \par\addvspace{#5}%
    \switchcolumn
    \par\addvspace{#5}%
    \switchcolumn*%
  \else
    \ifshoworiginal{#4 #6\par}\fi
    \ifshowtranslation{#4\biheadingfont #7\par}\fi
    \addvspace{#5}%
  \fi
}
\newcommand{\termsectionheading}[1]{%
  \par\Needspace{9\baselineskip}%
  \vspace{0.65em}%
  {\songfont\bfseries #1\par}%
  \vspace{0.2em}%
}
% 章節換頁：
%   雙語 paracol 欄內不要用 \clearpage；它會在同步兩欄時吐出只有欄線/頁碼的空頁。
%   \newpage 足以讓下一個 section 起新頁，又不會強制清空所有待排材料。
%   單語模式不在 paracol 內，仍可用 \clearpage。
\newcommand{\bisectionpagebreak}{\ifbilingualcolumns\newpage\else\clearpage\fi}
\newcommand{\bisection}[2]{%
  \biheading{\iffirstbilingualsection\global\firstbilingualsectionfalse\else\bisectionpagebreak\fi}{section}{6}{\centering\Large\bfseries}{0.8em}{#1}{#2}%
}
\newcommand{\bisubsection}[2]{%
  \biheading{}{subsection}{5}{\large\bfseries}{0.55em}{#1}{#2}%
}
\newcommand{\bisubsubsection}[2]{%
  \biheading{}{subsubsection}{4}{\normalsize\bfseries}{0.45em}{#1}{#2}%
}
\newcommand{\bicompositeheading}[4]{%
  \biheading{}{subsection}{5}{\large\bfseries}{0.16em}{#1}{#2}%
  \biheading{}{subsubsection}{3}{\normalsize\bfseries}{0.45em}{#3}{#4}%
}
\newif\iflawboxfirstitem
\lawboxfirstitemfalse
\newcommand{\lawboxprespace}[1]{%
  \par
  \iflawboxfirstitem
    \global\lawboxfirstitemfalse
  \else
    \addvspace{#1}%
  \fi
}
\newtcolorbox{lawboxinner}{
  caseboxbase,
  colframe=sourceline!85!black,
  colback=white,
  left=1.05em,
  right=1.05em,
  top=0.62em,
  bottom=0.62em,
  before upper={\global\inlawboxtrue\global\lawboxfirstitemtrue\global\lawboxhaspendingrnfalse\ifintranslation\lawboxtransfont\else\lawboxsourcefont\fi\RaggedRight\setlength{\parindent}{0pt}\setlength{\parskip}{0pt}\ifintranslation\setstretch{\lawboxtransstretch}\else\setstretch{\lawboxsourcestretch}\fi},
  after upper={\global\lawboxhaspendingrnfalse\global\inlawboxfalse}
}
\tcbset{
  lawboxpartbase/.style={
    enhanced,
    breakable=false,
    width=0.985\linewidth,
    colback=white,
    frame hidden,
    boxrule=0pt,
    arc=0pt,
    left=1.05em,
    right=1.05em,
    boxsep=1.5mm,
    before skip=0pt,
    after skip=0pt,
    before upper={\global\inlawboxtrue\global\lawboxfirstitemtrue\global\lawboxhaspendingrnfalse\ifintranslation\lawboxtransfont\else\lawboxsourcefont\fi\RaggedRight\setlength{\parindent}{0pt}\setlength{\parskip}{0pt}\ifintranslation\setstretch{\lawboxtransstretch}\else\setstretch{\lawboxsourcestretch}\fi},
    after upper={\global\lawboxhaspendingrnfalse\global\inlawboxfalse},
    borderline west={0.28pt}{0pt}{sourceline!85!black},
    borderline east={0.28pt}{0pt}{sourceline!85!black},
  },
  lawboxpartfirst/.style={
    lawboxpartbase,
    top=0.62em,
    bottom=0.04em,
    before skip=\lawboxblockbeforeskip,
    borderline north={0.28pt}{0pt}{sourceline!85!black},
  },
  lawboxpartmiddle/.style={
    lawboxpartbase,
    top=0.04em,
    bottom=0.04em,
  },
  lawboxpartlast/.style={
    lawboxpartbase,
    top=0.04em,
    bottom=0.62em,
    after skip=0pt,
    borderline south={0.28pt}{0pt}{sourceline!85!black},
  }
}
\newtcolorbox{lawboxpartinner}[1]{#1}
\newtcolorbox{termboxinner}{
  caseboxbase,
  colframe=noteblue!70!black,
  colback=casepaper,
  before upper={\songfont\setlength{\parindent}{0pt}\setlength{\parskip}{0.35em}}
}
\newtcolorbox{deboxinner}{
  caseboxbase,
  colframe=notegray!72!black,
  colback=white,
  left=1.05em,
  right=1.05em,
  top=0.62em,
  bottom=0.62em,
  before upper={\global\inlawboxtrue\global\lawboxfirstitemtrue\global\lawboxhaspendingrnfalse\small\setlength{\parindent}{0pt}\setlength{\parskip}{0.35em}},
  after upper={\global\lawboxhaspendingrnfalse\global\inlawboxfalse}
}
\NewEnviron{lawbox}{
  \begin{lawboxinner}\BODY\end{lawboxinner}
}
\NewEnviron{lawboxpart}[2]{
  \begin{lawboxpartinner}{lawboxpart#1,equal height group=#2}\BODY\end{lawboxpartinner}%
  \ifstrequal{#1}{last}{\global\lawboxpartendedblocktrue}{}%
}
\NewEnviron{termbox}{
  \par\addvspace{0.55em}
  {\color{noteblue!70!black}\rule{\linewidth}{0.28pt}}\par
  \begingroup
    \songfont\small\color{caseink}
    \setlength{\parindent}{0pt}
    \setlength{\parskip}{0.24em}
    \BODY
    \par
  \endgroup
  {\color{noteblue!70!black}\rule{\linewidth}{0.28pt}}\par
  \addvspace{0.55em}
}
\NewEnviron{debox}{
  \begin{deboxinner}\BODY\end{deboxinner}
}
\providecommand{\paragraphmark}[1]{}
\newcommand{\lawboxtitleafter}{\addvspace{0.06em}}
\newcommand{\lawtitle}[1]{\lawboxprespace{0.22em}{\lawboxtransfont\bfseries\lawboxuserandnumber #1}\par\lawboxtitleafter}
\newcommand{\absno}[2]{\lawboxprespace{0.04em}\noindent\lawboxuserandnumber\hangindent=2.6em\hangafter=1\makebox[2.6em][l]{(#1)}#2\par}
\newcommand{\nrno}[2]{\lawboxprespace{0pt}\noindent\lawboxuserandnumber\hspace*{2.6em}\hangindent=4.4em\hangafter=1\makebox[1.8em][l]{#1.}#2\par}
\newcommand{\letterno}[2]{\lawboxprespace{0pt}\noindent\lawboxuserandnumber\hspace*{4.4em}\hangindent=6.0em\hangafter=1\makebox[1.6em][l]{#1)}#2\par}
\newcommand{\sourcelawtitle}[1]{\lawboxprespace{0.22em}{\lawboxsourcefont\bfseries\lawboxuserandnumber #1}\par\lawboxtitleafter}
\newcommand{\sourcelawsubtitle}[1]{\lawboxprespace{0.04em}{\lawboxsourcefont\bfseries\lawboxuserandnumber #1}\par\lawboxtitleafter}
\newcommand{\sourceabsno}[2]{\lawboxprespace{0.04em}\noindent\lawboxuserandnumber\hangindent=2.45em\hangafter=1\makebox[2.45em][l]{(#1)}#2\par}
\newcommand{\sourcenrno}[2]{\lawboxprespace{0pt}\noindent\lawboxuserandnumber\hspace*{2.45em}\hangindent=4.25em\hangafter=1\makebox[1.8em][l]{#1.}#2\par}
\newcommand{\sourceletterno}[2]{\lawboxprespace{0pt}\noindent\lawboxuserandnumber\hspace*{4.25em}\hangindent=5.85em\hangafter=1\makebox[1.6em][l]{#1)}#2\par}
\newcommand{\sourcequotepara}[1]{\lawboxprespace{0.04em}\noindent\lawboxuserandnumber\hangindent=1.2em\hangafter=1 #1\par}
\newcommand{\sourcelawline}[1]{\lawboxprespace{0.04em}\noindent\lawboxuserandnumber #1\par}
\newcommand{\lawline}[1]{\lawboxprespace{0.04em}\noindent\lawboxuserandnumber #1\par}
\newcommand{\pendingtranslation}{\par{\songfont\footnotesize\color{pendingink}（中文譯文待補。）}\par}
\newcommand{\tocpart}[1]{\par\addvspace{0.52em}{\footnotesize\bfseries\color{pendingink}#1\par}\addvspace{0.16em}}
\newcommand{\tocdashline}{\par\addvspace{0.48em}\noindent{\color{notegray}\xleaders\hbox to 0.82em{\hfil\rule{0.42em}{0.28pt}\hfil}\hskip\linewidth}\par\addvspace{0.38em}}
% \tocpanel{font-cmd}{content}：將目錄置中於所在欄寬（雙語欄適用）。
% A3 橫式使用 0.58\linewidth，A4 橫式使用 0.84\linewidth，
% 使兩種紙張下目錄欄內容實際寬度相近（均約 100–105 mm）。
\newcommand{\tocpanel}[2]{%
  \makebox[\linewidth][c]{%
    \ifx\bilingualpaper\bilingualpaperaThree
      \begin{minipage}{0.58\linewidth}%
    \else
      \begin{minipage}{0.84\linewidth}%
    \fi
    \toccolstyle
    #1#2%
    \end{minipage}%
  }%
}
% ===== 首頁簡目錄的兩層 description list 縮排 =====
% enumitem 的 description list 可把每一列拆成：
%   [label]  labelsep  body text
% 這裡的「起點」都以所在 minipage/欄位的左緣為 0。
%
% 核心規則：
%   labelindent = label 欄位從左緣開始的位置。
%   labelwidth  = label 欄位本身寬度；標籤太長（如 Abw. Meinung）就加大它。
%   labelsep    = label 欄位右緣到正文的距離。
%   leftmargin  = 正文 body text 的起點。判斷層級視覺時主要看這個。
%
% 所以：
%   想讓「Begründung / 判決理由」整列正文往右或往左：改 tocitems 的 leftmargin。
%   想讓「A. Verfahren und Gegenstand / A. 程序與審查標的」正文往右或往左：
%     改 tocsubitems 的 leftmargin。
%   想讓 A./B./C. 這些標籤本身往右或往左，但正文不動：改 tocsubitems 的 labelindent。
%   想讓 A./B./C. 與正文間距變大或變小：改 tocsubitems 的 labelsep。
%
% 層級原則：
%   tocsubitems.leftmargin 必須大於 tocitems.leftmargin，
%   否則子層級正文會看起來比「Gründe / 理由」還靠左。
\newlist{tocitems}{description}{1}
\setlist[tocitems]{%
  labelindent=0pt,    % 第一層 label 欄位起點；0pt 表示貼齊目錄欄左緣。
  leftmargin=2.54cm,  % 第一層正文起點；例如 Begründung / 判決理由 的左緣。
  labelwidth=2.16cm,  % 第一層 label 欄位寬度；需容納 Leitsätze、Abw. Meinung、不同意見等。
  labelsep=0.32cm,    % 第一層 label 與正文的水平間距；只改間距，不改正文起點。
  itemsep=0.28em,     % 第一層各項之間的垂直距離。
  topsep=0.20em,      % 第一層 list 上下與前後內容的垂直距離。
  parsep=0pt,         % 同一 item 內段落之間的距離；目錄項通常不需要。
  partopsep=0pt,      % list 若緊接段落開始時的額外距離；關閉以免目錄高度飄動。
  align=left,         % label 欄內左對齊；長短標籤不靠右貼齊。
  font=\normalfont\bfseries % label 的字型；正文不受這行影響。
}
\newlist{tocsubitems}{description}{1}
\setlist[tocsubitems]{%
  labelindent=1cm,    % 第二層 label 起點；只控制 A./B./C. 本身的位置。
  leftmargin=3.00cm,  % 第二層正文起點；必須大於 2.54cm，才會比 Begründung / 判決理由更右。
  labelwidth=0.5cm,   % 第二層 label 欄寬；A./B./C. 很短，可小於第一層。
  labelsep=1.5cm,    % 第二層 label 與正文間距；A. 與 Verfahren / 程序 的距離看這裡。
  itemsep=0.12em,     % 第二層各項較密，避免 A.–E. 佔太高。
  topsep=0.04em,      % 第二層 list 與 Gründe / 理由之間的垂直距離。
  parsep=0pt,         % 同一第二層 item 內段落距離；目錄項通常不需要。
  partopsep=0pt,      % 關閉額外段落起始距離，避免版面不穩。
  align=left,         % A./B./C. label 欄內左對齊。
  font=\normalfont\bfseries, % A./B./C. label 加粗；正文不受這行影響。
  before={}           % 不在第二層 list 前自動插入額外內容。
}
\newlist{termitems}{description}{1}
\ifbilingualcolumns
  \setlist[termitems]{%
    leftmargin=5.2cm,
    labelwidth=4.8cm,
    labelsep=0.35cm,
    itemsep=0.16em,
    topsep=0.2em,
    parsep=0pt,
    partopsep=0pt,
    font=\normalfont\bfseries,
    style=nextline
  }
\else
  \setlist[termitems]{%
    leftmargin=0pt,
    labelwidth=0pt,
    labelsep=0pt,
    itemsep=0.24em,
    topsep=0.2em,
    parsep=0pt,
    partopsep=0pt,
    font=\normalfont\bfseries,
    style=nextline
  }
\fi
\newlist{abbritems}{description}{1}
\setlist[abbritems]{%
  leftmargin=2.38cm,
  labelwidth=2.02cm,
  labelsep=0.28cm,
  itemsep=0.16em,
  topsep=0.2em,
  parsep=0pt,
  partopsep=0pt,
  align=left,
  font=\normalfont\bfseries
}
\ifbilingualcolumns
  \newenvironment{termcolumns}{\begin{multicols}{2}}{\end{multicols}}
\else
  \newenvironment{termcolumns}{}{}
\fi

% ===== 編者註腳開關 =====
% 一般版預設不顯示編者註，避免正文腳註過密。
% 若開啟 \classroommodeon，GG/條文編者註仍會自動顯示，無需另開本開關。
% 若一般版也要顯示編者註，可在單案檔 \input 後加 \showeditornotestrue。
\newif\ifshoweditornotes
\showeditornotesfalse

% ===== 目錄排版輔助 =====
\newcommand{\toccolstyle}{\raggedright\arraybackslash}
\newcommand{\tocsingleindent}{\leftskip=1.45cm\rightskip=1.2cm}

% ===== 行內引文環境（德文原文與中文譯文各一，帶左側灰色邊線）=====
\newcommand{\quoteparbreak}{\par\addvspace{0.48em}}
\newcommand{\sourcequoteinner}[1]{%
  \begin{tcolorbox}[
    enhanced, breakable, width=\dimexpr\linewidth-0.8em\relax, frame hidden,
    colback=white, coltext=caseink, boxrule=0pt,
    borderline west={0.55pt}{0.88em}{notegray},
    left=1.78em, right=0.72em, top=0.18em, bottom=0.18em,
    boxsep=0pt, before skip=0.24em, after skip=0.24em,
    before upper={\normalfont\footnotesize\color{caseink}\setlength{\parindent}{0pt}\setlength{\parskip}{0.34em}\raggedright\setlength{\rightskip}{0pt plus 1em}}
  ]%
  #1\par
  \end{tcolorbox}%
}
\newcommand{\transquoteinner}[1]{%
  \begin{tcolorbox}[
    enhanced, breakable, width=\dimexpr\linewidth-0.8em\relax, frame hidden,
    colback=white, coltext=caseink, boxrule=0pt,
    borderline west={0.55pt}{0.88em}{notegray},
    left=1.78em, right=0.72em, top=0.18em, bottom=0.18em,
    boxsep=0pt, before skip=0.24em, after skip=0.24em,
    before upper={\songfont\small\color{caseink}\setlength{\parindent}{0pt}\setlength{\parskip}{0.34em}\raggedright\setlength{\rightskip}{0pt plus 1em}}
  ]%
  #1\par
  \end{tcolorbox}%
}

% ===== 大綱（Gliederung）Rn 輔助命令 =====
\newcommand{\outrn}[2]{\enspace{\normalfont\footnotesize\color{pendingink}(Rn\,#1–#2)}}
\newcommand{\outrnsingle}[1]{\enspace{\normalfont\footnotesize\color{pendingink}(Rn\,#1)}}
\newcommand{\outrnzh}[2]{\enspace{\normalfont\footnotesize\color{pendingink}（第\,#1–#2\,段）}}
\newcommand{\outrnzhs}[1]{\enspace{\normalfont\footnotesize\color{pendingink}（第\,#1\,段）}}
% ===== 大綱雙語對照表格輔助命令（三欄：段碼 │ 德文 │ 中文）=====
% 欄 1：段碼（由欄定義自動格式化：小字、等寬、右對齊、pendingink）
% 欄 2：德文標籤＋正文
% 欄 3：中文標籤＋正文
%
% 無段碼的列（\omainrow、\osecrow），欄 1 以 & 空置。
% 有段碼的列（\osecrowrn、\osubrow、\ointrow），#1 為預格式化字串
%   例：\osubrow{2–49}{I.}{...}{I.}{...}
%       \ointrow{107}{...}{...}
%
% \opartrow：段組標題（橫跨三欄）
\newcommand{\opartrow}[2]{%
  \noalign{\vspace{0.30em}}%
  \multicolumn{3}{@{}l@{}}{%
    \footnotesize\bfseries\color{pendingink}#1\hfill\songfont#2%
  }\\[0.06em]%
}
% \odashrow：虛線分隔（橫跨三欄）
\newcommand{\odashrow}{%
  \noalign{\vspace{0.28em}}%
  \multicolumn{3}{@{}l@{}}{\color{notegray}%
    \xleaders\hbox to 0.82em{\hfil\rule{0.42em}{0.28pt}\hfil}\hskip 0.88\linewidth}%
  \\[0.24em]%
}
% \odissentpart：不同意見書區段標頭。比一般 \opartrow 更像附錄性轉場。
\newcommand{\odissentpart}[2]{%
  \noalign{\vspace{0.42em}}%
  \multicolumn{3}{@{}p{0.88\textwidth}@{}}{%
    \color{notegray}\rule{\linewidth}{0.18pt}%
  }\\[-0.18em]%
  \multicolumn{3}{@{}p{0.88\textwidth}@{}}{%
    \makebox[\linewidth][s]{%
      \footnotesize\bfseries\color{pendingink}#1%
      \hfill{\songfont #2}%
    }%
  }\\[-0.02em]%
}
% \odissentopinion：不同意見書的署名與一行摘要，右欄保留段碼範圍。
\newcommand{\odissentopinion}[5]{%
  \footnotesize\hangindent=0.52cm\hangafter=1%
    {\bfseries\color{caseink}#2}\enspace{\color{notegray}\textperiodcentered}\enspace\color{caseink}#3%
  &\footnotesize\hangindent=0.52cm\hangafter=1%
    {\songfont\bfseries\color{caseink}#4}\enspace{\color{notegray}\textperiodcentered}\enspace\songfont\color{caseink}#5%
  &#1%
  \\[0.12em]%
}
% \omainrow：頂層項目，無段碼（Leitsätze、Urteil、Gründe …）
\newcommand{\omainrow}[4]{%
  \footnotesize\hangindent=1.92cm\hangafter=1%
    \makebox[1.92cm][l]{\bfseries\color{caseink}#1}\color{caseink}#2%
  &\footnotesize\hangindent=1.92cm\hangafter=1%
    \makebox[1.92cm][l]{\bfseries\color{caseink}#3}\songfont\color{caseink}#4%
  &%
  \\[0.15em]%
}
% \omainrowrn：頂層項目，有段碼（不同意見等標籤寬於 \osecrow 框者）
\newcommand{\omainrowrn}[5]{%
  \footnotesize\hangindent=1.92cm\hangafter=1%
    \makebox[1.92cm][l]{\bfseries\color{caseink}#2}\color{caseink}#3%
  &\footnotesize\hangindent=1.92cm\hangafter=1%
    \makebox[1.92cm][l]{\bfseries\color{caseink}#4}\songfont\color{caseink}#5%
  &#1%
  \\[0.15em]%
}
% \osecrow：一級子項，無段碼（A.、C.、D. 等有子項者）
\newcommand{\osecrow}[4]{%
  \footnotesize\hspace{1.10cm}\hangindent=1.98cm\hangafter=1%
    \makebox[0.86cm][l]{\bfseries\color{caseink}#1}\color{caseink}#2%
  &\footnotesize\hspace{1.10cm}\hangindent=1.98cm\hangafter=1%
    \makebox[0.86cm][l]{\bfseries\color{caseink}#3}\songfont\color{caseink}#4%
  &%
  \\[0.11em]%
}
% \osecrowrn：一級子項，有段碼（B.、E. 等完整段落範圍者）
% #1=段碼字串（如 101–106）, #2=DE標籤, #3=DE正文, #4=ZH標籤, #5=ZH正文
\newcommand{\osecrowrn}[5]{%
  \footnotesize\hspace{1.10cm}\hangindent=1.98cm\hangafter=1%
    \makebox[0.86cm][l]{\bfseries\color{caseink}#2}\color{caseink}#3%
  &\footnotesize\hspace{1.10cm}\hangindent=1.98cm\hangafter=1%
    \makebox[0.86cm][l]{\bfseries\color{caseink}#4}\songfont\color{caseink}#5%
  &#1%
  \\[0.11em]%
}
% \osubrow：二級子項，有段碼範圍（I.、II.、A.I. …）
% #1=段碼字串, #2=DE標籤, #3=DE正文, #4=ZH標籤, #5=ZH正文
\newcommand{\osubrow}[5]{%
  \footnotesize\hspace{1.40cm}\hangindent=2.30cm\hangafter=1%
    \makebox[0.90cm][l]{\bfseries\color{caseink}#2}\color{caseink}#3%
  &\footnotesize\hspace{1.40cm}\hangindent=2.30cm\hangafter=1%
    \makebox[0.90cm][l]{\bfseries\color{caseink}#4}\songfont\color{caseink}#5%
  &#1%
  \\[0.09em]%
}
% \ointrow：引入段落，有單一段碼，無標籤
% #1=段碼字串, #2=DE正文, #3=ZH正文
\newcommand{\ointrow}[3]{%
  \footnotesize\hspace{0.52cm}\hangindent=0.52cm\hangafter=1\color{caseink}#2%
  &\footnotesize\hspace{0.52cm}\hangindent=0.52cm\hangafter=1\songfont\color{caseink}#3%
  &#1%
  \\[0.09em]%
}
% \osubsubrow：三級子項（1.、2.、3.）
% #1=段碼字串, #2=DE標籤, #3=DE正文, #4=ZH標籤, #5=ZH正文
\newcommand{\osubsubrow}[5]{%
  \footnotesize\hspace{1.80cm}\hangindent=2.48cm\hangafter=1%
    \makebox[0.68cm][l]{\bfseries\color{caseink}#2}\color{caseink}#3%
  &\footnotesize\hspace{1.80cm}\hangindent=2.48cm\hangafter=1%
    \makebox[0.68cm][l]{\bfseries\color{caseink}#4}\songfont\color{caseink}#5%
  &#1%
  \\[0.08em]%
}
% \osubsubsubrow：四級子項（a）、b））
% #1=段碼字串, #2=DE標籤, #3=DE正文, #4=ZH標籤, #5=ZH正文
\newcommand{\osubsubsubrow}[5]{%
  \footnotesize\hspace{2.10cm}\hangindent=2.68cm\hangafter=1%
    \makebox[0.58cm][l]{\bfseries\color{caseink}#2}\color{caseink}#3%
  &\footnotesize\hspace{2.10cm}\hangindent=2.68cm\hangafter=1%
    \makebox[0.58cm][l]{\bfseries\color{caseink}#4}\songfont\color{caseink}#5%
  &#1%
  \\[0.07em]%
}
% \osubsubsubsubrow：五級子項（aa）、bb））
% 標籤框 0.62cm：足以容納「aa)」在 footnotesize bold 下（約 0.50cm）並留有間距
% #1=段碼字串, #2=DE標籤, #3=DE正文, #4=ZH標籤, #5=ZH正文
\newcommand{\osubsubsubsubrow}[5]{%
  \footnotesize\hspace{2.38cm}\hangindent=3.06cm\hangafter=1%
    \makebox[0.68cm][l]{\bfseries\color{caseink}#2}\color{caseink}#3%
  &\footnotesize\hspace{2.38cm}\hangindent=3.06cm\hangafter=1%
    \makebox[0.68cm][l]{\bfseries\color{caseink}#4}\songfont\color{caseink}#5%
  &#1%
  \\[0.06em]%
}
% \outlinepage：雙語對照大綱頁包裝環境（longtable 自動跨頁）
% 三欄：德文欄 + 中文欄 + 段碼欄（最右，不折行）
% 段碼欄寬度：1.80cm，足以容納最長段碼如「309–348」
% DE/ZH 欄寬：(0.87\textwidth - 2.08cm) / 2
%   驗算：2×欄 + 0.05\textwidth + 0.28cm gap + 1.80cm = 0.87tw - 2.08cm + 0.05tw + 2.08cm = 0.92tw ✓
\newenvironment{outlinepage}{%
  \vspace*{0.40cm}%
  \setlength{\LTleft}{0.04\textwidth}%
  \setlength{\LTright}{0.04\textwidth}%
  \setlength{\tabcolsep}{0pt}%
  % 大綱列距由 longtable 的 \arraystretch 與各 row macro 結尾的 \\[...em] 共同決定。
  % 這裡控制「每一列本身的表格 strut 高度」；數值越大，A./I./1. 各項之間越像有空行。
  % A3 原本用 0.62，視覺上沒有空行；A4 也沿用同值，避免切到 A4 後項目被撐開。
  % 若日後覺得大綱太擠，只微調這個值（例如 0.68），不要改各 \omainrow/\osubrow 的縮排。
  \renewcommand{\arraystretch}{0.62}%
  \begin{longtable}{%
    >{\vspace{0pt}\raggedright\arraybackslash}p{\dimexpr(0.87\textwidth-2.08cm)/2\relax}%
    @{\hspace{0.025\textwidth}}%
    !{\color{notegray}\vrule width 0.12pt}%
    @{\hspace{0.025\textwidth}}%
    >{\vspace{0pt}\raggedright\arraybackslash}p{\dimexpr(0.87\textwidth-2.08cm)/2\relax}%
    @{\hspace{0.28cm}}%
    >{\vspace{0pt}\footnotesize\normalfont\color{pendingink}\raggedright\arraybackslash}p{1.80cm}%
    @{}%
  }%
  % 首頁欄標籤：不要橫跨置中；分別放在德文欄與中文欄的實際欄位起點。
  \footnotesize\bfseries\color{sourceink}Gliederung%
  &\footnotesize\bfseries\songfont\color{sourceink}大綱%
  &%
  \\[0.36em]%
  \endfirsthead%
  % 續頁欄標籤：同樣分欄定位，以灰色細體區分；無需標注「續」。
  \footnotesize\color{pendingink}Gliederung%
  &\footnotesize\songfont\color{pendingink}大綱%
  &%
  \\[0.24em]%
  \endhead%
}{%
  \end{longtable}%
}
 之前設定；
%   預設 chinese / a4）
\ifdefined\docmode\else\def\docmode{chinese}\fi
\ifdefined\bilingualpaper\else\def\bilingualpaper{a4}\fi
\newif\ifshoworiginal
\newif\ifshowtranslation
\newif\ifbilinguallandscape
\newif\ifbilingualcolumns
\newif\ifintranslation
\newif\iffirstbilingualsection

\showoriginalfalse
\showtranslationfalse
\bilinguallandscapefalse
\bilingualcolumnsfalse
\intranslationfalse
\firstbilingualsectionfalse

\def\modechinese{chinese}
\def\modegerman{german}
\def\modebilingual{bilingual}
\def\bilingualpaperaThree{a3}
\def\bilingualpaperaFour{a4}

\ifx\docmode\modechinese
  \showtranslationtrue
\fi

\ifx\docmode\modegerman
  \showoriginaltrue
\fi

\ifx\docmode\modebilingual
  \showoriginaltrue
  \showtranslationtrue
  \bilinguallandscapetrue
  \bilingualcolumnstrue
\fi

% 編排模式說明：
% chinese  → \showoriginalfalse  \showtranslationtrue  （A4 直式，純中文）
% german   → \showoriginaltrue   \showtranslationfalse （A4 直式，純德文）
% bilingual→ \showoriginaltrue   \showtranslationtrue  \bilingualcolumnstrue（A3/A4 橫式對照）

\ifbilingualcolumns
  \ifx\bilingualpaper\bilingualpaperaFour
    \geometry{
      a4paper,
      landscape,
      left=1.25cm,
      right=2.25cm,
      top=1.25cm,
      bottom=1.75cm,
      footskip=8.5mm,
      marginparwidth=0.75cm,
      marginparsep=0.25cm
    }
  \else
    \geometry{
      a3paper,
      landscape,
      left=1.55cm,
      right=2.75cm,
      top=1.55cm,
      bottom=2.05cm,
      footskip=9.5mm,
      marginparwidth=0.9cm,
      marginparsep=0.35cm
    }
  \fi
\else
\geometry{
  a4paper,
  portrait,
  left=2.2cm,
  right=2.6cm,
  top=2.0cm,
  bottom=2.0cm,
  marginparwidth=0.8cm,
  marginparsep=0.3cm
}
\fi
% ===== 時間戳基礎設施 =====
\newcount\stamptimeval
\newcount\stampminuteval
\newcommand{\stamphh}{%
  \stamptimeval=\time
  \divide\stamptimeval by 60
  \ifnum\stamptimeval<10 0\fi
  \the\stamptimeval}
\newcommand{\stampmm}{%
  \stamptimeval=\time
  \stampminuteval=\time
  \divide\stampminuteval by 60
  \multiply\stampminuteval by 60
  \advance\stamptimeval by -\stampminuteval
  \ifnum\stamptimeval<10 0\fi
  \the\stamptimeval}
\newcommand{\stampwkde}[1]{%
  \ifcase#1Montag\or Dienstag\or Mittwoch\or Donnerstag\or Freitag\or Samstag\or Sonntag\fi}
\newcommand{\stampwkzh}[1]{%
  \ifcase#1 一\or 二\or 三\or 四\or 五\or 六\or 日\fi}
\newcommand{\stampmonthde}[1]{%
  \ifcase#1\or Januar\or Februar\or März\or April\or Mai\or Juni\or
  Juli\or August\or September\or Oktober\or November\or Dezember\fi}
\newcount\stampjdval
\newcount\stampwdval
\pgfcalendardatetojulian{\the\year-\the\month-\the\day}{\stampjdval}
\pgfcalendarjuliantoweekday{\the\stampjdval}{\stampwdval}
\newcommand{\timestampzh}{%
  \songfont 最後修改：\the\year 年\the\month 月\the\day 日%
  % （星期\stampwkzh{\the\stampwdval}）%
  % \space\stamphh:\stampmm
  }

\newcommand{\documentdatezh}{\the\year 年 \the\month 月 \the\day 日}
\newcommand{\documentdatede}{\the\day. \stampmonthde{\the\month} \the\year}
\newcommand{\bverfgetranslationnote}{%
  {\songfont 翻譯說明：初譯使用 Claude Code 與 GPT 協助迭代；仍請以德文原文為準。}%
}
\newcommand{\bverfgecourseinfo}{%
  {\songfont 課程：114 學年度第 2 學期；憲法專題研究（四）；授課老師：湯德宗教授；導讀日期：115 年 6 月 1 日。}%
}
\newcommand{\bverfgeeditorinfo}{%
  {\songfont 編者：王逸帆（R10A21126），國立台灣大學法律系研究所碩士班。}%
}
\newcommand{\bverfgetitleversionline}{%
  \ifbilingualcolumns
    Fassung \documentversion\enspace\textperiodcentered\enspace Stand: \documentdatede
    \quad
    {\songfont 版本 \documentversion\enspace\textperiodcentered\enspace 修訂日期：\documentdatezh\par}%
  \else
    \ifshowtranslation
      {\songfont 版本 \documentversion\enspace\textperiodcentered\enspace 修訂日期：\documentdatezh\par}%
    \else
      Fassung \documentversion\enspace\textperiodcentered\enspace Stand: \documentdatede\par
    \fi
  \fi
}
\newcommand{\bverfgetitlecornerinfo}{%
  \ifshowtranslation
    \vspace{0.72em}%
    \noindent\hfill
    \begin{minipage}{0.66\textwidth}
      \raggedleft\tiny\color{pendingink}\songfont
      \bverfgecourseinfo\par
      \bverfgeeditorinfo
    \end{minipage}%
  \fi
}
\newcommand{\bverfgetitlemeta}{%
  \vfill
  \begin{center}%
    \footnotesize\color{pendingink}%
    \bverfgetitleversionline
    \ifshowtranslation
      \vspace{0.28em}{\scriptsize\bverfgetranslationnote\par}%
    \fi
  \end{center}%
  \bverfgetitlecornerinfo
}

\pagestyle{fancy}
\fancyhf{}
\renewcommand{\headrulewidth}{0pt}
\renewcommand{\footrulewidth}{0pt}
\fancyfoot[C]{\thepage}
\ifbilingualcolumns
  \fancyfoot[R]{\scriptsize\color{pendingink}\timestampzh}%
\else
  \ifshoworiginal
  \else
    \fancyfoot[R]{\scriptsize\color{pendingink}\timestampzh\hspace{0.45cm}}%
  \fi
\fi
\setmainfont{Times New Roman}
\setsansfont{Helvetica Neue}
\setmonofont{Menlo}
\IfFileExists{/Users/iw/Library/Fonts/kaiu.ttf}{
  \setCJKmainfont[Path=/Users/iw/Library/Fonts/,AutoFakeBold=4]{kaiu.ttf}
}{
  \IfFontExistsTF{標楷體}{
    \setCJKmainfont[AutoFakeBold=4]{標楷體}
  }{
    \setCJKmainfont{Songti TC}
  }
}
\IfFontExistsTF{Songti TC}{
  \setCJKfamilyfont{song}{Songti TC}
  \setCJKmonofont{Songti TC}
}{
  \setCJKfamilyfont{song}[Path=/Users/iw/Library/Fonts/,AutoFakeBold=4]{kaiu.ttf}
  \setCJKmonofont[Path=/Users/iw/Library/Fonts/,AutoFakeBold=4]{kaiu.ttf}
}
\newcommand{\songfont}{\CJKfamily{song}}

% ===== 封面／目錄專用打字機字體（僅限封面、目錄頁使用，不影響正文）=====
\IfFontExistsTF{Radio Newsman}{%
  \newfontfamily\bverfgetitlede{Radio Newsman}%
}{%
  \newfontfamily\bverfgetitlede{Courier New}%
}
\IfFontExistsTF{Erikas Farbband}{%
  \newfontfamily\bverfgebodyde{Erikas Farbband}%
}{%
  \let\bverfgebodyde\bverfgetitlede
}
\IfFontExistsTF{Maoken Heavy Labourer}{%
  \setCJKfamilyfont{bverfgetitlecjk}{Maoken Heavy Labourer}%
}{%
  \setCJKfamilyfont{bverfgetitlecjk}{Songti TC}%
}
\IfFontExistsTF{Huiwen-mincho}{%
  \setCJKfamilyfont{bverfgebodycjk}{Huiwen-mincho}%
}{%
  \setCJKfamilyfont{bverfgebodycjk}{Songti TC}%
}
\newcommand{\bverfgetitlezh}{\bverfgetitlede\CJKfamily{bverfgetitlecjk}}
\newcommand{\bverfgebodyzh}{\bverfgebodyde\CJKfamily{bverfgebodycjk}}

\setstretch{1.35}
\setlist{itemsep=0.2em, topsep=0.4em}
\setlength{\parindent}{0pt}
\setlength{\parskip}{0.45em}
\raggedbottom
\setcounter{secnumdepth}{0}
\setcounter{tocdepth}{2}

\newcommand{\lawboxsourcefont}{\footnotesize}
\newcommand{\lawboxtransfont}{\footnotesize}
\newcommand{\lawboxsourcestretch}{1.02}
\newcommand{\lawboxtransstretch}{0.92}

\titleformat{\section}{\centering\Large\bfseries\songfont}{\thesection}{1em}{}
\titleformat{\subsection}{\large\bfseries\songfont}{\thesubsection}{1em}{}
\titleformat{\subsubsection}{\normalsize\bfseries\songfont}{\thesubsubsection}{1em}{}
\titleformat{\paragraph}{\normalsize\bfseries\songfont}{\theparagraph}{1em}{}
\titlespacing*{\section}{0pt}{1.8em}{1.0em}
\titlespacing*{\subsection}{0pt}{1.2em}{0.6em}
\titlespacing*{\subsubsection}{0pt}{1.0em}{0.45em}
\titlespacing*{\paragraph}{0pt}{0.6em}{0.4em}

\definecolor{noteblue}{HTML}{A9C9F5}
\definecolor{noteteal}{HTML}{AEE1D6}
\definecolor{notegray}{HTML}{D7DADF}
\definecolor{caseink}{HTML}{000000}
\definecolor{casepaper}{HTML}{FCFEFD}
\definecolor{sourcepaper}{HTML}{FAFBF8}
\definecolor{sourceline}{HTML}{D7DED8}
\definecolor{sourceink}{HTML}{000000}
\definecolor{pendingink}{HTML}{000000}

\tcbset{
  caseboxbase/.style={
    enhanced,
    breakable,
    width=0.985\linewidth,
    colback=casepaper,
    coltitle=caseink,
    fonttitle=\bfseries\songfont,
    boxrule=0.28pt,
    arc=1.2mm,
    left=2.4mm,
    right=2.4mm,
    top=1.5mm,
    bottom=1.5mm,
    boxsep=1.5mm,
    before skip=0.65em,
    after skip=0.65em,
    before upper={\setlength{\parindent}{0pt}\setlength{\parskip}{0.35em}},
  }
}

\newcommand{\bilingualstart}{}
\newcommand{\bilingualstop}{}
\ifbilingualcolumns
  \renewcommand{\bilingualstart}{\small\setlength{\columnsep}{1.25cm}\setlength{\columnseprule}{0.12pt}\columnratio{0.5,0.5}\multicolumnfootnotes\flushbottom\sloppy\interlinepenalty=80\clubpenalty=800\widowpenalty=800\displaywidowpenalty=800\begin{paracol}{2}\global\intranslationfalse\global\firstbilingualsectiontrue{\footnotesize\color{sourceink}Deutsch\par}\switchcolumn{\footnotesize\songfont\color{sourceink}中文譯文\par}\switchcolumn*}
  \renewcommand{\bilingualstop}{\ifintranslation\switchcolumn*\global\intranslationfalse\fi\end{paracol}\clearpage\raggedbottom}
\fi
\newif\ifrandnumbers
\randnumbersfalse
\newcounter{randnumber}
\newcounter{randnumberlocal}
\newcounter{lastsourcerandstart}
\newif\iflastsourcehasrandnumbers
\lastsourcehasrandnumbersfalse
\newif\ifinlawbox
\inlawboxfalse
\newif\iflawboxhaspendingrn
\lawboxhaspendingrnfalse
\newcommand{\lawboxpendingrn}{}
\newcommand{\startrandnumbers}{\global\randnumberstrue\setcounter{randnumber}{1}\global\lastsourcehasrandnumbersfalse}
\newcommand{\stoprandnumbers}{\global\randnumbersfalse\global\lastsourcehasrandnumbersfalse}
\newlength{\randnumberwidth}
\setlength{\randnumberwidth}{0.95cm}
\newlength{\randnumberrightoffset}
\setlength{\randnumberrightoffset}{0.85cm}
\newlength{\randnumberlawboxrightoffset}
\setlength{\randnumberlawboxrightoffset}{1.40cm}
\newcommand{\randnumberbox}[1]{%
  \leavevmode
  \makebox[0pt][l]{%
    \smash{%
      \tikz[remember picture,overlay,baseline=(rnmark.base)]{%
        \coordinate (rnmark) at (0,0);%
        \ifinlawbox
          \path let \p1=(current page.east), \p2=(rnmark.base) in
            node[
              anchor=base east,
              inner sep=0pt,
              outer sep=0pt,
              text width=\randnumberwidth,
              align=right,
              text=pendingink,
              font=\normalfont\mdseries\scriptsize\ttfamily
            ] at (\x1-\randnumberlawboxrightoffset,\y2) {{\normalfont\mdseries\scriptsize\ttfamily #1}};%
        \else
          \path let \p1=(current page.east), \p2=(rnmark.base) in
            node[
              anchor=base east,
              inner sep=0pt,
              outer sep=0pt,
              text width=\randnumberwidth,
              align=right,
              text=pendingink,
              font=\normalfont\mdseries\scriptsize\ttfamily
            ] at (\x1-\randnumberrightoffset,\y2) {{\normalfont\mdseries\scriptsize\ttfamily #1}};%
        \fi
      }%
    }%
  }%
  \ignorespaces
}
\newcommand{\lawboxstorerandnumber}[1]{%
  \gdef\lawboxpendingrn{#1}%
  \global\lawboxhaspendingrntrue
  \ignorespaces
}
\newcommand{\lawboxuserandnumber}{%
  \iflawboxhaspendingrn
    \randnumberbox{\lawboxpendingrn}%
    \global\lawboxhaspendingrnfalse
  \fi
}
\newcommand{\sourcern}[1]{%
  \ifrandnumbers
    \ifshoworiginal
      \ifshowtranslation\else
        \ifinlawbox\lawboxstorerandnumber{#1}\else\randnumberbox{#1}\fi
      \fi
    \fi
  \fi
}
\newcommand{\transrn}[1]{%
  \ifrandnumbers
    \ifshowtranslation
      \ifinlawbox\lawboxstorerandnumber{#1}\else\randnumberbox{#1}\fi
    \fi
  \fi
}
% ===== 課堂模式：德文原文縮寫註解 =====
% 日常切換只改各判決檔的一行：
%   \classroommodeon        開啟課堂版；註解掉這行即回到一般版。
% 模板預設：
%   1. 課堂模式關閉。
%   2. 課堂模式開啟時，縮寫註解包含「德文全稱＋中文譯名」。
%   3. 課堂模式開啟時，GG/條文編者註仍會自動顯示。
% 進階微調才需要在單案檔覆寫：
%   \classroomabbrzhoff     縮寫註解僅顯示德文全稱。
%   \showeditornotestrue    一般版也顯示編者註。
% 註解策略：同一實際 PDF 頁面中，同一縮寫只在第一次出現處加註。
% 這裡用 zref-abspage 的絕對頁序，不用 \thepage，避免文件內頁碼重置時誤判。
\newif\ifclassroommode
\newif\ifclassroomabbrzh
\classroommodefalse
\classroomabbrzhtrue
\newcommand{\classroommodeon}{\classroommodetrue}
\newcommand{\classroommodeoff}{\classroommodefalse}
\newcommand{\classroomabbrzhon}{\classroomabbrzhtrue}
\newcommand{\classroomabbrzhoff}{\classroomabbrzhfalse}
\newcommand{\abbrnotelabel}{\emph{Abk.}: }
\newcommand{\abbrnote}[3]{%
  \ifclassroommode
    \ifshoworiginal
      \ifintranslation\else
        \footnote{\normalfont\footnotesize\abbrnotelabel #2\ifclassroomabbrzh；{\songfont #3}\fi}%
      \fi
    \fi
  \fi
}
\newcommand{\abbrnoteonce}[4]{%
  #2%
  \ifclassroommode
    \ifshoworiginal
      \ifintranslation\else
        \begingroup
          \edef\abbrcurrentabspage{\number\numexpr\value{abspage}+1\relax}%
          \ifcsname abbrnoteseen@#1@\abbrcurrentabspage\endcsname\else
            \global\expandafter\let\csname abbrnoteseen@#1@\abbrcurrentabspage\endcsname\empty
            \abbrnote{#1}{#3}{#4}%
          \fi
        \endgroup
      \fi
    \fi
  \fi
}
\newcommand{\abbrAbs}{\abbrnoteonce{abs}{Abs.}{Absatz}{項}}
\newcommand{\abbrAAO}{\abbrnoteonce{aao}{a.a.O.}{am angegebenen Ort}{前引處}}
\newcommand{\abbrAE}{\abbrnoteonce{ae}{AE}{Alternativ-Entwurf}{替代草案}}
\newcommand{\abbrArt}{\abbrnoteonce{art}{Art.}{Artikel}{條}}
\newcommand{\abbrAufl}{\abbrnoteonce{aufl}{Aufl.}{Auflage}{版}}
\newcommand{\abbrBd}{\abbrnoteonce{bd}{Bd.}{Band}{卷}}
\newcommand{\abbrBGBl}{\abbrnoteonce{bgbl}{BGBl.}{Bundesgesetzblatt}{聯邦法律公報}}
\newcommand{\abbrBGHSt}{\abbrnoteonce{bghst}{BGHSt}{Entscheidungen des Bundesgerichtshofes in Strafsachen}{聯邦普通法院刑事判決彙編}}
\newcommand{\abbrBRDrucks}{\abbrnoteonce{brdrucks}{BRDrucks.}{Bundesrats-Drucksache}{聯邦參議院文件}}
\newcommand{\abbrBTDrucks}{\abbrnoteonce{btdrucks}{BTDrucks.}{Bundestags-Drucksache}{聯邦議院文件}}
\newcommand{\abbrBundesgesetzbl}{\abbrnoteonce{bundesgesetzbl}{Bundesgesetzbl.}{Bundesgesetzblatt}{聯邦法律公報}}
\newcommand{\abbrBvF}{\abbrnoteonce{bvf}{BvF}{Bundesverfassungsgerichtsverfahren}{聯邦憲法法院程序案號}}
\newcommand{\abbrBvL}{\abbrnoteonce{bvl}{BvL}{Bundesverfassungsgerichtsverfahren}{聯邦憲法法院程序案號}}
\newcommand{\abbrBvR}{\abbrnoteonce{bvr}{BvR}{Bundesverfassungsgerichtsverfahren}{聯邦憲法法院程序案號}}
\newcommand{\abbrBVerfG}{\abbrnoteonce{bverfg}{BVerfG}{Bundesverfassungsgericht}{聯邦憲法法院}}
\newcommand{\abbrBVerfGE}{\abbrnoteonce{bverfge}{BVerfGE}{Entscheidungen des Bundesverfassungsgerichts}{聯邦憲法法院判決集}}
\newcommand{\abbrBVerfGG}{\abbrnoteonce{bverfgg}{BVerfGG}{Bundesverfassungsgerichtsgesetz}{聯邦憲法法院法}}
\newcommand{\abbrDDR}{\abbrnoteonce{ddr}{DDR}{Deutsche Demokratische Republik}{德意志民主共和國；東德}}
\newcommand{\abbrDJT}{\abbrnoteonce{djt}{DJT}{Deutscher Juristentag}{德國法學家大會}}
\newcommand{\abbrDr}{\abbrnoteonce{dr}{Dr.}{Doktor}{Dr.}}
\newcommand{\abbrEuGRZ}{\abbrnoteonce{eugrz}{EuGRZ}{Europäische Grundrechte-Zeitschrift}{歐洲基本權期刊}}
\newcommand{\abbrEV}{\abbrnoteonce{ev}{EV}{Einigungsvertrag}{統一條約}}
\newcommand{\abbrF}{\abbrnoteonce{f}{f.}{folgende}{以下一頁；下一段}}
\newcommand{\abbrFF}{\abbrnoteonce{ff}{ff.}{fortfolgende}{以下數頁；以下各段}}
\newcommand{\abbrGBlDDR}{\abbrnoteonce{gblddr}{GBl. DDR}{Gesetzblatt der Deutschen Demokratischen Republik}{德意志民主共和國法律公報}}
\newcommand{\abbrGG}{\abbrnoteonce{gg}{GG}{Grundgesetz}{基本法}}
\newcommand{\abbrGVBl}{\abbrnoteonce{gvbl}{GVBl.}{Gesetz- und Verordnungsblatt}{法律及命令公報}}
\newcommand{\abbrJR}{\abbrnoteonce{jr}{JR}{Juristische Rundschau}{法學評論}}
\newcommand{\abbrJZ}{\abbrnoteonce{jz}{JZ}{JuristenZeitung}{法學家報}}
\newcommand{\abbrIVm}{\abbrnoteonce{ivm}{i.V.m.}{in Verbindung mit}{結合；連同}}
\newcommand{\abbrMdB}{\abbrnoteonce{mdb}{MdB}{Mitglied des Bundestages}{聯邦議院議員}}
\newcommand{\abbrMwN}{\abbrnoteonce{mwn}{m.w.N.}{mit weiteren Nachweisen}{附進一步文獻／判例出處}}
\newcommand{\abbrNF}{\abbrnoteonce{nf}{n.F.}{neue Fassung}{新版本}}
\newcommand{\abbrAF}{\abbrnoteonce{af}{a.F.}{alte Fassung}{舊版本}}
\newcommand{\abbrNr}{\abbrnoteonce{nr}{Nr.}{Nummer}{款、目或號，依語境}}
\newcommand{\abbrProf}{\abbrnoteonce{prof}{Prof.}{Professor}{教授}}
\newcommand{\abbrRdnr}{\abbrnoteonce{rdnr}{Rdnr.}{Randnummer}{邊碼；段碼}}
\newcommand{\abbrRdnrn}{\abbrnoteonce{rdnrn}{Rdnrn.}{Randnummern}{邊碼；段碼}}
\newcommand{\abbrRGBl}{\abbrnoteonce{rgbl}{RGBl.}{Reichsgesetzblatt}{帝國法律公報}}
\newcommand{\abbrRGSt}{\abbrnoteonce{rgst}{RGSt}{Entscheidungen des Reichsgerichts in Strafsachen}{帝國法院刑事判決彙編}}
\newcommand{\abbrRspr}{\abbrnoteonce{rspr}{Rspr.}{Rechtsprechung}{判例；司法實務}}
\newcommand{\abbrRVO}{\abbrnoteonce{rvo}{RVO}{Reichsversicherungsordnung}{帝國保險法}}
\newcommand{\abbrS}{\abbrnoteonce{s}{S.}{Seite}{頁}}
\newcommand{\abbrSFHG}{\abbrnoteonce{sfhg}{SFHG}{Schwangeren- und Familienhilfegesetz}{孕婦及家庭扶助法}}
\newcommand{\abbrSGBV}{\abbrnoteonce{sgbv}{SGB V}{Fünftes Buch Sozialgesetzbuch}{社會法典第五編}}
\newcommand{\abbrStAG}{\abbrnoteonce{staeg}{StÄG}{Strafrechtsänderungsgesetz}{刑法修正法}}
\newcommand{\abbrStenBer}{\abbrnoteonce{stenber}{StenBer.}{Stenographischer Bericht}{速記紀錄}}
\newcommand{\abbrStGB}{\abbrnoteonce{stgb}{StGB}{Strafgesetzbuch}{刑法}}
\newcommand{\abbrStREG}{\abbrnoteonce{streg}{StREG}{Strafrechtsreform-Ergänzungsgesetz}{刑法改革補充法}}
\newcommand{\abbrStrRG}{\abbrnoteonce{strrg}{StrRG}{Strafrechtsreformgesetz}{刑法改革法}}
\newcommand{\abbrUS}{\abbrnoteonce{us}{U.S.}{United States Reports}{美國最高法院判決彙編}}
\newcommand{\abbrVgl}{\abbrnoteonce{vgl}{vgl.}{vergleiche}{參見}}
\newcommand{\abbrVVDStRL}{\abbrnoteonce{vvdstrl}{VVDStRL}{Veröffentlichungen der Vereinigung der Deutschen Staatsrechtslehrer}{德國公法教師協會論文集}}
\newcommand{\abbrWP}{\abbrnoteonce{wp}{WP.}{Wahlperiode}{屆期}}
\newcommand{\abbrWp}{\abbrnoteonce{wp}{Wp.}{Wahlperiode}{屆期}}
\newcommand{\abbrZStrW}{\abbrnoteonce{zstrw}{ZStrW}{Zeitschrift für die gesamte Strafrechtswissenschaft}{整體刑法學期刊}}
\newcommand{\recordsourcerandstart}{%
  \ifrandnumbers
    \setcounter{lastsourcerandstart}{\value{randnumber}}%
    \global\lastsourcehasrandnumberstrue
  \else
    \global\lastsourcehasrandnumbersfalse
  \fi
}
\newlength{\biparagraphskip}
\setlength{\biparagraphskip}{0.52em}
\newlength{\lawboxblockbeforeskip}
\setlength{\lawboxblockbeforeskip}{0.72em}
\newlength{\lawboxblockafterskip}
\setlength{\lawboxblockafterskip}{0.92em}
\newif\iflawboxpartendedblock
\lawboxpartendedblockfalse
\newcommand{\sourceparstyle}{\footnotesize\color{sourceink}\setstretch{1.12}\RaggedRight\emergencystretch=3em\setlength{\parskip}{0.42em}\obeylines\selectfont}
\newcommand{\transparstyle}{\songfont\color{caseink}\setstretch{1.12}\setlength{\parindent}{0pt}\setlength{\parskip}{0.42em}}
\newcommand{\bilingualpairskip}{%
  \par\addvspace{\biparagraphskip}%
  \switchcolumn
  \par\addvspace{\biparagraphskip}%
  \switchcolumn*%
  \global\intranslationfalse
}
\newcommand{\bilingualboxpairskip}{%
  \par
  \switchcolumn
  \par
  \switchcolumn*%
  \global\intranslationfalse
}
\newcommand{\bilinguallawboxblockskip}{%
  \switchcolumn*%
  \par\addvspace{\lawboxblockafterskip}%
  \switchcolumn
  \par\addvspace{\lawboxblockafterskip}%
  \switchcolumn*%
  \global\intranslationfalse
}
\newcommand{\bikeepwithnext}[1][4]{%
  \ifbilingualcolumns
    \syncleft
    \Needspace{#1\baselineskip}%
    \switchcolumn
    \Needspace{#1\baselineskip}%
    \switchcolumn*%
    \global\intranslationfalse
  \else
    \Needspace{#1\baselineskip}%
  \fi
}
\NewEnviron{sourcepara}[1][]{
  \recordsourcerandstart
  \ifshoworiginal
    \ifbilingualcolumns
      \ifintranslation\switchcolumn*\global\intranslationfalse\fi
    \else
      \Needspace{5\baselineskip}
    \fi
    \ifbilingualcolumns\else\par\addvspace{0.35em}\fi
    {\sourceparstyle
    \noindent\BODY\par}
    \ifbilingualcolumns\else\addvspace{0.25em}\fi
    \ifbilingualcolumns
      \switchcolumn
      \global\intranslationtrue
    \fi
  \else
    \ifrandnumbers
      \begingroup
        \setbox0=\vbox{\hsize=\linewidth\sourceparstyle\noindent\BODY\par}%
      \endgroup
    \fi
  \fi
}
\NewEnviron{transpara}{
  \ifshowtranslation
    \setcounter{randnumberlocal}{\value{lastsourcerandstart}}%
    \ifbilingualcolumns
      \ifintranslation\else\switchcolumn\global\intranslationtrue\fi
      {\transparstyle\setlength{\leftskip}{0.55em}\setlength{\rightskip}{0.15em plus 1em}\addtolength{\linewidth}{-0.55em}\BODY\par}
      \switchcolumn*
      \bilingualpairskip
    \else
      {\transparstyle\BODY\par}
    \fi
  \else
    \ifbilingualcolumns
      \ifintranslation\switchcolumn*\global\intranslationfalse\fi
    \fi
  \fi
}
\NewEnviron{sourceboxpara}[1][]{
  \global\lawboxpartendedblockfalse
  \recordsourcerandstart
  \ifshoworiginal
    \ifbilingualcolumns
      \ifintranslation\switchcolumn*\global\intranslationfalse\fi
    \else
      \Needspace{3\baselineskip}
    \fi
    {\sourceparstyle
    \BODY\par}
    \ifbilingualcolumns\else
      \iflawboxpartendedblock\addvspace{\lawboxblockafterskip}\fi
    \fi
    \ifbilingualcolumns
      \switchcolumn
      \global\intranslationtrue
    \fi
  \fi
}
\NewEnviron{transboxpara}{
  \global\lawboxpartendedblockfalse
  \ifshowtranslation
    \setcounter{randnumberlocal}{\value{lastsourcerandstart}}%
    \ifbilingualcolumns
      \ifintranslation\else\switchcolumn\global\intranslationtrue\fi
      {\transparstyle\BODY\par}
      \iflawboxpartendedblock
        \bilinguallawboxblockskip
      \else
        \switchcolumn*
        \bilingualboxpairskip
      \fi
    \else
      {\transparstyle\BODY\par}
      \iflawboxpartendedblock\addvspace{\lawboxblockafterskip}\fi
    \fi
  \else
    \ifbilingualcolumns
      \ifintranslation\switchcolumn*\global\intranslationfalse\fi
    \fi
  \fi
}
\newcommand{\leitsatznum}[1]{\noindent\hangindent=3.0em\hangafter=1\makebox[2.25em][r]{\bfseries #1.}\hspace{0.55em}\ignorespaces}
\newcommand{\syncleft}{\ifbilingualcolumns\ifintranslation\switchcolumn*\global\intranslationfalse\fi\fi}
\pretocmd{\section}{\syncleft\clearpage}{}{}
\pretocmd{\subsection}{\syncleft}{}{}
\pretocmd{\subsubsection}{\syncleft}{}{}
\newcommand{\biheadingfont}{\songfont}
\newcommand{\bitocentry}[2]{%
  \ifshoworiginal
    \ifshowtranslation
      \ifstrequal{#1}{#2}{#1}{#1\space #2}%
    \else
      #1%
    \fi
  \else
    #2%
  \fi
}
\pdfstringdefDisableCommands{%
  \def\bitocentry#1#2{#1}%
}
\newcommand{\biheading}[7]{%
  \syncleft#1\bikeepwithnext[#3]\phantomsection\addcontentsline{toc}{#2}{\bitocentry{#6}{#7}}%
  \ifbilingualcolumns
    {#4 #6\par}%
    \switchcolumn
    {#4\biheadingfont #7\par}%
    \switchcolumn*%
    \global\intranslationfalse
    \par\addvspace{#5}%
    \switchcolumn
    \par\addvspace{#5}%
    \switchcolumn*%
  \else
    \ifshoworiginal{#4 #6\par}\fi
    \ifshowtranslation{#4\biheadingfont #7\par}\fi
    \addvspace{#5}%
  \fi
}
\newcommand{\termsectionheading}[1]{%
  \par\Needspace{9\baselineskip}%
  \vspace{0.65em}%
  {\songfont\bfseries #1\par}%
  \vspace{0.2em}%
}
% 章節換頁：
%   雙語 paracol 欄內不要用 \clearpage；它會在同步兩欄時吐出只有欄線/頁碼的空頁。
%   \newpage 足以讓下一個 section 起新頁，又不會強制清空所有待排材料。
%   單語模式不在 paracol 內，仍可用 \clearpage。
\newcommand{\bisectionpagebreak}{\ifbilingualcolumns\newpage\else\clearpage\fi}
\newcommand{\bisection}[2]{%
  \biheading{\iffirstbilingualsection\global\firstbilingualsectionfalse\else\bisectionpagebreak\fi}{section}{6}{\centering\Large\bfseries}{0.8em}{#1}{#2}%
}
\newcommand{\bisubsection}[2]{%
  \biheading{}{subsection}{5}{\large\bfseries}{0.55em}{#1}{#2}%
}
\newcommand{\bisubsubsection}[2]{%
  \biheading{}{subsubsection}{4}{\normalsize\bfseries}{0.45em}{#1}{#2}%
}
\newcommand{\bicompositeheading}[4]{%
  \biheading{}{subsection}{5}{\large\bfseries}{0.16em}{#1}{#2}%
  \biheading{}{subsubsection}{3}{\normalsize\bfseries}{0.45em}{#3}{#4}%
}
\newif\iflawboxfirstitem
\lawboxfirstitemfalse
\newcommand{\lawboxprespace}[1]{%
  \par
  \iflawboxfirstitem
    \global\lawboxfirstitemfalse
  \else
    \addvspace{#1}%
  \fi
}
\newtcolorbox{lawboxinner}{
  caseboxbase,
  colframe=sourceline!85!black,
  colback=white,
  left=1.05em,
  right=1.05em,
  top=0.62em,
  bottom=0.62em,
  before upper={\global\inlawboxtrue\global\lawboxfirstitemtrue\global\lawboxhaspendingrnfalse\ifintranslation\lawboxtransfont\else\lawboxsourcefont\fi\RaggedRight\setlength{\parindent}{0pt}\setlength{\parskip}{0pt}\ifintranslation\setstretch{\lawboxtransstretch}\else\setstretch{\lawboxsourcestretch}\fi},
  after upper={\global\lawboxhaspendingrnfalse\global\inlawboxfalse}
}
\tcbset{
  lawboxpartbase/.style={
    enhanced,
    breakable=false,
    width=0.985\linewidth,
    colback=white,
    frame hidden,
    boxrule=0pt,
    arc=0pt,
    left=1.05em,
    right=1.05em,
    boxsep=1.5mm,
    before skip=0pt,
    after skip=0pt,
    before upper={\global\inlawboxtrue\global\lawboxfirstitemtrue\global\lawboxhaspendingrnfalse\ifintranslation\lawboxtransfont\else\lawboxsourcefont\fi\RaggedRight\setlength{\parindent}{0pt}\setlength{\parskip}{0pt}\ifintranslation\setstretch{\lawboxtransstretch}\else\setstretch{\lawboxsourcestretch}\fi},
    after upper={\global\lawboxhaspendingrnfalse\global\inlawboxfalse},
    borderline west={0.28pt}{0pt}{sourceline!85!black},
    borderline east={0.28pt}{0pt}{sourceline!85!black},
  },
  lawboxpartfirst/.style={
    lawboxpartbase,
    top=0.62em,
    bottom=0.04em,
    before skip=\lawboxblockbeforeskip,
    borderline north={0.28pt}{0pt}{sourceline!85!black},
  },
  lawboxpartmiddle/.style={
    lawboxpartbase,
    top=0.04em,
    bottom=0.04em,
  },
  lawboxpartlast/.style={
    lawboxpartbase,
    top=0.04em,
    bottom=0.62em,
    after skip=0pt,
    borderline south={0.28pt}{0pt}{sourceline!85!black},
  }
}
\newtcolorbox{lawboxpartinner}[1]{#1}
\newtcolorbox{termboxinner}{
  caseboxbase,
  colframe=noteblue!70!black,
  colback=casepaper,
  before upper={\songfont\setlength{\parindent}{0pt}\setlength{\parskip}{0.35em}}
}
\newtcolorbox{deboxinner}{
  caseboxbase,
  colframe=notegray!72!black,
  colback=white,
  left=1.05em,
  right=1.05em,
  top=0.62em,
  bottom=0.62em,
  before upper={\global\inlawboxtrue\global\lawboxfirstitemtrue\global\lawboxhaspendingrnfalse\small\setlength{\parindent}{0pt}\setlength{\parskip}{0.35em}},
  after upper={\global\lawboxhaspendingrnfalse\global\inlawboxfalse}
}
\NewEnviron{lawbox}{
  \begin{lawboxinner}\BODY\end{lawboxinner}
}
\NewEnviron{lawboxpart}[2]{
  \begin{lawboxpartinner}{lawboxpart#1,equal height group=#2}\BODY\end{lawboxpartinner}%
  \ifstrequal{#1}{last}{\global\lawboxpartendedblocktrue}{}%
}
\NewEnviron{termbox}{
  \par\addvspace{0.55em}
  {\color{noteblue!70!black}\rule{\linewidth}{0.28pt}}\par
  \begingroup
    \songfont\small\color{caseink}
    \setlength{\parindent}{0pt}
    \setlength{\parskip}{0.24em}
    \BODY
    \par
  \endgroup
  {\color{noteblue!70!black}\rule{\linewidth}{0.28pt}}\par
  \addvspace{0.55em}
}
\NewEnviron{debox}{
  \begin{deboxinner}\BODY\end{deboxinner}
}
\providecommand{\paragraphmark}[1]{}
\newcommand{\lawboxtitleafter}{\addvspace{0.06em}}
\newcommand{\lawtitle}[1]{\lawboxprespace{0.22em}{\lawboxtransfont\bfseries\lawboxuserandnumber #1}\par\lawboxtitleafter}
\newcommand{\absno}[2]{\lawboxprespace{0.04em}\noindent\lawboxuserandnumber\hangindent=2.6em\hangafter=1\makebox[2.6em][l]{(#1)}#2\par}
\newcommand{\nrno}[2]{\lawboxprespace{0pt}\noindent\lawboxuserandnumber\hspace*{2.6em}\hangindent=4.4em\hangafter=1\makebox[1.8em][l]{#1.}#2\par}
\newcommand{\letterno}[2]{\lawboxprespace{0pt}\noindent\lawboxuserandnumber\hspace*{4.4em}\hangindent=6.0em\hangafter=1\makebox[1.6em][l]{#1)}#2\par}
\newcommand{\sourcelawtitle}[1]{\lawboxprespace{0.22em}{\lawboxsourcefont\bfseries\lawboxuserandnumber #1}\par\lawboxtitleafter}
\newcommand{\sourcelawsubtitle}[1]{\lawboxprespace{0.04em}{\lawboxsourcefont\bfseries\lawboxuserandnumber #1}\par\lawboxtitleafter}
\newcommand{\sourceabsno}[2]{\lawboxprespace{0.04em}\noindent\lawboxuserandnumber\hangindent=2.45em\hangafter=1\makebox[2.45em][l]{(#1)}#2\par}
\newcommand{\sourcenrno}[2]{\lawboxprespace{0pt}\noindent\lawboxuserandnumber\hspace*{2.45em}\hangindent=4.25em\hangafter=1\makebox[1.8em][l]{#1.}#2\par}
\newcommand{\sourceletterno}[2]{\lawboxprespace{0pt}\noindent\lawboxuserandnumber\hspace*{4.25em}\hangindent=5.85em\hangafter=1\makebox[1.6em][l]{#1)}#2\par}
\newcommand{\sourcequotepara}[1]{\lawboxprespace{0.04em}\noindent\lawboxuserandnumber\hangindent=1.2em\hangafter=1 #1\par}
\newcommand{\sourcelawline}[1]{\lawboxprespace{0.04em}\noindent\lawboxuserandnumber #1\par}
\newcommand{\lawline}[1]{\lawboxprespace{0.04em}\noindent\lawboxuserandnumber #1\par}
\newcommand{\pendingtranslation}{\par{\songfont\footnotesize\color{pendingink}（中文譯文待補。）}\par}
\newcommand{\tocpart}[1]{\par\addvspace{0.52em}{\footnotesize\bfseries\color{pendingink}#1\par}\addvspace{0.16em}}
\newcommand{\tocdashline}{\par\addvspace{0.48em}\noindent{\color{notegray}\xleaders\hbox to 0.82em{\hfil\rule{0.42em}{0.28pt}\hfil}\hskip\linewidth}\par\addvspace{0.38em}}
% \tocpanel{font-cmd}{content}：將目錄置中於所在欄寬（雙語欄適用）。
% A3 橫式使用 0.58\linewidth，A4 橫式使用 0.84\linewidth，
% 使兩種紙張下目錄欄內容實際寬度相近（均約 100–105 mm）。
\newcommand{\tocpanel}[2]{%
  \makebox[\linewidth][c]{%
    \ifx\bilingualpaper\bilingualpaperaThree
      \begin{minipage}{0.58\linewidth}%
    \else
      \begin{minipage}{0.84\linewidth}%
    \fi
    \toccolstyle
    #1#2%
    \end{minipage}%
  }%
}
% ===== 首頁簡目錄的兩層 description list 縮排 =====
% enumitem 的 description list 可把每一列拆成：
%   [label]  labelsep  body text
% 這裡的「起點」都以所在 minipage/欄位的左緣為 0。
%
% 核心規則：
%   labelindent = label 欄位從左緣開始的位置。
%   labelwidth  = label 欄位本身寬度；標籤太長（如 Abw. Meinung）就加大它。
%   labelsep    = label 欄位右緣到正文的距離。
%   leftmargin  = 正文 body text 的起點。判斷層級視覺時主要看這個。
%
% 所以：
%   想讓「Begründung / 判決理由」整列正文往右或往左：改 tocitems 的 leftmargin。
%   想讓「A. Verfahren und Gegenstand / A. 程序與審查標的」正文往右或往左：
%     改 tocsubitems 的 leftmargin。
%   想讓 A./B./C. 這些標籤本身往右或往左，但正文不動：改 tocsubitems 的 labelindent。
%   想讓 A./B./C. 與正文間距變大或變小：改 tocsubitems 的 labelsep。
%
% 層級原則：
%   tocsubitems.leftmargin 必須大於 tocitems.leftmargin，
%   否則子層級正文會看起來比「Gründe / 理由」還靠左。
\newlist{tocitems}{description}{1}
\setlist[tocitems]{%
  labelindent=0pt,    % 第一層 label 欄位起點；0pt 表示貼齊目錄欄左緣。
  leftmargin=2.54cm,  % 第一層正文起點；例如 Begründung / 判決理由 的左緣。
  labelwidth=2.16cm,  % 第一層 label 欄位寬度；需容納 Leitsätze、Abw. Meinung、不同意見等。
  labelsep=0.32cm,    % 第一層 label 與正文的水平間距；只改間距，不改正文起點。
  itemsep=0.28em,     % 第一層各項之間的垂直距離。
  topsep=0.20em,      % 第一層 list 上下與前後內容的垂直距離。
  parsep=0pt,         % 同一 item 內段落之間的距離；目錄項通常不需要。
  partopsep=0pt,      % list 若緊接段落開始時的額外距離；關閉以免目錄高度飄動。
  align=left,         % label 欄內左對齊；長短標籤不靠右貼齊。
  font=\normalfont\bfseries % label 的字型；正文不受這行影響。
}
\newlist{tocsubitems}{description}{1}
\setlist[tocsubitems]{%
  labelindent=1cm,    % 第二層 label 起點；只控制 A./B./C. 本身的位置。
  leftmargin=3.00cm,  % 第二層正文起點；必須大於 2.54cm，才會比 Begründung / 判決理由更右。
  labelwidth=0.5cm,   % 第二層 label 欄寬；A./B./C. 很短，可小於第一層。
  labelsep=1.5cm,    % 第二層 label 與正文間距；A. 與 Verfahren / 程序 的距離看這裡。
  itemsep=0.12em,     % 第二層各項較密，避免 A.–E. 佔太高。
  topsep=0.04em,      % 第二層 list 與 Gründe / 理由之間的垂直距離。
  parsep=0pt,         % 同一第二層 item 內段落距離；目錄項通常不需要。
  partopsep=0pt,      % 關閉額外段落起始距離，避免版面不穩。
  align=left,         % A./B./C. label 欄內左對齊。
  font=\normalfont\bfseries, % A./B./C. label 加粗；正文不受這行影響。
  before={}           % 不在第二層 list 前自動插入額外內容。
}
\newlist{termitems}{description}{1}
\ifbilingualcolumns
  \setlist[termitems]{%
    leftmargin=5.2cm,
    labelwidth=4.8cm,
    labelsep=0.35cm,
    itemsep=0.16em,
    topsep=0.2em,
    parsep=0pt,
    partopsep=0pt,
    font=\normalfont\bfseries,
    style=nextline
  }
\else
  \setlist[termitems]{%
    leftmargin=0pt,
    labelwidth=0pt,
    labelsep=0pt,
    itemsep=0.24em,
    topsep=0.2em,
    parsep=0pt,
    partopsep=0pt,
    font=\normalfont\bfseries,
    style=nextline
  }
\fi
\newlist{abbritems}{description}{1}
\setlist[abbritems]{%
  leftmargin=2.38cm,
  labelwidth=2.02cm,
  labelsep=0.28cm,
  itemsep=0.16em,
  topsep=0.2em,
  parsep=0pt,
  partopsep=0pt,
  align=left,
  font=\normalfont\bfseries
}
\ifbilingualcolumns
  \newenvironment{termcolumns}{\begin{multicols}{2}}{\end{multicols}}
\else
  \newenvironment{termcolumns}{}{}
\fi

% ===== 編者註腳開關 =====
% 一般版預設不顯示編者註，避免正文腳註過密。
% 若開啟 \classroommodeon，GG/條文編者註仍會自動顯示，無需另開本開關。
% 若一般版也要顯示編者註，可在單案檔 \input 後加 \showeditornotestrue。
\newif\ifshoweditornotes
\showeditornotesfalse

% ===== 目錄排版輔助 =====
\newcommand{\toccolstyle}{\raggedright\arraybackslash}
\newcommand{\tocsingleindent}{\leftskip=1.45cm\rightskip=1.2cm}

% ===== 行內引文環境（德文原文與中文譯文各一，帶左側灰色邊線）=====
\newcommand{\quoteparbreak}{\par\addvspace{0.48em}}
\newcommand{\sourcequoteinner}[1]{%
  \begin{tcolorbox}[
    enhanced, breakable, width=\dimexpr\linewidth-0.8em\relax, frame hidden,
    colback=white, coltext=caseink, boxrule=0pt,
    borderline west={0.55pt}{0.88em}{notegray},
    left=1.78em, right=0.72em, top=0.18em, bottom=0.18em,
    boxsep=0pt, before skip=0.24em, after skip=0.24em,
    before upper={\normalfont\footnotesize\color{caseink}\setlength{\parindent}{0pt}\setlength{\parskip}{0.34em}\raggedright\setlength{\rightskip}{0pt plus 1em}}
  ]%
  #1\par
  \end{tcolorbox}%
}
\newcommand{\transquoteinner}[1]{%
  \begin{tcolorbox}[
    enhanced, breakable, width=\dimexpr\linewidth-0.8em\relax, frame hidden,
    colback=white, coltext=caseink, boxrule=0pt,
    borderline west={0.55pt}{0.88em}{notegray},
    left=1.78em, right=0.72em, top=0.18em, bottom=0.18em,
    boxsep=0pt, before skip=0.24em, after skip=0.24em,
    before upper={\songfont\small\color{caseink}\setlength{\parindent}{0pt}\setlength{\parskip}{0.34em}\raggedright\setlength{\rightskip}{0pt plus 1em}}
  ]%
  #1\par
  \end{tcolorbox}%
}

% ===== 大綱（Gliederung）Rn 輔助命令 =====
\newcommand{\outrn}[2]{\enspace{\normalfont\footnotesize\color{pendingink}(Rn\,#1–#2)}}
\newcommand{\outrnsingle}[1]{\enspace{\normalfont\footnotesize\color{pendingink}(Rn\,#1)}}
\newcommand{\outrnzh}[2]{\enspace{\normalfont\footnotesize\color{pendingink}（第\,#1–#2\,段）}}
\newcommand{\outrnzhs}[1]{\enspace{\normalfont\footnotesize\color{pendingink}（第\,#1\,段）}}
% ===== 大綱雙語對照表格輔助命令（三欄：段碼 │ 德文 │ 中文）=====
% 欄 1：段碼（由欄定義自動格式化：小字、等寬、右對齊、pendingink）
% 欄 2：德文標籤＋正文
% 欄 3：中文標籤＋正文
%
% 無段碼的列（\omainrow、\osecrow），欄 1 以 & 空置。
% 有段碼的列（\osecrowrn、\osubrow、\ointrow），#1 為預格式化字串
%   例：\osubrow{2–49}{I.}{...}{I.}{...}
%       \ointrow{107}{...}{...}
%
% \opartrow：段組標題（橫跨三欄）
\newcommand{\opartrow}[2]{%
  \noalign{\vspace{0.30em}}%
  \multicolumn{3}{@{}l@{}}{%
    \footnotesize\bfseries\color{pendingink}#1\hfill\songfont#2%
  }\\[0.06em]%
}
% \odashrow：虛線分隔（橫跨三欄）
\newcommand{\odashrow}{%
  \noalign{\vspace{0.28em}}%
  \multicolumn{3}{@{}l@{}}{\color{notegray}%
    \xleaders\hbox to 0.82em{\hfil\rule{0.42em}{0.28pt}\hfil}\hskip 0.88\linewidth}%
  \\[0.24em]%
}
% \odissentpart：不同意見書區段標頭。比一般 \opartrow 更像附錄性轉場。
\newcommand{\odissentpart}[2]{%
  \noalign{\vspace{0.42em}}%
  \multicolumn{3}{@{}p{0.88\textwidth}@{}}{%
    \color{notegray}\rule{\linewidth}{0.18pt}%
  }\\[-0.18em]%
  \multicolumn{3}{@{}p{0.88\textwidth}@{}}{%
    \makebox[\linewidth][s]{%
      \footnotesize\bfseries\color{pendingink}#1%
      \hfill{\songfont #2}%
    }%
  }\\[-0.02em]%
}
% \odissentopinion：不同意見書的署名與一行摘要，右欄保留段碼範圍。
\newcommand{\odissentopinion}[5]{%
  \footnotesize\hangindent=0.52cm\hangafter=1%
    {\bfseries\color{caseink}#2}\enspace{\color{notegray}\textperiodcentered}\enspace\color{caseink}#3%
  &\footnotesize\hangindent=0.52cm\hangafter=1%
    {\songfont\bfseries\color{caseink}#4}\enspace{\color{notegray}\textperiodcentered}\enspace\songfont\color{caseink}#5%
  &#1%
  \\[0.12em]%
}
% \omainrow：頂層項目，無段碼（Leitsätze、Urteil、Gründe …）
\newcommand{\omainrow}[4]{%
  \footnotesize\hangindent=1.92cm\hangafter=1%
    \makebox[1.92cm][l]{\bfseries\color{caseink}#1}\color{caseink}#2%
  &\footnotesize\hangindent=1.92cm\hangafter=1%
    \makebox[1.92cm][l]{\bfseries\color{caseink}#3}\songfont\color{caseink}#4%
  &%
  \\[0.15em]%
}
% \omainrowrn：頂層項目，有段碼（不同意見等標籤寬於 \osecrow 框者）
\newcommand{\omainrowrn}[5]{%
  \footnotesize\hangindent=1.92cm\hangafter=1%
    \makebox[1.92cm][l]{\bfseries\color{caseink}#2}\color{caseink}#3%
  &\footnotesize\hangindent=1.92cm\hangafter=1%
    \makebox[1.92cm][l]{\bfseries\color{caseink}#4}\songfont\color{caseink}#5%
  &#1%
  \\[0.15em]%
}
% \osecrow：一級子項，無段碼（A.、C.、D. 等有子項者）
\newcommand{\osecrow}[4]{%
  \footnotesize\hspace{1.10cm}\hangindent=1.98cm\hangafter=1%
    \makebox[0.86cm][l]{\bfseries\color{caseink}#1}\color{caseink}#2%
  &\footnotesize\hspace{1.10cm}\hangindent=1.98cm\hangafter=1%
    \makebox[0.86cm][l]{\bfseries\color{caseink}#3}\songfont\color{caseink}#4%
  &%
  \\[0.11em]%
}
% \osecrowrn：一級子項，有段碼（B.、E. 等完整段落範圍者）
% #1=段碼字串（如 101–106）, #2=DE標籤, #3=DE正文, #4=ZH標籤, #5=ZH正文
\newcommand{\osecrowrn}[5]{%
  \footnotesize\hspace{1.10cm}\hangindent=1.98cm\hangafter=1%
    \makebox[0.86cm][l]{\bfseries\color{caseink}#2}\color{caseink}#3%
  &\footnotesize\hspace{1.10cm}\hangindent=1.98cm\hangafter=1%
    \makebox[0.86cm][l]{\bfseries\color{caseink}#4}\songfont\color{caseink}#5%
  &#1%
  \\[0.11em]%
}
% \osubrow：二級子項，有段碼範圍（I.、II.、A.I. …）
% #1=段碼字串, #2=DE標籤, #3=DE正文, #4=ZH標籤, #5=ZH正文
\newcommand{\osubrow}[5]{%
  \footnotesize\hspace{1.40cm}\hangindent=2.30cm\hangafter=1%
    \makebox[0.90cm][l]{\bfseries\color{caseink}#2}\color{caseink}#3%
  &\footnotesize\hspace{1.40cm}\hangindent=2.30cm\hangafter=1%
    \makebox[0.90cm][l]{\bfseries\color{caseink}#4}\songfont\color{caseink}#5%
  &#1%
  \\[0.09em]%
}
% \ointrow：引入段落，有單一段碼，無標籤
% #1=段碼字串, #2=DE正文, #3=ZH正文
\newcommand{\ointrow}[3]{%
  \footnotesize\hspace{0.52cm}\hangindent=0.52cm\hangafter=1\color{caseink}#2%
  &\footnotesize\hspace{0.52cm}\hangindent=0.52cm\hangafter=1\songfont\color{caseink}#3%
  &#1%
  \\[0.09em]%
}
% \osubsubrow：三級子項（1.、2.、3.）
% #1=段碼字串, #2=DE標籤, #3=DE正文, #4=ZH標籤, #5=ZH正文
\newcommand{\osubsubrow}[5]{%
  \footnotesize\hspace{1.80cm}\hangindent=2.48cm\hangafter=1%
    \makebox[0.68cm][l]{\bfseries\color{caseink}#2}\color{caseink}#3%
  &\footnotesize\hspace{1.80cm}\hangindent=2.48cm\hangafter=1%
    \makebox[0.68cm][l]{\bfseries\color{caseink}#4}\songfont\color{caseink}#5%
  &#1%
  \\[0.08em]%
}
% \osubsubsubrow：四級子項（a）、b））
% #1=段碼字串, #2=DE標籤, #3=DE正文, #4=ZH標籤, #5=ZH正文
\newcommand{\osubsubsubrow}[5]{%
  \footnotesize\hspace{2.10cm}\hangindent=2.68cm\hangafter=1%
    \makebox[0.58cm][l]{\bfseries\color{caseink}#2}\color{caseink}#3%
  &\footnotesize\hspace{2.10cm}\hangindent=2.68cm\hangafter=1%
    \makebox[0.58cm][l]{\bfseries\color{caseink}#4}\songfont\color{caseink}#5%
  &#1%
  \\[0.07em]%
}
% \osubsubsubsubrow：五級子項（aa）、bb））
% 標籤框 0.62cm：足以容納「aa)」在 footnotesize bold 下（約 0.50cm）並留有間距
% #1=段碼字串, #2=DE標籤, #3=DE正文, #4=ZH標籤, #5=ZH正文
\newcommand{\osubsubsubsubrow}[5]{%
  \footnotesize\hspace{2.38cm}\hangindent=3.06cm\hangafter=1%
    \makebox[0.68cm][l]{\bfseries\color{caseink}#2}\color{caseink}#3%
  &\footnotesize\hspace{2.38cm}\hangindent=3.06cm\hangafter=1%
    \makebox[0.68cm][l]{\bfseries\color{caseink}#4}\songfont\color{caseink}#5%
  &#1%
  \\[0.06em]%
}
% \outlinepage：雙語對照大綱頁包裝環境（longtable 自動跨頁）
% 三欄：德文欄 + 中文欄 + 段碼欄（最右，不折行）
% 段碼欄寬度：1.80cm，足以容納最長段碼如「309–348」
% DE/ZH 欄寬：(0.87\textwidth - 2.08cm) / 2
%   驗算：2×欄 + 0.05\textwidth + 0.28cm gap + 1.80cm = 0.87tw - 2.08cm + 0.05tw + 2.08cm = 0.92tw ✓
\newenvironment{outlinepage}{%
  \vspace*{0.40cm}%
  \setlength{\LTleft}{0.04\textwidth}%
  \setlength{\LTright}{0.04\textwidth}%
  \setlength{\tabcolsep}{0pt}%
  % 大綱列距由 longtable 的 \arraystretch 與各 row macro 結尾的 \\[...em] 共同決定。
  % 這裡控制「每一列本身的表格 strut 高度」；數值越大，A./I./1. 各項之間越像有空行。
  % A3 原本用 0.62，視覺上沒有空行；A4 也沿用同值，避免切到 A4 後項目被撐開。
  % 若日後覺得大綱太擠，只微調這個值（例如 0.68），不要改各 \omainrow/\osubrow 的縮排。
  \renewcommand{\arraystretch}{0.62}%
  \begin{longtable}{%
    >{\vspace{0pt}\raggedright\arraybackslash}p{\dimexpr(0.87\textwidth-2.08cm)/2\relax}%
    @{\hspace{0.025\textwidth}}%
    !{\color{notegray}\vrule width 0.12pt}%
    @{\hspace{0.025\textwidth}}%
    >{\vspace{0pt}\raggedright\arraybackslash}p{\dimexpr(0.87\textwidth-2.08cm)/2\relax}%
    @{\hspace{0.28cm}}%
    >{\vspace{0pt}\footnotesize\normalfont\color{pendingink}\raggedright\arraybackslash}p{1.80cm}%
    @{}%
  }%
  % 首頁欄標籤：不要橫跨置中；分別放在德文欄與中文欄的實際欄位起點。
  \footnotesize\bfseries\color{sourceink}Gliederung%
  &\footnotesize\bfseries\songfont\color{sourceink}大綱%
  &%
  \\[0.36em]%
  \endfirsthead%
  % 續頁欄標籤：同樣分欄定位，以灰色細體區分；無需標注「續」。
  \footnotesize\color{pendingink}Gliederung%
  &\footnotesize\songfont\color{pendingink}大綱%
  &%
  \\[0.24em]%
  \endhead%
}{%
  \end{longtable}%
}
 之前設定；
%   預設 chinese / a4）
\ifdefined\docmode\else\def\docmode{chinese}\fi
\ifdefined\bilingualpaper\else\def\bilingualpaper{a4}\fi
\newif\ifshoworiginal
\newif\ifshowtranslation
\newif\ifbilinguallandscape
\newif\ifbilingualcolumns
\newif\ifintranslation
\newif\iffirstbilingualsection

\showoriginalfalse
\showtranslationfalse
\bilinguallandscapefalse
\bilingualcolumnsfalse
\intranslationfalse
\firstbilingualsectionfalse

\def\modechinese{chinese}
\def\modegerman{german}
\def\modebilingual{bilingual}
\def\bilingualpaperaThree{a3}
\def\bilingualpaperaFour{a4}

\ifx\docmode\modechinese
  \showtranslationtrue
\fi

\ifx\docmode\modegerman
  \showoriginaltrue
\fi

\ifx\docmode\modebilingual
  \showoriginaltrue
  \showtranslationtrue
  \bilinguallandscapetrue
  \bilingualcolumnstrue
\fi

% 編排模式說明：
% chinese  → \showoriginalfalse  \showtranslationtrue  （A4 直式，純中文）
% german   → \showoriginaltrue   \showtranslationfalse （A4 直式，純德文）
% bilingual→ \showoriginaltrue   \showtranslationtrue  \bilingualcolumnstrue（A3/A4 橫式對照）

\ifbilingualcolumns
  \ifx\bilingualpaper\bilingualpaperaFour
    \geometry{
      a4paper,
      landscape,
      left=1.25cm,
      right=2.25cm,
      top=1.25cm,
      bottom=1.75cm,
      footskip=8.5mm,
      marginparwidth=0.75cm,
      marginparsep=0.25cm
    }
  \else
    \geometry{
      a3paper,
      landscape,
      left=1.55cm,
      right=2.75cm,
      top=1.55cm,
      bottom=2.05cm,
      footskip=9.5mm,
      marginparwidth=0.9cm,
      marginparsep=0.35cm
    }
  \fi
\else
\geometry{
  a4paper,
  portrait,
  left=2.2cm,
  right=2.6cm,
  top=2.0cm,
  bottom=2.0cm,
  marginparwidth=0.8cm,
  marginparsep=0.3cm
}
\fi
% ===== 時間戳基礎設施 =====
\newcount\stamptimeval
\newcount\stampminuteval
\newcommand{\stamphh}{%
  \stamptimeval=\time
  \divide\stamptimeval by 60
  \ifnum\stamptimeval<10 0\fi
  \the\stamptimeval}
\newcommand{\stampmm}{%
  \stamptimeval=\time
  \stampminuteval=\time
  \divide\stampminuteval by 60
  \multiply\stampminuteval by 60
  \advance\stamptimeval by -\stampminuteval
  \ifnum\stamptimeval<10 0\fi
  \the\stamptimeval}
\newcommand{\stampwkde}[1]{%
  \ifcase#1Montag\or Dienstag\or Mittwoch\or Donnerstag\or Freitag\or Samstag\or Sonntag\fi}
\newcommand{\stampwkzh}[1]{%
  \ifcase#1 一\or 二\or 三\or 四\or 五\or 六\or 日\fi}
\newcommand{\stampmonthde}[1]{%
  \ifcase#1\or Januar\or Februar\or März\or April\or Mai\or Juni\or
  Juli\or August\or September\or Oktober\or November\or Dezember\fi}
\newcount\stampjdval
\newcount\stampwdval
\pgfcalendardatetojulian{\the\year-\the\month-\the\day}{\stampjdval}
\pgfcalendarjuliantoweekday{\the\stampjdval}{\stampwdval}
\newcommand{\timestampzh}{%
  \songfont 最後修改：\the\year 年\the\month 月\the\day 日%
  % （星期\stampwkzh{\the\stampwdval}）%
  % \space\stamphh:\stampmm
  }

\newcommand{\documentdatezh}{\the\year 年 \the\month 月 \the\day 日}
\newcommand{\documentdatede}{\the\day. \stampmonthde{\the\month} \the\year}
\newcommand{\bverfgetranslationnote}{%
  {\songfont 翻譯說明：初譯使用 Claude Code 與 GPT 協助迭代；仍請以德文原文為準。}%
}
\newcommand{\bverfgecourseinfo}{%
  {\songfont 課程：114 學年度第 2 學期；憲法專題研究（四）；授課老師：湯德宗教授；導讀日期：115 年 6 月 1 日。}%
}
\newcommand{\bverfgeeditorinfo}{%
  {\songfont 編者：王逸帆（R10A21126），國立台灣大學法律系研究所碩士班。}%
}
\newcommand{\bverfgetitleversionline}{%
  \ifbilingualcolumns
    Fassung \documentversion\enspace\textperiodcentered\enspace Stand: \documentdatede
    \quad
    {\songfont 版本 \documentversion\enspace\textperiodcentered\enspace 修訂日期：\documentdatezh\par}%
  \else
    \ifshowtranslation
      {\songfont 版本 \documentversion\enspace\textperiodcentered\enspace 修訂日期：\documentdatezh\par}%
    \else
      Fassung \documentversion\enspace\textperiodcentered\enspace Stand: \documentdatede\par
    \fi
  \fi
}
\newcommand{\bverfgetitlecornerinfo}{%
  \ifshowtranslation
    \vspace{0.72em}%
    \noindent\hfill
    \begin{minipage}{0.66\textwidth}
      \raggedleft\tiny\color{pendingink}\songfont
      \bverfgecourseinfo\par
      \bverfgeeditorinfo
    \end{minipage}%
  \fi
}
\newcommand{\bverfgetitlemeta}{%
  \vfill
  \begin{center}%
    \footnotesize\color{pendingink}%
    \bverfgetitleversionline
    \ifshowtranslation
      \vspace{0.28em}{\scriptsize\bverfgetranslationnote\par}%
    \fi
  \end{center}%
  \bverfgetitlecornerinfo
}

\pagestyle{fancy}
\fancyhf{}
\renewcommand{\headrulewidth}{0pt}
\renewcommand{\footrulewidth}{0pt}
\fancyfoot[C]{\thepage}
\ifbilingualcolumns
  \fancyfoot[R]{\scriptsize\color{pendingink}\timestampzh}%
\else
  \ifshoworiginal
  \else
    \fancyfoot[R]{\scriptsize\color{pendingink}\timestampzh\hspace{0.45cm}}%
  \fi
\fi
\setmainfont{Times New Roman}
\setsansfont{Helvetica Neue}
\setmonofont{Menlo}
\IfFileExists{/Users/iw/Library/Fonts/kaiu.ttf}{
  \setCJKmainfont[Path=/Users/iw/Library/Fonts/,AutoFakeBold=4]{kaiu.ttf}
}{
  \IfFontExistsTF{標楷體}{
    \setCJKmainfont[AutoFakeBold=4]{標楷體}
  }{
    \setCJKmainfont{Songti TC}
  }
}
\IfFontExistsTF{Songti TC}{
  \setCJKfamilyfont{song}{Songti TC}
  \setCJKmonofont{Songti TC}
}{
  \setCJKfamilyfont{song}[Path=/Users/iw/Library/Fonts/,AutoFakeBold=4]{kaiu.ttf}
  \setCJKmonofont[Path=/Users/iw/Library/Fonts/,AutoFakeBold=4]{kaiu.ttf}
}
\newcommand{\songfont}{\CJKfamily{song}}

% ===== 封面／目錄專用打字機字體（僅限封面、目錄頁使用，不影響正文）=====
\IfFontExistsTF{Radio Newsman}{%
  \newfontfamily\bverfgetitlede{Radio Newsman}%
}{%
  \newfontfamily\bverfgetitlede{Courier New}%
}
\IfFontExistsTF{Erikas Farbband}{%
  \newfontfamily\bverfgebodyde{Erikas Farbband}%
}{%
  \let\bverfgebodyde\bverfgetitlede
}
\IfFontExistsTF{Maoken Heavy Labourer}{%
  \setCJKfamilyfont{bverfgetitlecjk}{Maoken Heavy Labourer}%
}{%
  \setCJKfamilyfont{bverfgetitlecjk}{Songti TC}%
}
\IfFontExistsTF{Huiwen-mincho}{%
  \setCJKfamilyfont{bverfgebodycjk}{Huiwen-mincho}%
}{%
  \setCJKfamilyfont{bverfgebodycjk}{Songti TC}%
}
\newcommand{\bverfgetitlezh}{\bverfgetitlede\CJKfamily{bverfgetitlecjk}}
\newcommand{\bverfgebodyzh}{\bverfgebodyde\CJKfamily{bverfgebodycjk}}

\setstretch{1.35}
\setlist{itemsep=0.2em, topsep=0.4em}
\setlength{\parindent}{0pt}
\setlength{\parskip}{0.45em}
\raggedbottom
\setcounter{secnumdepth}{0}
\setcounter{tocdepth}{2}

\newcommand{\lawboxsourcefont}{\footnotesize}
\newcommand{\lawboxtransfont}{\footnotesize}
\newcommand{\lawboxsourcestretch}{1.02}
\newcommand{\lawboxtransstretch}{0.92}

\titleformat{\section}{\centering\Large\bfseries\songfont}{\thesection}{1em}{}
\titleformat{\subsection}{\large\bfseries\songfont}{\thesubsection}{1em}{}
\titleformat{\subsubsection}{\normalsize\bfseries\songfont}{\thesubsubsection}{1em}{}
\titleformat{\paragraph}{\normalsize\bfseries\songfont}{\theparagraph}{1em}{}
\titlespacing*{\section}{0pt}{1.8em}{1.0em}
\titlespacing*{\subsection}{0pt}{1.2em}{0.6em}
\titlespacing*{\subsubsection}{0pt}{1.0em}{0.45em}
\titlespacing*{\paragraph}{0pt}{0.6em}{0.4em}

\definecolor{noteblue}{HTML}{A9C9F5}
\definecolor{noteteal}{HTML}{AEE1D6}
\definecolor{notegray}{HTML}{D7DADF}
\definecolor{caseink}{HTML}{000000}
\definecolor{casepaper}{HTML}{FCFEFD}
\definecolor{sourcepaper}{HTML}{FAFBF8}
\definecolor{sourceline}{HTML}{D7DED8}
\definecolor{sourceink}{HTML}{000000}
\definecolor{pendingink}{HTML}{000000}

\tcbset{
  caseboxbase/.style={
    enhanced,
    breakable,
    width=0.985\linewidth,
    colback=casepaper,
    coltitle=caseink,
    fonttitle=\bfseries\songfont,
    boxrule=0.28pt,
    arc=1.2mm,
    left=2.4mm,
    right=2.4mm,
    top=1.5mm,
    bottom=1.5mm,
    boxsep=1.5mm,
    before skip=0.65em,
    after skip=0.65em,
    before upper={\setlength{\parindent}{0pt}\setlength{\parskip}{0.35em}},
  }
}

\newcommand{\bilingualstart}{}
\newcommand{\bilingualstop}{}
\ifbilingualcolumns
  \renewcommand{\bilingualstart}{\small\setlength{\columnsep}{1.25cm}\setlength{\columnseprule}{0.12pt}\columnratio{0.5,0.5}\multicolumnfootnotes\flushbottom\sloppy\interlinepenalty=80\clubpenalty=800\widowpenalty=800\displaywidowpenalty=800\begin{paracol}{2}\global\intranslationfalse\global\firstbilingualsectiontrue{\footnotesize\color{sourceink}Deutsch\par}\switchcolumn{\footnotesize\songfont\color{sourceink}中文譯文\par}\switchcolumn*}
  \renewcommand{\bilingualstop}{\ifintranslation\switchcolumn*\global\intranslationfalse\fi\end{paracol}\clearpage\raggedbottom}
\fi
\newif\ifrandnumbers
\randnumbersfalse
\newcounter{randnumber}
\newcounter{randnumberlocal}
\newcounter{lastsourcerandstart}
\newif\iflastsourcehasrandnumbers
\lastsourcehasrandnumbersfalse
\newif\ifinlawbox
\inlawboxfalse
\newif\iflawboxhaspendingrn
\lawboxhaspendingrnfalse
\newcommand{\lawboxpendingrn}{}
\newcommand{\startrandnumbers}{\global\randnumberstrue\setcounter{randnumber}{1}\global\lastsourcehasrandnumbersfalse}
\newcommand{\stoprandnumbers}{\global\randnumbersfalse\global\lastsourcehasrandnumbersfalse}
\newlength{\randnumberwidth}
\setlength{\randnumberwidth}{0.95cm}
\newlength{\randnumberrightoffset}
\setlength{\randnumberrightoffset}{0.85cm}
\newlength{\randnumberlawboxrightoffset}
\setlength{\randnumberlawboxrightoffset}{1.40cm}
\newcommand{\randnumberbox}[1]{%
  \leavevmode
  \makebox[0pt][l]{%
    \smash{%
      \tikz[remember picture,overlay,baseline=(rnmark.base)]{%
        \coordinate (rnmark) at (0,0);%
        \ifinlawbox
          \path let \p1=(current page.east), \p2=(rnmark.base) in
            node[
              anchor=base east,
              inner sep=0pt,
              outer sep=0pt,
              text width=\randnumberwidth,
              align=right,
              text=pendingink,
              font=\normalfont\mdseries\scriptsize\ttfamily
            ] at (\x1-\randnumberlawboxrightoffset,\y2) {{\normalfont\mdseries\scriptsize\ttfamily #1}};%
        \else
          \path let \p1=(current page.east), \p2=(rnmark.base) in
            node[
              anchor=base east,
              inner sep=0pt,
              outer sep=0pt,
              text width=\randnumberwidth,
              align=right,
              text=pendingink,
              font=\normalfont\mdseries\scriptsize\ttfamily
            ] at (\x1-\randnumberrightoffset,\y2) {{\normalfont\mdseries\scriptsize\ttfamily #1}};%
        \fi
      }%
    }%
  }%
  \ignorespaces
}
\newcommand{\lawboxstorerandnumber}[1]{%
  \gdef\lawboxpendingrn{#1}%
  \global\lawboxhaspendingrntrue
  \ignorespaces
}
\newcommand{\lawboxuserandnumber}{%
  \iflawboxhaspendingrn
    \randnumberbox{\lawboxpendingrn}%
    \global\lawboxhaspendingrnfalse
  \fi
}
\newcommand{\sourcern}[1]{%
  \ifrandnumbers
    \ifshoworiginal
      \ifshowtranslation\else
        \ifinlawbox\lawboxstorerandnumber{#1}\else\randnumberbox{#1}\fi
      \fi
    \fi
  \fi
}
\newcommand{\transrn}[1]{%
  \ifrandnumbers
    \ifshowtranslation
      \ifinlawbox\lawboxstorerandnumber{#1}\else\randnumberbox{#1}\fi
    \fi
  \fi
}
% ===== 課堂模式：德文原文縮寫註解 =====
% 日常切換只改各判決檔的一行：
%   \classroommodeon        開啟課堂版；註解掉這行即回到一般版。
% 模板預設：
%   1. 課堂模式關閉。
%   2. 課堂模式開啟時，縮寫註解包含「德文全稱＋中文譯名」。
%   3. 課堂模式開啟時，GG/條文編者註仍會自動顯示。
% 進階微調才需要在單案檔覆寫：
%   \classroomabbrzhoff     縮寫註解僅顯示德文全稱。
%   \showeditornotestrue    一般版也顯示編者註。
% 註解策略：同一實際 PDF 頁面中，同一縮寫只在第一次出現處加註。
% 這裡用 zref-abspage 的絕對頁序，不用 \thepage，避免文件內頁碼重置時誤判。
\newif\ifclassroommode
\newif\ifclassroomabbrzh
\classroommodefalse
\classroomabbrzhtrue
\newcommand{\classroommodeon}{\classroommodetrue}
\newcommand{\classroommodeoff}{\classroommodefalse}
\newcommand{\classroomabbrzhon}{\classroomabbrzhtrue}
\newcommand{\classroomabbrzhoff}{\classroomabbrzhfalse}
\newcommand{\abbrnotelabel}{\emph{Abk.}: }
\newcommand{\abbrnote}[3]{%
  \ifclassroommode
    \ifshoworiginal
      \ifintranslation\else
        \footnote{\normalfont\footnotesize\abbrnotelabel #2\ifclassroomabbrzh；{\songfont #3}\fi}%
      \fi
    \fi
  \fi
}
\newcommand{\abbrnoteonce}[4]{%
  #2%
  \ifclassroommode
    \ifshoworiginal
      \ifintranslation\else
        \begingroup
          \edef\abbrcurrentabspage{\number\numexpr\value{abspage}+1\relax}%
          \ifcsname abbrnoteseen@#1@\abbrcurrentabspage\endcsname\else
            \global\expandafter\let\csname abbrnoteseen@#1@\abbrcurrentabspage\endcsname\empty
            \abbrnote{#1}{#3}{#4}%
          \fi
        \endgroup
      \fi
    \fi
  \fi
}
\newcommand{\abbrAbs}{\abbrnoteonce{abs}{Abs.}{Absatz}{項}}
\newcommand{\abbrAAO}{\abbrnoteonce{aao}{a.a.O.}{am angegebenen Ort}{前引處}}
\newcommand{\abbrAE}{\abbrnoteonce{ae}{AE}{Alternativ-Entwurf}{替代草案}}
\newcommand{\abbrArt}{\abbrnoteonce{art}{Art.}{Artikel}{條}}
\newcommand{\abbrAufl}{\abbrnoteonce{aufl}{Aufl.}{Auflage}{版}}
\newcommand{\abbrBd}{\abbrnoteonce{bd}{Bd.}{Band}{卷}}
\newcommand{\abbrBGBl}{\abbrnoteonce{bgbl}{BGBl.}{Bundesgesetzblatt}{聯邦法律公報}}
\newcommand{\abbrBGHSt}{\abbrnoteonce{bghst}{BGHSt}{Entscheidungen des Bundesgerichtshofes in Strafsachen}{聯邦普通法院刑事判決彙編}}
\newcommand{\abbrBRDrucks}{\abbrnoteonce{brdrucks}{BRDrucks.}{Bundesrats-Drucksache}{聯邦參議院文件}}
\newcommand{\abbrBTDrucks}{\abbrnoteonce{btdrucks}{BTDrucks.}{Bundestags-Drucksache}{聯邦議院文件}}
\newcommand{\abbrBundesgesetzbl}{\abbrnoteonce{bundesgesetzbl}{Bundesgesetzbl.}{Bundesgesetzblatt}{聯邦法律公報}}
\newcommand{\abbrBvF}{\abbrnoteonce{bvf}{BvF}{Bundesverfassungsgerichtsverfahren}{聯邦憲法法院程序案號}}
\newcommand{\abbrBvL}{\abbrnoteonce{bvl}{BvL}{Bundesverfassungsgerichtsverfahren}{聯邦憲法法院程序案號}}
\newcommand{\abbrBvR}{\abbrnoteonce{bvr}{BvR}{Bundesverfassungsgerichtsverfahren}{聯邦憲法法院程序案號}}
\newcommand{\abbrBVerfG}{\abbrnoteonce{bverfg}{BVerfG}{Bundesverfassungsgericht}{聯邦憲法法院}}
\newcommand{\abbrBVerfGE}{\abbrnoteonce{bverfge}{BVerfGE}{Entscheidungen des Bundesverfassungsgerichts}{聯邦憲法法院判決集}}
\newcommand{\abbrBVerfGG}{\abbrnoteonce{bverfgg}{BVerfGG}{Bundesverfassungsgerichtsgesetz}{聯邦憲法法院法}}
\newcommand{\abbrDDR}{\abbrnoteonce{ddr}{DDR}{Deutsche Demokratische Republik}{德意志民主共和國；東德}}
\newcommand{\abbrDJT}{\abbrnoteonce{djt}{DJT}{Deutscher Juristentag}{德國法學家大會}}
\newcommand{\abbrDr}{\abbrnoteonce{dr}{Dr.}{Doktor}{Dr.}}
\newcommand{\abbrEuGRZ}{\abbrnoteonce{eugrz}{EuGRZ}{Europäische Grundrechte-Zeitschrift}{歐洲基本權期刊}}
\newcommand{\abbrEV}{\abbrnoteonce{ev}{EV}{Einigungsvertrag}{統一條約}}
\newcommand{\abbrF}{\abbrnoteonce{f}{f.}{folgende}{以下一頁；下一段}}
\newcommand{\abbrFF}{\abbrnoteonce{ff}{ff.}{fortfolgende}{以下數頁；以下各段}}
\newcommand{\abbrGBlDDR}{\abbrnoteonce{gblddr}{GBl. DDR}{Gesetzblatt der Deutschen Demokratischen Republik}{德意志民主共和國法律公報}}
\newcommand{\abbrGG}{\abbrnoteonce{gg}{GG}{Grundgesetz}{基本法}}
\newcommand{\abbrGVBl}{\abbrnoteonce{gvbl}{GVBl.}{Gesetz- und Verordnungsblatt}{法律及命令公報}}
\newcommand{\abbrJR}{\abbrnoteonce{jr}{JR}{Juristische Rundschau}{法學評論}}
\newcommand{\abbrJZ}{\abbrnoteonce{jz}{JZ}{JuristenZeitung}{法學家報}}
\newcommand{\abbrIVm}{\abbrnoteonce{ivm}{i.V.m.}{in Verbindung mit}{結合；連同}}
\newcommand{\abbrMdB}{\abbrnoteonce{mdb}{MdB}{Mitglied des Bundestages}{聯邦議院議員}}
\newcommand{\abbrMwN}{\abbrnoteonce{mwn}{m.w.N.}{mit weiteren Nachweisen}{附進一步文獻／判例出處}}
\newcommand{\abbrNF}{\abbrnoteonce{nf}{n.F.}{neue Fassung}{新版本}}
\newcommand{\abbrAF}{\abbrnoteonce{af}{a.F.}{alte Fassung}{舊版本}}
\newcommand{\abbrNr}{\abbrnoteonce{nr}{Nr.}{Nummer}{款、目或號，依語境}}
\newcommand{\abbrProf}{\abbrnoteonce{prof}{Prof.}{Professor}{教授}}
\newcommand{\abbrRdnr}{\abbrnoteonce{rdnr}{Rdnr.}{Randnummer}{邊碼；段碼}}
\newcommand{\abbrRdnrn}{\abbrnoteonce{rdnrn}{Rdnrn.}{Randnummern}{邊碼；段碼}}
\newcommand{\abbrRGBl}{\abbrnoteonce{rgbl}{RGBl.}{Reichsgesetzblatt}{帝國法律公報}}
\newcommand{\abbrRGSt}{\abbrnoteonce{rgst}{RGSt}{Entscheidungen des Reichsgerichts in Strafsachen}{帝國法院刑事判決彙編}}
\newcommand{\abbrRspr}{\abbrnoteonce{rspr}{Rspr.}{Rechtsprechung}{判例；司法實務}}
\newcommand{\abbrRVO}{\abbrnoteonce{rvo}{RVO}{Reichsversicherungsordnung}{帝國保險法}}
\newcommand{\abbrS}{\abbrnoteonce{s}{S.}{Seite}{頁}}
\newcommand{\abbrSFHG}{\abbrnoteonce{sfhg}{SFHG}{Schwangeren- und Familienhilfegesetz}{孕婦及家庭扶助法}}
\newcommand{\abbrSGBV}{\abbrnoteonce{sgbv}{SGB V}{Fünftes Buch Sozialgesetzbuch}{社會法典第五編}}
\newcommand{\abbrStAG}{\abbrnoteonce{staeg}{StÄG}{Strafrechtsänderungsgesetz}{刑法修正法}}
\newcommand{\abbrStenBer}{\abbrnoteonce{stenber}{StenBer.}{Stenographischer Bericht}{速記紀錄}}
\newcommand{\abbrStGB}{\abbrnoteonce{stgb}{StGB}{Strafgesetzbuch}{刑法}}
\newcommand{\abbrStREG}{\abbrnoteonce{streg}{StREG}{Strafrechtsreform-Ergänzungsgesetz}{刑法改革補充法}}
\newcommand{\abbrStrRG}{\abbrnoteonce{strrg}{StrRG}{Strafrechtsreformgesetz}{刑法改革法}}
\newcommand{\abbrUS}{\abbrnoteonce{us}{U.S.}{United States Reports}{美國最高法院判決彙編}}
\newcommand{\abbrVgl}{\abbrnoteonce{vgl}{vgl.}{vergleiche}{參見}}
\newcommand{\abbrVVDStRL}{\abbrnoteonce{vvdstrl}{VVDStRL}{Veröffentlichungen der Vereinigung der Deutschen Staatsrechtslehrer}{德國公法教師協會論文集}}
\newcommand{\abbrWP}{\abbrnoteonce{wp}{WP.}{Wahlperiode}{屆期}}
\newcommand{\abbrWp}{\abbrnoteonce{wp}{Wp.}{Wahlperiode}{屆期}}
\newcommand{\abbrZStrW}{\abbrnoteonce{zstrw}{ZStrW}{Zeitschrift für die gesamte Strafrechtswissenschaft}{整體刑法學期刊}}
\newcommand{\recordsourcerandstart}{%
  \ifrandnumbers
    \setcounter{lastsourcerandstart}{\value{randnumber}}%
    \global\lastsourcehasrandnumberstrue
  \else
    \global\lastsourcehasrandnumbersfalse
  \fi
}
\newlength{\biparagraphskip}
\setlength{\biparagraphskip}{0.52em}
\newlength{\lawboxblockbeforeskip}
\setlength{\lawboxblockbeforeskip}{0.72em}
\newlength{\lawboxblockafterskip}
\setlength{\lawboxblockafterskip}{0.92em}
\newif\iflawboxpartendedblock
\lawboxpartendedblockfalse
\newcommand{\sourceparstyle}{\footnotesize\color{sourceink}\setstretch{1.12}\RaggedRight\emergencystretch=3em\setlength{\parskip}{0.42em}\obeylines\selectfont}
\newcommand{\transparstyle}{\songfont\color{caseink}\setstretch{1.12}\setlength{\parindent}{0pt}\setlength{\parskip}{0.42em}}
\newcommand{\bilingualpairskip}{%
  \par\addvspace{\biparagraphskip}%
  \switchcolumn
  \par\addvspace{\biparagraphskip}%
  \switchcolumn*%
  \global\intranslationfalse
}
\newcommand{\bilingualboxpairskip}{%
  \par
  \switchcolumn
  \par
  \switchcolumn*%
  \global\intranslationfalse
}
\newcommand{\bilinguallawboxblockskip}{%
  \switchcolumn*%
  \par\addvspace{\lawboxblockafterskip}%
  \switchcolumn
  \par\addvspace{\lawboxblockafterskip}%
  \switchcolumn*%
  \global\intranslationfalse
}
\newcommand{\bikeepwithnext}[1][4]{%
  \ifbilingualcolumns
    \syncleft
    \Needspace{#1\baselineskip}%
    \switchcolumn
    \Needspace{#1\baselineskip}%
    \switchcolumn*%
    \global\intranslationfalse
  \else
    \Needspace{#1\baselineskip}%
  \fi
}
\NewEnviron{sourcepara}[1][]{
  \recordsourcerandstart
  \ifshoworiginal
    \ifbilingualcolumns
      \ifintranslation\switchcolumn*\global\intranslationfalse\fi
    \else
      \Needspace{5\baselineskip}
    \fi
    \ifbilingualcolumns\else\par\addvspace{0.35em}\fi
    {\sourceparstyle
    \noindent\BODY\par}
    \ifbilingualcolumns\else\addvspace{0.25em}\fi
    \ifbilingualcolumns
      \switchcolumn
      \global\intranslationtrue
    \fi
  \else
    \ifrandnumbers
      \begingroup
        \setbox0=\vbox{\hsize=\linewidth\sourceparstyle\noindent\BODY\par}%
      \endgroup
    \fi
  \fi
}
\NewEnviron{transpara}{
  \ifshowtranslation
    \setcounter{randnumberlocal}{\value{lastsourcerandstart}}%
    \ifbilingualcolumns
      \ifintranslation\else\switchcolumn\global\intranslationtrue\fi
      {\transparstyle\setlength{\leftskip}{0.55em}\setlength{\rightskip}{0.15em plus 1em}\addtolength{\linewidth}{-0.55em}\BODY\par}
      \switchcolumn*
      \bilingualpairskip
    \else
      {\transparstyle\BODY\par}
    \fi
  \else
    \ifbilingualcolumns
      \ifintranslation\switchcolumn*\global\intranslationfalse\fi
    \fi
  \fi
}
\NewEnviron{sourceboxpara}[1][]{
  \global\lawboxpartendedblockfalse
  \recordsourcerandstart
  \ifshoworiginal
    \ifbilingualcolumns
      \ifintranslation\switchcolumn*\global\intranslationfalse\fi
    \else
      \Needspace{3\baselineskip}
    \fi
    {\sourceparstyle
    \BODY\par}
    \ifbilingualcolumns\else
      \iflawboxpartendedblock\addvspace{\lawboxblockafterskip}\fi
    \fi
    \ifbilingualcolumns
      \switchcolumn
      \global\intranslationtrue
    \fi
  \fi
}
\NewEnviron{transboxpara}{
  \global\lawboxpartendedblockfalse
  \ifshowtranslation
    \setcounter{randnumberlocal}{\value{lastsourcerandstart}}%
    \ifbilingualcolumns
      \ifintranslation\else\switchcolumn\global\intranslationtrue\fi
      {\transparstyle\BODY\par}
      \iflawboxpartendedblock
        \bilinguallawboxblockskip
      \else
        \switchcolumn*
        \bilingualboxpairskip
      \fi
    \else
      {\transparstyle\BODY\par}
      \iflawboxpartendedblock\addvspace{\lawboxblockafterskip}\fi
    \fi
  \else
    \ifbilingualcolumns
      \ifintranslation\switchcolumn*\global\intranslationfalse\fi
    \fi
  \fi
}
\newcommand{\leitsatznum}[1]{\noindent\hangindent=3.0em\hangafter=1\makebox[2.25em][r]{\bfseries #1.}\hspace{0.55em}\ignorespaces}
\newcommand{\syncleft}{\ifbilingualcolumns\ifintranslation\switchcolumn*\global\intranslationfalse\fi\fi}
\pretocmd{\section}{\syncleft\clearpage}{}{}
\pretocmd{\subsection}{\syncleft}{}{}
\pretocmd{\subsubsection}{\syncleft}{}{}
\newcommand{\biheadingfont}{\songfont}
\newcommand{\bitocentry}[2]{%
  \ifshoworiginal
    \ifshowtranslation
      \ifstrequal{#1}{#2}{#1}{#1\space #2}%
    \else
      #1%
    \fi
  \else
    #2%
  \fi
}
\pdfstringdefDisableCommands{%
  \def\bitocentry#1#2{#1}%
}
\newcommand{\biheading}[7]{%
  \syncleft#1\bikeepwithnext[#3]\phantomsection\addcontentsline{toc}{#2}{\bitocentry{#6}{#7}}%
  \ifbilingualcolumns
    {#4 #6\par}%
    \switchcolumn
    {#4\biheadingfont #7\par}%
    \switchcolumn*%
    \global\intranslationfalse
    \par\addvspace{#5}%
    \switchcolumn
    \par\addvspace{#5}%
    \switchcolumn*%
  \else
    \ifshoworiginal{#4 #6\par}\fi
    \ifshowtranslation{#4\biheadingfont #7\par}\fi
    \addvspace{#5}%
  \fi
}
\newcommand{\termsectionheading}[1]{%
  \par\Needspace{9\baselineskip}%
  \vspace{0.65em}%
  {\songfont\bfseries #1\par}%
  \vspace{0.2em}%
}
% 章節換頁：
%   雙語 paracol 欄內不要用 \clearpage；它會在同步兩欄時吐出只有欄線/頁碼的空頁。
%   \newpage 足以讓下一個 section 起新頁，又不會強制清空所有待排材料。
%   單語模式不在 paracol 內，仍可用 \clearpage。
\newcommand{\bisectionpagebreak}{\ifbilingualcolumns\newpage\else\clearpage\fi}
\newcommand{\bisection}[2]{%
  \biheading{\iffirstbilingualsection\global\firstbilingualsectionfalse\else\bisectionpagebreak\fi}{section}{6}{\centering\Large\bfseries}{0.8em}{#1}{#2}%
}
\newcommand{\bisubsection}[2]{%
  \biheading{}{subsection}{5}{\large\bfseries}{0.55em}{#1}{#2}%
}
\newcommand{\bisubsubsection}[2]{%
  \biheading{}{subsubsection}{4}{\normalsize\bfseries}{0.45em}{#1}{#2}%
}
\newcommand{\bicompositeheading}[4]{%
  \biheading{}{subsection}{5}{\large\bfseries}{0.16em}{#1}{#2}%
  \biheading{}{subsubsection}{3}{\normalsize\bfseries}{0.45em}{#3}{#4}%
}
\newif\iflawboxfirstitem
\lawboxfirstitemfalse
\newcommand{\lawboxprespace}[1]{%
  \par
  \iflawboxfirstitem
    \global\lawboxfirstitemfalse
  \else
    \addvspace{#1}%
  \fi
}
\newtcolorbox{lawboxinner}{
  caseboxbase,
  colframe=sourceline!85!black,
  colback=white,
  left=1.05em,
  right=1.05em,
  top=0.62em,
  bottom=0.62em,
  before upper={\global\inlawboxtrue\global\lawboxfirstitemtrue\global\lawboxhaspendingrnfalse\ifintranslation\lawboxtransfont\else\lawboxsourcefont\fi\RaggedRight\setlength{\parindent}{0pt}\setlength{\parskip}{0pt}\ifintranslation\setstretch{\lawboxtransstretch}\else\setstretch{\lawboxsourcestretch}\fi},
  after upper={\global\lawboxhaspendingrnfalse\global\inlawboxfalse}
}
\tcbset{
  lawboxpartbase/.style={
    enhanced,
    breakable=false,
    width=0.985\linewidth,
    colback=white,
    frame hidden,
    boxrule=0pt,
    arc=0pt,
    left=1.05em,
    right=1.05em,
    boxsep=1.5mm,
    before skip=0pt,
    after skip=0pt,
    before upper={\global\inlawboxtrue\global\lawboxfirstitemtrue\global\lawboxhaspendingrnfalse\ifintranslation\lawboxtransfont\else\lawboxsourcefont\fi\RaggedRight\setlength{\parindent}{0pt}\setlength{\parskip}{0pt}\ifintranslation\setstretch{\lawboxtransstretch}\else\setstretch{\lawboxsourcestretch}\fi},
    after upper={\global\lawboxhaspendingrnfalse\global\inlawboxfalse},
    borderline west={0.28pt}{0pt}{sourceline!85!black},
    borderline east={0.28pt}{0pt}{sourceline!85!black},
  },
  lawboxpartfirst/.style={
    lawboxpartbase,
    top=0.62em,
    bottom=0.04em,
    before skip=\lawboxblockbeforeskip,
    borderline north={0.28pt}{0pt}{sourceline!85!black},
  },
  lawboxpartmiddle/.style={
    lawboxpartbase,
    top=0.04em,
    bottom=0.04em,
  },
  lawboxpartlast/.style={
    lawboxpartbase,
    top=0.04em,
    bottom=0.62em,
    after skip=0pt,
    borderline south={0.28pt}{0pt}{sourceline!85!black},
  }
}
\newtcolorbox{lawboxpartinner}[1]{#1}
\newtcolorbox{termboxinner}{
  caseboxbase,
  colframe=noteblue!70!black,
  colback=casepaper,
  before upper={\songfont\setlength{\parindent}{0pt}\setlength{\parskip}{0.35em}}
}
\newtcolorbox{deboxinner}{
  caseboxbase,
  colframe=notegray!72!black,
  colback=white,
  left=1.05em,
  right=1.05em,
  top=0.62em,
  bottom=0.62em,
  before upper={\global\inlawboxtrue\global\lawboxfirstitemtrue\global\lawboxhaspendingrnfalse\small\setlength{\parindent}{0pt}\setlength{\parskip}{0.35em}},
  after upper={\global\lawboxhaspendingrnfalse\global\inlawboxfalse}
}
\NewEnviron{lawbox}{
  \begin{lawboxinner}\BODY\end{lawboxinner}
}
\NewEnviron{lawboxpart}[2]{
  \begin{lawboxpartinner}{lawboxpart#1,equal height group=#2}\BODY\end{lawboxpartinner}%
  \ifstrequal{#1}{last}{\global\lawboxpartendedblocktrue}{}%
}
\NewEnviron{termbox}{
  \par\addvspace{0.55em}
  {\color{noteblue!70!black}\rule{\linewidth}{0.28pt}}\par
  \begingroup
    \songfont\small\color{caseink}
    \setlength{\parindent}{0pt}
    \setlength{\parskip}{0.24em}
    \BODY
    \par
  \endgroup
  {\color{noteblue!70!black}\rule{\linewidth}{0.28pt}}\par
  \addvspace{0.55em}
}
\NewEnviron{debox}{
  \begin{deboxinner}\BODY\end{deboxinner}
}
\providecommand{\paragraphmark}[1]{}
\newcommand{\lawboxtitleafter}{\addvspace{0.06em}}
\newcommand{\lawtitle}[1]{\lawboxprespace{0.22em}{\lawboxtransfont\bfseries\lawboxuserandnumber #1}\par\lawboxtitleafter}
\newcommand{\absno}[2]{\lawboxprespace{0.04em}\noindent\lawboxuserandnumber\hangindent=2.6em\hangafter=1\makebox[2.6em][l]{(#1)}#2\par}
\newcommand{\nrno}[2]{\lawboxprespace{0pt}\noindent\lawboxuserandnumber\hspace*{2.6em}\hangindent=4.4em\hangafter=1\makebox[1.8em][l]{#1.}#2\par}
\newcommand{\letterno}[2]{\lawboxprespace{0pt}\noindent\lawboxuserandnumber\hspace*{4.4em}\hangindent=6.0em\hangafter=1\makebox[1.6em][l]{#1)}#2\par}
\newcommand{\sourcelawtitle}[1]{\lawboxprespace{0.22em}{\lawboxsourcefont\bfseries\lawboxuserandnumber #1}\par\lawboxtitleafter}
\newcommand{\sourcelawsubtitle}[1]{\lawboxprespace{0.04em}{\lawboxsourcefont\bfseries\lawboxuserandnumber #1}\par\lawboxtitleafter}
\newcommand{\sourceabsno}[2]{\lawboxprespace{0.04em}\noindent\lawboxuserandnumber\hangindent=2.45em\hangafter=1\makebox[2.45em][l]{(#1)}#2\par}
\newcommand{\sourcenrno}[2]{\lawboxprespace{0pt}\noindent\lawboxuserandnumber\hspace*{2.45em}\hangindent=4.25em\hangafter=1\makebox[1.8em][l]{#1.}#2\par}
\newcommand{\sourceletterno}[2]{\lawboxprespace{0pt}\noindent\lawboxuserandnumber\hspace*{4.25em}\hangindent=5.85em\hangafter=1\makebox[1.6em][l]{#1)}#2\par}
\newcommand{\sourcequotepara}[1]{\lawboxprespace{0.04em}\noindent\lawboxuserandnumber\hangindent=1.2em\hangafter=1 #1\par}
\newcommand{\sourcelawline}[1]{\lawboxprespace{0.04em}\noindent\lawboxuserandnumber #1\par}
\newcommand{\lawline}[1]{\lawboxprespace{0.04em}\noindent\lawboxuserandnumber #1\par}
\newcommand{\pendingtranslation}{\par{\songfont\footnotesize\color{pendingink}（中文譯文待補。）}\par}
\newcommand{\tocpart}[1]{\par\addvspace{0.52em}{\footnotesize\bfseries\color{pendingink}#1\par}\addvspace{0.16em}}
\newcommand{\tocdashline}{\par\addvspace{0.48em}\noindent{\color{notegray}\xleaders\hbox to 0.82em{\hfil\rule{0.42em}{0.28pt}\hfil}\hskip\linewidth}\par\addvspace{0.38em}}
% \tocpanel{font-cmd}{content}：將目錄置中於所在欄寬（雙語欄適用）。
% A3 橫式使用 0.58\linewidth，A4 橫式使用 0.84\linewidth，
% 使兩種紙張下目錄欄內容實際寬度相近（均約 100–105 mm）。
\newcommand{\tocpanel}[2]{%
  \makebox[\linewidth][c]{%
    \ifx\bilingualpaper\bilingualpaperaThree
      \begin{minipage}{0.58\linewidth}%
    \else
      \begin{minipage}{0.84\linewidth}%
    \fi
    \toccolstyle
    #1#2%
    \end{minipage}%
  }%
}
% ===== 首頁簡目錄的兩層 description list 縮排 =====
% enumitem 的 description list 可把每一列拆成：
%   [label]  labelsep  body text
% 這裡的「起點」都以所在 minipage/欄位的左緣為 0。
%
% 核心規則：
%   labelindent = label 欄位從左緣開始的位置。
%   labelwidth  = label 欄位本身寬度；標籤太長（如 Abw. Meinung）就加大它。
%   labelsep    = label 欄位右緣到正文的距離。
%   leftmargin  = 正文 body text 的起點。判斷層級視覺時主要看這個。
%
% 所以：
%   想讓「Begründung / 判決理由」整列正文往右或往左：改 tocitems 的 leftmargin。
%   想讓「A. Verfahren und Gegenstand / A. 程序與審查標的」正文往右或往左：
%     改 tocsubitems 的 leftmargin。
%   想讓 A./B./C. 這些標籤本身往右或往左，但正文不動：改 tocsubitems 的 labelindent。
%   想讓 A./B./C. 與正文間距變大或變小：改 tocsubitems 的 labelsep。
%
% 層級原則：
%   tocsubitems.leftmargin 必須大於 tocitems.leftmargin，
%   否則子層級正文會看起來比「Gründe / 理由」還靠左。
\newlist{tocitems}{description}{1}
\setlist[tocitems]{%
  labelindent=0pt,    % 第一層 label 欄位起點；0pt 表示貼齊目錄欄左緣。
  leftmargin=2.54cm,  % 第一層正文起點；例如 Begründung / 判決理由 的左緣。
  labelwidth=2.16cm,  % 第一層 label 欄位寬度；需容納 Leitsätze、Abw. Meinung、不同意見等。
  labelsep=0.32cm,    % 第一層 label 與正文的水平間距；只改間距，不改正文起點。
  itemsep=0.28em,     % 第一層各項之間的垂直距離。
  topsep=0.20em,      % 第一層 list 上下與前後內容的垂直距離。
  parsep=0pt,         % 同一 item 內段落之間的距離；目錄項通常不需要。
  partopsep=0pt,      % list 若緊接段落開始時的額外距離；關閉以免目錄高度飄動。
  align=left,         % label 欄內左對齊；長短標籤不靠右貼齊。
  font=\normalfont\bfseries % label 的字型；正文不受這行影響。
}
\newlist{tocsubitems}{description}{1}
\setlist[tocsubitems]{%
  labelindent=1cm,    % 第二層 label 起點；只控制 A./B./C. 本身的位置。
  leftmargin=3.00cm,  % 第二層正文起點；必須大於 2.54cm，才會比 Begründung / 判決理由更右。
  labelwidth=0.5cm,   % 第二層 label 欄寬；A./B./C. 很短，可小於第一層。
  labelsep=1.5cm,    % 第二層 label 與正文間距；A. 與 Verfahren / 程序 的距離看這裡。
  itemsep=0.12em,     % 第二層各項較密，避免 A.–E. 佔太高。
  topsep=0.04em,      % 第二層 list 與 Gründe / 理由之間的垂直距離。
  parsep=0pt,         % 同一第二層 item 內段落距離；目錄項通常不需要。
  partopsep=0pt,      % 關閉額外段落起始距離，避免版面不穩。
  align=left,         % A./B./C. label 欄內左對齊。
  font=\normalfont\bfseries, % A./B./C. label 加粗；正文不受這行影響。
  before={}           % 不在第二層 list 前自動插入額外內容。
}
\newlist{termitems}{description}{1}
\ifbilingualcolumns
  \setlist[termitems]{%
    leftmargin=5.2cm,
    labelwidth=4.8cm,
    labelsep=0.35cm,
    itemsep=0.16em,
    topsep=0.2em,
    parsep=0pt,
    partopsep=0pt,
    font=\normalfont\bfseries,
    style=nextline
  }
\else
  \setlist[termitems]{%
    leftmargin=0pt,
    labelwidth=0pt,
    labelsep=0pt,
    itemsep=0.24em,
    topsep=0.2em,
    parsep=0pt,
    partopsep=0pt,
    font=\normalfont\bfseries,
    style=nextline
  }
\fi
\newlist{abbritems}{description}{1}
\setlist[abbritems]{%
  leftmargin=2.38cm,
  labelwidth=2.02cm,
  labelsep=0.28cm,
  itemsep=0.16em,
  topsep=0.2em,
  parsep=0pt,
  partopsep=0pt,
  align=left,
  font=\normalfont\bfseries
}
\ifbilingualcolumns
  \newenvironment{termcolumns}{\begin{multicols}{2}}{\end{multicols}}
\else
  \newenvironment{termcolumns}{}{}
\fi

% ===== 編者註腳開關 =====
% 一般版預設不顯示編者註，避免正文腳註過密。
% 若開啟 \classroommodeon，GG/條文編者註仍會自動顯示，無需另開本開關。
% 若一般版也要顯示編者註，可在單案檔 \input 後加 \showeditornotestrue。
\newif\ifshoweditornotes
\showeditornotesfalse

% ===== 目錄排版輔助 =====
\newcommand{\toccolstyle}{\raggedright\arraybackslash}
\newcommand{\tocsingleindent}{\leftskip=1.45cm\rightskip=1.2cm}

% ===== 行內引文環境（德文原文與中文譯文各一，帶左側灰色邊線）=====
\newcommand{\quoteparbreak}{\par\addvspace{0.48em}}
\newcommand{\sourcequoteinner}[1]{%
  \begin{tcolorbox}[
    enhanced, breakable, width=\dimexpr\linewidth-0.8em\relax, frame hidden,
    colback=white, coltext=caseink, boxrule=0pt,
    borderline west={0.55pt}{0.88em}{notegray},
    left=1.78em, right=0.72em, top=0.18em, bottom=0.18em,
    boxsep=0pt, before skip=0.24em, after skip=0.24em,
    before upper={\normalfont\footnotesize\color{caseink}\setlength{\parindent}{0pt}\setlength{\parskip}{0.34em}\raggedright\setlength{\rightskip}{0pt plus 1em}}
  ]%
  #1\par
  \end{tcolorbox}%
}
\newcommand{\transquoteinner}[1]{%
  \begin{tcolorbox}[
    enhanced, breakable, width=\dimexpr\linewidth-0.8em\relax, frame hidden,
    colback=white, coltext=caseink, boxrule=0pt,
    borderline west={0.55pt}{0.88em}{notegray},
    left=1.78em, right=0.72em, top=0.18em, bottom=0.18em,
    boxsep=0pt, before skip=0.24em, after skip=0.24em,
    before upper={\songfont\small\color{caseink}\setlength{\parindent}{0pt}\setlength{\parskip}{0.34em}\raggedright\setlength{\rightskip}{0pt plus 1em}}
  ]%
  #1\par
  \end{tcolorbox}%
}

% ===== 大綱（Gliederung）Rn 輔助命令 =====
\newcommand{\outrn}[2]{\enspace{\normalfont\footnotesize\color{pendingink}(Rn\,#1–#2)}}
\newcommand{\outrnsingle}[1]{\enspace{\normalfont\footnotesize\color{pendingink}(Rn\,#1)}}
\newcommand{\outrnzh}[2]{\enspace{\normalfont\footnotesize\color{pendingink}（第\,#1–#2\,段）}}
\newcommand{\outrnzhs}[1]{\enspace{\normalfont\footnotesize\color{pendingink}（第\,#1\,段）}}
% ===== 大綱雙語對照表格輔助命令（三欄：段碼 │ 德文 │ 中文）=====
% 欄 1：段碼（由欄定義自動格式化：小字、等寬、右對齊、pendingink）
% 欄 2：德文標籤＋正文
% 欄 3：中文標籤＋正文
%
% 無段碼的列（\omainrow、\osecrow），欄 1 以 & 空置。
% 有段碼的列（\osecrowrn、\osubrow、\ointrow），#1 為預格式化字串
%   例：\osubrow{2–49}{I.}{...}{I.}{...}
%       \ointrow{107}{...}{...}
%
% \opartrow：段組標題（橫跨三欄）
\newcommand{\opartrow}[2]{%
  \noalign{\vspace{0.30em}}%
  \multicolumn{3}{@{}l@{}}{%
    \footnotesize\bfseries\color{pendingink}#1\hfill\songfont#2%
  }\\[0.06em]%
}
% \odashrow：虛線分隔（橫跨三欄）
\newcommand{\odashrow}{%
  \noalign{\vspace{0.28em}}%
  \multicolumn{3}{@{}l@{}}{\color{notegray}%
    \xleaders\hbox to 0.82em{\hfil\rule{0.42em}{0.28pt}\hfil}\hskip 0.88\linewidth}%
  \\[0.24em]%
}
% \odissentpart：不同意見書區段標頭。比一般 \opartrow 更像附錄性轉場。
\newcommand{\odissentpart}[2]{%
  \noalign{\vspace{0.42em}}%
  \multicolumn{3}{@{}p{0.88\textwidth}@{}}{%
    \color{notegray}\rule{\linewidth}{0.18pt}%
  }\\[-0.18em]%
  \multicolumn{3}{@{}p{0.88\textwidth}@{}}{%
    \makebox[\linewidth][s]{%
      \footnotesize\bfseries\color{pendingink}#1%
      \hfill{\songfont #2}%
    }%
  }\\[-0.02em]%
}
% \odissentopinion：不同意見書的署名與一行摘要，右欄保留段碼範圍。
\newcommand{\odissentopinion}[5]{%
  \footnotesize\hangindent=0.52cm\hangafter=1%
    {\bfseries\color{caseink}#2}\enspace{\color{notegray}\textperiodcentered}\enspace\color{caseink}#3%
  &\footnotesize\hangindent=0.52cm\hangafter=1%
    {\songfont\bfseries\color{caseink}#4}\enspace{\color{notegray}\textperiodcentered}\enspace\songfont\color{caseink}#5%
  &#1%
  \\[0.12em]%
}
% \omainrow：頂層項目，無段碼（Leitsätze、Urteil、Gründe …）
\newcommand{\omainrow}[4]{%
  \footnotesize\hangindent=1.92cm\hangafter=1%
    \makebox[1.92cm][l]{\bfseries\color{caseink}#1}\color{caseink}#2%
  &\footnotesize\hangindent=1.92cm\hangafter=1%
    \makebox[1.92cm][l]{\bfseries\color{caseink}#3}\songfont\color{caseink}#4%
  &%
  \\[0.15em]%
}
% \omainrowrn：頂層項目，有段碼（不同意見等標籤寬於 \osecrow 框者）
\newcommand{\omainrowrn}[5]{%
  \footnotesize\hangindent=1.92cm\hangafter=1%
    \makebox[1.92cm][l]{\bfseries\color{caseink}#2}\color{caseink}#3%
  &\footnotesize\hangindent=1.92cm\hangafter=1%
    \makebox[1.92cm][l]{\bfseries\color{caseink}#4}\songfont\color{caseink}#5%
  &#1%
  \\[0.15em]%
}
% \osecrow：一級子項，無段碼（A.、C.、D. 等有子項者）
\newcommand{\osecrow}[4]{%
  \footnotesize\hspace{1.10cm}\hangindent=1.98cm\hangafter=1%
    \makebox[0.86cm][l]{\bfseries\color{caseink}#1}\color{caseink}#2%
  &\footnotesize\hspace{1.10cm}\hangindent=1.98cm\hangafter=1%
    \makebox[0.86cm][l]{\bfseries\color{caseink}#3}\songfont\color{caseink}#4%
  &%
  \\[0.11em]%
}
% \osecrowrn：一級子項，有段碼（B.、E. 等完整段落範圍者）
% #1=段碼字串（如 101–106）, #2=DE標籤, #3=DE正文, #4=ZH標籤, #5=ZH正文
\newcommand{\osecrowrn}[5]{%
  \footnotesize\hspace{1.10cm}\hangindent=1.98cm\hangafter=1%
    \makebox[0.86cm][l]{\bfseries\color{caseink}#2}\color{caseink}#3%
  &\footnotesize\hspace{1.10cm}\hangindent=1.98cm\hangafter=1%
    \makebox[0.86cm][l]{\bfseries\color{caseink}#4}\songfont\color{caseink}#5%
  &#1%
  \\[0.11em]%
}
% \osubrow：二級子項，有段碼範圍（I.、II.、A.I. …）
% #1=段碼字串, #2=DE標籤, #3=DE正文, #4=ZH標籤, #5=ZH正文
\newcommand{\osubrow}[5]{%
  \footnotesize\hspace{1.40cm}\hangindent=2.30cm\hangafter=1%
    \makebox[0.90cm][l]{\bfseries\color{caseink}#2}\color{caseink}#3%
  &\footnotesize\hspace{1.40cm}\hangindent=2.30cm\hangafter=1%
    \makebox[0.90cm][l]{\bfseries\color{caseink}#4}\songfont\color{caseink}#5%
  &#1%
  \\[0.09em]%
}
% \ointrow：引入段落，有單一段碼，無標籤
% #1=段碼字串, #2=DE正文, #3=ZH正文
\newcommand{\ointrow}[3]{%
  \footnotesize\hspace{0.52cm}\hangindent=0.52cm\hangafter=1\color{caseink}#2%
  &\footnotesize\hspace{0.52cm}\hangindent=0.52cm\hangafter=1\songfont\color{caseink}#3%
  &#1%
  \\[0.09em]%
}
% \osubsubrow：三級子項（1.、2.、3.）
% #1=段碼字串, #2=DE標籤, #3=DE正文, #4=ZH標籤, #5=ZH正文
\newcommand{\osubsubrow}[5]{%
  \footnotesize\hspace{1.80cm}\hangindent=2.48cm\hangafter=1%
    \makebox[0.68cm][l]{\bfseries\color{caseink}#2}\color{caseink}#3%
  &\footnotesize\hspace{1.80cm}\hangindent=2.48cm\hangafter=1%
    \makebox[0.68cm][l]{\bfseries\color{caseink}#4}\songfont\color{caseink}#5%
  &#1%
  \\[0.08em]%
}
% \osubsubsubrow：四級子項（a）、b））
% #1=段碼字串, #2=DE標籤, #3=DE正文, #4=ZH標籤, #5=ZH正文
\newcommand{\osubsubsubrow}[5]{%
  \footnotesize\hspace{2.10cm}\hangindent=2.68cm\hangafter=1%
    \makebox[0.58cm][l]{\bfseries\color{caseink}#2}\color{caseink}#3%
  &\footnotesize\hspace{2.10cm}\hangindent=2.68cm\hangafter=1%
    \makebox[0.58cm][l]{\bfseries\color{caseink}#4}\songfont\color{caseink}#5%
  &#1%
  \\[0.07em]%
}
% \osubsubsubsubrow：五級子項（aa）、bb））
% 標籤框 0.62cm：足以容納「aa)」在 footnotesize bold 下（約 0.50cm）並留有間距
% #1=段碼字串, #2=DE標籤, #3=DE正文, #4=ZH標籤, #5=ZH正文
\newcommand{\osubsubsubsubrow}[5]{%
  \footnotesize\hspace{2.38cm}\hangindent=3.06cm\hangafter=1%
    \makebox[0.68cm][l]{\bfseries\color{caseink}#2}\color{caseink}#3%
  &\footnotesize\hspace{2.38cm}\hangindent=3.06cm\hangafter=1%
    \makebox[0.68cm][l]{\bfseries\color{caseink}#4}\songfont\color{caseink}#5%
  &#1%
  \\[0.06em]%
}
% \outlinepage：雙語對照大綱頁包裝環境（longtable 自動跨頁）
% 三欄：德文欄 + 中文欄 + 段碼欄（最右，不折行）
% 段碼欄寬度：1.80cm，足以容納最長段碼如「309–348」
% DE/ZH 欄寬：(0.87\textwidth - 2.08cm) / 2
%   驗算：2×欄 + 0.05\textwidth + 0.28cm gap + 1.80cm = 0.87tw - 2.08cm + 0.05tw + 2.08cm = 0.92tw ✓
\newenvironment{outlinepage}{%
  \vspace*{0.40cm}%
  \setlength{\LTleft}{0.04\textwidth}%
  \setlength{\LTright}{0.04\textwidth}%
  \setlength{\tabcolsep}{0pt}%
  % 大綱列距由 longtable 的 \arraystretch 與各 row macro 結尾的 \\[...em] 共同決定。
  % 這裡控制「每一列本身的表格 strut 高度」；數值越大，A./I./1. 各項之間越像有空行。
  % A3 原本用 0.62，視覺上沒有空行；A4 也沿用同值，避免切到 A4 後項目被撐開。
  % 若日後覺得大綱太擠，只微調這個值（例如 0.68），不要改各 \omainrow/\osubrow 的縮排。
  \renewcommand{\arraystretch}{0.62}%
  \begin{longtable}{%
    >{\vspace{0pt}\raggedright\arraybackslash}p{\dimexpr(0.87\textwidth-2.08cm)/2\relax}%
    @{\hspace{0.025\textwidth}}%
    !{\color{notegray}\vrule width 0.12pt}%
    @{\hspace{0.025\textwidth}}%
    >{\vspace{0pt}\raggedright\arraybackslash}p{\dimexpr(0.87\textwidth-2.08cm)/2\relax}%
    @{\hspace{0.28cm}}%
    >{\vspace{0pt}\footnotesize\normalfont\color{pendingink}\raggedright\arraybackslash}p{1.80cm}%
    @{}%
  }%
  % 首頁欄標籤：不要橫跨置中；分別放在德文欄與中文欄的實際欄位起點。
  \footnotesize\bfseries\color{sourceink}Gliederung%
  &\footnotesize\bfseries\songfont\color{sourceink}大綱%
  &%
  \\[0.36em]%
  \endfirsthead%
  % 續頁欄標籤：同樣分欄定位，以灰色細體區分；無需標注「續」。
  \footnotesize\color{pendingink}Gliederung%
  &\footnotesize\songfont\color{pendingink}大綱%
  &%
  \\[0.24em]%
  \endhead%
}{%
  \end{longtable}%
}


% 課堂模式：取消下行註解即開啟；註解掉即回到一般版。
% 開啟時會在德文原文縮寫後加註，並自動顯示 GG/條文編者註。
% \classroommodeon

\newcommand{\documentversion}{v1.1}
\newcommand{\ggversionde}{\ifggversionnotede\else GG-Fassung: 25. Februar 1975, nach der 31. GG-Aenderung (BGBl. I 1972, S. 1305) und vor der 32. GG-Aenderung (BGBl. I 1975, S. 1901). \global\ggversionnotedetrue\fi}
\newcommand{\ggversionzh}{\ifggversionnotezh\else 以下引用《基本法》1975年2月25日判決時有效版本；即第31次《基本法》修正（BGBl. I 1972, S. 1305）後、第32次《基本法》修正（BGBl. I 1975, S. 1901）前之版本。 \global\ggversionnotezhtrue\fi}
% bverfge_gg.tex — 德國墮胎案 GG 條文腳註共用定義
% 使用方式：在各案 .tex 中先定義 \ggversionde/\ggversionzh，再 % bverfge_gg.tex — 德國墮胎案 GG 條文腳註共用定義
% 使用方式：在各案 .tex 中先定義 \ggversionde/\ggversionzh，再 % bverfge_gg.tex — 德國墮胎案 GG 條文腳註共用定義
% 使用方式：在各案 .tex 中先定義 \ggversionde/\ggversionzh，再 \input{../../共用LaTeX/bverfge_gg}

% ===== GG 版本旗標（\ggversionde/zh 由各案在 \input{../../共用LaTeX/bverfge_gg} 之前定義）=====
\newif\ifggversionnotede
\newif\ifggversionnotezh

% ===== 編者註腳基礎設施（首次標記模式）=====
\newif\ifeditornotelabelde
\newif\ifeditornotelabelzh
\newcommand{\editornotelabelde}{\ifeditornotelabelde\else\emph{Hrsg.-Anm. (nachfolgend ohne Kennzeichnung)}: \global\editornotelabeldetrue\fi}
\newcommand{\editornotelabelzh}{\ifeditornotelabelzh\else 編者註（下同）：\global\editornotelabelzhtrue\fi}
\newcommand{\editornotede}[1]{\ifshoweditornotes\footnote{\normalfont\footnotesize\editornotelabelde #1}\else\ifclassroommode\footnote{\normalfont\footnotesize\editornotelabelde #1}\fi\fi}
\newcommand{\editornotezh}[1]{\ifshoweditornotes\footnote{\normalfont\footnotesize\editornotelabelzh #1}\else\ifclassroommode\footnote{\normalfont\footnotesize\editornotelabelzh #1}\fi\fi}

% ===== GG 引用系統（\ggversionde/zh 由各案在 \input 前定義）=====
\newcommand{\ggmarknotede}[1]{\global\expandafter\let\csname ggnoteseen@de@#1\endcsname\empty}
\newcommand{\ggmarknotezh}[1]{\global\expandafter\let\csname ggnoteseen@zh@#1\endcsname\empty}
\newcommand{\ggnotedeonce}[2]{\ifcsname ggnoteseen@de@#1\endcsname\else\global\expandafter\let\csname ggnoteseen@de@#1\endcsname\empty\ggnotede{#2}\fi}
\newcommand{\ggnotezhonce}[2]{\ifcsname ggnoteseen@zh@#1\endcsname\else\global\expandafter\let\csname ggnoteseen@zh@#1\endcsname\empty\ggnotezh{#2}\fi}
\newcommand{\ggnotede}[1]{\editornotede{\ggversionde #1}}
\newcommand{\ggnotezh}[1]{\editornotezh{\ggversionzh #1}}

% ===== Art. 1 =====
\newcommand{\ggfnArtOneOneDE}{\ggnotedeonce{art-1-1}{Art. 1 Abs. 1: ``Die Wuerde des Menschen ist unantastbar. Sie zu achten und zu schuetzen ist Verpflichtung aller staatlichen Gewalt.''}}
\newcommand{\ggfnArtOneOneZH}{\ggnotezhonce{art-1-1}{第1條第1項：「人之尊嚴不可侵犯。尊重及保護之，為一切國家權力之義務。」}}
\newcommand{\ggfnArtOneOneTwoDE}{\ggnotedeonce{art-1-1-2}{Art. 1 Abs. 1 Satz 2: ``Sie zu achten und zu schuetzen ist Verpflichtung aller staatlichen Gewalt.''}}
\newcommand{\ggfnArtOneOneTwoZH}{\ggnotezhonce{art-1-1-2}{第1條第1項第2句：「尊重及保護之，為一切國家權力之義務。」}}

% ===== Art. 2 =====
\newcommand{\ggfnArtTwoOneDE}{\ggnotedeonce{art-2-1}{Art. 2 Abs. 1: ``Jeder hat das Recht auf die freie Entfaltung seiner Persoenlichkeit, soweit er nicht die Rechte anderer verletzt und nicht gegen die verfassungsmaessige Ordnung oder das Sittengesetz verstoesst.''}}
\newcommand{\ggfnArtTwoOneZH}{\ggnotezhonce{art-2-1}{第2條第1項：「人人有自由發展其人格之權利，但不得侵害他人權利，亦不得違反合憲秩序或道德律。」}}
% 簡易預設版（墮胎一案以 \renewcommand 覆寫為含條件邏輯之版本）
\newcommand{\ggfnArtTwoTwoDE}{\ggnotedeonce{art-2-2}{Art. 2 Abs. 2: ``Jeder hat das Recht auf Leben und koerperliche Unversehrtheit. Die Freiheit der Person ist unverletzlich. In diese Rechte darf nur auf Grund eines Gesetzes eingegriffen werden.''}\ggmarknotede{art-2-2-1}}
\newcommand{\ggfnArtTwoTwoZH}{\ggnotezhonce{art-2-2}{第2條第2項：「人人有生命與身體完整性之權利。人身自由不可侵犯。此等權利僅得依法律加以干預。」}\ggmarknotezh{art-2-2-1}}
\newcommand{\ggfnArtTwoTwoOneDE}{\ggnotedeonce{art-2-2-1}{Art. 2 Abs. 2 Satz 1: ``Jeder hat das Recht auf Leben und koerperliche Unversehrtheit.''}}
\newcommand{\ggfnArtTwoTwoOneZH}{\ggnotezhonce{art-2-2-1}{第2條第2項第1句：「人人有生命與身體完整性之權利。」}}
\newcommand{\ggfnArtTwoTwoThreeDE}{\ggnotedeonce{art-2-2-3}{Art. 2 Abs. 2 Satz 3: ``In diese Rechte darf nur auf Grund eines Gesetzes eingegriffen werden.''}}
\newcommand{\ggfnArtTwoTwoThreeZH}{\ggnotezhonce{art-2-2-3}{第2條第2項第3句：「此等權利僅得依法律加以干預。」}}
% 墮胎一案特有：第2條第2項第2–3句（Art. 2(2) 去掉第1句後的「餘文」）
\newcommand{\ggfnArtTwoTwoRestDE}{\ggnotedeonce{art-2-2-rest}{Art. 2 Abs. 2 Saetze 2-3: ``Die Freiheit der Person ist unverletzlich. In diese Rechte darf nur auf Grund eines Gesetzes eingegriffen werden.''}\ggmarknotede{art-2-2}}
\newcommand{\ggfnArtTwoTwoRestZH}{\ggnotezhonce{art-2-2-rest}{第2條第2項第2、3句：「人身自由不可侵犯。此等權利僅得依法律加以干預。」}\ggmarknotezh{art-2-2}}

% ===== Art. 3 =====
\newcommand{\ggfnArtThreeDE}{\ggnotedeonce{art-3}{Art. 3: ``Alle Menschen sind vor dem Gesetz gleich. Maenner und Frauen sind gleichberechtigt. Niemand darf wegen seines Geschlechtes, seiner Abstammung, seiner Rasse, seiner Sprache, seiner Heimat und Herkunft, seines Glaubens, seiner religioesen oder politischen Anschauungen benachteiligt oder bevorzugt werden.''}}
\newcommand{\ggfnArtThreeZH}{\ggnotezhonce{art-3}{第3條：「所有人在法律前一律平等。男女享有平等權利。任何人不得因其性別、血統、種族、語言、故鄉及出身、信仰、宗教或政治見解而受不利益或優惠。」}}
\newcommand{\ggfnArtThreeTwoDE}{\ggnotedeonce{art-3-2}{Art. 3 Abs. 2: ``Maenner und Frauen sind gleichberechtigt.''}}
\newcommand{\ggfnArtThreeTwoZH}{\ggnotezhonce{art-3-2}{第3條第2項：「男女享有平等權利。」}}

% ===== Art. 4 =====
\newcommand{\ggfnArtFourOneDE}{\ggnotedeonce{art-4-1}{Art. 4 Abs. 1: ``Die Freiheit des Glaubens, des Gewissens und die Freiheit des religioesen und weltanschaulichen Bekenntnisses sind unverletzlich.''}}
\newcommand{\ggfnArtFourOneZH}{\ggnotezhonce{art-4-1}{第4條第1項：「信仰自由、良心自由，以及宗教與世界觀信念自由，不可侵犯。」}}

% ===== Art. 5 =====
\newcommand{\ggfnArtFiveOneDE}{\ggnotedeonce{art-5-1}{Art. 5 Abs. 1: ``Jeder hat das Recht, seine Meinung in Wort, Schrift und Bild frei zu aeussern und zu verbreiten und sich aus allgemein zugaenglichen Quellen ungehindert zu unterrichten. Die Pressefreiheit und die Freiheit der Berichterstattung durch Rundfunk und Film werden gewaehrleistet. Eine Zensur findet nicht statt.''}}
\newcommand{\ggfnArtFiveOneZH}{\ggnotezhonce{art-5-1}{第5條第1項：「人人有以言語、文字及圖像自由表達及傳播其意見，並由一般可接近來源不受阻礙獲取資訊之權利。新聞自由及廣播、電影報導自由受保障。不設審查制度。」}}

% ===== Art. 6 =====
\newcommand{\ggfnArtSixDE}{\ggnotedeonce{art-6}{Art. 6 Abs. 1: ``Ehe und Familie stehen unter dem besonderen Schutze der staatlichen Ordnung.'' Abs. 2 Satz 1: ``Pflege und Erziehung der Kinder sind das natuerliche Recht der Eltern und die zuvoerderst ihnen obliegende Pflicht.'' Abs. 4: ``Jede Mutter hat Anspruch auf den Schutz und die Fuersorge der Gemeinschaft.''}}
\newcommand{\ggfnArtSixZH}{\ggnotezhonce{art-6}{第6條第1項：「婚姻與家庭受國家秩序之特別保護。」第2項第1句：「照護及教育子女，為父母之自然權利，並為首先由其承擔之義務。」第4項：「每一母親均有請求共同體保護與照顧之權利。」}}
\newcommand{\ggfnArtSixOneDE}{\ggnotedeonce{art-6-1}{Art. 6 Abs. 1: ``Ehe und Familie stehen unter dem besonderen Schutze der staatlichen Ordnung.''}}
\newcommand{\ggfnArtSixOneZH}{\ggnotezhonce{art-6-1}{第6條第1項：「婚姻與家庭受國家秩序之特別保護。」}}
\newcommand{\ggfnArtSixTwoDE}{\ggnotedeonce{art-6-2}{Art. 6 Abs. 2: ``Pflege und Erziehung der Kinder sind das natuerliche Recht der Eltern und die zuvoerderst ihnen obliegende Pflicht. Ueber ihre Betaetigung wacht die staatliche Gemeinschaft.''}}
\newcommand{\ggfnArtSixTwoZH}{\ggnotezhonce{art-6-2}{第6條第2項：「照護及教育子女，為父母之自然權利，並為首先由其承擔之義務。國家共同體監督其行使。」}}
\newcommand{\ggfnArtSixTwoOneDE}{\ggnotedeonce{art-6-2-1}{Art. 6 Abs. 2 Satz 1: ``Pflege und Erziehung der Kinder sind das natuerliche Recht der Eltern und die zuvoerderst ihnen obliegende Pflicht.''}}
\newcommand{\ggfnArtSixTwoOneZH}{\ggnotezhonce{art-6-2-1}{第6條第2項第1句：「照護及教育子女，為父母之自然權利，並為首先由其承擔之義務。」}}
\newcommand{\ggfnArtSixFourDE}{\ggnotedeonce{art-6-4}{Art. 6 Abs. 4: ``Jede Mutter hat Anspruch auf den Schutz und die Fuersorge der Gemeinschaft.''}}
\newcommand{\ggfnArtSixFourZH}{\ggnotezhonce{art-6-4}{第6條第4項：「每一母親均有請求共同體保護與照顧之權利。」}}

% ===== Art. 12 =====
\newcommand{\ggfnArtTwelveOneDE}{\ggnotedeonce{art-12-1}{Art. 12 Abs. 1: ``Alle Deutschen haben das Recht, Beruf, Arbeitsplatz und Ausbildungsstaette frei zu waehlen. Die Berufsausuebung kann durch Gesetz oder auf Grund eines Gesetzes geregelt werden.''}}
\newcommand{\ggfnArtTwelveOneZH}{\ggnotezhonce{art-12-1}{第12條第1項：「所有德國人均有自由選擇職業、工作場所及訓練場所之權利。職業行使得以法律或基於法律加以規範。」}}

% ===== Art. 19 =====
\newcommand{\ggfnArtNineteenTwoDE}{\ggnotedeonce{art-19-2}{Art. 19 Abs. 2: ``In keinem Falle darf ein Grundrecht in seinem Wesensgehalt angetastet werden.''}}
\newcommand{\ggfnArtNineteenTwoZH}{\ggnotezhonce{art-19-2}{第19條第2項：「任何情形下，基本權利之本質內容均不得受到侵害。」}}

% ===== Art. 20 =====
\newcommand{\ggfnArtTwentyOneDE}{\ggnotedeonce{art-20-1}{Art. 20 Abs. 1: ``Die Bundesrepublik Deutschland ist ein demokratischer und sozialer Bundesstaat.''}}
\newcommand{\ggfnArtTwentyOneZH}{\ggnotezhonce{art-20-1}{第20條第1項：「德意志聯邦共和國為民主及社會之聯邦國。」}}
\newcommand{\ggfnArtTwentyThreeDE}{\ggnotedeonce{art-20-3}{Art. 20 Abs. 3: ``Die Gesetzgebung ist an die verfassungsmaessige Ordnung, die vollziehende Gewalt und die Rechtsprechung sind an Gesetz und Recht gebunden.''}}
\newcommand{\ggfnArtTwentyThreeZH}{\ggnotezhonce{art-20-3}{第20條第3項：「立法受合憲秩序拘束；行政權與司法權受法律與法拘束。」}}

% ===== Art. 26 + Art. 143（墮胎一案特有）=====
\newcommand{\ggfnArtTwoSixAndOneFourThreeDE}{\ggnotedeonce{art-26-143}{Art. 26 Abs. 1 Satz 2: ``Sie sind unter Strafe zu stellen.'' Art. 143 in der urspruenglichen Fassung (24. Mai 1949 bis 1. September 1951) enthielt Strafvorschriften zum Hochverrat; 1975 war Art. 143 bereits weggefallen.}}
\newcommand{\ggfnArtTwoSixAndOneFourThreeZH}{\ggnotezhonce{art-26-143}{第26條第1項第2句：「此等行為應規定為犯罪處罰。」第143條原始版本（1949年5月24日至1951年9月1日）含有關於叛國罪之刑罰規定；至1975年時，第143條已廢止。}}

% ===== Art. 28 =====
\newcommand{\ggfnArtTwentyEightOneDE}{\ggnotedeonce{art-28-1-1}{Art. 28 Abs. 1 Satz 1: ``Die verfassungsmaessige Ordnung in den Laendern muss den Grundsaetzen des republikanischen, demokratischen und sozialen Rechtsstaates im Sinne dieses Grundgesetzes entsprechen.''}}
\newcommand{\ggfnArtTwentyEightOneZH}{\ggnotezhonce{art-28-1-1}{第28條第1項第1句：「各邦之合憲秩序，須符合本基本法意義下共和、民主及社會法治國之原則。」}}

% ===== Art. 30 =====
\newcommand{\ggfnArtThirtyDE}{\ggnotedeonce{art-30}{Art. 30: ``Die Ausuebung der staatlichen Befugnisse und die Erfuellung der staatlichen Aufgaben ist Sache der Laender, soweit dieses Grundgesetz keine andere Regelung trifft oder zulaesst.''}}
\newcommand{\ggfnArtThirtyZH}{\ggnotezhonce{art-30}{第30條：「國家權限之行使與國家任務之履行，屬於各邦，但本基本法另有規定或容許者，不在此限。」}}

% ===== Art. 74 =====
\newcommand{\ggfnArtSeventyFourOneDE}{\ggnotedeonce{art-74-1}{Art. 74 Nr. 1: Die konkurrierende Gesetzgebung erstreckt sich auf ``das buergerliche Recht, das Strafrecht und den Strafvollzug, die Gerichtsverfassung, das gerichtliche Verfahren, die Rechtsanwaltschaft, das Notariat und die Rechtsberatung''.}}
\newcommand{\ggfnArtSeventyFourOneZH}{\ggnotezhonce{art-74-1}{第74條第1款：競合立法及於「民法、刑法及刑罰執行、法院組織、訴訟程序、律師制度、公證制度及法律諮詢」。}}
\newcommand{\ggfnArtSeventyFourSevenDE}{\ggnotedeonce{art-74-7}{Art. 74 Nr. 7: Die konkurrierende Gesetzgebung erstreckt sich auf ``die oeffentliche Fuersorge''.}}
\newcommand{\ggfnArtSeventyFourSevenZH}{\ggnotezhonce{art-74-7}{第74條第7款：競合立法及於「公共扶助」。}}
\newcommand{\ggfnArtSeventyFourTwelveDE}{\ggnotedeonce{art-74-12}{Art. 74 Nr. 12: Die konkurrierende Gesetzgebung erstreckt sich auf ``das Arbeitsrecht einschliesslich der Betriebsverfassung, des Arbeitsschutzes und der Arbeitsvermittlung sowie die Sozialversicherung einschliesslich der Arbeitslosenversicherung''.}}
\newcommand{\ggfnArtSeventyFourTwelveZH}{\ggnotezhonce{art-74-12}{第74條第12款：競合立法及於「勞動法，包括企業組織、勞動保護及就業媒合，以及社會保險，包括失業保險」。}}

% ===== Art. 77（墮胎一案特有）=====
\newcommand{\ggfnArtSevenSevenThreeDE}{\ggnotedeonce{art-77-3}{Art. 77 Abs. 3 Satz 1: ``Soweit zu einem Gesetze die Zustimmung des Bundesrates nicht erforderlich ist, kann der Bundesrat, wenn das Verfahren nach Absatz 2 beendigt ist, gegen ein vom Bundestage beschlossenes Gesetz binnen zwei Wochen Einspruch einlegen.''}}
\newcommand{\ggfnArtSevenSevenThreeZH}{\ggnotezhonce{art-77-3}{第77條第3項第1句：「法律無須聯邦參議院同意者，於第2項程序終了後，聯邦參議院得於二週內對聯邦議院通過之法律提出異議。」}}

% ===== Art. 80 =====
\newcommand{\ggfnArtEightyOneOneDE}{\ggnotedeonce{art-80-1-1}{Art. 80 Abs. 1 Satz 1: ``Durch Gesetz koennen die Bundesregierung, ein Bundesminister oder die Landesregierungen ermaechtigt werden, Rechtsverordnungen zu erlassen.''}}
\newcommand{\ggfnArtEightyOneOneZH}{\ggnotezhonce{art-80-1-1}{第80條第1項第1句：「得以法律授權聯邦政府、聯邦部長或邦政府發布法規命令。」}}

% ===== Art. 83 =====
\newcommand{\ggfnArtEightyThreeDE}{\ggnotedeonce{art-83}{Art. 83: ``Die Laender fuehren die Bundesgesetze als eigene Angelegenheit aus, soweit dieses Grundgesetz nichts anderes bestimmt oder zulaesst.''}}
\newcommand{\ggfnArtEightyThreeZH}{\ggnotezhonce{art-83}{第83條：「各邦以自己事務執行聯邦法律，但本基本法另有規定或容許者，不在此限。」}}

% ===== Art. 84 =====
\newcommand{\ggfnArtEightyFourOneDE}{\ggnotedeonce{art-84-1}{Art. 84 Abs. 1: ``Fuehren die Laender die Bundesgesetze als eigene Angelegenheit aus, so regeln sie die Einrichtung der Behoerden und das Verwaltungsverfahren, soweit nicht Bundesgesetze mit Zustimmung des Bundesrates etwas anderes bestimmen.''}}
\newcommand{\ggfnArtEightyFourOneZH}{\ggnotezhonce{art-84-1}{第84條第1項：「各邦以自己事務執行聯邦法律時，除聯邦法律經聯邦參議院同意另有規定外，由各邦規範機關之設置及行政程序。」}}
% 別名（墮胎一案使用之縮寫命名）
\let\ggfnArtEightFourOneDE\ggfnArtEightyFourOneDE
\let\ggfnArtEightFourOneZH\ggfnArtEightyFourOneZH

% ===== Art. 93 =====
\newcommand{\ggfnArtNinetyThreeOneTwoDE}{\ggnotedeonce{art-93-1-2}{Art. 93 Abs. 1 Nr. 2: Das Bundesverfassungsgericht entscheidet ``bei Meinungsverschiedenheiten oder Zweifeln ueber die foermliche und sachliche Vereinbarkeit von Bundesrecht oder Landesrecht mit diesem Grundgesetze ... auf Antrag der Bundesregierung, einer Landesregierung oder eines Drittels der Mitglieder des Bundestages''.}}
\newcommand{\ggfnArtNinetyThreeOneTwoZH}{\ggnotezhonce{art-93-1-2}{第93條第1項第2款：聯邦憲法法院就「聯邦法或邦法在形式及實質上是否與本基本法相容之意見分歧或疑義……依聯邦政府、邦政府或聯邦議院三分之一議員之聲請」作成裁判。}}
% 別名（墮胎一案使用之縮寫命名）
\let\ggfnArtNineThreeOneTwoDE\ggfnArtNinetyThreeOneTwoDE
\let\ggfnArtNineThreeOneTwoZH\ggfnArtNinetyThreeOneTwoZH

% ===== Art. 102 =====
\newcommand{\ggfnArtOneZeroTwoDE}{\ggnotedeonce{art-102}{Art. 102: ``Die Todesstrafe ist abgeschafft.''}}
\newcommand{\ggfnArtOneZeroTwoZH}{\ggnotezhonce{art-102}{第102條：「死刑廢止。」}}

% ===== Art. 103 =====
\newcommand{\ggfnArtOneZeroThreeTwoDE}{\ggnotedeonce{art-103-2}{Art. 103 Abs. 2: ``Eine Tat kann nur bestraft werden, wenn die Strafbarkeit gesetzlich bestimmt war, bevor die Tat begangen wurde.''}}
\newcommand{\ggfnArtOneZeroThreeTwoZH}{\ggnotezhonce{art-103-2}{第103條第2項：「行為之可罰性，須於行為前以法律定之，始得處罰。」}}

% ===== Art. 104 =====
\newcommand{\ggfnArtOneZeroFourOneDE}{\ggnotedeonce{art-104-1}{Art. 104 Abs. 1: ``Die Freiheit der Person kann nur auf Grund eines foermlichen Gesetzes und nur unter Beachtung der darin vorgeschriebenen Formen beschraenkt werden.''}}
\newcommand{\ggfnArtOneZeroFourOneZH}{\ggnotezhonce{art-104-1}{第104條第1項：「人身自由僅得基於形式法律，並遵守其中規定之形式加以限制。」}}

% ===== 組合腳註：Art. 3 + Art. 6（墮胎一案特有）=====
\newcommand{\ggfnArtThreeSixDE}{\ggnotedeonce{art-3-6}{Art. 3: ``Alle Menschen sind vor dem Gesetz gleich. Maenner und Frauen sind gleichberechtigt. Niemand darf wegen seines Geschlechtes, seiner Abstammung, seiner Rasse, seiner Sprache, seiner Heimat und Herkunft, seines Glaubens, seiner religioesen oder politischen Anschauungen benachteiligt oder bevorzugt werden.'' Art. 6 Abs. 1: ``Ehe und Familie stehen unter dem besonderen Schutze der staatlichen Ordnung.'' Abs. 2 Satz 1: ``Pflege und Erziehung der Kinder sind das natuerliche Recht der Eltern und die zuvoerderst ihnen obliegende Pflicht.'' Abs. 4: ``Jede Mutter hat Anspruch auf den Schutz und die Fuersorge der Gemeinschaft.''}\ggmarknotede{art-3}\ggmarknotede{art-6}\ggmarknotede{art-6-1-4}}
\newcommand{\ggfnArtThreeSixZH}{\ggnotezhonce{art-3-6}{第3條：「所有人在法律前一律平等。男女享有平等權利。任何人不得因其性別、血統、種族、語言、故鄉及出身、信仰、宗教或政治見解而受不利益或優惠。」第6條第1項：「婚姻與家庭受國家秩序之特別保護。」第2項第1句：「照護及教育子女，為父母之自然權利，並為首先由其承擔之義務。」第4項：「每一母親均有請求共同體保護與照顧之權利。」}\ggmarknotezh{art-3}\ggmarknotezh{art-6}\ggmarknotezh{art-6-1-4}}

% ===== 組合腳註：Art. 6(1) + Art. 6(4)（墮胎一案特有）=====
\newcommand{\ggfnArtSixOneFourDE}{\ggnotedeonce{art-6-1-4}{Art. 6 Abs. 1: ``Ehe und Familie stehen unter dem besonderen Schutze der staatlichen Ordnung.'' Art. 6 Abs. 4: ``Jede Mutter hat Anspruch auf den Schutz und die Fuersorge der Gemeinschaft.''}\ggmarknotede{art-6}}
\newcommand{\ggfnArtSixOneFourZH}{\ggnotezhonce{art-6-1-4}{第6條第1項：「婚姻與家庭受國家秩序之特別保護。」第4項：「每一母親均有請求共同體保護與照顧之權利。」}\ggmarknotezh{art-6}}

% ===== 組合腳註：Art. 1(1) + Art. 2(2)(1)（墮胎二案特有）=====
\newcommand{\ggfnArtOneOneTwoTwoOneDE}{\ggnotedeonce{art-1-1+2-2-1}{Art. 1 Abs. 1: ``Die Wuerde des Menschen ist unantastbar. Sie zu achten und zu schuetzen ist Verpflichtung aller staatlichen Gewalt.'' Art. 2 Abs. 2 Satz 1: ``Jeder hat das Recht auf Leben und koerperliche Unversehrtheit.''}\ggmarknotede{art-1-1}\ggmarknotede{art-2-2-1}}
\newcommand{\ggfnArtOneOneTwoTwoOneZH}{\ggnotezhonce{art-1-1+2-2-1}{第1條第1項：「人之尊嚴不可侵犯。尊重及保護之，為一切國家權力之義務。」第2條第2項第1句：「人人有生命與身體完整性之權利。」}\ggmarknotezh{art-1-1}\ggmarknotezh{art-2-2-1}}

% ===== 組合腳註：Art. 1(1) + Art. 2(2)（墮胎二案特有）=====
\newcommand{\ggfnArtOneOneTwoTwoDE}{\ggnotedeonce{art-1-1+2-2}{Art. 1 Abs. 1: ``Die Wuerde des Menschen ist unantastbar. Sie zu achten und zu schuetzen ist Verpflichtung aller staatlichen Gewalt.'' Art. 2 Abs. 2: ``Jeder hat das Recht auf Leben und koerperliche Unversehrtheit. Die Freiheit der Person ist unverletzlich. In diese Rechte darf nur auf Grund eines Gesetzes eingegriffen werden.''}\ggmarknotede{art-1-1}\ggmarknotede{art-2-2}\ggmarknotede{art-2-2-1}}
\newcommand{\ggfnArtOneOneTwoTwoZH}{\ggnotezhonce{art-1-1+2-2}{第1條第1項：「人之尊嚴不可侵犯。尊重及保護之，為一切國家權力之義務。」第2條第2項：「人人有生命與身體完整性之權利。人身自由不可侵犯。此等權利僅得依法律加以干預。」}\ggmarknotezh{art-1-1}\ggmarknotezh{art-2-2}\ggmarknotezh{art-2-2-1}}

% ===== 組合腳註：Art. 1(1) + Art. 2(1)（墮胎二案特有）=====
\newcommand{\ggfnArtOneOneTwoOneDE}{\ggnotedeonce{art-1-1+2-1}{Art. 1 Abs. 1: ``Die Wuerde des Menschen ist unantastbar. Sie zu achten und zu schuetzen ist Verpflichtung aller staatlichen Gewalt.'' Art. 2 Abs. 1: ``Jeder hat das Recht auf die freie Entfaltung seiner Persoenlichkeit, soweit er nicht die Rechte anderer verletzt und nicht gegen die verfassungsmaessige Ordnung oder das Sittengesetz verstoesst.''}\ggmarknotede{art-1-1}\ggmarknotede{art-2-1}}
\newcommand{\ggfnArtOneOneTwoOneZH}{\ggnotezhonce{art-1-1+2-1}{第1條第1項：「人之尊嚴不可侵犯。尊重及保護之，為一切國家權力之義務。」第2條第1項：「人人有自由發展其人格之權利，但不得侵害他人權利，亦不得違反合憲秩序或道德律。」}\ggmarknotezh{art-1-1}\ggmarknotezh{art-2-1}}

% ===== 組合腳註：Art. 2(2)(1) + Art. 1(1)（墮胎一案特有，含交叉標記）=====
\newcommand{\ggfnLifeDignityDE}{\ggnotedeonce{life-dignity}{Art. 2 Abs. 2 Satz 1: ``Jeder hat das Recht auf Leben und koerperliche Unversehrtheit.'' Art. 1 Abs. 1: ``Die Wuerde des Menschen ist unantastbar. Sie zu achten und zu schuetzen ist Verpflichtung aller staatlichen Gewalt.''}\ggmarknotede{art-2-2-1}\ggmarknotede{art-1-1}\ggmarknotede{art-1-1-2}}
\newcommand{\ggfnLifeDignityZH}{\ggnotezhonce{life-dignity}{第2條第2項第1句：「人人有生命與身體完整性之權利。」第1條第1項：「人之尊嚴不可侵犯。尊重及保護之，為一切國家權力之義務。」}\ggmarknotezh{art-2-2-1}\ggmarknotezh{art-1-1}\ggmarknotezh{art-1-1-2}}

% ===== 組合腳註：Art. 1(1) + Art. 2(2)（墮胎一案特有，條件判斷版）=====
\newcommand{\ggfnArtOneTwoTwoDE}{\ifcsname ggnoteseen@de@life-dignity\endcsname\ggfnArtTwoTwoRestDE{}\else\ggnotedeonce{art-1-2-2}{Art. 1 Abs. 1: ``Die Wuerde des Menschen ist unantastbar. Sie zu achten und zu schuetzen ist Verpflichtung aller staatlichen Gewalt.'' Art. 2 Abs. 2: ``Jeder hat das Recht auf Leben und koerperliche Unversehrtheit. Die Freiheit der Person ist unverletzlich. In diese Rechte darf nur auf Grund eines Gesetzes eingegriffen werden.''}\ggmarknotede{art-1-1}\ggmarknotede{art-1-1-2}\ggmarknotede{art-2-2}\ggmarknotede{art-2-2-1}\fi}
\newcommand{\ggfnArtOneTwoTwoZH}{\ifcsname ggnoteseen@zh@life-dignity\endcsname\ggfnArtTwoTwoRestZH{}\else\ggnotezhonce{art-1-2-2}{第1條第1項：「人之尊嚴不可侵犯。尊重及保護之，為一切國家權力之義務。」第2條第2項：「人人有生命與身體完整性之權利。人身自由不可侵犯。此等權利僅得依法律加以干預。」}\ggmarknotezh{art-1-1}\ggmarknotezh{art-1-1-2}\ggmarknotezh{art-2-2}\ggmarknotezh{art-2-2-1}\fi}

% ===== 組合腳註：Art. 1(1) + Art. 19(2)（墮胎一案特有，條件判斷版）=====
\newcommand{\ggfnArtOneNineteenDE}{\ifcsname ggnoteseen@de@art-1-1\endcsname\ggfnArtNineteenTwoDE{}\else\ggnotedeonce{art-1-19}{Art. 1 Abs. 1: ``Die Wuerde des Menschen ist unantastbar. Sie zu achten und zu schuetzen ist Verpflichtung aller staatlichen Gewalt.'' Art. 19 Abs. 2: ``In keinem Falle darf ein Grundrecht in seinem Wesensgehalt angetastet werden.''}\ggmarknotede{art-1-1}\ggmarknotede{art-1-1-2}\ggmarknotede{art-19-2}\fi}
\newcommand{\ggfnArtOneNineteenZH}{\ifcsname ggnoteseen@zh@art-1-1\endcsname\ggfnArtNineteenTwoZH{}\else\ggnotezhonce{art-1-19}{第1條第1項：「人之尊嚴不可侵犯。尊重及保護之，為一切國家權力之義務。」第19條第2項：「任何情形下，基本權利之本質內容均不得受到侵害。」}\ggmarknotezh{art-1-1}\ggmarknotezh{art-1-1-2}\ggmarknotezh{art-19-2}\fi}


% ===== GG 版本旗標（\ggversionde/zh 由各案在 % bverfge_gg.tex — 德國墮胎案 GG 條文腳註共用定義
% 使用方式：在各案 .tex 中先定義 \ggversionde/\ggversionzh，再 \input{../../共用LaTeX/bverfge_gg}

% ===== GG 版本旗標（\ggversionde/zh 由各案在 \input{../../共用LaTeX/bverfge_gg} 之前定義）=====
\newif\ifggversionnotede
\newif\ifggversionnotezh

% ===== 編者註腳基礎設施（首次標記模式）=====
\newif\ifeditornotelabelde
\newif\ifeditornotelabelzh
\newcommand{\editornotelabelde}{\ifeditornotelabelde\else\emph{Hrsg.-Anm. (nachfolgend ohne Kennzeichnung)}: \global\editornotelabeldetrue\fi}
\newcommand{\editornotelabelzh}{\ifeditornotelabelzh\else 編者註（下同）：\global\editornotelabelzhtrue\fi}
\newcommand{\editornotede}[1]{\ifshoweditornotes\footnote{\normalfont\footnotesize\editornotelabelde #1}\else\ifclassroommode\footnote{\normalfont\footnotesize\editornotelabelde #1}\fi\fi}
\newcommand{\editornotezh}[1]{\ifshoweditornotes\footnote{\normalfont\footnotesize\editornotelabelzh #1}\else\ifclassroommode\footnote{\normalfont\footnotesize\editornotelabelzh #1}\fi\fi}

% ===== GG 引用系統（\ggversionde/zh 由各案在 \input 前定義）=====
\newcommand{\ggmarknotede}[1]{\global\expandafter\let\csname ggnoteseen@de@#1\endcsname\empty}
\newcommand{\ggmarknotezh}[1]{\global\expandafter\let\csname ggnoteseen@zh@#1\endcsname\empty}
\newcommand{\ggnotedeonce}[2]{\ifcsname ggnoteseen@de@#1\endcsname\else\global\expandafter\let\csname ggnoteseen@de@#1\endcsname\empty\ggnotede{#2}\fi}
\newcommand{\ggnotezhonce}[2]{\ifcsname ggnoteseen@zh@#1\endcsname\else\global\expandafter\let\csname ggnoteseen@zh@#1\endcsname\empty\ggnotezh{#2}\fi}
\newcommand{\ggnotede}[1]{\editornotede{\ggversionde #1}}
\newcommand{\ggnotezh}[1]{\editornotezh{\ggversionzh #1}}

% ===== Art. 1 =====
\newcommand{\ggfnArtOneOneDE}{\ggnotedeonce{art-1-1}{Art. 1 Abs. 1: ``Die Wuerde des Menschen ist unantastbar. Sie zu achten und zu schuetzen ist Verpflichtung aller staatlichen Gewalt.''}}
\newcommand{\ggfnArtOneOneZH}{\ggnotezhonce{art-1-1}{第1條第1項：「人之尊嚴不可侵犯。尊重及保護之，為一切國家權力之義務。」}}
\newcommand{\ggfnArtOneOneTwoDE}{\ggnotedeonce{art-1-1-2}{Art. 1 Abs. 1 Satz 2: ``Sie zu achten und zu schuetzen ist Verpflichtung aller staatlichen Gewalt.''}}
\newcommand{\ggfnArtOneOneTwoZH}{\ggnotezhonce{art-1-1-2}{第1條第1項第2句：「尊重及保護之，為一切國家權力之義務。」}}

% ===== Art. 2 =====
\newcommand{\ggfnArtTwoOneDE}{\ggnotedeonce{art-2-1}{Art. 2 Abs. 1: ``Jeder hat das Recht auf die freie Entfaltung seiner Persoenlichkeit, soweit er nicht die Rechte anderer verletzt und nicht gegen die verfassungsmaessige Ordnung oder das Sittengesetz verstoesst.''}}
\newcommand{\ggfnArtTwoOneZH}{\ggnotezhonce{art-2-1}{第2條第1項：「人人有自由發展其人格之權利，但不得侵害他人權利，亦不得違反合憲秩序或道德律。」}}
% 簡易預設版（墮胎一案以 \renewcommand 覆寫為含條件邏輯之版本）
\newcommand{\ggfnArtTwoTwoDE}{\ggnotedeonce{art-2-2}{Art. 2 Abs. 2: ``Jeder hat das Recht auf Leben und koerperliche Unversehrtheit. Die Freiheit der Person ist unverletzlich. In diese Rechte darf nur auf Grund eines Gesetzes eingegriffen werden.''}\ggmarknotede{art-2-2-1}}
\newcommand{\ggfnArtTwoTwoZH}{\ggnotezhonce{art-2-2}{第2條第2項：「人人有生命與身體完整性之權利。人身自由不可侵犯。此等權利僅得依法律加以干預。」}\ggmarknotezh{art-2-2-1}}
\newcommand{\ggfnArtTwoTwoOneDE}{\ggnotedeonce{art-2-2-1}{Art. 2 Abs. 2 Satz 1: ``Jeder hat das Recht auf Leben und koerperliche Unversehrtheit.''}}
\newcommand{\ggfnArtTwoTwoOneZH}{\ggnotezhonce{art-2-2-1}{第2條第2項第1句：「人人有生命與身體完整性之權利。」}}
\newcommand{\ggfnArtTwoTwoThreeDE}{\ggnotedeonce{art-2-2-3}{Art. 2 Abs. 2 Satz 3: ``In diese Rechte darf nur auf Grund eines Gesetzes eingegriffen werden.''}}
\newcommand{\ggfnArtTwoTwoThreeZH}{\ggnotezhonce{art-2-2-3}{第2條第2項第3句：「此等權利僅得依法律加以干預。」}}
% 墮胎一案特有：第2條第2項第2–3句（Art. 2(2) 去掉第1句後的「餘文」）
\newcommand{\ggfnArtTwoTwoRestDE}{\ggnotedeonce{art-2-2-rest}{Art. 2 Abs. 2 Saetze 2-3: ``Die Freiheit der Person ist unverletzlich. In diese Rechte darf nur auf Grund eines Gesetzes eingegriffen werden.''}\ggmarknotede{art-2-2}}
\newcommand{\ggfnArtTwoTwoRestZH}{\ggnotezhonce{art-2-2-rest}{第2條第2項第2、3句：「人身自由不可侵犯。此等權利僅得依法律加以干預。」}\ggmarknotezh{art-2-2}}

% ===== Art. 3 =====
\newcommand{\ggfnArtThreeDE}{\ggnotedeonce{art-3}{Art. 3: ``Alle Menschen sind vor dem Gesetz gleich. Maenner und Frauen sind gleichberechtigt. Niemand darf wegen seines Geschlechtes, seiner Abstammung, seiner Rasse, seiner Sprache, seiner Heimat und Herkunft, seines Glaubens, seiner religioesen oder politischen Anschauungen benachteiligt oder bevorzugt werden.''}}
\newcommand{\ggfnArtThreeZH}{\ggnotezhonce{art-3}{第3條：「所有人在法律前一律平等。男女享有平等權利。任何人不得因其性別、血統、種族、語言、故鄉及出身、信仰、宗教或政治見解而受不利益或優惠。」}}
\newcommand{\ggfnArtThreeTwoDE}{\ggnotedeonce{art-3-2}{Art. 3 Abs. 2: ``Maenner und Frauen sind gleichberechtigt.''}}
\newcommand{\ggfnArtThreeTwoZH}{\ggnotezhonce{art-3-2}{第3條第2項：「男女享有平等權利。」}}

% ===== Art. 4 =====
\newcommand{\ggfnArtFourOneDE}{\ggnotedeonce{art-4-1}{Art. 4 Abs. 1: ``Die Freiheit des Glaubens, des Gewissens und die Freiheit des religioesen und weltanschaulichen Bekenntnisses sind unverletzlich.''}}
\newcommand{\ggfnArtFourOneZH}{\ggnotezhonce{art-4-1}{第4條第1項：「信仰自由、良心自由，以及宗教與世界觀信念自由，不可侵犯。」}}

% ===== Art. 5 =====
\newcommand{\ggfnArtFiveOneDE}{\ggnotedeonce{art-5-1}{Art. 5 Abs. 1: ``Jeder hat das Recht, seine Meinung in Wort, Schrift und Bild frei zu aeussern und zu verbreiten und sich aus allgemein zugaenglichen Quellen ungehindert zu unterrichten. Die Pressefreiheit und die Freiheit der Berichterstattung durch Rundfunk und Film werden gewaehrleistet. Eine Zensur findet nicht statt.''}}
\newcommand{\ggfnArtFiveOneZH}{\ggnotezhonce{art-5-1}{第5條第1項：「人人有以言語、文字及圖像自由表達及傳播其意見，並由一般可接近來源不受阻礙獲取資訊之權利。新聞自由及廣播、電影報導自由受保障。不設審查制度。」}}

% ===== Art. 6 =====
\newcommand{\ggfnArtSixDE}{\ggnotedeonce{art-6}{Art. 6 Abs. 1: ``Ehe und Familie stehen unter dem besonderen Schutze der staatlichen Ordnung.'' Abs. 2 Satz 1: ``Pflege und Erziehung der Kinder sind das natuerliche Recht der Eltern und die zuvoerderst ihnen obliegende Pflicht.'' Abs. 4: ``Jede Mutter hat Anspruch auf den Schutz und die Fuersorge der Gemeinschaft.''}}
\newcommand{\ggfnArtSixZH}{\ggnotezhonce{art-6}{第6條第1項：「婚姻與家庭受國家秩序之特別保護。」第2項第1句：「照護及教育子女，為父母之自然權利，並為首先由其承擔之義務。」第4項：「每一母親均有請求共同體保護與照顧之權利。」}}
\newcommand{\ggfnArtSixOneDE}{\ggnotedeonce{art-6-1}{Art. 6 Abs. 1: ``Ehe und Familie stehen unter dem besonderen Schutze der staatlichen Ordnung.''}}
\newcommand{\ggfnArtSixOneZH}{\ggnotezhonce{art-6-1}{第6條第1項：「婚姻與家庭受國家秩序之特別保護。」}}
\newcommand{\ggfnArtSixTwoDE}{\ggnotedeonce{art-6-2}{Art. 6 Abs. 2: ``Pflege und Erziehung der Kinder sind das natuerliche Recht der Eltern und die zuvoerderst ihnen obliegende Pflicht. Ueber ihre Betaetigung wacht die staatliche Gemeinschaft.''}}
\newcommand{\ggfnArtSixTwoZH}{\ggnotezhonce{art-6-2}{第6條第2項：「照護及教育子女，為父母之自然權利，並為首先由其承擔之義務。國家共同體監督其行使。」}}
\newcommand{\ggfnArtSixTwoOneDE}{\ggnotedeonce{art-6-2-1}{Art. 6 Abs. 2 Satz 1: ``Pflege und Erziehung der Kinder sind das natuerliche Recht der Eltern und die zuvoerderst ihnen obliegende Pflicht.''}}
\newcommand{\ggfnArtSixTwoOneZH}{\ggnotezhonce{art-6-2-1}{第6條第2項第1句：「照護及教育子女，為父母之自然權利，並為首先由其承擔之義務。」}}
\newcommand{\ggfnArtSixFourDE}{\ggnotedeonce{art-6-4}{Art. 6 Abs. 4: ``Jede Mutter hat Anspruch auf den Schutz und die Fuersorge der Gemeinschaft.''}}
\newcommand{\ggfnArtSixFourZH}{\ggnotezhonce{art-6-4}{第6條第4項：「每一母親均有請求共同體保護與照顧之權利。」}}

% ===== Art. 12 =====
\newcommand{\ggfnArtTwelveOneDE}{\ggnotedeonce{art-12-1}{Art. 12 Abs. 1: ``Alle Deutschen haben das Recht, Beruf, Arbeitsplatz und Ausbildungsstaette frei zu waehlen. Die Berufsausuebung kann durch Gesetz oder auf Grund eines Gesetzes geregelt werden.''}}
\newcommand{\ggfnArtTwelveOneZH}{\ggnotezhonce{art-12-1}{第12條第1項：「所有德國人均有自由選擇職業、工作場所及訓練場所之權利。職業行使得以法律或基於法律加以規範。」}}

% ===== Art. 19 =====
\newcommand{\ggfnArtNineteenTwoDE}{\ggnotedeonce{art-19-2}{Art. 19 Abs. 2: ``In keinem Falle darf ein Grundrecht in seinem Wesensgehalt angetastet werden.''}}
\newcommand{\ggfnArtNineteenTwoZH}{\ggnotezhonce{art-19-2}{第19條第2項：「任何情形下，基本權利之本質內容均不得受到侵害。」}}

% ===== Art. 20 =====
\newcommand{\ggfnArtTwentyOneDE}{\ggnotedeonce{art-20-1}{Art. 20 Abs. 1: ``Die Bundesrepublik Deutschland ist ein demokratischer und sozialer Bundesstaat.''}}
\newcommand{\ggfnArtTwentyOneZH}{\ggnotezhonce{art-20-1}{第20條第1項：「德意志聯邦共和國為民主及社會之聯邦國。」}}
\newcommand{\ggfnArtTwentyThreeDE}{\ggnotedeonce{art-20-3}{Art. 20 Abs. 3: ``Die Gesetzgebung ist an die verfassungsmaessige Ordnung, die vollziehende Gewalt und die Rechtsprechung sind an Gesetz und Recht gebunden.''}}
\newcommand{\ggfnArtTwentyThreeZH}{\ggnotezhonce{art-20-3}{第20條第3項：「立法受合憲秩序拘束；行政權與司法權受法律與法拘束。」}}

% ===== Art. 26 + Art. 143（墮胎一案特有）=====
\newcommand{\ggfnArtTwoSixAndOneFourThreeDE}{\ggnotedeonce{art-26-143}{Art. 26 Abs. 1 Satz 2: ``Sie sind unter Strafe zu stellen.'' Art. 143 in der urspruenglichen Fassung (24. Mai 1949 bis 1. September 1951) enthielt Strafvorschriften zum Hochverrat; 1975 war Art. 143 bereits weggefallen.}}
\newcommand{\ggfnArtTwoSixAndOneFourThreeZH}{\ggnotezhonce{art-26-143}{第26條第1項第2句：「此等行為應規定為犯罪處罰。」第143條原始版本（1949年5月24日至1951年9月1日）含有關於叛國罪之刑罰規定；至1975年時，第143條已廢止。}}

% ===== Art. 28 =====
\newcommand{\ggfnArtTwentyEightOneDE}{\ggnotedeonce{art-28-1-1}{Art. 28 Abs. 1 Satz 1: ``Die verfassungsmaessige Ordnung in den Laendern muss den Grundsaetzen des republikanischen, demokratischen und sozialen Rechtsstaates im Sinne dieses Grundgesetzes entsprechen.''}}
\newcommand{\ggfnArtTwentyEightOneZH}{\ggnotezhonce{art-28-1-1}{第28條第1項第1句：「各邦之合憲秩序，須符合本基本法意義下共和、民主及社會法治國之原則。」}}

% ===== Art. 30 =====
\newcommand{\ggfnArtThirtyDE}{\ggnotedeonce{art-30}{Art. 30: ``Die Ausuebung der staatlichen Befugnisse und die Erfuellung der staatlichen Aufgaben ist Sache der Laender, soweit dieses Grundgesetz keine andere Regelung trifft oder zulaesst.''}}
\newcommand{\ggfnArtThirtyZH}{\ggnotezhonce{art-30}{第30條：「國家權限之行使與國家任務之履行，屬於各邦，但本基本法另有規定或容許者，不在此限。」}}

% ===== Art. 74 =====
\newcommand{\ggfnArtSeventyFourOneDE}{\ggnotedeonce{art-74-1}{Art. 74 Nr. 1: Die konkurrierende Gesetzgebung erstreckt sich auf ``das buergerliche Recht, das Strafrecht und den Strafvollzug, die Gerichtsverfassung, das gerichtliche Verfahren, die Rechtsanwaltschaft, das Notariat und die Rechtsberatung''.}}
\newcommand{\ggfnArtSeventyFourOneZH}{\ggnotezhonce{art-74-1}{第74條第1款：競合立法及於「民法、刑法及刑罰執行、法院組織、訴訟程序、律師制度、公證制度及法律諮詢」。}}
\newcommand{\ggfnArtSeventyFourSevenDE}{\ggnotedeonce{art-74-7}{Art. 74 Nr. 7: Die konkurrierende Gesetzgebung erstreckt sich auf ``die oeffentliche Fuersorge''.}}
\newcommand{\ggfnArtSeventyFourSevenZH}{\ggnotezhonce{art-74-7}{第74條第7款：競合立法及於「公共扶助」。}}
\newcommand{\ggfnArtSeventyFourTwelveDE}{\ggnotedeonce{art-74-12}{Art. 74 Nr. 12: Die konkurrierende Gesetzgebung erstreckt sich auf ``das Arbeitsrecht einschliesslich der Betriebsverfassung, des Arbeitsschutzes und der Arbeitsvermittlung sowie die Sozialversicherung einschliesslich der Arbeitslosenversicherung''.}}
\newcommand{\ggfnArtSeventyFourTwelveZH}{\ggnotezhonce{art-74-12}{第74條第12款：競合立法及於「勞動法，包括企業組織、勞動保護及就業媒合，以及社會保險，包括失業保險」。}}

% ===== Art. 77（墮胎一案特有）=====
\newcommand{\ggfnArtSevenSevenThreeDE}{\ggnotedeonce{art-77-3}{Art. 77 Abs. 3 Satz 1: ``Soweit zu einem Gesetze die Zustimmung des Bundesrates nicht erforderlich ist, kann der Bundesrat, wenn das Verfahren nach Absatz 2 beendigt ist, gegen ein vom Bundestage beschlossenes Gesetz binnen zwei Wochen Einspruch einlegen.''}}
\newcommand{\ggfnArtSevenSevenThreeZH}{\ggnotezhonce{art-77-3}{第77條第3項第1句：「法律無須聯邦參議院同意者，於第2項程序終了後，聯邦參議院得於二週內對聯邦議院通過之法律提出異議。」}}

% ===== Art. 80 =====
\newcommand{\ggfnArtEightyOneOneDE}{\ggnotedeonce{art-80-1-1}{Art. 80 Abs. 1 Satz 1: ``Durch Gesetz koennen die Bundesregierung, ein Bundesminister oder die Landesregierungen ermaechtigt werden, Rechtsverordnungen zu erlassen.''}}
\newcommand{\ggfnArtEightyOneOneZH}{\ggnotezhonce{art-80-1-1}{第80條第1項第1句：「得以法律授權聯邦政府、聯邦部長或邦政府發布法規命令。」}}

% ===== Art. 83 =====
\newcommand{\ggfnArtEightyThreeDE}{\ggnotedeonce{art-83}{Art. 83: ``Die Laender fuehren die Bundesgesetze als eigene Angelegenheit aus, soweit dieses Grundgesetz nichts anderes bestimmt oder zulaesst.''}}
\newcommand{\ggfnArtEightyThreeZH}{\ggnotezhonce{art-83}{第83條：「各邦以自己事務執行聯邦法律，但本基本法另有規定或容許者，不在此限。」}}

% ===== Art. 84 =====
\newcommand{\ggfnArtEightyFourOneDE}{\ggnotedeonce{art-84-1}{Art. 84 Abs. 1: ``Fuehren die Laender die Bundesgesetze als eigene Angelegenheit aus, so regeln sie die Einrichtung der Behoerden und das Verwaltungsverfahren, soweit nicht Bundesgesetze mit Zustimmung des Bundesrates etwas anderes bestimmen.''}}
\newcommand{\ggfnArtEightyFourOneZH}{\ggnotezhonce{art-84-1}{第84條第1項：「各邦以自己事務執行聯邦法律時，除聯邦法律經聯邦參議院同意另有規定外，由各邦規範機關之設置及行政程序。」}}
% 別名（墮胎一案使用之縮寫命名）
\let\ggfnArtEightFourOneDE\ggfnArtEightyFourOneDE
\let\ggfnArtEightFourOneZH\ggfnArtEightyFourOneZH

% ===== Art. 93 =====
\newcommand{\ggfnArtNinetyThreeOneTwoDE}{\ggnotedeonce{art-93-1-2}{Art. 93 Abs. 1 Nr. 2: Das Bundesverfassungsgericht entscheidet ``bei Meinungsverschiedenheiten oder Zweifeln ueber die foermliche und sachliche Vereinbarkeit von Bundesrecht oder Landesrecht mit diesem Grundgesetze ... auf Antrag der Bundesregierung, einer Landesregierung oder eines Drittels der Mitglieder des Bundestages''.}}
\newcommand{\ggfnArtNinetyThreeOneTwoZH}{\ggnotezhonce{art-93-1-2}{第93條第1項第2款：聯邦憲法法院就「聯邦法或邦法在形式及實質上是否與本基本法相容之意見分歧或疑義……依聯邦政府、邦政府或聯邦議院三分之一議員之聲請」作成裁判。}}
% 別名（墮胎一案使用之縮寫命名）
\let\ggfnArtNineThreeOneTwoDE\ggfnArtNinetyThreeOneTwoDE
\let\ggfnArtNineThreeOneTwoZH\ggfnArtNinetyThreeOneTwoZH

% ===== Art. 102 =====
\newcommand{\ggfnArtOneZeroTwoDE}{\ggnotedeonce{art-102}{Art. 102: ``Die Todesstrafe ist abgeschafft.''}}
\newcommand{\ggfnArtOneZeroTwoZH}{\ggnotezhonce{art-102}{第102條：「死刑廢止。」}}

% ===== Art. 103 =====
\newcommand{\ggfnArtOneZeroThreeTwoDE}{\ggnotedeonce{art-103-2}{Art. 103 Abs. 2: ``Eine Tat kann nur bestraft werden, wenn die Strafbarkeit gesetzlich bestimmt war, bevor die Tat begangen wurde.''}}
\newcommand{\ggfnArtOneZeroThreeTwoZH}{\ggnotezhonce{art-103-2}{第103條第2項：「行為之可罰性，須於行為前以法律定之，始得處罰。」}}

% ===== Art. 104 =====
\newcommand{\ggfnArtOneZeroFourOneDE}{\ggnotedeonce{art-104-1}{Art. 104 Abs. 1: ``Die Freiheit der Person kann nur auf Grund eines foermlichen Gesetzes und nur unter Beachtung der darin vorgeschriebenen Formen beschraenkt werden.''}}
\newcommand{\ggfnArtOneZeroFourOneZH}{\ggnotezhonce{art-104-1}{第104條第1項：「人身自由僅得基於形式法律，並遵守其中規定之形式加以限制。」}}

% ===== 組合腳註：Art. 3 + Art. 6（墮胎一案特有）=====
\newcommand{\ggfnArtThreeSixDE}{\ggnotedeonce{art-3-6}{Art. 3: ``Alle Menschen sind vor dem Gesetz gleich. Maenner und Frauen sind gleichberechtigt. Niemand darf wegen seines Geschlechtes, seiner Abstammung, seiner Rasse, seiner Sprache, seiner Heimat und Herkunft, seines Glaubens, seiner religioesen oder politischen Anschauungen benachteiligt oder bevorzugt werden.'' Art. 6 Abs. 1: ``Ehe und Familie stehen unter dem besonderen Schutze der staatlichen Ordnung.'' Abs. 2 Satz 1: ``Pflege und Erziehung der Kinder sind das natuerliche Recht der Eltern und die zuvoerderst ihnen obliegende Pflicht.'' Abs. 4: ``Jede Mutter hat Anspruch auf den Schutz und die Fuersorge der Gemeinschaft.''}\ggmarknotede{art-3}\ggmarknotede{art-6}\ggmarknotede{art-6-1-4}}
\newcommand{\ggfnArtThreeSixZH}{\ggnotezhonce{art-3-6}{第3條：「所有人在法律前一律平等。男女享有平等權利。任何人不得因其性別、血統、種族、語言、故鄉及出身、信仰、宗教或政治見解而受不利益或優惠。」第6條第1項：「婚姻與家庭受國家秩序之特別保護。」第2項第1句：「照護及教育子女，為父母之自然權利，並為首先由其承擔之義務。」第4項：「每一母親均有請求共同體保護與照顧之權利。」}\ggmarknotezh{art-3}\ggmarknotezh{art-6}\ggmarknotezh{art-6-1-4}}

% ===== 組合腳註：Art. 6(1) + Art. 6(4)（墮胎一案特有）=====
\newcommand{\ggfnArtSixOneFourDE}{\ggnotedeonce{art-6-1-4}{Art. 6 Abs. 1: ``Ehe und Familie stehen unter dem besonderen Schutze der staatlichen Ordnung.'' Art. 6 Abs. 4: ``Jede Mutter hat Anspruch auf den Schutz und die Fuersorge der Gemeinschaft.''}\ggmarknotede{art-6}}
\newcommand{\ggfnArtSixOneFourZH}{\ggnotezhonce{art-6-1-4}{第6條第1項：「婚姻與家庭受國家秩序之特別保護。」第4項：「每一母親均有請求共同體保護與照顧之權利。」}\ggmarknotezh{art-6}}

% ===== 組合腳註：Art. 1(1) + Art. 2(2)(1)（墮胎二案特有）=====
\newcommand{\ggfnArtOneOneTwoTwoOneDE}{\ggnotedeonce{art-1-1+2-2-1}{Art. 1 Abs. 1: ``Die Wuerde des Menschen ist unantastbar. Sie zu achten und zu schuetzen ist Verpflichtung aller staatlichen Gewalt.'' Art. 2 Abs. 2 Satz 1: ``Jeder hat das Recht auf Leben und koerperliche Unversehrtheit.''}\ggmarknotede{art-1-1}\ggmarknotede{art-2-2-1}}
\newcommand{\ggfnArtOneOneTwoTwoOneZH}{\ggnotezhonce{art-1-1+2-2-1}{第1條第1項：「人之尊嚴不可侵犯。尊重及保護之，為一切國家權力之義務。」第2條第2項第1句：「人人有生命與身體完整性之權利。」}\ggmarknotezh{art-1-1}\ggmarknotezh{art-2-2-1}}

% ===== 組合腳註：Art. 1(1) + Art. 2(2)（墮胎二案特有）=====
\newcommand{\ggfnArtOneOneTwoTwoDE}{\ggnotedeonce{art-1-1+2-2}{Art. 1 Abs. 1: ``Die Wuerde des Menschen ist unantastbar. Sie zu achten und zu schuetzen ist Verpflichtung aller staatlichen Gewalt.'' Art. 2 Abs. 2: ``Jeder hat das Recht auf Leben und koerperliche Unversehrtheit. Die Freiheit der Person ist unverletzlich. In diese Rechte darf nur auf Grund eines Gesetzes eingegriffen werden.''}\ggmarknotede{art-1-1}\ggmarknotede{art-2-2}\ggmarknotede{art-2-2-1}}
\newcommand{\ggfnArtOneOneTwoTwoZH}{\ggnotezhonce{art-1-1+2-2}{第1條第1項：「人之尊嚴不可侵犯。尊重及保護之，為一切國家權力之義務。」第2條第2項：「人人有生命與身體完整性之權利。人身自由不可侵犯。此等權利僅得依法律加以干預。」}\ggmarknotezh{art-1-1}\ggmarknotezh{art-2-2}\ggmarknotezh{art-2-2-1}}

% ===== 組合腳註：Art. 1(1) + Art. 2(1)（墮胎二案特有）=====
\newcommand{\ggfnArtOneOneTwoOneDE}{\ggnotedeonce{art-1-1+2-1}{Art. 1 Abs. 1: ``Die Wuerde des Menschen ist unantastbar. Sie zu achten und zu schuetzen ist Verpflichtung aller staatlichen Gewalt.'' Art. 2 Abs. 1: ``Jeder hat das Recht auf die freie Entfaltung seiner Persoenlichkeit, soweit er nicht die Rechte anderer verletzt und nicht gegen die verfassungsmaessige Ordnung oder das Sittengesetz verstoesst.''}\ggmarknotede{art-1-1}\ggmarknotede{art-2-1}}
\newcommand{\ggfnArtOneOneTwoOneZH}{\ggnotezhonce{art-1-1+2-1}{第1條第1項：「人之尊嚴不可侵犯。尊重及保護之，為一切國家權力之義務。」第2條第1項：「人人有自由發展其人格之權利，但不得侵害他人權利，亦不得違反合憲秩序或道德律。」}\ggmarknotezh{art-1-1}\ggmarknotezh{art-2-1}}

% ===== 組合腳註：Art. 2(2)(1) + Art. 1(1)（墮胎一案特有，含交叉標記）=====
\newcommand{\ggfnLifeDignityDE}{\ggnotedeonce{life-dignity}{Art. 2 Abs. 2 Satz 1: ``Jeder hat das Recht auf Leben und koerperliche Unversehrtheit.'' Art. 1 Abs. 1: ``Die Wuerde des Menschen ist unantastbar. Sie zu achten und zu schuetzen ist Verpflichtung aller staatlichen Gewalt.''}\ggmarknotede{art-2-2-1}\ggmarknotede{art-1-1}\ggmarknotede{art-1-1-2}}
\newcommand{\ggfnLifeDignityZH}{\ggnotezhonce{life-dignity}{第2條第2項第1句：「人人有生命與身體完整性之權利。」第1條第1項：「人之尊嚴不可侵犯。尊重及保護之，為一切國家權力之義務。」}\ggmarknotezh{art-2-2-1}\ggmarknotezh{art-1-1}\ggmarknotezh{art-1-1-2}}

% ===== 組合腳註：Art. 1(1) + Art. 2(2)（墮胎一案特有，條件判斷版）=====
\newcommand{\ggfnArtOneTwoTwoDE}{\ifcsname ggnoteseen@de@life-dignity\endcsname\ggfnArtTwoTwoRestDE{}\else\ggnotedeonce{art-1-2-2}{Art. 1 Abs. 1: ``Die Wuerde des Menschen ist unantastbar. Sie zu achten und zu schuetzen ist Verpflichtung aller staatlichen Gewalt.'' Art. 2 Abs. 2: ``Jeder hat das Recht auf Leben und koerperliche Unversehrtheit. Die Freiheit der Person ist unverletzlich. In diese Rechte darf nur auf Grund eines Gesetzes eingegriffen werden.''}\ggmarknotede{art-1-1}\ggmarknotede{art-1-1-2}\ggmarknotede{art-2-2}\ggmarknotede{art-2-2-1}\fi}
\newcommand{\ggfnArtOneTwoTwoZH}{\ifcsname ggnoteseen@zh@life-dignity\endcsname\ggfnArtTwoTwoRestZH{}\else\ggnotezhonce{art-1-2-2}{第1條第1項：「人之尊嚴不可侵犯。尊重及保護之，為一切國家權力之義務。」第2條第2項：「人人有生命與身體完整性之權利。人身自由不可侵犯。此等權利僅得依法律加以干預。」}\ggmarknotezh{art-1-1}\ggmarknotezh{art-1-1-2}\ggmarknotezh{art-2-2}\ggmarknotezh{art-2-2-1}\fi}

% ===== 組合腳註：Art. 1(1) + Art. 19(2)（墮胎一案特有，條件判斷版）=====
\newcommand{\ggfnArtOneNineteenDE}{\ifcsname ggnoteseen@de@art-1-1\endcsname\ggfnArtNineteenTwoDE{}\else\ggnotedeonce{art-1-19}{Art. 1 Abs. 1: ``Die Wuerde des Menschen ist unantastbar. Sie zu achten und zu schuetzen ist Verpflichtung aller staatlichen Gewalt.'' Art. 19 Abs. 2: ``In keinem Falle darf ein Grundrecht in seinem Wesensgehalt angetastet werden.''}\ggmarknotede{art-1-1}\ggmarknotede{art-1-1-2}\ggmarknotede{art-19-2}\fi}
\newcommand{\ggfnArtOneNineteenZH}{\ifcsname ggnoteseen@zh@art-1-1\endcsname\ggfnArtNineteenTwoZH{}\else\ggnotezhonce{art-1-19}{第1條第1項：「人之尊嚴不可侵犯。尊重及保護之，為一切國家權力之義務。」第19條第2項：「任何情形下，基本權利之本質內容均不得受到侵害。」}\ggmarknotezh{art-1-1}\ggmarknotezh{art-1-1-2}\ggmarknotezh{art-19-2}\fi}
 之前定義）=====
\newif\ifggversionnotede
\newif\ifggversionnotezh

% ===== 編者註腳基礎設施（首次標記模式）=====
\newif\ifeditornotelabelde
\newif\ifeditornotelabelzh
\newcommand{\editornotelabelde}{\ifeditornotelabelde\else\emph{Hrsg.-Anm. (nachfolgend ohne Kennzeichnung)}: \global\editornotelabeldetrue\fi}
\newcommand{\editornotelabelzh}{\ifeditornotelabelzh\else 編者註（下同）：\global\editornotelabelzhtrue\fi}
\newcommand{\editornotede}[1]{\ifshoweditornotes\footnote{\normalfont\footnotesize\editornotelabelde #1}\else\ifclassroommode\footnote{\normalfont\footnotesize\editornotelabelde #1}\fi\fi}
\newcommand{\editornotezh}[1]{\ifshoweditornotes\footnote{\normalfont\footnotesize\editornotelabelzh #1}\else\ifclassroommode\footnote{\normalfont\footnotesize\editornotelabelzh #1}\fi\fi}

% ===== GG 引用系統（\ggversionde/zh 由各案在 \input 前定義）=====
\newcommand{\ggmarknotede}[1]{\global\expandafter\let\csname ggnoteseen@de@#1\endcsname\empty}
\newcommand{\ggmarknotezh}[1]{\global\expandafter\let\csname ggnoteseen@zh@#1\endcsname\empty}
\newcommand{\ggnotedeonce}[2]{\ifcsname ggnoteseen@de@#1\endcsname\else\global\expandafter\let\csname ggnoteseen@de@#1\endcsname\empty\ggnotede{#2}\fi}
\newcommand{\ggnotezhonce}[2]{\ifcsname ggnoteseen@zh@#1\endcsname\else\global\expandafter\let\csname ggnoteseen@zh@#1\endcsname\empty\ggnotezh{#2}\fi}
\newcommand{\ggnotede}[1]{\editornotede{\ggversionde #1}}
\newcommand{\ggnotezh}[1]{\editornotezh{\ggversionzh #1}}

% ===== Art. 1 =====
\newcommand{\ggfnArtOneOneDE}{\ggnotedeonce{art-1-1}{Art. 1 Abs. 1: ``Die Wuerde des Menschen ist unantastbar. Sie zu achten und zu schuetzen ist Verpflichtung aller staatlichen Gewalt.''}}
\newcommand{\ggfnArtOneOneZH}{\ggnotezhonce{art-1-1}{第1條第1項：「人之尊嚴不可侵犯。尊重及保護之，為一切國家權力之義務。」}}
\newcommand{\ggfnArtOneOneTwoDE}{\ggnotedeonce{art-1-1-2}{Art. 1 Abs. 1 Satz 2: ``Sie zu achten und zu schuetzen ist Verpflichtung aller staatlichen Gewalt.''}}
\newcommand{\ggfnArtOneOneTwoZH}{\ggnotezhonce{art-1-1-2}{第1條第1項第2句：「尊重及保護之，為一切國家權力之義務。」}}

% ===== Art. 2 =====
\newcommand{\ggfnArtTwoOneDE}{\ggnotedeonce{art-2-1}{Art. 2 Abs. 1: ``Jeder hat das Recht auf die freie Entfaltung seiner Persoenlichkeit, soweit er nicht die Rechte anderer verletzt und nicht gegen die verfassungsmaessige Ordnung oder das Sittengesetz verstoesst.''}}
\newcommand{\ggfnArtTwoOneZH}{\ggnotezhonce{art-2-1}{第2條第1項：「人人有自由發展其人格之權利，但不得侵害他人權利，亦不得違反合憲秩序或道德律。」}}
% 簡易預設版（墮胎一案以 \renewcommand 覆寫為含條件邏輯之版本）
\newcommand{\ggfnArtTwoTwoDE}{\ggnotedeonce{art-2-2}{Art. 2 Abs. 2: ``Jeder hat das Recht auf Leben und koerperliche Unversehrtheit. Die Freiheit der Person ist unverletzlich. In diese Rechte darf nur auf Grund eines Gesetzes eingegriffen werden.''}\ggmarknotede{art-2-2-1}}
\newcommand{\ggfnArtTwoTwoZH}{\ggnotezhonce{art-2-2}{第2條第2項：「人人有生命與身體完整性之權利。人身自由不可侵犯。此等權利僅得依法律加以干預。」}\ggmarknotezh{art-2-2-1}}
\newcommand{\ggfnArtTwoTwoOneDE}{\ggnotedeonce{art-2-2-1}{Art. 2 Abs. 2 Satz 1: ``Jeder hat das Recht auf Leben und koerperliche Unversehrtheit.''}}
\newcommand{\ggfnArtTwoTwoOneZH}{\ggnotezhonce{art-2-2-1}{第2條第2項第1句：「人人有生命與身體完整性之權利。」}}
\newcommand{\ggfnArtTwoTwoThreeDE}{\ggnotedeonce{art-2-2-3}{Art. 2 Abs. 2 Satz 3: ``In diese Rechte darf nur auf Grund eines Gesetzes eingegriffen werden.''}}
\newcommand{\ggfnArtTwoTwoThreeZH}{\ggnotezhonce{art-2-2-3}{第2條第2項第3句：「此等權利僅得依法律加以干預。」}}
% 墮胎一案特有：第2條第2項第2–3句（Art. 2(2) 去掉第1句後的「餘文」）
\newcommand{\ggfnArtTwoTwoRestDE}{\ggnotedeonce{art-2-2-rest}{Art. 2 Abs. 2 Saetze 2-3: ``Die Freiheit der Person ist unverletzlich. In diese Rechte darf nur auf Grund eines Gesetzes eingegriffen werden.''}\ggmarknotede{art-2-2}}
\newcommand{\ggfnArtTwoTwoRestZH}{\ggnotezhonce{art-2-2-rest}{第2條第2項第2、3句：「人身自由不可侵犯。此等權利僅得依法律加以干預。」}\ggmarknotezh{art-2-2}}

% ===== Art. 3 =====
\newcommand{\ggfnArtThreeDE}{\ggnotedeonce{art-3}{Art. 3: ``Alle Menschen sind vor dem Gesetz gleich. Maenner und Frauen sind gleichberechtigt. Niemand darf wegen seines Geschlechtes, seiner Abstammung, seiner Rasse, seiner Sprache, seiner Heimat und Herkunft, seines Glaubens, seiner religioesen oder politischen Anschauungen benachteiligt oder bevorzugt werden.''}}
\newcommand{\ggfnArtThreeZH}{\ggnotezhonce{art-3}{第3條：「所有人在法律前一律平等。男女享有平等權利。任何人不得因其性別、血統、種族、語言、故鄉及出身、信仰、宗教或政治見解而受不利益或優惠。」}}
\newcommand{\ggfnArtThreeTwoDE}{\ggnotedeonce{art-3-2}{Art. 3 Abs. 2: ``Maenner und Frauen sind gleichberechtigt.''}}
\newcommand{\ggfnArtThreeTwoZH}{\ggnotezhonce{art-3-2}{第3條第2項：「男女享有平等權利。」}}

% ===== Art. 4 =====
\newcommand{\ggfnArtFourOneDE}{\ggnotedeonce{art-4-1}{Art. 4 Abs. 1: ``Die Freiheit des Glaubens, des Gewissens und die Freiheit des religioesen und weltanschaulichen Bekenntnisses sind unverletzlich.''}}
\newcommand{\ggfnArtFourOneZH}{\ggnotezhonce{art-4-1}{第4條第1項：「信仰自由、良心自由，以及宗教與世界觀信念自由，不可侵犯。」}}

% ===== Art. 5 =====
\newcommand{\ggfnArtFiveOneDE}{\ggnotedeonce{art-5-1}{Art. 5 Abs. 1: ``Jeder hat das Recht, seine Meinung in Wort, Schrift und Bild frei zu aeussern und zu verbreiten und sich aus allgemein zugaenglichen Quellen ungehindert zu unterrichten. Die Pressefreiheit und die Freiheit der Berichterstattung durch Rundfunk und Film werden gewaehrleistet. Eine Zensur findet nicht statt.''}}
\newcommand{\ggfnArtFiveOneZH}{\ggnotezhonce{art-5-1}{第5條第1項：「人人有以言語、文字及圖像自由表達及傳播其意見，並由一般可接近來源不受阻礙獲取資訊之權利。新聞自由及廣播、電影報導自由受保障。不設審查制度。」}}

% ===== Art. 6 =====
\newcommand{\ggfnArtSixDE}{\ggnotedeonce{art-6}{Art. 6 Abs. 1: ``Ehe und Familie stehen unter dem besonderen Schutze der staatlichen Ordnung.'' Abs. 2 Satz 1: ``Pflege und Erziehung der Kinder sind das natuerliche Recht der Eltern und die zuvoerderst ihnen obliegende Pflicht.'' Abs. 4: ``Jede Mutter hat Anspruch auf den Schutz und die Fuersorge der Gemeinschaft.''}}
\newcommand{\ggfnArtSixZH}{\ggnotezhonce{art-6}{第6條第1項：「婚姻與家庭受國家秩序之特別保護。」第2項第1句：「照護及教育子女，為父母之自然權利，並為首先由其承擔之義務。」第4項：「每一母親均有請求共同體保護與照顧之權利。」}}
\newcommand{\ggfnArtSixOneDE}{\ggnotedeonce{art-6-1}{Art. 6 Abs. 1: ``Ehe und Familie stehen unter dem besonderen Schutze der staatlichen Ordnung.''}}
\newcommand{\ggfnArtSixOneZH}{\ggnotezhonce{art-6-1}{第6條第1項：「婚姻與家庭受國家秩序之特別保護。」}}
\newcommand{\ggfnArtSixTwoDE}{\ggnotedeonce{art-6-2}{Art. 6 Abs. 2: ``Pflege und Erziehung der Kinder sind das natuerliche Recht der Eltern und die zuvoerderst ihnen obliegende Pflicht. Ueber ihre Betaetigung wacht die staatliche Gemeinschaft.''}}
\newcommand{\ggfnArtSixTwoZH}{\ggnotezhonce{art-6-2}{第6條第2項：「照護及教育子女，為父母之自然權利，並為首先由其承擔之義務。國家共同體監督其行使。」}}
\newcommand{\ggfnArtSixTwoOneDE}{\ggnotedeonce{art-6-2-1}{Art. 6 Abs. 2 Satz 1: ``Pflege und Erziehung der Kinder sind das natuerliche Recht der Eltern und die zuvoerderst ihnen obliegende Pflicht.''}}
\newcommand{\ggfnArtSixTwoOneZH}{\ggnotezhonce{art-6-2-1}{第6條第2項第1句：「照護及教育子女，為父母之自然權利，並為首先由其承擔之義務。」}}
\newcommand{\ggfnArtSixFourDE}{\ggnotedeonce{art-6-4}{Art. 6 Abs. 4: ``Jede Mutter hat Anspruch auf den Schutz und die Fuersorge der Gemeinschaft.''}}
\newcommand{\ggfnArtSixFourZH}{\ggnotezhonce{art-6-4}{第6條第4項：「每一母親均有請求共同體保護與照顧之權利。」}}

% ===== Art. 12 =====
\newcommand{\ggfnArtTwelveOneDE}{\ggnotedeonce{art-12-1}{Art. 12 Abs. 1: ``Alle Deutschen haben das Recht, Beruf, Arbeitsplatz und Ausbildungsstaette frei zu waehlen. Die Berufsausuebung kann durch Gesetz oder auf Grund eines Gesetzes geregelt werden.''}}
\newcommand{\ggfnArtTwelveOneZH}{\ggnotezhonce{art-12-1}{第12條第1項：「所有德國人均有自由選擇職業、工作場所及訓練場所之權利。職業行使得以法律或基於法律加以規範。」}}

% ===== Art. 19 =====
\newcommand{\ggfnArtNineteenTwoDE}{\ggnotedeonce{art-19-2}{Art. 19 Abs. 2: ``In keinem Falle darf ein Grundrecht in seinem Wesensgehalt angetastet werden.''}}
\newcommand{\ggfnArtNineteenTwoZH}{\ggnotezhonce{art-19-2}{第19條第2項：「任何情形下，基本權利之本質內容均不得受到侵害。」}}

% ===== Art. 20 =====
\newcommand{\ggfnArtTwentyOneDE}{\ggnotedeonce{art-20-1}{Art. 20 Abs. 1: ``Die Bundesrepublik Deutschland ist ein demokratischer und sozialer Bundesstaat.''}}
\newcommand{\ggfnArtTwentyOneZH}{\ggnotezhonce{art-20-1}{第20條第1項：「德意志聯邦共和國為民主及社會之聯邦國。」}}
\newcommand{\ggfnArtTwentyThreeDE}{\ggnotedeonce{art-20-3}{Art. 20 Abs. 3: ``Die Gesetzgebung ist an die verfassungsmaessige Ordnung, die vollziehende Gewalt und die Rechtsprechung sind an Gesetz und Recht gebunden.''}}
\newcommand{\ggfnArtTwentyThreeZH}{\ggnotezhonce{art-20-3}{第20條第3項：「立法受合憲秩序拘束；行政權與司法權受法律與法拘束。」}}

% ===== Art. 26 + Art. 143（墮胎一案特有）=====
\newcommand{\ggfnArtTwoSixAndOneFourThreeDE}{\ggnotedeonce{art-26-143}{Art. 26 Abs. 1 Satz 2: ``Sie sind unter Strafe zu stellen.'' Art. 143 in der urspruenglichen Fassung (24. Mai 1949 bis 1. September 1951) enthielt Strafvorschriften zum Hochverrat; 1975 war Art. 143 bereits weggefallen.}}
\newcommand{\ggfnArtTwoSixAndOneFourThreeZH}{\ggnotezhonce{art-26-143}{第26條第1項第2句：「此等行為應規定為犯罪處罰。」第143條原始版本（1949年5月24日至1951年9月1日）含有關於叛國罪之刑罰規定；至1975年時，第143條已廢止。}}

% ===== Art. 28 =====
\newcommand{\ggfnArtTwentyEightOneDE}{\ggnotedeonce{art-28-1-1}{Art. 28 Abs. 1 Satz 1: ``Die verfassungsmaessige Ordnung in den Laendern muss den Grundsaetzen des republikanischen, demokratischen und sozialen Rechtsstaates im Sinne dieses Grundgesetzes entsprechen.''}}
\newcommand{\ggfnArtTwentyEightOneZH}{\ggnotezhonce{art-28-1-1}{第28條第1項第1句：「各邦之合憲秩序，須符合本基本法意義下共和、民主及社會法治國之原則。」}}

% ===== Art. 30 =====
\newcommand{\ggfnArtThirtyDE}{\ggnotedeonce{art-30}{Art. 30: ``Die Ausuebung der staatlichen Befugnisse und die Erfuellung der staatlichen Aufgaben ist Sache der Laender, soweit dieses Grundgesetz keine andere Regelung trifft oder zulaesst.''}}
\newcommand{\ggfnArtThirtyZH}{\ggnotezhonce{art-30}{第30條：「國家權限之行使與國家任務之履行，屬於各邦，但本基本法另有規定或容許者，不在此限。」}}

% ===== Art. 74 =====
\newcommand{\ggfnArtSeventyFourOneDE}{\ggnotedeonce{art-74-1}{Art. 74 Nr. 1: Die konkurrierende Gesetzgebung erstreckt sich auf ``das buergerliche Recht, das Strafrecht und den Strafvollzug, die Gerichtsverfassung, das gerichtliche Verfahren, die Rechtsanwaltschaft, das Notariat und die Rechtsberatung''.}}
\newcommand{\ggfnArtSeventyFourOneZH}{\ggnotezhonce{art-74-1}{第74條第1款：競合立法及於「民法、刑法及刑罰執行、法院組織、訴訟程序、律師制度、公證制度及法律諮詢」。}}
\newcommand{\ggfnArtSeventyFourSevenDE}{\ggnotedeonce{art-74-7}{Art. 74 Nr. 7: Die konkurrierende Gesetzgebung erstreckt sich auf ``die oeffentliche Fuersorge''.}}
\newcommand{\ggfnArtSeventyFourSevenZH}{\ggnotezhonce{art-74-7}{第74條第7款：競合立法及於「公共扶助」。}}
\newcommand{\ggfnArtSeventyFourTwelveDE}{\ggnotedeonce{art-74-12}{Art. 74 Nr. 12: Die konkurrierende Gesetzgebung erstreckt sich auf ``das Arbeitsrecht einschliesslich der Betriebsverfassung, des Arbeitsschutzes und der Arbeitsvermittlung sowie die Sozialversicherung einschliesslich der Arbeitslosenversicherung''.}}
\newcommand{\ggfnArtSeventyFourTwelveZH}{\ggnotezhonce{art-74-12}{第74條第12款：競合立法及於「勞動法，包括企業組織、勞動保護及就業媒合，以及社會保險，包括失業保險」。}}

% ===== Art. 77（墮胎一案特有）=====
\newcommand{\ggfnArtSevenSevenThreeDE}{\ggnotedeonce{art-77-3}{Art. 77 Abs. 3 Satz 1: ``Soweit zu einem Gesetze die Zustimmung des Bundesrates nicht erforderlich ist, kann der Bundesrat, wenn das Verfahren nach Absatz 2 beendigt ist, gegen ein vom Bundestage beschlossenes Gesetz binnen zwei Wochen Einspruch einlegen.''}}
\newcommand{\ggfnArtSevenSevenThreeZH}{\ggnotezhonce{art-77-3}{第77條第3項第1句：「法律無須聯邦參議院同意者，於第2項程序終了後，聯邦參議院得於二週內對聯邦議院通過之法律提出異議。」}}

% ===== Art. 80 =====
\newcommand{\ggfnArtEightyOneOneDE}{\ggnotedeonce{art-80-1-1}{Art. 80 Abs. 1 Satz 1: ``Durch Gesetz koennen die Bundesregierung, ein Bundesminister oder die Landesregierungen ermaechtigt werden, Rechtsverordnungen zu erlassen.''}}
\newcommand{\ggfnArtEightyOneOneZH}{\ggnotezhonce{art-80-1-1}{第80條第1項第1句：「得以法律授權聯邦政府、聯邦部長或邦政府發布法規命令。」}}

% ===== Art. 83 =====
\newcommand{\ggfnArtEightyThreeDE}{\ggnotedeonce{art-83}{Art. 83: ``Die Laender fuehren die Bundesgesetze als eigene Angelegenheit aus, soweit dieses Grundgesetz nichts anderes bestimmt oder zulaesst.''}}
\newcommand{\ggfnArtEightyThreeZH}{\ggnotezhonce{art-83}{第83條：「各邦以自己事務執行聯邦法律，但本基本法另有規定或容許者，不在此限。」}}

% ===== Art. 84 =====
\newcommand{\ggfnArtEightyFourOneDE}{\ggnotedeonce{art-84-1}{Art. 84 Abs. 1: ``Fuehren die Laender die Bundesgesetze als eigene Angelegenheit aus, so regeln sie die Einrichtung der Behoerden und das Verwaltungsverfahren, soweit nicht Bundesgesetze mit Zustimmung des Bundesrates etwas anderes bestimmen.''}}
\newcommand{\ggfnArtEightyFourOneZH}{\ggnotezhonce{art-84-1}{第84條第1項：「各邦以自己事務執行聯邦法律時，除聯邦法律經聯邦參議院同意另有規定外，由各邦規範機關之設置及行政程序。」}}
% 別名（墮胎一案使用之縮寫命名）
\let\ggfnArtEightFourOneDE\ggfnArtEightyFourOneDE
\let\ggfnArtEightFourOneZH\ggfnArtEightyFourOneZH

% ===== Art. 93 =====
\newcommand{\ggfnArtNinetyThreeOneTwoDE}{\ggnotedeonce{art-93-1-2}{Art. 93 Abs. 1 Nr. 2: Das Bundesverfassungsgericht entscheidet ``bei Meinungsverschiedenheiten oder Zweifeln ueber die foermliche und sachliche Vereinbarkeit von Bundesrecht oder Landesrecht mit diesem Grundgesetze ... auf Antrag der Bundesregierung, einer Landesregierung oder eines Drittels der Mitglieder des Bundestages''.}}
\newcommand{\ggfnArtNinetyThreeOneTwoZH}{\ggnotezhonce{art-93-1-2}{第93條第1項第2款：聯邦憲法法院就「聯邦法或邦法在形式及實質上是否與本基本法相容之意見分歧或疑義……依聯邦政府、邦政府或聯邦議院三分之一議員之聲請」作成裁判。}}
% 別名（墮胎一案使用之縮寫命名）
\let\ggfnArtNineThreeOneTwoDE\ggfnArtNinetyThreeOneTwoDE
\let\ggfnArtNineThreeOneTwoZH\ggfnArtNinetyThreeOneTwoZH

% ===== Art. 102 =====
\newcommand{\ggfnArtOneZeroTwoDE}{\ggnotedeonce{art-102}{Art. 102: ``Die Todesstrafe ist abgeschafft.''}}
\newcommand{\ggfnArtOneZeroTwoZH}{\ggnotezhonce{art-102}{第102條：「死刑廢止。」}}

% ===== Art. 103 =====
\newcommand{\ggfnArtOneZeroThreeTwoDE}{\ggnotedeonce{art-103-2}{Art. 103 Abs. 2: ``Eine Tat kann nur bestraft werden, wenn die Strafbarkeit gesetzlich bestimmt war, bevor die Tat begangen wurde.''}}
\newcommand{\ggfnArtOneZeroThreeTwoZH}{\ggnotezhonce{art-103-2}{第103條第2項：「行為之可罰性，須於行為前以法律定之，始得處罰。」}}

% ===== Art. 104 =====
\newcommand{\ggfnArtOneZeroFourOneDE}{\ggnotedeonce{art-104-1}{Art. 104 Abs. 1: ``Die Freiheit der Person kann nur auf Grund eines foermlichen Gesetzes und nur unter Beachtung der darin vorgeschriebenen Formen beschraenkt werden.''}}
\newcommand{\ggfnArtOneZeroFourOneZH}{\ggnotezhonce{art-104-1}{第104條第1項：「人身自由僅得基於形式法律，並遵守其中規定之形式加以限制。」}}

% ===== 組合腳註：Art. 3 + Art. 6（墮胎一案特有）=====
\newcommand{\ggfnArtThreeSixDE}{\ggnotedeonce{art-3-6}{Art. 3: ``Alle Menschen sind vor dem Gesetz gleich. Maenner und Frauen sind gleichberechtigt. Niemand darf wegen seines Geschlechtes, seiner Abstammung, seiner Rasse, seiner Sprache, seiner Heimat und Herkunft, seines Glaubens, seiner religioesen oder politischen Anschauungen benachteiligt oder bevorzugt werden.'' Art. 6 Abs. 1: ``Ehe und Familie stehen unter dem besonderen Schutze der staatlichen Ordnung.'' Abs. 2 Satz 1: ``Pflege und Erziehung der Kinder sind das natuerliche Recht der Eltern und die zuvoerderst ihnen obliegende Pflicht.'' Abs. 4: ``Jede Mutter hat Anspruch auf den Schutz und die Fuersorge der Gemeinschaft.''}\ggmarknotede{art-3}\ggmarknotede{art-6}\ggmarknotede{art-6-1-4}}
\newcommand{\ggfnArtThreeSixZH}{\ggnotezhonce{art-3-6}{第3條：「所有人在法律前一律平等。男女享有平等權利。任何人不得因其性別、血統、種族、語言、故鄉及出身、信仰、宗教或政治見解而受不利益或優惠。」第6條第1項：「婚姻與家庭受國家秩序之特別保護。」第2項第1句：「照護及教育子女，為父母之自然權利，並為首先由其承擔之義務。」第4項：「每一母親均有請求共同體保護與照顧之權利。」}\ggmarknotezh{art-3}\ggmarknotezh{art-6}\ggmarknotezh{art-6-1-4}}

% ===== 組合腳註：Art. 6(1) + Art. 6(4)（墮胎一案特有）=====
\newcommand{\ggfnArtSixOneFourDE}{\ggnotedeonce{art-6-1-4}{Art. 6 Abs. 1: ``Ehe und Familie stehen unter dem besonderen Schutze der staatlichen Ordnung.'' Art. 6 Abs. 4: ``Jede Mutter hat Anspruch auf den Schutz und die Fuersorge der Gemeinschaft.''}\ggmarknotede{art-6}}
\newcommand{\ggfnArtSixOneFourZH}{\ggnotezhonce{art-6-1-4}{第6條第1項：「婚姻與家庭受國家秩序之特別保護。」第4項：「每一母親均有請求共同體保護與照顧之權利。」}\ggmarknotezh{art-6}}

% ===== 組合腳註：Art. 1(1) + Art. 2(2)(1)（墮胎二案特有）=====
\newcommand{\ggfnArtOneOneTwoTwoOneDE}{\ggnotedeonce{art-1-1+2-2-1}{Art. 1 Abs. 1: ``Die Wuerde des Menschen ist unantastbar. Sie zu achten und zu schuetzen ist Verpflichtung aller staatlichen Gewalt.'' Art. 2 Abs. 2 Satz 1: ``Jeder hat das Recht auf Leben und koerperliche Unversehrtheit.''}\ggmarknotede{art-1-1}\ggmarknotede{art-2-2-1}}
\newcommand{\ggfnArtOneOneTwoTwoOneZH}{\ggnotezhonce{art-1-1+2-2-1}{第1條第1項：「人之尊嚴不可侵犯。尊重及保護之，為一切國家權力之義務。」第2條第2項第1句：「人人有生命與身體完整性之權利。」}\ggmarknotezh{art-1-1}\ggmarknotezh{art-2-2-1}}

% ===== 組合腳註：Art. 1(1) + Art. 2(2)（墮胎二案特有）=====
\newcommand{\ggfnArtOneOneTwoTwoDE}{\ggnotedeonce{art-1-1+2-2}{Art. 1 Abs. 1: ``Die Wuerde des Menschen ist unantastbar. Sie zu achten und zu schuetzen ist Verpflichtung aller staatlichen Gewalt.'' Art. 2 Abs. 2: ``Jeder hat das Recht auf Leben und koerperliche Unversehrtheit. Die Freiheit der Person ist unverletzlich. In diese Rechte darf nur auf Grund eines Gesetzes eingegriffen werden.''}\ggmarknotede{art-1-1}\ggmarknotede{art-2-2}\ggmarknotede{art-2-2-1}}
\newcommand{\ggfnArtOneOneTwoTwoZH}{\ggnotezhonce{art-1-1+2-2}{第1條第1項：「人之尊嚴不可侵犯。尊重及保護之，為一切國家權力之義務。」第2條第2項：「人人有生命與身體完整性之權利。人身自由不可侵犯。此等權利僅得依法律加以干預。」}\ggmarknotezh{art-1-1}\ggmarknotezh{art-2-2}\ggmarknotezh{art-2-2-1}}

% ===== 組合腳註：Art. 1(1) + Art. 2(1)（墮胎二案特有）=====
\newcommand{\ggfnArtOneOneTwoOneDE}{\ggnotedeonce{art-1-1+2-1}{Art. 1 Abs. 1: ``Die Wuerde des Menschen ist unantastbar. Sie zu achten und zu schuetzen ist Verpflichtung aller staatlichen Gewalt.'' Art. 2 Abs. 1: ``Jeder hat das Recht auf die freie Entfaltung seiner Persoenlichkeit, soweit er nicht die Rechte anderer verletzt und nicht gegen die verfassungsmaessige Ordnung oder das Sittengesetz verstoesst.''}\ggmarknotede{art-1-1}\ggmarknotede{art-2-1}}
\newcommand{\ggfnArtOneOneTwoOneZH}{\ggnotezhonce{art-1-1+2-1}{第1條第1項：「人之尊嚴不可侵犯。尊重及保護之，為一切國家權力之義務。」第2條第1項：「人人有自由發展其人格之權利，但不得侵害他人權利，亦不得違反合憲秩序或道德律。」}\ggmarknotezh{art-1-1}\ggmarknotezh{art-2-1}}

% ===== 組合腳註：Art. 2(2)(1) + Art. 1(1)（墮胎一案特有，含交叉標記）=====
\newcommand{\ggfnLifeDignityDE}{\ggnotedeonce{life-dignity}{Art. 2 Abs. 2 Satz 1: ``Jeder hat das Recht auf Leben und koerperliche Unversehrtheit.'' Art. 1 Abs. 1: ``Die Wuerde des Menschen ist unantastbar. Sie zu achten und zu schuetzen ist Verpflichtung aller staatlichen Gewalt.''}\ggmarknotede{art-2-2-1}\ggmarknotede{art-1-1}\ggmarknotede{art-1-1-2}}
\newcommand{\ggfnLifeDignityZH}{\ggnotezhonce{life-dignity}{第2條第2項第1句：「人人有生命與身體完整性之權利。」第1條第1項：「人之尊嚴不可侵犯。尊重及保護之，為一切國家權力之義務。」}\ggmarknotezh{art-2-2-1}\ggmarknotezh{art-1-1}\ggmarknotezh{art-1-1-2}}

% ===== 組合腳註：Art. 1(1) + Art. 2(2)（墮胎一案特有，條件判斷版）=====
\newcommand{\ggfnArtOneTwoTwoDE}{\ifcsname ggnoteseen@de@life-dignity\endcsname\ggfnArtTwoTwoRestDE{}\else\ggnotedeonce{art-1-2-2}{Art. 1 Abs. 1: ``Die Wuerde des Menschen ist unantastbar. Sie zu achten und zu schuetzen ist Verpflichtung aller staatlichen Gewalt.'' Art. 2 Abs. 2: ``Jeder hat das Recht auf Leben und koerperliche Unversehrtheit. Die Freiheit der Person ist unverletzlich. In diese Rechte darf nur auf Grund eines Gesetzes eingegriffen werden.''}\ggmarknotede{art-1-1}\ggmarknotede{art-1-1-2}\ggmarknotede{art-2-2}\ggmarknotede{art-2-2-1}\fi}
\newcommand{\ggfnArtOneTwoTwoZH}{\ifcsname ggnoteseen@zh@life-dignity\endcsname\ggfnArtTwoTwoRestZH{}\else\ggnotezhonce{art-1-2-2}{第1條第1項：「人之尊嚴不可侵犯。尊重及保護之，為一切國家權力之義務。」第2條第2項：「人人有生命與身體完整性之權利。人身自由不可侵犯。此等權利僅得依法律加以干預。」}\ggmarknotezh{art-1-1}\ggmarknotezh{art-1-1-2}\ggmarknotezh{art-2-2}\ggmarknotezh{art-2-2-1}\fi}

% ===== 組合腳註：Art. 1(1) + Art. 19(2)（墮胎一案特有，條件判斷版）=====
\newcommand{\ggfnArtOneNineteenDE}{\ifcsname ggnoteseen@de@art-1-1\endcsname\ggfnArtNineteenTwoDE{}\else\ggnotedeonce{art-1-19}{Art. 1 Abs. 1: ``Die Wuerde des Menschen ist unantastbar. Sie zu achten und zu schuetzen ist Verpflichtung aller staatlichen Gewalt.'' Art. 19 Abs. 2: ``In keinem Falle darf ein Grundrecht in seinem Wesensgehalt angetastet werden.''}\ggmarknotede{art-1-1}\ggmarknotede{art-1-1-2}\ggmarknotede{art-19-2}\fi}
\newcommand{\ggfnArtOneNineteenZH}{\ifcsname ggnoteseen@zh@art-1-1\endcsname\ggfnArtNineteenTwoZH{}\else\ggnotezhonce{art-1-19}{第1條第1項：「人之尊嚴不可侵犯。尊重及保護之，為一切國家權力之義務。」第19條第2項：「任何情形下，基本權利之本質內容均不得受到侵害。」}\ggmarknotezh{art-1-1}\ggmarknotezh{art-1-1-2}\ggmarknotezh{art-19-2}\fi}


% ===== GG 版本旗標（\ggversionde/zh 由各案在 % bverfge_gg.tex — 德國墮胎案 GG 條文腳註共用定義
% 使用方式：在各案 .tex 中先定義 \ggversionde/\ggversionzh，再 % bverfge_gg.tex — 德國墮胎案 GG 條文腳註共用定義
% 使用方式：在各案 .tex 中先定義 \ggversionde/\ggversionzh，再 \input{../../共用LaTeX/bverfge_gg}

% ===== GG 版本旗標（\ggversionde/zh 由各案在 \input{../../共用LaTeX/bverfge_gg} 之前定義）=====
\newif\ifggversionnotede
\newif\ifggversionnotezh

% ===== 編者註腳基礎設施（首次標記模式）=====
\newif\ifeditornotelabelde
\newif\ifeditornotelabelzh
\newcommand{\editornotelabelde}{\ifeditornotelabelde\else\emph{Hrsg.-Anm. (nachfolgend ohne Kennzeichnung)}: \global\editornotelabeldetrue\fi}
\newcommand{\editornotelabelzh}{\ifeditornotelabelzh\else 編者註（下同）：\global\editornotelabelzhtrue\fi}
\newcommand{\editornotede}[1]{\ifshoweditornotes\footnote{\normalfont\footnotesize\editornotelabelde #1}\else\ifclassroommode\footnote{\normalfont\footnotesize\editornotelabelde #1}\fi\fi}
\newcommand{\editornotezh}[1]{\ifshoweditornotes\footnote{\normalfont\footnotesize\editornotelabelzh #1}\else\ifclassroommode\footnote{\normalfont\footnotesize\editornotelabelzh #1}\fi\fi}

% ===== GG 引用系統（\ggversionde/zh 由各案在 \input 前定義）=====
\newcommand{\ggmarknotede}[1]{\global\expandafter\let\csname ggnoteseen@de@#1\endcsname\empty}
\newcommand{\ggmarknotezh}[1]{\global\expandafter\let\csname ggnoteseen@zh@#1\endcsname\empty}
\newcommand{\ggnotedeonce}[2]{\ifcsname ggnoteseen@de@#1\endcsname\else\global\expandafter\let\csname ggnoteseen@de@#1\endcsname\empty\ggnotede{#2}\fi}
\newcommand{\ggnotezhonce}[2]{\ifcsname ggnoteseen@zh@#1\endcsname\else\global\expandafter\let\csname ggnoteseen@zh@#1\endcsname\empty\ggnotezh{#2}\fi}
\newcommand{\ggnotede}[1]{\editornotede{\ggversionde #1}}
\newcommand{\ggnotezh}[1]{\editornotezh{\ggversionzh #1}}

% ===== Art. 1 =====
\newcommand{\ggfnArtOneOneDE}{\ggnotedeonce{art-1-1}{Art. 1 Abs. 1: ``Die Wuerde des Menschen ist unantastbar. Sie zu achten und zu schuetzen ist Verpflichtung aller staatlichen Gewalt.''}}
\newcommand{\ggfnArtOneOneZH}{\ggnotezhonce{art-1-1}{第1條第1項：「人之尊嚴不可侵犯。尊重及保護之，為一切國家權力之義務。」}}
\newcommand{\ggfnArtOneOneTwoDE}{\ggnotedeonce{art-1-1-2}{Art. 1 Abs. 1 Satz 2: ``Sie zu achten und zu schuetzen ist Verpflichtung aller staatlichen Gewalt.''}}
\newcommand{\ggfnArtOneOneTwoZH}{\ggnotezhonce{art-1-1-2}{第1條第1項第2句：「尊重及保護之，為一切國家權力之義務。」}}

% ===== Art. 2 =====
\newcommand{\ggfnArtTwoOneDE}{\ggnotedeonce{art-2-1}{Art. 2 Abs. 1: ``Jeder hat das Recht auf die freie Entfaltung seiner Persoenlichkeit, soweit er nicht die Rechte anderer verletzt und nicht gegen die verfassungsmaessige Ordnung oder das Sittengesetz verstoesst.''}}
\newcommand{\ggfnArtTwoOneZH}{\ggnotezhonce{art-2-1}{第2條第1項：「人人有自由發展其人格之權利，但不得侵害他人權利，亦不得違反合憲秩序或道德律。」}}
% 簡易預設版（墮胎一案以 \renewcommand 覆寫為含條件邏輯之版本）
\newcommand{\ggfnArtTwoTwoDE}{\ggnotedeonce{art-2-2}{Art. 2 Abs. 2: ``Jeder hat das Recht auf Leben und koerperliche Unversehrtheit. Die Freiheit der Person ist unverletzlich. In diese Rechte darf nur auf Grund eines Gesetzes eingegriffen werden.''}\ggmarknotede{art-2-2-1}}
\newcommand{\ggfnArtTwoTwoZH}{\ggnotezhonce{art-2-2}{第2條第2項：「人人有生命與身體完整性之權利。人身自由不可侵犯。此等權利僅得依法律加以干預。」}\ggmarknotezh{art-2-2-1}}
\newcommand{\ggfnArtTwoTwoOneDE}{\ggnotedeonce{art-2-2-1}{Art. 2 Abs. 2 Satz 1: ``Jeder hat das Recht auf Leben und koerperliche Unversehrtheit.''}}
\newcommand{\ggfnArtTwoTwoOneZH}{\ggnotezhonce{art-2-2-1}{第2條第2項第1句：「人人有生命與身體完整性之權利。」}}
\newcommand{\ggfnArtTwoTwoThreeDE}{\ggnotedeonce{art-2-2-3}{Art. 2 Abs. 2 Satz 3: ``In diese Rechte darf nur auf Grund eines Gesetzes eingegriffen werden.''}}
\newcommand{\ggfnArtTwoTwoThreeZH}{\ggnotezhonce{art-2-2-3}{第2條第2項第3句：「此等權利僅得依法律加以干預。」}}
% 墮胎一案特有：第2條第2項第2–3句（Art. 2(2) 去掉第1句後的「餘文」）
\newcommand{\ggfnArtTwoTwoRestDE}{\ggnotedeonce{art-2-2-rest}{Art. 2 Abs. 2 Saetze 2-3: ``Die Freiheit der Person ist unverletzlich. In diese Rechte darf nur auf Grund eines Gesetzes eingegriffen werden.''}\ggmarknotede{art-2-2}}
\newcommand{\ggfnArtTwoTwoRestZH}{\ggnotezhonce{art-2-2-rest}{第2條第2項第2、3句：「人身自由不可侵犯。此等權利僅得依法律加以干預。」}\ggmarknotezh{art-2-2}}

% ===== Art. 3 =====
\newcommand{\ggfnArtThreeDE}{\ggnotedeonce{art-3}{Art. 3: ``Alle Menschen sind vor dem Gesetz gleich. Maenner und Frauen sind gleichberechtigt. Niemand darf wegen seines Geschlechtes, seiner Abstammung, seiner Rasse, seiner Sprache, seiner Heimat und Herkunft, seines Glaubens, seiner religioesen oder politischen Anschauungen benachteiligt oder bevorzugt werden.''}}
\newcommand{\ggfnArtThreeZH}{\ggnotezhonce{art-3}{第3條：「所有人在法律前一律平等。男女享有平等權利。任何人不得因其性別、血統、種族、語言、故鄉及出身、信仰、宗教或政治見解而受不利益或優惠。」}}
\newcommand{\ggfnArtThreeTwoDE}{\ggnotedeonce{art-3-2}{Art. 3 Abs. 2: ``Maenner und Frauen sind gleichberechtigt.''}}
\newcommand{\ggfnArtThreeTwoZH}{\ggnotezhonce{art-3-2}{第3條第2項：「男女享有平等權利。」}}

% ===== Art. 4 =====
\newcommand{\ggfnArtFourOneDE}{\ggnotedeonce{art-4-1}{Art. 4 Abs. 1: ``Die Freiheit des Glaubens, des Gewissens und die Freiheit des religioesen und weltanschaulichen Bekenntnisses sind unverletzlich.''}}
\newcommand{\ggfnArtFourOneZH}{\ggnotezhonce{art-4-1}{第4條第1項：「信仰自由、良心自由，以及宗教與世界觀信念自由，不可侵犯。」}}

% ===== Art. 5 =====
\newcommand{\ggfnArtFiveOneDE}{\ggnotedeonce{art-5-1}{Art. 5 Abs. 1: ``Jeder hat das Recht, seine Meinung in Wort, Schrift und Bild frei zu aeussern und zu verbreiten und sich aus allgemein zugaenglichen Quellen ungehindert zu unterrichten. Die Pressefreiheit und die Freiheit der Berichterstattung durch Rundfunk und Film werden gewaehrleistet. Eine Zensur findet nicht statt.''}}
\newcommand{\ggfnArtFiveOneZH}{\ggnotezhonce{art-5-1}{第5條第1項：「人人有以言語、文字及圖像自由表達及傳播其意見，並由一般可接近來源不受阻礙獲取資訊之權利。新聞自由及廣播、電影報導自由受保障。不設審查制度。」}}

% ===== Art. 6 =====
\newcommand{\ggfnArtSixDE}{\ggnotedeonce{art-6}{Art. 6 Abs. 1: ``Ehe und Familie stehen unter dem besonderen Schutze der staatlichen Ordnung.'' Abs. 2 Satz 1: ``Pflege und Erziehung der Kinder sind das natuerliche Recht der Eltern und die zuvoerderst ihnen obliegende Pflicht.'' Abs. 4: ``Jede Mutter hat Anspruch auf den Schutz und die Fuersorge der Gemeinschaft.''}}
\newcommand{\ggfnArtSixZH}{\ggnotezhonce{art-6}{第6條第1項：「婚姻與家庭受國家秩序之特別保護。」第2項第1句：「照護及教育子女，為父母之自然權利，並為首先由其承擔之義務。」第4項：「每一母親均有請求共同體保護與照顧之權利。」}}
\newcommand{\ggfnArtSixOneDE}{\ggnotedeonce{art-6-1}{Art. 6 Abs. 1: ``Ehe und Familie stehen unter dem besonderen Schutze der staatlichen Ordnung.''}}
\newcommand{\ggfnArtSixOneZH}{\ggnotezhonce{art-6-1}{第6條第1項：「婚姻與家庭受國家秩序之特別保護。」}}
\newcommand{\ggfnArtSixTwoDE}{\ggnotedeonce{art-6-2}{Art. 6 Abs. 2: ``Pflege und Erziehung der Kinder sind das natuerliche Recht der Eltern und die zuvoerderst ihnen obliegende Pflicht. Ueber ihre Betaetigung wacht die staatliche Gemeinschaft.''}}
\newcommand{\ggfnArtSixTwoZH}{\ggnotezhonce{art-6-2}{第6條第2項：「照護及教育子女，為父母之自然權利，並為首先由其承擔之義務。國家共同體監督其行使。」}}
\newcommand{\ggfnArtSixTwoOneDE}{\ggnotedeonce{art-6-2-1}{Art. 6 Abs. 2 Satz 1: ``Pflege und Erziehung der Kinder sind das natuerliche Recht der Eltern und die zuvoerderst ihnen obliegende Pflicht.''}}
\newcommand{\ggfnArtSixTwoOneZH}{\ggnotezhonce{art-6-2-1}{第6條第2項第1句：「照護及教育子女，為父母之自然權利，並為首先由其承擔之義務。」}}
\newcommand{\ggfnArtSixFourDE}{\ggnotedeonce{art-6-4}{Art. 6 Abs. 4: ``Jede Mutter hat Anspruch auf den Schutz und die Fuersorge der Gemeinschaft.''}}
\newcommand{\ggfnArtSixFourZH}{\ggnotezhonce{art-6-4}{第6條第4項：「每一母親均有請求共同體保護與照顧之權利。」}}

% ===== Art. 12 =====
\newcommand{\ggfnArtTwelveOneDE}{\ggnotedeonce{art-12-1}{Art. 12 Abs. 1: ``Alle Deutschen haben das Recht, Beruf, Arbeitsplatz und Ausbildungsstaette frei zu waehlen. Die Berufsausuebung kann durch Gesetz oder auf Grund eines Gesetzes geregelt werden.''}}
\newcommand{\ggfnArtTwelveOneZH}{\ggnotezhonce{art-12-1}{第12條第1項：「所有德國人均有自由選擇職業、工作場所及訓練場所之權利。職業行使得以法律或基於法律加以規範。」}}

% ===== Art. 19 =====
\newcommand{\ggfnArtNineteenTwoDE}{\ggnotedeonce{art-19-2}{Art. 19 Abs. 2: ``In keinem Falle darf ein Grundrecht in seinem Wesensgehalt angetastet werden.''}}
\newcommand{\ggfnArtNineteenTwoZH}{\ggnotezhonce{art-19-2}{第19條第2項：「任何情形下，基本權利之本質內容均不得受到侵害。」}}

% ===== Art. 20 =====
\newcommand{\ggfnArtTwentyOneDE}{\ggnotedeonce{art-20-1}{Art. 20 Abs. 1: ``Die Bundesrepublik Deutschland ist ein demokratischer und sozialer Bundesstaat.''}}
\newcommand{\ggfnArtTwentyOneZH}{\ggnotezhonce{art-20-1}{第20條第1項：「德意志聯邦共和國為民主及社會之聯邦國。」}}
\newcommand{\ggfnArtTwentyThreeDE}{\ggnotedeonce{art-20-3}{Art. 20 Abs. 3: ``Die Gesetzgebung ist an die verfassungsmaessige Ordnung, die vollziehende Gewalt und die Rechtsprechung sind an Gesetz und Recht gebunden.''}}
\newcommand{\ggfnArtTwentyThreeZH}{\ggnotezhonce{art-20-3}{第20條第3項：「立法受合憲秩序拘束；行政權與司法權受法律與法拘束。」}}

% ===== Art. 26 + Art. 143（墮胎一案特有）=====
\newcommand{\ggfnArtTwoSixAndOneFourThreeDE}{\ggnotedeonce{art-26-143}{Art. 26 Abs. 1 Satz 2: ``Sie sind unter Strafe zu stellen.'' Art. 143 in der urspruenglichen Fassung (24. Mai 1949 bis 1. September 1951) enthielt Strafvorschriften zum Hochverrat; 1975 war Art. 143 bereits weggefallen.}}
\newcommand{\ggfnArtTwoSixAndOneFourThreeZH}{\ggnotezhonce{art-26-143}{第26條第1項第2句：「此等行為應規定為犯罪處罰。」第143條原始版本（1949年5月24日至1951年9月1日）含有關於叛國罪之刑罰規定；至1975年時，第143條已廢止。}}

% ===== Art. 28 =====
\newcommand{\ggfnArtTwentyEightOneDE}{\ggnotedeonce{art-28-1-1}{Art. 28 Abs. 1 Satz 1: ``Die verfassungsmaessige Ordnung in den Laendern muss den Grundsaetzen des republikanischen, demokratischen und sozialen Rechtsstaates im Sinne dieses Grundgesetzes entsprechen.''}}
\newcommand{\ggfnArtTwentyEightOneZH}{\ggnotezhonce{art-28-1-1}{第28條第1項第1句：「各邦之合憲秩序，須符合本基本法意義下共和、民主及社會法治國之原則。」}}

% ===== Art. 30 =====
\newcommand{\ggfnArtThirtyDE}{\ggnotedeonce{art-30}{Art. 30: ``Die Ausuebung der staatlichen Befugnisse und die Erfuellung der staatlichen Aufgaben ist Sache der Laender, soweit dieses Grundgesetz keine andere Regelung trifft oder zulaesst.''}}
\newcommand{\ggfnArtThirtyZH}{\ggnotezhonce{art-30}{第30條：「國家權限之行使與國家任務之履行，屬於各邦，但本基本法另有規定或容許者，不在此限。」}}

% ===== Art. 74 =====
\newcommand{\ggfnArtSeventyFourOneDE}{\ggnotedeonce{art-74-1}{Art. 74 Nr. 1: Die konkurrierende Gesetzgebung erstreckt sich auf ``das buergerliche Recht, das Strafrecht und den Strafvollzug, die Gerichtsverfassung, das gerichtliche Verfahren, die Rechtsanwaltschaft, das Notariat und die Rechtsberatung''.}}
\newcommand{\ggfnArtSeventyFourOneZH}{\ggnotezhonce{art-74-1}{第74條第1款：競合立法及於「民法、刑法及刑罰執行、法院組織、訴訟程序、律師制度、公證制度及法律諮詢」。}}
\newcommand{\ggfnArtSeventyFourSevenDE}{\ggnotedeonce{art-74-7}{Art. 74 Nr. 7: Die konkurrierende Gesetzgebung erstreckt sich auf ``die oeffentliche Fuersorge''.}}
\newcommand{\ggfnArtSeventyFourSevenZH}{\ggnotezhonce{art-74-7}{第74條第7款：競合立法及於「公共扶助」。}}
\newcommand{\ggfnArtSeventyFourTwelveDE}{\ggnotedeonce{art-74-12}{Art. 74 Nr. 12: Die konkurrierende Gesetzgebung erstreckt sich auf ``das Arbeitsrecht einschliesslich der Betriebsverfassung, des Arbeitsschutzes und der Arbeitsvermittlung sowie die Sozialversicherung einschliesslich der Arbeitslosenversicherung''.}}
\newcommand{\ggfnArtSeventyFourTwelveZH}{\ggnotezhonce{art-74-12}{第74條第12款：競合立法及於「勞動法，包括企業組織、勞動保護及就業媒合，以及社會保險，包括失業保險」。}}

% ===== Art. 77（墮胎一案特有）=====
\newcommand{\ggfnArtSevenSevenThreeDE}{\ggnotedeonce{art-77-3}{Art. 77 Abs. 3 Satz 1: ``Soweit zu einem Gesetze die Zustimmung des Bundesrates nicht erforderlich ist, kann der Bundesrat, wenn das Verfahren nach Absatz 2 beendigt ist, gegen ein vom Bundestage beschlossenes Gesetz binnen zwei Wochen Einspruch einlegen.''}}
\newcommand{\ggfnArtSevenSevenThreeZH}{\ggnotezhonce{art-77-3}{第77條第3項第1句：「法律無須聯邦參議院同意者，於第2項程序終了後，聯邦參議院得於二週內對聯邦議院通過之法律提出異議。」}}

% ===== Art. 80 =====
\newcommand{\ggfnArtEightyOneOneDE}{\ggnotedeonce{art-80-1-1}{Art. 80 Abs. 1 Satz 1: ``Durch Gesetz koennen die Bundesregierung, ein Bundesminister oder die Landesregierungen ermaechtigt werden, Rechtsverordnungen zu erlassen.''}}
\newcommand{\ggfnArtEightyOneOneZH}{\ggnotezhonce{art-80-1-1}{第80條第1項第1句：「得以法律授權聯邦政府、聯邦部長或邦政府發布法規命令。」}}

% ===== Art. 83 =====
\newcommand{\ggfnArtEightyThreeDE}{\ggnotedeonce{art-83}{Art. 83: ``Die Laender fuehren die Bundesgesetze als eigene Angelegenheit aus, soweit dieses Grundgesetz nichts anderes bestimmt oder zulaesst.''}}
\newcommand{\ggfnArtEightyThreeZH}{\ggnotezhonce{art-83}{第83條：「各邦以自己事務執行聯邦法律，但本基本法另有規定或容許者，不在此限。」}}

% ===== Art. 84 =====
\newcommand{\ggfnArtEightyFourOneDE}{\ggnotedeonce{art-84-1}{Art. 84 Abs. 1: ``Fuehren die Laender die Bundesgesetze als eigene Angelegenheit aus, so regeln sie die Einrichtung der Behoerden und das Verwaltungsverfahren, soweit nicht Bundesgesetze mit Zustimmung des Bundesrates etwas anderes bestimmen.''}}
\newcommand{\ggfnArtEightyFourOneZH}{\ggnotezhonce{art-84-1}{第84條第1項：「各邦以自己事務執行聯邦法律時，除聯邦法律經聯邦參議院同意另有規定外，由各邦規範機關之設置及行政程序。」}}
% 別名（墮胎一案使用之縮寫命名）
\let\ggfnArtEightFourOneDE\ggfnArtEightyFourOneDE
\let\ggfnArtEightFourOneZH\ggfnArtEightyFourOneZH

% ===== Art. 93 =====
\newcommand{\ggfnArtNinetyThreeOneTwoDE}{\ggnotedeonce{art-93-1-2}{Art. 93 Abs. 1 Nr. 2: Das Bundesverfassungsgericht entscheidet ``bei Meinungsverschiedenheiten oder Zweifeln ueber die foermliche und sachliche Vereinbarkeit von Bundesrecht oder Landesrecht mit diesem Grundgesetze ... auf Antrag der Bundesregierung, einer Landesregierung oder eines Drittels der Mitglieder des Bundestages''.}}
\newcommand{\ggfnArtNinetyThreeOneTwoZH}{\ggnotezhonce{art-93-1-2}{第93條第1項第2款：聯邦憲法法院就「聯邦法或邦法在形式及實質上是否與本基本法相容之意見分歧或疑義……依聯邦政府、邦政府或聯邦議院三分之一議員之聲請」作成裁判。}}
% 別名（墮胎一案使用之縮寫命名）
\let\ggfnArtNineThreeOneTwoDE\ggfnArtNinetyThreeOneTwoDE
\let\ggfnArtNineThreeOneTwoZH\ggfnArtNinetyThreeOneTwoZH

% ===== Art. 102 =====
\newcommand{\ggfnArtOneZeroTwoDE}{\ggnotedeonce{art-102}{Art. 102: ``Die Todesstrafe ist abgeschafft.''}}
\newcommand{\ggfnArtOneZeroTwoZH}{\ggnotezhonce{art-102}{第102條：「死刑廢止。」}}

% ===== Art. 103 =====
\newcommand{\ggfnArtOneZeroThreeTwoDE}{\ggnotedeonce{art-103-2}{Art. 103 Abs. 2: ``Eine Tat kann nur bestraft werden, wenn die Strafbarkeit gesetzlich bestimmt war, bevor die Tat begangen wurde.''}}
\newcommand{\ggfnArtOneZeroThreeTwoZH}{\ggnotezhonce{art-103-2}{第103條第2項：「行為之可罰性，須於行為前以法律定之，始得處罰。」}}

% ===== Art. 104 =====
\newcommand{\ggfnArtOneZeroFourOneDE}{\ggnotedeonce{art-104-1}{Art. 104 Abs. 1: ``Die Freiheit der Person kann nur auf Grund eines foermlichen Gesetzes und nur unter Beachtung der darin vorgeschriebenen Formen beschraenkt werden.''}}
\newcommand{\ggfnArtOneZeroFourOneZH}{\ggnotezhonce{art-104-1}{第104條第1項：「人身自由僅得基於形式法律，並遵守其中規定之形式加以限制。」}}

% ===== 組合腳註：Art. 3 + Art. 6（墮胎一案特有）=====
\newcommand{\ggfnArtThreeSixDE}{\ggnotedeonce{art-3-6}{Art. 3: ``Alle Menschen sind vor dem Gesetz gleich. Maenner und Frauen sind gleichberechtigt. Niemand darf wegen seines Geschlechtes, seiner Abstammung, seiner Rasse, seiner Sprache, seiner Heimat und Herkunft, seines Glaubens, seiner religioesen oder politischen Anschauungen benachteiligt oder bevorzugt werden.'' Art. 6 Abs. 1: ``Ehe und Familie stehen unter dem besonderen Schutze der staatlichen Ordnung.'' Abs. 2 Satz 1: ``Pflege und Erziehung der Kinder sind das natuerliche Recht der Eltern und die zuvoerderst ihnen obliegende Pflicht.'' Abs. 4: ``Jede Mutter hat Anspruch auf den Schutz und die Fuersorge der Gemeinschaft.''}\ggmarknotede{art-3}\ggmarknotede{art-6}\ggmarknotede{art-6-1-4}}
\newcommand{\ggfnArtThreeSixZH}{\ggnotezhonce{art-3-6}{第3條：「所有人在法律前一律平等。男女享有平等權利。任何人不得因其性別、血統、種族、語言、故鄉及出身、信仰、宗教或政治見解而受不利益或優惠。」第6條第1項：「婚姻與家庭受國家秩序之特別保護。」第2項第1句：「照護及教育子女，為父母之自然權利，並為首先由其承擔之義務。」第4項：「每一母親均有請求共同體保護與照顧之權利。」}\ggmarknotezh{art-3}\ggmarknotezh{art-6}\ggmarknotezh{art-6-1-4}}

% ===== 組合腳註：Art. 6(1) + Art. 6(4)（墮胎一案特有）=====
\newcommand{\ggfnArtSixOneFourDE}{\ggnotedeonce{art-6-1-4}{Art. 6 Abs. 1: ``Ehe und Familie stehen unter dem besonderen Schutze der staatlichen Ordnung.'' Art. 6 Abs. 4: ``Jede Mutter hat Anspruch auf den Schutz und die Fuersorge der Gemeinschaft.''}\ggmarknotede{art-6}}
\newcommand{\ggfnArtSixOneFourZH}{\ggnotezhonce{art-6-1-4}{第6條第1項：「婚姻與家庭受國家秩序之特別保護。」第4項：「每一母親均有請求共同體保護與照顧之權利。」}\ggmarknotezh{art-6}}

% ===== 組合腳註：Art. 1(1) + Art. 2(2)(1)（墮胎二案特有）=====
\newcommand{\ggfnArtOneOneTwoTwoOneDE}{\ggnotedeonce{art-1-1+2-2-1}{Art. 1 Abs. 1: ``Die Wuerde des Menschen ist unantastbar. Sie zu achten und zu schuetzen ist Verpflichtung aller staatlichen Gewalt.'' Art. 2 Abs. 2 Satz 1: ``Jeder hat das Recht auf Leben und koerperliche Unversehrtheit.''}\ggmarknotede{art-1-1}\ggmarknotede{art-2-2-1}}
\newcommand{\ggfnArtOneOneTwoTwoOneZH}{\ggnotezhonce{art-1-1+2-2-1}{第1條第1項：「人之尊嚴不可侵犯。尊重及保護之，為一切國家權力之義務。」第2條第2項第1句：「人人有生命與身體完整性之權利。」}\ggmarknotezh{art-1-1}\ggmarknotezh{art-2-2-1}}

% ===== 組合腳註：Art. 1(1) + Art. 2(2)（墮胎二案特有）=====
\newcommand{\ggfnArtOneOneTwoTwoDE}{\ggnotedeonce{art-1-1+2-2}{Art. 1 Abs. 1: ``Die Wuerde des Menschen ist unantastbar. Sie zu achten und zu schuetzen ist Verpflichtung aller staatlichen Gewalt.'' Art. 2 Abs. 2: ``Jeder hat das Recht auf Leben und koerperliche Unversehrtheit. Die Freiheit der Person ist unverletzlich. In diese Rechte darf nur auf Grund eines Gesetzes eingegriffen werden.''}\ggmarknotede{art-1-1}\ggmarknotede{art-2-2}\ggmarknotede{art-2-2-1}}
\newcommand{\ggfnArtOneOneTwoTwoZH}{\ggnotezhonce{art-1-1+2-2}{第1條第1項：「人之尊嚴不可侵犯。尊重及保護之，為一切國家權力之義務。」第2條第2項：「人人有生命與身體完整性之權利。人身自由不可侵犯。此等權利僅得依法律加以干預。」}\ggmarknotezh{art-1-1}\ggmarknotezh{art-2-2}\ggmarknotezh{art-2-2-1}}

% ===== 組合腳註：Art. 1(1) + Art. 2(1)（墮胎二案特有）=====
\newcommand{\ggfnArtOneOneTwoOneDE}{\ggnotedeonce{art-1-1+2-1}{Art. 1 Abs. 1: ``Die Wuerde des Menschen ist unantastbar. Sie zu achten und zu schuetzen ist Verpflichtung aller staatlichen Gewalt.'' Art. 2 Abs. 1: ``Jeder hat das Recht auf die freie Entfaltung seiner Persoenlichkeit, soweit er nicht die Rechte anderer verletzt und nicht gegen die verfassungsmaessige Ordnung oder das Sittengesetz verstoesst.''}\ggmarknotede{art-1-1}\ggmarknotede{art-2-1}}
\newcommand{\ggfnArtOneOneTwoOneZH}{\ggnotezhonce{art-1-1+2-1}{第1條第1項：「人之尊嚴不可侵犯。尊重及保護之，為一切國家權力之義務。」第2條第1項：「人人有自由發展其人格之權利，但不得侵害他人權利，亦不得違反合憲秩序或道德律。」}\ggmarknotezh{art-1-1}\ggmarknotezh{art-2-1}}

% ===== 組合腳註：Art. 2(2)(1) + Art. 1(1)（墮胎一案特有，含交叉標記）=====
\newcommand{\ggfnLifeDignityDE}{\ggnotedeonce{life-dignity}{Art. 2 Abs. 2 Satz 1: ``Jeder hat das Recht auf Leben und koerperliche Unversehrtheit.'' Art. 1 Abs. 1: ``Die Wuerde des Menschen ist unantastbar. Sie zu achten und zu schuetzen ist Verpflichtung aller staatlichen Gewalt.''}\ggmarknotede{art-2-2-1}\ggmarknotede{art-1-1}\ggmarknotede{art-1-1-2}}
\newcommand{\ggfnLifeDignityZH}{\ggnotezhonce{life-dignity}{第2條第2項第1句：「人人有生命與身體完整性之權利。」第1條第1項：「人之尊嚴不可侵犯。尊重及保護之，為一切國家權力之義務。」}\ggmarknotezh{art-2-2-1}\ggmarknotezh{art-1-1}\ggmarknotezh{art-1-1-2}}

% ===== 組合腳註：Art. 1(1) + Art. 2(2)（墮胎一案特有，條件判斷版）=====
\newcommand{\ggfnArtOneTwoTwoDE}{\ifcsname ggnoteseen@de@life-dignity\endcsname\ggfnArtTwoTwoRestDE{}\else\ggnotedeonce{art-1-2-2}{Art. 1 Abs. 1: ``Die Wuerde des Menschen ist unantastbar. Sie zu achten und zu schuetzen ist Verpflichtung aller staatlichen Gewalt.'' Art. 2 Abs. 2: ``Jeder hat das Recht auf Leben und koerperliche Unversehrtheit. Die Freiheit der Person ist unverletzlich. In diese Rechte darf nur auf Grund eines Gesetzes eingegriffen werden.''}\ggmarknotede{art-1-1}\ggmarknotede{art-1-1-2}\ggmarknotede{art-2-2}\ggmarknotede{art-2-2-1}\fi}
\newcommand{\ggfnArtOneTwoTwoZH}{\ifcsname ggnoteseen@zh@life-dignity\endcsname\ggfnArtTwoTwoRestZH{}\else\ggnotezhonce{art-1-2-2}{第1條第1項：「人之尊嚴不可侵犯。尊重及保護之，為一切國家權力之義務。」第2條第2項：「人人有生命與身體完整性之權利。人身自由不可侵犯。此等權利僅得依法律加以干預。」}\ggmarknotezh{art-1-1}\ggmarknotezh{art-1-1-2}\ggmarknotezh{art-2-2}\ggmarknotezh{art-2-2-1}\fi}

% ===== 組合腳註：Art. 1(1) + Art. 19(2)（墮胎一案特有，條件判斷版）=====
\newcommand{\ggfnArtOneNineteenDE}{\ifcsname ggnoteseen@de@art-1-1\endcsname\ggfnArtNineteenTwoDE{}\else\ggnotedeonce{art-1-19}{Art. 1 Abs. 1: ``Die Wuerde des Menschen ist unantastbar. Sie zu achten und zu schuetzen ist Verpflichtung aller staatlichen Gewalt.'' Art. 19 Abs. 2: ``In keinem Falle darf ein Grundrecht in seinem Wesensgehalt angetastet werden.''}\ggmarknotede{art-1-1}\ggmarknotede{art-1-1-2}\ggmarknotede{art-19-2}\fi}
\newcommand{\ggfnArtOneNineteenZH}{\ifcsname ggnoteseen@zh@art-1-1\endcsname\ggfnArtNineteenTwoZH{}\else\ggnotezhonce{art-1-19}{第1條第1項：「人之尊嚴不可侵犯。尊重及保護之，為一切國家權力之義務。」第19條第2項：「任何情形下，基本權利之本質內容均不得受到侵害。」}\ggmarknotezh{art-1-1}\ggmarknotezh{art-1-1-2}\ggmarknotezh{art-19-2}\fi}


% ===== GG 版本旗標（\ggversionde/zh 由各案在 % bverfge_gg.tex — 德國墮胎案 GG 條文腳註共用定義
% 使用方式：在各案 .tex 中先定義 \ggversionde/\ggversionzh，再 \input{../../共用LaTeX/bverfge_gg}

% ===== GG 版本旗標（\ggversionde/zh 由各案在 \input{../../共用LaTeX/bverfge_gg} 之前定義）=====
\newif\ifggversionnotede
\newif\ifggversionnotezh

% ===== 編者註腳基礎設施（首次標記模式）=====
\newif\ifeditornotelabelde
\newif\ifeditornotelabelzh
\newcommand{\editornotelabelde}{\ifeditornotelabelde\else\emph{Hrsg.-Anm. (nachfolgend ohne Kennzeichnung)}: \global\editornotelabeldetrue\fi}
\newcommand{\editornotelabelzh}{\ifeditornotelabelzh\else 編者註（下同）：\global\editornotelabelzhtrue\fi}
\newcommand{\editornotede}[1]{\ifshoweditornotes\footnote{\normalfont\footnotesize\editornotelabelde #1}\else\ifclassroommode\footnote{\normalfont\footnotesize\editornotelabelde #1}\fi\fi}
\newcommand{\editornotezh}[1]{\ifshoweditornotes\footnote{\normalfont\footnotesize\editornotelabelzh #1}\else\ifclassroommode\footnote{\normalfont\footnotesize\editornotelabelzh #1}\fi\fi}

% ===== GG 引用系統（\ggversionde/zh 由各案在 \input 前定義）=====
\newcommand{\ggmarknotede}[1]{\global\expandafter\let\csname ggnoteseen@de@#1\endcsname\empty}
\newcommand{\ggmarknotezh}[1]{\global\expandafter\let\csname ggnoteseen@zh@#1\endcsname\empty}
\newcommand{\ggnotedeonce}[2]{\ifcsname ggnoteseen@de@#1\endcsname\else\global\expandafter\let\csname ggnoteseen@de@#1\endcsname\empty\ggnotede{#2}\fi}
\newcommand{\ggnotezhonce}[2]{\ifcsname ggnoteseen@zh@#1\endcsname\else\global\expandafter\let\csname ggnoteseen@zh@#1\endcsname\empty\ggnotezh{#2}\fi}
\newcommand{\ggnotede}[1]{\editornotede{\ggversionde #1}}
\newcommand{\ggnotezh}[1]{\editornotezh{\ggversionzh #1}}

% ===== Art. 1 =====
\newcommand{\ggfnArtOneOneDE}{\ggnotedeonce{art-1-1}{Art. 1 Abs. 1: ``Die Wuerde des Menschen ist unantastbar. Sie zu achten und zu schuetzen ist Verpflichtung aller staatlichen Gewalt.''}}
\newcommand{\ggfnArtOneOneZH}{\ggnotezhonce{art-1-1}{第1條第1項：「人之尊嚴不可侵犯。尊重及保護之，為一切國家權力之義務。」}}
\newcommand{\ggfnArtOneOneTwoDE}{\ggnotedeonce{art-1-1-2}{Art. 1 Abs. 1 Satz 2: ``Sie zu achten und zu schuetzen ist Verpflichtung aller staatlichen Gewalt.''}}
\newcommand{\ggfnArtOneOneTwoZH}{\ggnotezhonce{art-1-1-2}{第1條第1項第2句：「尊重及保護之，為一切國家權力之義務。」}}

% ===== Art. 2 =====
\newcommand{\ggfnArtTwoOneDE}{\ggnotedeonce{art-2-1}{Art. 2 Abs. 1: ``Jeder hat das Recht auf die freie Entfaltung seiner Persoenlichkeit, soweit er nicht die Rechte anderer verletzt und nicht gegen die verfassungsmaessige Ordnung oder das Sittengesetz verstoesst.''}}
\newcommand{\ggfnArtTwoOneZH}{\ggnotezhonce{art-2-1}{第2條第1項：「人人有自由發展其人格之權利，但不得侵害他人權利，亦不得違反合憲秩序或道德律。」}}
% 簡易預設版（墮胎一案以 \renewcommand 覆寫為含條件邏輯之版本）
\newcommand{\ggfnArtTwoTwoDE}{\ggnotedeonce{art-2-2}{Art. 2 Abs. 2: ``Jeder hat das Recht auf Leben und koerperliche Unversehrtheit. Die Freiheit der Person ist unverletzlich. In diese Rechte darf nur auf Grund eines Gesetzes eingegriffen werden.''}\ggmarknotede{art-2-2-1}}
\newcommand{\ggfnArtTwoTwoZH}{\ggnotezhonce{art-2-2}{第2條第2項：「人人有生命與身體完整性之權利。人身自由不可侵犯。此等權利僅得依法律加以干預。」}\ggmarknotezh{art-2-2-1}}
\newcommand{\ggfnArtTwoTwoOneDE}{\ggnotedeonce{art-2-2-1}{Art. 2 Abs. 2 Satz 1: ``Jeder hat das Recht auf Leben und koerperliche Unversehrtheit.''}}
\newcommand{\ggfnArtTwoTwoOneZH}{\ggnotezhonce{art-2-2-1}{第2條第2項第1句：「人人有生命與身體完整性之權利。」}}
\newcommand{\ggfnArtTwoTwoThreeDE}{\ggnotedeonce{art-2-2-3}{Art. 2 Abs. 2 Satz 3: ``In diese Rechte darf nur auf Grund eines Gesetzes eingegriffen werden.''}}
\newcommand{\ggfnArtTwoTwoThreeZH}{\ggnotezhonce{art-2-2-3}{第2條第2項第3句：「此等權利僅得依法律加以干預。」}}
% 墮胎一案特有：第2條第2項第2–3句（Art. 2(2) 去掉第1句後的「餘文」）
\newcommand{\ggfnArtTwoTwoRestDE}{\ggnotedeonce{art-2-2-rest}{Art. 2 Abs. 2 Saetze 2-3: ``Die Freiheit der Person ist unverletzlich. In diese Rechte darf nur auf Grund eines Gesetzes eingegriffen werden.''}\ggmarknotede{art-2-2}}
\newcommand{\ggfnArtTwoTwoRestZH}{\ggnotezhonce{art-2-2-rest}{第2條第2項第2、3句：「人身自由不可侵犯。此等權利僅得依法律加以干預。」}\ggmarknotezh{art-2-2}}

% ===== Art. 3 =====
\newcommand{\ggfnArtThreeDE}{\ggnotedeonce{art-3}{Art. 3: ``Alle Menschen sind vor dem Gesetz gleich. Maenner und Frauen sind gleichberechtigt. Niemand darf wegen seines Geschlechtes, seiner Abstammung, seiner Rasse, seiner Sprache, seiner Heimat und Herkunft, seines Glaubens, seiner religioesen oder politischen Anschauungen benachteiligt oder bevorzugt werden.''}}
\newcommand{\ggfnArtThreeZH}{\ggnotezhonce{art-3}{第3條：「所有人在法律前一律平等。男女享有平等權利。任何人不得因其性別、血統、種族、語言、故鄉及出身、信仰、宗教或政治見解而受不利益或優惠。」}}
\newcommand{\ggfnArtThreeTwoDE}{\ggnotedeonce{art-3-2}{Art. 3 Abs. 2: ``Maenner und Frauen sind gleichberechtigt.''}}
\newcommand{\ggfnArtThreeTwoZH}{\ggnotezhonce{art-3-2}{第3條第2項：「男女享有平等權利。」}}

% ===== Art. 4 =====
\newcommand{\ggfnArtFourOneDE}{\ggnotedeonce{art-4-1}{Art. 4 Abs. 1: ``Die Freiheit des Glaubens, des Gewissens und die Freiheit des religioesen und weltanschaulichen Bekenntnisses sind unverletzlich.''}}
\newcommand{\ggfnArtFourOneZH}{\ggnotezhonce{art-4-1}{第4條第1項：「信仰自由、良心自由，以及宗教與世界觀信念自由，不可侵犯。」}}

% ===== Art. 5 =====
\newcommand{\ggfnArtFiveOneDE}{\ggnotedeonce{art-5-1}{Art. 5 Abs. 1: ``Jeder hat das Recht, seine Meinung in Wort, Schrift und Bild frei zu aeussern und zu verbreiten und sich aus allgemein zugaenglichen Quellen ungehindert zu unterrichten. Die Pressefreiheit und die Freiheit der Berichterstattung durch Rundfunk und Film werden gewaehrleistet. Eine Zensur findet nicht statt.''}}
\newcommand{\ggfnArtFiveOneZH}{\ggnotezhonce{art-5-1}{第5條第1項：「人人有以言語、文字及圖像自由表達及傳播其意見，並由一般可接近來源不受阻礙獲取資訊之權利。新聞自由及廣播、電影報導自由受保障。不設審查制度。」}}

% ===== Art. 6 =====
\newcommand{\ggfnArtSixDE}{\ggnotedeonce{art-6}{Art. 6 Abs. 1: ``Ehe und Familie stehen unter dem besonderen Schutze der staatlichen Ordnung.'' Abs. 2 Satz 1: ``Pflege und Erziehung der Kinder sind das natuerliche Recht der Eltern und die zuvoerderst ihnen obliegende Pflicht.'' Abs. 4: ``Jede Mutter hat Anspruch auf den Schutz und die Fuersorge der Gemeinschaft.''}}
\newcommand{\ggfnArtSixZH}{\ggnotezhonce{art-6}{第6條第1項：「婚姻與家庭受國家秩序之特別保護。」第2項第1句：「照護及教育子女，為父母之自然權利，並為首先由其承擔之義務。」第4項：「每一母親均有請求共同體保護與照顧之權利。」}}
\newcommand{\ggfnArtSixOneDE}{\ggnotedeonce{art-6-1}{Art. 6 Abs. 1: ``Ehe und Familie stehen unter dem besonderen Schutze der staatlichen Ordnung.''}}
\newcommand{\ggfnArtSixOneZH}{\ggnotezhonce{art-6-1}{第6條第1項：「婚姻與家庭受國家秩序之特別保護。」}}
\newcommand{\ggfnArtSixTwoDE}{\ggnotedeonce{art-6-2}{Art. 6 Abs. 2: ``Pflege und Erziehung der Kinder sind das natuerliche Recht der Eltern und die zuvoerderst ihnen obliegende Pflicht. Ueber ihre Betaetigung wacht die staatliche Gemeinschaft.''}}
\newcommand{\ggfnArtSixTwoZH}{\ggnotezhonce{art-6-2}{第6條第2項：「照護及教育子女，為父母之自然權利，並為首先由其承擔之義務。國家共同體監督其行使。」}}
\newcommand{\ggfnArtSixTwoOneDE}{\ggnotedeonce{art-6-2-1}{Art. 6 Abs. 2 Satz 1: ``Pflege und Erziehung der Kinder sind das natuerliche Recht der Eltern und die zuvoerderst ihnen obliegende Pflicht.''}}
\newcommand{\ggfnArtSixTwoOneZH}{\ggnotezhonce{art-6-2-1}{第6條第2項第1句：「照護及教育子女，為父母之自然權利，並為首先由其承擔之義務。」}}
\newcommand{\ggfnArtSixFourDE}{\ggnotedeonce{art-6-4}{Art. 6 Abs. 4: ``Jede Mutter hat Anspruch auf den Schutz und die Fuersorge der Gemeinschaft.''}}
\newcommand{\ggfnArtSixFourZH}{\ggnotezhonce{art-6-4}{第6條第4項：「每一母親均有請求共同體保護與照顧之權利。」}}

% ===== Art. 12 =====
\newcommand{\ggfnArtTwelveOneDE}{\ggnotedeonce{art-12-1}{Art. 12 Abs. 1: ``Alle Deutschen haben das Recht, Beruf, Arbeitsplatz und Ausbildungsstaette frei zu waehlen. Die Berufsausuebung kann durch Gesetz oder auf Grund eines Gesetzes geregelt werden.''}}
\newcommand{\ggfnArtTwelveOneZH}{\ggnotezhonce{art-12-1}{第12條第1項：「所有德國人均有自由選擇職業、工作場所及訓練場所之權利。職業行使得以法律或基於法律加以規範。」}}

% ===== Art. 19 =====
\newcommand{\ggfnArtNineteenTwoDE}{\ggnotedeonce{art-19-2}{Art. 19 Abs. 2: ``In keinem Falle darf ein Grundrecht in seinem Wesensgehalt angetastet werden.''}}
\newcommand{\ggfnArtNineteenTwoZH}{\ggnotezhonce{art-19-2}{第19條第2項：「任何情形下，基本權利之本質內容均不得受到侵害。」}}

% ===== Art. 20 =====
\newcommand{\ggfnArtTwentyOneDE}{\ggnotedeonce{art-20-1}{Art. 20 Abs. 1: ``Die Bundesrepublik Deutschland ist ein demokratischer und sozialer Bundesstaat.''}}
\newcommand{\ggfnArtTwentyOneZH}{\ggnotezhonce{art-20-1}{第20條第1項：「德意志聯邦共和國為民主及社會之聯邦國。」}}
\newcommand{\ggfnArtTwentyThreeDE}{\ggnotedeonce{art-20-3}{Art. 20 Abs. 3: ``Die Gesetzgebung ist an die verfassungsmaessige Ordnung, die vollziehende Gewalt und die Rechtsprechung sind an Gesetz und Recht gebunden.''}}
\newcommand{\ggfnArtTwentyThreeZH}{\ggnotezhonce{art-20-3}{第20條第3項：「立法受合憲秩序拘束；行政權與司法權受法律與法拘束。」}}

% ===== Art. 26 + Art. 143（墮胎一案特有）=====
\newcommand{\ggfnArtTwoSixAndOneFourThreeDE}{\ggnotedeonce{art-26-143}{Art. 26 Abs. 1 Satz 2: ``Sie sind unter Strafe zu stellen.'' Art. 143 in der urspruenglichen Fassung (24. Mai 1949 bis 1. September 1951) enthielt Strafvorschriften zum Hochverrat; 1975 war Art. 143 bereits weggefallen.}}
\newcommand{\ggfnArtTwoSixAndOneFourThreeZH}{\ggnotezhonce{art-26-143}{第26條第1項第2句：「此等行為應規定為犯罪處罰。」第143條原始版本（1949年5月24日至1951年9月1日）含有關於叛國罪之刑罰規定；至1975年時，第143條已廢止。}}

% ===== Art. 28 =====
\newcommand{\ggfnArtTwentyEightOneDE}{\ggnotedeonce{art-28-1-1}{Art. 28 Abs. 1 Satz 1: ``Die verfassungsmaessige Ordnung in den Laendern muss den Grundsaetzen des republikanischen, demokratischen und sozialen Rechtsstaates im Sinne dieses Grundgesetzes entsprechen.''}}
\newcommand{\ggfnArtTwentyEightOneZH}{\ggnotezhonce{art-28-1-1}{第28條第1項第1句：「各邦之合憲秩序，須符合本基本法意義下共和、民主及社會法治國之原則。」}}

% ===== Art. 30 =====
\newcommand{\ggfnArtThirtyDE}{\ggnotedeonce{art-30}{Art. 30: ``Die Ausuebung der staatlichen Befugnisse und die Erfuellung der staatlichen Aufgaben ist Sache der Laender, soweit dieses Grundgesetz keine andere Regelung trifft oder zulaesst.''}}
\newcommand{\ggfnArtThirtyZH}{\ggnotezhonce{art-30}{第30條：「國家權限之行使與國家任務之履行，屬於各邦，但本基本法另有規定或容許者，不在此限。」}}

% ===== Art. 74 =====
\newcommand{\ggfnArtSeventyFourOneDE}{\ggnotedeonce{art-74-1}{Art. 74 Nr. 1: Die konkurrierende Gesetzgebung erstreckt sich auf ``das buergerliche Recht, das Strafrecht und den Strafvollzug, die Gerichtsverfassung, das gerichtliche Verfahren, die Rechtsanwaltschaft, das Notariat und die Rechtsberatung''.}}
\newcommand{\ggfnArtSeventyFourOneZH}{\ggnotezhonce{art-74-1}{第74條第1款：競合立法及於「民法、刑法及刑罰執行、法院組織、訴訟程序、律師制度、公證制度及法律諮詢」。}}
\newcommand{\ggfnArtSeventyFourSevenDE}{\ggnotedeonce{art-74-7}{Art. 74 Nr. 7: Die konkurrierende Gesetzgebung erstreckt sich auf ``die oeffentliche Fuersorge''.}}
\newcommand{\ggfnArtSeventyFourSevenZH}{\ggnotezhonce{art-74-7}{第74條第7款：競合立法及於「公共扶助」。}}
\newcommand{\ggfnArtSeventyFourTwelveDE}{\ggnotedeonce{art-74-12}{Art. 74 Nr. 12: Die konkurrierende Gesetzgebung erstreckt sich auf ``das Arbeitsrecht einschliesslich der Betriebsverfassung, des Arbeitsschutzes und der Arbeitsvermittlung sowie die Sozialversicherung einschliesslich der Arbeitslosenversicherung''.}}
\newcommand{\ggfnArtSeventyFourTwelveZH}{\ggnotezhonce{art-74-12}{第74條第12款：競合立法及於「勞動法，包括企業組織、勞動保護及就業媒合，以及社會保險，包括失業保險」。}}

% ===== Art. 77（墮胎一案特有）=====
\newcommand{\ggfnArtSevenSevenThreeDE}{\ggnotedeonce{art-77-3}{Art. 77 Abs. 3 Satz 1: ``Soweit zu einem Gesetze die Zustimmung des Bundesrates nicht erforderlich ist, kann der Bundesrat, wenn das Verfahren nach Absatz 2 beendigt ist, gegen ein vom Bundestage beschlossenes Gesetz binnen zwei Wochen Einspruch einlegen.''}}
\newcommand{\ggfnArtSevenSevenThreeZH}{\ggnotezhonce{art-77-3}{第77條第3項第1句：「法律無須聯邦參議院同意者，於第2項程序終了後，聯邦參議院得於二週內對聯邦議院通過之法律提出異議。」}}

% ===== Art. 80 =====
\newcommand{\ggfnArtEightyOneOneDE}{\ggnotedeonce{art-80-1-1}{Art. 80 Abs. 1 Satz 1: ``Durch Gesetz koennen die Bundesregierung, ein Bundesminister oder die Landesregierungen ermaechtigt werden, Rechtsverordnungen zu erlassen.''}}
\newcommand{\ggfnArtEightyOneOneZH}{\ggnotezhonce{art-80-1-1}{第80條第1項第1句：「得以法律授權聯邦政府、聯邦部長或邦政府發布法規命令。」}}

% ===== Art. 83 =====
\newcommand{\ggfnArtEightyThreeDE}{\ggnotedeonce{art-83}{Art. 83: ``Die Laender fuehren die Bundesgesetze als eigene Angelegenheit aus, soweit dieses Grundgesetz nichts anderes bestimmt oder zulaesst.''}}
\newcommand{\ggfnArtEightyThreeZH}{\ggnotezhonce{art-83}{第83條：「各邦以自己事務執行聯邦法律，但本基本法另有規定或容許者，不在此限。」}}

% ===== Art. 84 =====
\newcommand{\ggfnArtEightyFourOneDE}{\ggnotedeonce{art-84-1}{Art. 84 Abs. 1: ``Fuehren die Laender die Bundesgesetze als eigene Angelegenheit aus, so regeln sie die Einrichtung der Behoerden und das Verwaltungsverfahren, soweit nicht Bundesgesetze mit Zustimmung des Bundesrates etwas anderes bestimmen.''}}
\newcommand{\ggfnArtEightyFourOneZH}{\ggnotezhonce{art-84-1}{第84條第1項：「各邦以自己事務執行聯邦法律時，除聯邦法律經聯邦參議院同意另有規定外，由各邦規範機關之設置及行政程序。」}}
% 別名（墮胎一案使用之縮寫命名）
\let\ggfnArtEightFourOneDE\ggfnArtEightyFourOneDE
\let\ggfnArtEightFourOneZH\ggfnArtEightyFourOneZH

% ===== Art. 93 =====
\newcommand{\ggfnArtNinetyThreeOneTwoDE}{\ggnotedeonce{art-93-1-2}{Art. 93 Abs. 1 Nr. 2: Das Bundesverfassungsgericht entscheidet ``bei Meinungsverschiedenheiten oder Zweifeln ueber die foermliche und sachliche Vereinbarkeit von Bundesrecht oder Landesrecht mit diesem Grundgesetze ... auf Antrag der Bundesregierung, einer Landesregierung oder eines Drittels der Mitglieder des Bundestages''.}}
\newcommand{\ggfnArtNinetyThreeOneTwoZH}{\ggnotezhonce{art-93-1-2}{第93條第1項第2款：聯邦憲法法院就「聯邦法或邦法在形式及實質上是否與本基本法相容之意見分歧或疑義……依聯邦政府、邦政府或聯邦議院三分之一議員之聲請」作成裁判。}}
% 別名（墮胎一案使用之縮寫命名）
\let\ggfnArtNineThreeOneTwoDE\ggfnArtNinetyThreeOneTwoDE
\let\ggfnArtNineThreeOneTwoZH\ggfnArtNinetyThreeOneTwoZH

% ===== Art. 102 =====
\newcommand{\ggfnArtOneZeroTwoDE}{\ggnotedeonce{art-102}{Art. 102: ``Die Todesstrafe ist abgeschafft.''}}
\newcommand{\ggfnArtOneZeroTwoZH}{\ggnotezhonce{art-102}{第102條：「死刑廢止。」}}

% ===== Art. 103 =====
\newcommand{\ggfnArtOneZeroThreeTwoDE}{\ggnotedeonce{art-103-2}{Art. 103 Abs. 2: ``Eine Tat kann nur bestraft werden, wenn die Strafbarkeit gesetzlich bestimmt war, bevor die Tat begangen wurde.''}}
\newcommand{\ggfnArtOneZeroThreeTwoZH}{\ggnotezhonce{art-103-2}{第103條第2項：「行為之可罰性，須於行為前以法律定之，始得處罰。」}}

% ===== Art. 104 =====
\newcommand{\ggfnArtOneZeroFourOneDE}{\ggnotedeonce{art-104-1}{Art. 104 Abs. 1: ``Die Freiheit der Person kann nur auf Grund eines foermlichen Gesetzes und nur unter Beachtung der darin vorgeschriebenen Formen beschraenkt werden.''}}
\newcommand{\ggfnArtOneZeroFourOneZH}{\ggnotezhonce{art-104-1}{第104條第1項：「人身自由僅得基於形式法律，並遵守其中規定之形式加以限制。」}}

% ===== 組合腳註：Art. 3 + Art. 6（墮胎一案特有）=====
\newcommand{\ggfnArtThreeSixDE}{\ggnotedeonce{art-3-6}{Art. 3: ``Alle Menschen sind vor dem Gesetz gleich. Maenner und Frauen sind gleichberechtigt. Niemand darf wegen seines Geschlechtes, seiner Abstammung, seiner Rasse, seiner Sprache, seiner Heimat und Herkunft, seines Glaubens, seiner religioesen oder politischen Anschauungen benachteiligt oder bevorzugt werden.'' Art. 6 Abs. 1: ``Ehe und Familie stehen unter dem besonderen Schutze der staatlichen Ordnung.'' Abs. 2 Satz 1: ``Pflege und Erziehung der Kinder sind das natuerliche Recht der Eltern und die zuvoerderst ihnen obliegende Pflicht.'' Abs. 4: ``Jede Mutter hat Anspruch auf den Schutz und die Fuersorge der Gemeinschaft.''}\ggmarknotede{art-3}\ggmarknotede{art-6}\ggmarknotede{art-6-1-4}}
\newcommand{\ggfnArtThreeSixZH}{\ggnotezhonce{art-3-6}{第3條：「所有人在法律前一律平等。男女享有平等權利。任何人不得因其性別、血統、種族、語言、故鄉及出身、信仰、宗教或政治見解而受不利益或優惠。」第6條第1項：「婚姻與家庭受國家秩序之特別保護。」第2項第1句：「照護及教育子女，為父母之自然權利，並為首先由其承擔之義務。」第4項：「每一母親均有請求共同體保護與照顧之權利。」}\ggmarknotezh{art-3}\ggmarknotezh{art-6}\ggmarknotezh{art-6-1-4}}

% ===== 組合腳註：Art. 6(1) + Art. 6(4)（墮胎一案特有）=====
\newcommand{\ggfnArtSixOneFourDE}{\ggnotedeonce{art-6-1-4}{Art. 6 Abs. 1: ``Ehe und Familie stehen unter dem besonderen Schutze der staatlichen Ordnung.'' Art. 6 Abs. 4: ``Jede Mutter hat Anspruch auf den Schutz und die Fuersorge der Gemeinschaft.''}\ggmarknotede{art-6}}
\newcommand{\ggfnArtSixOneFourZH}{\ggnotezhonce{art-6-1-4}{第6條第1項：「婚姻與家庭受國家秩序之特別保護。」第4項：「每一母親均有請求共同體保護與照顧之權利。」}\ggmarknotezh{art-6}}

% ===== 組合腳註：Art. 1(1) + Art. 2(2)(1)（墮胎二案特有）=====
\newcommand{\ggfnArtOneOneTwoTwoOneDE}{\ggnotedeonce{art-1-1+2-2-1}{Art. 1 Abs. 1: ``Die Wuerde des Menschen ist unantastbar. Sie zu achten und zu schuetzen ist Verpflichtung aller staatlichen Gewalt.'' Art. 2 Abs. 2 Satz 1: ``Jeder hat das Recht auf Leben und koerperliche Unversehrtheit.''}\ggmarknotede{art-1-1}\ggmarknotede{art-2-2-1}}
\newcommand{\ggfnArtOneOneTwoTwoOneZH}{\ggnotezhonce{art-1-1+2-2-1}{第1條第1項：「人之尊嚴不可侵犯。尊重及保護之，為一切國家權力之義務。」第2條第2項第1句：「人人有生命與身體完整性之權利。」}\ggmarknotezh{art-1-1}\ggmarknotezh{art-2-2-1}}

% ===== 組合腳註：Art. 1(1) + Art. 2(2)（墮胎二案特有）=====
\newcommand{\ggfnArtOneOneTwoTwoDE}{\ggnotedeonce{art-1-1+2-2}{Art. 1 Abs. 1: ``Die Wuerde des Menschen ist unantastbar. Sie zu achten und zu schuetzen ist Verpflichtung aller staatlichen Gewalt.'' Art. 2 Abs. 2: ``Jeder hat das Recht auf Leben und koerperliche Unversehrtheit. Die Freiheit der Person ist unverletzlich. In diese Rechte darf nur auf Grund eines Gesetzes eingegriffen werden.''}\ggmarknotede{art-1-1}\ggmarknotede{art-2-2}\ggmarknotede{art-2-2-1}}
\newcommand{\ggfnArtOneOneTwoTwoZH}{\ggnotezhonce{art-1-1+2-2}{第1條第1項：「人之尊嚴不可侵犯。尊重及保護之，為一切國家權力之義務。」第2條第2項：「人人有生命與身體完整性之權利。人身自由不可侵犯。此等權利僅得依法律加以干預。」}\ggmarknotezh{art-1-1}\ggmarknotezh{art-2-2}\ggmarknotezh{art-2-2-1}}

% ===== 組合腳註：Art. 1(1) + Art. 2(1)（墮胎二案特有）=====
\newcommand{\ggfnArtOneOneTwoOneDE}{\ggnotedeonce{art-1-1+2-1}{Art. 1 Abs. 1: ``Die Wuerde des Menschen ist unantastbar. Sie zu achten und zu schuetzen ist Verpflichtung aller staatlichen Gewalt.'' Art. 2 Abs. 1: ``Jeder hat das Recht auf die freie Entfaltung seiner Persoenlichkeit, soweit er nicht die Rechte anderer verletzt und nicht gegen die verfassungsmaessige Ordnung oder das Sittengesetz verstoesst.''}\ggmarknotede{art-1-1}\ggmarknotede{art-2-1}}
\newcommand{\ggfnArtOneOneTwoOneZH}{\ggnotezhonce{art-1-1+2-1}{第1條第1項：「人之尊嚴不可侵犯。尊重及保護之，為一切國家權力之義務。」第2條第1項：「人人有自由發展其人格之權利，但不得侵害他人權利，亦不得違反合憲秩序或道德律。」}\ggmarknotezh{art-1-1}\ggmarknotezh{art-2-1}}

% ===== 組合腳註：Art. 2(2)(1) + Art. 1(1)（墮胎一案特有，含交叉標記）=====
\newcommand{\ggfnLifeDignityDE}{\ggnotedeonce{life-dignity}{Art. 2 Abs. 2 Satz 1: ``Jeder hat das Recht auf Leben und koerperliche Unversehrtheit.'' Art. 1 Abs. 1: ``Die Wuerde des Menschen ist unantastbar. Sie zu achten und zu schuetzen ist Verpflichtung aller staatlichen Gewalt.''}\ggmarknotede{art-2-2-1}\ggmarknotede{art-1-1}\ggmarknotede{art-1-1-2}}
\newcommand{\ggfnLifeDignityZH}{\ggnotezhonce{life-dignity}{第2條第2項第1句：「人人有生命與身體完整性之權利。」第1條第1項：「人之尊嚴不可侵犯。尊重及保護之，為一切國家權力之義務。」}\ggmarknotezh{art-2-2-1}\ggmarknotezh{art-1-1}\ggmarknotezh{art-1-1-2}}

% ===== 組合腳註：Art. 1(1) + Art. 2(2)（墮胎一案特有，條件判斷版）=====
\newcommand{\ggfnArtOneTwoTwoDE}{\ifcsname ggnoteseen@de@life-dignity\endcsname\ggfnArtTwoTwoRestDE{}\else\ggnotedeonce{art-1-2-2}{Art. 1 Abs. 1: ``Die Wuerde des Menschen ist unantastbar. Sie zu achten und zu schuetzen ist Verpflichtung aller staatlichen Gewalt.'' Art. 2 Abs. 2: ``Jeder hat das Recht auf Leben und koerperliche Unversehrtheit. Die Freiheit der Person ist unverletzlich. In diese Rechte darf nur auf Grund eines Gesetzes eingegriffen werden.''}\ggmarknotede{art-1-1}\ggmarknotede{art-1-1-2}\ggmarknotede{art-2-2}\ggmarknotede{art-2-2-1}\fi}
\newcommand{\ggfnArtOneTwoTwoZH}{\ifcsname ggnoteseen@zh@life-dignity\endcsname\ggfnArtTwoTwoRestZH{}\else\ggnotezhonce{art-1-2-2}{第1條第1項：「人之尊嚴不可侵犯。尊重及保護之，為一切國家權力之義務。」第2條第2項：「人人有生命與身體完整性之權利。人身自由不可侵犯。此等權利僅得依法律加以干預。」}\ggmarknotezh{art-1-1}\ggmarknotezh{art-1-1-2}\ggmarknotezh{art-2-2}\ggmarknotezh{art-2-2-1}\fi}

% ===== 組合腳註：Art. 1(1) + Art. 19(2)（墮胎一案特有，條件判斷版）=====
\newcommand{\ggfnArtOneNineteenDE}{\ifcsname ggnoteseen@de@art-1-1\endcsname\ggfnArtNineteenTwoDE{}\else\ggnotedeonce{art-1-19}{Art. 1 Abs. 1: ``Die Wuerde des Menschen ist unantastbar. Sie zu achten und zu schuetzen ist Verpflichtung aller staatlichen Gewalt.'' Art. 19 Abs. 2: ``In keinem Falle darf ein Grundrecht in seinem Wesensgehalt angetastet werden.''}\ggmarknotede{art-1-1}\ggmarknotede{art-1-1-2}\ggmarknotede{art-19-2}\fi}
\newcommand{\ggfnArtOneNineteenZH}{\ifcsname ggnoteseen@zh@art-1-1\endcsname\ggfnArtNineteenTwoZH{}\else\ggnotezhonce{art-1-19}{第1條第1項：「人之尊嚴不可侵犯。尊重及保護之，為一切國家權力之義務。」第19條第2項：「任何情形下，基本權利之本質內容均不得受到侵害。」}\ggmarknotezh{art-1-1}\ggmarknotezh{art-1-1-2}\ggmarknotezh{art-19-2}\fi}
 之前定義）=====
\newif\ifggversionnotede
\newif\ifggversionnotezh

% ===== 編者註腳基礎設施（首次標記模式）=====
\newif\ifeditornotelabelde
\newif\ifeditornotelabelzh
\newcommand{\editornotelabelde}{\ifeditornotelabelde\else\emph{Hrsg.-Anm. (nachfolgend ohne Kennzeichnung)}: \global\editornotelabeldetrue\fi}
\newcommand{\editornotelabelzh}{\ifeditornotelabelzh\else 編者註（下同）：\global\editornotelabelzhtrue\fi}
\newcommand{\editornotede}[1]{\ifshoweditornotes\footnote{\normalfont\footnotesize\editornotelabelde #1}\else\ifclassroommode\footnote{\normalfont\footnotesize\editornotelabelde #1}\fi\fi}
\newcommand{\editornotezh}[1]{\ifshoweditornotes\footnote{\normalfont\footnotesize\editornotelabelzh #1}\else\ifclassroommode\footnote{\normalfont\footnotesize\editornotelabelzh #1}\fi\fi}

% ===== GG 引用系統（\ggversionde/zh 由各案在 \input 前定義）=====
\newcommand{\ggmarknotede}[1]{\global\expandafter\let\csname ggnoteseen@de@#1\endcsname\empty}
\newcommand{\ggmarknotezh}[1]{\global\expandafter\let\csname ggnoteseen@zh@#1\endcsname\empty}
\newcommand{\ggnotedeonce}[2]{\ifcsname ggnoteseen@de@#1\endcsname\else\global\expandafter\let\csname ggnoteseen@de@#1\endcsname\empty\ggnotede{#2}\fi}
\newcommand{\ggnotezhonce}[2]{\ifcsname ggnoteseen@zh@#1\endcsname\else\global\expandafter\let\csname ggnoteseen@zh@#1\endcsname\empty\ggnotezh{#2}\fi}
\newcommand{\ggnotede}[1]{\editornotede{\ggversionde #1}}
\newcommand{\ggnotezh}[1]{\editornotezh{\ggversionzh #1}}

% ===== Art. 1 =====
\newcommand{\ggfnArtOneOneDE}{\ggnotedeonce{art-1-1}{Art. 1 Abs. 1: ``Die Wuerde des Menschen ist unantastbar. Sie zu achten und zu schuetzen ist Verpflichtung aller staatlichen Gewalt.''}}
\newcommand{\ggfnArtOneOneZH}{\ggnotezhonce{art-1-1}{第1條第1項：「人之尊嚴不可侵犯。尊重及保護之，為一切國家權力之義務。」}}
\newcommand{\ggfnArtOneOneTwoDE}{\ggnotedeonce{art-1-1-2}{Art. 1 Abs. 1 Satz 2: ``Sie zu achten und zu schuetzen ist Verpflichtung aller staatlichen Gewalt.''}}
\newcommand{\ggfnArtOneOneTwoZH}{\ggnotezhonce{art-1-1-2}{第1條第1項第2句：「尊重及保護之，為一切國家權力之義務。」}}

% ===== Art. 2 =====
\newcommand{\ggfnArtTwoOneDE}{\ggnotedeonce{art-2-1}{Art. 2 Abs. 1: ``Jeder hat das Recht auf die freie Entfaltung seiner Persoenlichkeit, soweit er nicht die Rechte anderer verletzt und nicht gegen die verfassungsmaessige Ordnung oder das Sittengesetz verstoesst.''}}
\newcommand{\ggfnArtTwoOneZH}{\ggnotezhonce{art-2-1}{第2條第1項：「人人有自由發展其人格之權利，但不得侵害他人權利，亦不得違反合憲秩序或道德律。」}}
% 簡易預設版（墮胎一案以 \renewcommand 覆寫為含條件邏輯之版本）
\newcommand{\ggfnArtTwoTwoDE}{\ggnotedeonce{art-2-2}{Art. 2 Abs. 2: ``Jeder hat das Recht auf Leben und koerperliche Unversehrtheit. Die Freiheit der Person ist unverletzlich. In diese Rechte darf nur auf Grund eines Gesetzes eingegriffen werden.''}\ggmarknotede{art-2-2-1}}
\newcommand{\ggfnArtTwoTwoZH}{\ggnotezhonce{art-2-2}{第2條第2項：「人人有生命與身體完整性之權利。人身自由不可侵犯。此等權利僅得依法律加以干預。」}\ggmarknotezh{art-2-2-1}}
\newcommand{\ggfnArtTwoTwoOneDE}{\ggnotedeonce{art-2-2-1}{Art. 2 Abs. 2 Satz 1: ``Jeder hat das Recht auf Leben und koerperliche Unversehrtheit.''}}
\newcommand{\ggfnArtTwoTwoOneZH}{\ggnotezhonce{art-2-2-1}{第2條第2項第1句：「人人有生命與身體完整性之權利。」}}
\newcommand{\ggfnArtTwoTwoThreeDE}{\ggnotedeonce{art-2-2-3}{Art. 2 Abs. 2 Satz 3: ``In diese Rechte darf nur auf Grund eines Gesetzes eingegriffen werden.''}}
\newcommand{\ggfnArtTwoTwoThreeZH}{\ggnotezhonce{art-2-2-3}{第2條第2項第3句：「此等權利僅得依法律加以干預。」}}
% 墮胎一案特有：第2條第2項第2–3句（Art. 2(2) 去掉第1句後的「餘文」）
\newcommand{\ggfnArtTwoTwoRestDE}{\ggnotedeonce{art-2-2-rest}{Art. 2 Abs. 2 Saetze 2-3: ``Die Freiheit der Person ist unverletzlich. In diese Rechte darf nur auf Grund eines Gesetzes eingegriffen werden.''}\ggmarknotede{art-2-2}}
\newcommand{\ggfnArtTwoTwoRestZH}{\ggnotezhonce{art-2-2-rest}{第2條第2項第2、3句：「人身自由不可侵犯。此等權利僅得依法律加以干預。」}\ggmarknotezh{art-2-2}}

% ===== Art. 3 =====
\newcommand{\ggfnArtThreeDE}{\ggnotedeonce{art-3}{Art. 3: ``Alle Menschen sind vor dem Gesetz gleich. Maenner und Frauen sind gleichberechtigt. Niemand darf wegen seines Geschlechtes, seiner Abstammung, seiner Rasse, seiner Sprache, seiner Heimat und Herkunft, seines Glaubens, seiner religioesen oder politischen Anschauungen benachteiligt oder bevorzugt werden.''}}
\newcommand{\ggfnArtThreeZH}{\ggnotezhonce{art-3}{第3條：「所有人在法律前一律平等。男女享有平等權利。任何人不得因其性別、血統、種族、語言、故鄉及出身、信仰、宗教或政治見解而受不利益或優惠。」}}
\newcommand{\ggfnArtThreeTwoDE}{\ggnotedeonce{art-3-2}{Art. 3 Abs. 2: ``Maenner und Frauen sind gleichberechtigt.''}}
\newcommand{\ggfnArtThreeTwoZH}{\ggnotezhonce{art-3-2}{第3條第2項：「男女享有平等權利。」}}

% ===== Art. 4 =====
\newcommand{\ggfnArtFourOneDE}{\ggnotedeonce{art-4-1}{Art. 4 Abs. 1: ``Die Freiheit des Glaubens, des Gewissens und die Freiheit des religioesen und weltanschaulichen Bekenntnisses sind unverletzlich.''}}
\newcommand{\ggfnArtFourOneZH}{\ggnotezhonce{art-4-1}{第4條第1項：「信仰自由、良心自由，以及宗教與世界觀信念自由，不可侵犯。」}}

% ===== Art. 5 =====
\newcommand{\ggfnArtFiveOneDE}{\ggnotedeonce{art-5-1}{Art. 5 Abs. 1: ``Jeder hat das Recht, seine Meinung in Wort, Schrift und Bild frei zu aeussern und zu verbreiten und sich aus allgemein zugaenglichen Quellen ungehindert zu unterrichten. Die Pressefreiheit und die Freiheit der Berichterstattung durch Rundfunk und Film werden gewaehrleistet. Eine Zensur findet nicht statt.''}}
\newcommand{\ggfnArtFiveOneZH}{\ggnotezhonce{art-5-1}{第5條第1項：「人人有以言語、文字及圖像自由表達及傳播其意見，並由一般可接近來源不受阻礙獲取資訊之權利。新聞自由及廣播、電影報導自由受保障。不設審查制度。」}}

% ===== Art. 6 =====
\newcommand{\ggfnArtSixDE}{\ggnotedeonce{art-6}{Art. 6 Abs. 1: ``Ehe und Familie stehen unter dem besonderen Schutze der staatlichen Ordnung.'' Abs. 2 Satz 1: ``Pflege und Erziehung der Kinder sind das natuerliche Recht der Eltern und die zuvoerderst ihnen obliegende Pflicht.'' Abs. 4: ``Jede Mutter hat Anspruch auf den Schutz und die Fuersorge der Gemeinschaft.''}}
\newcommand{\ggfnArtSixZH}{\ggnotezhonce{art-6}{第6條第1項：「婚姻與家庭受國家秩序之特別保護。」第2項第1句：「照護及教育子女，為父母之自然權利，並為首先由其承擔之義務。」第4項：「每一母親均有請求共同體保護與照顧之權利。」}}
\newcommand{\ggfnArtSixOneDE}{\ggnotedeonce{art-6-1}{Art. 6 Abs. 1: ``Ehe und Familie stehen unter dem besonderen Schutze der staatlichen Ordnung.''}}
\newcommand{\ggfnArtSixOneZH}{\ggnotezhonce{art-6-1}{第6條第1項：「婚姻與家庭受國家秩序之特別保護。」}}
\newcommand{\ggfnArtSixTwoDE}{\ggnotedeonce{art-6-2}{Art. 6 Abs. 2: ``Pflege und Erziehung der Kinder sind das natuerliche Recht der Eltern und die zuvoerderst ihnen obliegende Pflicht. Ueber ihre Betaetigung wacht die staatliche Gemeinschaft.''}}
\newcommand{\ggfnArtSixTwoZH}{\ggnotezhonce{art-6-2}{第6條第2項：「照護及教育子女，為父母之自然權利，並為首先由其承擔之義務。國家共同體監督其行使。」}}
\newcommand{\ggfnArtSixTwoOneDE}{\ggnotedeonce{art-6-2-1}{Art. 6 Abs. 2 Satz 1: ``Pflege und Erziehung der Kinder sind das natuerliche Recht der Eltern und die zuvoerderst ihnen obliegende Pflicht.''}}
\newcommand{\ggfnArtSixTwoOneZH}{\ggnotezhonce{art-6-2-1}{第6條第2項第1句：「照護及教育子女，為父母之自然權利，並為首先由其承擔之義務。」}}
\newcommand{\ggfnArtSixFourDE}{\ggnotedeonce{art-6-4}{Art. 6 Abs. 4: ``Jede Mutter hat Anspruch auf den Schutz und die Fuersorge der Gemeinschaft.''}}
\newcommand{\ggfnArtSixFourZH}{\ggnotezhonce{art-6-4}{第6條第4項：「每一母親均有請求共同體保護與照顧之權利。」}}

% ===== Art. 12 =====
\newcommand{\ggfnArtTwelveOneDE}{\ggnotedeonce{art-12-1}{Art. 12 Abs. 1: ``Alle Deutschen haben das Recht, Beruf, Arbeitsplatz und Ausbildungsstaette frei zu waehlen. Die Berufsausuebung kann durch Gesetz oder auf Grund eines Gesetzes geregelt werden.''}}
\newcommand{\ggfnArtTwelveOneZH}{\ggnotezhonce{art-12-1}{第12條第1項：「所有德國人均有自由選擇職業、工作場所及訓練場所之權利。職業行使得以法律或基於法律加以規範。」}}

% ===== Art. 19 =====
\newcommand{\ggfnArtNineteenTwoDE}{\ggnotedeonce{art-19-2}{Art. 19 Abs. 2: ``In keinem Falle darf ein Grundrecht in seinem Wesensgehalt angetastet werden.''}}
\newcommand{\ggfnArtNineteenTwoZH}{\ggnotezhonce{art-19-2}{第19條第2項：「任何情形下，基本權利之本質內容均不得受到侵害。」}}

% ===== Art. 20 =====
\newcommand{\ggfnArtTwentyOneDE}{\ggnotedeonce{art-20-1}{Art. 20 Abs. 1: ``Die Bundesrepublik Deutschland ist ein demokratischer und sozialer Bundesstaat.''}}
\newcommand{\ggfnArtTwentyOneZH}{\ggnotezhonce{art-20-1}{第20條第1項：「德意志聯邦共和國為民主及社會之聯邦國。」}}
\newcommand{\ggfnArtTwentyThreeDE}{\ggnotedeonce{art-20-3}{Art. 20 Abs. 3: ``Die Gesetzgebung ist an die verfassungsmaessige Ordnung, die vollziehende Gewalt und die Rechtsprechung sind an Gesetz und Recht gebunden.''}}
\newcommand{\ggfnArtTwentyThreeZH}{\ggnotezhonce{art-20-3}{第20條第3項：「立法受合憲秩序拘束；行政權與司法權受法律與法拘束。」}}

% ===== Art. 26 + Art. 143（墮胎一案特有）=====
\newcommand{\ggfnArtTwoSixAndOneFourThreeDE}{\ggnotedeonce{art-26-143}{Art. 26 Abs. 1 Satz 2: ``Sie sind unter Strafe zu stellen.'' Art. 143 in der urspruenglichen Fassung (24. Mai 1949 bis 1. September 1951) enthielt Strafvorschriften zum Hochverrat; 1975 war Art. 143 bereits weggefallen.}}
\newcommand{\ggfnArtTwoSixAndOneFourThreeZH}{\ggnotezhonce{art-26-143}{第26條第1項第2句：「此等行為應規定為犯罪處罰。」第143條原始版本（1949年5月24日至1951年9月1日）含有關於叛國罪之刑罰規定；至1975年時，第143條已廢止。}}

% ===== Art. 28 =====
\newcommand{\ggfnArtTwentyEightOneDE}{\ggnotedeonce{art-28-1-1}{Art. 28 Abs. 1 Satz 1: ``Die verfassungsmaessige Ordnung in den Laendern muss den Grundsaetzen des republikanischen, demokratischen und sozialen Rechtsstaates im Sinne dieses Grundgesetzes entsprechen.''}}
\newcommand{\ggfnArtTwentyEightOneZH}{\ggnotezhonce{art-28-1-1}{第28條第1項第1句：「各邦之合憲秩序，須符合本基本法意義下共和、民主及社會法治國之原則。」}}

% ===== Art. 30 =====
\newcommand{\ggfnArtThirtyDE}{\ggnotedeonce{art-30}{Art. 30: ``Die Ausuebung der staatlichen Befugnisse und die Erfuellung der staatlichen Aufgaben ist Sache der Laender, soweit dieses Grundgesetz keine andere Regelung trifft oder zulaesst.''}}
\newcommand{\ggfnArtThirtyZH}{\ggnotezhonce{art-30}{第30條：「國家權限之行使與國家任務之履行，屬於各邦，但本基本法另有規定或容許者，不在此限。」}}

% ===== Art. 74 =====
\newcommand{\ggfnArtSeventyFourOneDE}{\ggnotedeonce{art-74-1}{Art. 74 Nr. 1: Die konkurrierende Gesetzgebung erstreckt sich auf ``das buergerliche Recht, das Strafrecht und den Strafvollzug, die Gerichtsverfassung, das gerichtliche Verfahren, die Rechtsanwaltschaft, das Notariat und die Rechtsberatung''.}}
\newcommand{\ggfnArtSeventyFourOneZH}{\ggnotezhonce{art-74-1}{第74條第1款：競合立法及於「民法、刑法及刑罰執行、法院組織、訴訟程序、律師制度、公證制度及法律諮詢」。}}
\newcommand{\ggfnArtSeventyFourSevenDE}{\ggnotedeonce{art-74-7}{Art. 74 Nr. 7: Die konkurrierende Gesetzgebung erstreckt sich auf ``die oeffentliche Fuersorge''.}}
\newcommand{\ggfnArtSeventyFourSevenZH}{\ggnotezhonce{art-74-7}{第74條第7款：競合立法及於「公共扶助」。}}
\newcommand{\ggfnArtSeventyFourTwelveDE}{\ggnotedeonce{art-74-12}{Art. 74 Nr. 12: Die konkurrierende Gesetzgebung erstreckt sich auf ``das Arbeitsrecht einschliesslich der Betriebsverfassung, des Arbeitsschutzes und der Arbeitsvermittlung sowie die Sozialversicherung einschliesslich der Arbeitslosenversicherung''.}}
\newcommand{\ggfnArtSeventyFourTwelveZH}{\ggnotezhonce{art-74-12}{第74條第12款：競合立法及於「勞動法，包括企業組織、勞動保護及就業媒合，以及社會保險，包括失業保險」。}}

% ===== Art. 77（墮胎一案特有）=====
\newcommand{\ggfnArtSevenSevenThreeDE}{\ggnotedeonce{art-77-3}{Art. 77 Abs. 3 Satz 1: ``Soweit zu einem Gesetze die Zustimmung des Bundesrates nicht erforderlich ist, kann der Bundesrat, wenn das Verfahren nach Absatz 2 beendigt ist, gegen ein vom Bundestage beschlossenes Gesetz binnen zwei Wochen Einspruch einlegen.''}}
\newcommand{\ggfnArtSevenSevenThreeZH}{\ggnotezhonce{art-77-3}{第77條第3項第1句：「法律無須聯邦參議院同意者，於第2項程序終了後，聯邦參議院得於二週內對聯邦議院通過之法律提出異議。」}}

% ===== Art. 80 =====
\newcommand{\ggfnArtEightyOneOneDE}{\ggnotedeonce{art-80-1-1}{Art. 80 Abs. 1 Satz 1: ``Durch Gesetz koennen die Bundesregierung, ein Bundesminister oder die Landesregierungen ermaechtigt werden, Rechtsverordnungen zu erlassen.''}}
\newcommand{\ggfnArtEightyOneOneZH}{\ggnotezhonce{art-80-1-1}{第80條第1項第1句：「得以法律授權聯邦政府、聯邦部長或邦政府發布法規命令。」}}

% ===== Art. 83 =====
\newcommand{\ggfnArtEightyThreeDE}{\ggnotedeonce{art-83}{Art. 83: ``Die Laender fuehren die Bundesgesetze als eigene Angelegenheit aus, soweit dieses Grundgesetz nichts anderes bestimmt oder zulaesst.''}}
\newcommand{\ggfnArtEightyThreeZH}{\ggnotezhonce{art-83}{第83條：「各邦以自己事務執行聯邦法律，但本基本法另有規定或容許者，不在此限。」}}

% ===== Art. 84 =====
\newcommand{\ggfnArtEightyFourOneDE}{\ggnotedeonce{art-84-1}{Art. 84 Abs. 1: ``Fuehren die Laender die Bundesgesetze als eigene Angelegenheit aus, so regeln sie die Einrichtung der Behoerden und das Verwaltungsverfahren, soweit nicht Bundesgesetze mit Zustimmung des Bundesrates etwas anderes bestimmen.''}}
\newcommand{\ggfnArtEightyFourOneZH}{\ggnotezhonce{art-84-1}{第84條第1項：「各邦以自己事務執行聯邦法律時，除聯邦法律經聯邦參議院同意另有規定外，由各邦規範機關之設置及行政程序。」}}
% 別名（墮胎一案使用之縮寫命名）
\let\ggfnArtEightFourOneDE\ggfnArtEightyFourOneDE
\let\ggfnArtEightFourOneZH\ggfnArtEightyFourOneZH

% ===== Art. 93 =====
\newcommand{\ggfnArtNinetyThreeOneTwoDE}{\ggnotedeonce{art-93-1-2}{Art. 93 Abs. 1 Nr. 2: Das Bundesverfassungsgericht entscheidet ``bei Meinungsverschiedenheiten oder Zweifeln ueber die foermliche und sachliche Vereinbarkeit von Bundesrecht oder Landesrecht mit diesem Grundgesetze ... auf Antrag der Bundesregierung, einer Landesregierung oder eines Drittels der Mitglieder des Bundestages''.}}
\newcommand{\ggfnArtNinetyThreeOneTwoZH}{\ggnotezhonce{art-93-1-2}{第93條第1項第2款：聯邦憲法法院就「聯邦法或邦法在形式及實質上是否與本基本法相容之意見分歧或疑義……依聯邦政府、邦政府或聯邦議院三分之一議員之聲請」作成裁判。}}
% 別名（墮胎一案使用之縮寫命名）
\let\ggfnArtNineThreeOneTwoDE\ggfnArtNinetyThreeOneTwoDE
\let\ggfnArtNineThreeOneTwoZH\ggfnArtNinetyThreeOneTwoZH

% ===== Art. 102 =====
\newcommand{\ggfnArtOneZeroTwoDE}{\ggnotedeonce{art-102}{Art. 102: ``Die Todesstrafe ist abgeschafft.''}}
\newcommand{\ggfnArtOneZeroTwoZH}{\ggnotezhonce{art-102}{第102條：「死刑廢止。」}}

% ===== Art. 103 =====
\newcommand{\ggfnArtOneZeroThreeTwoDE}{\ggnotedeonce{art-103-2}{Art. 103 Abs. 2: ``Eine Tat kann nur bestraft werden, wenn die Strafbarkeit gesetzlich bestimmt war, bevor die Tat begangen wurde.''}}
\newcommand{\ggfnArtOneZeroThreeTwoZH}{\ggnotezhonce{art-103-2}{第103條第2項：「行為之可罰性，須於行為前以法律定之，始得處罰。」}}

% ===== Art. 104 =====
\newcommand{\ggfnArtOneZeroFourOneDE}{\ggnotedeonce{art-104-1}{Art. 104 Abs. 1: ``Die Freiheit der Person kann nur auf Grund eines foermlichen Gesetzes und nur unter Beachtung der darin vorgeschriebenen Formen beschraenkt werden.''}}
\newcommand{\ggfnArtOneZeroFourOneZH}{\ggnotezhonce{art-104-1}{第104條第1項：「人身自由僅得基於形式法律，並遵守其中規定之形式加以限制。」}}

% ===== 組合腳註：Art. 3 + Art. 6（墮胎一案特有）=====
\newcommand{\ggfnArtThreeSixDE}{\ggnotedeonce{art-3-6}{Art. 3: ``Alle Menschen sind vor dem Gesetz gleich. Maenner und Frauen sind gleichberechtigt. Niemand darf wegen seines Geschlechtes, seiner Abstammung, seiner Rasse, seiner Sprache, seiner Heimat und Herkunft, seines Glaubens, seiner religioesen oder politischen Anschauungen benachteiligt oder bevorzugt werden.'' Art. 6 Abs. 1: ``Ehe und Familie stehen unter dem besonderen Schutze der staatlichen Ordnung.'' Abs. 2 Satz 1: ``Pflege und Erziehung der Kinder sind das natuerliche Recht der Eltern und die zuvoerderst ihnen obliegende Pflicht.'' Abs. 4: ``Jede Mutter hat Anspruch auf den Schutz und die Fuersorge der Gemeinschaft.''}\ggmarknotede{art-3}\ggmarknotede{art-6}\ggmarknotede{art-6-1-4}}
\newcommand{\ggfnArtThreeSixZH}{\ggnotezhonce{art-3-6}{第3條：「所有人在法律前一律平等。男女享有平等權利。任何人不得因其性別、血統、種族、語言、故鄉及出身、信仰、宗教或政治見解而受不利益或優惠。」第6條第1項：「婚姻與家庭受國家秩序之特別保護。」第2項第1句：「照護及教育子女，為父母之自然權利，並為首先由其承擔之義務。」第4項：「每一母親均有請求共同體保護與照顧之權利。」}\ggmarknotezh{art-3}\ggmarknotezh{art-6}\ggmarknotezh{art-6-1-4}}

% ===== 組合腳註：Art. 6(1) + Art. 6(4)（墮胎一案特有）=====
\newcommand{\ggfnArtSixOneFourDE}{\ggnotedeonce{art-6-1-4}{Art. 6 Abs. 1: ``Ehe und Familie stehen unter dem besonderen Schutze der staatlichen Ordnung.'' Art. 6 Abs. 4: ``Jede Mutter hat Anspruch auf den Schutz und die Fuersorge der Gemeinschaft.''}\ggmarknotede{art-6}}
\newcommand{\ggfnArtSixOneFourZH}{\ggnotezhonce{art-6-1-4}{第6條第1項：「婚姻與家庭受國家秩序之特別保護。」第4項：「每一母親均有請求共同體保護與照顧之權利。」}\ggmarknotezh{art-6}}

% ===== 組合腳註：Art. 1(1) + Art. 2(2)(1)（墮胎二案特有）=====
\newcommand{\ggfnArtOneOneTwoTwoOneDE}{\ggnotedeonce{art-1-1+2-2-1}{Art. 1 Abs. 1: ``Die Wuerde des Menschen ist unantastbar. Sie zu achten und zu schuetzen ist Verpflichtung aller staatlichen Gewalt.'' Art. 2 Abs. 2 Satz 1: ``Jeder hat das Recht auf Leben und koerperliche Unversehrtheit.''}\ggmarknotede{art-1-1}\ggmarknotede{art-2-2-1}}
\newcommand{\ggfnArtOneOneTwoTwoOneZH}{\ggnotezhonce{art-1-1+2-2-1}{第1條第1項：「人之尊嚴不可侵犯。尊重及保護之，為一切國家權力之義務。」第2條第2項第1句：「人人有生命與身體完整性之權利。」}\ggmarknotezh{art-1-1}\ggmarknotezh{art-2-2-1}}

% ===== 組合腳註：Art. 1(1) + Art. 2(2)（墮胎二案特有）=====
\newcommand{\ggfnArtOneOneTwoTwoDE}{\ggnotedeonce{art-1-1+2-2}{Art. 1 Abs. 1: ``Die Wuerde des Menschen ist unantastbar. Sie zu achten und zu schuetzen ist Verpflichtung aller staatlichen Gewalt.'' Art. 2 Abs. 2: ``Jeder hat das Recht auf Leben und koerperliche Unversehrtheit. Die Freiheit der Person ist unverletzlich. In diese Rechte darf nur auf Grund eines Gesetzes eingegriffen werden.''}\ggmarknotede{art-1-1}\ggmarknotede{art-2-2}\ggmarknotede{art-2-2-1}}
\newcommand{\ggfnArtOneOneTwoTwoZH}{\ggnotezhonce{art-1-1+2-2}{第1條第1項：「人之尊嚴不可侵犯。尊重及保護之，為一切國家權力之義務。」第2條第2項：「人人有生命與身體完整性之權利。人身自由不可侵犯。此等權利僅得依法律加以干預。」}\ggmarknotezh{art-1-1}\ggmarknotezh{art-2-2}\ggmarknotezh{art-2-2-1}}

% ===== 組合腳註：Art. 1(1) + Art. 2(1)（墮胎二案特有）=====
\newcommand{\ggfnArtOneOneTwoOneDE}{\ggnotedeonce{art-1-1+2-1}{Art. 1 Abs. 1: ``Die Wuerde des Menschen ist unantastbar. Sie zu achten und zu schuetzen ist Verpflichtung aller staatlichen Gewalt.'' Art. 2 Abs. 1: ``Jeder hat das Recht auf die freie Entfaltung seiner Persoenlichkeit, soweit er nicht die Rechte anderer verletzt und nicht gegen die verfassungsmaessige Ordnung oder das Sittengesetz verstoesst.''}\ggmarknotede{art-1-1}\ggmarknotede{art-2-1}}
\newcommand{\ggfnArtOneOneTwoOneZH}{\ggnotezhonce{art-1-1+2-1}{第1條第1項：「人之尊嚴不可侵犯。尊重及保護之，為一切國家權力之義務。」第2條第1項：「人人有自由發展其人格之權利，但不得侵害他人權利，亦不得違反合憲秩序或道德律。」}\ggmarknotezh{art-1-1}\ggmarknotezh{art-2-1}}

% ===== 組合腳註：Art. 2(2)(1) + Art. 1(1)（墮胎一案特有，含交叉標記）=====
\newcommand{\ggfnLifeDignityDE}{\ggnotedeonce{life-dignity}{Art. 2 Abs. 2 Satz 1: ``Jeder hat das Recht auf Leben und koerperliche Unversehrtheit.'' Art. 1 Abs. 1: ``Die Wuerde des Menschen ist unantastbar. Sie zu achten und zu schuetzen ist Verpflichtung aller staatlichen Gewalt.''}\ggmarknotede{art-2-2-1}\ggmarknotede{art-1-1}\ggmarknotede{art-1-1-2}}
\newcommand{\ggfnLifeDignityZH}{\ggnotezhonce{life-dignity}{第2條第2項第1句：「人人有生命與身體完整性之權利。」第1條第1項：「人之尊嚴不可侵犯。尊重及保護之，為一切國家權力之義務。」}\ggmarknotezh{art-2-2-1}\ggmarknotezh{art-1-1}\ggmarknotezh{art-1-1-2}}

% ===== 組合腳註：Art. 1(1) + Art. 2(2)（墮胎一案特有，條件判斷版）=====
\newcommand{\ggfnArtOneTwoTwoDE}{\ifcsname ggnoteseen@de@life-dignity\endcsname\ggfnArtTwoTwoRestDE{}\else\ggnotedeonce{art-1-2-2}{Art. 1 Abs. 1: ``Die Wuerde des Menschen ist unantastbar. Sie zu achten und zu schuetzen ist Verpflichtung aller staatlichen Gewalt.'' Art. 2 Abs. 2: ``Jeder hat das Recht auf Leben und koerperliche Unversehrtheit. Die Freiheit der Person ist unverletzlich. In diese Rechte darf nur auf Grund eines Gesetzes eingegriffen werden.''}\ggmarknotede{art-1-1}\ggmarknotede{art-1-1-2}\ggmarknotede{art-2-2}\ggmarknotede{art-2-2-1}\fi}
\newcommand{\ggfnArtOneTwoTwoZH}{\ifcsname ggnoteseen@zh@life-dignity\endcsname\ggfnArtTwoTwoRestZH{}\else\ggnotezhonce{art-1-2-2}{第1條第1項：「人之尊嚴不可侵犯。尊重及保護之，為一切國家權力之義務。」第2條第2項：「人人有生命與身體完整性之權利。人身自由不可侵犯。此等權利僅得依法律加以干預。」}\ggmarknotezh{art-1-1}\ggmarknotezh{art-1-1-2}\ggmarknotezh{art-2-2}\ggmarknotezh{art-2-2-1}\fi}

% ===== 組合腳註：Art. 1(1) + Art. 19(2)（墮胎一案特有，條件判斷版）=====
\newcommand{\ggfnArtOneNineteenDE}{\ifcsname ggnoteseen@de@art-1-1\endcsname\ggfnArtNineteenTwoDE{}\else\ggnotedeonce{art-1-19}{Art. 1 Abs. 1: ``Die Wuerde des Menschen ist unantastbar. Sie zu achten und zu schuetzen ist Verpflichtung aller staatlichen Gewalt.'' Art. 19 Abs. 2: ``In keinem Falle darf ein Grundrecht in seinem Wesensgehalt angetastet werden.''}\ggmarknotede{art-1-1}\ggmarknotede{art-1-1-2}\ggmarknotede{art-19-2}\fi}
\newcommand{\ggfnArtOneNineteenZH}{\ifcsname ggnoteseen@zh@art-1-1\endcsname\ggfnArtNineteenTwoZH{}\else\ggnotezhonce{art-1-19}{第1條第1項：「人之尊嚴不可侵犯。尊重及保護之，為一切國家權力之義務。」第19條第2項：「任何情形下，基本權利之本質內容均不得受到侵害。」}\ggmarknotezh{art-1-1}\ggmarknotezh{art-1-1-2}\ggmarknotezh{art-19-2}\fi}
 之前定義）=====
\newif\ifggversionnotede
\newif\ifggversionnotezh

% ===== 編者註腳基礎設施（首次標記模式）=====
\newif\ifeditornotelabelde
\newif\ifeditornotelabelzh
\newcommand{\editornotelabelde}{\ifeditornotelabelde\else\emph{Hrsg.-Anm. (nachfolgend ohne Kennzeichnung)}: \global\editornotelabeldetrue\fi}
\newcommand{\editornotelabelzh}{\ifeditornotelabelzh\else 編者註（下同）：\global\editornotelabelzhtrue\fi}
\newcommand{\editornotede}[1]{\ifshoweditornotes\footnote{\normalfont\footnotesize\editornotelabelde #1}\else\ifclassroommode\footnote{\normalfont\footnotesize\editornotelabelde #1}\fi\fi}
\newcommand{\editornotezh}[1]{\ifshoweditornotes\footnote{\normalfont\footnotesize\editornotelabelzh #1}\else\ifclassroommode\footnote{\normalfont\footnotesize\editornotelabelzh #1}\fi\fi}

% ===== GG 引用系統（\ggversionde/zh 由各案在 \input 前定義）=====
\newcommand{\ggmarknotede}[1]{\global\expandafter\let\csname ggnoteseen@de@#1\endcsname\empty}
\newcommand{\ggmarknotezh}[1]{\global\expandafter\let\csname ggnoteseen@zh@#1\endcsname\empty}
\newcommand{\ggnotedeonce}[2]{\ifcsname ggnoteseen@de@#1\endcsname\else\global\expandafter\let\csname ggnoteseen@de@#1\endcsname\empty\ggnotede{#2}\fi}
\newcommand{\ggnotezhonce}[2]{\ifcsname ggnoteseen@zh@#1\endcsname\else\global\expandafter\let\csname ggnoteseen@zh@#1\endcsname\empty\ggnotezh{#2}\fi}
\newcommand{\ggnotede}[1]{\editornotede{\ggversionde #1}}
\newcommand{\ggnotezh}[1]{\editornotezh{\ggversionzh #1}}

% ===== Art. 1 =====
\newcommand{\ggfnArtOneOneDE}{\ggnotedeonce{art-1-1}{Art. 1 Abs. 1: ``Die Wuerde des Menschen ist unantastbar. Sie zu achten und zu schuetzen ist Verpflichtung aller staatlichen Gewalt.''}}
\newcommand{\ggfnArtOneOneZH}{\ggnotezhonce{art-1-1}{第1條第1項：「人之尊嚴不可侵犯。尊重及保護之，為一切國家權力之義務。」}}
\newcommand{\ggfnArtOneOneTwoDE}{\ggnotedeonce{art-1-1-2}{Art. 1 Abs. 1 Satz 2: ``Sie zu achten und zu schuetzen ist Verpflichtung aller staatlichen Gewalt.''}}
\newcommand{\ggfnArtOneOneTwoZH}{\ggnotezhonce{art-1-1-2}{第1條第1項第2句：「尊重及保護之，為一切國家權力之義務。」}}

% ===== Art. 2 =====
\newcommand{\ggfnArtTwoOneDE}{\ggnotedeonce{art-2-1}{Art. 2 Abs. 1: ``Jeder hat das Recht auf die freie Entfaltung seiner Persoenlichkeit, soweit er nicht die Rechte anderer verletzt und nicht gegen die verfassungsmaessige Ordnung oder das Sittengesetz verstoesst.''}}
\newcommand{\ggfnArtTwoOneZH}{\ggnotezhonce{art-2-1}{第2條第1項：「人人有自由發展其人格之權利，但不得侵害他人權利，亦不得違反合憲秩序或道德律。」}}
% 簡易預設版（墮胎一案以 \renewcommand 覆寫為含條件邏輯之版本）
\newcommand{\ggfnArtTwoTwoDE}{\ggnotedeonce{art-2-2}{Art. 2 Abs. 2: ``Jeder hat das Recht auf Leben und koerperliche Unversehrtheit. Die Freiheit der Person ist unverletzlich. In diese Rechte darf nur auf Grund eines Gesetzes eingegriffen werden.''}\ggmarknotede{art-2-2-1}}
\newcommand{\ggfnArtTwoTwoZH}{\ggnotezhonce{art-2-2}{第2條第2項：「人人有生命與身體完整性之權利。人身自由不可侵犯。此等權利僅得依法律加以干預。」}\ggmarknotezh{art-2-2-1}}
\newcommand{\ggfnArtTwoTwoOneDE}{\ggnotedeonce{art-2-2-1}{Art. 2 Abs. 2 Satz 1: ``Jeder hat das Recht auf Leben und koerperliche Unversehrtheit.''}}
\newcommand{\ggfnArtTwoTwoOneZH}{\ggnotezhonce{art-2-2-1}{第2條第2項第1句：「人人有生命與身體完整性之權利。」}}
\newcommand{\ggfnArtTwoTwoThreeDE}{\ggnotedeonce{art-2-2-3}{Art. 2 Abs. 2 Satz 3: ``In diese Rechte darf nur auf Grund eines Gesetzes eingegriffen werden.''}}
\newcommand{\ggfnArtTwoTwoThreeZH}{\ggnotezhonce{art-2-2-3}{第2條第2項第3句：「此等權利僅得依法律加以干預。」}}
% 墮胎一案特有：第2條第2項第2–3句（Art. 2(2) 去掉第1句後的「餘文」）
\newcommand{\ggfnArtTwoTwoRestDE}{\ggnotedeonce{art-2-2-rest}{Art. 2 Abs. 2 Saetze 2-3: ``Die Freiheit der Person ist unverletzlich. In diese Rechte darf nur auf Grund eines Gesetzes eingegriffen werden.''}\ggmarknotede{art-2-2}}
\newcommand{\ggfnArtTwoTwoRestZH}{\ggnotezhonce{art-2-2-rest}{第2條第2項第2、3句：「人身自由不可侵犯。此等權利僅得依法律加以干預。」}\ggmarknotezh{art-2-2}}

% ===== Art. 3 =====
\newcommand{\ggfnArtThreeDE}{\ggnotedeonce{art-3}{Art. 3: ``Alle Menschen sind vor dem Gesetz gleich. Maenner und Frauen sind gleichberechtigt. Niemand darf wegen seines Geschlechtes, seiner Abstammung, seiner Rasse, seiner Sprache, seiner Heimat und Herkunft, seines Glaubens, seiner religioesen oder politischen Anschauungen benachteiligt oder bevorzugt werden.''}}
\newcommand{\ggfnArtThreeZH}{\ggnotezhonce{art-3}{第3條：「所有人在法律前一律平等。男女享有平等權利。任何人不得因其性別、血統、種族、語言、故鄉及出身、信仰、宗教或政治見解而受不利益或優惠。」}}
\newcommand{\ggfnArtThreeTwoDE}{\ggnotedeonce{art-3-2}{Art. 3 Abs. 2: ``Maenner und Frauen sind gleichberechtigt.''}}
\newcommand{\ggfnArtThreeTwoZH}{\ggnotezhonce{art-3-2}{第3條第2項：「男女享有平等權利。」}}

% ===== Art. 4 =====
\newcommand{\ggfnArtFourOneDE}{\ggnotedeonce{art-4-1}{Art. 4 Abs. 1: ``Die Freiheit des Glaubens, des Gewissens und die Freiheit des religioesen und weltanschaulichen Bekenntnisses sind unverletzlich.''}}
\newcommand{\ggfnArtFourOneZH}{\ggnotezhonce{art-4-1}{第4條第1項：「信仰自由、良心自由，以及宗教與世界觀信念自由，不可侵犯。」}}

% ===== Art. 5 =====
\newcommand{\ggfnArtFiveOneDE}{\ggnotedeonce{art-5-1}{Art. 5 Abs. 1: ``Jeder hat das Recht, seine Meinung in Wort, Schrift und Bild frei zu aeussern und zu verbreiten und sich aus allgemein zugaenglichen Quellen ungehindert zu unterrichten. Die Pressefreiheit und die Freiheit der Berichterstattung durch Rundfunk und Film werden gewaehrleistet. Eine Zensur findet nicht statt.''}}
\newcommand{\ggfnArtFiveOneZH}{\ggnotezhonce{art-5-1}{第5條第1項：「人人有以言語、文字及圖像自由表達及傳播其意見，並由一般可接近來源不受阻礙獲取資訊之權利。新聞自由及廣播、電影報導自由受保障。不設審查制度。」}}

% ===== Art. 6 =====
\newcommand{\ggfnArtSixDE}{\ggnotedeonce{art-6}{Art. 6 Abs. 1: ``Ehe und Familie stehen unter dem besonderen Schutze der staatlichen Ordnung.'' Abs. 2 Satz 1: ``Pflege und Erziehung der Kinder sind das natuerliche Recht der Eltern und die zuvoerderst ihnen obliegende Pflicht.'' Abs. 4: ``Jede Mutter hat Anspruch auf den Schutz und die Fuersorge der Gemeinschaft.''}}
\newcommand{\ggfnArtSixZH}{\ggnotezhonce{art-6}{第6條第1項：「婚姻與家庭受國家秩序之特別保護。」第2項第1句：「照護及教育子女，為父母之自然權利，並為首先由其承擔之義務。」第4項：「每一母親均有請求共同體保護與照顧之權利。」}}
\newcommand{\ggfnArtSixOneDE}{\ggnotedeonce{art-6-1}{Art. 6 Abs. 1: ``Ehe und Familie stehen unter dem besonderen Schutze der staatlichen Ordnung.''}}
\newcommand{\ggfnArtSixOneZH}{\ggnotezhonce{art-6-1}{第6條第1項：「婚姻與家庭受國家秩序之特別保護。」}}
\newcommand{\ggfnArtSixTwoDE}{\ggnotedeonce{art-6-2}{Art. 6 Abs. 2: ``Pflege und Erziehung der Kinder sind das natuerliche Recht der Eltern und die zuvoerderst ihnen obliegende Pflicht. Ueber ihre Betaetigung wacht die staatliche Gemeinschaft.''}}
\newcommand{\ggfnArtSixTwoZH}{\ggnotezhonce{art-6-2}{第6條第2項：「照護及教育子女，為父母之自然權利，並為首先由其承擔之義務。國家共同體監督其行使。」}}
\newcommand{\ggfnArtSixTwoOneDE}{\ggnotedeonce{art-6-2-1}{Art. 6 Abs. 2 Satz 1: ``Pflege und Erziehung der Kinder sind das natuerliche Recht der Eltern und die zuvoerderst ihnen obliegende Pflicht.''}}
\newcommand{\ggfnArtSixTwoOneZH}{\ggnotezhonce{art-6-2-1}{第6條第2項第1句：「照護及教育子女，為父母之自然權利，並為首先由其承擔之義務。」}}
\newcommand{\ggfnArtSixFourDE}{\ggnotedeonce{art-6-4}{Art. 6 Abs. 4: ``Jede Mutter hat Anspruch auf den Schutz und die Fuersorge der Gemeinschaft.''}}
\newcommand{\ggfnArtSixFourZH}{\ggnotezhonce{art-6-4}{第6條第4項：「每一母親均有請求共同體保護與照顧之權利。」}}

% ===== Art. 12 =====
\newcommand{\ggfnArtTwelveOneDE}{\ggnotedeonce{art-12-1}{Art. 12 Abs. 1: ``Alle Deutschen haben das Recht, Beruf, Arbeitsplatz und Ausbildungsstaette frei zu waehlen. Die Berufsausuebung kann durch Gesetz oder auf Grund eines Gesetzes geregelt werden.''}}
\newcommand{\ggfnArtTwelveOneZH}{\ggnotezhonce{art-12-1}{第12條第1項：「所有德國人均有自由選擇職業、工作場所及訓練場所之權利。職業行使得以法律或基於法律加以規範。」}}

% ===== Art. 19 =====
\newcommand{\ggfnArtNineteenTwoDE}{\ggnotedeonce{art-19-2}{Art. 19 Abs. 2: ``In keinem Falle darf ein Grundrecht in seinem Wesensgehalt angetastet werden.''}}
\newcommand{\ggfnArtNineteenTwoZH}{\ggnotezhonce{art-19-2}{第19條第2項：「任何情形下，基本權利之本質內容均不得受到侵害。」}}

% ===== Art. 20 =====
\newcommand{\ggfnArtTwentyOneDE}{\ggnotedeonce{art-20-1}{Art. 20 Abs. 1: ``Die Bundesrepublik Deutschland ist ein demokratischer und sozialer Bundesstaat.''}}
\newcommand{\ggfnArtTwentyOneZH}{\ggnotezhonce{art-20-1}{第20條第1項：「德意志聯邦共和國為民主及社會之聯邦國。」}}
\newcommand{\ggfnArtTwentyThreeDE}{\ggnotedeonce{art-20-3}{Art. 20 Abs. 3: ``Die Gesetzgebung ist an die verfassungsmaessige Ordnung, die vollziehende Gewalt und die Rechtsprechung sind an Gesetz und Recht gebunden.''}}
\newcommand{\ggfnArtTwentyThreeZH}{\ggnotezhonce{art-20-3}{第20條第3項：「立法受合憲秩序拘束；行政權與司法權受法律與法拘束。」}}

% ===== Art. 26 + Art. 143（墮胎一案特有）=====
\newcommand{\ggfnArtTwoSixAndOneFourThreeDE}{\ggnotedeonce{art-26-143}{Art. 26 Abs. 1 Satz 2: ``Sie sind unter Strafe zu stellen.'' Art. 143 in der urspruenglichen Fassung (24. Mai 1949 bis 1. September 1951) enthielt Strafvorschriften zum Hochverrat; 1975 war Art. 143 bereits weggefallen.}}
\newcommand{\ggfnArtTwoSixAndOneFourThreeZH}{\ggnotezhonce{art-26-143}{第26條第1項第2句：「此等行為應規定為犯罪處罰。」第143條原始版本（1949年5月24日至1951年9月1日）含有關於叛國罪之刑罰規定；至1975年時，第143條已廢止。}}

% ===== Art. 28 =====
\newcommand{\ggfnArtTwentyEightOneDE}{\ggnotedeonce{art-28-1-1}{Art. 28 Abs. 1 Satz 1: ``Die verfassungsmaessige Ordnung in den Laendern muss den Grundsaetzen des republikanischen, demokratischen und sozialen Rechtsstaates im Sinne dieses Grundgesetzes entsprechen.''}}
\newcommand{\ggfnArtTwentyEightOneZH}{\ggnotezhonce{art-28-1-1}{第28條第1項第1句：「各邦之合憲秩序，須符合本基本法意義下共和、民主及社會法治國之原則。」}}

% ===== Art. 30 =====
\newcommand{\ggfnArtThirtyDE}{\ggnotedeonce{art-30}{Art. 30: ``Die Ausuebung der staatlichen Befugnisse und die Erfuellung der staatlichen Aufgaben ist Sache der Laender, soweit dieses Grundgesetz keine andere Regelung trifft oder zulaesst.''}}
\newcommand{\ggfnArtThirtyZH}{\ggnotezhonce{art-30}{第30條：「國家權限之行使與國家任務之履行，屬於各邦，但本基本法另有規定或容許者，不在此限。」}}

% ===== Art. 74 =====
\newcommand{\ggfnArtSeventyFourOneDE}{\ggnotedeonce{art-74-1}{Art. 74 Nr. 1: Die konkurrierende Gesetzgebung erstreckt sich auf ``das buergerliche Recht, das Strafrecht und den Strafvollzug, die Gerichtsverfassung, das gerichtliche Verfahren, die Rechtsanwaltschaft, das Notariat und die Rechtsberatung''.}}
\newcommand{\ggfnArtSeventyFourOneZH}{\ggnotezhonce{art-74-1}{第74條第1款：競合立法及於「民法、刑法及刑罰執行、法院組織、訴訟程序、律師制度、公證制度及法律諮詢」。}}
\newcommand{\ggfnArtSeventyFourSevenDE}{\ggnotedeonce{art-74-7}{Art. 74 Nr. 7: Die konkurrierende Gesetzgebung erstreckt sich auf ``die oeffentliche Fuersorge''.}}
\newcommand{\ggfnArtSeventyFourSevenZH}{\ggnotezhonce{art-74-7}{第74條第7款：競合立法及於「公共扶助」。}}
\newcommand{\ggfnArtSeventyFourTwelveDE}{\ggnotedeonce{art-74-12}{Art. 74 Nr. 12: Die konkurrierende Gesetzgebung erstreckt sich auf ``das Arbeitsrecht einschliesslich der Betriebsverfassung, des Arbeitsschutzes und der Arbeitsvermittlung sowie die Sozialversicherung einschliesslich der Arbeitslosenversicherung''.}}
\newcommand{\ggfnArtSeventyFourTwelveZH}{\ggnotezhonce{art-74-12}{第74條第12款：競合立法及於「勞動法，包括企業組織、勞動保護及就業媒合，以及社會保險，包括失業保險」。}}

% ===== Art. 77（墮胎一案特有）=====
\newcommand{\ggfnArtSevenSevenThreeDE}{\ggnotedeonce{art-77-3}{Art. 77 Abs. 3 Satz 1: ``Soweit zu einem Gesetze die Zustimmung des Bundesrates nicht erforderlich ist, kann der Bundesrat, wenn das Verfahren nach Absatz 2 beendigt ist, gegen ein vom Bundestage beschlossenes Gesetz binnen zwei Wochen Einspruch einlegen.''}}
\newcommand{\ggfnArtSevenSevenThreeZH}{\ggnotezhonce{art-77-3}{第77條第3項第1句：「法律無須聯邦參議院同意者，於第2項程序終了後，聯邦參議院得於二週內對聯邦議院通過之法律提出異議。」}}

% ===== Art. 80 =====
\newcommand{\ggfnArtEightyOneOneDE}{\ggnotedeonce{art-80-1-1}{Art. 80 Abs. 1 Satz 1: ``Durch Gesetz koennen die Bundesregierung, ein Bundesminister oder die Landesregierungen ermaechtigt werden, Rechtsverordnungen zu erlassen.''}}
\newcommand{\ggfnArtEightyOneOneZH}{\ggnotezhonce{art-80-1-1}{第80條第1項第1句：「得以法律授權聯邦政府、聯邦部長或邦政府發布法規命令。」}}

% ===== Art. 83 =====
\newcommand{\ggfnArtEightyThreeDE}{\ggnotedeonce{art-83}{Art. 83: ``Die Laender fuehren die Bundesgesetze als eigene Angelegenheit aus, soweit dieses Grundgesetz nichts anderes bestimmt oder zulaesst.''}}
\newcommand{\ggfnArtEightyThreeZH}{\ggnotezhonce{art-83}{第83條：「各邦以自己事務執行聯邦法律，但本基本法另有規定或容許者，不在此限。」}}

% ===== Art. 84 =====
\newcommand{\ggfnArtEightyFourOneDE}{\ggnotedeonce{art-84-1}{Art. 84 Abs. 1: ``Fuehren die Laender die Bundesgesetze als eigene Angelegenheit aus, so regeln sie die Einrichtung der Behoerden und das Verwaltungsverfahren, soweit nicht Bundesgesetze mit Zustimmung des Bundesrates etwas anderes bestimmen.''}}
\newcommand{\ggfnArtEightyFourOneZH}{\ggnotezhonce{art-84-1}{第84條第1項：「各邦以自己事務執行聯邦法律時，除聯邦法律經聯邦參議院同意另有規定外，由各邦規範機關之設置及行政程序。」}}
% 別名（墮胎一案使用之縮寫命名）
\let\ggfnArtEightFourOneDE\ggfnArtEightyFourOneDE
\let\ggfnArtEightFourOneZH\ggfnArtEightyFourOneZH

% ===== Art. 93 =====
\newcommand{\ggfnArtNinetyThreeOneTwoDE}{\ggnotedeonce{art-93-1-2}{Art. 93 Abs. 1 Nr. 2: Das Bundesverfassungsgericht entscheidet ``bei Meinungsverschiedenheiten oder Zweifeln ueber die foermliche und sachliche Vereinbarkeit von Bundesrecht oder Landesrecht mit diesem Grundgesetze ... auf Antrag der Bundesregierung, einer Landesregierung oder eines Drittels der Mitglieder des Bundestages''.}}
\newcommand{\ggfnArtNinetyThreeOneTwoZH}{\ggnotezhonce{art-93-1-2}{第93條第1項第2款：聯邦憲法法院就「聯邦法或邦法在形式及實質上是否與本基本法相容之意見分歧或疑義……依聯邦政府、邦政府或聯邦議院三分之一議員之聲請」作成裁判。}}
% 別名（墮胎一案使用之縮寫命名）
\let\ggfnArtNineThreeOneTwoDE\ggfnArtNinetyThreeOneTwoDE
\let\ggfnArtNineThreeOneTwoZH\ggfnArtNinetyThreeOneTwoZH

% ===== Art. 102 =====
\newcommand{\ggfnArtOneZeroTwoDE}{\ggnotedeonce{art-102}{Art. 102: ``Die Todesstrafe ist abgeschafft.''}}
\newcommand{\ggfnArtOneZeroTwoZH}{\ggnotezhonce{art-102}{第102條：「死刑廢止。」}}

% ===== Art. 103 =====
\newcommand{\ggfnArtOneZeroThreeTwoDE}{\ggnotedeonce{art-103-2}{Art. 103 Abs. 2: ``Eine Tat kann nur bestraft werden, wenn die Strafbarkeit gesetzlich bestimmt war, bevor die Tat begangen wurde.''}}
\newcommand{\ggfnArtOneZeroThreeTwoZH}{\ggnotezhonce{art-103-2}{第103條第2項：「行為之可罰性，須於行為前以法律定之，始得處罰。」}}

% ===== Art. 104 =====
\newcommand{\ggfnArtOneZeroFourOneDE}{\ggnotedeonce{art-104-1}{Art. 104 Abs. 1: ``Die Freiheit der Person kann nur auf Grund eines foermlichen Gesetzes und nur unter Beachtung der darin vorgeschriebenen Formen beschraenkt werden.''}}
\newcommand{\ggfnArtOneZeroFourOneZH}{\ggnotezhonce{art-104-1}{第104條第1項：「人身自由僅得基於形式法律，並遵守其中規定之形式加以限制。」}}

% ===== 組合腳註：Art. 3 + Art. 6（墮胎一案特有）=====
\newcommand{\ggfnArtThreeSixDE}{\ggnotedeonce{art-3-6}{Art. 3: ``Alle Menschen sind vor dem Gesetz gleich. Maenner und Frauen sind gleichberechtigt. Niemand darf wegen seines Geschlechtes, seiner Abstammung, seiner Rasse, seiner Sprache, seiner Heimat und Herkunft, seines Glaubens, seiner religioesen oder politischen Anschauungen benachteiligt oder bevorzugt werden.'' Art. 6 Abs. 1: ``Ehe und Familie stehen unter dem besonderen Schutze der staatlichen Ordnung.'' Abs. 2 Satz 1: ``Pflege und Erziehung der Kinder sind das natuerliche Recht der Eltern und die zuvoerderst ihnen obliegende Pflicht.'' Abs. 4: ``Jede Mutter hat Anspruch auf den Schutz und die Fuersorge der Gemeinschaft.''}\ggmarknotede{art-3}\ggmarknotede{art-6}\ggmarknotede{art-6-1-4}}
\newcommand{\ggfnArtThreeSixZH}{\ggnotezhonce{art-3-6}{第3條：「所有人在法律前一律平等。男女享有平等權利。任何人不得因其性別、血統、種族、語言、故鄉及出身、信仰、宗教或政治見解而受不利益或優惠。」第6條第1項：「婚姻與家庭受國家秩序之特別保護。」第2項第1句：「照護及教育子女，為父母之自然權利，並為首先由其承擔之義務。」第4項：「每一母親均有請求共同體保護與照顧之權利。」}\ggmarknotezh{art-3}\ggmarknotezh{art-6}\ggmarknotezh{art-6-1-4}}

% ===== 組合腳註：Art. 6(1) + Art. 6(4)（墮胎一案特有）=====
\newcommand{\ggfnArtSixOneFourDE}{\ggnotedeonce{art-6-1-4}{Art. 6 Abs. 1: ``Ehe und Familie stehen unter dem besonderen Schutze der staatlichen Ordnung.'' Art. 6 Abs. 4: ``Jede Mutter hat Anspruch auf den Schutz und die Fuersorge der Gemeinschaft.''}\ggmarknotede{art-6}}
\newcommand{\ggfnArtSixOneFourZH}{\ggnotezhonce{art-6-1-4}{第6條第1項：「婚姻與家庭受國家秩序之特別保護。」第4項：「每一母親均有請求共同體保護與照顧之權利。」}\ggmarknotezh{art-6}}

% ===== 組合腳註：Art. 1(1) + Art. 2(2)(1)（墮胎二案特有）=====
\newcommand{\ggfnArtOneOneTwoTwoOneDE}{\ggnotedeonce{art-1-1+2-2-1}{Art. 1 Abs. 1: ``Die Wuerde des Menschen ist unantastbar. Sie zu achten und zu schuetzen ist Verpflichtung aller staatlichen Gewalt.'' Art. 2 Abs. 2 Satz 1: ``Jeder hat das Recht auf Leben und koerperliche Unversehrtheit.''}\ggmarknotede{art-1-1}\ggmarknotede{art-2-2-1}}
\newcommand{\ggfnArtOneOneTwoTwoOneZH}{\ggnotezhonce{art-1-1+2-2-1}{第1條第1項：「人之尊嚴不可侵犯。尊重及保護之，為一切國家權力之義務。」第2條第2項第1句：「人人有生命與身體完整性之權利。」}\ggmarknotezh{art-1-1}\ggmarknotezh{art-2-2-1}}

% ===== 組合腳註：Art. 1(1) + Art. 2(2)（墮胎二案特有）=====
\newcommand{\ggfnArtOneOneTwoTwoDE}{\ggnotedeonce{art-1-1+2-2}{Art. 1 Abs. 1: ``Die Wuerde des Menschen ist unantastbar. Sie zu achten und zu schuetzen ist Verpflichtung aller staatlichen Gewalt.'' Art. 2 Abs. 2: ``Jeder hat das Recht auf Leben und koerperliche Unversehrtheit. Die Freiheit der Person ist unverletzlich. In diese Rechte darf nur auf Grund eines Gesetzes eingegriffen werden.''}\ggmarknotede{art-1-1}\ggmarknotede{art-2-2}\ggmarknotede{art-2-2-1}}
\newcommand{\ggfnArtOneOneTwoTwoZH}{\ggnotezhonce{art-1-1+2-2}{第1條第1項：「人之尊嚴不可侵犯。尊重及保護之，為一切國家權力之義務。」第2條第2項：「人人有生命與身體完整性之權利。人身自由不可侵犯。此等權利僅得依法律加以干預。」}\ggmarknotezh{art-1-1}\ggmarknotezh{art-2-2}\ggmarknotezh{art-2-2-1}}

% ===== 組合腳註：Art. 1(1) + Art. 2(1)（墮胎二案特有）=====
\newcommand{\ggfnArtOneOneTwoOneDE}{\ggnotedeonce{art-1-1+2-1}{Art. 1 Abs. 1: ``Die Wuerde des Menschen ist unantastbar. Sie zu achten und zu schuetzen ist Verpflichtung aller staatlichen Gewalt.'' Art. 2 Abs. 1: ``Jeder hat das Recht auf die freie Entfaltung seiner Persoenlichkeit, soweit er nicht die Rechte anderer verletzt und nicht gegen die verfassungsmaessige Ordnung oder das Sittengesetz verstoesst.''}\ggmarknotede{art-1-1}\ggmarknotede{art-2-1}}
\newcommand{\ggfnArtOneOneTwoOneZH}{\ggnotezhonce{art-1-1+2-1}{第1條第1項：「人之尊嚴不可侵犯。尊重及保護之，為一切國家權力之義務。」第2條第1項：「人人有自由發展其人格之權利，但不得侵害他人權利，亦不得違反合憲秩序或道德律。」}\ggmarknotezh{art-1-1}\ggmarknotezh{art-2-1}}

% ===== 組合腳註：Art. 2(2)(1) + Art. 1(1)（墮胎一案特有，含交叉標記）=====
\newcommand{\ggfnLifeDignityDE}{\ggnotedeonce{life-dignity}{Art. 2 Abs. 2 Satz 1: ``Jeder hat das Recht auf Leben und koerperliche Unversehrtheit.'' Art. 1 Abs. 1: ``Die Wuerde des Menschen ist unantastbar. Sie zu achten und zu schuetzen ist Verpflichtung aller staatlichen Gewalt.''}\ggmarknotede{art-2-2-1}\ggmarknotede{art-1-1}\ggmarknotede{art-1-1-2}}
\newcommand{\ggfnLifeDignityZH}{\ggnotezhonce{life-dignity}{第2條第2項第1句：「人人有生命與身體完整性之權利。」第1條第1項：「人之尊嚴不可侵犯。尊重及保護之，為一切國家權力之義務。」}\ggmarknotezh{art-2-2-1}\ggmarknotezh{art-1-1}\ggmarknotezh{art-1-1-2}}

% ===== 組合腳註：Art. 1(1) + Art. 2(2)（墮胎一案特有，條件判斷版）=====
\newcommand{\ggfnArtOneTwoTwoDE}{\ifcsname ggnoteseen@de@life-dignity\endcsname\ggfnArtTwoTwoRestDE{}\else\ggnotedeonce{art-1-2-2}{Art. 1 Abs. 1: ``Die Wuerde des Menschen ist unantastbar. Sie zu achten und zu schuetzen ist Verpflichtung aller staatlichen Gewalt.'' Art. 2 Abs. 2: ``Jeder hat das Recht auf Leben und koerperliche Unversehrtheit. Die Freiheit der Person ist unverletzlich. In diese Rechte darf nur auf Grund eines Gesetzes eingegriffen werden.''}\ggmarknotede{art-1-1}\ggmarknotede{art-1-1-2}\ggmarknotede{art-2-2}\ggmarknotede{art-2-2-1}\fi}
\newcommand{\ggfnArtOneTwoTwoZH}{\ifcsname ggnoteseen@zh@life-dignity\endcsname\ggfnArtTwoTwoRestZH{}\else\ggnotezhonce{art-1-2-2}{第1條第1項：「人之尊嚴不可侵犯。尊重及保護之，為一切國家權力之義務。」第2條第2項：「人人有生命與身體完整性之權利。人身自由不可侵犯。此等權利僅得依法律加以干預。」}\ggmarknotezh{art-1-1}\ggmarknotezh{art-1-1-2}\ggmarknotezh{art-2-2}\ggmarknotezh{art-2-2-1}\fi}

% ===== 組合腳註：Art. 1(1) + Art. 19(2)（墮胎一案特有，條件判斷版）=====
\newcommand{\ggfnArtOneNineteenDE}{\ifcsname ggnoteseen@de@art-1-1\endcsname\ggfnArtNineteenTwoDE{}\else\ggnotedeonce{art-1-19}{Art. 1 Abs. 1: ``Die Wuerde des Menschen ist unantastbar. Sie zu achten und zu schuetzen ist Verpflichtung aller staatlichen Gewalt.'' Art. 19 Abs. 2: ``In keinem Falle darf ein Grundrecht in seinem Wesensgehalt angetastet werden.''}\ggmarknotede{art-1-1}\ggmarknotede{art-1-1-2}\ggmarknotede{art-19-2}\fi}
\newcommand{\ggfnArtOneNineteenZH}{\ifcsname ggnoteseen@zh@art-1-1\endcsname\ggfnArtNineteenTwoZH{}\else\ggnotezhonce{art-1-19}{第1條第1項：「人之尊嚴不可侵犯。尊重及保護之，為一切國家權力之義務。」第19條第2項：「任何情形下，基本權利之本質內容均不得受到侵害。」}\ggmarknotezh{art-1-1}\ggmarknotezh{art-1-1-2}\ggmarknotezh{art-19-2}\fi}

% 本案 Art. 2(2) 需依前文是否已引 Art. 2(2)(1) 決定是否僅引餘文，覆寫共用版型
\renewcommand{\ggfnArtTwoTwoDE}{\ifcsname ggnoteseen@de@art-2-2-1\endcsname\ggfnArtTwoTwoRestDE{}\else\ggnotedeonce{art-2-2}{Art. 2 Abs. 2: ``Jeder hat das Recht auf Leben und koerperliche Unversehrtheit. Die Freiheit der Person ist unverletzlich. In diese Rechte darf nur auf Grund eines Gesetzes eingegriffen werden.''}\ggmarknotede{art-2-2-1}\fi}
\renewcommand{\ggfnArtTwoTwoZH}{\ifcsname ggnoteseen@zh@art-2-2-1\endcsname\ggfnArtTwoTwoRestZH{}\else\ggnotezhonce{art-2-2}{第2條第2項：「人人有生命與身體完整性之權利。人身自由不可侵犯。此等權利僅得依法律加以干預。」}\ggmarknotezh{art-2-2-1}\fi}
\newcommand{\tocde}{%
  \section*{\mdseries\bverfgetitlede Inhaltsverzeichnis}
  \tocpart{Entscheidungstext}
  \begin{tocitems}
  \item[Leitsätze] Amtliche Leitsätze
  \item[Urteil] Entscheidung
  \item[Tenor] Entscheidungsformel
  \item[Gründe] Begründung
  \end{tocitems}
  \begin{tocsubitems}
  \item[A.] Verfahren und Gegenstand
  \item[B.] Zustimmung des Bundesrates
  \item[C.] Zulässigkeit
  \item[D.] Materielle Prüfung
  \item[E.] Ergebnis
  \end{tocsubitems}
  \begin{tocitems}
  \item[Abw. Meinung] Rupp-v. Brünneck/Simon
  \end{tocitems}
  \tocdashline
  \begin{tocitems}
  \item[Glossar] Terminologie
  \end{tocitems}
}
\newcommand{\toczh}{%
  \section*{\mdseries\bverfgetitlezh 目錄}
  \tocpart{判決原文部分}
  \begin{tocitems}
  \item[裁判要旨] 官方裁判要旨
  \item[判決] 判決
  \item[主文] 判決主文
  \item[理由] 判決理由
  \end{tocitems}
  \begin{tocsubitems}
  \item[A.] 程序與審查標的
  \item[B.] 聯邦參議院同意
  \item[C.] 合法性
  \item[D.] 實體審查
  \item[E.] 結論
  \end{tocsubitems}
  \begin{tocitems}
  \item[不同意見] Rupp-v. Brünneck/Simon
  \end{tocitems}
  \tocdashline
  \begin{tocitems}
  \item[譯語表] 術語對照
  \end{tocitems}
}
% 大綱雙語對照表格內容（由 outlinepage 環境包裝後插入文件）
\newcommand{\outlinepaired}{%
\opartrow{Entscheidungstext}{判決原文部分}%
\omainrow{Leitsätze}{Sechs amtliche Leitsätze: werdende Leben als selbständiges Rechtsgut; staatliche Schutzpflicht auch gegenüber der Mutter; Lebensvorrang für die gesamte Schwangerschaft; Strafrecht als ultima ratio; Zumutbarkeitsgrenzen; Urteilsergebnis.}{裁判要旨}{六條裁判要旨：胎兒生命為獨立法益；國家對母親亦負全面保護義務；生命保護於整個妊娠期間優先；刑法為最後手段；不合理負擔之例外；違憲結論。}%
\omainrow{Urteil}{Erster Senat, 25.\,Feb.\,1975; Parteien und Verfahrensangaben.}{判決}{第一庭，1975年2月25日；當事人及程序標識。}%
\omainrow{Tenor}{§\,218a StGB (5.\,StrRG) mit Art.\,2 Abs.\,2 S.\,1 i.\,V.\,m.\ Art.\,1 Abs.\,1 GG unvereinbar und nichtig; drei Übergangsanordnungen nach § 35 BVerfGG.}{主文}{《第五刑法改革法》版本§\,218a 違憲無效；三項依《聯邦憲法法院法》第35條所為之臨時命令。}%
\omainrow{Gründe}{Begründung der Entscheidung (A–E).}{理由}{判決理由（A–E）。}%
\opartrow{Gliederung der Gründe}{理由之大綱}%
\osecrow{A.}{\textit{Verfahren und Gegenstand}}{A.}{程序與審查標的}%
\ointrow{1}{Gegenstand: Vereinbarkeit der Fristenregelung (§\,218a StGB) mit dem GG.}{標的：期限規定（§\,218a StGB）與《基本法》相容性。}%
\osubrow{2–49}{I.}{Inhalt des 5.\,StrRG; Gesetzgebungsgeschichte § 218 StGB (1871–1974).}{I.}{《第五刑法改革法》內容；§\,218 StGB 立法沿革（1871–1974）。}%
\osubrow{50–67}{II.}{Verfassungsrügen der Antragsteller (193 MdB und 5 Landesregierungen).}{II.}{聲請人之合憲性指摘（193名聯邦議員及5邦政府）。}%
\osubrow{68–97}{III.}{Äußerungen von Bundestag, Bundesregierung und weiteren Beteiligten.}{III.}{聯邦議院、聯邦政府及其他參與者之意見陳述。}%
\osubrow{98–100}{IV.}{Sachverständigenaussagen in der mündlichen Verhandlung.}{IV.}{言詞辯論中鑑定人之陳述。}%
\osecrowrn{101–106}{B.}{\textit{Zustimmung des Bundesrates} — 5.\,StrRG bedurfte nicht der Zustimmung des Bundesrates.}{B.}{聯邦參議院同意 — 《第五刑法改革法》無需聯邦參議院同意。}%
\osecrow{C.}{\textit{Zulässigkeit}}{C.}{合法性}%
\ointrow{107}{Komplexität der Abbruchsproblematik; verfassungsrechtliche Prüfungsmaßstäbe.}{終止妊娠問題之複雜性；憲法審查標準概述。}%
\osubrow{108–123}{I.}{Art.\,2 Abs.\,2 S.\,1 GG schützt das werdende Leben als selbständiges Rechtsgut; staatliche Schutzpflicht aus objektivem Grundrechtsgehalt.}{I.}{《基本法》第2條第2項第1句保護發展中生命為獨立法益；保護義務源自基本權之客觀法內容。}%
\osubsubrow{108–121}{1.}{Art.\,2 Abs.\,2 S.\,1 GG schützt werdende Leben als selbständiges Rechtsgut (fünffache Begründung).}{1.}{第2條第2項第1句保護發展中生命為獨立法益（五重論據）。}%
\osubsubsubrow{109}{a)}{Historische Auslegung: Reaktion auf nationalsozialistische Vernichtungsmaßnahmen.}{a)}{歷史解釋：對國家社會主義消滅生命措施之回應。}%
\osubsubsubrow{110}{b)}{Wortlautauslegung: „jeder" in Art.\,2 Abs.\,2 S.\,1 GG erfasst auch das Ungeborene.}{b)}{文義解釋：第2條第2項第1句「人人」包含未出生者。}%
\osubsubsubrow{111–112}{c)}{Systematische Auslegung: Einwand, „jeder" bezeichne nur Geborene, greift nicht durch.}{c)}{體系解釋：「人人」僅指已出生者之異議不成立。}%
\osubsubsubrow{113–118}{d)}{Entstehungsgeschichte: Beratungen des Parlamentarischen Rates bestätigen Schutz des Ungeborenen.}{d)}{立法原意，歷史解釋（制定史）：議會委員會審議確認保護未出生者之立法意旨。}%
\osubsubsubrow{119–121}{e)}{Gesetzgebungsgeschichte des 5.\,StrRG sowie Sachverständige bestätigen Lebensschutz.}{e)}{《第五刑法改革法》立法史及鑑定意見進一步確認生命保護。}%
\osubsubrow{122}{2.}{Schutzpflicht aus Art.\,2 Abs.\,2 GG unmittelbar ableitbar; objektiv-rechtliche Dimension der Grundrechte.}{2.}{保護義務可直接由第2條第2項導出；基本權之客觀法向度。}%
\osubsubrow{123}{3.}{Frage der Rechtssubjektivität des Nasciturus nicht entscheidungserheblich.}{3.}{胎兒之法律主體資格問題對本案不具決定性。}%
\osubrow{124–127}{II.}{Umfassende staatliche Schutzpflicht auch gegenüber der Mutter; Lebensvorrang für die gesamte Schwangerschaft.}{II.}{國家保護義務對母親亦存在；胎兒生命保護於整個妊娠期間優先。}%
\osubsubrow{124}{1.}{Schutzpflicht umfassend: nicht nur Abwehrrecht, sondern aktives Gebot zum Schutz auch vor Dritten.}{1.}{保護義務全面：不僅為防禦權，亦為積極保護命令，包含對第三人之保護。}%
\osubsubrow{125–126}{2.}{Schutzpflicht besteht grundsätzlich auch gegenüber der Mutter; Vorrang des Lebens des Nasciturus.}{2.}{保護義務原則上亦存在於對母親之關係；胎兒生命權原則上具優先性。}%
\osubsubrow{127}{3.}{Rechtsordnung muss Schwangerschaftsabbruch grundsätzlich als Unrecht kennzeichnen.}{3.}{法律秩序必須原則上將終止妊娠標識為不法。}%
\osubrow{128–139}{III.}{Rechtliche Mißbilligung des Abbruchs geboten; Strafrecht als ultima ratio; Gesetzgeber hat Spielraum bei Wahl der Schutzmittel (Vorrang der Prävention).}{III.}{法律否定評價義務；刑法為最後手段；立法者在保護手段之選擇上有形成空間（預防優先）。}%
\osubsubrow{129}{1.}{Vorrang der Prävention: sozialpolitische und fürsorgerische Mittel haben grundsätzlich Vorrang vor dem Strafrecht.}{1.}{預防優先：社會政策及扶助手段原則上優先於刑法。}%
\osubsubrow{130–134}{2.}{Strafrecht als ultima ratio: geboten, wenn andere Mittel keinen hinreichenden Schutz bieten.}{2.}{刑法為最後手段：於其他手段不能提供充分保護時不可或缺。}%
\osubsubsubrow{132}{a)}{Strafrecht schützt elementare Gemeinschaftswerte; Abbruch als Tötungshandlung rechtlich eindeutig zu kennzeichnen.}{a)}{刑法保護社會基本價值；終止妊娠作為殺戮行為在法律上須明確標識。}%
\osubsubsubrow{133–134}{b)}{Strafe kein Selbstzweck; Verzicht auf strafrechtliche Mißbilligung gleichwohl verfassungswidrig.}{b)}{刑罰非目的本身；但放棄刑法上的否定評價在憲法上仍不被允許。}%
\osubsubrow{135–139}{3.}{Schutzpflicht auch gegenüber der Mutter; Unzumutbarkeitsgrenzen als verfassungsrechtlich gebotene Ausnahmen.}{3.}{保護義務亦存在於對母親；不合理負擔構成憲法上要求之例外界限。}%
\osecrow{D.}{\textit{Materielle Prüfung}}{D.}{實體審查}%
\ointrow{140}{Fristenregelung genügt den verfassungsrechtlichen Schutzanforderungen nicht.}{期限規定未充分履行憲法保護義務。}%
\osubrow{141}{I.}{Prüfungsbefugnis des BVerfG; Spielraum des Gesetzgebers.}{I.}{聯邦憲法法院之審查職權；立法形成空間。}%
\osubrow{142–171}{II.}{Fristenregelung bietet keinen hinreichenden Schutz des werdenden Lebens; Leitgedanke „Hilfe statt Strafe" verfassungsrechtlich unzureichend umgesetzt.}{II.}{期限規定未能提供憲法要求之生命保護；「協助而非刑罰」指導方針在憲法上未充分實現。}%
\osubsubrow{145–150}{1.}{Rechtliche Mißbilligung muss auch im Sozialrecht sichtbar bleiben; Fristenregelung gibt falsche Signale.}{1.}{法律否定評價亦必須在社會法領域保持可見；期限規定傳遞錯誤訊號。}%
\osubsubrow{151–165}{2.}{Formale Mißbilligung ohne effektive Schutzwirkung genügt der Verfassung nicht.}{2.}{不具實際保護效果之形式否定評價不符合憲法要求。}%
\osubsubsubrow{153–155}{a)}{Strafrecht als Träger des allgemeinen Rechtsgutsbewusstseins; Fristenlösung untergräbt dieses.}{a)}{刑法作為一般法益意識之承載；期限規定損害此一效果。}%
\osubsubsubrow{156–158}{b)}{Pauschale Abwägung Leben gegen Leben verkehrt den Schutzgehalt des Art.\,2 Abs.\,2 GG.}{b)}{生命對生命之概括衡量曲解第2條第2項之保護內容。}%
\osubsubsubrow{159–160}{c)}{Gesamtrechnung fehlt verlässliche tatsächliche Grundlage; Rechtsvergleich bestätigt Ineffektivität.}{c)}{總體計算缺乏可靠事實基礎；比較法考察確認刑罰替代之無效性。}%
\osubsubrow{161–171}{3.}{§\,218c Abs.\,1 StGB-Beratungsregelung als Ersatz unzureichend.}{3.}{§\,218c 第1項諮詢規定不足以作為充分替代保護。}%
\osubsubsubrow{166}{a)}{Unterrichtung über Hilfen unzureichend; bloße Informationsvermittlung ohne individuelle Beratung.}{a)}{關於協助措施之告知不足；僅提供資訊而無個別諮詢。}%
\osubsubsubrow{167–169}{b)}{Aufklärung durch abbruchbereiten Arzt nicht zielführend; strukturelles Interessenkonfliktproblem.}{b)}{由願意實施終止妊娠之醫師進行說明不具目標導向；結構性利益衝突問題。}%
\osubsubsubrow{171}{c)}{Fehlende Einbindung in Hilfssystem beeinträchtigt Wirksamkeit der Beratungsregelung.}{c)}{未納入協助體系損害諮詢規定之實際保護效果。}%
\osubrow{172–175}{III.}{Zulässige Indikationen: medizinisch, eugenisch, kriminologisch, allgemeine Notlage.}{III.}{合憲例外事由：醫學、優生、犯罪學及一般困境指標事由。}%
\osubrow{176–178}{IV.}{Übergangsanordnungen nach § 35 BVerfGG.}{IV.}{依《聯邦憲法法院法》第35條所為之臨時命令。}%
\osecrowrn{179–180}{E.}{\textit{Ergebnis} — 5.\,StrRG verletzt die Schutzpflicht; § 218a nichtig; Übergangsregelung gilt bis zur Neuregelung.}{E.}{結論 — §\,218a 違憲無效；臨時規制繼續有效至新立法為止。}%
\odissentpart{Abweichende Meinung}{不同意見}%
\odissentopinion{1–50}{Rupp-v.\,Brünneck / Simon}{Keine verfassungsrechtliche Pönalisierungspflicht; Strafrecht als ungeeignetes und historisch belastetes Schutzmittel.}{Rupp-v.\,Brünneck / Simon}{憲法不課予刑事處罰義務；刑法為不適當且歷史上有重負之保護手段。}%
\osecrow{A.}{\textit{Keine verfassungsrechtliche Pönalisierungspflicht}}{A.}{不課予刑事處罰義務}%
\osubrow{1–9}{I.}{Keine verfassungsrechtliche Pönalisierungspflicht; legislativer Gestaltungsspielraum.}{I.}{憲法不課予刑事處罰義務；立法形成空間。}%
\osubrow{10–17}{II.}{Strafrecht als ungeeignetes und historisch belastetes Schutzmittel.}{II.}{刑法為不適當且歷史上有重負之保護手段。}%
\osecrow{B.}{\textit{Kein Verfassungsverstoß bei unterstellter Strafpflicht}}{B.}{即使承認刑罰義務，立法者亦未違憲}%
\ointrow{18}{Einführung: Mehrheitsbegründung begegnet erheblichen Einwänden.}{導言：多數意見論據面臨重大反論。}%
\osubrow{19–44}{I.}{Singularität des Abbruchsproblems; ineffektive Strafnorm; sozialpolitischer Vorrang.}{I.}{墮胎問題之特殊性；刑罰規範之無效性；社會政策優先。}%
\osubrow{45–48}{II.}{Verfassung schreibt keine „Mißbilligung" ohne tatsächliche Schutzwirkung vor.}{II.}{憲法不要求無實際保護效果之法律「否定評價」。}%
\osubrow{49}{III.}{Rechtsvergleich: Reformen in westlichen Kulturstaaten.}{III.}{比較法考察：西方文化國家之改革。}%
\osubrow{50}{IV.}{Schlußfolgerung.}{IV.}{綜合結論。}%
}
\begin{document}
\pagenumbering{roman}
\hypersetup{pageanchor=false}

\begin{titlepage}
\thispagestyle{empty}
\vspace*{2.2cm}
\ifbilingualcolumns
  \noindent
  \begin{tabular}{@{}>{\centering\arraybackslash}p{0.485\textwidth}|>{\centering\arraybackslash}p{0.485\textwidth}@{}}
  {\bverfgetitlede\LARGE BVerfGE 39, 1\par}
  \vspace{0.35cm}
  {\bverfgetitlede\Large Schwangerschaftsabbruch I\par}
  \vspace{0.8cm}
  {\bverfgetitlede\Large Synoptische Ausgabe\par}
  \vspace{1.4cm}
  {\bverfgebodyde Bundesverfassungsgericht, Erster Senat\par}
  \vspace{0.35cm}
  {\bverfgebodyde Urteil vom 25. Februar 1975\par}
  \vspace{0.35cm}
  {\bverfgebodyde Aktenzeichen: 1 BvF 1, 2, 3, 4, 5, 6/74\par}
  &
  {\bverfgetitlezh\LARGE BVerfGE 39, 1\par}
  \vspace{0.35cm}
  {\bverfgetitlezh\Large Schwangerschaftsabbruch I\par}
  \vspace{0.8cm}
  {\bverfgetitlezh\Large 德中對照本\par}
  \vspace{1.4cm}
  {\bverfgebodyzh 德國聯邦憲法法院第一庭\par}
  \vspace{0.35cm}
  {\bverfgebodyzh 判決，1975年2月25日\par}
  \vspace{0.35cm}
  {\bverfgebodyzh 案號：1 BvF 1, 2, 3, 4, 5, 6/74\par}
  \end{tabular}
\else
  \begin{center}
  {\bverfgetitlede\LARGE BVerfGE 39, 1\par}
  \vspace{0.35cm}
  {\bverfgetitlede\Large Schwangerschaftsabbruch I\par}
  \vspace{0.8cm}
  \ifshowtranslation
    {\bverfgetitlezh\Large 中文譯本\par}
    \vspace{1.4cm}
    {\bverfgebodyzh 德國聯邦憲法法院第一庭\par}
    \vspace{0.35cm}
    {\bverfgebodyzh 判決，1975年2月25日\par}
    \vspace{0.35cm}
    {\bverfgebodyzh 案號：1 BvF 1, 2, 3, 4, 5, 6/74\par}
  \else
    {\bverfgetitlede\Large Amtliche Textfassung\par}
    \vspace{1.4cm}
    {\bverfgebodyde Bundesverfassungsgericht, Erster Senat\par}
    \vspace{0.35cm}
    {\bverfgebodyde Urteil vom 25. Februar 1975\par}
    \vspace{0.35cm}
    {\bverfgebodyde Aktenzeichen: 1 BvF 1, 2, 3, 4, 5, 6/74\par}
  \fi
  \end{center}
\fi
\bverfgetitlemeta
\end{titlepage}

\hypersetup{pageanchor=true}
\vspace*{1.0cm}
\begin{center}
\begin{minipage}{0.92\textwidth}
\ifbilingualcolumns
  \noindent
  \begin{tabular}{@{}p{0.47\textwidth}@{\hspace{0.025\textwidth}}!{\color{notegray}\vrule width 0.12pt}@{\hspace{0.025\textwidth}}p{0.47\textwidth}@{}}
  \tocpanel{\bverfgebodyde}{\tocde}
  &
  \tocpanel{\bverfgebodyzh}{\toczh}
  \end{tabular}
\else
  \tocsingleindent
  \ifshowtranslation
    \tocpanel{\bverfgebodyzh}{\toczh}
  \else
    \tocpanel{\bverfgebodyde}{\tocde}
  \fi
\fi
\end{minipage}
\end{center}
\clearpage
\ifbilingualcolumns
  \begin{outlinepage}\outlinepaired\end{outlinepage}
\fi
\clearpage
\pagenumbering{arabic}
\bilingualstart

\bisection{Leitsätze}{裁判要旨}
\syncleft
\begin{sourcepara}
\leitsatznum{1}Das sich im Mutterleib entwickelnde Leben steht als selbständiges Rechtsgut unter dem Schutz der Verfassung (\abbrArt{} 2 \abbrAbs{} 2 Satz 1, \abbrArt{} 1 \abbrAbs{} 1 \abbrGG{}\ggfnLifeDignityDE{}). Die Schutzpflicht des Staates verbietet nicht nur unmittelbare staatliche Eingriffe in das sich entwickelnde Leben, sondern gebietet dem Staat auch, sich schützend und fördernd vor dieses Leben zu stellen.
\end{sourcepara}
\begin{transpara}
\leitsatznum{1}在母胎內發展中的生命，作為獨立法益，受憲法保護（《基本法》第 2 條第 2 項第 1 句、第 1 條第 1 項\ggfnLifeDignityZH{}）。國家的保護義務不僅禁止國家直接干預發展中的生命，也要求國家以保護並促進的方式站在此一生命之前。
\end{transpara}

\syncleft
\begin{sourcepara}
\leitsatznum{2}Die Verpflichtung des Staates, das sich entwickelnde Leben in Schutz zu nehmen, besteht auch gegenüber der Mutter.
\end{sourcepara}
\begin{transpara}
\leitsatznum{2}國家保護正在發展中生命的義務，亦存在於對母親的關係中。
\end{transpara}

\syncleft
\begin{sourcepara}
\leitsatznum{3}Der Lebensschutz der Leibesfrucht genießt grundsätzlich für die gesamte Dauer der Schwangerschaft Vorrang vor dem Selbstbestimmungsrecht der Schwangeren und darf nicht für eine bestimmte Frist in Frage gestellt werden.
\end{sourcepara}
\begin{transpara}
\leitsatznum{3}胎兒的生命保護，原則上於整個妊娠期間均優先於懷孕婦女的自決權，且不得因設定某一期限而使之受到質疑。
\end{transpara}

\syncleft
\begin{sourcepara}
\leitsatznum{4}Der Gesetzgeber kann die grundgesetzlich gebotene rechtliche Mißbilligung des Schwangerschaftsabbruchs auch auf andere Weise zum Ausdruck bringen als mit dem Mittel der Strafdrohung. Entscheidend ist, ob die Gesamtheit der dem Schutz des ungeborenen Lebens dienenden Maßnahmen einen der Bedeutung des zu sichernden Rechtsgutes entsprechenden tatsächlichen Schutz gewährleistet. Im äußersten Falle, wenn der von der Verfassung gebotene Schutz auf keine andere Weise erreicht werden kann, ist der Gesetzgeber verpflichtet, zur Sicherung des sich entwickelnden Lebens das Mittel des Strafrechts einzusetzen.
\end{sourcepara}
\begin{transpara}
\leitsatznum{4}立法機關也可以透過刑罰威嚇以外的其他方式表達《基本法》所要求的對終止妊娠的法律否定評價。關鍵在於，旨在保護未出生生命的全部措施是否能夠保證與所要保障的法益的重要性相對應的實際保護。在極端情況下，當《基本法》所要求的保護無法透過其他方式實現時，立法機關就有義務運用刑法來保障正在發展中的生命。
\end{transpara}

\syncleft
\begin{sourcepara}
\leitsatznum{5}Eine Fortsetzung der Schwangerschaft ist unzumutbar, wenn der Abbruch erforderlich ist, um von der Schwangeren eine Gefahr für ihr Leben oder die Gefahr einer schwerwiegenden Beeinträchtigung ihres Gesundheitszustandes abzuwenden. Darüber hinaus steht es dem Gesetzgeber frei, andere außergewöhnliche Belastungen für die Schwangere, die ähnlich schwerwiegen, als unzumutbar zu werten und in diesen Fällen den Schwangerschaftsabbruch straffrei zu lassen.
\end{sourcepara}
\begin{transpara}
\leitsatznum{5}如果為了防止懷孕婦女生命受到威脅或健康受到嚴重損害而必須終止妊娠，則繼續妊娠是不合理的。此外，立法機關可以自由地考慮對懷孕婦女造成的其他同樣沉重的不合理負擔，並在這些情況下讓終止妊娠不罰。
\end{transpara}

\syncleft
\begin{sourcepara}
\leitsatznum{6}Das Fünfte Gesetz zur Reform des Strafrechts vom 18. Juni 1974 (\abbrBGBl{} I \abbrS{} 1297) ist der verfassungsrechtlichen Verpflichtung, das werdende Leben zu schützen, nicht in dem gebotenen Umfang gerecht geworden.
\end{sourcepara}
\begin{transpara}
\leitsatznum{6}1974 年 6 月 18 日頒布的《刑法改革第五法》（BGBl. I S. 1297）沒有在《基本法》要求的範圍內履行保護正在發展中生命的義務。
\end{transpara}

\bisection{Urteil}{判決}
\syncleft
\begin{sourcepara}
{\centering
\small des Ersten Senats\par
\vspace{0.18em}
\normalsize vom 25. Februar 1975\par
\vspace{0.18em}
\footnotesize auf die mündliche Verhandlung\par
vom 18./19. November 1974\par}
\end{sourcepara}
\begin{transpara}
{\centering
\small 第一庭\par
\vspace{0.18em}
\normalsize 1975年2月25日判決\par
\vspace{0.18em}
\footnotesize 經1974年11月18日、19日言詞辯論後作成\par}
\end{transpara}

\syncleft
\begin{sourcepara}
{\centering\bfseries -- 1 \abbrBvF{} 1, 2, 3, 4, 5, 6/74 --\par}
\end{sourcepara}
\begin{transpara}
{\centering\bfseries -- 1 BvF 1, 2, 3, 4, 5, 6/74 --\par}
\end{transpara}

\syncleft
\begin{sourcepara}
{\footnotesize
\renewcommand{\arraystretch}{1.08}%
\begin{tabular}{@{}>{\RaggedRight\arraybackslash}p{0.21\linewidth}>{\RaggedRight\arraybackslash}p{0.74\linewidth}@{}}
in den Verfahren &
wegen verfassungsrechtlicher Prüfung des Fünften Gesetzes zur Reform des Strafrechts
(5. \abbrStrRG{}) vom 18. Juni 1974 (\abbrBGBl{} I \abbrS{} 1297).
\tabularnewline[0.36em]
Antragsteller: &
\begin{tabular}[t]{@{}>{\bfseries}p{1.7em}@{}>{\RaggedRight\arraybackslash}p{\dimexpr\linewidth-1.7em\relax}@{}}
I. & 193 Mitglieder des Deutschen Bundestages;\\
II. & die Regierung des Landes Baden-Württemberg;\\
III. & die Regierung des Saarlandes;\\
IV. & die Regierung des Freistaates Bayern;\\
V. & die Regierung des Landes Schleswig-Holstein;\\
VI. & die Regierung des Landes Rheinland-Pfalz.
\end{tabular}
\end{tabular}
\par}
\end{sourcepara}
\begin{transpara}
{\footnotesize
\renewcommand{\arraystretch}{1.08}%
\begin{tabular}{@{}>{\RaggedRight\arraybackslash}p{0.21\linewidth}>{\RaggedRight\arraybackslash}p{0.74\linewidth}@{}}
程序事由 &
審查1974年6月18日第五刑法改革法（5. StrRG）（BGBl. I S. 1297）是否合憲。
\tabularnewline[0.36em]
聲請人 &
\begin{tabular}[t]{@{}>{\bfseries}p{1.7em}@{}>{\RaggedRight\arraybackslash}p{\dimexpr\linewidth-1.7em\relax}@{}}
I. & 德國聯邦議院193名議員；\\
II. & 巴登-符騰堡邦政府；\\
III. & 薩爾邦政府；\\
IV. & 巴伐利亞自由邦政府；\\
V. & 石勒蘇益格-荷爾斯泰因邦政府；\\
VI. & 萊茵蘭-普法爾茨邦政府。
\end{tabular}
\end{tabular}
\par}
\end{transpara}

\bisection{Entscheidungsformel}{主文}
\syncleft
\begin{sourcepara}
{\bfseries I. § 218a des Strafgesetzbuches in der Fassung des Fünften Gesetzes zur Reform des Strafrechts (5. \abbrStrRG{}) vom 18. Juni 1974 (\abbrBundesgesetzbl{} I \abbrS{} 1297) ist mit Artikel 2 Absatz 2 Satz 1 in Verbindung mit Artikel 1 Absatz 1 des Grundgesetzes insoweit unvereinbar und nichtig, als er den Schwangerschaftsabbruch auch dann von der Strafbarkeit ausnimmt, wenn keine Gründe vorliegen, die - im Sinne der Entscheidungsgründe - vor der Wertordnung des Grundgesetzes Bestand haben.}
\end{sourcepara}
\begin{transpara}
{\bfseries I. 1974年6月18日《第五刑法改革法》（5. StrRG，BGBl. I S. 1297）所定《刑法典》第218a條，於其使終止妊娠即使欠缺依本判決理由所示、足以通過基本法價值秩序審查之事由時，仍排除刑事可罰性的範圍內，與基本法第2條第2項第1句連同第1條第1項不相容，並為無效。}
\end{transpara}

\syncleft
\begin{sourcepara}
{\bfseries II. Bis zum Inkrafttreten einer gesetzlichen Neuregelung wird gemäß § 35 des Gesetzes über das Bundesverfassungsgericht angeordnet:}
\end{sourcepara}
\begin{transpara}
{\bfseries II. 於新的法律規定生效前，依《聯邦憲法法院法》第35條命令如下：}
\end{transpara}

\syncleft
\begin{sourcepara}
{\bfseries 1. § 218b und § 219 des Strafgesetzbuches in der Fassung des Fünften Gesetzes zur Reform des Strafrechts (5. \abbrStrRG{}) vom 18. Juni 1974 (\abbrBundesgesetzbl{} I \abbrS{} 1297) sind auch auf Schwangerschaftsabbrüche in den ersten zwölf Wochen seit der Empfängnis anzuwenden.}
\end{sourcepara}
\begin{transpara}
{\bfseries 1. 1974年6月18日《第五刑法改革法》（5. StrRG，\abbrBundesgesetzbl{} I S. 1297）所定《刑法典》第218b條及第219條，亦適用於受孕後最初十二週內之終止妊娠。}
\end{transpara}

\syncleft
\begin{sourcepara}
{\bfseries 2. Der mit Einwilligung der Schwangeren von einem Arzt innerhalb der ersten zwölf Wochen seit der Empfängnis vorgenommene Schwangerschaftsabbruch ist nicht nach § 218 des Strafgesetzbuches strafbar, wenn an der Schwangeren eine rechtswidrige Tat nach den §§ 176 bis 179 des Strafgesetzbuches vorgenommen worden ist und dringende Gründe für die Annahme sprechen, dass die Schwangerschaft auf der Tat beruht.}
\end{sourcepara}
\begin{transpara}
{\bfseries 2. 醫師經懷孕婦女同意，於受孕後最初十二週內所實施之終止妊娠，如該懷孕婦女曾遭受《刑法典》第176條至第179條所定不法行為，且有急迫理由足認該妊娠係基於該行為而發生者，不依《刑法典》第218條處罰。}
\end{transpara}

\syncleft
\begin{sourcepara}
{\bfseries 3. Ist der Abbruch der Schwangerschaft in den ersten zwölf Wochen seit der Empfängnis von einem Arzt mit Einwilligung der Schwangeren vorgenommen worden, um von der Schwangeren die auf andere ihr zumutbare Weise nicht abzuwendende Gefahr eine schwerwiegenden Notlage abzuwenden, so kann das Gericht von einer Bestrafung nach § 218 des Strafgesetzbuches absehen.}
\end{sourcepara}
\begin{transpara}
{\bfseries 3. 醫師經懷孕婦女同意，於受孕後最初十二週內實施終止妊娠，以避免懷孕婦女陷入無法以其他對其可期待之方式排除的嚴重困境之危險者，法院得免依《刑法典》第218條處罰。}
\end{transpara}

\bisection{Gründe}{理由}
\startrandnumbers
\bisubsection{A.}{A.}
\syncleft
\begin{sourcepara}
\sourcern{1}Gegenstand des Verfahrens ist die Frage, ob die sogenannte Fristenregelung des Fünften Strafrechtsreformgesetzes, wonach der Schwangerschaftsabbruch in den ersten zwölf Wochen seit der Empfängnis unter bestimmten Voraussetzungen straffrei bleibt, mit dem Grundgesetz vereinbar ist.
\end{sourcepara}
\begin{transpara}
\transrn{1}本程序之標的，是《第五刑法改革法》所謂期限規定是否與《基本法》相容；依該規定，受孕後最初十二週內之終止妊娠，在一定要件下不受處罰。
\end{transpara}

\bisubsubsection{I.}{I.}
\syncleft
\begin{sourcepara}
\sourcern{2}1. Das Fünfte Gesetz zur Reform des Strafrechts (5. \abbrStrRG{}) vom 18. Juni 1974 (\abbrBGBl{} I \abbrS{} 1297) hat die Strafbarkeit des Schwangerschaftsabbruchs neu geregelt. Die §§ 218 bis 220 \abbrStGB{} sind durch Bestimmungen ersetzt worden, die gegenüber dem bisherigen Rechtszustand hauptsächlich folgende Änderungen enthalten:
\end{sourcepara}
\begin{transpara}
\transrn{2}1. 1974年6月18日《第五刑法改革法》（5. StrRG，BGBl. I S. 1297）重新規定終止妊娠之可罰性。《刑法典》第218條至第220條由新的規定取代；相較於既有法狀態，其主要變更如下：
\end{transpara}

\syncleft
\begin{sourcepara}
\sourcern{3}Bestraft wird zwar grundsätzlich, wer eine Schwangerschaft später als am dreizehnten Tag nach der Empfängnis abbricht (§ 218 \abbrAbs{} 1). Jedoch ist der mit Einwilligung der Schwangeren von einem Arzt vorgenommene Schwangerschaftsabbruch nicht nach § 218 strafbar, wenn seit der Empfängnis nicht mehr als zwölf Wochen verstrichen sind (§ 218a - Fristenregelung). Ferner ist der mit Einwilligung der Schwangeren von einem Arzt nach Ablauf der Zwölfwochenfrist vorgenommene Schwangerschaftsabbruch nicht nach § 218 strafbar, wenn er nach den Erkenntnissen der medizinischen Wissenschaft indiziert ist, entweder um von der Schwangeren eine Gefahr für ihr Leben oder die Gefahr einer schwerwiegenden Beeinträchtigung ihres Gesundheitszustandes abzuwenden, sofern diese nicht auf andere zumutbare Weise abgewendet werden kann (§ 218b \abbrNr{}1 - medizinische Indikation), oder weil dringende Gründe für die Annahme sprechen, daß das Kind wegen einer Erbanlage oder schädlicher Einflüsse vor der Geburt an einer nicht behebbaren Schädigung seines Gesundheitszustandes leiden würde, die so schwer wiegt, daß von der Schwangeren die Fortsetzung der Schwangerschaft nicht verlangt werden kann, und wenn seit der Empfängnis nicht mehr als zweiundzwanzig Wochen verstrichen sind (§ 218b \abbrNr{} 2 - eugenische Indikation). Wer eine Schwangerschaft abbricht, ohne daß die Schwangere zuvor von einer Beratungsstelle oder von einem Arzt sozial und ärztlich beraten worden ist, wird bestraft (§ 218c). Ebenso macht sich strafbar, wer nach Ablauf von zwölf Wochen seit der Empfängnis eine Schwangerschaft abbricht, ohne daß eine zuständige Stelle vorher bestätigt hat, daß die Voraussetzungen des § 218b (medizinische oder eugenische Indikation) vorliegen (§ 219). Die Schwangere selbst wird nicht nach § 218c oder § 219 bestraft.
\end{sourcepara}
\begin{transpara}
\transrn{3}原則上，受孕後第十三日以後終止妊娠者，應受處罰（第218條第1項）。然而，醫師經懷孕婦女同意所實施之終止妊娠，如自受孕起尚未逾十二週，不依第218條處罰（第218a條──期限規定）。此外，十二週期限屆滿後，醫師經懷孕婦女同意所實施之終止妊娠，如依醫學科學知見具有指標事由，亦不依第218條處罰：其一，為避免懷孕婦女生命危險或健康狀態受嚴重損害之危險，且該危險不能以其他對其可期待之方式排除者（第218b條第1款──醫學指標事由）；其二，有急迫理由足認該子女因遺傳因素或出生前有害影響，將罹受無法治癒之健康損害，且其嚴重程度足以使人不得要求懷孕婦女繼續妊娠，並且自受孕起尚未逾二十二週者（第218b條第2款──優生指標事由）。終止妊娠者，如懷孕婦女事前未由諮詢機構或醫師接受社會及醫療諮詢，應受處罰（第218c條）。同樣地，受孕後十二週屆滿後終止妊娠者，如未經主管機關事前確認第218b條之要件（醫學或優生指標事由）存在，亦成立犯罪（第219條）。懷孕婦女本人不依第218c條或第219條處罰。
\end{transpara}

\syncleft
\begin{sourcepara}
\sourcern{4}Im einzelnen haben die für das vorliegende Verfahren wesentlichen Vorschriften des Fünften Strafrechtsreformgesetzes folgenden Wortlaut:
\end{sourcepara}
\begin{transpara}
\transrn{4}詳言之，《第五刑法改革法》中對本程序具有重要性之規定，文字如下：
\end{transpara}

\syncleft
\begin{sourceboxpara}
\begin{lawboxpart}{first}{bverfge39-law-rn4-title218}
\sourcelawtitle{§ 218\\Abbruch der Schwangerschaft}
\end{lawboxpart}
\end{sourceboxpara}
\begin{transboxpara}
\begin{lawboxpart}{first}{bverfge39-law-rn4-title218}
\lawtitle{第218條\\終止妊娠}
\end{lawboxpart}
\end{transboxpara}

\syncleft
\begin{sourceboxpara}
\begin{lawboxpart}{middle}{bverfge39-law-rn5}
\sourcern{5}\sourceabsno{1}{Wer eine Schwangerschaft später als am dreizehnten Tage nach der Empfängnis abbricht, wird mit Freiheitsstrafe bis zu drei Jahren oder mit Geldstrafe bestraft.}
\end{lawboxpart}
\end{sourceboxpara}
\begin{transboxpara}
\begin{lawboxpart}{middle}{bverfge39-law-rn5}
\transrn{5}\absno{1}{受孕後第十三日以後終止妊娠者，處三年以下自由刑或罰金。}
\end{lawboxpart}
\end{transboxpara}

\syncleft
\begin{sourceboxpara}
\begin{lawboxpart}{middle}{bverfge39-law-rn6-abs2}
\sourcern{6}\sourceabsno{2}{Die Strafe ist Freiheitsstrafe von sechs Monaten bis zu fünf Jahren, wenn der Täter}
\end{lawboxpart}
\end{sourceboxpara}
\begin{transboxpara}
\begin{lawboxpart}{middle}{bverfge39-law-rn6-abs2}
\transrn{6}\absno{2}{行為人有下列情形之一者，處六月以上五年以下自由刑：}
\end{lawboxpart}
\end{transboxpara}

\syncleft
\begin{sourceboxpara}
\begin{lawboxpart}{middle}{bverfge39-law-rn6-no1}
\sourcenrno{1}{gegen den Willen der Schwangeren handelt oder}
\end{lawboxpart}
\end{sourceboxpara}
\begin{transboxpara}
\begin{lawboxpart}{middle}{bverfge39-law-rn6-no1}
\nrno{1}{違背懷孕婦女意願而行為；或}
\end{lawboxpart}
\end{transboxpara}

\syncleft
\begin{sourceboxpara}
\begin{lawboxpart}{middle}{bverfge39-law-rn6-no2}
\sourcenrno{2}{leichtfertig die Gefahr des Todes oder einer schweren Gesundheitsschädigung der Schwangeren verursacht.}
\end{lawboxpart}
\end{sourceboxpara}
\begin{transboxpara}
\begin{lawboxpart}{middle}{bverfge39-law-rn6-no2}
\nrno{2}{輕率造成懷孕婦女死亡或健康嚴重受損之危險。}
\end{lawboxpart}
\end{transboxpara}

\syncleft
\begin{sourceboxpara}
\begin{lawboxpart}{middle}{bverfge39-law-rn6-supervision}
\sourcelawline{Das Gericht kann Führungsaufsicht anordnen (§ 68 \abbrAbs{} 1 \abbrNr{} 2).}
\end{lawboxpart}
\end{sourceboxpara}
\begin{transboxpara}
\begin{lawboxpart}{middle}{bverfge39-law-rn6-supervision}
\lawline{法院可以命令保護管束（第68條第1項第2款）。}
\end{lawboxpart}
\end{transboxpara}

\syncleft
\begin{sourceboxpara}
\begin{lawboxpart}{middle}{bverfge39-law-rn6-abs3}
\sourceabsno{3}{Begeht die Schwangere die Tat, so ist die Strafe Freiheitsstrafe bis zu einem Jahr oder Geldstrafe.}
\end{lawboxpart}
\end{sourceboxpara}
\begin{transboxpara}
\begin{lawboxpart}{middle}{bverfge39-law-rn6-abs3}
\absno{3}{懷孕婦女實施該行為者，處一年以下有期徒刑或罰金。}
\end{lawboxpart}
\end{transboxpara}

\syncleft
\begin{sourceboxpara}
\begin{lawboxpart}{middle}{bverfge39-law-rn6-abs4}
\sourceabsno{4}{Der Versuch ist strafbar. Die Frau wird nicht wegen Versuchs bestraft.}
\end{lawboxpart}
\end{sourceboxpara}
\begin{transboxpara}
\begin{lawboxpart}{middle}{bverfge39-law-rn6-abs4}
\absno{4}{未遂犯罰之。婦女不因未遂受處罰。}
\end{lawboxpart}
\end{transboxpara}

\syncleft
\begin{sourceboxpara}
\begin{lawboxpart}{middle}{bverfge39-law-title218a}
\sourcelawtitle{§ 218a\\Straflosigkeit des Schwangerschaftsabbruchs in den ersten zwölf Wochen}
\end{lawboxpart}
\end{sourceboxpara}
\begin{transboxpara}
\begin{lawboxpart}{middle}{bverfge39-law-title218a}
\lawtitle{第218a條\\前十二週內終止妊娠不罰}
\end{lawboxpart}
\end{transboxpara}

\syncleft
\begin{sourceboxpara}
\begin{lawboxpart}{middle}{bverfge39-law-rn7}
\sourcern{7}\sourcelawline{Der mit Einwilligung der Schwangeren von einem Arzt vorgenommene Schwangerschaftsabbruch ist nicht nach § 218 strafbar, wenn seit der Empfängnis nicht mehr als zwölf Wochen verstrichen sind.}
\end{lawboxpart}
\end{sourceboxpara}
\begin{transboxpara}
\begin{lawboxpart}{middle}{bverfge39-law-rn7}
\transrn{7}\lawline{醫師經懷孕婦女同意所實施之終止妊娠，如自受孕起尚未逾十二週，不依第218條處罰。}
\end{lawboxpart}
\end{transboxpara}

\syncleft
\begin{sourceboxpara}
\begin{lawboxpart}{middle}{bverfge39-law-title218b}
\sourcelawtitle{§ 218b\\Indikation zum Schwangerschaftsabbruch nach zwölf Wochen}
\end{lawboxpart}
\end{sourceboxpara}
\begin{transboxpara}
\begin{lawboxpart}{middle}{bverfge39-law-title218b}
\lawtitle{第218b條\\十二週後終止妊娠的指標事由}
\end{lawboxpart}
\end{transboxpara}

\syncleft
\begin{sourceboxpara}
\begin{lawboxpart}{middle}{bverfge39-law-rn8-head}
\sourcern{8}\sourcelawline{Der mit Einwilligung der Schwangeren von einem Arzt nach Ablauf von zwölf Wochen seit der Empfängnis vorgenommene Schwangerschaftsabbruch ist nicht nach § 218 strafbar, wenn nach den Erkenntnissen der medizinischen Wissenschaft}
\end{lawboxpart}
\end{sourceboxpara}
\begin{transboxpara}
\begin{lawboxpart}{middle}{bverfge39-law-rn8-head}
\transrn{8}\lawline{受孕後十二週屆滿後，醫師經懷孕婦女同意所實施之終止妊娠，如依醫學科學知見具有下列情形，不依第218條處罰：}
\end{lawboxpart}
\end{transboxpara}

\syncleft
\begin{sourceboxpara}
\begin{lawboxpart}{middle}{bverfge39-law-rn8-no1}
\sourcenrno{1}{der Schwangerschaftsabbruch angezeigt ist, um von der Schwangeren eine Gefahr für ihr Leben oder die Gefahr einer schwerwiegenden Beeinträchtigung ihres Gesundheitszustandes abzuwenden, sofern die Gefahr nicht auf eine andere für sie zumutbare Weise abgewendet werden kann, oder}
\end{lawboxpart}
\end{sourceboxpara}
\begin{transboxpara}
\begin{lawboxpart}{middle}{bverfge39-law-rn8-no1}
\nrno{1}{終止妊娠係為避免懷孕婦女生命危險或健康狀態受嚴重損害之危險，且該危險不能以其他對其可期待之方式排除；或}
\end{lawboxpart}
\end{transboxpara}

\syncleft
\begin{sourceboxpara}
\begin{lawboxpart}{middle}{bverfge39-law-rn8-no2}
\sourcenrno{2}{dringende Gründe für die Annahme sprechen, daß das Kind infolge einer Erbanlage oder schädlicher Einflüsse vor der Geburt an einer nicht behebbaren Schädigung seines Gesundheitszustandes leiden würde, die so schwer wiegt, daß von der Schwangeren die Fortsetzung der Schwangerschaft nicht verlangt werden kann, und seit der Empfängnis nicht mehr als zweiundzwanzig Wochen verstrichen sind.}
\end{lawboxpart}
\end{sourceboxpara}
\begin{transboxpara}
\begin{lawboxpart}{middle}{bverfge39-law-rn8-no2}
\nrno{2}{有急迫理由足認該子女因遺傳因素或出生前有害影響，將罹受無法治癒之健康損害，且其嚴重程度足以使人不得要求懷孕婦女繼續妊娠，並且自受孕起尚未逾二十二週。}
\end{lawboxpart}
\end{transboxpara}

\syncleft
\begin{sourceboxpara}
\begin{lawboxpart}{middle}{bverfge39-law-title218c}
\sourcelawtitle{§ 218c\\Abbruch der Schwangerschaft ohne Unterrichtung und Beratung der Schwangeren}
\end{lawboxpart}
\end{sourceboxpara}
\begin{transboxpara}
\begin{lawboxpart}{middle}{bverfge39-law-title218c}
\lawtitle{第218c條\\未告知及諮詢懷孕婦女而終止妊娠}
\end{lawboxpart}
\end{transboxpara}

\syncleft
\begin{sourceboxpara}
\begin{lawboxpart}{middle}{bverfge39-law-rn9-abs1}
\sourcern{9}\sourceabsno{1}{Wer eine Schwangerschaft abbricht, ohne daß die Schwangere}
\end{lawboxpart}
\end{sourceboxpara}
\begin{transboxpara}
\begin{lawboxpart}{middle}{bverfge39-law-rn9-abs1}
\transrn{9}\absno{1}{終止妊娠者，若懷孕婦女未}
\end{lawboxpart}
\end{transboxpara}

\syncleft
\begin{sourceboxpara}
\begin{lawboxpart}{middle}{bverfge39-law-rn9-no1}
\sourcenrno{1}{sich wegen der Frage des Abbruchs ihrer Schwangerschaft vorher an einen Arzt oder eine hierzu ermächtigte Beratungsstelle gewandt hat und dort über die zur Verfügung stehenden öffentlichen und privaten Hilfen für Schwangere, Mütter und Kinder unterrichtet worden ist, insbesondere über solche Hilfen, die die Fortsetzung der Schwangerschaft und die Lage von Mutter und Kind erleichtern, und}
\end{lawboxpart}
\end{sourceboxpara}
\begin{transboxpara}
\begin{lawboxpart}{middle}{bverfge39-law-rn9-no1}
\nrno{1}{事前就終止妊娠問題洽詢醫師或受授權之諮詢機構，並於該處獲告知可供懷孕婦女、母親及兒童利用之公私協助，尤其是有助於繼續妊娠並改善母親與兒童處境之協助；且}
\end{lawboxpart}
\end{transboxpara}

\syncleft
\begin{sourceboxpara}
\begin{lawboxpart}{middle}{bverfge39-law-rn9-no2}
\sourcenrno{2}{ärztlich beraten worden ist, wird mit Freiheitsstrafe bis zu einem Jahr oder mit Geldstrafe bestraft, wenn die Tat nicht nach § 218 strafbar ist.}
\end{lawboxpart}
\end{sourceboxpara}
\begin{transboxpara}
\begin{lawboxpart}{middle}{bverfge39-law-rn9-no2}
\nrno{2}{接受醫師諮詢，且該行為不依第218條處罰時，處一年以下有期徒刑或罰金。}
\end{lawboxpart}
\end{transboxpara}

\syncleft
\begin{sourceboxpara}
\begin{lawboxpart}{middle}{bverfge39-law-rn9-abs2}
\sourceabsno{2}{Die Frau, an der der Eingriff vorgenommen wird, ist nicht nach Absatz 1 strafbar.}
\end{lawboxpart}
\end{sourceboxpara}
\begin{transboxpara}
\begin{lawboxpart}{middle}{bverfge39-law-rn9-abs2}
\absno{2}{接受介入之婦女，不依第1項處罰。}
\end{lawboxpart}
\end{transboxpara}

\syncleft
\begin{sourceboxpara}
\begin{lawboxpart}{middle}{bverfge39-law-title219}
\sourcelawtitle{§ 219\\Abbruch der Schwangerschaft ohne Begutachtung}
\end{lawboxpart}
\end{sourceboxpara}
\begin{transboxpara}
\begin{lawboxpart}{middle}{bverfge39-law-title219}
\lawtitle{第219條\\未經鑑定而終止妊娠}
\end{lawboxpart}
\end{transboxpara}

\syncleft
\begin{sourceboxpara}
\begin{lawboxpart}{middle}{bverfge39-law-rn10-abs1}
\sourcern{10}\sourceabsno{1}{Wer nach Ablauf von zwölf Wochen seit der Empfängnis eine Schwangerschaft abbricht, ohne daß eine zuständige Stelle vorher bestätigt hat, daß die Voraussetzungen des § 218b \abbrNr{} 1 oder \abbrNr{} 2 vorliegen, wird mit Freiheitsstrafe bis zu einem Jahr oder mit Geldstrafe bestraft, wenn die Tat nicht nach § 218 strafbar ist.}
\end{lawboxpart}
\end{sourceboxpara}
\begin{transboxpara}
\begin{lawboxpart}{middle}{bverfge39-law-rn10-abs1}
\transrn{10}\absno{1}{受孕後十二週屆滿後終止妊娠者，如未經主管機關事前確認第218b條第1款或第2款之要件存在，且該行為不依第218條處罰，處一年以下自由刑或罰金。}
\end{lawboxpart}
\end{transboxpara}

\syncleft
\begin{sourceboxpara}
\begin{lawboxpart}{last}{bverfge39-law-rn10-abs2}
\sourceabsno{2}{Die Frau, an der der Eingriff vorgenommen wird, ist nicht nach Absatz 1 strafbar.}
\end{lawboxpart}
\end{sourceboxpara}
\begin{transboxpara}
\begin{lawboxpart}{last}{bverfge39-law-rn10-abs2}
\absno{2}{接受介入之婦女，不依第1項處罰。}
\end{lawboxpart}
\end{transboxpara}

\syncleft
\begin{sourcepara}
\sourcern{11}2. Nach bisherigem Recht war die Abtötung der Leibesfrucht einer Schwangeren generell eine strafbare Handlung (§ 218 \abbrStGB{}). Allerdings wurde spätestens seit der Entscheidung des Reichsgerichts vom 11. März 1927 (\abbrRGSt{} 61, 242) von der Rechtsprechung im Falle der sogenannten medizinischen Indikation der Rechtfertigungsgrund des übergesetzlichen Notstands nach den Grundsätzen der Pflichten- und Güterabwägung anerkannt. Danach entfiel die Rechtswidrigkeit der Tat bei einer ernsten, auf andere Weise nicht abwendbaren Gefahr für das Leben oder die Gesundheit der Schwangeren, sofern der Eingriff von einem Arzt mit Einwilligung der Schwangeren nach den Regeln der ärztlichen Kunst vorgenommen wurde. Durch § 14 \abbrAbs{} 1 des Gesetzes zur Verhütung erbkranken Nachwuchses in der Fassung des Änderungsgesetzes vom 26. Juni 1935 (\abbrRGBl{} I \abbrS{} 773) wurden diese Voraussetzungen für die aus medizinischen Gründen zulässige Schwangerschaftsbeendigung gesetzlich festgelegt. Diese Vorschrift galt auch nach 1945 in einzelnen Bundesländern fort; wo sie aufgehoben wurde, waren die darin angeführten Voraussetzungen nach der Entscheidung des Bundesgerichtshofs vom 15. Januar 1952 (\abbrBGHSt{} 2, 111) als Mindestvoraussetzungen für die Zulässigkeit der Schwangerschaftsunterbrechung nach den Grundsätzen des übergesetzlichen Notstandes zu beachten.
\end{sourcepara}
\begin{transpara}
\transrn{11}2. 依既有法律，殺害懷孕婦女之胎兒通常為可罰行為（《刑法典》第218條）。不過，最晚自帝國法院1927年3月11日判決（RGSt 61, 242）以來，判例即於所謂醫學指標事由之情形，依義務與法益衡量之原則，承認超法律緊急避難作為阻卻違法事由。據此，懷孕婦女之生命或健康如有嚴重且不能以其他方式排除之危險，而該介入係由醫師經懷孕婦女同意並依醫術規則實施者，行為之違法性即告消滅。1935年6月26日修正法（RGBl. I S. 773）所定《預防遺傳病後代法》第14條第1項，將基於醫學理由得為終止妊娠之此等要件明文化。該規定於1945年後仍在個別聯邦各邦繼續有效；在其遭廢止之處，依聯邦最高法院1952年1月15日判決（BGHSt 2, 111），其中所列要件仍須作為依超法律緊急避難原則許可終止妊娠之最低要件加以遵守。
\end{transpara}

\syncleft
\begin{sourcepara}
\sourcern{12}3. Die Strafvorschrift des § 218 \abbrStGB{} geht in ihrem Kern auf die §§ 181, 182 des Strafgesetzbuchs für die Preußischen Staaten vom 14. April 1851 (Gesetz-Sammlung \abbrS{} 101) zurück; denn diese Vorschriften dienten als Vorbild für die Regelung im Strafgesetzbuch des Norddeutschen Bundes vom 31. Mai 1870 (Bundes-Gesetzblatt des Norddeutschen Bundes \abbrS{} 197), die wörtlich in das Strafgesetzbuch für das Deutsche Reich vom 15. Mai 1871 (\abbrRGBl{} \abbrS{} 127) übernommen wurde. Die Bestimmung lautete in ihrer ursprünglichen Fassung:
\end{sourcepara}
\begin{transpara}
\transrn{12}3. 《刑法典》第218條刑事條項的核心，可追溯至1851年4月14日《普魯士諸國刑法典》第181、182條（Gesetz-Sammlung S. 101）；因為這些規定是1870年5月31日《北德意志邦聯刑法典》（Bundes-Gesetzblatt des Norddeutschen Bundes S. 197）相關規制之範本，並被逐字納入1871年5月15日《德意志帝國刑法典》（RGBl. S. 127）。該條項原始版本如下：
\end{transpara}

\syncleft
\begin{sourceboxpara}
\begin{lawboxpart}{first}{bverfge39-law-rn12-title218}
\sourcelawtitle{§ 218}
\end{lawboxpart}
\end{sourceboxpara}
\begin{transboxpara}
\begin{lawboxpart}{first}{bverfge39-law-rn12-title218}
\lawtitle{第218條}
\end{lawboxpart}
\end{transboxpara}

\syncleft
\begin{sourceboxpara}
\begin{lawboxpart}{last}{bverfge39-law-rn13}
\sourcern{13}\sourcelawline{Eine Schwangere, welche ihre Frucht vorsätzlich abtreibt oder im Mutterleibe tötet, wird mit Zuchthaus bis zu fünf Jahren bestraft.
Sind mildernde Umstände vorhanden, so tritt Gefängnisstrafe nicht unter sechs Monaten ein.
Dieselben Strafvorschriften finden auf denjenigen Anwendung, welcher mit Einwilligung der Schwangeren die Mittel zu der Abtreibung oder Tötung bei ihr angewendet oder ihr beigebracht hat.}
\end{lawboxpart}
\end{sourceboxpara}
\begin{transboxpara}
\begin{lawboxpart}{last}{bverfge39-law-rn13}
\transrn{13}\lawline{故意終止妊娠或殺死子宮內胎兒的懷孕婦女，將被處以最高五年的重懲役刑（Zuchthaus）。若有情節減輕情形，則處六個月以上一般監禁（Gefängnisstrafe）。同樣的刑法規定適用於任何經懷孕婦女同意使用或向其提供終止妊娠或殺害手段的人。}
\end{lawboxpart}
\end{transboxpara}

\syncleft
\begin{sourcepara}
\sourcern{14}Die Strafvorschrift blieb mehr als 50 Jahre lang unverändert. Erst das Gesetz zur Abänderung des Strafgesetzbuchs vom 18. Mai 1926 (\abbrRGBl{} I \abbrS{} 239) milderte die Strafandrohungen (grundsätzlich Gefängnisstrafe, jedoch Zuchthausstrafe für den gewerbsmäßigen Fremdabtreiber).
\end{sourcepara}
\begin{transpara}
\transrn{14}該刑罰規定五十餘年間未經變更。直到1926年5月18日《刑法典修正法》（RGBl. I S. 239），始減輕其刑罰威嚇（原則上為一般監禁刑，但對業務性為他人墮胎者仍為重懲役刑）。
\end{transpara}

\syncleft
\begin{sourcepara}
\sourcern{15}Durch die Verordnung zur Durchführung der Verordnung zum Schutz von Ehe, Familie und Mutterschaft vom 18. März 1943 (\abbrRGBl{} I \abbrS{} 169) wurden die Strafandrohungen wieder erheblich verschärft.
\end{sourcepara}
\begin{transpara}
\transrn{15}1943年3月18日《婚姻、家庭及母性保護條例施行條例》（RGBl. I S. 169）又大幅加重刑罰威嚇。
\end{transpara}

\syncleft
\begin{sourcepara}
\sourcern{16}Das erste Gesetz zur Reform des Strafrechts (1. \abbrStrRG{}) vom 25. Juni 1969 (\abbrBGBl{} I \abbrS{} 645) milderte die Strafdrohung bei Selbstabtreibung, indem der besonders schwere Fall entfiel. Die Fremdabtreibung wurde zu einem Vergehen herabgestuft.
\end{sourcepara}
\begin{transpara}
\transrn{16}1969年6月25日《第一刑法改革法》（1. StrRG，BGBl. I S. 645）刪除特別嚴重情節，因而減輕自行墮胎之刑罰威嚇；他人墮胎則降為輕罪。
\end{transpara}

\syncleft
\begin{sourcepara}
\sourcern{17}4. a) Das generelle Abtreibungsverbot war von Anfang an Gegenstand von Angriffen. Insbesondere nach der Jahrhundertwende begann auch innerhalb der Rechtswissenschaft eine lebhafte Diskussion über die Strafwürdigkeit der Abtreibung. Die Zweifel setzten schon bei der Frage ein, welches Rechtsgut durch das Abtreibungsverbot geschützt werden solle. Aber auch die Zulassung von Ausnahmen von dem uneingeschränkten Verbot und die richtige Bemessung der anzudrohenden Strafen wurden erörtert.
\end{sourcepara}
\begin{transpara}
\transrn{17}4. a) 一般性墮胎禁止自始即受到批判。尤其世紀之交以後，法學界內部亦開始熱烈討論墮胎是否具有應罰性。疑義首先出現在墮胎禁令究竟應保護何種法益的問題上；此外，對無限制禁令是否容許例外，以及應如何正確衡量所威嚇之刑罰，也均有所討論。
\end{transpara}

\syncleft
\begin{sourcepara}
\sourcern{18}Der 1909 vom Reichsjustizamt veröffentlichte Vorentwurf für ein deutsches Strafgesetzbuch, der Gegenentwurf der Professoren Kahl, v. Lilienthal, Franz v. Liszt und Goldschmidt von 1911, der Entwurf nach den Beschlüssen der Strafrechtskommission von 1913 und der Entwurf von 1919 wollten lediglich die Strafandrohung mildern. In der Begründung zum Entwurf von 1919 wurde die Straflosigkeit des Schwangerschaftsabbruchs "im Hinblick auf die schweren Schädigungen, die aus dem Umsichgreifen der Abtreibung für das Volkswohl erwachsen", abgelehnt.
\end{sourcepara}
\begin{transpara}
\transrn{18}帝國司法局1909年公布之德國刑法典預備草案、Kahl、v. Lilienthal、Franz v. Liszt 及 Goldschmidt 等教授1911年提出之對案、依1913年刑法委員會決議所作草案，以及1919年草案，均僅欲減輕刑罰威嚇。1919年草案理由以「鑑於墮胎蔓延對民族福祉所生之嚴重損害」為由，拒絕使終止妊娠不罰。
\end{transpara}

\syncleft
\begin{sourcepara}
\sourcern{19}b) Während der Weimarer Republik war im Rahmen der - allerdings nicht zum Abschluß gekommenen - Arbeiten an einer großen Strafrechtsreform auch § 218 \abbrStGB{} Gegenstand lebhafter Erörterungen. Eine große Zahl von Gesetzesinitiativen und Interpellationen zur Reform dieser Strafvorschrift wurden im Reichstag eingebracht. Sie verfolgten teils das Ziel, die §§ 218 bis 220 \abbrStGB{} ersatzlos zu streichen, teils schlugen sie vor, die Strafandrohung für die ersten drei Schwangerschaftsmonate zurückzunehmen. Ein Antrag, den Frau Schuch, Professor Radbruch und 53 andere Mitglieder der SPD-Fraktion am 31. Juli 1920 im Reichstag stellten, sah die Straflosigkeit der Abtreibung vor, "wenn sie von der Schwangeren oder einem staatlich anerkannten (approbierten) Arzt innerhalb der ersten drei Monate der Schwangerschaft vorgenommen worden" ist (vgl. auch zur Begründung Radbruch in: Grotjahn-Radbruch, Die Abtreibung der Leibesfrucht, 1921). Sämtliche Anträge hatten letztlich keinen Erfolg.
\end{sourcepara}
\begin{transpara}
\transrn{19}b) 魏瑪共和國時期，在大型刑法改革之工作脈絡中，雖然該改革最終未能完成，《刑法典》第218條仍成為熱烈討論之對象。國會提出大量改革此一刑罰規定之立法倡議與質詢；其部分目的在於無替代地刪除《刑法典》第218條至第220條，部分則建議撤回妊娠最初三個月之刑罰威嚇。Schuch 女士、Radbruch 教授及社民黨黨團其他53名成員於1920年7月31日在帝國議會提出之動議，主張墮胎「如係由懷孕婦女或國家承認（取得醫師資格）之醫師於妊娠最初三個月內實施」即不罰（理由亦參 Radbruch in: Grotjahn-Radbruch, Die Abtreibung der Leibesfrucht, 1921）。所有動議最終均未成功。
\end{transpara}

\syncleft
\begin{sourcepara}
\sourcern{20}Der Entwurf eines Allgemeinen Deutschen Strafgesetzbuchs, den Radbruch als damaliger Reichsjustizminister im Jahre 1922 vorlegte und der zur Grundlage der weiteren Strafrechtsreformarbeiten wurde, drohte für die Abtreibung Gefängnisstrafe an (vgl. "Gustav Radbruchs Entwurf eines Allgemeinen Deutschen Strafgesetzbuches [1922]", Tübingen 1952, \abbrS{} 28, § 225). Ähnliche Regelungen sahen die Entwürfe von 1925, 1927 und 1930 vor.
\end{sourcepara}
\begin{transpara}
\transrn{20}Radbruch 於1922年以當時帝國司法部長身分提出之《一般德國刑法典草案》，後來成為後續刑法改革工作之基礎；該草案對墮胎威嚇一般監禁刑（參見「Gustav Radbruchs Entwurf eines Allgemeinen Deutschen Strafgesetzbuches [1922]」，Tübingen 1952, S. 28, § 225）。1925年、1927年及1930年之草案亦有類似規定。
\end{transpara}

\syncleft
\begin{sourcepara}
\sourcern{21}c) In der nationalsozialistischen Zeit wurde die Abtreibung vornehmlich unter den Gesichtspunkten "Schutz der Volkskraft", "Angriffe auf die Lebenskraft des Volkes", "Angriffe auf Rasse und Erbgut" gesehen. Die Bestrebungen zielten, abgesehen von einigen Indikationsfällen, die straflos bleiben sollten (vgl. auch § 10a des Gesetzes zur Verhütung des erbkranken Nachwuchses), auf strenge Strafen.
\end{sourcepara}
\begin{transpara}
\transrn{21}c) 國家社會主義時期，墮胎主要從「保護民族力量」、「對民族生命力之攻擊」、「對種族及遺傳資質之攻擊」等觀點加以理解。相關努力除少數應維持不罰之指標事由情形外（亦參《預防遺傳病後代法》第10a條），均以嚴厲處罰為目標。
\end{transpara}

\syncleft
\begin{sourcepara}
\sourcern{22}5. Als vorläufiges Ergebnis der nach 1945 wiederaufgenommenen Reformarbeiten am Strafgesetzbuch stellte das Bundesjustizministerium im Jahre 1960 auf der Grundlage der Beschlüsse der Großen Strafrechtskommission einen Gesamtentwurf mit Begründung zusammen, der auch die Beratungen einer Länderkommission verwertete (Entwurf eines Strafgesetzbuches [\abbrStGB{}] - E 1960 - mit Begründung, Bonn 1960).
\end{sourcepara}
\begin{transpara}
\transrn{22}5. 作為1945年後重啟《刑法典》改革工作之初步成果，聯邦司法部於1960年依大刑法委員會之決議，彙整一份附理由之總草案，並採納一邦委員會之審議成果（Entwurf eines Strafgesetzbuches [StGB] - E 1960 - mit Begründung, Bonn 1960）。
\end{transpara}

\syncleft
\begin{sourcepara}
\sourcern{23}Danach blieb die Abtreibung grundsätzlich strafbar (§§ 140, 141); jedoch wurde für den Fall der medizinischen Indikation (§ 157) Straffreiheit gewährt. Ferner sollte nach § 160 die Strafbarkeit einer von einem Arzt mit Einwilligung der Schwangeren vorgenommenen Abtötung der Leibesfrucht entfallen, wenn das Gericht festgestellt hat, daß jemand an der Frau eine Notzucht, eine schwere Schändung, während sie geisteskrank, willenlos, bewußtlos oder zum Widerstand körperlich unfähig war, begangen hat oder ihr Samen von einem anderen als ihrem Ehemann ohne ihre Einwilligung übertragen hat und dringende Gründe für die Annahme sprechen, daß die Schwangerschaft auf der Tat beruht, sofern seit dem Ende des Monats, in den der Beginn der Schwangerschaft fällt, nicht mehr als zwölf Wochen verstrichen sind.
\end{sourcepara}
\begin{transpara}
\transrn{23}依此，墮胎原則上仍受處罰（§§ 140, 141）；但於醫學指標事由之情形（§ 157），則給予不罰。此外，依第160條，若法院認定：有人對該婦女實施強姦、嚴重猥褻，或趁其精神疾病、無意志、無意識或身體上不能抗拒時實施行為，或未經其同意將其丈夫以外之人精液移入其體內，且有急迫理由足以認為妊娠係基於該行為而發生，則由醫師經懷孕婦女同意所實施之殺害胎兒行為應不具可罰性；但以自妊娠開始所屬月份之終了起，尚未經過十二週為限。
\end{transpara}

\syncleft
\begin{sourcepara}
\sourcern{24}Diese ethische Indikation war jedoch in der von der Bundesregierung den gesetzgebenden Körperschaften vorgelegten Fassung des Entwurfs 1960 nicht mehr enthalten (vgl. \abbrBRDrucks{} 270/60, \abbrS{} 38 und \abbrS{} 278). Zu einer Beratung dieses Entwurfs kam es in der 3. Wahlperiode des Bundestages nicht mehr.
\end{sourcepara}
\begin{transpara}
\transrn{24}然而，聯邦政府提交各立法機關之1960年草案版本，已不再包含此一倫理指標事由（參見 BRDrucks. 270/60, S. 38 und S. 278）。該草案於聯邦議院第3屆任期內未再進入審議。
\end{transpara}

\syncleft
\begin{sourcepara}
\sourcern{25}Im Jahre 1962 wurde den Gesetzgebungsorganen ein neuer Regierungsentwurf zugeleitet, der die §§ 140, 141, 157 des Entwurfs 1960 im wesentlichen unverändert übernommen hatte (vgl. \abbrBRDrucks{} 200/62, \abbrS{} 35/36, 38). Auch dieser Entwurf konnte in der 4. Wahlperiode des Bundestages nicht verabschiedet werden.
\end{sourcepara}
\begin{transpara}
\transrn{25}1962年，新的政府草案送交各立法機關；該草案基本上未經變更地承繼1960年草案第140條、第141條及第157條（參見 BRDrucks. 200/62, S. 35/36, 38）。此一草案亦未能於聯邦議院第4屆任期內通過。
\end{transpara}

\syncleft
\begin{sourcepara}
\sourcern{26}Im November 1965 wurde der Entwurf 1962 von einer Gruppe von Bundestagsabgeordneten als Initiativgesetzentwurf eingebracht (\abbrBTDrucks{} V/32). Der Bundestag überwies den Entwurf an den Sonderausschuß für die Strafrechtsreform, der unter Berücksichtigung des von deutschen und schweizerischen Strafrechtslehrern im Jahre 1966 veröffentlichten sogenannten Alternativ-Entwurfs (AE) eines Strafgesetzbuchs (Allgemeiner Teil) zwei Teilentwürfe zur Reform des Strafrechts vorlegte, die vom Bundestag im Jahre 1969 als Erstes und Zweites Gesetz zur Reform des Strafrechts verabschiedet wurden. Durch das erste Gesetz erhielten die Strafandrohungen des § 218 \abbrStGB{} den bereits erwähnten milderen Strafrahmen. Man war sich bei der Beratung im Strafrechtssonderausschuß jedoch darin einig, daß mit der Korrektur des Strafmaßes die Problematik des § 218 \abbrStGB{} nicht gelöst sei, daß vielmehr eine umfassende Reform dieses Bereichs erfolgen müsse (vgl. die Ausführungen des Abgeordneten \abbrDr{} Müller-Emmert, Deutscher Bundestag, 5. \abbrWp{}, 144. Sitzung des Sonderausschusses für die Strafrechtsreform, \abbrStenBer{} \abbrS{} 3195).
\end{sourcepara}
\begin{transpara}
\transrn{26}1965年11月，一群聯邦議院議員將1962年草案作為主動提出之法律草案提交（BTDrucks. V/32）。聯邦議院將該草案交付刑法改革特別委員會；該委員會納入德國及瑞士刑法學者於1966年公布之所謂《刑法典替代草案》（AE，總則）後，提出兩份刑法改革部分草案，並由聯邦議院於1969年作為《第一刑法改革法》及《第二刑法改革法》通過。透過第一法，《刑法典》第218條之刑罰威嚇取得前述較輕刑度。然而，刑法特別委員會審議時一致認為，僅修正刑度並未解決《刑法典》第218條之問題，毋寧必須對此領域作全面改革（參見 Müller-Emmert 議員之發言，Deutscher Bundestag, 5. Wp., 144. Sitzung des Sonderausschusses für die Strafrechtsreform, StenBer. S. 3195）。
\end{transpara}

\syncleft
\begin{sourcepara}
\sourcern{27}Von dieser Auffassung ging auch der im Frühjahr 1970 veröffentlichte Besondere Teil des Alternativ-Entwurfs aus (AE, Bes. Teil, Straftaten gegen die Person, Erster Halbband, Tübingen 1970, \abbrS{} 25 \abbrFF{}). Schwangerschaftsabbrüche seien zwar - so wurde dargelegt - fast ausnahmslos verboten und strafbar. Die soziale Wirklichkeit habe sich jedoch von diesen Rechtsnormen so weit entfernt, daß die Strafandrohungen kaum noch eine Wirkung ausüben könnten. Dies sei in hohem Maße schädlich und unzureichend; denn die Vernichtung werdenden Lebens sei - von Ausnahmesituationen abgesehen - nicht nur ethisch verwerflich, sondern stelle auch die Vernichtung eines Rechtsgutes dar. Über die gesetzgeberischen Maßnahmen, mit denen ein wirksamer Schutz des werdenden Lebens zu erreichen sei, konnten sich die Verfasser indes nicht einigen.
\end{sourcepara}
\begin{transpara}
\transrn{27}1970年春公布之《替代草案》分則亦以此見解為出發點（AE, Bes. Teil, Straftaten gegen die Person, Erster Halbband, Tübingen 1970, S. 25 ff.）。其說明指出，終止妊娠固然幾乎毫無例外地受到禁止並具有可罰性；然而，社會現實已與此等法規範相距甚遠，以致刑罰威嚇幾乎無法再發揮作用。此一狀態在高度上有害且不足；因為毀滅正在形成之生命，除例外情況外，不僅在倫理上可非難，也構成對法益之毀滅。惟作者們無法就何種立法措施足以有效保護正在形成之生命達成一致。
\end{transpara}

\syncleft
\begin{sourcepara}
\sourcern{28}Die Mehrheit entschied sich dafür, in den ersten drei Monaten der Schwangerschaft den Abbruch straffrei zu lassen, und zwar innerhalb von vier Wochen nach der Empfängnis generell, im zweiten und dritten Monat unter der Voraussetzung, daß der Abbruch von einem Arzt vorgenommen wird, nachdem die Schwangere eine Beratungsstelle aufgesucht hat (§ 105 AE). Dazu wurden folgende Erwägungen vorgetragen:
\end{sourcepara}
\begin{transpara}
\transrn{28}多數意見決定，妊娠最初三個月內之終止妊娠應不罰；其中受孕後四週內一律不罰，第二及第三個月則以懷孕婦女曾尋求諮詢機構協助，且由醫師實施終止妊娠為要件（§ 105 AE）。其提出下列考量：
\end{transpara}

\syncleft
\begin{sourcepara}
\sourcern{29}\sourcequoteinner{"Der Grundgedanke dieses Vorschlags ist der, daß dem Entschluß einer Frau zur Schwangerschaftsunterbrechung und seiner Verwirklichung nur dadurch entgegengewirkt werden kann, daß der Schwangeren in den Grenzen des Möglichen Hilfe bei der Behebung der materiellen, sozialen und familiären Schwierigkeiten gewährt wird, die sie zur Schwangerschaftsunterbrechung drängen, und daß ihr durch eine persönliche Beratung und offene Aussprache eine überlegte und verantwortliche Entscheidung ermöglicht wird. Diesem Zweck soll die Institution der Beratungsstellen dienen, ... .\quoteparbreak
Die Beratungsstellen sollen deshalb die Möglichkeit besitzen, finanzielle, soziale und familiäre Hilfe zu leisten. Sie sollen ferner der Schwangeren und ihren Angehörigen durch geeignete Mitarbeiter seelische Betreuung gewähren und dabei den Beteiligten namentlich verdeutlichen, daß die Schwangerschaftsunterbrechung auch medizinisch keine Bagatelle darstellt, sondern ein schwerwiegender und u.U. folgenreicher Eingriff ist und daß die Unterbrechung selbst bei noch so drängenden Motiven eben die Vernichtung werdenden Lebens darstellt, also eine hohe ethische Verantwortung berührt und verletzt" (\abbrAAO{}, \abbrS{} 27)."}
\end{sourcepara}
\begin{transpara}
\transrn{29}\transquoteinner{「此項建議之基本思想在於：對婦女終止妊娠之決意及其實現，唯有在可能範圍內協助懷孕婦女排除迫使其終止妊娠之物質、社會及家庭困難，並透過個別諮詢與坦誠對話，使其得以作成經審慎考慮且負責任之決定，始能加以對抗。諮詢機構此一制度即應服務於此目的，……。\quoteparbreak
因此，諮詢機構應具備提供財務、社會及家庭協助之可能性。其並應透過適當工作人員，對懷孕婦女及其親屬提供心理照護，並特別向參與者闡明：終止妊娠即使在醫學上亦非小事，而是一項重大且在一定情形下後果深遠之介入；且此種中止，即使基於再迫切之動機，終究仍是毀滅正在形成之生命，因而觸及並侵害高度倫理責任」（a.a.O., S. 27）。}
\end{transpara}

\syncleft
\begin{sourcepara}
\sourcern{30}Die zum Schwangerschaftsabbruch neigende Frau solle die Beratungsstelle aufsuchen können, ohne befürchten zu müssen, daß ihr damit die Verwirklichung ihrer Absicht rechtlich unmöglich gemacht werde. Eine später als drei Monate nach der Empfängnis vorgenommene Schwangerschaftsunterbrechung solle nur bei Vorliegen medizinischer oder eugenischer Indikation straflos sein, wobei die Voraussetzungen dieser Indikationen durch eine ärztliche Gutachterstelle festzustellen seien (§ 106 AE).
\end{sourcepara}
\begin{transpara}
\transrn{30}傾向終止妊娠之婦女應能尋求諮詢機構協助，而無須擔心此舉將在法律上使其意圖無法實現。受孕後三個月以後所實施之終止妊娠，則僅於具備醫學或優生指標事由時始應不罰；此等指標事由之要件，並應由醫學鑑定機構認定（§ 106 AE）。
\end{transpara}

\syncleft
\begin{sourcepara}
\sourcern{31}Die Minderheit sah in diesem Vorschlag keinen wirksamen, sondern allenfalls einen mittelbaren Schutz des werdenden Lebens, der in all den Fällen aufgegeben sei, in denen die Schwangere von der Beratungsstelle nicht überzeugt werden könne. Deshalb hielten die Vertreter dieser Meinung grundsätzlich an der Strafbarkeit des Schwangerschaftsabbruchs mit Ausnahme der ersten vier Wochen fest, sahen jedoch Straflosigkeit vor, "wenn der Schwangeren die Austragung der Schwangerschaft unter Berücksichtigung der gesamten Lebensumstände nicht zumutbar ist", und konkretisierten diese Generalklausel durch einen Katalog von fünf Indikationen. Weitere Voraussetzung der Straflosigkeit sollte sein, daß die Einwilligung der Schwangeren sowie die Genehmigung einer ärztlichen Gutachterstelle vorliege und daß der Abbruch in den ersten drei Monaten seit der Empfängnis vorgenommen werde.
\end{sourcepara}
\begin{transpara}
\transrn{31}少數意見認為，此一建議並非對正在形成之生命之有效保護，至多僅為間接保護；凡諮詢機構無法說服懷孕婦女之情形，保護即告放棄。因此，此一見解之代表原則上堅持終止妊娠具有可罰性，惟最初四週除外；但其同時規定，若「考量懷孕婦女全部生活情況後，不能期待其將妊娠懷至足月」，則不罰，並以五項指標事由之列舉具體化此一概括條款。不罰之其他要件，則應為懷孕婦女之同意、醫學鑑定機構之許可，以及終止妊娠係於受孕後最初三個月內實施。
\end{transpara}

\syncleft
\begin{sourcepara}
\sourcern{32}6. Anfang 1972 legte die Bundesregierung den Entwurf eines Fünften Gesetzes zur Reform des Strafrechts (5. \abbrStrRG{}) vor (\abbrBRDrucks{} 58/72). Er hielt an der grundsätzlichen Strafbarkeit des Schwangerschaftsabbruchs fest. In der Begründung wurde ausgeführt, menschliches Leben sei auch vor der Geburt ein schutzwürdiges und schutzbedürftiges Rechtsgut. Das Grundgesetz habe in den \abbrArt{} 1 und \abbrArt{} 2 \abbrAbs{} 2\ggfnArtOneTwoTwoDE{} eine Wertentscheidung für das Leben getroffen. Bei der Reform der Abtreibungsvorschriften handele es sich demnach nicht um die Beseitigung von Straftatbeständen, die kein sozialschädliches Verhalten zum Gegenstand hätten. Vielmehr müsse eine an der Grundrechtsordnung orientierte Strafrechtsreform die Vorschriften über den Schwangerschaftsabbruch so ausgestalten, daß der Schutz des werdenden Lebens nach Lage der Dinge am ehesten gewährleistet sei. Zu diesem Zweck müsse die Neuregelung dem Grundsatz der Unantastbarkeit des werdenden Lebens gerecht werden, gleichzeitig aber einen Ausgleich zwischen dem Recht des ungeborenen Kindes und der Menschenwürde der Schwangeren sowie ihrem Recht auf freie Persönlichkeitsentfaltung herstellen. Dabei könne weder dem einen noch dem anderen Recht ein absoluter Vorrang eingeräumt werden. Es komme darauf an, für besondere, schwerwiegende Konfliktsituationen Lösungen zu finden, die den Wertentscheidungen der Verfassung Rechnung trügen (\abbrAAO{}, \abbrS{} 8).
\end{sourcepara}
\begin{transpara}
\transrn{32}6. 1972年初，聯邦政府提出《第五刑法改革法》草案（5. StrRG，BRDrucks. 58/72）。該草案維持終止妊娠原則上具有可罰性。其理由指出，人的生命即使在出生前，亦為值得保護且需要保護之法益。《基本法》於第1條及第2條第2項\ggfnArtOneTwoTwoZH{}已對生命作成價值決定。因此，墮胎規定之改革並非廢除以非社會有害行為為對象之犯罪構成要件。毋寧，依基本權秩序取向之刑法改革，必須如此形塑終止妊娠之規定，使正在形成之生命依事物狀態最可能獲得保障。為此，新規定必須符合正在形成之生命不可侵犯之原則，同時也須在未出生子女之權利與懷孕婦女之人性尊嚴及其自由發展人格之權利間建立衡平。在此，任一權利均不得享有絕對優先。關鍵在於，為特殊且嚴重之衝突情境尋得符合憲法價值決定之解決方案（a.a.O., S. 8）。
\end{transpara}

\syncleft
\begin{sourcepara}
\sourcern{33}Die Fristenregelung könnte die mit ihr verbundenen gesundheitspolitischen Erwartungen nur erfüllen, wenn jeder Schwangerschaftsabbruch in den ersten drei Monaten als rechtlich gebilligt erscheine. Das sei mit der Wertordnung der Verfassung unvereinbar. Wenn die Gesellschaft das werdende Leben als schutzwürdiges Rechtsgut von vergleichsweise hohem Rang anerkenne, dann könne sie nicht, ohne in Widerspruch zu dieser Prämisse zu geraten, die Vernichtung dieses Rechtsgutes von dem freien Belieben des Einzelnen abhängig machen (\abbrAAO{}, \abbrS{} 9).
\end{sourcepara}
\begin{transpara}
\transrn{33}期限規定唯有在最初三個月內之每一次終止妊娠均呈現為法律上被肯認時，始可能滿足與其相連之衛生政策期待。此與憲法之價值秩序不相容。若社會承認正在形成之生命係值得保護且位階相對較高之法益，則其不能在不違反此一前提之情況下，將此一法益之毀滅委諸個人任意（a.a.O., S. 9）。
\end{transpara}

\syncleft
\begin{sourcepara}
\sourcern{34}Hiervon ausgehend lehnte der Entwurf die "in der Zeit der Weimarer Republik und neuerdings wieder lebhaft erörterte Fristenlösung" ab und entschied sich dafür, daß Ausnahmen von dem grundsätzlichen Verbot des Schwangerschaftsabbruchs nur beim Vorliegen einer gesetzlichen Indikation gelten sollten. Eine Indikation sollte angenommen werden,
\end{sourcepara}
\begin{transpara}
\transrn{34}基於此，草案拒絕「魏瑪共和國時期以及近來又被熱烈討論之期限解決方案」，並決定：對終止妊娠原則禁止之例外，僅於具備法定指標事由時始得成立。指標事由應於下列情形認定：
\end{transpara}

\syncleft
\begin{sourcepara}
\sourcern{35}a) wenn der Abbruch der Schwangerschaft nach den Erkenntnissen der medizinischen Wissenschaft angezeigt ist, um von der Schwangeren eine Gefahr für ihr Leben oder die Gefahr einer schwerwiegenden Beeinträchtigung ihres Gesundheitszustandes abzuwenden, sofern die Gefahr nicht auf eine andere für sie zumutbare Weise abgewendet werden kann (§ 219 - medizinische Indikation);
\end{sourcepara}
\begin{transpara}
\transrn{35}a) 依醫學科學知見，終止妊娠係為避免懷孕婦女生命危險或健康狀態受嚴重損害之危險所必要，且該危險不能以其他對其可期待之方式排除者（第219條──醫學指標事由）；
\end{transpara}

\syncleft
\begin{sourcepara}
\sourcern{36}b) wenn nach den Erkenntnissen der medizinischen Wissenschaft dringende Gründe für die Annahme sprechen, daß das Kind infolge einer Erbanlage oder infolge schädlicher Einflüsse vor der Geburt an einer nicht behebbaren Schädigung seines Gesundheitszustandes leiden würde, die so schwer wiegt, daß von der Schwangeren die Fortsetzung der Schwangerschaft nicht verlangt werden kann, sofern seit dem Beginn der Schwangerschaft nicht mehr als zwanzig Wochen verstrichen sind (§ 219b - eugenische oder kindliche Indikation);
\end{sourcepara}
\begin{transpara}
\transrn{36}b) 依醫學科學知見，有急迫理由足認該子女因遺傳因素或出生前有害影響，將罹受無法治癒之健康損害，且其嚴重程度足以使人不得要求懷孕婦女繼續妊娠，並且自妊娠開始起尚未逾二十週者（§ 219b──優生或子女指標事由）；
\end{transpara}

\syncleft
\begin{sourcepara}
\sourcern{37}c) wenn an der Schwangeren eine rechtswidrige Tat nach § 176 (sexueller Mißbrauch von Kindern), § 177 (Vergewaltigung) oder § 179 \abbrAbs{} 1 (sexueller Mißbrauch Widerstandsunfähiger) vorgenommen worden ist und dringende Gründe für die Annahme sprechen, daß die Schwangerschaft auf der Tat beruht, sofern seit dem Beginn der Schwangerschaft nicht mehr als zwölf Wochen verstrichen sind (§ 219c - ethische oder kriminologische Indikation);
\end{sourcepara}
\begin{transpara}
\transrn{37}c) 懷孕婦女曾遭受第176條（對兒童之性虐待）、第177條（強姦）或第179條第1項（對不能抗拒者之性虐待）所定不法行為，且有急迫理由足認妊娠係基於該行為而發生，並且自妊娠開始起尚未逾十二週者（第219c條──倫理或犯罪學指標事由）；
\end{transpara}

\syncleft
\begin{sourcepara}
\sourcern{38}d) wenn der Abbruch der Schwangerschaft angezeigt ist, um von der Schwangeren die Gefahr einer schwerwiegenden Notlage abzuwenden, sofern die Gefahr nicht auf eine andere für die Schwangere zumutbare Weise abgewendet werden kann und wenn seit dem Beginn der Schwangerschaft nicht mehr als zwölf Wochen verstrichen sind (§ 219d - soziale oder Notlagen-Indikation).
\end{sourcepara}
\begin{transpara}
\transrn{38}d) 為避免懷孕婦女陷入嚴重困境之危險而有終止妊娠之指標事由，且該危險不能以其他對懷孕婦女可期待之方式排除，並且自妊娠開始起尚未逾十二週者（第219d條──社會或困境指標事由）。
\end{transpara}

\syncleft
\begin{sourcepara}
\sourcern{39}Beim Vorliegen eines dieser Indikationsfälle sollte der von einem Arzt mit Einwilligung der Schwangeren vorgenommene Schwangerschaftsabbruch nicht nach § 218 strafbar sein. Als Beginn der Schwangerschaft im Sinne des Gesetzes wurde der Abschluß der Einnistung des befruchteten Eies in der Gebärmutter bestimmt (§ 218 \abbrAbs{} 5).
\end{sourcepara}
\begin{transpara}
\transrn{39}具備上述任一指標事由時，由醫師經懷孕婦女同意所實施之終止妊娠，不應依第218條處罰。法律意義下之妊娠開始，定為受精卵在子宮內完成著床之時（第218條第5項）。
\end{transpara}

\syncleft
\begin{sourcepara}
\sourcern{40}Zur gleichen Zeit legten die Abgeordneten \abbrDr{} de With und Genossen den Entwurf eines Gesetzes zur Änderung des § 218 des Strafgesetzbuches vor (\abbrBTDrucks{} VI/3137), der einen Schwangerschaftsabbruch innerhalb der ersten drei Monate straflos lassen wollte, falls er mit Einwilligung der Schwangeren nach ärztlicher Beratung von einem Arzt vorgenommen werde (Fristenregelung).
\end{sourcepara}
\begin{transpara}
\transrn{40}同時，Dr. de With 等議員提出修正《刑法》第218條之法律草案（BTDrucks. VI/3137）；該草案主張，若終止妊娠係於最初三個月內，由醫師於懷孕婦女同意並經醫師諮詢後實施，應不具可罰性（期限規定）。
\end{transpara}

\syncleft
\begin{sourcepara}
\sourcern{41}Beide Entwürfe wurden im Sonderausschuß für die Strafrechtsreform gemeinsam behandelt. Am 10., 11. und 12. April 1972 fand eine öffentliche Anhörung von Sachverständigen und Auskunftspersonen aller einschlägigen Fachrichtungen statt, in der alle mit der Reform der Abtreibungsvorschriften zusammenhängenden Fragen umfassend erörtert wurden (vgl. Deutscher Bundestag, 6. \abbrWp{}, 74., 75., 76. Sitzung des Sonderausschusses für Strafrechtsreform vom 10., 11. und 12. April 1972, \abbrStenBer{} \abbrS{} 2141 bis \abbrS{} 2361). Die vorzeitige Auflösung des 6. Deutschen Bundestages führte zum Abbruch der Beratungen.
\end{sourcepara}
\begin{transpara}
\transrn{41}兩份草案均由刑法改革特別委員會共同審議。1972年4月10日、11日及12日，該委員會舉行公開聽證，聽取各相關專業領域之專家及知情人士意見，並全面討論與終止妊娠規定改革有關之所有問題（參見 Deutscher Bundestag, 6. Wp., 74., 75., 76. Sitzung des Sonderausschusses für Strafrechtsreform vom 10., 11. und 12. April 1972, StenBer. S. 2141 bis S. 2361）。第六屆德國聯邦議院提前解散，導致審議中止。
\end{transpara}

\syncleft
\begin{sourcepara}
\sourcern{42}In der 7. Wahlperiode brachte die Bundesregierung keinen eigenen Gesetzentwurf ein. Statt dessen wurden vier Gesetzentwürfe aus der Mitte des Bundestages vorgelegt. Der Entwurf der Abgeordneten \abbrDr{} Müller-Emmert und Genossen hatte eine Indikationenregelung zum Inhalt und sah zusätzlich vor, daß die Schwangere stets straffrei bleiben sollte (\abbrBTDrucks{} 7/443). Der Entwurf der Fraktionen der SPD, FDP (\abbrBTDrucks{} 7/375) schlug dagegen eine Fristenregelung vor. In einem Entwurf der CDU/CSU-Fraktion war eine engere Indikationenregelung vorgesehen (medizinische - die eugenische einschließend - und ehtische Indikation; Absehen von Strafe bei Unzumutbarkeit - \abbrBTDrucks{} 7/554 -). Ein von den Abgeordneten \abbrDr{} Heck und Genossen eingebrachter Entwurf wollte die Straflosigkeit des Schwangerschaftsabbruchs im wesentlichen auf den Fall einer - weitgefaßten - medizinischen Indikation begrenzen (\abbrBTDrucks{} 7/561).
\end{sourcepara}
\begin{transpara}
\transrn{42}在第七屆任期內，聯邦政府未提出自身法案；取而代之者，是由聯邦議院內部提出四項法律草案。Dr. Müller-Emmert 等議員之草案採指標事由規定，並另規定懷孕婦女本人始終不罰（BTDrucks. 7/443）。相對地，SPD、FDP 黨團之草案（BTDrucks. 7/375）提出期限規定。CDU/CSU 黨團草案則設想較窄之指標事由規定（醫學指標事由，包含優生指標事由，以及倫理指標事由；於不可期待時免除刑罰，BTDrucks. 7/554）。Dr. Heck 等議員提出之草案，基本上欲將終止妊娠不罰限於一項廣義醫學指標事由之情形（BTDrucks. 7/561）。
\end{transpara}

\syncleft
\begin{sourcepara}
\sourcern{43}Alle vier Entwürfe wurden an den Sonderausschuß für die Strafrechtsreform zur gemeinsamen Beratung überwiesen. Keiner dieser Entwürfe erhielt jedoch hier die erforderliche Mehrheit (vgl. \abbrBTDrucks{} 7/1982 \abbrS{} 4). Der Ausschuß sah sich deshalb nicht in der Lage, dem Plenum "bestimmte Beschlüsse zu empfehlen", sondern legte alle vier Entwürfe in getrennten Berichten zur Entscheidung vor - vgl. die \abbrBTDrucks{} 7/1982, 7/1981 (neu), 7/1983, 7/1984 (neu). Für die später vom Bundestag angenommene, von den Fraktionen der SPD, FDP vorgeschlagene Fristenregelung waren nach dem Bericht des Sonderausschusses - \abbrBTDrucks{} 7/1981 (neu) \abbrS{} 9/10 - insbesondere folgende Erwägungen maßgeblich
\end{sourcepara}
\begin{transpara}
\transrn{43}所有四份草案均提交刑法改革特別委員會共同審議。然而，這些草案均未在委員會取得所需多數（參見 BTDrucks. 7/1982 S. 4）。因此，委員會認為自己無法向全體會議「建議特定決議」，而是以個別報告將四份草案全部提交決定，參見 BTDrucks. 7/1982, 7/1981 (neu), 7/1983, 7/1984 (neu)。依特別委員會報告 BTDrucks. 7/1981 (neu) S. 9/10，對於後來由聯邦議院採納、由 SPD 與 FDP 黨團提出之期限規定，尤其具有決定意義者為下列考量：
\end{transpara}

\syncleft
\begin{sourcepara}
\sourcequoteinner{"Im Bereich des Strafrechts schlagen die Befürworter dieses Entwurfs für die ersten drei Schwangerschaftsmonate vor, die Strafandrohung im Interesse einer Verbesserung der Beratungssituation zurückzunehmen. Das heißt, strafrechtlich gesichert ist nur noch die Pflicht, sich einer umfassenden Beratung zu unterziehen und den Eingriff von einem Arzt vornehmen zu lassen. Das heißt weiter, daß in den ersten drei Monaten der Schutz des werdenden Lebens nicht mehr von der durchgängigen Strafdrohung, sondern von einem Beratungssystem gewährleistet wird, dessen Benutzung durch eine Strafdrohung verlangt wird. Die Befürworter der Fristenregelung gehen davon aus, daß die Strafdrohung erst nach dem dritten Monat wirklich 'greift'. Es hat sich gezeigt, daß ein durchgängiges strafrechtliches Verbot nicht geeignet ist, den Schutz des ungeborenen Lebens zu gewährleisten. Eine schwangere Frau, die ihre Schwangerschaft abbrechen lassen will, wird dies in aller Regel tun ohne Rücksicht auf das Strafgesetz, sie wird in jedem Fall einen Weg finden, einen Abbruch zu erreichen. Die Ursachen der Wirkungslosigkeit der Strafvorschrift wurden u.a. von Sachverständigen in der Öffentlichen Anhörung überzeugend damit erklärt, daß die Entscheidung zum Schwangerschaftsabbruch in aller Regel einer schwerwiegenden Konfliktsituation der Schwangeren entspringt und in den Tiefen der Persönlichkeit getroffen wird, die eine Strafdrohung nicht zu erreichen vermag (Rolinski AP VI \abbrS{} 2219, 2225; Schulte AP VI \abbrS{} 2200; Brocher AP VI \abbrS{} 2209) ... .\quoteparbreak
Die Fristenregelung gibt nicht den Gedanken von der Schutzwürdigkeit und Schutzbedürftigkeit des ungeborenen Lebens auf. Die Befürworter dieses Entwurfs sind nur der Auffassung, daß das Strafrecht in seiner geltenden Form nicht das geeignete Mittel ist ... .\quoteparbreak
Bei einem Schwangerschaftsabbruch müssen (und können) Beratung und Hilfe einsetzen, bevor die schwangere Frau den entscheidenden Schritt getan hat. Solange jedoch die Frau irgendwelche strafrechtlichen Sanktionen befürchten muß, wird sie kaum Beratung und Hilfe in Anspruch nehmen. Eine Frau, die, aus welchen Gründen auch immer, einen Schwangerschaftsabbruch will, wird vielmehr den Eingriff entweder selbst vornehmen oder einen Arzt oder eine andere Person suchen, von der sie weiß, daß diese den Eingriff vornimmt, ohne viel zu fragen ... . Solange eine Strafvorschrift besteht, sind diese Frauen für eine Beratung und Hilfe schwerlich erreichbar, weil sie sich meist nur an solche Personen um 'Hilfe' wenden, von denen sie sicher sind, daß sie sie dem gewünschten Schwangerschaftsabbruch näher bringen. Sie kommen meist gar nicht in den Bereich irgendeiner Person, die ihnen echte Hilfe bieten könnte und wollte."}
\end{sourcepara}
\begin{transpara}
\transquoteinner{「在刑法領域，本草案之支持者建議，為改善諮詢情境，於妊娠最初三個月撤回刑罰威嚇。這表示，刑法所擔保者，僅餘接受全面諮詢之義務，以及由醫師實施介入之義務。進一步而言，這也表示，在最初三個月，正在形成之生命之保護，不再由貫穿始終的刑罰威嚇加以確保，而是由一套諮詢制度加以確保；使用此一制度本身，則以刑罰威嚇予以要求。期限規定之支持者認為，刑罰威嚇唯有在第三個月以後才真正產生規範作用。事實已顯示，貫穿始終的刑事禁止並不適於保障未出生生命。欲終止妊娠之懷孕婦女，通常會不顧刑法而為之；無論如何，她都會找到達成終止妊娠之途徑。刑罰規定之所以欠缺效能，其原因已由專家在公開聽證中令人信服地說明：終止妊娠之決定通常源自懷孕婦女之嚴重衝突情境，並在人格深層作成，而刑罰威嚇無法觸及該處（Rolinski AP VI S. 2219, 2225; Schulte AP VI S. 2200; Brocher AP VI S. 2209）……。\quoteparbreak
期限規定並未放棄未出生生命值得保護且需要保護之思想。本草案之支持者僅認為，刑法以其現行形式並非適當手段……。\quoteparbreak
於終止妊娠之情形，諮詢與協助必須（且能夠）在懷孕婦女採取決定性步驟以前介入。然而，只要婦女仍須擔心任何刑事制裁，她幾乎就不會尋求諮詢與協助。毋寧，無論基於何種理由而欲終止妊娠之婦女，會自行實施介入，或尋找她知道不會多問、即會實施介入之醫師或其他人……。只要刑罰規定仍然存在，諮詢與協助即難以接觸這些婦女；因為她們通常只會向那些她們確信能使其更接近所欲終止妊娠之人尋求『幫助』。她們通常根本不會進入任何能夠且願意提供真正幫助之人的接觸範圍。」}
\end{transpara}

\syncleft
\begin{sourcepara}
\sourcern{44}Bei der Abstimmung in der zweiten Beratung des Deutschen Bundestages erhielt keiner der Entwürfe die als erforderlich vereinbarte Mehrheit der Stimmen. Daraufhin wurden die beiden Fraktionsentwürfe, welche die höchsten Stimmenzahlen erhalten hatten, zur Entscheidung gestellt. 245 Abgeordnete stimmten für den Entwurf der SPD, FDP-Fraktionen, 219 Abgeordnete für den Antrag der CDU/CSU-Fraktion (vgl. im einzelnen Deutscher Bundestag, 7. \abbrWp{}, 95. Sitzung, \abbrStenBer{} \abbrS{} 6443).
\end{sourcepara}
\begin{transpara}
\transrn{44}德國聯邦議院第二次審議投票時，所有草案均未獲得事先約定為必要之多數票。隨後，獲得最高票數的兩份黨團草案被提交表決。245名議員投票支持 SPD、FDP 黨團草案，219名議員投票支持 CDU/CSU 黨團提案（詳參德國聯邦議院，第7屆選期，第95次會議，StenBer. S. 6443）。
\end{transpara}

\syncleft
\begin{sourcepara}
\sourcern{45}Bei der namentlichen Schlußabstimmung über den Entwurf der Fraktionen der SPD, FDP in der dritten Beratung stimmten von 489 voll stimmberechtigten Abgeordneten 247 mit Ja, 233 mit Nein, 9 enthielten sich der Stimme (Deutscher Bundestag, 7. \abbrWp{}, 96. Sitzung, \abbrStenBer{} \abbrS{} 6503).
\end{sourcepara}
\begin{transpara}
\transrn{45}第三次審議中，對 SPD、FDP 黨團草案所作最後唱名表決，489名具完全投票權之議員中，247人投贊成票，233人投反對票，9人棄權（德國聯邦議院，第7屆選期，第96次會議，StenBer. S. 6503）。
\end{transpara}

\syncleft
\begin{sourcepara}
\sourcern{46}Der Bundesrat bezeichnete den Gesetzesbeschluß als zustimmungsbedürftig, versagte ihm nach erfolgloser Anrufung des Vermittlungsausschusses die Zustimmung und legte vorsorglich Einspruch gemäß \abbrArt{} 77 \abbrAbs{} 3 \abbrGG{}\ggfnArtSevenSevenThreeDE{} ein (Bundesrat, 406. Sitzung vom 31. Mai 1974, \abbrStenBer{} \abbrS{} 214). Diesen Einspruch wies der Bundestag, der das Gesetz nicht als zustimmungsbedürftig ansah, am 5.Juni 1974 mit 260 gegen 218 Stimmen bei 4 Enthaltungen zurück (Deutscher Bundestag, 7. \abbrWp{}, 104. Sitzung, \abbrStenBer{} \abbrS{} 6947).
\end{sourcepara}
\begin{transpara}
\transrn{46}聯邦參議院認為該法律決議須經其同意；在請求召集調解委員會未果後，拒絕給予同意，並依《基本法》第77條第3項\ggfnArtSevenSevenThreeZH{}預防性提出異議（聯邦參議院，1974年5月31日第406次會議，StenBer. S. 214）。聯邦議院則認為該法不須經同意，並於1974年6月5日以260票對218票、4票棄權駁回此項異議（Deutscher Bundestag, 7. Wp., 104. Sitzung, StenBer. S. 6947）。
\end{transpara}

\syncleft
\begin{sourcepara}
\sourcern{47}7. Zur Unterstützung der strafrechtlichen Reform durch sozialpolitische Maßnahmen beschloß der Deutsche Bundestag am 21. März 1974 auf einen Initiativantrag der Fraktionen der SPD, FDP (\abbrBTDrucks{} 7/376) das Gesetz über ergänzende Maßnahmen zum Fünften Strafrechtsreformgesetz (Strafrechtsreform-Ergänzungsgesetz - \abbrStREG{}). Darin sind u.a. Ansprüche auf ärztliche Beratung über Fragen der Empfängnisregelung sowie auf ärztliche Hilfe bei straffreiem Schwangerschaftsabbruch als Leistungen der gesetzlichen Krankenversicherung und der Sozialhilfe vorgesehen (vgl. \abbrBTDrucks{} 7/1753 und Deutscher Bundestag, 7. \abbrWp{}, 88. Sitzung, \abbrStenBer{} \abbrS{} 5769).
\end{sourcepara}
\begin{transpara}
\transrn{47}7. 為以社會政策措施支持刑法改革，德國聯邦議院於1974年3月21日依 SPD、FDP 黨團之倡議提案（BTDrucks. 7/376），通過《第五刑法改革法補充措施法》（刑法改革補充法，StREG）。其中包括將受胎調節問題之醫師諮詢請求權，以及不罰終止妊娠時之醫師協助請求權，規定為法定醫療保險及社會扶助給付（參見 BTDrucks. 7/1753 及 Deutscher Bundestag, 7. Wp., 88. Sitzung, StenBer. S. 5769）。
\end{transpara}

\syncleft
\begin{sourcepara}
\sourcern{48}Diesem Gesetzbeschluß verweigerte der Bundesrat nach erfolgloser Anrufung des Vermittlungsausschusses seine Zustimmung (Bundesrat, 410. Sitzung vom 12. Juli 1974, \abbrStenBer{} \abbrS{} 324). Daraufhin rief die Bundesregierung ihrerseits den Vermittlungsausschuß an. Dieser hat bisher noch keine Entschließung getroffen.
\end{sourcepara}
\begin{transpara}
\transrn{48}聯邦參議院在請求召集調解委員會未果後，拒絕同意此項法律決議（聯邦參議院，1974年7月12日第410次會議，StenBer. S. 324）。其後，聯邦政府亦請求召集調解委員會；該委員會迄今尚未作成決議。
\end{transpara}

\syncleft
\begin{sourcepara}
\sourcern{49}8. Am 21. Juni 1974 ordnete das Bundesverfassungsgericht auf Antrag der Regierung des Landes Baden-Württemberg im Wege der einstweiligen Anordnung u.a. an, daß § 218a Strafgesetzbuch in der Fassung des Fünften Gesetzes zur Reform des Strafrechts (5. \abbrStrRG{}) einstweilen nicht in Kraft tritt, jedoch der medizinisch, eugenisch oder der ethisch indizierte Schwangerschaftsabbruch innerhalb der ersten zwölf Wochen seit der Empfängnis straffrei bleibt (\abbrBVerfGE{} 37, 324; \abbrBGBl{} 1974 I \abbrS{} 1309). Die einstweilige Anordnung wurde bis zur Verkündung dieses Urteils verlängert.
\end{sourcepara}
\begin{transpara}
\transrn{49}8. 1974年6月21日，聯邦憲法法院應巴登-符騰堡邦政府聲請，以臨時命令等方式命令：《第五刑法改革法》（5. StrRG）版本《刑法》第218a條暫不生效；但受孕後最初十二週內，具醫學、優生或倫理指標事由之終止妊娠，仍不受刑罰處罰（BVerfGE 37, 324；BGBl. 1974 I S. 1309）。該臨時命令延長至本判決宣示為止。
\end{transpara}

\bisubsubsection{II.}{II.}
\syncleft
\begin{sourcepara}
\sourcern{50}193 Mitglieder des Deutschen Bundestages sowie die Landesregierungen von Baden-Württemberg, Bayern, Rheinland-Pfalz, des Saarlandes und von Schleswig-Holstein haben gemäß \abbrArt{} 93 \abbrAbs{} 1 \abbrNr{} 2 \abbrGG{}\ggfnArtNineThreeOneTwoDE{}, § 13 \abbrNr{} 6 \abbrBVerfGG{} den Antrag auf verfassungsrechtliche Überprüfung des § 218a \abbrStGB{} in der Fassung des Fünften Strafrechtsreformgesetzes gestellt. Sie halten die Vorschrift für unvereinbar mit dem Grundgesetz, weil die darin enthaltene Freigabe des Schwangerschaftsabbruchs während der ersten zwölf Wochen nach der Empfängnis (Fristenregelung) vor allem gegen \abbrArt{} 2 \abbrAbs{} 2 Satz 1 in Verbindung mit \abbrArt{} 1 \abbrAbs{} 1 \abbrGG{}\ggfnLifeDignityDE{} und außerdem gegen die \abbrArt{} 3 sowie 6 \abbrAbs{} 1, 2 und 4 \abbrGG{}\ggfnArtThreeSixDE{} sowie gegen das Rechtsstaatsprinzip verstoße. Die antragstellenden Landesregierungen sind ferner der Meinung, für das Fünfte Strafrechtsreformgesetz sei die Zustimmung des Bundesrates erforderlich gewesen.
\end{sourcepara}
\begin{transpara}
\transrn{50}德國聯邦議院193名議員，以及巴登-符騰堡、巴伐利亞、萊茵蘭-普法爾茨、薩爾及石勒蘇益格-荷爾斯泰因各邦政府，依《基本法》第93條第1項第2款\ggfnArtNineThreeOneTwoZH{}及《聯邦憲法法院法》第13條第6款，聲請就《第五刑法改革法》版本《刑法》第218a條進行憲法審查。其認為該規定與《基本法》不相容，因其中所含受孕後最初十二週內放開終止妊娠之規制（期限規定），尤其違反《基本法》第2條第2項第1句結合第1條第1項\ggfnLifeDignityZH{}，另亦違反第3條及第6條第1項、第2項、第4項\ggfnArtThreeSixZH{}以及法治國原則。聲請之各邦政府並認為，《第五刑法改革法》須經聯邦參議院同意。
\end{transpara}

\syncleft
\begin{sourcepara}
\sourcern{51}Zur Begründung sind im wesentlichen vorgetragen:
\end{sourcepara}
\begin{transpara}
\transrn{51}其理由主要陳述如下：
\end{transpara}

\syncleft
\begin{sourcepara}
\sourcern{52}1. Das Gesetz enthalte in \abbrArt{} 6 und 7 Änderungen der Strafprozeßordnung und des Einführungsgesetzes zum Strafgesetzbuch, also von Gesetzen, die ihrerseits mit Zustimmung des Bundesrates ergangen seien. Dies allein begründe bereits die Zustimmungsbedürftigkeit. Die gegenteilige Auffassung des Bundesverfassungsgerichts (\abbrBVerfGE{} 37, 363) möge nochmals überprüft werden.
\end{sourcepara}
\begin{transpara}
\transrn{52}1. 該法第6條及第7條修正《刑事訴訟法》及《刑法施行法》；此二法本身均係經聯邦參議院同意而制定。僅此即足以成立須經同意之理由。聯邦憲法法院相反見解（BVerfGE 37, 363）應再予審查。
\end{transpara}

\syncleft
\begin{sourcepara}
\sourcern{53}Abgesehen davon enthalte aber das Fünfte Strafrechtsreformgesetz selbst Vorschriften, welche die Zustimmungsbedürftigkeit gemäß \abbrArt{} 84 \abbrAbs{} 1 \abbrGG{}\ggfnArtEightFourOneDE{} auslösten; denn in § 218c \abbrAbs{} 1 \abbrNr{} 1 \abbrStGB{} sei die Einschaltung einer "ermächtigten Beratungsstelle" und in § 219 \abbrAbs{} 1 \abbrStGB{} die Bestätigung der Indikation durch eine "zuständige Stelle" vorgesehen. Ferner sei auch nach der erwähnten Entscheidung (\abbrAAO{}, \abbrS{} 383) ein Änderungsgesetz dann zustimmungsbedürftig, wenn es sich zwar auf die Regelung materiell-rechtlicher Fragen beschränke, in diesem Bereich jedoch Neuerungen in Kraft setze, die den nicht ausdrücklich geänderten Vorschriften über das Verwaltungsverfahren eine wesentlich andere Bedeutung und Tragweite verliehen. Das müsse zwangsläufig auch dann zur Zustimmungsbedürftigkeit eines Änderungsgesetzes führen, wenn das Ursprungsgesetz zwar noch keine Vorschriften über das Verwaltungsverfahren enthalte, das Änderungsgesetz aber durch die Ausgestaltung seiner materiell-rechtlichen Regelungen den notwendigen Vorschriften über das Verwaltungsverfahren derart vorgreife, daß die Länder in entscheidenden Punkten festgelegt würden. Dies sei hier der Fall; denn den Ländern werde bei der Ausgestaltung des Verwaltungsverfahrens nur noch ein geringer Spielraum belassen.
\end{sourcepara}
\begin{transpara}
\transrn{53}除此之外，《第五刑法改革法》本身亦包含依《基本法》第84條第1項\ggfnArtEightFourOneZH{}引發須經同意之規定；因《刑法》第218c條第1項第1款規定須納入「經授權之諮詢機構」，第219條第1項則規定由「主管機關」確認指標事由。此外，依前述裁判（a.a.O., S. 383），修正法即使限於規範實體法問題，但若其於該領域施行之新規制，使未明文修正之行政程序規定取得實質上不同之意義及射程，仍須經同意。即使原法尚無行政程序規定，但修正法透過其實體法規制之設計，已如此預先決定必要行政程序規定，以致各邦在關鍵點上被固定，亦必然導致修正法須經同意。本案即屬此情形；因為各邦於行政程序之形成上僅餘狹小空間。
\end{transpara}

\syncleft
\begin{sourcepara}
\sourcern{54}Schließlich müßten dieses Gesetz und das Strafrechtsreform-Ergänzungsgesetz rechtlich als eine Einheit angesehen werden, weil diese Gesetze der Sache nach unlöslich zusammengehörten. Ziel des Reformvorhabens sei eine "Fristenregelung mit Beratung"; die Beratung ihrerseits werde aber für alle im Mittelpunkt der Reform stehenden Fälle durch das Strafrechtsreform-Ergänzungsgesetz erst ermöglicht. Beide Gesetze stellten eine einheitliche politische Entscheidung dar. Es sei deshalb sachwidrig, wenn die Aufhebung der Strafandrohung und die Vorschriften über die Ermöglichung der Beratung in zwei verschiedenen Gesetzesvorlagen untergebracht worden seien. Offenbar sei dies deshalb geschehen, um die Zustimmung des Bundesrates zu umgehen; denn das Strafrechtsreform-Ergänzungsgesetz bedürfe zweifellos der Zustimmung des Bundesrates. Der Bundestag habe damit sein ihm prinzipiell zustehendes Ermessen, den zu regelnden Rechtsstoff auf mehrere Gesetze zu verteilen, überschritten.
\end{sourcepara}
\begin{transpara}
\transrn{54}最後，該法與《刑法改革補充法》在法律上必須被視為一個整體，因為二者就事物本質而言不可分割。改革計畫之目標是「附諮詢之期限規定」；但諮詢本身，對於改革核心所涉及之所有案件，乃是透過《刑法改革補充法》才成為可能。兩部法律構成一項統一的政治決定。因此，將刑罰威嚇之廢止與使諮詢成為可能之規定分置於兩項不同法律草案中，並不合乎事理。顯然，此舉是為規避聯邦參議院之同意；因為《刑法改革補充法》無疑須經聯邦參議院同意。聯邦議院因此已逾越其原則上固有的裁量，即將待規範之法律素材分配於數部法律之裁量。
\end{transpara}

\syncleft
\begin{sourcepara}
\sourcern{55}2. Das Grundrecht aus \abbrArt{} 2 \abbrAbs{} 2 Satz 1 \abbrGG{}\ggfnArtTwoTwoOneDE{} als das fundamentalste und ursprünglichste Menschenrecht schütze in umfassender Weise auch das ungeborene Leben. Diese Rechtsauffassung stimme mit der Entstehungsgeschichte und der herrschenden Meinung überein, stehe in der Tradition deutscher Rechtskultur und könne sich auf den Wortlaut der Verfassung stützen. Vor allem aber werde nur diese Rechtsansicht der erkennbaren Funktion der Verfassungsnorm gerecht.
\end{sourcepara}
\begin{transpara}
\transrn{55}2. 《基本法》第二條第二項第一句\ggfnArtTwoTwoOneZH{}的基本權利作為最基本、最原始的人權，也全面保障未出生的生命。這項法律意見符合淵源和主流意見，符合德國法律文化傳統，可以以《基本法》的文義為依據。然而，最重要的是，只有這種法律觀點才能公正地對待該基本法規範的公認功能。
\end{transpara}

\syncleft
\begin{sourcepara}
\sourcern{56}\abbrArt{} 2 \abbrAbs{} 2 Satz 1 \abbrGG{}\ggfnArtTwoTwoOneDE{} enthalte nicht nur ein Abwehrrecht gegen unmittelbare staatliche Eingriffe, sondern bilde zugleich die Grundlage für einen positiven Schutzanspruch gegenüber dem Staat. Die Schutzpflicht könne der Grundwertentscheidung der Verfassung für das werdende Leben entnommen werden, deren spezifische Funktion es sei, die in \abbrArt{} 2 \abbrAbs{} 2 Satz 1 \abbrGG{}\ggfnArtTwoTwoOneDE{} zum Ausdruck kommende Wertung auch für das Verhältnis zu Dritten fruchtbar werden zu lassen. Diese Verpflichtung ergebe sich aber auch unmittelbar aus \abbrArt{} 1 \abbrAbs{} 1 Satz 2 \abbrGG{}\ggfnArtOneOneTwoDE{}.
\end{sourcepara}
\begin{transpara}
\transrn{56}《基本法》第二條第二項第一句\ggfnArtTwoTwoOneZH{}不僅載有針對國家直接干預的抗辯權，也構成積極主張國家保護的基礎。保護義務可取自《基本法》對正在發展中生命所作的基本價值決定，其具體功能是使《基本法》第二條第二項第一句\ggfnArtTwoTwoOneZH{}所表達的評價對於與第三方的關係也富有成效。不過，這項義務也直接源自於《基本法》第一條第一項第二句\ggfnArtOneOneTwoZH{}。
\end{transpara}

\syncleft
\begin{sourcepara}
\sourcern{57}Der Schutz entspreche verfassungsrechtlichen Anforderungen nur dann, wenn er prinzipiell umfassend ausgestaltet sei. Dies beruhe nicht nur auf dem besonders hohen Rang der Rechtsgüter des \abbrArt{} 2 \abbrAbs{} 2 \abbrGG{}\ggfnArtTwoTwoDE{} und auf der Tatsache, daß jedes einzelne Leben den Schutz des Grundrechts genieße, vielmehr sei entscheidend, daß Verletzungen speziell des Grundrechts auf (biologisches) Leben zur totalen Vernichtung der Basis menschlicher Existenz führten.
\end{sourcepara}
\begin{transpara}
\transrn{57}只有從根本上全面的保護才符合基本法要求。這不僅基於《基本法》第二條第二項\ggfnArtTwoTwoZH{}的法益的特別崇高地位和每個個體的生命都享有基本權利的保護，而且至關重要的是，特別是對（生物）生命基本權利的侵犯會導致人類生存基礎的徹底破壞。
\end{transpara}

\syncleft
\begin{sourcepara}
\sourcern{58}Allerdings dürfe der Gesetzgeber bei der Ausgestaltung des Schutzes für das werdende Leben etwaige kollidierende Rechtssphären Dritter, insbesondere der Mutter, berücksichtigen, soweit diese Drittinteressen ihrerseits von der Verfassung getragen würden. Auf der Hand liege dies im Falle einer Gefährdung des Lebens der Schwangeren. Aber auch bei andersartigen schwerwiegenden Gefahren für die Frau könne der Gesetzgeber zu dem Ergebnis gelangen, daß ihr die unfreiwillige Austragung der Leibesfrucht nicht zugemutet werden könne. Das besage nicht, daß das werdende Leben als Rechtsgut geringer geachtet werde als etwa die Gesundheit oder sonstige Rechtsgüter der Frau. Wohl aber stoße die vollständige Verwirklichung des Schutzes des werdenden Lebens an die Grenzen des rechtlich Durchsetzbaren. Dem Gesetzgeber sei deshalb die Befugnis zuzusprechen, einzelne (Indikations-) Tatbestände näher zu konkretisieren. Damit sei ausreichend Raum gegeben, um einem etwaigen Wandel der gesellschaftlichen Wertvorstellungen in den Bahnen der Verfassung Rechnung zu tragen. Dieser Gesichtspunkt müsse aber dort versagen, wo keinerlei relevante Rechtspositionen der Mutter tangiert seien, wo der Schwangerschaftsabbruch aus Gleichgültigkeit oder reiner Bequemlichkeit erfolge.
\end{sourcepara}
\begin{transpara}
\transrn{58}不過，立法者在形成對正在發展中生命之保護時，得考量可能與之衝突之第三人法領域，尤其是母親之法領域；其前提是，此等第三人利益本身亦受憲法承認。懷孕婦女生命受危害之情形，顯然如此。即使婦女面臨其他種類之重大危險，立法者亦可能得出結論：不能期待其非自願地將胎兒懷至足月。這並不表示，形成中生命作為法益，較婦女健康或其他法益受到較低評價；毋寧是，對形成中生命保護之完全實現，會碰到法律可強制執行之界限。因此，應承認立法者有權進一步具體化個別（指標事由）構成要件。如此即有充分空間，在憲法軌道內顧及社會價值觀可能發生之變遷。然而，若母親並無任何相關法地位受到觸及，終止妊娠只是出於冷漠或純粹方便，此一觀點即不能成立。
\end{transpara}

\syncleft
\begin{sourcepara}
\sourcern{59}Die Pflicht des Staates, das ungeborene Leben zu schützen, gewinne dann einen besonderen Akzent, wenn es darum gehe, einen seit 100 Jahren bestehenden strafrechtlichen Schutz zu beseitigen. Dann bedürfe es einer besonders strengen Prüfung, ob die Schutzbeseitigung sich nicht mit der grundgesetzlichen Wertentscheidung für das Leben in Widerspruch setze.
\end{sourcepara}
\begin{transpara}
\transrn{59}當涉及消除已經存在了 100 年的刑事保護時，國家保護未出生生命的義務尤其受到重視。然後需要進行特別嚴格的審查，以確定取消保護是否與《基本法》對生命價值的決定不相矛盾。
\end{transpara}

\syncleft
\begin{sourcepara}
\sourcern{60}Mit der Fristenlösung verstoße der Staat in mehrfacher Hinsicht gegen die ihm obliegende Schutzverpflichtung:
\end{sourcepara}
\begin{transpara}
\transrn{60}透過期限規定，國家在幾個方面違反了其保護義務：
\end{transpara}

\syncleft
\begin{sourcepara}
\sourcern{61}a) Der Gesetzgeber verletze seine Pflicht schon dadurch, daß er die Vernichtung ungeborenen Lebens innerhalb der ersten zwölf Schwangerschaftswochen rechtlich erlaube, sofern sie nur von einem Arzt mit Einwilligung der Schwangeren vorgenommen werde. Anders als im Sinne rechtlicher Billigung könne die strafrechtliche Freigabe des Schwangerschaftsabbruchs nicht gedeutet werden.
\end{sourcepara}
\begin{transpara}
\transrn{61}a) 立法機關在法律上允許在懷孕前十二週內摧毀未出生的生命，這是違反其職責的，前提是必須由醫師在懷孕婦女同意的情況下進行。終止妊娠的刑事授權只能在法律同意的意義上解釋。
\end{transpara}

\syncleft
\begin{sourcepara}
\sourcern{62}b) Ferner entziehe der Gesetzgeber durch die Aufhebung der Strafbarkeit des Schwangerschaftsabbruchs in den ersten zwölf Wochen dem werdenden Leben für die Zukunft die sozialethische Wertschätzung in der Bevölkerung. Es entspreche gesicherter rechtssoziologischer Erkenntnis, daß strafrechtliche Normen normbildende Kraft für das sozialethische Urteil der Bürger besäßen.
\end{sourcepara}
\begin{transpara}
\transrn{62}b) 此外，立法機關在前十二週內廢除了對終止妊娠的刑事定罪，從而剝奪了人們未來的社會倫理價值。刑法規範對公民的社會倫理判斷具有規範形成力，這符合既定的法律社會學知識。
\end{transpara}

\syncleft
\begin{sourcepara}
\sourcern{63}c) Aber selbst wenn man unterstelle, daß der Staat den Schwangerschaftsabbruch weiterhin diskriminiere, verletze er seine verfassungsrechtliche Schutzpflicht schon dadurch, daß er die bisherige Strafandrohung für die ersten zwölf Schwangerschaftswochen ausnahmslos aufhebe. Sozialfürsorgerische Maßnahmen allein könnten nicht verhindern, daß abtreibungswillige Schwangere ihre Absicht realisierten. Denn wenn sie das Hilfsangebot nicht anzunehmen gewillt seien, so könne auf diesem Wege der Schutz des werdenden Lebens nicht gewährleistet werden. Der Wegfall der Strafsanktion hinterlasse daher eine relevante Schutzlücke. Hinzu komme, daß die als einzige Hilfsmaßnahme vorgeschriebene Beratung, so wie sie im Gesetz geregelt sei, sich in der Praxis als wirkungslos erweisen werde. Wenn die Schwangere wisse, daß die Entscheidung letztlich von ihr allein abhänge, werde sie wenig Neigung verspüren, sich diese Entscheidung durch Ermahnungen des Beraters erschweren zu lassen. Aber auch die Berater seien in einer widersprüchlichen Situation: Erkläre eine Schwangere, sie sei zur Abtreibung entschlossen, so stünden sie vor der Alternative, entweder nur formularmäßig Abtreibungserlaubnisse auszuschreiben oder der Schwangeren gegen ihren Willen Gewissensskrupel aufzuzwingen, die ihr das spätere Leben nur noch mehr belasten würden. Im übrigen seien auch die Ärzte, die die Hauptlast der medizinischen Beratung zu tragen hätten, damit überfordert.
\end{sourcepara}
\begin{transpara}
\transrn{63}c) 即使認為國家仍持續非難終止妊娠，完全撤回迄今對妊娠最初十二週之刑罰威嚇，仍違反憲法保護義務。僅以社會福利措施，不能阻止欲終止妊娠之懷孕婦女實現其意願；因為若其不願接受協助，便不能以此方式確保正在發展中生命之保護。因此，撤回刑事制裁將留下具有保護相關性之缺口。此外，法律僅規定諮詢作為唯一協助措施，在實務上將證明無效。若懷孕婦女知道最後決定取決於自己，即不會願意讓諮詢者之警告加重其決定困難。諮詢者亦處於矛盾境地：若懷孕婦女表示其已決定終止妊娠，諮詢者將面臨兩難，不是僅以形式方式發給終止妊娠之通行證，就是違反其意願迫使其產生良心疑慮，而此等疑慮只會成為其日後生活之負擔。此外，承擔醫療諮詢主要負荷之醫師亦將過度負擔。
\end{transpara}

\syncleft
\begin{sourcepara}
\sourcern{64}Zwar habe der Gesetzgeber bei der Abwägung, ob eine Strafsanktion eher schädlich als nützlich sei, einen gewissen Spielraum. Dieser werde aber durch das beherrschende rechtsstaatliche Prinzip der Erforderlichkeit (Übermaßverbot) eng begrenzt. Folge man der inzwischen herrschend gewordenen Lehre von der praktischen Konkordanz und des nach beiden Seiten hin schonendsten Ausgleichs, so müßten kollidierende Rechtsgüter einander so zugeordnet werden, daß jedes von ihnen Wirklichkeit gewinne. Es laufe auf die totale Verkehrung dieser Grundsätze hinaus, wenn einem der kollidierenden Güter generell der absolute Vorrang zugesprochen werde. Hierauf aber basiere die Fristenlösung, die der Selbstbestimmungsmacht der Frau, die in Wahrheit insoweit eine Fremdbestimmungsmacht sei, den ausschließlichen Vorrang einräume. Hinzu komme, daß hier einzelne Leben gegen andere einzelne Leben (oder Gesundheitsgüter) abgewogen werden sollten. Dabei gehe es um eine Art abstrakter staatlicher Abzählung und Aufrechnung von Leben gegen Leben. Dadurch werde aber der beherrschende und primäre individualrechtliche Gehalt des Grundrechts aus \abbrArt{} 2 \abbrAbs{} 2 Satz 1 \abbrGG{}\ggfnArtTwoTwoOneDE{} in Frage gestellt; es gebe kein anderes Grundrecht, das unmittelbarer auf den Schutz des Einzelnen als solchen abziele.
\end{sourcepara}
\begin{transpara}
\transrn{64}立法機關在權衡刑罰性制裁是否弊大於利時，確實有一定的空間。然而，這受到主導的法治國必要性原則（禁止過度）的嚴格限制。如果遵循目前流行的實踐一致性原則和雙方最溫和的平衡，則必須將相互衝突的法律利益分配給彼此，使每個利益都成為現實。如果一種衝突商品通常被給予絕對優先權，則相當於完全顛倒了這些原則。但期限規定正是基於此，將女性的自決權排他性地放在第一位，而這其實是一種外在的決定權。此外，個人生命應與其他個人生命（或健康產品）進行權衡。這是一種抽象狀態的生命與生命的數數與抵銷。但這對《基本法》第二條第二項第一句\ggfnArtTwoTwoOneZH{}所規定的基本權利的主導性、首要的個人法律內容提出了質疑；沒有其他基本權利更直接地旨在保護個人本身。
\end{transpara}

\syncleft
\begin{sourcepara}
\sourcern{65}Im übrigen könne man eine solche Abwägung nur mit der Annahme zu legitimieren versuchen, daß die Zahl der geretteten Leben erwartungsgemäß die der preisgegebenen deutlich übersteige. Indes stehe das Gegenteil fest. Während der parlamentarischen Beratungen habe der Regierungsvertreter aufgrund umfangreicher Untersuchungen und Vergleiche erklärt, selbst bei niedriger Schätzung sei mit einer Steigerung der Gesamtzahl legaler und illegaler Aborte um 40\% zu rechnen.
\end{sourcepara}
\begin{transpara}
\transrn{65}此外，人們只能試圖使這種權衡合法化，並假設，正如預期的那樣，拯救的生命數量大大超過了犧牲的生命數量。然而，相反的情况是肯定的。在議會審議期間，政府代表表示，根據廣泛的研究和比較，即使是較低的估計，合法和非法終止妊娠總數也預計將增加40\%。
\end{transpara}

\syncleft
\begin{sourcepara}
\sourcern{66}Schließlich sei auch zu beachten, daß die Fristenregelung gerade der bedrängten Frau jede Möglichkeit nehme, einen auf Abtreibung gerichteten Druck des Erzeugers oder sonstiger Beteiligter dadurch zu mildern, daß sie auf das Unrecht, die Strafbarkeit, das Risiko des ihr angesonnenen Verhaltens hinweise. Daran könne keine Beratung etwas ändern, am allerwenigsten eine solche, die wirkliche Gegendruck-Mittel nicht einsetzen dürfe. So sei es gerade die Fristenlösung, die die Schwangere in die Isolation treibe und sie unsachlichen Pressionen ausliefere. Sie sei die in Wahrheit unsoziale Regelung.
\end{sourcepara}
\begin{transpara}
\transrn{66}最後，還應該指出的是，期限規定剝奪了陷入困境的婦女透過指出她正在考慮的行為的不公正、刑事責任和風險來減輕來自生產者或其他相關方的終止妊娠壓力的機會。沒有任何諮詢可以改變這一點，尤其是不允許使用真正的反壓力措施的諮詢。正是期限規定使懷孕婦女陷入孤立，並使她面臨無關的壓力。這其實是一種反社會的規定。
\end{transpara}

\syncleft
\begin{sourcepara}
\sourcern{67}3. Die Freigabe des Schwangerschaftsabbruchs in den ersten zwölf Wochen verstoße auch gegen \abbrArt{} 3 \abbrGG{}\ggfnArtThreeDE{}; denn unter keinem sachlichen Gesichtspunkt lasse sich eine Fristenregelung, innerhalb derer kein Rechtsschutz für das werdende Leben bestehe, rechtfertigen. Ferner verletze die Fristenregelung \abbrArt{} 6 \abbrAbs{} 1 und \abbrAbs{} 4 \abbrGG{}\ggfnArtSixOneFourDE{} sowie das Rechtsstaatsprinzip.
\end{sourcepara}
\begin{transpara}
\transrn{67}3. 允許在首十二週內終止妊娠，亦違反《基本法》第三條\ggfnArtThreeZH{}；因為從客觀的角度來看，對發育中的生命沒有法律保護的期限是合理的。此外，期限規定亦違反《基本法》第六條第一項及第四項\ggfnArtSixOneFourZH{}以及法治原則。
\end{transpara}

\bisubsubsection{III.}{III.}
\syncleft
\begin{sourcepara}
\sourcern{68}Von den Verfassungsorganen, denen gemäß § 77 \abbrBVerfGG{} Gelegenheit zur Äußerung gegeben worden ist, haben der Bundestag und die Bundesregierung Stellung genommen. Die Regierung des Landes Nordrhein-Westfalen hat erklärt, daß sie sich der Stellungnahme der Bundesregierung anschließe.
\end{sourcepara}
\begin{transpara}
\transrn{68}在有機會根據 BVerfGG 第 77 條發表評論的憲法機構中，聯邦議院和聯邦政府已表明立場。北萊茵-威斯特法倫邦政府表示支持聯邦政府的聲明。
\end{transpara}

\syncleft
\begin{sourcepara}
\sourcern{69}1. Der Bundesminister der Justiz, der sich namens der Bundesregierung geäußert hat, hält die Anträge für unbegründet.
\end{sourcepara}
\begin{transpara}
\transrn{69}1. 代表聯邦政府發言的聯邦司法部長認為這些聲請沒有根據。
\end{transpara}

\syncleft
\begin{sourcepara}
\sourcern{70}a) Nach der Entscheidung des Bundesverfassungsgerichts vom 25. Juni 1974 (\abbrBVerfGE{} 37, 363) habe das Fünfte Strafrechtsreformgesetz nicht der Zustimmung des Bundesrates bedurft. Es enthalte weder neue Vorschriften, welche die Zustimmungsbedürftigkeit begründeten, noch ändere es Regelungen, die der Zustimmungsbedürftigkeit unterlegen hätten. Der Gesetzgeber sei nicht gehindert gewesen, das Fünfte Strafrechtsreformgesetz und das Strafrechtsreform-Ergänzungsgesetz getrennt zu beraten und zu verabschieden.
\end{sourcepara}
\begin{transpara}
\transrn{70}a) 根據聯邦憲法法院1974年6月25日的判決（BVerfGE 37, 363），第五刑法改革法案不需要聯邦參議院的同意。它不包含任何證明需要獲得同意的新法規，也不會改變原本需要獲得同意的法規。不妨礙立法機關分別討論通過《刑法改革第五法》和《刑法改革補充法》。
\end{transpara}

\syncleft
\begin{sourcepara}
\sourcern{71}b) Die im Fünften Strafrechtsreformgesetz vorgesehene Fristenregelung sei mit \abbrArt{} 2 \abbrAbs{} 2 Satz 1 \abbrGG{}\ggfnArtTwoTwoOneDE{} vereinbar.
\end{sourcepara}
\begin{transpara}
\transrn{71}b) 刑法改革第五法所規定的期限規定，與《基本法》第二條第二項第一句\ggfnArtTwoTwoOneZH{}相符。
\end{transpara}

\syncleft
\begin{sourcepara}
\sourcern{72}Die Bundesregierung sei stets davon ausgegangen, daß das ungeborene Leben in den Schutz des \abbrArt{} 2 \abbrAbs{} 2 Satz 1 \abbrGG{}\ggfnArtTwoTwoOneDE{} einbezogen sei. Auch in den ausführlichen Beratungen im Bundestag und Bundesrat seien der Verfassungsrang und die Verpflichtung zum Schutz des ungeborenen Lebens außer Streit gewesen.
\end{sourcepara}
\begin{transpara}
\transrn{72}聯邦政府一直認為未出生的生命受到《基本法》第二條第二項第一句\ggfnArtTwoTwoOneZH{}的保護。在聯邦議院和聯邦參議院的詳細討論中，憲法地位和保護未出生生命的義務也沒有爭議。
\end{transpara}

\syncleft
\begin{sourcepara}
\sourcern{73}Ungeachtet der Verpflichtung zu einem angemessenen und wirksamen Schutz des ungeborenen Lebens verlange \abbrArt{} 2 \abbrAbs{} 2 Satz 1 \abbrGG{}\ggfnArtTwoTwoOneDE{} jedoch keinen Schutz des werdenden Lebens, der dem des geborenen gleichartig sei. Der Gesetzgeber müsse sich deshalb nicht für einen durchgängigen strafrechtlichen Schutz des ungeborenen Lebens entscheiden, sofern der Schutz auf andere Weise den Anforderungen der Schutzwürdigkeit und Schutzbedürftigkeit des ungeborenen Lebens entspreche.
\end{sourcepara}
\begin{transpara}
\transrn{73}儘管有充分、有效保護未出生生命的義務，《基本法》第二條第二項第一句\ggfnArtTwoTwoOneZH{}並沒有要求保護與出生生命類似的正在發展中的生命。因此，立法機關不必決定對未出生生命進行一致的刑法保護，只要該保護滿足未出生生命的價值要求和以另一種方式保護的需要即可。
\end{transpara}

\syncleft
\begin{sourcepara}
\sourcern{74}Diese Auffassung werde zunächst durch die Entstehungsgeschichte des \abbrArt{} 2 \abbrAbs{} 2 Satz 1 \abbrGG{}\ggfnArtTwoTwoOneDE{} bestätigt. Aus ihr ergebe sich eindeutig, daß nicht von einem durchgehenden Pönalisierungszwang zum Schutz des ungeborenen Lebens ausgegangen werden müsse.
\end{sourcepara}
\begin{transpara}
\transrn{74}這一觀點最初為《基本法》第二條第二項第一句\ggfnArtTwoTwoOneZH{}的歷史所證實。由此可見，沒有必要假設會存在持續的強制刑罰來保護未出生的生命。
\end{transpara}

\syncleft
\begin{sourcepara}
\sourcern{75}Gleichartigkeit des Schutzes sei allerdings erforderlich, soweit staatliche Angriffe abzuwenden seien. Es gehe jedoch weder darum noch um die Abwehr von Angriffen Dritter; denn das Verhältnis der Leibesfrucht zur Mutter sei von besonderer \abbrArt{} Die Leibesfrucht sei mit Leib und Leben der Mutter in denkbar engster Weise verbunden, auch wenn sie als selbständiges Rechtsgut anerkannt werde. Den Schutz habe bereits die Natur in die unmittelbare Obhut der Mutter gelegt. Die Möglichkeiten der Rechtsordnung, das ungeborene Leben auch gegenüber der Mutter zu schützen, seien von der Sache her begrenzt. Strafrechtliche Sanktionen hätten nach den bisherigen Erfahrungen nur in beschränktem Maße Schwangere zum Austragen einer Leibesfrucht bewegen können, wenn bei ihnen dazu die Bereitschaft nicht vorhanden sei. Die Schutzpflicht könne deshalb für den Staat nicht eine durchgehende Strafpflicht begründen. Mit der Zurücknahme der Strafdrohung gegenüber einem Abbruch der Schwangerschaft in den ersten Schwangerschaftswochen verleihe der Staat der Mutter kein Recht zum Eingriff. Der Gesetzgeber begrenze lediglich die Strafdrohung im Hinblick darauf, daß andere Schutzvorkehrungen als geeigneter und wirksamer angesehen würden oder um schutzwürdigen Belangen der Schwangeren Rechnung zu tragen.
\end{sourcepara}
\begin{transpara}
\transrn{75}然而，如果要避免國家攻擊，統一的保護是必要的。然而，這與此無關，也與防禦第三方攻擊無關；因為胎兒和母親之間的關係是一種特殊的關係。胎兒與母親的身體和生命以盡可能密切的方式聯繫在一起，即使它被承認為獨立的法律資產。大自然已經將保護置於母親的直接照顧中。法律體系保護未出生生命（包括針對母親的生命）的可能性本質上是有限的。根據以往的經驗，刑事制裁只能說服懷孕婦女在不願意的情況下生下胎兒。因此，保護義務不能成為國家持續刑罰義務的理由。透過撤銷對懷孕頭幾週終止妊娠的刑罰威嚇，國家並沒有賦予母親干預的權利。立法機關限制刑罰威嚇只是為了確保其他保護措施被認為更合適、更有效，或是為了考慮到值得保護的懷孕婦女利益。
\end{transpara}

\syncleft
\begin{sourcepara}
\sourcern{76}Verfassungsrechtlich könne von einer generellen Strafpflicht des Staates auch deshalb nicht ausgegangen werden, weil der Gesetzgeber gehalten sei, entgegenstehende Grundrechte der Schwangeren selbst und verfassungsrechtliche Wertentscheidungen zu berücksichtigen. \abbrArt{} 2 \abbrAbs{} 2 Satz 1 \abbrGG{}\ggfnArtTwoTwoOneDE{} wirke zugunsten schwangerer Frauen als ein auch ihnen zustehendes Grundrecht auf Leben und körperliche Unversehrtheit. Dem Schutz des Lebens und der Gesundheit Schwangerer komme im Hinblick auf die große Zahl der nicht von Ärzten vorgenommenen Schwangerschaftsabbrüche ein beachtliches Gewicht zu. Das Kurpfuschertum sei zwar im Vergleich zu früheren Zeiten zurückgegangen, begründe aber immer noch für das Leben und die Gesundheit schwangerer Frauen eine erhebliche Gefahr. Auch bei rechtswidrigen Schwangerschaftsabbrüchen durch Ärzte sei in kritischen Fällen die erforderliche Not- sowie die Nachversorgung nicht unbedingt gewährleistet. Außerdem sei die Selbstverantwortung schwangerer Frauen zu beachten. Die Wertentscheidung zugunsten der Selbstverantwortung des Menschen habe hier insofern eine entscheidende verfassungsrechtliche Bedeutung, als der Gesetzgeber sie gerade bei der Regelung eines Lebensbereichs, der in starkem Maße von der natürlichen Verantwortung der Frau für ihre Leibesfrucht geprägt sei, mit zugrunde legen müsse. Der Gesetzgeber habe angesichts der sich gegenseitig beeinflussenden und begrenzenden Wertentscheidungen und Grundrechte die Aufgabe, alle Rechtspositionen zu berücksichtigen und zum Ausgleich zu bringen.
\end{sourcepara}
\begin{transpara}
\transrn{76}從憲法上看，也不能因此出發認為國家負有一般性的刑罰義務，因為立法者有義務考量孕婦自身相對立的基本權利與憲法上的價值決定。《基本法》第 2 條第 2 項第 1 句\ggfnArtTwoTwoOneZH{}，對孕婦而言，亦作為其所享有的生命與身體完整性基本權而發生作用。鑑於大量終止妊娠並非由醫師實施，對孕婦生命與健康的保護具有相當重要性。密醫行為雖較過去減少，但對孕婦生命與健康仍構成重大危險。即使是由醫師實施的違法終止妊娠，在危急案件中，必要的緊急照護與後續照護也未必能獲得保障。此外，也必須注意孕婦的自主負責。支持人的自主負責之價值決定，在此具有決定性的憲法意義；因為立法者在規範一個強烈受到婦女對其胎兒之自然責任所形塑的生活領域時，尤其必須將此一價值決定一併作為基礎。面對彼此相互影響並相互限制的價值決定與基本權利，立法者的任務乃是考量所有法地位，並使其達成平衡。
\end{transpara}

\syncleft
\begin{sourcepara}
\sourcern{77}Beim Erlaß von Strafrechtsnormen sei überdies zu beachten, daß dem Strafrecht aus dem Gebot des sinn- und maßvollen Strafens Grenzen gesetzt seien. Nicht überall dort, wo verfassungsrechtliche Wertentscheidungen vorlägen, sei ihre Sicherung durch Strafsanktionen geboten. Eine durchgehende Parallelität von verfassungsrechtlicher und strafrechtlicher Wertordnung lasse sich nicht herstellen; beide Bereiche seien nicht identisch.
\end{sourcepara}
\begin{transpara}
\transrn{77}在製定刑法規範時，也應考慮刑法受到合理、適度處罰要求的限制。並非所有做出憲法價值決定的地方都需要透過刑罰性制裁來確保這些決定。憲法價值秩序與刑事價值秩序之間無法建立一致的平行關係；兩個區域並不相同。
\end{transpara}

\syncleft
\begin{sourcepara}
\sourcern{78}\abbrArt{} 2 \abbrAbs{} 2 Satz 1 \abbrGG{}\ggfnArtTwoTwoOneDE{} könne hiernach grundsätzlich kein weitergehender Auftrag für die Ausgestaltung der allgemeinen Rechtsordnung entnommen werden, als der, einen dem ungeborenen Leben angemessenen und wirksamen Schutz zu gewährleisten. Der Gesetzgeber habe die von einer verfassungsrechtlichen Grundentscheidung ausgehenden Richtlinien und Impulse zu beachten. Im übrigen sei er jedoch in der Ausgestaltung dieser Maßnahmen frei.
\end{sourcepara}
\begin{transpara}
\transrn{78}原則上，不能由此推斷《基本法》第二條第二項第一句\ggfnArtTwoTwoOneZH{}為一般法律制度的設計提供比確保未出生生命得到適當和有效保護更深遠的授權。立法機關必須遵守基本憲法決定所產生的指導方針和動力。然而，除此之外，他可以自由地設計這些措施。
\end{transpara}

\syncleft
\begin{sourcepara}
\sourcern{79}c) Die Fristenregelung des Fünften Strafrechtsreformgesetzes genüge in Inhalt und Ausgestaltung den verfassungsrechtlichen Anforderungen des \abbrArt{} 2 \abbrAbs{} 2 Satz 1 \abbrGG{}\ggfnArtTwoTwoOneDE{}. Die Beratung schwangerer Frauen während der ersten zwölf Schwangerschaftswochen, das Kernstück der Fristenregelung, gewährleiste auch ohne Strafdrohung den erforderlichen Schutz des ungeborenen Lebens.
\end{sourcepara}
\begin{transpara}
\transrn{79}（三）刑法改革第五法的期限規定，無論在內容上或設計上都符合《基本法》第二條第二項第一句\ggfnArtTwoTwoOneZH{}的憲法要求。期限規定的核心是在懷孕前十二週為懷孕婦女提供諮詢，確保未出生的生命得到必要的保護，即使沒有刑罰威嚇。
\end{transpara}

\syncleft
\begin{sourcepara}
\sourcern{80}Der Gesetzgeber sei dabei davon ausgegangen, daß die Entscheidung zum Schwangerschaftsabbruch in aller Regel einer schwerwiegenden Konfliktsituation der Schwangeren entspringe und in den Tiefen der Persönlichkeit getroffen werde. Eine Strafvorschrift vermöge deshalb zum Schwangerschaftsabbruch geneigte oder bereite Frauen regelmäßig nicht zu erreichen. Andererseits liege, wenn eine Frau trotz des Risikos für ihre eigene Gesundheit oder gar das eigene Leben einen Schwangerschaftsabbruch erwäge, eine Situation vor, die Beratung und Hilfestellung erfordere. Beides erscheine auch nicht als aussichtslos, da viele dieser Frauen noch in ihrem Entschluß schwankten und keineswegs von vornherein auf den Schwangerschaftsabbruch fixiert seien. Solche Frauen könnten durch eine eingehende Beratung zu einer positiven Einstellung zur Schwangerschaft gebracht werden. Eine durchgehende Strafandrohung auch für die erste Zeit der Schwangerschaft, wie sie bei Indikationsregelungen vorgesehen sei, schwäche die Effektivität von Beratungs- und Hilfsangeboten entscheidend. Da eine Indikationsregelung zwangsläufig irgendeine Form eines Begutachtungssystems voraussetze, unterwerfe sie die Schwangere dem Zwang, sich dieser Begutachtung zu unterziehen. Dies sei aber einer offenen Beteiligung der Schwangeren an der Beratung abträglich und führe bei abtreibungsbereiten und -geneigten Frauen häufig dazu, daß der Weg zur Beratung und zur Gutachterstelle von vornherein vermieden werde.
\end{sourcepara}
\begin{transpara}
\transrn{80}立法者於此係基於下列認識：終止妊娠之決定通常源自懷孕婦女之重大衝突處境，且是在其人格深處作成。因此，刑罰規範通常無法觸及傾向或準備終止妊娠之婦女。另一方面，若婦女即使冒著自身健康、甚至生命之風險，仍考慮終止妊娠，即存在需要諮詢與協助之處境。此二者並非毫無希望，因為許多此類婦女仍在決意中動搖，並非一開始即固定於終止妊娠。藉由深入諮詢，此類婦女可能被引導至對妊娠之正面態度。即使在妊娠初期仍一貫維持刑罰威嚇，如指標規制所設想者，將決定性地削弱諮詢與協助服務之效果。由於指標規制必然預設某種鑑定制度，懷孕婦女即被迫接受此種鑑定。這不利於懷孕婦女開放地參與諮詢，且常使準備或傾向墮胎之婦女從一開始即避開諮詢及鑑定機構。
\end{transpara}

\syncleft
\begin{sourcepara}
\sourcern{81}Im übrigen ließen die im Zusammenhang mit der Reform des § 218 \abbrStGB{} angestellten Prognosen keine verfassungsrechtlichen Schlußfolgerungen zu. Sie müßten auf unbekannte Größen zurückgreifen und erlaubten deshalb keine zuverlässigen Rückschlüsse. Die Prognosefehler kumulierten sich überdies, wenn die dem Fünften Strafrechtsreformgesetz zugrunde liegende Regelung als Entkriminalisierung erklärt und mit der Rechtslage in Ländern verglichen werde, in denen diese tatsächlich so verstanden werden müsse. Der Zweck dieses Gesetzes sei nicht die Entkriminalisierung des Schwangerschaftsabbruchs, sondern die Gewährleistung eines angemessenen und wirksamen Schutzes des ungeborenen Lebens, auch wenn anstelle einer durchgehenden Bestrafung des Schwangerschaftsabbruchs eine verbundene Regelung vorgesehen sei, die für die ersten zwölf Schwangerschaftswochen der Beratung als vorbeugender Maßnahme eine größere Eignung als der Strafandrohung beimesse.
\end{sourcepara}
\begin{transpara}
\transrn{81}此外，與《StGB》第218條改革相關的預測不允許得出任何憲法結論。他們必須依賴未知的數量，因此不允許任何可靠的結論。當第五刑法改革法案所依據的法規被宣布非刑事化並與實際上必須以這種方式理解的國家的法律情況進行比較時，預測誤差也會累積。該法的目的不是使終止妊娠非刑罪化，而是確保對未出生生命的充分和有效保護，即使不是對終止妊娠直接科以刑罰，而是製定了相關規定，在懷孕的前十二週內，提出諮詢作為預防措施，比刑罰威嚇更合適。
\end{transpara}

\syncleft
\begin{sourcepara}
\sourcern{82}Der Gesetzgeber habe die Beratung wirksam gestaltet. Neben einer Beratung, bei der ärztliche Gesichtspunkte im Vordergrund stünden, stehe die Unterrichtung über die zur Verfügung stehenden öffentlichen und privaten Hilfen für Schwangere, Mütter und Kinder sowie insbesondere über solche Hilfen, die die Fortsetzung der Schwangerschaft und die Lage von Mutter und Kind erleichtern. Unzweideutig sei, daß die Beratung auf Fortsetzung der Schwangerschaft ausgerichtet sein müsse. Der Gesetzgeber habe die Beratung durch eine Strafandrohung abgesichert. Durch das Strafrechtsreform-Ergänzungsgesetz werde sichergestellt, daß der Weg zum Arzt oder zur Beratungsstelle nicht aus Kostengründen erschwert werde. Die Befürchtung, die Fristenregelung führe zu einem "Dammbruch" im Rechtsbewußtsein unseres Volkes, sei bereits in der wissenschaftlichen Literatur überzeugend widerlegt.
\end{sourcepara}
\begin{transpara}
\transrn{82}立法機關使磋商取得了成效。除了專注於醫療方面的諮詢外，還提供有關懷孕婦女、母親和兒童可獲得的公共和私人援助的信息，特別是有關使繼續妊娠和母親和兒童的狀況更容易的援助的信息。毫無疑問，該諮詢必須旨在繼續妊娠。立法機關以刑罰相威嚇來確保磋商。刑法改革補充法確保就醫或諮詢中心不會因費用原因而變得更加困難。對於期限規定會導致我國人民法律意識「決堤」的擔憂，已經在科學文獻中得到了令人信服的反駁。
\end{transpara}

\syncleft
\begin{sourcepara}
\sourcern{83}d) Die bedingte Rücknahme der Strafdrohung für den Schwangerschaftsabbruch in den ersten zwölf Schwangerschaftswochen verstoße auch nicht gegen \abbrArt{} 3 \abbrGG{}\ggfnArtThreeDE{}. Die unterschiedlichen Schutzvorkehrungen beruhten nicht auf einer verschiedenen Bewertung der Leibesfrucht nach ihrem Alter, sondern seien Teil einer einheitlichen Regelung, die nicht nur insgesamt, sondern auch in ihren Teilen auf einen gleichwertigen Schutz der Leibesfrucht ausgerichtet sei und diesen Schutz in einer dem jeweiligen Entwicklungsstadium angemessenen und wirksamen Weise ausgestalte. Auch \abbrArt{} 6 \abbrAbs{} 1 und 4 \abbrGG{} sowie das Rechtsstaatsprinzip seien nicht verletzt.
\end{sourcepara}
\begin{transpara}
\transrn{83}d) 有條件撤銷對懷孕前十二週內墮胎的刑罰威嚇，並不違反《基本法》第三條\ggfnArtThreeZH{}。不同的保護措施並非基於根據年齡對胎兒進行不同的評估，而是統一規定的一部分，該規定不僅總體而且部分地旨在對胎兒進行同等保護，並針對各個發育階段以適當和有效的方式提供這種保護。亦不違反《基本法》第六條第一項和第四項\ggfnArtSixOneFourZH{} 以及法治原則。
\end{transpara}

\syncleft
\begin{sourcepara}
\sourcern{84}2. Für den Deutschen Bundestag hat der Bundestagsabgeordnete Rechtsanwalt Professor \abbrDr{} Ehmke eine Stellungnahme abgegeben. Er ist ebenfalls der Auffassung, die Fristenregelung des Fünften Strafrechtsreformgesetzes sei mit dem Grundgesetz vereinbar. Er verneint die Zustimmungsbedürftigkeit dieses Gesetzes aus den gleichen Gründen wie die Bundesregierung.
\end{sourcepara}
\begin{transpara}
\transrn{84}2.德國聯邦議院方面，聯邦議院議員、律師Prof. Dr. Ehmke 發表聲明。他亦認為刑法改革第五法的期限規定與《基本法》一致。他基於與聯邦政府相同的理由，否認本法須經聯邦參議院同意。
\end{transpara}

\syncleft
\begin{sourcepara}
\sourcern{85}Weiter führt er aus:
\end{sourcepara}
\begin{transpara}
\transrn{85}他繼續說：
\end{transpara}

\syncleft
\begin{sourcepara}
\sourcern{86}Das Gesetz enthalte keine "perfekte", aber eine sinnvolle Lösung des seit Jahrzehnten diskutierten Reformproblems. Die hohen Dunkelziffern, die geringe Zahl der - nur auf leichte Strafen lautenden - Verurteilungen, die sich negativ auf die soziale Geltungskraft des § 218 \abbrStGB{} ausgewirkt habe, die Benachteiligung der Frauen aus sozial schwächeren Schichten, die sich Kurpfuschern hätten anvertrauen müssen, sowie die Folgekriminalität hätten dringend eine Reform des § 218 \abbrStGB{} verlangt. Gesundheitspolitisch wie gesellschaftspolitisch sei die Fristenlösung vorzuziehen; sie werde insbesondere mit hoher Wahrscheinlichkeit die Zahl der illegalen Schwangerschaftsabbrüche und die mit ihnen verbundenen gesundheitlichen Gefahren für die Mütter vermindern und im übrigen auch die Folgekriminalität eindämmen können. Im Gegensatz zu einer Indikationslösung schließe die Fristenlösung sowohl regionale wie individuelle Ungleichheit und damit Ungerechtigkeit bei der Anwendung der Indikationstatbestände als auch die soziale Ungleichheit hinsichtlich der Ausweichmöglichkeiten ins Ausland aus. Sie beseitige die Angst der Frauen vor dem Abgewiesenwerden, die sie wieder in die Illegalität treibe, und eröffne damit einer ohne Zeitdruck erfolgenden Beratung und selbstverantwortlichen Überlegung der Mutter überhaupt erst eine Chance. Es verspreche mehr Erfolg, die Einsicht und die Selbstverantwortung der Frauen zu stärken, als auf ihre Angst vor der Strafdrohung zu bauen.
\end{sourcepara}
\begin{transpara}
\transrn{86}該法律並非提供一項「完美」解決方案，但對已爭論數十年之改革問題而言，仍是一項合理方案。高犯罪黑數、僅宣告輕刑之有罪判決數量偏低，已對《刑法》第218條之社會效力造成負面影響；社會弱勢階層婦女被迫求助於密醫所受之不利益，以及隨之而來之後續犯罪，均迫切要求改革《刑法》第218條。從衛生政策及社會政策觀點而言，期限方案較可取；尤其可以高度可能地減少非法終止妊娠之數量，以及與之相連之母親健康危險，並且也能抑制後續犯罪。相較於指標方案，期限規定排除了適用指標事由構成要件時之地區性與個別性不平等及其不公，也排除了關於赴外國迴避可能性之社會不平等。期限規定消除婦女對遭拒絕、因而再度被推入非法領域之恐懼，並由此才使無時間壓力之諮詢及母親自主負責之思考獲得機會。強化婦女之理解與自主責任，較諸建立在其對刑罰威嚇之恐懼上，更有成功可能。
\end{transpara}

\syncleft
\begin{sourcepara}
\sourcern{87}Die Frage der Verfassungsmäßigkeit der Fristenlösung lasse sich nicht allein aus dem Wortlaut des \abbrArt{} 2 \abbrAbs{} 2 Satz 1 \abbrGG{}\ggfnArtTwoTwoOneDE{} und der Entstehungsgeschichte dieser Vorschrift beantworten. Aus den Verhandlungen des Parlamentarischen Rates ergebe sich nur, daß er eine Reform der Strafvorschriften über den Schwangerschaftsabbruch bewußt nicht habe präjudizieren wollen und deshalb die Frage, ob und wieweit \abbrArt{} 2 \abbrAbs{} 2 Satz 1 \abbrGG{}\ggfnArtTwoTwoOneDE{} auch ungeborenes Leben schütze, nicht ausdrücklich entschieden habe. Daher sei auf anderem Wege über die Auslegung des \abbrArt{} 2 \abbrAbs{} 2 Satz 1 \abbrGG{}\ggfnArtTwoTwoOneDE{} Klarheit zu gewinnen.
\end{sourcepara}
\begin{transpara}
\transrn{87}時限方案是否合憲的問題，不能只從《基本法》第二條第二項第一句\ggfnArtTwoTwoOneZH{}的措詞和這條規定的歷史來回答。議會理事會談判中唯一的表現是，它有意識地不想預先判斷終止妊娠刑法改革，因此沒有明確決定《基本法》第二條第二項第一句\ggfnArtTwoTwoOneZH{}是否以及在多大程度上也保護未出生的生命。因此，《基本法》第二條第二項第一句\ggfnArtTwoTwoOneZH{}的解釋還有另一種方式。
\end{transpara}

\syncleft
\begin{sourcepara}
\sourcern{88}Die Verwendung des Wortes "jeder" spreche gegen die Annahme eines Grundrechts des ungeborenen Lebens, da sowohl in der Umgangs- wie auch in der Rechtssprache mit "jeder" eindeutig eine menschliche Person bezeichnet werde. Auch im Rechtssinne beginne das Person- und Menschsein erst mit der Geburt.
\end{sourcepara}
\begin{transpara}
\transrn{88}「每個人」一詞的使用違背了對未出生生命的基本權利的假設，因為在日常和法律語言中，「每個人」顯然指的是一個人。即使在法律意義上，成為一個人、成為人類也只是從出生開始。
\end{transpara}

\syncleft
\begin{sourcepara}
\sourcern{89}Damit sei allerdings noch nicht entschieden, ob und inwieweit ungeborenes Leben ein durch \abbrArt{} 2 \abbrAbs{} 2 Satz 1 \abbrGG{}\ggfnArtTwoTwoOneDE{} geschütztes Rechtsgut darstelle. In Übereinstimmung mit der in der Lehre überwiegend und im Deutschen Bundestag einhellig vertretenen Meinung sei \abbrArt{} 2 \abbrAbs{} 2 Satz 1 \abbrGG{}\ggfnArtTwoTwoOneDE{} nach seinem Sinn und Zweck dahin auszulegen, daß er als Grundsatznorm der Verfassung auch das ungeborene Leben als Vorstufe des Menschenlebens schütze. Damit sei aber noch völlig offen, wann der Rechtsschutz beginne, welcher Art er zu sein habe und wie er gegenüber verschiedenen potentiellen Verletzern zu differenzieren sei.
\end{sourcepara}
\begin{transpara}
\transrn{89}然而，胎兒生命是否及在何種範圍內構成受《基本法》第2條第2項第1句\ggfnArtTwoTwoOneZH{}保護之法益，尚未因此決定。依德國聯邦議院之主流且一致見解，《基本法》第2條第2項第1句\ggfnArtTwoTwoOneZH{}應按其意義與目的解釋為一項憲法基本規範，亦保護作為人類生命前階段之未出生生命。然而，法律保護自何時開始、應採何種類型，以及如何依不同可能侵害者加以區分，仍完全開放。
\end{transpara}

\syncleft
\begin{sourcepara}
\sourcern{90}Diese Frage könne weder durch Rückgriff auf kirchliche Lehren oder religiöse Überzeugungen noch mit den Methoden der Naturwissenschaft beantwortet werden. Auch eine Berufung auf \abbrArt{} 1 \abbrAbs{} 1 und \abbrArt{} 19 \abbrAbs{} 2 \abbrGG{}\ggfnArtOneNineteenDE{} könne keine zusätzlichen Erkenntnisse bringen. Bei der Interpretation des Begriffs "Mensch" stellten sich die gleichen noch nicht beantworteten Auslegungsfragen. \abbrArt{} 19 \abbrAbs{} 2 \abbrGG{} bestimme nicht den Umfang eines Grundrechts, sondern seinen unverzichtbaren "Kern"; bevor diese Bestimmung herangezogen werden könne, müsse also zunächst geprüft werden, ob und inwieweit ungeborenes Leben durch \abbrArt{} 2 \abbrAbs{} 2 Satz 1 \abbrGG{}\ggfnArtTwoTwoOneDE{} überhaupt geschützt werde.
\end{sourcepara}
\begin{transpara}
\transrn{90}此問題既不能由教會教義或宗教信念回答，亦不能由自然科學方法回答。即使援引《基本法》第1條第1項及第19條第2項\ggfnArtOneNineteenZH{}，亦不能提供進一步說明。解釋「人」此一概念時，仍會出現尚未回答之同一解釋問題。《基本法》第19條第2項並不確定基本權之射程，而是確定其不可侵犯之「核心」；在適用該條前，須先審查未出生生命是否以及在何種程度上確由《基本法》第2條第2項第1句\ggfnArtTwoTwoOneZH{}保護。
\end{transpara}

\syncleft
\begin{sourcepara}
\sourcern{91}Die Frage des Strafrechtsschutzes ungeborenen Lebens verlange differenzierte Antworten, je nachdem gegen welchen potentiellen Täter er sich richte und in welchem Entwicklungsstadium sich das ungeborene Leben befinde. Auch die Rechtsgeschichte zeige, daß das ungeborene Leben - sowohl im weltlichen wie im kirchlichen Recht - immer anders, und zwar schwächer geschützt worden sei als das geborene Leben.
\end{sourcepara}
\begin{transpara}
\transrn{91}刑法對未出生生命的保護問題需要不同的答案，這取決於它所針對的潛在犯罪者以及未出生生命所處的發展階段。法律史也表明，無論在世俗法或教會法中，未出生生命受到的保護始終與出生生命不同，而且保護力度也更弱。
\end{transpara}

\syncleft
\begin{sourcepara}
\sourcern{92}Nach kirchlichem Recht sei bis zum Ende des 19. Jahrhunderts nur die Abtötung einer "beseelten" Leibesfrucht strafbar gewesen, wobei als Zeitpunkt der Beseelung in der Praxis der 80. Tag nach der Empfängnis angenommen worden sei. Das weltliche Recht sei zunächst nachhaltig von der kirchlichen Lehre und vom kirchlichen Recht beeinflußt worden; seit dem Aufkommen naturwissenschaftlichen Denkens habe es auf den Beginn der Kindesbewegungen (im Common Law: quickening) abgestellt, bis zu deren Einsetzen der Schwangerschaftsabbruch wesentlich milder bestraft worden sei. Aus bevölkerungspolitischen Gründen habe das Preußische Allgemeine Landrecht dann zwar den Schwangerschaftsabbruch schon von der Empfängnis an bestraft, die Strafdrohung aber nach der 30. Schwangerschaftswoche, also nach der Lebensfähigkeit des Kindes, abgestuft.
\end{sourcepara}
\begin{transpara}
\transrn{92}根據教會法，直到19世紀末，只有殺死「有靈魂」的胎兒才會受處罰，儘管在實踐中，賦予靈魂的時間被假定為受孕後第80天。世俗法最初深受教會教義和教會法的影響。自然科學思想興起後，則以胎兒運動（Kindesbewegungen，普通法稱 quickening）之開始為判斷基準，在胎動出現之前，終止妊娠所受的刑罰相對較輕。出於人口政策的原因，普魯士普通土地法從受孕開始就對終止妊娠施以刑罰，但在懷孕30週後，即根據孩子的生存能力，逐漸加重刑罰威嚇。
\end{transpara}

\syncleft
\begin{sourcepara}
\sourcern{93}Daß ungeborenes Leben strafrechtlich nicht wie ein Menschenleben zu bewerten sei, zeige sich auch darin, daß selbst die Gegner der Fristenlösung über die medizinische Indikation hinausgehende weitere Indikationstatbestände vorgeschlagen hätten. Von dem irrigen Ausgangspunkt her, ungeborenes Leben sei strafrechtlich ein dem Menschenleben gleichwertiges Rechtsgut, käme diese Strafrechtslehre bei der eugenischen und der ethischen Indikation in Bedrängnis. Letztlich gehe es entscheidend um die Wahrung mütterlicher Interessen, wobei man gleichzeitig bemüht sei, die des nasciturus soweit als möglich zu schützen. Man müsse deshalb die Frage von den Rechten der Mutter her betrachten und die Strafvorschrift des § 218 \abbrStGB{} dogmatisch richtig als ein Gesetz verstehen, daß zum Schutz des objektiven Rechtsguts ungeborenes Leben in Grundrechte der Mutter eingreife. Diese Bewertung des Schwangerschaftsabbruchs sei allgemein im Rechtsbewußtsein der freien Welt zu beobachten, wie ein rechtsvergleichender Blick auf das westliche Ausland zeige.
\end{sourcepara}
\begin{transpara}
\transrn{93}即使是期限規定的反對者也提出了超越醫學指標事由的進一步指標事由，這一事實也表明，根據刑法，未出生的生命不應像人的生命一樣受到重視。基於胎兒生命在刑法上與人的生命同等的法律財產的錯誤出發點，這種刑法理論在涉及優生學和倫理學時就會遇到麻煩。最終，保護母親的利益至關重要，同時盡可能保護胎兒（nasciturus）的利益。因此，必須從母親權利的角度來看待問題，教條式地正確理解《刑法》第218條的刑事規定，認為這是一部侵犯母親基本權利的法律，以保護未出生生命的客觀法益。正如對西方國家法律的比較所顯示的那樣，這種對終止妊娠的評估通常可以在自由世界的法律意識中觀察到。
\end{transpara}

\syncleft
\begin{sourcepara}
\sourcern{94}Sehe man die Lage so, dann stärke die Fristenregelung die Grundrechte der Frau aus \abbrArt{} 1 \abbrAbs{} 1 und \abbrArt{} 2 \abbrAbs{} 1 \abbrGG{}. Das heiße zwar nicht, daß ihre individuelle Freiheit unbeschränkt sein könnte. Es sei jedoch für die Frage des Schwangerschaftsabbruchs zunächst von der Selbstverantwortung der Frau auszugehen, die allerdings nicht im Sinne eines "Verfügungsrechts" der Frau über die Leibesfrucht, die nicht Teil des Körpers der Frau sei, mißverstanden werden dürfe. Eine Strafvorschrift gegen den Schwangerschaftsabbruch verlange von der Frau nicht nur, wie von Dritten, eine bloße Unterlassung, sondern vielmehr das Aufsichnehmen von Gefahren für Leib und Leben, von körperlichen und psychischen Belastungen sowie der Pflichten, die sich später aus dem Muttersein und der Verantwortung für die Erziehung und Betreuung ihrer Kinder ergäben. Ferner garantiere \abbrArt{} 6 \abbrAbs{} 1 \abbrGG{} einen Bereich privater, familiärer Lebensgestaltung, der dem Eingriff des Staates grundsätzlich entzogen sei. \abbrArt{} 6 \abbrAbs{} 1 \abbrGG{} sei heute im Lichte des werdenden Menschenrechts auf Familienplanung zu interpretieren. Weiter ergebe sich aus \abbrArt{} 6 \abbrAbs{} 4 \abbrGG{} eine positive soziale Schutzpflicht des Staates. Schließlich seien in \abbrArt{} 2 \abbrAbs{} 2 Satz 1 \abbrGG{}\ggfnArtTwoTwoOneDE{} Leben und Gesundheit der Mutter im Gegensatz zum ungeborenen Leben nicht bloß als objektive Rechtsgüter, sondern als echte Grundrechte geschützt.
\end{sourcepara}
\begin{transpara}
\transrn{94}若如此觀察情勢，則期限規定會強化婦女依《基本法》第 1 條第 1 項與第 2 條第 1 項所享有的基本權。這固然並不意味著其個人自由可以不受限制。然而，就終止妊娠問題而言，首先應從婦女的自主負責出發；但此點不得被誤解為婦女對胎兒享有「處分權」，因為胎兒並非婦女身體的一部分。反對終止妊娠的刑罰規定，要求婦女承擔的並非如同第三人一般僅是不作為，而毋寧是承受對身體與生命的危險、身體與心理負擔，以及日後由母親身分和對其子女教育、照護責任所產生的義務。此外，《基本法》第 6 條第 1 項保障一個原則上不受國家干預的私人、家庭生活形成領域。今日，《基本法》第 6 條第 1 項應在正在形成中的家庭計畫人權之光照下加以解釋。再者，《基本法》第 6 條第 4 項產生國家的積極社會保護義務。最後，依《基本法》第 2 條第 2 項第 1 句\ggfnArtTwoTwoOneZH{}，母親的生命與健康不同於未出生生命，不僅作為客觀法益受到保護，而且作為真正的基本權受到保護。
\end{transpara}

\syncleft
\begin{sourcepara}
\sourcern{95}Bei Abwägung aller Umstände könne \abbrArt{} 2 \abbrAbs{} 2 Satz 1 \abbrGG{}\ggfnArtTwoTwoOneDE{}, der als Grundsatznorm auch ungeborenes Leben schütze, nicht dahin ausgelegt werden, daß er die durchgehende Pönalisierung des Schwangerschaftsabbruchs vom Beginn des Lebens an vorschreibe.
\end{sourcepara}
\begin{transpara}
\transrn{95}綜合考慮各種情況，《基本法》第二條第二項第一句\ggfnArtTwoTwoOneZH{}作為保護未出生生命的基本規範，不能被解釋為要求從生命之初就將終止妊娠定為刑事犯罪。
\end{transpara}

\syncleft
\begin{sourcepara}
\sourcern{96}Der Gesetzgeber habe endlich den Weg eines positiven Schutzes ungeborenen Lebens betreten. Die vorgesehene Beratung solle der Frau einerseits zeigen, welche Hilfen sie bei Austragung der Schwangerschaft und für die Mutterschaft von der Gesellschaft erwarten könne, und solle ihr andererseits alle medizinischen Gesichtspunkte der Schwangerschaft wie des Schwangerschaftsabbruchs vor Augen führen. Das Strafrechtsreform-Ergänzungsgesetz solle die Beratung durch einen Beratungsanspruch absichern und im übrigen die soziale Sicherung für die werdende Mutter und damit für das werdende Leben verbessern.
\end{sourcepara}
\begin{transpara}
\transrn{96}立法機關終於走上了積極保護未出生生命的道路。一方面，預期的諮詢應該向婦女展示她在懷孕至足月和成為母親方面可以從社會獲得什麼幫助，另一方面，應該讓她了解懷孕和終止妊娠的所有醫學方面。刑法改革補充法應透過諮詢權確保諮詢，並改善準媽媽及其未來生活的社會保障。
\end{transpara}

\syncleft
\begin{sourcepara}
\sourcern{97}Prognosen über die zukünftige Wirkung eines Gesetzes seien naturgemäß mit erheblicher Unsicherheit belastet. Dies gelte insbesondere in einem Bereich, in dem man angesichts der ungewöhnlich hohen Dunkelziffer nur von geschätzten Zahlen ausgehen könne. Auch hier müsse der vom Bundesverfassungsgericht aufgestellte Grundsatz gelten, daß eine gesetzliche Maßnahme nicht schon deshalb als verfassungswidrig angesehen werden dürfe, weil sie auf einer Fehlprognose beruhen könnte. Das Bundesverfassungsgericht habe vielmehr nur zu prüfen, ob die Annahmen des Gesetzgebers so verfehlt seien, daß sie als nicht sachgerecht und und unvertretbar bezeichnet werden müßten.
\end{sourcepara}
\begin{transpara}
\transrn{97}對法律未來影響的預測自然會受到相當大的不確定性的影響。在一個由於未報告病例數量異常多而只能假設估計數字的地區尤其如此。在這裡，聯邦憲法法院制定的原則也必須適用，即一項法律措施不應僅僅因為它可能基於不正確的預測而被視為違憲。相反，聯邦憲法法院只需審查立法機關的假設是否錯誤到必須被描述為不恰當和不合理的程度。
\end{transpara}

\bisubsubsection{IV.}{IV.}
\syncleft
\begin{sourcepara}
\sourcern{98}In der mündlichen Verhandlung am 18. und 19. November haben sich zu den verfassungsrechtlichen Fragen geäußert:
\end{sourcepara}
\begin{transpara}
\transrn{98}在11月18日至19日的言詞辯論會上，以下人士就憲法問題發表評論：
\end{transpara}

\syncleft
\begin{sourcepara}
\sourcern{99}Für die Antragsteller der Justizminister des Landes Baden-Württemberg, \abbrDr{} Bender, der Minister für Familie, Gesundheit und Sozialordnung des Saarlandes, Frau Rita Waschbuesch, die Professoren \abbrDr{} Lerche, \abbrDr{} Ossenbühl und \abbrDr{} Rudolphi, Rechtsanwalt Erhard (MdB) sowie die Ministerialdirigenten Professor \abbrDr{} Odersky (Bayern) und \abbrDr{} Fischer (Rheinland-Pfalz), für den Deutschen Bundestag dessen Vizepräsidentin Frau Liselotte Funcke, Rechtsanwalt Professor \abbrDr{} Ehmke (MdB) und Professor \abbrDr{} Stratenwerth, für die Bundesregierung der Bundesminister der Justiz \abbrDr{} Vogel, die Professoren \abbrDr{} Baumann und \abbrDr{} Peter Schneider sowie Ministerialdirektor Bahlmann und Ministerialrat Harsdorf.
\end{sourcepara}
\begin{transpara}
\transrn{99}代表聲請人：巴登-符騰堡邦司法部長 Dr. Bender，薩爾邦家庭、衛生及社會事務部長 Rita Waschbuesch 女士，Prof. Dr. Lerche、Prof. Dr. Ossenbühl、Prof. Dr. Rudolphi，律師 Erhard（聯邦議院議員），以及部長級主任 Prof. Dr. Odersky（巴伐利亞）和 Dr. Fischer（萊茵蘭-普法爾茨）；代表德國聯邦議院：副議長 Liselotte Funcke 女士，律師 Prof. Dr. Ehmke（聯邦議院議員），Prof. Dr. Stratenwerth；代表聯邦政府：聯邦司法部長 Dr. Vogel、Prof. Dr. Baumann、Prof. Dr. Peter Schneider，以及部長級主任 Bahlmann 和部長級參事 Harsdorf。
\end{transpara}

\syncleft
\begin{sourcepara}
\sourcern{100}Auf Anregung der Bundesregierung ist ferner Professor \abbrDr{}\abbrDr{} Jürgens als Auskunftsperson angehört worden.
\end{sourcepara}
\begin{transpara}
\transrn{100}應聯邦政府建議，Prof. Dr. Jürgens 亦以提供資訊之人的身分接受聽取意見。
\end{transpara}

\bisubsection{B.}{B.}
\syncleft
\begin{sourcepara}
\sourcern{101}Das Fünfte Strafrechtsreformgesetz bedurfte nicht der Zustimmung des Bundesrates.
\end{sourcepara}
\begin{transpara}
\transrn{101}《第五刑法改革法》無須經聯邦參議院同意。
\end{transpara}

\syncleft
\begin{sourcepara}
\sourcern{102}1. Das Gesetz ändert zwar in den Artikeln 6 und 7 die Strafprozeßordnung und das Einführungsgesetz zum Strafgesetzbuch, die ihrerseits mit Zustimmung des Bundesrates ergangen sind. Deswegen allein ist es aber noch nicht zustimmungsbedürftig (\abbrBVerfGE{} 37, 363). Das Gesetz ändert ferner keine gesetzlichen Bestimmungen, die ihrerseits zustimmungsbedürftig waren.
\end{sourcepara}
\begin{transpara}
\transrn{102}1. 該法固然於第6條及第7條修正《刑事訴訟法》及《刑法施行法》，而此二法本身均係經聯邦參議院同意而制定。然而，僅因此尚不使該法須經同意（BVerfGE 37, 363）。此外，該法亦未修正任何本身須經同意之法律規定。
\end{transpara}

\syncleft
\begin{sourcepara}
\sourcern{103}2. Das Fünfte Strafrechtsreformgesetz enthält selbst keine gemäß \abbrArt{} 84 \abbrAbs{} 1 oder einer anderen Bestimmung des Grundgesetzes zustimmungsbedürftigen Vorschriften. Weder § 218c noch § 219 \abbrStGB{} \abbrNF{} regeln die Einrichtung von Behörden oder das Verwaltungsverfahren. Sie setzen vielmehr lediglich die materiell-rechtlich Voraussetzungen eines nicht strafbaren Schwangerschaftsabbruchs fest. Das gilt auch, soweit § 218c \abbrAbs{} 1 \abbrNr{} 1 \abbrStGB{} verlangt, daß die Schwangere sich vor dem Abbruch an eine ermächtigte Beratungsstelle gewandt hat, und den Gegenstand der Beratung umschreibt. Die Errichtung und Einrichtung von Beratungsstellen sowie der Erlaß von Verwaltungsvorschriften für das von diesen Stellen zu praktizierende Verfahren bleiben in vollem Umfang den Ländern überlassen. Aus den gleichen Gründen löst § 219 \abbrStGB{} nicht das Zustimmungserfordernis des Bundesrates aus, wenn diese Vorschrift vor der Durchführung eines nach § 218b indizierten Schwangerschaftsabbruchs die Bestätigung der sachlichen Voraussetzungen durch "eine zuständige Stelle" fordert.
\end{sourcepara}
\begin{transpara}
\transrn{103}2. 《第五刑法改革法》本身並不包含任何依《基本法》第84條第1項\ggfnArtEightFourOneZH{}或其他規定須經同意之規定。新版本《刑法》第218c條及第219條均未規範機關設置或行政程序。毋寧，其僅確定不罰終止妊娠之實體法要件。即使第218c條第1項第1款要求懷孕婦女於終止妊娠前洽詢經授權之諮詢機構，並描述諮詢標的，亦同。諮詢機構之設立與配置，以及該等機構所實施程序之行政規則發布，完全留由各邦處理。基於同一理由，第219條於要求依第218b條具指標事由之終止妊娠實施前，須由「主管機關」確認實體要件時，並不引發聯邦參議院之同意要求。
\end{transpara}

\syncleft
\begin{sourcepara}
\sourcern{104}3. Den antragstellenden Landesregierungen kann auch nicht gefolgt werden, soweit sie bei der Beurteilung der Zustimmungsbedürftigkeit an den in der Entscheidung vom 25. Juni 1974 aufgestellten Rechtsgrundsatz anknüpfen, wonach ein Änderungsgesetz dann der Zustimmung des Bundesrates bedarf, "wenn durch die Änderung materiell- rechtlicher Normen die nicht ausdrücklich geänderten Vorschriften über das Verwaltungsverfahren bei sinnorientierter Auslegung ihrerseits eine wesentlich andere Bedeutung und Tragweite erfahren" (\abbrBVerfGE{} 37, 363, 4. Leitsatz und [383]). Die tatsächlichen Voraussetzungen für eine unmittelbare Anwendung dieses Leitsatzes liegen hier unbestrittenermaßen nicht vor. Ob eine Fortentwicklung dieses Rechtsgrundsatzes in dem von der Landesregierung von Rheinland-Pfalz vorgetragenen Sinn überhaupt in Betracht kommen könnte, mag dahinstehen. Auch nach dieser Rechtsansicht ergäbe sich keine Zustimmungsbedürftigkeit des Fünften Strafrechtsreformgesetzes, da den Ländern nach den materiell-rechtlichen Vorschriften für die ihnen obliegenden Verwaltungsregelungen noch ein weiter Gestaltungsspielraum verbleibt.
\end{sourcepara}
\begin{transpara}
\transrn{104}3. 聲請之各邦政府在判斷是否須經同意時，援引1974年6月25日裁判所建立之法律原則；依該原則，修正法於「因實體法規範之變更，使未明文修正之行政程序規定在目的取向解釋下，本身取得實質上不同之意義及射程」時，須經聯邦參議院同意（BVerfGE 37, 363，第4項裁判要旨及[383]）。就此亦不能採其見解。此項原則直接適用之前提，在本案無爭議地並不存在。此項法律原則是否可能依萊茵蘭-普法爾茨邦政府所主張之方向進一步發展，得暫置不論。即使依此法律見解，《第五刑法改革法》亦不生須經同意之要求，因為依實體法規定，各邦對其所負行政規制仍保有廣泛形成空間。
\end{transpara}

\syncleft
\begin{sourcepara}
\sourcern{105}4. Schließlich kann auch aus dem engen Zusammenhang des Fünften Strafrechtsreformgesetzes mit dem Strafrechtsreform-Ergänzungsgesetz, das mit dem vom Bundestag beschlossenen Inhalt für zustimmungsbedürftig angesehen wird, kein Zustimmungserfordernis abgeleitet werden. Abgesehen davon, daß das Strafrechtsreform-Ergänzungsgesetz noch nicht zustande gekommen ist, ist der Gesetzgeber in Ausübung seiner gesetzgeberischen Freiheit grundsätzlich nicht gehindert, ein Gesetzesvorhaben in mehrere Einzelgesetze aufzuteilen. Das Bundesverfassungsgericht ist bisher von der Zulässigkeit solcher Aufspaltungen ausgegangen (vgl. \abbrBVerfGE{} 34, 9 [28]; 37, 363 [382]). In der Entscheidung \abbrBVerfGE{} 24, 184 (199 \abbrF{}) - Apostille - hat es das Gericht offengelassen, ob der Aufspaltungsbefugnis verfassungsrechtliche Grenzen gezogen sind und wo diese Grenzen verlaufen. Hier sind solche Grenzen jedenfalls nicht überschritten.
\end{sourcepara}
\begin{transpara}
\transrn{105}4. 最後，亦不能由《第五刑法改革法》與《刑法改革補充法》之緊密關聯，推導出同意要求；後者以聯邦議院決議之內容而言，被認為須經同意。且不論《刑法改革補充法》尚未成立，立法者原則上不因其行使立法自由而被禁止將一項法律計畫分割為數部單行法律。聯邦憲法法院迄今亦以此種分割為容許（參見 BVerfGE 34, 9 [28]; 37, 363 [382]）。在 BVerfGE 24, 184 (199 f.)「Apostille」裁判中，本院未決定分割權限是否受憲法界限拘束，以及該界限何在。無論如何，本案並未逾越此等界限。
\end{transpara}

\syncleft
\begin{sourcepara}
\sourcern{106}Das Fünfte Strafrechtsreformgesetz und das geplante Strafrechtsreform-Ergänzungsgesetz sind zwar aufeinander abgestimmt, sie müssen aber nicht notwendig zu einer gesetzestechnischen Einheit miteinander verbunden werden. Das erstgenannte Gesetz enthält im wesentlichen nur Straf- und Strafverfahrensrecht. Demgegenüber hat das Strafrechtsreform-Ergänzungsgesetz sozial- und arbeitsrechtliche Maßnahmen zum Inhalt. Die inhaltliche Unabhängigkeit des Strafrechtsreform-Ergänzungsgesetzes vom Fünften Strafrechtsreformgesetz ergibt sich augenfällig schon daraus, daß das Strafrechtsreform-Ergänzungsgesetz seinem Wortlaut nach auf alle für die Neuregelung des Schwangerschaftsabbruchs vorgeschlagenen Lösungen, nämlich die "Fristenlösung" und die drei "Indikationslösungen", anwendbar wäre.
\end{sourcepara}
\begin{transpara}
\transrn{106}《刑法改革第五法》與擬定的《刑法改革補充法》是相互協調的，但不一定要合併形成一個法律單位。第一條法律實質上僅包含刑法和刑事訴訟法。相較之下，刑法改革補充法涵蓋了社會法和勞動法措施。刑法改革補充法內容相對於刑法改革第五法的獨立性，從其措詞來看，刑法改革補充法將適用於終止妊娠新規提出的所有方案，即「時限方案」和三個「指標方案」。
\end{transpara}

\bisubsection{C.}{C.}
\syncleft
\begin{sourcepara}
\sourcern{107}Die Frage der rechtlichen Behandlung des Schwangerschaftsabbruchs wird in der Öffentlichkeit seit Jahrzehnten unter mannigfachen Gesichtspunkten diskutiert. In der Tat wirft dieses Phänomen des Soziallebens vielfältige Probleme biologischer, insbesondere humangenetischer, anthropologischer, ferner medizinischer, psychologischer, sozialer, gesellschaftspolitischer und nicht zuletzt ethischer und moraltheologischer Art auf, die Grundfragen menschlicher Existenz berühren. Aufgabe des Gesetzgebers ist es, die aus diesen verschiedenen Sichtweisen entwickelten, unter sich vielseitig verschränkten Argumente zu würdigen, sie durch spezifisch rechtspolitische Überlegungen sowie durch die praktischen Erfahrungen des Rechtslebens zu ergänzen und auf dieser Grundlage die Entscheidung zu gewinnen, in welcher Weise die Rechtsordnung auf diesen sozialen Vorgang reagieren soll. Die nach außergewöhnlich umfangreichen Vorarbeiten im Fünften Strafrechtsreformgesetz getroffene gesetzliche Regelung kann vom Bundesverfassungsgericht allein unter dem Gesichtspunkt geprüft werden, ob sie mit dem Grundgesetz als dem höchsten in der Bundesrepublik geltenden Recht vereinbar ist. Gewicht und Ernst der verfassungsrechtlichen Fragestellung werden deutlich, wenn bedacht wird, daß es hier um den Schutz menschlichen Lebens geht, eines zentralen Wertes jeder rechtlichen Ordnung. Die Entscheidung über Maßstäbe und Grenzen der gesetzgeberischen Entscheidungsfreiheit erfordert eine Gesamtschau des verfassungsrechtlichen Normenbestandes und der in ihm beschlossenen Wertordnung.
\end{sourcepara}
\begin{transpara}
\transrn{107}幾十年來，人們從不同的角度公開討論了終止妊娠的法律待遇問題。事實上，這種社會生活現象引發了各種生物學問題，特別是人類遺傳、人類學、醫學、心理、社會、社會政治以及最後但並非最不重要的倫理和道德神學性質的問題，這些問題觸及人類存在的基本問題。立法機關的任務是理解從這些不同角度產生的多重相互關聯的論點，並以具體的法律政策考慮和法律生活的實踐經驗來補充它們，並在此基礎上決定法律制度應如何應對這一社會進程。經過極其廣泛的準備工作後通過的第五刑法改革法案所通過的法律規定，可以由聯邦憲法法院僅從是否符合作為聯邦共和國適用的最高法律的《基本法》的角度進行審查。當考慮到這裡所涉問題是保護人類生命（任何法律秩序的核心價值）時，憲法問題的重要性和嚴重性就變得清晰起來。決定立法自由的標準和限制需要全面了解憲法規範的主體以及其中所決定的價值秩序。
\end{transpara}

\bisubsubsection{I.}{I.}
\syncleft
\begin{sourcepara}
\sourcern{108}1. \abbrArt{} 2 \abbrAbs{} 2 Satz 1 \abbrGG{}\ggfnArtTwoTwoOneDE{} schützt auch das sich im Mutterleib entwickelnde Leben als selbständiges Rechtsgut.
\end{sourcepara}
\begin{transpara}
\transrn{108}1. 《基本法》第二條第二項第一句\ggfnArtTwoTwoOneZH{}也保護子宮內發育的生命作為獨立的法益。
\end{transpara}

\syncleft
\begin{sourcepara}
\sourcern{109}a) Die ausdrückliche Aufnahme des an sich selbstverständlichen Rechts auf Leben in das Grundgesetz - anders als etwa in der Weimarer Verfassung - erklärt sich hauptsächlich als Reaktion auf die "Vernichtung lebensunwerten Lebens", auf "Endlösung" und "Liquidierung", die vom nationalsozialistischen Regime als staatliche Maßnahmen durchgeführt wurden. \abbrArt{} 2 \abbrAbs{} 2 Satz 1 \abbrGG{}\ggfnArtTwoTwoOneDE{} enthält ebenso wie die Abschaffung der Todesstrafe durch \abbrArt{} 102\ggfnArtOneZeroTwoDE{} \abbrGG{} "ein Bekenntnis zum grundsätzlichen Wert des Menschenlebens und zu einer Staatsauffassung, die sich in betonten Gegensatz zu den Anschauungen eines politischen Regimes stellt, dem das einzelne Leben wenig bedeutete und das deshalb mit dem angemaßten Recht über Leben und Tod des Bürgers schrankenlosen Mißbrauch trieb" (\abbrBVerfGE{} 18, 112 [117]).
\end{sourcepara}
\begin{transpara}
\transrn{109}a) 將不言而喻的生命權明確納入《基本法》——與魏瑪憲法等不同——主要可以解釋為對「毀滅不值得生命的生命」、「最終解決」和「清算」的反應，這些行為是國家社會主義政權作為國家措施實施的。《基本法》第2條第2項第1句，如同透過《基本法》第102條\ggfnArtOneZeroTwoZH{}廢除死刑，包含「對人類生命基本價值和一種國家觀念的承諾，此種國家觀念與以下政治政權的觀點形成鮮明對比——在那個政權中個人生命毫無意義，因此無限制地濫用其所僭奪的對公民生死之支配權」（BVerfGE 18, 112 [117]）。
\end{transpara}

\syncleft
\begin{sourcepara}
\sourcern{110}b) Bei der Auslegung des \abbrArt{} 2 \abbrAbs{} 2 Satz 1 \abbrGG{}\ggfnArtTwoTwoOneDE{} ist auszugehen von seinem Wortlaut: "Jeder hat das Recht auf Leben ... ". Leben im Sinne der geschichtlichen Existenz eines menschlichen Individuums besteht nach gesicherter biologisch-physiologischer Erkenntnis jedenfalls vom 14. Tage nach der Empfängnis (Nidation, Individuation) an (vgl. hierzu die Ausführungen von Hinrichsen vor dem Sonderausschuß für die Strafrechtsreform, 6. \abbrWp{}, 74. Sitzung, \abbrStenBer{} \abbrS{} 2142 \abbrFF{}). Der damit begonnene Entwicklungsprozeß ist ein kontinuierlicher Vorgang, der keine scharfen Einschnitte aufweist und eine genaue Abgrenzung der verschiedenen Entwicklungsstufen des menschlichen Lebens nicht zuläßt. Er ist auch nicht mit der Geburt beendet; die für die menschliche Persönlichkeit spezifischen Bewußtseinsphänomene z.B. treten erst längere Zeit nach der Geburt auf. Deshalb kann der Schutz des \abbrArt{} 2 \abbrAbs{} 2 Satz 1 \abbrGG{}\ggfnArtTwoTwoOneDE{} weder auf den "fertigen" Menschen nach der Geburt noch auf den selbständig lebensfähigen nasciturus beschränkt werden. Das Recht auf Leben wird jedem gewährleistet, der "lebt"; zwischen einzelnen Abschnitten des sich entwickelnden Lebens vor der Geburt oder zwischen ungeborenem und geborenem Leben kann hier kein Unterschied gemacht werden. "Jeder" im Sinne des \abbrArt{} 2 \abbrAbs{} 2 Satz 1 \abbrGG{}\ggfnArtTwoTwoOneDE{} ist "jeder Lebende", anders ausgedrückt: jedes Leben besitzende menschliche Individuum; "jeder" ist daher auch das noch ungeborene menschliche Wesen.
\end{sourcepara}
\begin{transpara}
\transrn{110}b) 解釋《基本法》第2條第2項第1句\ggfnArtTwoTwoOneZH{}時，應從其文義出發：「人人有生命權……」。依確定之生物學與生理學知見，作為人類個體歷史性存在之生命，無論如何自受孕後第14日起（著床、個體化）即已存在（參見 Hinrichsen 於刑法改革特別委員會之陳述，第6屆選期，第74次會議，StenBer. S. 2142 ff.）。由此開始之發展過程是一項連續過程，並無銳利切割點，亦不容精確劃分人類生命之不同發展階段。該過程亦不以出生為終點；例如，人格所特有之意識現象，直至出生後相當時間始出現。因此，《基本法》第2條第2項第1句\ggfnArtTwoTwoOneZH{}之保護，既不能限於出生後「完成」之人，也不能限於具有獨立存活能力之胎兒（nasciturus）。生命權保障每一個「活著」者；此處不能在出生前發展中生命之個別階段之間，或在未出生生命與已出生生命之間作區分。《基本法》第2條第2項第1句\ggfnArtTwoTwoOneZH{}意義下之「人人」即「每一活著者」，換言之，即每一具生命之人類個體；因此，「人人」亦包括尚未出生之人類存在。
\end{transpara}

\syncleft
\begin{sourcepara}
\sourcern{111}c) Gegenüber dem Einwand, "jeder" bezeichne sowohl in der Umgangs- als auch in der Rechtssprache gemeinhin eine "fertige" menschliche Person, eine reine Wortinterpretation spreche daher gegen die Einbeziehung des ungeborenen Lebens in den Wirkungsbereich des \abbrArt{} 2 \abbrAbs{} 2 Satz 1 \abbrGG{}\ggfnArtTwoTwoOneDE{}, ist zu betonen, daß jedenfalls Sinn und Zweck dieser Grundgesetzbestimmung es erfordern, den Lebensschutz auch auf das sich entwickelnde Leben auszudehnen. Die Sicherung der menschlichen Existenz gegenüber staatlichen Übergriffen wäre unvollständig, wenn sie nicht auch die Vorstufe des "fertigen Lebens", das ungeborene Leben, umfaßte.
\end{sourcepara}
\begin{transpara}
\transrn{111}c) 對於「每個人」在日常及法律語言中通常指「已完成」之人，因此單從該詞解釋反對將未出生生命納入《基本法》第2條第2項第1句\ggfnArtTwoTwoOneZH{}保護範圍之意見，必須指出：無論如何，該基本法規定之意義與目的，要求將生命保護延伸至正在發展中之生命。若不包含「已完成生命」之前階段，也就是未出生生命，對人類存在免受國家侵害之保障即不完整。
\end{transpara}

\syncleft
\begin{sourcepara}
\sourcern{112}Diese extensive Auslegung entspricht dem in der Rechtsprechung des Bundesverfassungsgerichts aufgestellten Grundsatz, "wonach in Zweifelsfällen diejenige Auslegung zu wählen ist, welche die juristische Wirkungskraft der Grundrechtsnorm am stärksten entfaltet" (\abbrBVerfGE{} 32, 54 [71]; 6, 55 [72]).
\end{sourcepara}
\begin{transpara}
\transrn{112}這種廣泛的解釋符合聯邦憲法法院判例中確立的原則，「根據該原則，在有疑問的情況下，必須選擇最充分發揮基本權利規範的法律效力的解釋」（BVerfGE 32, 54 [71]; 6, 55 [72]）。
\end{transpara}

\syncleft
\begin{sourcepara}
\sourcern{113}d) Zur Begründung dieses Ergebnisses läßt sich auch die Entstehungsgeschichte des \abbrArt{} 2 \abbrAbs{} 2 Satz 1 \abbrGG{}\ggfnArtTwoTwoOneDE{} heranziehen.
\end{sourcepara}
\begin{transpara}
\transrn{113}d) 為了證明這個結果的合理性，也可以利用《基本法》第二條第二項第一句\ggfnArtTwoTwoOneZH{}的歷史淵源。
\end{transpara}

\syncleft
\begin{sourcepara}
\sourcern{114}Nachdem die Fraktion der Deutschen Partei (DP) wiederholt den Antrag gestellt hatte, im Zusammenhang mit dem Recht auf Leben und körperliche Unversehrtheit auch das "keimende Leben" ausdrücklich zu erwähnen (Drucks. des Parlamentarischen Rates 11.48 - 298 und 12.48 - 398), beriet der Parlamentarische Rat erstmalig in der 32. Sitzung seines Ausschusses für Grundsatzfragen am 11. Januar 1949 diesen Problemkreis. Bei der Erörterung der Frage, ob in das Grundgesetz eine Vorschrift aufgenommen werden solle, die ärztliche Eingriffe verbietet, die nicht der Heilung dienen, erklärte der Abgeordnete \abbrDr{} Heuss (FDP), ohne damit auf Widerspruch zu stoßen, gedacht sei hier an die Zwangssterilisation und beim Recht auf Leben an die Abtreibung. Der Hauptausschuß des Parlamentarischen Rates befaßte sich in seiner 42. Sitzung am 18. Januar 1949 bei der zweiten Lesung über die Grundrechte (Verhandlungen des Hauptausschusses des Parlamentarischen Rates, \abbrStenBer{} \abbrS{} 529 \abbrFF{}) eingehender mit der Frage nach der Einbeziehung des werdenden Lebens in den Schutz der Verfassung. Der Abgeordnete \abbrDr{} Seebohm (DP) beantragte, dem damaligen \abbrArt{} 2 \abbrAbs{} 1 \abbrGG{} die beiden Sätze anzufügen: "Das keimende Leben wird geschützt" und "Die Todesstrafe wird abgeschafft". Dazu führte \abbrDr{} Seebohm (\abbrAAO{}, \abbrS{} 533 \abbrF{}) aus, das Recht auf Leben und körperliche Unversehrtheit umfasse möglicherweise nicht unbedingt auch das keimende Leben. Deshalb müsse es hier besonders erwähnt werden. Mindestens müsse man aber, wenn eine andere Auffassung vorliegen sollte, hierzu ausdrücklich zu Protokoll geben, daß bei dem Recht auf Leben und körperliche Unversehrtheit das keimende Leben ausdrücklich eingeschlossen sei.
\end{sourcepara}
\begin{transpara}
\transrn{114}在德國黨（DP）派多次要求在生命權和身體完整權中明確提及「萌芽生命」（議會理事會印本11.48 - 298和12.48 - 398）後，議會理事會在1949年1月11日的基本問題委員會第32次會議上首次討論了這一問題領域。為了將這一問題納入《基本法》，該法禁止無助於治癒目的的醫療干預措施，國會議員Dr. 赫斯（自民黨）在沒有遇到任何矛盾的情況下，在這裡考慮強制絕育，以及在生命權的情況下考慮終止妊娠。在 1949 年 1 月 18 日舉行的第 42 次會議上，議會理事會主要委員會在基本權利二讀期間更詳細地討論了將正在發展中的生命納入憲法保護的問題（議會理事會主要委員會的談判，StenBer，第 529 頁及以下）。副議員Dr. 西博姆（民主黨）要求在當時的《基本法》第二條第一項中加入「保護萌芽的生命」和「廢除死刑」這兩句話。這導致 Dr. Seebohm（前述，第 533 頁）指出，生命權和身體健全權不一定包括生命的萌芽。這就是為什麼它值得在這裡特別提及。然而，如果有不同的觀點，至少必須明確指出，生命權和身體完整權明確包括生命的萌芽。
\end{transpara}

\syncleft
\begin{sourcepara}
\sourcern{115}Die Abgeordnete Frau \abbrDr{} Weber erklärte namens der CDU/CSU, daß ihre Fraktion, wenn sie für das Recht auf Leben eintrete, das Leben schlechthin meine und für sie auch das keimende Leben und vor allem der Schutz des keimenden Lebens darin enthalten sei (\abbrAAO{}, \abbrS{} 534). \abbrDr{} Heuss (FDP) stimmte an sich mit Frau \abbrDr{} Weber überein, daß der Begriff des Lebens auch das keimende Leben mit umfasse; in die Verfassung sollten aber nicht Dinge hineingenommen werden, die im Strafgesetz geregelt seien. Infolgedessen halte er die Erwähnung sowohl des keimenden Lebens als auch der Todesstrafe als Sonderfrage für überflüssig (\abbrAAO{}, \abbrS{} 535).
\end{sourcepara}
\begin{transpara}
\transrn{115}國會議員Dr. 韋伯代表基民盟/基社盟解釋說，當她的派系主張生命權時，這意味著生命本身，對她來說，它還包括萌芽的生命，最重要的是保護萌芽的生命（loc.cit.，第534頁）。Dr. 赫斯（自民黨）同意 Dr. 韋伯的觀點，認為生命的概念還包括萌芽生命；但是，刑法規定的事項不應寫入憲法。因此，他認為將萌芽生命和死刑作為一個特殊問題提及是多餘的（loc. cit.，第 535 頁）。
\end{transpara}

\syncleft
\begin{sourcepara}
\sourcern{116}"Nach den unwidersprochenen Erklärungen, wonach das keimende Leben in das Recht auf Leben und körperliche Unversehrtheit einbezogen ist", wollte \abbrDr{} Seebohm seinen Antrag zurückziehen (\abbrAAO{}, \abbrS{} 535). Indessen erklärte der Abgeordnete \abbrDr{} Greve (SPD): "Ich muß hier ausdrücklich zu Protokoll geben, daß zum mindesten, was mich angeht, ich unter dem Recht auf Leben nicht auch das Recht auf das keimende Leben verstehe. Ich darf auch für meine Freunde, zum mindesten in ihrer sehr großen Mehrzahl, eine Erklärung gleichen Inhaltes abgeben, um protokollarisch festzuhalten, daß der Hauptausschuß des Parlamentarischen Rates in seiner Gesamtheit nicht auf dem von Herrn Kollegen \abbrDr{} Seebohm soeben zum Ausdruck gebrachten Standpunkt steht". Der darauf von \abbrDr{} Seebohm wieder aufgenommene Antrag wurde zwar mit 11 gegen 7 Stimmen abgelehnt (\abbrAAO{}, \abbrS{} 535). In dem Schriftlichen Bericht des Hauptausschusses (\abbrS{} 7) führte der Abgeordnete \abbrDr{} von Mangoldt (CDU) jedoch zu \abbrArt{} 2 \abbrGG{} aus: "Dabei hat mit der Gewährleistung des Rechts auf Leben auch das keimende Leben geschützt werden sollen. Von der Deutschen Partei im Hauptausschuß eingebrachte Anträge, einen besonderen Satz über den Schutz des keimenden Lebens einzufügen, haben nur deshalb keine Mehrheit gefunden, weil nach der im Ausschuß vorherrschenden Auffassung das zu schützende Gut bereits durch die gegenwärtige Fassung gesichert war".
\end{sourcepara}
\begin{transpara}
\transrn{116}在「關於萌芽中生命已被納入生命權與身體完整權之無異議說明」之後，Dr. Seebohm 本欲撤回其提案（同前，第535頁）。然而，Dr. Greve（SPD）議員表示：「我必須在此明確記入議事錄：至少就我而言，我所理解的生命權並不也包括萌芽中生命的權利。我也可以代表我的同僚，至少代表其中非常大多數，作成同一內容之聲明，以便在議事錄中確認，議會委員會主委員會全體並非站在 Dr. Seebohm 同僚剛才所表達的立場上。」其後 Dr. Seebohm 重新提出的提案，固然以11票對7票遭否決（同前，第535頁）。然而，在主委員會書面報告（第7頁）中，Dr. von Mangoldt（CDU）議員就《基本法》第2條說明：「在保障生命權時，也應保護萌芽中生命。德國黨在主委員會提出、要求加入一個關於保護萌芽中生命之特別句子的提案，之所以未能獲得多數，僅是因為依委員會中占支配地位的見解，應受保護之法益已由現行文字加以保障。」
\end{transpara}

\syncleft
\begin{sourcepara}
\sourcern{117}Das Plenum des Parlamentarischen Rates stimmte dem \abbrArt{} 2 \abbrAbs{} 2 \abbrGG{}\ggfnArtTwoTwoDE{} am 6. Mai 1949 in zweiter Lesung bei 2 Gegenstimmen zu. Bei der dritten Lesung am 8. Mai 1949 brachten sowohl der Abgeordnete \abbrDr{} Seebohm als auch die Abgeordnete \abbrDr{} Weber zum Ausdruck, daß nach ihrer Auffassung \abbrArt{} 2 \abbrAbs{} 2 \abbrGG{}\ggfnArtTwoTwoDE{} auch das keimende Leben in den Schutz dieses Grundrechts einbeziehe (Verhandlungen des Parlamentarischen Rates, \abbrStenBer{} \abbrS{} 218, 223). Beide Redner blieben mit ihren Ausführungen ohne Widerspruch.
\end{sourcepara}
\begin{transpara}
\transrn{117}1949年5月6日，議會全體會議二讀以2票反對通過《基本法》第二條第二項\ggfnArtTwoTwoZH{}。在1949年5月8日的三讀中，議員Dr. 西博姆和議員 Dr. 韋伯均表示，她認為《基本法》第二條第二項\ggfnArtTwoTwoZH{}也包括保護這一基本權利的萌芽生命（《議會理事會的談判》，StenBer.第218、223頁）。兩位發言者的發言均無矛盾。
\end{transpara}

\syncleft
\begin{sourcepara}
\sourcern{118}Die Entstehungsgeschichte des \abbrArt{} 2 \abbrAbs{} 2 Satz 1 \abbrGG{}\ggfnArtTwoTwoOneDE{} legt es somit nahe, daß die Formulierung "jeder hat das Recht auf Leben" auch das "keimende" Leben einschließen sollte. Jedenfalls kann aus den Materialien noch weniger für die gegenteilige Ansicht abgeleitet werden. Andererseits ergibt sich aus ihr kein Anhaltspunkt für die Beantwortung der Frage, ob das ungeborene Leben strafrechtlich geschützt werden muß.
\end{sourcepara}
\begin{transpara}
\transrn{118}因此，《基本法》第二條第二項第一句\ggfnArtTwoTwoOneZH{}的淵源表明，「人人有生命權」也應包括生命的「萌芽」。無論如何，從材料中可以推斷出更少的東西來支持相反的觀點。另一方面，它也沒有為回答未出生生命是否需要刑法保護的問題提供任何依據。
\end{transpara}

\syncleft
\begin{sourcepara}
\sourcern{119}e) Bei den Beratungen des Fünften Strafrechtsreformgesetzes bestand im übrigen Einigkeit über die Schutzwürdigkeit des ungeborenen Lebens, wobei allerdings die verfassungsrechtliche Problematik nicht abschließend behandelt wurde. In dem Bericht des Sonderausschusses für die Strafrechtsreform zu dem von den Fraktionen der SPD und FDP eingebrachten Gesetzentwurf heißt es hierzu u.a.:
\end{sourcepara}
\begin{transpara}
\transrn{119}e) 在第五刑法改革法案的審議過程中，人們一致認為需要保護未出生的生命，儘管憲法問題沒有最終解決。刑法改革特別委員會關於社民黨和自民黨兩派提交的法律草案的報告除其他外指出：
\end{transpara}

\syncleft
\begin{sourcepara}
\sourcern{120}\sourcequoteinner{"Das ungeborene Leben ist ein Rechtsgut, das geborenem grundsätzlich gleich zu achten ist.\quoteparbreak
Diese Feststellung versteht sich für das Stadium, in dem das ungeborene Leben auch außerhalb des Mutterleibes lebensfähig wäre, von selbst. Sie ist aber bereits für das frühere, etwa mit dem 14. Tag nach der Empfängnis beginnende Entwicklungsstadium gerechtfertigt, wie u.a. Hinrichsen in der Öffentlichen Anhörung (AP VI \abbrS{} 2142 \abbrFF{}) überzeugend begründet hat ... . Daß in der gesamten späteren Entwicklung keine diesem Vorgang entsprechende Zäsur mehr festzustellen sei, ist die ganz überwiegende Auffassung in der medizinischen, anthropologischen und theologischen Wissenschaft ... .\quoteparbreak
Damit verbietet es sich, das ungeborene Leben ab dem Ende der Nidation zu negieren oder auch nur mit Indifferenz zu betrachten. Dabei braucht die in der Literatur umstrittene Frage, ob und ggf. inwieweit das Grundgesetz es in seinen Schutz einbeziehe, an dieser Stelle nicht beantwortet zu werden. Jedenfalls entspricht es, sieht man von extremen Auffassungen einzelner Gruppen ab, dem Rechtsverständnis der Allgemeinheit, das ungeborene Leben als Rechtsgut von hohem Rang zu bewerten. Dieses Rechtsverständnis liegt auch dem Entwurf zu Grunde." (\abbrBTDrucks{}7/1981 neu, \abbrS{} 5)}
\end{sourcepara}
\begin{transpara}
\transrn{120}\transquoteinner{「未出生生命是一項法益，原則上應與已出生者同等尊重。\quoteparbreak
此一認定，對於未出生生命即使在母胎之外亦已具存活能力之階段而言，自不待言。然而，即使對於較早、約自受孕後第14日起算之發展階段，亦屬正當；Hinrichsen 等人於公開聽證中（AP VI S. 2142 ff.）已提出有說服力之論證……。在其後整體發展中，已無法再確認與此過程相當之切割點，此為醫學、人類學及神學界絕大多數見解……。\quoteparbreak
因此，自著床（Nidation）結束之時起，即不得否定未出生生命，甚至不得以漠然態度視之。至於文獻中有爭議之問題，即《基本法》是否以及在何種範圍內將其納入保護，於此無須回答。無論如何，除個別群體之極端見解外，將未出生生命評價為高位階法益，符合一般法律認知。此一法律認知亦為草案之基礎。」（BTDrucks. 7/1981 [neu], S. 5）}
\end{transpara}

\syncleft
\begin{sourcepara}
\sourcern{121}Nahezu gleichlautend sind insoweit die Ausschußberichte zu den übrigen Entwürfen (\abbrBTDrucks{} 7/1982 \abbrS{} 5; \abbrBTDrucks{} 7/1983 \abbrS{} 5; \abbrBTDrucks{} 7/1984 neu, \abbrS{} 4).
\end{sourcepara}
\begin{transpara}
\transrn{121}其他草案之委員會報告在此方面措辭幾乎相同（BTDrucks. 7/1982 S. 5；BTDrucks. 7/1983 S. 5；BTDrucks. 7/1984 [neu], S. 4）。
\end{transpara}

\syncleft
\begin{sourcepara}
\sourcern{122}2. Die Pflicht des Staates, jedes menschliche Leben zu schützen, läßt sich deshalb bereits unmittelbar aus \abbrArt{} 2 \abbrAbs{} 2 Satz 1 \abbrGG{}\ggfnArtTwoTwoOneDE{} ableiten. Sie ergibt sich darüber hinaus auch aus der ausdrücklichen Vorschrift des \abbrArt{} 1 \abbrAbs{} 1 Satz 2 \abbrGG{}\ggfnArtOneOneTwoDE{}; denn das sich entwickelnde Leben nimmt auch an dem Schutz teil, den \abbrArt{} 1 \abbrAbs{} 1 \abbrGG{} der Menschenwürde gewährt. Wo menschliches Leben existiert, kommt ihm Menschenwürde zu; es ist nicht entscheidend, ob der Träger sich dieser Würde bewußt ist und sie selbst zu wahren weiß. Die von Anfang an im menschlichen Sein angelegten potentiellen Fähigkeiten genügen, um die Menschenwürde zu begründen.
\end{sourcepara}
\begin{transpara}
\transrn{122}2. 因此，國家保護每一人類生命之義務，已可直接由《基本法》第2條第2項第1句\ggfnArtTwoTwoOneZH{}導出。此外，該義務亦來自《基本法》第1條第1項第2句之明文規定；因為正在發展中生命亦分享《基本法》第1條第1項對人性尊嚴所給予之保護。凡有人類生命存在之處，即具有人性尊嚴；其承載者是否意識到此一尊嚴並懂得自行維護，並非決定性。自人之存在初始即已內含之潛在能力，已足以奠定人性尊嚴。
\end{transpara}

\syncleft
\begin{sourcepara}
\sourcern{123}3. Hingegen braucht die im vorliegenden Verfahren wie auch in der Rechtsprechung und im wissenschaftlichen Schrifttum umstrittene Frage nicht entschieden zu werden, ob der nasciturus selbst Grundrechtsträger ist oder aber wegen mangelnder Rechts- und Grundrechtsfähigkeit "nur" von den objektiven Normen der Verfassung in seinem Recht auf Leben geschützt wird. Nach der ständigen Rechtsprechung des Bundesverfassungsgerichts enthalten die Grundrechtsnormen nicht nur subjektive Abwehrrechte des Einzelnen gegen den Staat, sondern sie verkörpern zugleich eine objektive Wertordnung, die als verfassungsrechtliche Grundentscheidung für alle Bereiche des Rechts gilt und Richtlinien und Impulse für Gesetzgebung, Verwaltung und Rechtsprechung gibt (\abbrBVerfGE{} 7, 198 [205] - Lüth -; 35, 79 [114] - Hochschulurteil - mit weiteren Nachweisen). Ob und gegebenenfalls in welchem Umfang der Staat zu rechtlichem Schutz des werdenden Lebens von Verfassungs wegen verpflichtet ist, kann deshalb schon aus dem objektiv-rechtlichen Gehalt der grundrechtlichen Normen erschlossen werden.
\end{sourcepara}
\begin{transpara}
\transrn{123}3. 另一方面，胎兒（nasciturus）本身是否為基本權主體，抑或因欠缺權利能力與基本權能力而「僅」由憲法的客觀規範保護其生命權，這一在本程序中以及判例與學說上均有爭議的問題，無須作成決定。依聯邦憲法法院一貫判例，基本權規範不僅包含個人對抗國家的主觀防禦權，同時也體現一套客觀價值秩序；此一價值秩序作為憲法上的基本決定，適用於法律的一切領域，並為立法、行政與司法提供準則與推動力（BVerfGE 7, 198 [205] - Lüth -；35, 79 [114] - Hochschulurteil - 並附進一步引注）。因此，國家是否以及在何種範圍內基於憲法負有法律上保護發展中生命的義務，已可從基本權規範的客觀法內容加以推導。
\end{transpara}

\bisubsubsection{II.}{II.}
\syncleft
\begin{sourcepara}
\sourcern{124}1. Die Schutzpflicht des Staates ist umfassend. Sie verbietet nicht nur - selbstverständlich - unmittelbare staatliche Eingriffe in das sich entwickelnde Leben, sondern gebietet dem Staat auch, sich schützend und fördernd vor diese Leben zu stellen, das heißt vor allem, es auch vor rechtswidrigen Eingriffen von seiten anderer zu bewahren. An diesem Gebot haben sich die einzelnen Bereiche der Rechtsordnung, je nach ihrer besonderen Aufgabenstellung, auszurichten. Die Schutzverpflichtung des Staates muß um so ernster genommen werden, je höher der Rang des in Frage stehenden Rechtsgutes innerhalb der Wertordnung des Grundgesetzes anzusetzen ist. Das menschliche leben stellt, wie nicht näher begründet werden muß, innerhalb der grundgesetzlichen Ordnung einen Höchstwert dar; es ist die vitale Basis der Menschenwürde und die Voraussetzung aller anderen Grundrechte.
\end{sourcepara}
\begin{transpara}
\transrn{124}1. 國家的保護義務是全面的。它不僅當然禁止國家直接干預發展中的生命，也命令國家以保護並促進的方式站在此一生命之前；也就是說，首先也要保護其免受他方的違法干預。法律秩序的各個領域，均須依其特殊任務取向，依循此一要求。系爭法益在《基本法》價值秩序中所應賦予的位階越高，國家的保護義務就越必須受到嚴肅對待。人的生命，毋庸詳加論證，在《基本法》秩序內構成最高價值；它是人的尊嚴之生命基礎，也是所有其他基本權的前提。
\end{transpara}

\syncleft
\begin{sourcepara}
\sourcern{125}2. die Verpflichtung des Staates, das sich entwickelnde Leben in Schutz zu nehmen, besteht grundsätzlich auch gegenüber der Mutter. Unzweifelhaft begründet die natürliche Verbindung des ungeborenen Lebens mit dem der Mutter eine besonders geartete Beziehung, für die es in anderen Lebenssachverhalten keine Parallele gibt. Die Schwangerschaft gehört zur Intimsphäre der Frau, deren Schutz durch \abbrArt{} 2 \abbrAbs{} 1 in Verbindung mit \abbrArt{} 1 \abbrAbs{} 1 \abbrGG{} verfassungsrechtlich verbürgt ist. Wäre der Embryo nur als Teil des mütterlichen Organismus anzusehen, so würde auch der Schwangerschaftsabbruch in dem Bereich privater Lebensgestaltung verbleiben, in den einzudringen dem Gesetzgeber verwehrt ist (\abbrBVerfGE{} 6, 32 [41]; 6, 389 [433]; 27, 344 [350]; 32, 373 [379]). Da indessen der nasciturus ein selbständiges menschliches Wesen ist, das unter dem Schutz der Verfassung steht, kommt dem Schwangerschaftsabbruch eine soziale Dimension zu, die ihn der Regelung durch den Staat zugänglich und bedürftig macht. Das Recht der Frau auf freie Entfaltung ihrer Persönlichkeit, welches die Handlungsfreiheit im umfassenden Sinn zum Inhalt hat und damit auch die Selbstverantwortung der Frau umfaßt, sich gegen eine Elternschaft und die daraus folgenden Pflichten zu entscheiden, kann zwar ebenfalls Anerkennung und Schutz beanspruchen. Dieses Recht ist aber nicht uneingeschränkt gewährt - die Rechte anderer, die verfassungsmäßige Ordnung, das Sittengesetz begrenzen es. Von vornherein kann es niemals die Befugnis umfassen, in die geschützte Rechtssphäre eines anderen ohne rechtfertigenden Grund einzugreifen oder sie gar mit dem Leben selbst zu zerstören, am wenigsten dann, wenn nach der Natur der Sache eine besondere Verantwortung gerade für dieses Leben besteht.
\end{sourcepara}
\begin{transpara}
\transrn{125}2. 國家保護發展中生命的義務，原則上也存在於對母親的關係中。毫無疑問，未出生生命與母親生命之間的自然連結，形成了一種特殊性質的關係，在其他生活事實中並無可相比擬之處。妊娠屬於婦女的私密領域，其保護由《基本法》第 2 條第 1 項連結第 1 條第 1 項在憲法上予以保障。倘若胚胎僅被視為母體有機體的一部分，則終止妊娠也會停留在私人生活形成的領域內，而立法者不得侵入該領域（BVerfGE 6, 32 [41]; 6, 389 [433]; 27, 344 [350]; 32, 373 [379]）。然而，由於胎兒（nasciturus）是一個獨立的人類存在，並受憲法保護，終止妊娠便具有一種社會面向，因而使其成為國家得以並且有必要加以規範的事項。婦女人格自由發展的權利，其內容即為廣義意義下的行動自由，並因此也包括婦女自行決定不成為父母、以及不承擔由此而生義務的自主責任；此一權利固然同樣可以請求承認與保護。然而，該權利並非毫無限制地受到保障──他人的權利、合憲秩序與道德律均對其加以限制。從根本上說，該權利絕不可能包含一種權限，使人得以在沒有正當化理由的情況下侵入他人受保護的權利領域，甚或以生命本身為代價將該領域予以摧毀；尤其是在依事情本質而言，正是對此一生命本身負有特殊責任時，更不可能如此。
\end{transpara}

\syncleft
\begin{sourcepara}
\sourcern{126}Ein Ausgleich, der sowohl den Lebensschutz des nasciturus gewährleistet als auch der Schwangeren die Freiheit des Schwangerschaftsabbruchs beläßt, ist nicht möglich, da Schwangerschaftsabbruch immer Vernichtung des ungeborenen Lebens bedeutet. Bei der deshalb erforderlichen Abwägung "sind beide Verfassungswerte in ihrer Beziehung zur Menschenwürde als dem Mittelpunkt des Wertsystems der Verfassung zu sehen" (\abbrBVerfGE{} 35, 202 [225]). Bei einer Orientierung an \abbrArt{} 1 \abbrAbs{} 1 \abbrGG{} muß die Entscheidung zugunsten des Vorrangs des Lebensschutzes für die Leibesfrucht vor dem Selbstbestimmungsrecht der Schwangeren fallen. Diese kann durch Schwangerschaft, Geburt und Kindeserziehung in manchen persönlichen Entfaltungsmöglichkeiten beeinträchtigt sein. Das ungeborene Leben hingegen wird durch den Schwangerschaftsabbruch vernichtet. Nach dem Prinzip des schonendsten Ausgleichs konkurrierender grundgesetzlich geschützter Positionen unter Berücksichtigung des Grundgedankens des \abbrArt{} 19 \abbrAbs{} 2 \abbrGG{} muß deshalb dem Lebensschutz des nasciturus der Vorzug gegeben werden. Dieser Vorrang gilt grundsätzlich für die gesamte Dauer der Schwangerschaft und darf auch nicht für eine bestimmte Frist in Frage gestellt werden. Die bei der dritten Beratung des Strafrechtsreformgesetzes im Bundestag geäußerte Meinung, es gehe darum, den Vorrang " des aus der Menschenwürde fließenden Selbstbestimmungsrechtes der Frau gegenüber allem anderen, auch dem Lebensrecht des Kindes, für eine bestimmte Frist herauszustellen" (Deutscher Bundestag, 7.\abbrWp{}, 96. Sitzung, \abbrStenBer{} \abbrS{} 6492), ist mit der grundgesetzlichen Wertordnung nicht vereinbar.
\end{sourcepara}
\begin{transpara}
\transrn{126}一種既能保障胎兒（nasciturus）的生命保護、同時又保留孕婦終止妊娠自由的平衡，是不可能存在的，因為終止妊娠始終意味著未出生生命的毀滅。在因此所必須進行的權衡中，「兩種憲法價值都必須在其與人的尊嚴之關係中被觀察，而人的尊嚴乃是憲法價值體系的中心」（BVerfGE 35, 202 [225]）。若以《基本法》第 1 條第 1 項作為價值導向，則決定必須作成有利於胎兒生命保護優先於孕婦自決權的結果。孕婦可能因妊娠、生產與子女養育，而在若干個人人格發展可能性上受到影響。相對地，未出生生命則因終止妊娠而遭到毀滅。因此，依競合之基本法保護地位間最溫和調和原則，並考量《基本法》第 19 條第 2 項的基本思想，必須給予胎兒（nasciturus）生命保護以優先地位。此一優先原則上適用於整個妊娠期間，也不得因某一特定期間而被置於疑問之中。聯邦議院在《刑法改革法》第三讀時所表達的見解，認為其目的在於就某一特定期間凸顯「源自人的尊嚴之婦女自決權，相對於其他一切，包括相對於兒童生命權在內，所具有的優先地位」（德國聯邦議院，第 7 屆選期，第 96 次會議，StenBer. S. 6492），與《基本法》的價值秩序並不相容。
\end{transpara}

\syncleft
\begin{sourcepara}
\sourcern{127}3. Von hier aus erschließt sich die von der Verfassung geforderte Grundhaltung der Rechtsordnung zum Schwangerschaftsabbruch: Die Rechtsordnung darf nicht das Selbstbestimmungsrecht der Frau zur alleinigen Richtschnur ihrer Regelungen machen. Der Staat muß grundsätzlich von einer Pflicht zur Austragung der Schwangerschaft ausgehen, ihren Abbruch also grundsätzlich als Unrecht ansehen. In der Rechtsordnung muß die Mißbilligung des Schwangerschaftsabbruchs klar zum Ausdruck kommen. Es muß der falsche Eindruck vermieden werden, als handle es sich beim Schwangerschaftsabbruch um den gleichen sozialen Vorgang wie etwa den Gang zum Arzt zwecks Heilung einer Krankheit oder gar um eine rechtlich irrelevante Alternative zur Empfängnisverhütung. Der Staat darf sich seiner Verantwortung auch nicht durch Anerkennung eines "rechtsfreien Raumes" entziehen, indem er sich der Wertung enthält und diese der eigenverantwortlichen Entscheidung des Einzelnen überläßt.
\end{sourcepara}
\begin{transpara}
\transrn{127}3. 由此可以推導出憲法對法律秩序就終止妊娠所要求的基本態度：法律秩序不得將婦女的自決權作為其規範的唯一準繩。國家原則上必須以懷孕應被繼續至分娩的義務為出發點，亦即原則上將終止妊娠視為不法。法律秩序中必須清楚表達對終止妊娠的否定評價。必須避免造成一種錯誤印象，彷彿終止妊娠與為治療疾病而就醫屬於同一類社會過程，甚至彷彿它只是避孕在法律上無關緊要的替代方案。國家也不得藉由承認一個「無法律空間」（rechtsfreier Raum）而逃避其責任，也就是不得不作評價，並將該評價交由個人自行負責決定。
\end{transpara}

\bisubsubsection{III.}{III.}
\syncleft
\begin{sourcepara}
\sourcern{128}Wie der Staat seine Verpflichtung zu einem effektiven Schutz des sich entwickelnden Lebens erfüllt, ist in erster Linie vom Gesetzgeber zu entscheiden. Er befindet darüber, welche Schutzmaßnahmen er für zweckdienlich und geboten hält, um einen wirksamen Lebensschutz zu gewährleisten.
\end{sourcepara}
\begin{transpara}
\transrn{128}國家如何履行有效保護發展中生命之義務，首先應由立法機關決定。由立法機關判斷，哪些保護措施為確保有效生命保護所合目的且屬必要。
\end{transpara}

\syncleft
\begin{sourcepara}
\sourcern{129}1. Dabei gilt auch und erst recht für den Schutz des ungeborenen Lebens der Leitgedanke des Vorranges der Prävention vor der Repression (vgl. \abbrBVerfGE{} 30, 336 [350]). Es ist daher Aufgabe des Staates, in erster Linie sozialpolitische und fürsorgerische Mittel zur Sicherung des werdenden Lebens einzusetzen. Was hier geschehen kann und wie die Hilfsmaßnahmen im einzelnen auszugestalten sind, bleibt weithin dem Gesetzgeber überlassen und entzieht sich im allgemeinen verfassungsgerichtlicher Beurteilung. Dabei wird es hauptsächlich darauf ankommen, die Bereitschaft der werdenden Mutter zu stärken, die Schwangerschaft eigenverantwortlich anzunehmen und die Leibesfrucht zum vollen Leben zu bringen. Bei aller Schutzverpflichtung des Staates darf nicht aus den Augen verloren werden, daß das sich entwickelnde Leben von Natur aus in erster Linie dem Schutz der Mutter anvertraut ist. Den mütterlichen Schutzwillen dort, wo er verlorengegangen ist, wieder zu erwecken und erforderlichenfalls zu stärken, sollte das vornehmste Ziel der staatlichen Bemühungen um Lebensschutz sein. Freilich sind die Einwirkungsmöglichkeiten des Gesetzgebers hier begrenzt. Von ihm eingeleitete Maßnahmen werden häufig nur mittelbar und mit zeitlicher Verzögerung durch eine umfassende Erziehungsarbeit und die dadurch erreichte Veränderung gesellschaftlicher Einstellungen und Anschauungen wirksam.
\end{sourcepara}
\begin{transpara}
\transrn{129}1. 對未出生生命之保護，亦且尤其適用預防優先於壓制之指導思想（參見 BVerfGE 30, 336 [350]）。因此，國家之任務，首先在於運用社會政策及扶助手段以保障正在發展中生命。此處可採取何種措施，以及各項協助措施應如何具體形塑，廣泛地委由立法機關決定，通常不屬憲法法院審查判斷之範圍。重點尤其在於強化準母親以自主責任接受妊娠、並使胎兒得以充分成長為生命之意願。即使國家負有保護義務，仍不得忽略：正在發展中生命依其本性，首先託付於母親之保護。喚醒已經喪失之母性保護意願，並於必要時加以強化，應是國家生命保護努力之最高目標。當然，立法機關在此之影響可能有限。其所採取之措施，往往只能透過全面教育工作及由此達成之社會態度與觀念變化，間接且經時間延遲後發揮效果。
\end{transpara}

\syncleft
\begin{sourcepara}
\sourcern{130}2. Die Frage, inwieweit der Staat von Verfassungs wegen verpflichtet ist, zum Schutz des ungeborenen Lebens auch das Mittel des Strafrechts als der schärfsten ihm zur Verfügung stehenden Waffe einzusetzen, kann nicht von der vereinfachten Fragestellung aus beantwortet werden, ob der Staat bestimmte Handlungen bestrafen muß. Notwendig ist eine Gesamtbetrachtung, die einerseits den Wert des verletzten Rechtsgutes und das Maß der Sozialschädlichkeit der Verletzungshandlung - auch im Vergleich mit anderen unter Strafe gestellten und sozialethisch etwa gleich bewerteten Handlungen - in den Blick nimmt, andererseits die traditionellen rechtlichen Regelungen dieses Lebensbereichs ebenso wie die Entwicklung der Vorstellungen über die Rolle des Strafrechts in der modernen Gesellschaft berücksichtigt und schließlich die praktische Wirksamkeit von Strafdrohungen und die Möglichkeit ihres Ersatzes durch andere rechtliche Sanktionen nicht außer acht läßt.
\end{sourcepara}
\begin{transpara}
\transrn{130}2. 國家在多大程度上有憲法義務使用刑法作為保護未出生生命的最強武器，這一問題不能從國家是否必須對某些行為科以刑罰這一簡單問題來回答。需要有一個整體的看法，一方面考慮到被侵犯的法益的價值和違法行為的社會危害程度——也與其他應受刑罰和同樣受到社會倫理重視的行為進行比較——另一方面考慮到這一生活領域的傳統法律規定以及刑法在現代社會中的作用觀念的發展，最後不忽視刑罰威嚇的實際效力和替代的可能性。他們受到其他法律制裁。
\end{transpara}

\syncleft
\begin{sourcepara}
\sourcern{131}Der Gesetzgeber ist grundsätzlich nicht verpflichtet, die gleichen Maßnahmen strafrechtlicher Art zum Schutze des ungeborenen Lebens zu ergreifen, wie er sie zur Sicherung des geborenen Lebens für zweckdienlich und geboten hält. Wie ein Blick in die Rechtsgeschichte zeigt, war dies bei der Anwendung strafrechtlicher Sanktionen nie der Fall und traf auch für die bis zum Fünften Strafrechtsreformgesetz gegebene Rechtslage nicht zu.
\end{sourcepara}
\begin{transpara}
\transrn{131}原則上，立法機關沒有義務採取其認為適當和必要的刑法措施來保護未出生的生命。正如法律史所示，在刑事制裁之運用上，此種等同對待從未存在過，在刑法改革第五法之前的法律狀況下亦不如此。
\end{transpara}

\syncleft
\begin{sourcepara}
\sourcern{132}a) Aufgabe des Strafrechts war es seit jeher, die elementaren Werte des Gemeinschaftslebens zu schützen. Daß das Leben jedes einzelnen Menschen zu den wichtigsten Rechtsgütern gehört, ist oben dargelegt worden. Der Abbruch einer Schwangerschaft zerstört unwiderruflich entstandenes menschliches Leben. Der Schwangerschaftsabbruch ist eine Tötungshandlung; das wird aufs deutlichste dadurch bezeugt, daß die ihn betreffende Strafdrohung - auch noch im Fünften Strafrechtsreformgesetz - im Abschnitt "Verbrechen und Vergehen wider das Leben" enthalten ist und im bisherigen Strafrecht als "Abtötung der Leibesfrucht" bezeichnet war - die jetzt übliche Bezeichnung als "Schwangerschaftsabbruch" kann diesen Sachverhalt nicht verschleiern. Keine rechtliche Regelung kann daran vorbeikommen, daß mit dieser Handlung gegen die in \abbrArt{} 2 \abbrAbs{} 2 Satz 1 \abbrGG{}\ggfnArtTwoTwoOneDE{} verbürgte grundsätzliche Unantastbarkeit und Unverfügbarkeit des menschlichen Lebens verstoßen wird. Von hier aus gesehen ist der Einsatz des Strafrechts zur Ahndung von "Abtreibungshandlungen" ohne Zweifel legitim; er ist in den meisten Kulturstaaten - unter verschieden gestalteten Voraussetzungen - geltendes Recht und entspricht insbesondere auch der deutschen Rechtstradition. Ebenso ergibt sich hieraus, daß auf eine klare rechtliche Kennzeichnung dieses Vorgangs als "Unrecht" nicht verzichtet werden kann.
\end{sourcepara}
\begin{transpara}
\transrn{132}a) 刑法向來負有保護共同體生活基本價值之任務。如前所述，每一個人的生命均屬最重要法益之一。終止妊娠不可逆地毀滅人類生命。終止妊娠是一項殺害行為；此點最明確地表現在，關於該行為之刑罰威嚇，即使在《第五刑法改革法》中，仍置於「妨害生命之重罪與輕罪」章節內，且在既有刑法中稱為「殺害胎兒」；現今通用之「終止妊娠」名稱，不能遮蔽此一事實。任何法律規制均不能迴避：此一行為違反《基本法》第2條第2項第1句\ggfnArtTwoTwoOneZH{}所保障之人類生命原則上的不可侵犯性與不可支配性。由此觀之，運用刑法制裁「墮胎行為」無疑具有正當性；其在多數文化國家中，在不同前提下均屬現行法，並且尤其符合德國法律傳統。由此亦可得出：不能放棄以法律清楚標示此一過程為「不法」。
\end{transpara}

\syncleft
\begin{sourcepara}
\sourcern{133}b) Indes kann Strafe niemals Selbstzweck sein. Ihr Einsatz unterliegt grundsätzlich der Entscheidung des Gesetzgebers. Er ist nicht gehindert, unter Beachtung der oben angegebenen Gesichtspunkte die grundgesetzlich gebotene rechtliche Mißbilligung des Schwangerschaftsabbruchs auch auf andere Weise zum Ausdruck zu bringen als mit dem Mittel der Strafdrohung. Entscheidend ist, ob die Gesamtheit der dem Schutz des ungeborenen Lebens dienenden Maßnahmen, seien sie bürgerlich-rechtlicher, öffentlichrechtlicher, insbesondere sozialrechtlicher oder strafrechtlicher Natur, einen der Bedeutung des zu sichernden Rechtsgutes entsprechenden tatsächlichen Schutz gewährleistet. Im äußersten Falle, wenn nämlich der von der Verfassung gebotene Schutz auf keine andere Weise zu erreichen ist, kann der Gesetzgeber verpflichtet sein, zum Schutze des sich entwickelnden Lebens das Mittel des Strafrechts einzusetzen. Die Strafnorm stellt gewissermaßen die "ultima ratio" im Instrumentarium des Gesetzgebers dar. Nach dem das ganze öffentliche Recht einschließlich des Verfassungsrechts beherrschenden rechtsstaatlichen Prinzip der Verhältnismäßigkeit darf er von diesem Mittel nur behutsam und zurückhaltend Gebrauch machen. Jedoch muß auch dieses letzte Mittel eingesetzt werden, wenn anders ein effektiver Lebensschutz nicht zu erreichen ist. Dies fordern der Wert und die Bedeutung des zu schützenden Rechtsgutes. Es handelt sich dann nicht um eine "absolute" Pflicht zu strafen, sondern um die aus der Einsicht in die Unzulänglichkeit aller anderen Mittel erwachsende "relative" Verpflichtung zur Benutzung der Strafdrohung.
\end{sourcepara}
\begin{transpara}
\transrn{133}b) 然而，刑罰絕非目的本身。刑罰之運用原則上取決於立法機關之決定。立法機關在遵守上述觀點下，並非不能以刑罰威嚇以外之其他方式，表達《基本法》所要求之對終止妊娠的法律否定評價。決定性者在於：服務於保護未出生生命之全部措施，不論其屬民法、公法，尤其是社會法或刑法性質，是否能保障與所欲保護法益之重要性相稱的實際保護。在極端情形，亦即憲法所要求之保護無法以其他方式達成時，立法機關可能負有義務，運用刑法手段保護正在發展中生命。刑罰規範可謂立法機關工具箱中之「最後手段」。依支配整個公法、包含憲法在內之法治國比例原則，立法機關僅得謹慎且克制地使用此一手段。然而，若無此最後手段即無法達成有效生命保護，仍必須加以使用。這是所保護法益之價值與意義所要求。此時所涉並非「絕對」處罰義務，而是基於認識到所有其他手段均不足而產生之使用刑罰威嚇的「相對」義務。
\end{transpara}

\syncleft
\begin{sourcepara}
\sourcern{134}Demgegenüber greift der Einwand nicht durch, aus einer Freiheit gewährenden Grundrechtsnorm könne niemals eine staatliche Verpflichtung zum Strafen abgeleitet werden. Wenn der Staat durch eine wertentscheidende Grundsatznorm verpflichtet ist, ein besonders wichtiges Rechtsgut auch gegen Angriffe Dritter wirksam zu schützen, so werden oft Maßnahmen unvermeidlich sein, durch welche die Freiheitsbereiche anderer Grundrechtsträger tangiert werden. Insofern ist die Rechtslage beim Einsatz sozialrechtlicher oder zivilrechtlicher Mittel grundsätzlich nicht anders als bei dem Erlaß einer Strafnorm. Unterschiede bestehen allenfalls hinsichtlich der Stärke des erforderlichen Eingriffes. Allerdings muß der Gesetzgeber den hierbei entstehenden Konflikt durch eine Abwägung der beiden einander gegenüberstehenden Grundwerte oder Freiheitsbereiche nach Maßgabe der grundgesetzlichen Wertordnung und unter Beachtung des rechtsstaatlichen Verhältnismäßigkeitsgrundsatzes lösen. Würde man die Pflicht generell verneinen, auch das Mittel des Strafrechts einzusetzen, so würde der zu gewährende Lebensschutz wesentlich eingeschränkt. Dem Wert des von Vernichtung bedrohten Rechtsgutes entspricht der Ernst der für die Vernichtung angedrohten Sanktion, dem elementaren Wert des Menschenlebens die strafrechtliche Ahndung seiner Vernichtung.
\end{sourcepara}
\begin{transpara}
\transrn{134}相較之下，主張不得從保障自由之基本權規範中導出國家處罰義務的反對意見，並不能成立。若國家因一項價值決定性基本規範而負有義務，須有效保護特別重要之法益免受第三人攻擊，則往往不可避免地必須採取會觸及其他基本權主體自由領域之措施。就此而言，運用社會法或民法手段時之法律狀況，原則上與制定刑罰規範並無不同。差異至多存在於所需干預之強度。然而，立法機關必須依《基本法》價值秩序，並遵守法治國比例原則，衡量相互對立之基本價值或自由領域，以解決由此產生之衝突。若一概否認使用刑法手段之義務，所應給予之生命保護將受到重大限制。受毀滅威脅之法益價值，對應於對其毀滅所威嚇制裁之嚴重性；人類生命之基本價值，則對應於以刑法制裁其毀滅。
\end{transpara}

\syncleft
\begin{sourcepara}
\sourcern{135}3. Die Verpflichtung des Staates zum Schutz des werdenden Lebens besteht - wie dargelegt - auch gegenüber der Mutter. Hier läßt jedoch der Einsatz des Strafrechts besondere Probleme entstehen, die sich aus der singulären Lage der schwangeren Frau ergeben. Die einschneidenden Wirkungen einer Schwangerschaft auf den körperlichen und seelischen Zustand der Frau sind unmittelbar einsichtig und bedürfen keiner näheren Darlegung. Sie bedeuten häufig eine erhebliche Änderung der gesamten Lebensführung und eine Einschränkung der persönlichen Entfaltungsmöglichkeiten. Nicht immer und nicht voll wird diese Belastung dadurch ausgeglichen, daß die Frau in ihrer Aufgabe als Mutter neue Erfüllung findet und daß die Schwangere Anspruch auf den Beistand der Gemeinschaft hat (\abbrArt{} 6 \abbrAbs{} 4 \abbrGG{}). Hier können sich im Einzelfall schwere, ja lebensbedrohende Konfliktsituationen ergeben. Das Lebensrecht des Ungeborenen kann zu einer Belastung der Frau führen, die wesentlich über das normalerweise mit einer Schwangerschaft verbundene Maß hinausgeht. Es ergibt sich hier die Frage der Zumutbarkeit, mit anderen Worten die Frage, ob der Staat auch in solchen Fällen mit dem Mittel des Strafrechts die Austragung der Schwangerschaft erzwingen darf. Achtung vor dem ungeborenen Leben und Recht der Frau, nicht über das zumutbare Maß hinaus zur Aufopferung eigener Lebenswerte im Interesse der Respektierung dieses Rechtsgutes gezwungen zu werden, treffen aufeinander. In einer solchen Konfliktslage, die im allgemeinen auch keine eindeutige moralische Beurteilung zuläßt und in der die Entscheidung zum Abbruch einer Schwangerschaft den Rang einer achtenswerten Gewissensentscheidung haben kann, ist der Gesetzgeber zur besonderen Zurückhaltung verpflichtet. Wenn er in diesen Fällen das Verhalten der Schwangeren nicht als strafwürdig ansieht und auf das Mittel der Kriminalstrafe verzichtet, so ist das jedenfalls als Ergebnis einer dem Gesetzgeber obliegenden Abwägung auch verfassungsrechtlich hinzunehmen.
\end{sourcepara}
\begin{transpara}
\transrn{135}3. 如前所述，國家保護正在形成中生命之義務亦存在於對母親之關係中。然而，在此運用刑法會產生特殊問題，這些問題源自懷孕婦女之獨特處境。妊娠對婦女身體與心理狀態之深刻影響，直接可見，無須詳述。其往往意味著整體生活方式之重大改變，以及個人人格發展可能性之限制。此一負擔，並不總是也不完全能由婦女在母親任務中獲得新的實現、以及懷孕婦女享有共同體扶助請求權（《基本法》第6條第4項）而獲得補償。在個案中，可能出現嚴重、甚至危及生命之衝突情境。未出生者生命權可能對婦女形成負擔，遠超通常妊娠所伴隨之程度。此處即產生可期待性問題，換言之，國家是否亦得在此等案件中以刑法手段強制婦女繼續妊娠至分娩。對未出生生命之尊重，與婦女不被迫為尊重此一法益而超過可期待程度犧牲自身生命價值之權利，在此相互衝突。在此種衝突情境中，通常也無法作成明確道德評價，且終止妊娠之決定可能具有值得尊重之良心決定位階，立法機關負有特別克制之義務。若立法機關在此等案件中不認為懷孕婦女之行為具有應罰性，並放棄刑罰手段，則此作為立法機關所負衡量之結果，無論如何在憲法上亦應被接受。
\end{transpara}

\syncleft
\begin{sourcepara}
\sourcern{136}Für die inhaltliche Ausfüllung des Unzumutbarkeitskriteriums müssen jedoch Umstände ausscheiden, die den Pflichtigen nicht schwerwiegend belasten, da sie die Normalsituation darstellen, mit der jeder fertig werden muß. Vielmehr müssen Umstände erheblichen Gewichts gegeben sein, die dem Betroffenen die Erfüllung seiner Pflicht außergewöhnlich erschweren, so daß sie von ihm billigerweise nicht erwartet werden kann. Sie liegen insbesondere dann vor, wenn der Betroffene durch die Pflichterfüllung in schwere innere Konflikte gestürzt wird. Die Lösung solcher Konflikte durch eine Strafdrohung erscheint im allgemeinen nicht als angemessen (vgl. \abbrBVerfGE{} 32, 98 [109] - Gesundbeter), da sie äußeren Zwang einsetzt, wo die Achtung vor der Persönlichkeitssphäre des Menschen volle innere Entscheidungsfreiheit fordert.
\end{sourcepara}
\begin{transpara}
\transrn{136}然而，填充不可期待性標準之內容時，必須排除未對義務人造成重大負擔之情事，因其僅屬每個人均須處理之正常情境。毋寧，必須存在具有相當重量之情事，使當事人履行其義務異常困難，以致依公平觀之不能期待其履行。尤其在當事人因履行義務而陷入嚴重內心衝突時，即屬此種情形。以刑罰威嚇解決此類衝突，通常並非適當（參見 BVerfGE 32, 98 [109] -- Gesundbeter），因其在尊重人格領域要求完全內在決定自由之處，施加外在強制。
\end{transpara}

\syncleft
\begin{sourcepara}
\sourcern{137}Unzumutbar erscheint die Fortsetzung der Schwangerschaft insbesondere, wenn sich erweist, daß der Abbruch erforderlich ist, um von der Schwangeren "eine Gefahr für ihr Leben oder die Gefahr einer schwerwiegenden Beeinträchtigung ihres Gesundheitszustandes abzuwenden" (§ 218b \abbrNr{} 1 \abbrStGB{} in der Fassung des Fünften Strafrechtsreformgesetzes). In diesem Fall steht ihr eigenes "Recht auf Leben und körperliche Unversehrtheit" (\abbrArt{} 2 \abbrAbs{} 2 Satz 1 \abbrGG{}\ggfnArtTwoTwoOneDE{}) auf dem Spiel, dessen Aufopferung für das ungeborene Leben von ihr nicht erwartet werden kann. Darüber hinaus steht es dem Gesetzgeber frei, auch bei anderen außergewöhnlichen Belastungen für die Schwangere, die unter dem Gesichtspunkt der Unzumutbarkeit ähnlich schwer wie die in § 218b \abbrNr{} 1 angeführten wiegen, den Schwangerschaftsabbruch straffrei zu lassen. Hierzu können insbesondere die in dem in der 6. Wahlperiode des Bundestages vorgelegten Entwurf der Bundesregierung enthaltenen und sowohl in der öffentlichen Diskussion wie auch im Verlauf des Gesetzgebungsverfahrens erörterten Fälle der eugenischen (vgl. § 218b \abbrNr{} 2 \abbrStGB{}), der ethischen (kriminologischen) und der sozialen oder Notlageindikation zum Schwangerschaftsabbruch gezählt werden. Bei den Beratungen des Sonderausschusses für die Strafrechtsreform (7. \abbrWp{}, 25. Sitzung, \abbrStenBer{} \abbrS{} 1470 \abbrFF{}) hat der Vertreter der Bundesregierung ausführlich und mit überzeugenden Gründen dargelegt, warum in diesen vier Indikationsfällen die Austragung der Schwangerschaft nicht als zumutbar erscheint. Der entscheidende Gesichtspunkt ist, daß in allen diesen Fällen ein anderes, vom Standpunkt der Verfassung aus ebenfalls schutzwürdiges Interesse sich mit solcher Dringlichkeit geltend macht, daß die staatliche Rechtsordnung nicht verlangen kann, die Schwangere müsse hier dem Recht des Ungeborenen unter allen Umständen den Vorrang einräumen.
\end{sourcepara}
\begin{transpara}
\transrn{137}尤其在下列情形，繼續妊娠顯得不可期待：終止妊娠被證明是「為避免懷孕婦女生命危險或健康狀態受嚴重損害之危險」所必要者（《第五刑法改革法》版本《刑法》第218b條第1款）。在此情形，其自身「生命權及身體完整性權」（《基本法》第2條第2項第1句）受到威脅，不能期待其為未出生生命而犧牲自己。此外，若懷孕婦女承受其他特殊負擔，且從不可期待性觀點而言，其重量與第218b條第1款所列情形同等重大，立法者亦得使終止妊娠不受刑罰處罰。這尤其包括優生情形（參見《刑法》第218b條第2款）、倫理（犯罪學）以及社會或困境情形；此等情形均已納入聯邦政府於聯邦議院第六屆任期提出之草案，並在公開討論與立法程序中受到討論。聯邦政府代表於刑法改革特別委員會審議中（第七屆任期，第25次會議，StenBer. S. 1470 ff.）詳細且有說服力地說明，為何在這四類指標事由中，將妊娠懷至足月顯得不可期待。關鍵在於，在所有此等情形中，另一項從憲法觀點亦值得保護之利益，以如此迫切之方式顯現，以致國家法律秩序不能要求懷孕婦女在任何情況下均必須使未出生者之權利優先。
\end{transpara}

\syncleft
\begin{sourcepara}
\sourcern{138}Auch die Indikation der allgemeinen Notlage (soziale Indikation) kann hierher eingeordnet werden. Denn die allgemeine soziale Lage der Schwangeren und ihrer Familie kann Konflikte von solcher Schwere erzeugen, daß von der Schwangeren über ein bestimmtes Maß hinaus Opfer zugunsten des ungeborenen Lebens mit den Mitteln des Strafrechts nicht erzwungen werden können. Bei der Regelung dieses Indikationsfalles muß der Gesetzgeber den straffreien Tatbestand so umschreiben, daß die Schwere des hier vorauszusetzenden sozialen Konflikts deutlich erkennbar wird und - unter dem Gesichtspunkt der Unzumutbarkeit betrachtet - die Kongruenz dieser Indikation mit den anderen Indikationsfällen gewahrt bleibt. Wenn der Gesetzgeber echte Konfliktsfälle dieser Art aus dem Strafrechtsschutz herausnimmt, verletzt er nicht seine Verpflichtung zum Lebensschutz. Auch in diesen Fällen darf der Staat sich nicht damit begnügen, bloß zu prüfen und gegebenenfalls zu bescheinigen, daß die gesetzlichen Voraussetzungen für einen straffreien Schwangerschaftsabbruch vorliegen. Vielmehr wird auch hier von ihm erwartet, daß er Beratung und Hilfe anbietet mit dem Ziel, die Schwangere an die grundsätzliche Pflicht zur Achtung des Lebensrechts des Ungeborenen zu mahnen, sie zur Fortsetzung der Schwangerschaft zu ermutigen und sie - vor allem in Fällen sozialer Not - durch praktische Hilfsmaßnahmen zu unterstützen.
\end{sourcepara}
\begin{transpara}
\transrn{138}一般困境指標事由（社會指標事由）亦可歸入此處。因為懷孕婦女及其家庭之一般社會處境，可能產生如此嚴重之衝突，以致不能以刑法手段強迫懷孕婦女為未出生生命之利益，承擔超過一定程度之犧牲。立法者在規範此一指標事由案件時，必須以使人清楚看出此處所預設社會衝突之嚴重性之方式，描述不罰構成要件，並且從不可期待性觀點觀察，維持此一指標事由與其他指標事由案件之一致性。若立法者將此類真正衝突案件排除於刑法保護之外，並不違反其生命保護義務。即使在此等情形中，國家亦不得僅滿足於審查並於必要時證明不罰終止妊娠之法定要件存在。毋寧，國家在此亦被期待提供諮詢與協助，其目標在於提醒懷孕婦女注意尊重未出生者生命權之原則義務，鼓勵其繼續妊娠，並特別在社會困境情形中透過實際協助措施支持她。
\end{transpara}

\syncleft
\begin{sourcepara}
\sourcern{139}In allen anderen Fällen bleibt der Schwangerschaftsabbruch strafwürdiges Unrecht; denn hier steht die Vernichtung eines Rechtsgutes von höchstem Rang im freien - nicht durch eine Notlage motivierten - Belieben eines anderen. Wollte der Gesetzgeber auch hier auf die strafrechtliche Ahndung verzichten, so wäre das nur unter der Voraussetzung mit dem Schutzgebot des \abbrArt{} 2 \abbrAbs{} 2 Satz 1 \abbrGG{}\ggfnArtTwoTwoOneDE{} vereinbar, daß ihm eine andere gleich wirksame rechtliche Sanktion zu Gebote stände, die den Unrechtscharakter der Handlung (die Mißbilligung durch die Rechtsordnung) deutlich erkennen läßt und Schwangerschaftsabbrüche ebenso wirksam verhindert wie eine Strafvorschrift.
\end{sourcepara}
\begin{transpara}
\transrn{139}在所有其他情況下，終止妊娠仍屬具有應罰性之不法；因為在此，最高位階法益之毀滅，取決於他人自由且非由困境所促成之任意。若立法機關在此亦欲放棄刑事制裁，則唯有在另有同等有效之法律制裁可供其使用，且該制裁能清楚顯示行為之不法性質（法律秩序之否定評價），並能如刑罰規範般有效防止終止妊娠時，始與《基本法》第2條第2項第1句\ggfnArtTwoTwoOneZH{}之保護命令相容。
\end{transpara}

\bisubsection{D.}{D.}
\syncleft
\begin{sourcepara}
\sourcern{140}Prüft man nach diesen Maßstäben die angegriffene Fristenregelung des Fünften Strafrechtsreformgesetzes, so ergibt sich, daß das Gesetz der Verpflichtung aus \abbrArt{} 2 \abbrAbs{} 2 Satz 1 in Verbindung mit \abbrArt{} 1 \abbrAbs{} 1 \abbrGG{}\ggfnLifeDignityDE{}, das werdende Leben wirksam zu schützen, nicht in dem gebotenen Umfang gerecht geworden ist.
\end{sourcepara}
\begin{transpara}
\transrn{140}如果依照這些標準來檢視刑法改革第五法中受到質疑的期限規定，就會發現法律並沒有充分公正地履行《基本法》第二條第二項第一句\ggfnArtTwoTwoOneZH{}並聯第一條第一項\ggfnLifeDignityZH{}所規定的有效保護正在發展中的生命的義務。
\end{transpara}

\bisubsubsection{I.}{I.}
\syncleft
\begin{sourcepara}
\sourcern{141}Das Verfassungsgebot, das sich entwickelnde Leben zu schützen, richtet sich zwar in erster Linie an den Gesetzgeber. Dem Bundesverfassungsgericht obliegt jedoch die Aufgabe, in Ausübung der ihm vom Grundgesetz zugewiesenen Funktion festzustellen, ob der Gesetzgeber dieses Gebot erfüllt hat. Zwar muß das Gericht den Spielraum des Gesetzgebers sorgfältig beachten, der diesem bei der Beurteilung der seiner Normierung zugrundeliegenden tatsächlichen Verhältnisse, der etwa erforderlichen Prognose und der Wahl der Mittel zukommt. Das Gericht darf sich nicht an die Stelle des Gesetzgebers setzen; es ist jedoch seine Aufgabe, zu überprüfen, ob der Gesetzgeber im Rahmen der ihm zur Verfügung stehenden Möglichkeiten das Erforderliche getan hat, um Gefahren von dem zu schützenden Rechtsgut abzuwenden. Das gilt grundsätzlich auch für die Frage, ob der Gesetzgeber zum Einsatz seines schärfsten Mittels, des Strafrechts, verpflichtet ist, wobei sich die Prüfung allerdings nicht etwa auf die einzelnen Modalitäten des Strafens erstrecken kann.
\end{sourcepara}
\begin{transpara}
\transrn{141}保護發展中生命的憲法要求主要是針對立法機關的。然而，聯邦憲法法院在履行《基本法》賦予的職能時，負責判定立法機關是否履行了這項要求。誠然，法院在評估其標準化所依據的實際情況、任何必要的預測和手段的選擇時，必須仔細考慮立法機關的靈活性。法院不得代替立法機關；然而，其任務是檢查立法機關是否在其可能的範圍內採取了必要的措施，以避免對受保護的法益造成危險。這原則上也適用於立法機關是否有義務使用其最嚴厲手段（刑法）的問題，儘管審查不能延伸到個別的刑罰方式。
\end{transpara}

\bisubsubsection{II.}{II.}
\syncleft
\begin{sourcepara}
\sourcern{142}Es ist allgemein anerkannt, daß der bisherige § 218 \abbrStGB{}, gerade weil er für nahezu alle Fälle des Schwangerschaftsabbruchs undifferenziert Strafe androhte, das sich entwickelnde Leben im Ergebnis nur unzureichend geschützt hat. Denn die Einsicht, daß es Fälle gibt, in denen die strafrechtliche Sanktion unangemessen ist, hat schließlich dazu geführt, daß auch wirklich strafwürdige Fälle nicht mehr mit der notwendigen Strenge verfolgt werden. Dazu kommt die bei diesem Delikt der Natur der Sache nach häufig schwierige Aufklärung des Sachverhalts. Gewiß gehen die Zahlen über die Dunkelziffer bei Schwangerschaftsabbrüchen weit auseinander und es mag auch kaum möglich sein, durch empirische Untersuchungen hierüber zuverlässige Angaben zu ermitteln. Jedenfalls war die Zahl der illegalen Schwangerschaftsabbrüche in der Bundesrepublik hoch. Das Bestehen einer generellen Strafnorm mag auch dazu beigetragen haben, daß der Staat es unterlassen hat, andere ausreichende Maßnahmen zum Schutze des werdenden Lebens zu ergreifen.
\end{sourcepara}
\begin{transpara}
\transrn{142}人們普遍認為，之前的《刑法》第218條，正是因為它威脅對幾乎所有終止妊娠案件進行一視同仁的刑罰，最終只能不充分地保護發育中的生命。有鑑於刑事制裁不當的情況，最終導致真正應受刑罰的案件不再受到必要的嚴厲起訴。此外，由於該犯罪的性質，案件事實往往難以查明。有關未報告終止妊娠數量的數字肯定差異很大，並且很難透過實證研究獲得有關這方面的可靠資訊。無論如何，聯邦共和國的非法終止妊娠數量很高。一般刑法規範的存在也可能導致國家未能採取其他適當措施來保護正在發展中的生命。
\end{transpara}

\syncleft
\begin{sourcepara}
\sourcern{143}Der Gesetzgeber ist bei der endgültigen Fassung des Fünften Strafrechtsreformgesetzes von dem Leitgedanken des Vorranges der präventiven Maßnahmen vor den repressiven Sanktionen ausgegangen (vgl. hierzu den vom Bundestag im Zusammenhang mit der Verabschiedung des Fünften Strafrechtsreformgesetzes angenommenen Entschließungsantrag der Fraktionen der SPD, FDP - \abbrBTDrucks{} 7/2042). Dem Gesetz liegt die Vorstellung zugrunde, das werdende Leben werde besser durch individuelle Beratung der Schwangeren geschützt als durch eine Strafandrohung, welche die Abtreibungswillige einer möglichen Beeinflussung entziehe, kriminalpolitisch verfehlt sei und sich ohnedies als wirkungslos erwiesen habe. Hieraus hat der Gesetzgeber die Folgerung gezogen, die Strafandrohung für die ersten zwölf Wochen der Schwangerschaft unter bestimmten Voraussetzungen überhaupt aufzugeben und statt dessen die präventive Beratung und Unterrichtung (§§ 218a und 218c) einzuführen.
\end{sourcepara}
\begin{transpara}
\transrn{143}在《第五刑法改革法案》的最終版本中，立法機關以預防性措施優先於鎮壓性制裁的指導原則為出發點（參見聯邦議院就第五刑法改革法案的通過而通過的社民黨和自民黨議會團體決議動議 - BTDrucks.7/2042）。該法律的基本理念是，透過對懷孕婦女的個人諮詢比透過刑罰威嚇能更好地保護正在發展中的生命，刑罰威嚇剝奪了那些希望終止妊娠的人可能產生的影響力，從刑事政策角度來看是錯誤的，而且已被證明是無效的。立法機關由此得出結論，在某些情況下應放棄對懷孕前十二週的刑罰威嚇，而應引入預防性諮詢和資訊（§§ 218a 和 218c）。
\end{transpara}

\syncleft
\begin{sourcepara}
\sourcern{144}Es ist verfassungsrechtlich unbedenklich und zu billigen, wenn der Gesetzgeber seine Pflichte zu einem besseren Schutz ungeborenen Lebens durch präventive Maßnahmen einschließlich einer die Eigenverantwortung der Frau stärkenden Beratung zu erfüllen versucht. Jedoch begegnet die getroffene Regelung in mehrfacher Hinsicht durchgreifenden verfassungsrechtlichen Bedenken.
\end{sourcepara}
\begin{transpara}
\transrn{144}若立法者試圖透過預防性措施，包括透過強化婦女自主負責的諮詢，來履行其更妥善保護未出生生命的義務，則此在憲法上並無疑義，且應予肯認。然而，所作成的規定在數個方面遭遇具有決定性的憲法疑慮。
\end{transpara}

\syncleft
\begin{sourcepara}
\sourcern{145}1. Die von der Verfassung geforderte rechtliche Mißbilligung des Schwangerschaftsabbruchs muß auch in der Rechtsordnung unterhalb der Verfassung deutlich in Erscheinung treten. Davon können - wie dargelegt - nur die Fälle ausgenommen werden, in denen die Fortsetzung der Schwangerschaft der Frau auch unter Berücksichtigung der in \abbrArt{} 2 \abbrAbs{} 2 Satz 1 \abbrGG{}\ggfnArtTwoTwoOneDE{} getroffenen Wertentscheidung nicht zumutbar ist. Dieses generelle Unwerturteil kommt in den Bestimmungen des Fünften Strafrechtsreformgesetzes über den Schwangerschaftsabbruch in den ersten zwölf Wochen nicht zum Ausdruck; denn das Gesetz läßt es im unklaren, ob der nicht indizierte Schwangerschaftsabbruch nach Aufhebung der Strafandrohung durch § 218a \abbrStGB{} noch Recht oder Unrecht ist. Das gilt ungeachtet der Tatsache, daß sich § 218a \abbrStGB{} gesetzestechnisch als eine Ausnahmeregelung im Verhältnis zu der allgemeinen Strafvorschrift des § 218 \abbrStGB{} darstellt, und auch unabhängig davon, welche Auffassung man in der Frage vertritt, ob die Vorschrift den § 218 \abbrStGB{} tatbestandlich einengt oder ob sie einen Rechtfertigungsgrund schafft oder endlich nur einen Schuld- oder Strafausschließungsgrund zum Inhalt hat. Beim unbefangenen Leser der Bestimmung muß der Eindruck entstehen, daß § 218a das rechtliche Unwerturteil durch die generelle Aufhebung der Strafbarkeit - ohne Ansehung der Gründe - vollständig zurücknimmt und den Schwangerschaftsabbruch unter den dort genannten Voraussetzungen rechtlich erlaubt. Der Tatbestand des § 218 \abbrStGB{} tritt demgegenüber um so mehr in den Hintergrund, als nach aller Erfahrung in den ersten zwölf Wochen die weitaus meisten Schwangerschaftsabbrüche - nach Angaben des Regierungsvertreters (\abbrAAO{}, \abbrS{} 1472) über 9/10 - vorgenommen werden. Es entsteht das Bild einer fast völligen - strafrechtlichen - Freigabe des Schwangerschaftsabbruchs (so auch Roxin in: J. Baumann [Hrsg.], Das Abtreibungsverbot des § 218, \abbrS{} 185). Auch in keiner anderen Bestimmung des Fünften Strafrechtsreformgesetzes kommt zum Ausdruck, daß der nicht indizierte Schwangerschaftsabbruch in den ersten zwölf Wochen weiterhin rechtlich mißbilligt wird. Insbesondere besagt \abbrArt{} 2 des Gesetzes, wonach grundsätzlich niemand verpflichtet ist, an einem Schwangerschaftsabbruch mitzuwirken, nichts über die Legalität oder Illegalität einer solchen Maßnahme; diese Vorschrift will in erster Linie der Gewissensfreiheit des Einzelnen Rechnung tragen und die Freiheit der ethischen Überzeugung dessen schützen, der sich vor die Frage gestellt sieht, ob er an einem nach § 218a \abbrStGB{} straffreien Schwangerschaftsabbruch durch aktives Handeln mitwirken kann und soll.
\end{sourcepara}
\begin{transpara}
\transrn{145}1. 憲法所要求、法律上對終止妊娠之否定評價，也必須在憲法以下之法律秩序中清楚顯現。如前所述，只有在即使考量《基本法》第2條第2項第1句\ggfnArtTwoTwoOneZH{}所作價值決定，繼續妊娠對婦女仍不可期待之情形，始得例外。此一一般性負面評價（\textit{allgemeines Unwerturteil}），並未在《第五刑法改革法》關於最初十二週內終止妊娠之規定中表達出來；因為法律未說明，在《刑法》第218a條取消刑罰威嚇後，未具指標事由之終止妊娠究竟仍屬合法或不法。即使第218a條在法律技術上呈現為相對於《刑法》第218條一般刑罰規定之例外規定，亦不影響此點；無論人們認為該規定是在構成要件上限縮第218條、創設正當化事由，抑或最終僅包含阻卻罪責或阻卻刑罰事由，亦同。對不帶偏見之讀者而言，該規定必然造成一種印象：第218a條透過概括取消可罰性，且不問理由，已完全收回法律上之負面評價，並在法律上允許符合該條要件之終止妊娠。相較之下，《刑法》第218條構成要件益發退居背景；因為依一切經驗，絕大多數終止妊娠均在最初十二週內實施，依政府代表說法（a.a.O., S. 1472）甚至超過十分之九。由此形成幾乎完全在刑法上開放終止妊娠之圖像（Roxin 亦同見於：J. Baumann [Hrsg.], Das Abtreibungsverbot des § 218, S. 185）。《第五刑法改革法》其他規定亦未表達：最初十二週內未具指標事由之終止妊娠，在法律上仍受否定評價。尤其是，該法第2條規定原則上無人有義務參與終止妊娠，並未說明此類措施之合法性或違法性；該規定首要係顧及個人良心自由，並保護面臨下列問題者之倫理信念自由：其是否能夠、且是否應當以積極作為參與依《刑法》第218a條不罰之終止妊娠。
\end{transpara}

\syncleft
\begin{sourcepara}
\sourcern{146}Ein Blick in die im Strafrechtsreform-Ergänzungsgesetz für das Gebiet des Sozialrechts vorgesehenen Regelungen zwingt darüber hinaus zu dem Schluß, daß es sich bei dem Schwangerschaftsabbruch in den ersten zwölf Wochen um einen Vorgang handelt, dem rechtlich nichts Verwerfliches anhaftet und der deshalb auch sozialrechtlich gefördert und erleichtert werden darf. Denn gesetzliche Rechtsansprüche auf soziale Leistungen setzen voraus, daß der Tatbestand, bei dessen Erfüllung sie gewährt werden, keine rechtlich verbotene (mißbilligte) Handlung darstellt. Die vorgesehene Gesamtregelung kann deshalb nur so gedeutet werden, daß der vom Arzt vorgenommene Schwangerschaftsabbruch während der ersten zwölf Wochen nicht rechtswidrig, also (vom Recht) erlaubt sein soll.
\end{sourcepara}
\begin{transpara}
\transrn{146}觀察《刑法改革補充法》中社會法領域之規定，亦可得出一項結論：最初十二週內之終止妊娠，是一項在法律上不受否定評價之過程，因而亦得在社會法上受到促進及便利。社會給付之法律請求權，以取得該等給付所針對之行為並非法律禁止（受否定評價）之行為為前提。因此，所設想之整體規制只能被理解為：醫師於最初十二週內實施之終止妊娠，不應屬違法，而係法律所容許。
\end{transpara}

\syncleft
\begin{sourcepara}
\sourcern{147}Diese Auffassung vertrat auch die Bundesregierung bei der in der 6. Wahlperiode des Deutschen Bundestages eingebrachten Gesetzesvorlage; dort heißt es in der Begründung zu \abbrArt{} 1 (\abbrBTDrucks{} VI/3434 \abbrS{} 9):
\end{sourcepara}
\begin{transpara}
\transrn{147}聯邦政府在德國聯邦議院第六屆選舉期間提出的法案中也表達了這一觀點；它在藝術的理由中說。 1（BTDrucks.VI/3434 第 9 頁）：
\end{transpara}

\syncleft
\begin{sourcepara}
\sourcern{148}\sourcequoteinner{"Während sich der Gesetzgeber in anderen Bereichen darauf verlassen darf, daß eine Aufhebung strafrechtlicher Verbote nicht als rechtliche Billigung des bisher strafbaren Verhaltens verstanden werden wird, sind bei der Neuregelung des Schwangerschaftsabbruchs besondere Gesichtspunkte zu beachten: Die Fristenlösung könnte die von ihr erwartete gesundheitspolitische Aufgabe nur erfüllen, wenn jeder Schwangerschaftsabbruch in den ersten drei Monaten als rechtlich gebilligt erscheint. Der Eingriff müßte im Rahmen der allgemeinen ärztlichen Versorgung vorgenommen werden. Der ärztliche Behandlungsvertrag müßte wirksam sein. Nicht zuletzt wegen der Unanwendbarkeit der §§ 134, 138 BGB könnten diese und andere Umstände nur in dem Sinne interpretiert werden, daß die Rechtsordnung den Eingriff vor Ablauf der Dreimonatsfrist in jedem Fall als einen normalen sozialen Vorgang anerkennt".}
\end{sourcepara}
\begin{transpara}
\transrn{148}\transquoteinner{「雖然在其他領域，立法機關可以信賴：刑事禁令之廢止不會被理解為對迄今可罰行為之法律認可；但在重新規範終止妊娠時，必須注意特殊觀點：期限規定只有在妊娠前三個月內每一次終止妊娠均呈現為法律上受到認可時，才可能完成其所期待之衛生政策任務。該手術必須在一般醫療照護之框架內實施。醫療契約必須有效。尤其因《民法》第134條、第138條不適用，此等及其他情況只能被解釋為：法秩序在三個月期限屆滿前，於任何情形均承認該手術為正常的社會過程。」}
\end{transpara}

\syncleft
\begin{sourcepara}
\sourcern{149}Ähnlich äußerte sich der Regierungsvertreter vor dem Sonderausschuß für die Strafrechtsreform (7. \abbrWp{}, 25. Sitzung, \abbrStenBer{} \abbrS{} 1473):
\end{sourcepara}
\begin{transpara}
\transrn{149}政府代表在刑法改革特別委員會上發表了類似聲明（第 7 屆會議，第 25 屆會議，StenBer. 第 1473 頁）：
\end{transpara}

\syncleft
\begin{sourcepara}
\sourcern{150}\sourcequoteinner{"So viel ist wichtig festzuhalten: Der ärztliche Schwangerschaftsabbruch im ersten Trimester der Schwangerschaft ist im Rahmen der Fristenregelung nicht rechtswidrig; er ist erlaubt. Nur so läßt sich seine Integration in das System des Strafrechts - mit Straffreiheit auch der Teilnehmer, Ausschluß der Nothilfe - rechtfertigen, und nur so lassen sich die zivilrechtlichen Implikationen - Gültigkeit des Behandlungsvertrags trotz § 134 BGB -, die gesundheitsrechtliche Förderung des Vorganges und vor allen Dingen der mit dem Strafrechtsreform-Ergänzungsgesetz geplante Einbau in die Sozialversicherung begründen."}
\end{sourcepara}
\begin{transpara}
\transrn{150}\transquoteinner{「重要的是必須確認：在期限規定之框架內，妊娠第一三分期之醫師實施終止妊娠並非違法；它是被允許的。唯有如此，才能正當化其納入刑法體系之中：參與者亦不受罰，且排除緊急救助；也唯有如此，才能說明其民法上的意涵，例如治療契約儘管有《民法》第134條仍屬有效，以及衛生法上對此過程之促進，尤其是計畫透過《刑法改革補充法》將其納入社會保險。」}
\end{transpara}

\syncleft
\begin{sourcepara}
\sourcern{151}2. Eine - formale - gesetzliche Mißbilligung des Schwangerschaftsabbruchs würde im übrigen nicht ausreichen; denn darüber wird sich die zum Schwangerschaftsabbruch entschlossene Frau hinwegsetzen. Der Gesetzgeber des Fünften Strafrechtsreformgesetzes hat aus der Einsicht heraus, daß auch positive Maßnahmen zum Schutze des werdenden Lebens erforderlich sind, hinsichtlich des mit Einwilligung der Schwangeren von einem Arzt vorgenommenen Schwangerschaftsabbruchs die Strafnorm durch ein Beratungssystem nach § 218c \abbrStGB{} ersetzt. Durch die völlige Aufhebung der Strafbarkeit ist jedoch eine Schutzlücke entstanden, welche die Sicherung des sich entwickelnden Lebens in einer nicht geringen Anzahl von Fällen gänzlich beseitigt, indem es dieses Leben der völlig freien Verfügungsgewalt der Frau ausliefert. Es gibt viele Frauen, die von vornherein zum Schwangerschaftsabbruch entschlossen und einer Beratung, wie sie § 218c \abbrAbs{} 1 vorsieht, nicht zugänglich sind, ohne daß ein nach der Wertordnung der Verfassung achtenswerter Grund für den Abbruch vorliegt. Diese Frauen befinden sich weder in einer materiellen Notlage noch in einer schwerwiegenden seelischen Konfliktsituation. Sie lehnen die Schwangerschaft ab, weil sie nicht willens sind, den damit verbundenen Verzicht und die natürlichen mütterlichen Pflichten zu übernehmen. Sie mögen ernstliche Gründe für ihre Haltung gegenüber dem werdenden Leben haben; es sind aber keine Gründe, die gegenüber dem Gebot des Schutzes menschlichen Lebens Bestand haben können. Die Schwangerschaft ist diesen Frauen nach den oben wiedergegebenen Grundsätzen zumutbar. Das Verhalten auch dieser durch keine verfassungsrechtlich erheblichen Beweggründe zum Schwangerschaftsabbruch rechtlich legitimierten Gruppe von Frauen ist nach § 218a des Fünften Strafrechtsreformgesetzes voll gedeckt. Das sich entwickelnde Leben ist ihrer willkürlichen Entschließung schutzlos preisgegeben.
\end{sourcepara}
\begin{transpara}
\transrn{151}2. 僅僅在法律上正式反對終止妊娠是不夠的；因為決定終止妊娠的婦女會忽略這一點。有鑑於保護預期壽命也需要採取積極措施，第五刑法改革法的立法機關根據《刑法》第218c條關於經懷孕婦女同意由醫師進行終止妊娠的規定，以諮詢制度取代了刑法。然而，刑事責任的完全廢除造成了保護方面的差距，在相當多的情況下，這種差距完全消除了生命發育的保障，讓婦女完全自行決定生命的發展。許多婦女從一開始就決定終止妊娠，但如果沒有符合憲法價值秩序的正當終止妊娠理由，就無法獲得第 218c 條第1項規定的諮詢。這些婦女既沒有物質上的困難，也沒有嚴重的心理衝突。她們拒絕懷孕，因為她們不願意接受隨之而來的犧牲和自然的母親責任。他們對正在發展中的生命的態度可能有嚴肅的原因；但它們並不是能夠應對保護人類生命的迫切需求的理由。根據上述原則，懷孕對這些女性來說是合理的。這群婦女的行為，在憲法上沒有重大終止妊娠的動機，完全受到刑法改革第五法第218a條的管轄。正在發育的生命對於它們的任意決定毫無抵抗力。
\end{transpara}

\syncleft
\begin{sourcepara}
\sourcern{152}Hiergegen wird eingewandt, unbeeinflußbare Frauen verständen es erfahrungsgemäß zumeist, sich der Strafe zu entziehen, so daß die Strafandrohung ohnehin weitgehend leerlaufe. Außerdem stehe der Gesetzgeber vor dem Dilemma, daß sich präventive Beratung und repressive Strafdrohung in ihrer lebensschützenden Wirkung notwendig teilweise ausschlössen: Die Strafandrohung der Indikationenlösung vermöge zwar durch ihre Abschreckungswirkung in einem nicht genau feststellbaren Umfang unmotivierte Schwangerschaftsabbrüche zu verhindern. Zugleich verhindere aber die Strafandrohung, daß durch Beratung beeinflußbarer Frauen in anderen Fällen Leben gerettet würde; denn gerade Frauen, bei denen die Voraussetzungen einer Indikation fehlten, und darüber hinaus auch solche, die dem Ausgang eines Indikationsfeststellungsverfahrens nicht trauten, würden ihre Schwangerschaft angesichts der Strafandrohung vorsorglich geheimhalten und sich daher weithin einer helfenden Beeinflussung durch Umgebung und Beratungsstellen entziehen. Von dieser Einsicht her könne es einen lückenlosen Schutz ungeborenen Lebens nicht geben. Der Gesetzgeber habe keine andere Wahl als Leben gegen Leben abzuwägen, nämlich das Leben, das durch eine bestimmte Regelung der Abtreibungsfrage voraussichtlich gerettet werden könne, gegen das Leben, das durch ebendieselbe Regelung voraussichtlich geopfert werde; denn auch die Strafdrohung schütze nicht nur, sondern vernichte zugleich ungeborenes Leben. Da keine Lösung dem Schutz individuellen Lebens eindeutig besser diene, habe der Gesetzgeber mit der Fristenregelung die ihm verfassungsrechtlich gezogenen Grenzen nicht überschritten.
\end{sourcepara}
\begin{transpara}
\transrn{152}對此有意見主張，依經驗，不可受影響之婦女通常懂得逃避刑罰，因此刑罰威嚇無論如何大多落空。此外，立法者面臨一項兩難：預防性諮詢與壓制性刑罰威嚇，在生命保護效果上必然部分相互排斥。指標方案之刑罰威嚇，固然能藉其嚇阻效果，在無法精確確認之範圍內防止無動機之終止妊娠；但同時，刑罰威嚇也會阻止在其他案件中，藉由諮詢可受影響之婦女而挽救生命。因為正是欠缺指標事由要件之婦女，以及不信任指標事由確認程序結果之婦女，會鑑於刑罰威嚇而預先隱瞞妊娠，因而在很大程度上避開其環境及諮詢機構之協助性影響。由此觀點出發，不可能存在毫無漏洞之未出生生命保護。立法者別無選擇，只能以生命衡量生命，也就是將某一終止妊娠規制預期可拯救之生命，與同一規制預期會犧牲之生命相衡量；因為刑罰威嚇不僅保護，也同時毀滅未出生生命。既然沒有任何方案明確較有助於保護個別生命，立法者採取期限規定並未逾越憲法為其劃定之界限。
\end{transpara}

\syncleft
\begin{sourcepara}
\sourcern{153}a) Diese Auffassung wird zunächst dem Wesen und der Funktion des Strafrechts nicht gerecht. Die Strafnorm richtet sich grundsätzlich an alle Rechtsunterworfenen und verpflichtet sie in gleicher Weise. Zwar gelingt es den Strafverfolgungsbehörden praktisch nie, alle Täter, die gegen das Strafgesetz verstoßen, einer Bestrafung zuzuführen. Die Dunkelziffern sind bei den verschiedenen Strafdelikten verschieden hoch. Unbestritten sind sie bei Abtreibungstaten besonders erheblich. Indessen darf darüber die generalpräventive Funktion des Strafrechts nicht vergessen werden. Sieht man die Aufgabe des Strafrechts in dem Schutz besonders wichtiger Rechtsgüter und elementarer Werte der Gemeinschaft, so kommt gerade dieser Funktion eine hohe Bedeutung zu. Ebenso wichtig wie die sichtbare Reaktion im Einzelfall ist die Fernwirkung einer Strafnorm, die in ihrem prinzipiellen normativen Inhalt ("die Abtreibung ist strafbar") nunmehr seit sehr langer Zeit besteht. Schon die bloße Existenz einer solchen Strafandrohung hat Einfluß auf die Wertvorstellungen und die Verhaltensweisen der Bevölkerung (vgl. Bericht des Sonderausschusses für die Strafrechtsreform, \abbrBTDrucks{} 7/1981 (neu) \abbrS{} 10). Das Wissen um die Rechtsfolgen im Falle ihrer Übertretung bildet eine Schwelle, vor deren Überschreitung viele zurückschrecken. Diese Wirkung wird ins Gegenteil verkehrt, wenn durch eine generelle Aufhebung der Strafbarkeit auch zweifellos strafwürdiges Verhalten für rechtlich einwandfrei erklärt wird. Dies muß die in der Bevölkerung herrschenden Auffassungen von "Recht" und "Unrecht" verwirren. Die rein theoretische Verlautbarung, der Schwangerschaftsabbruch werde "toleriert", aber nicht "gebilligt", muß wirkungslos bleiben, solange keine rechtliche Sanktion erkennbar ist, die die gerechtfertigten Fälle des Schwangerschaftsabbruchs von den verwerflichen klar scheidet. Entfällt die Drohung mit Strafe ganz allgemein, so wird notwendig im Bewußtsein der Staatsbürger der Eindruck entstehen, in allen Fällen sei der Schwangerschaftsabbruch rechtlich erlaubt und darum auch sozialethisch nicht mehr zu mißbilligen. Der "gefährliche Schluß von der rechtlichen Sanktionslosigkeit auf das moralische Erlaubtsein" (Engisch, Auf der Suche nach der Gerechtigkeit, 1971, \abbrS{} 104) liegt zu nahe, als daß er nicht von einer großen Anzahl Rechtsunterworfener gezogen würde.
\end{sourcepara}
\begin{transpara}
\transrn{153}a) 此一見解首先未能正確掌握刑法之本質與功能。刑罰規範原則上針對所有受法律拘束者，並以相同方式課予其義務。固然，刑事追訴機關實務上幾乎不可能將所有違反刑法之行為人均繩之以法。各類刑事犯罪之黑數高低不同；墮胎犯罪之黑數特別龐大，則無爭議。然而，不得因此忘卻刑法之一般預防功能。若將刑法任務理解為保護共同體特別重要之法益與基本價值，則此一功能正具有重大意義。與個案中可見反應同樣重要者，是刑罰規範之遠期效果；該規範之原則性規範內容（「墮胎應受刑罰」）已存在甚久。單是此種刑罰威嚇之存在，即會影響人民之價值觀與行為方式（參見刑法改革特別委員會報告，BTDrucks. 7/1981 [neu], S. 10）。知悉違反時之法律後果，形成許多人不敢逾越之門檻。若透過一般性廢除可罰性，使毫無疑問具有應罰性之行為被宣告為法律上無瑕疵，此一效果即轉為相反。這必然混淆人民對「法」與「不法」之通行觀念。只要看不出有法律制裁能清楚區分正當化之終止妊娠案件與可非難案件，單純理論上宣稱終止妊娠係受「容忍」而非受「肯認」，必然不生效果。若刑罰威嚇被全面取消，公民意識中勢必產生印象：所有情況下終止妊娠均為法律所允許，因而在社會倫理上也不再受到否定評價。Engisch 所謂「由法律無制裁性推論至道德允許性之錯誤推論」（Auf der Suche nach der Gerechtigkeit, 1971, S. 104）過於接近，以致大量受法律拘束者難免作此推論。
\end{transpara}

\syncleft
\begin{sourcepara}
\sourcern{154}Dem entspricht auch die Auffassung der Bundesregierung in der Begründung zu dem in der 6. Wahlperiode des Deutschen Bundestages eingebrachten Gesetzentwurf (\abbrBTDrucks{} VI/3434 \abbrS{} 9):
\end{sourcepara}
\begin{transpara}
\transrn{154}聯邦政府在其於德國聯邦議院第6屆選期提出之法律草案理由中，亦持相同見解（BTDrucks. VI/3434 S. 9）：
\end{transpara}

\syncleft
\begin{sourcepara}
\sourcern{155}\sourcequoteinner{"Die Fristenlösung würde dazu führen, daß das allgemeine Bewußtsein von der Schutzwürdigkeit des ungeborenen Lebens während der ersten drei Schwangerschaftsmonate schwindet. Sie würde der Ansicht Vorschub leisten, daß der Schwangerschaftsabbruch, jedenfalls im Frühstadium der Schwangerschaft, ebenso dem freien Verfügungsrecht der Schwangeren unterliegt wie die Verhütung der Schwangerschaft. Eine solche Auffassung ist mit der Wertordnung der Verfassung unvereinbar".}
\end{sourcepara}
\begin{transpara}
\transrn{155}\transquoteinner{「期限規定將導致人們普遍喪失對在妊娠前三個月保護未出生生命之必要性的認識。這將助長這樣一種觀點，即終止妊娠，至少在妊娠早期階段，如同避孕一樣，均委由孕婦自由處置。這種觀點與憲法的價值體系不相容。」}
\end{transpara}

\syncleft
\begin{sourcepara}
\sourcern{156}b) Die pauschale Abwägung von Leben gegen Leben, die zur Freigabe der Vernichtung der vermeintlich geringeren Zahl im Interesse der Erhaltung der angeblich größeren Zahl führt, ist nicht vereinbar mit der Verpflichtung zum individuellen Schutz jedes einzelnen konkreten Lebens.
\end{sourcepara}
\begin{transpara}
\transrn{156}b) 對生命與生命的普遍權衡，導致同意為了保護所謂的較大數量而毀滅所謂的較小數量，這與單獨保護每個具體生命的義務是不相容的。
\end{transpara}

\syncleft
\begin{sourcepara}
\sourcern{157}In der Rechtsprechung des Bundesverfassungsgerichts ist der Grundsatz entwickelt worden, daß die Verfassungswidrigkeit einer gesetzlichen Vorschrift, die ihrer Struktur und tatsächlichen Wirkung nach einen bestimmten Personenkreis benachteiligt, nicht mit dem Hinweis darauf widerlegt werden kann, daß diese Vorschrift oder andere Bestimmungen des Gesetzes einen anderen Kreis von Personen begünstigen. Noch weniger genügt hierfür die Betonung der allgemein rechtsschutzfreundlichen Tendenz des Gesetzes im ganzen. Dieses Prinzip (vgl. \abbrBVerfGE{} 12, 151 [168]; 15, 328 [333]; 18, 97 [108]; 32, 260 [269]) muß in besonderem Maße für das höchstpersönliche Rechtsgut "Leben" gelten. Der Schutz des einzelnen Lebens darf nicht deswegen aufgegeben werden, weil das an sich achtenswerte Ziel verfolgt wird, andere Leben zu retten. Jedes menschliche Leben - auch das erst sich entwickelnde Leben - ist als solches gleich wertvoll und kann deshalb keiner irgendwie gearteten unterschiedlichen Bewertung oder gar zahlenmäßigen Abwägung unterworfen werden.
\end{sourcepara}
\begin{transpara}
\transrn{157}聯邦憲法法院的判例法制定了這樣的原則：如果法律條項的結構和實際效果對某一群體不利，則不能透過指出該條項或法律的其他條項有利於另一群體來反駁其違憲性。僅僅強調整個法律普遍的法律保護傾向還不夠。這項原則（參見 BVerfGE 12, 151 [168]；15, 328 [333]；18, 97 [108]；32, 260 [269]）必須特別適用於「生命」的高度個人法益。絕不能放棄對個人生命的保護，因為我們正在追求拯救其他生命這一本質上令人尊敬的目標。每一個人類生命——包括剛發育的生命——本身都是同樣有價值的，因此不能進行任何不同的評估，甚至不能進行數字衡量。
\end{transpara}

\syncleft
\begin{sourcepara}
\sourcern{158}In der rechtspolitischen Grundkonzeption des Fünften Strafrechtsreformgesetzes wird auch eine Auffassung von der Funktion des rechtsstaatlichen Gesetzes erkennbar, der nicht gefolgt werden kann. Der von der Verfassung geforderte Rechtsschutz für das konkrete einzelne Menschenleben wird zurückgestellt zugunsten einer mehr "sozialtechnischen" Verwendung des Gesetzes als einer gezielten Aktion des Gesetzgebers zur Erreichung eines bestimmten gesellschaftspolitisch erwünschten Zieles, der "Eindämmung der Abtreibungsseuche". Der Gesetzgeber darf aber nicht nur das Ziel im Auge haben, sei es auch noch so erstrebenswert; er muß beachten, daß auch jeder Schritt auf dem Wege dahin sich vor der Verfassung und ihren unverzichtbaren Postulaten zu rechtfertigen hat. Der Effizienz der Regelung im ganzen darf der Grundrechtsschutz im einzelnen nicht geopfert werden. Das Gesetz ist nicht nur Instrument zur Steuerung gesellschaftlicher Prozesse nach soziologischen Erkenntnissen und Prognosen, es ist auch bleibender Ausdruck sozialethischer und - ihr folgend - rechtlicher Bewertung menschlicher Handlungen; es soll sagen, was für den Einzelnen Recht und Unrecht ist.
\end{sourcepara}
\begin{transpara}
\transrn{158}刑法改革第五法的基本法律政策理念也揭示了一種不可遵循的憲法功能觀。憲法所要求的對個人生命的法律保護被擱置一邊，以支持對法律進行更加「社會技術」的運用，而不是立法機關採取有針對性的行動，以實現特定的社會政治理想目標，即「遏制終止妊娠流行病」。然而，立法者不僅必須牢記目標，無論目標多麼令人嚮往；還必須牢記目標。他必須牢記，一路走來的每一步都必須在憲法及其不可或缺的假設面前證明自己的合理性。不能為了保護個人的基本權利而犧牲整個監管的效率。法律不僅是根據社會學發現和預測控制社會進程的工具，也是社會倫理的持久表達，以及隨之而來的對人類行為的法律評估；它的目的是說明對個人來說什麼是對的，什麼是錯的。
\end{transpara}

\syncleft
\begin{sourcepara}
\sourcern{159}c) Einer - prinzipiell abzulehnenden - "Gesamtrechnung" fehlt übrigens auch eine verläßliche tatsächliche Grundlage. Es fehlen ausreichende Anhaltspunkte dafür, daß die Zahl der Schwangerschaftsabbrüche in Zukunft erheblich geringer sein werde als bei der bisherigen gesetzlichen Regelung. Der Regierungsvertreter ist vielmehr vor dem Sonderausschuß für die Strafrechtsreform (7. \abbrWp{}, 25. Sitzung, \abbrStenBer{} \abbrS{} 1451) aufgrund sehr eingehender Erwägungen und Vergleiche zu dem Ergebnis gekommen, daß in der Bundesrepublik nach Einführung der Fristenregelung eine Steigerung der Gesamtzahl legaler und illegaler Aborte um 40\% zu erwarten sei. Diese Berechnung ist allerdings von dem in der mündlichen Verhandlung gehörten Professor \abbrDr{}\abbrDr{} Jürgens in Zweifel gezogen worden. Jedoch läßt das aus dem Ausland, insbesondere aus England nach Inkrafttreten der Abortion Act von 1967 (vgl. die Angaben im Report of the Committee on the Working of the Abortion Act - Lane-Report) und aus der \abbrDDR{} nach Erlaß des Gesetzes über die Unterbrechung der Schwangerschaft vom 9. März 1972 (vgl. Deutsches Ärzteblatt 1974, \abbrS{} 2/65) vorliegende Zahlenmaterial keinen sicheren Schluß auf einen wesentlichen Rückgang der Schwangerschaftsabbrüche zu. Experimente sind aber bei dem hohen Wert des zu schützenden Rechtsgutes nicht zulässig.
\end{sourcepara}
\begin{transpara}
\transrn{159}c) 此外，此種原則上即應拒絕之「總體計算」，也欠缺可靠的事實基礎。並無充分根據可認為，未來終止妊娠之數量會比現行法律規制下顯著降低。相反，政府代表在刑法改革特別委員會前（7. Wp., 25. Sitzung, StenBer. S. 1451），基於相當詳盡之考量與比較，得出結論：在聯邦共和國引入期限規定後，合法及非法墮胎總數預期將增加40\%。此一計算固然受到言詞辯論中聽取意見之 Professor Dr. Dr. Jürgens 質疑。然而，來自外國之數據，尤其是英國於1967年 Abortion Act 生效後之資料（參見 Report of the Committee on the Working of the Abortion Act，Lane-Report 所載資料），以及德意志民主共和國於1972年3月9日制定《妊娠中止法》後之資料（參見 Deutsches Ärzteblatt 1974, S. 2/65），均不容許可靠推論終止妊娠將顯著下降。鑑於受保護法益具有高度價值，並不容許以此作實驗。
\end{transpara}

\syncleft
\begin{sourcepara}
\sourcern{160}Indessen haben es Vertreter aller Parteien im Sonderausschuß für die Strafrechtsreform abgelehnt, ausländische Zahlen über die Schwangerschaftsabbrüche schematisch auf die Bundesrepublik Deutschland umzurechnen (7. \abbrWp{}, 20. Sitzung, \abbrStenBer{} \abbrS{} 1286 \abbrFF{}); die Auswirkungen unterschiedlicher Sozialstrukturen, Mentalitäten, religiöser Bindungen und Verhaltensweisen ließen sich kaum kalkulieren. Selbst wenn man aber alle Besonderheiten der Verhältnisse in der Bundesrepublik Deutschland nur zugunsten der Fristenregelung in Rechnung stellt, ist mit einem Anstieg der Schwangerschaftsabbrüche zu rechnen, weil - wie dargelegt - schon das bloße Bestehen der Strafnorm des § 218 \abbrStGB{} Einfluß auf Wertvorstellungen und Verhaltensweisen der Bevölkerung gehabt hat. Dabei fällt ins Gewicht, daß infolge der Strafbarkeit die Möglichkeit, einen Schwangerschaftsabbruch überhaupt oder gar lege artis zu erlangen, bisher erheblich (u.a. finanziell) eingeschränkt war. Daß von der Fristenlösung eine auch nur quantitative Verstärkung des Lebensschutzes ausgehen könnte, ist jedenfalls nicht ersichtlich.
\end{sourcepara}
\begin{transpara}
\transrn{160}同時，刑法改革特別委員會中的各方代表拒絕將外國終止妊娠數據示意性地換算為德意志聯邦共和國（第 7 屆會議，第 20 屆會議，StenBer. 第 1286 頁及以下）；不同的社會結構、心態、宗教聯繫和行為的影響難以計算。然而，即使僅考慮到德意志聯邦共和國條件的所有特殊特徵，僅支持期限規定，終止妊娠率的增加也是可以預期的，因為正如所解釋的，《StGB》第218條刑法規範的存在本身就對人們的價值觀和行為產生了影響。這裡重要的是，由於刑事責任，迄今為止獲得終止妊娠甚至法律的可能性都受到相當大的限制（包括經濟上）。無論如何，時間限制解決方案能否帶來人壽保護的數量上的增加並不明顯。
\end{transpara}

\syncleft
\begin{sourcepara}
\sourcern{161}3. Die in § 218c \abbrAbs{} 1 \abbrStGB{} vorgesehene Beratung und Unterrichtung der Schwangeren kann - auch für sich betrachtet - nicht als geeignet angesehen werden, auf eine Fortsetzung der Schwangerschaft hinzuwirken.
\end{sourcepara}
\begin{transpara}
\transrn{161}3. 《刑法》第 218c 條第1項中為懷孕婦女提供的諮詢和資訊－即使單獨考慮－也不能被視為適合幫助繼續懷孕。
\end{transpara}

\syncleft
\begin{sourcepara}
\sourcern{162}Die in dieser Bestimmung vorgesehenen Maßnahmen bleiben hinter den Vorstellungen des Alternativ-Entwurfs der 16 Strafrechtswissenschaftler zurück, auf dem die Konzeption des Fünften Strafrechtsreformgesetzes im übrigen weitgehend beruht. Die dort (in § 105 \abbrAbs{} 1 \abbrNr{} 2) vorgesehenen Beratungsstellen sollten die Möglichkeit haben, selbst finanzielle, soziale und familiäre Hilfe zu leisten. Sie sollten ferner der Schwangeren und ihren Angehörigen durch geeignete Mitarbeiter seelische Betreuung gewähren und intensiv auf die Fortsetzung der Schwangerschaft hinwirken (vgl. im einzelnen oben \abbrS{} 11 \abbrF{}).
\end{sourcepara}
\begin{transpara}
\transrn{162}該條項規定的措施與16位刑法學者的替代草案中的想法存在差距，而刑法改革第五法的概念很大程度上是基於該草案。在那裡提供的諮詢中心（§ 105 Abs. 1 Nr. 2）應有機會自行提供經濟、社會和家庭援助。他們也應由適當的員工為懷孕婦女及其親屬提供心理支持，並加緊工作以確保妊娠持續（詳情請參閱上文第 11 f. 頁）。
\end{transpara}

\syncleft
\begin{sourcepara}
\sourcern{163}Die Beratungsstellen etwa im Sinne dieser oder ähnlicher Vorschläge so auszustatten, daß sie unmittelbare Hilfe vermitteln können, hätte um so näher gelegen, als nach dem Bericht des Sonderausschusses für die Strafrechtsreform (\abbrBTDrucks{} 7/1981 [neu] \abbrS{} 7 mit Nachweisen aus dem Anhörungsverfahren) die ungünstige Wohnsituation, die Unmöglichkeit, neben einer Ausbildung oder Erwerbstätigkeit ein Kind zu versorgen, sowie wirtschaftliche Not und sonstige materielle Gründe, außerdem bei ledigen Schwangeren die Angst vor gesellschaftlichen Sanktionen zu den am häufigsten genannten Ursachen und Motiven für den Wunsch nach Schwangerschaftsabbruch gehören sollen.
\end{sourcepara}
\begin{transpara}
\transrn{163}根據刑法改革特別委員會的報告（BTDrucks.7/1981[新]第7頁，附有來自聽證程序的證據），本著這些或類似諮詢的精神為諮詢中心配備使其能夠提供立即幫助的方式就更加明顯了，最常見的是不利的生活條件，在培訓或工作期間無法照顧孩子，以及經濟困難和其他物質原因。由於對社會制裁的恐懼在未婚懷孕婦女中最為常見，因此應包括所提到的希望終止妊娠的原因和動機。
\end{transpara}

\syncleft
\begin{sourcepara}
\sourcern{164}Demgegenüber sollen die Beratungsstellen über die "zur Verfügung stehenden öffentlichen und privaten Hilfen für Schwangere, Mütter und Kinder" unterrichten, "insbesondere über solche Hilfen, die die Fortsetzung der Schwangerschaft und die Lage von Mutter und Kind erleichtern". Dies könnte dahin gedeutet werden, daß die Beratungsstellen nur informieren sollen, ohne auf den Motivationsprozeß gezielt Einfluß zu nehmen. Ob die neutrale Umschreibung der Aufgabe der Beratungsstellen darauf zurückzuführen ist, daß im Sonderausschuß für die Strafrechtsreform die Meinung vertreten wurde, die Schwangere solle durch die Beratung nicht in ihrem Entschluß beeinflußt werden (so Abg. von Schoeler, FDP, 7. \abbrWp{}, 25. Sitzung, \abbrStenBer{} \abbrS{} 1473), kann offenbleiben. Auf eine solche Einflußnahme kommt es jedenfalls entscheidend an, wenn der Beratung ein Schutzeffekt zugunsten des werdenden Lebens zukommen soll. § 218c \abbrAbs{} 1 \abbrNr{} 1 und 2 lassen allerdings auch die Auslegung zu, daß Beratung und Unterrichtung die Schwangere zur Austragung ihrer Schwangerschaft veranlassen sollen. In diesem Sinne ist wohl der Bericht des Sonderausschusses (\abbrBTDrucks{} 7/1981 [neu] \abbrS{} 16) zu verstehen; danach soll die Beratung die gesamten Lebensumstände der Schwangeren berücksichtigen und persönlich-individuell erfolgen, nicht telefonisch oder durch Aushändigung gedruckten Materials (vgl. auch die bereits erwähnte Entschließung des Bundestages, \abbrBTDrucks{} 7/2042).
\end{sourcepara}
\begin{transpara}
\transrn{164}相較之下，諮詢機構應告知「可供懷孕婦女、母親及兒童利用之公共與私人協助」，「尤其是有助於繼續妊娠並改善母親與子女處境之協助」。此可解釋為諮詢機構僅應提供資訊，而不應有意識地影響動機形成過程。諮詢機構任務之中性表述，是否源自刑法改革特別委員會中曾有人主張，懷孕婦女不應因諮詢而在其決定上受到影響（von Schoeler 議員，FDP，第7屆選期，第25次會議，StenBer. S. 1473），可暫置不論。無論如何，若諮詢應對正在發展中生命產生保護效果，此種影響即具有決定性。當然，第218c條第1項第1款及第2款亦容許解釋為：諮詢與告知旨在促使懷孕婦女繼續妊娠。特別委員會報告（BTDrucks. 7/1981 [neu] S. 16）或許即應如此理解；據此，諮詢應考量懷孕婦女之全部生活情況，並以親自、個別方式進行，而非以電話或交付印刷材料方式為之（另參前述聯邦議院決議，BTDrucks. 7/2042）。
\end{transpara}

\syncleft
\begin{sourcepara}
\sourcern{165}Selbst wenn man es für denkbar halten mag, daß eine derartige Beratung eine gewisse Wirkung im Sinne einer Abkehr von dem Entschluß zum Schwangerschaftsabbruch ausüben könnte, so weist jedenfalls ihre Ausgestaltung im einzelnen Mängel auf, die keinen effektiven Schutz des werdenden Lebens erwarten lassen.
\end{sourcepara}
\begin{transpara}
\transrn{165}即使人們認為這樣的諮詢可以在放棄終止妊娠的決定方面產生一定的效果，但其設計在細節上有缺陷，不會引起對正在發展中的生命的有效保護的任何期望。
\end{transpara}

\syncleft
\begin{sourcepara}
\sourcern{166}a) Die Unterrichtung über die zur Verfügung stehenden öffentlichen und privaten Hilfen für Schwangere, Mütter und Kinder nach § 218c \abbrAbs{} 1 \abbrNr{} 1 kann auch von jedem Arzt vorgenommen werden. Sozialrecht und Sozialwesen sind jedoch schon für den fachlich Vorgebildeten schwer zu überblicken. Von einem Arzt kann eine zuverlässige Unterrichtung über die gerade im Einzelfall bestehenden Ansprüche und Möglichkeiten nicht erwartet werden, zumal dafür häufig individuelle Bedürftigkeitsermittlungen erforderlich sind (z.B. für Mietbeihilfe oder Sozialhilfe). Die Ärzte sind für eine solche Beratungstätigkeit weder nach ihrer Berufsausbildung qualifiziert noch steht ihnen im allgemeinen die für eine individuelle Beratung erforderliche Zeit zur Verfügung.
\end{sourcepara}
\begin{transpara}
\transrn{166}a) 任何醫師也可以提供有關根據第218c條第1項第1款向懷孕婦女、母親和兒童提供的公共和私人援助的資訊。然而，即使對於受過專業訓練的人來說，社會法和社會服務也很難理解。不能指望醫師提供有關個別情況下可用的權利和選擇的可靠信息，特別是因為這通常需要進行個人需求評估（例如租金援助或社會援助）。醫師的專業訓練既不具備從事此類諮詢工作的資格，通常也沒有進行個別諮詢所需的時間。
\end{transpara}

\syncleft
\begin{sourcepara}
\sourcern{167}b) Besonders bedenklich ist, daß die Unterrichtung über soziale Hilfen von demselben Arzt vorgenommen werden kann, der den Schwangerschaftsabbruch ausführen soll. Dadurch wird auch die ärztliche Beratung nach § 218c \abbrAbs{} 1 \abbrNr{} 2 entwertet, die an sich in den ärztlichen Aufgabenbereich fällt. Sie soll sich nach den Vorstellungen des Sonderausschusses für die Strafrechtsreform wie folgt gestalten:
\end{sourcepara}
\begin{transpara}
\transrn{167}b) 尤其令人疑慮者，是關於社會協助之告知得由同一名將實施終止妊娠之醫師為之。這也削弱第218c條第1項第2款醫師諮詢之價值；該諮詢本身原屬醫師任務範圍。依刑法改革特別委員會之設想，其應如下形成：
\end{transpara}

\syncleft
\begin{sourcepara}
\sourcern{168}\sourcequoteinner{"Damit ist einmal die Beratung über die Art des Eingriffs und dessen mögliche gesundheitliche Folgen gemeint. Sie darf sich jedoch - was durch die bewußte Wahl des Begriffs "ärztlich" ausgedrückt wird - nicht auf diesen rein medizinischen Aspekt beschränken. Vielmehr muß sie sich in dem jeweils möglichen und angemessenen Umfang auf die gegenwärtige und künftige Gesamtsituation der Schwangeren, soweit sie durch den Schwangerschaftsabbruch tangiert werden kann, erstrecken und zugleich, dem anderen Auftrag des Arztes entsprechend, den Schutz des ungeborenen Lebens miteinbeziehen. Der Arzt muß also die Schwangere auch darüber aufklären, daß durch den Eingriff menschliches Leben vernichtet wird und in welchem Entwicklungsstadium es sich befindet. Die - in der Öffentlichen Anhörung z.B. von Pross (AP VI \abbrS{} 2255, 2256) und Rolinski (AP VI \abbrS{} 2221) bestätigte - Erfahrung zeigt, daß viele Frauen in dieser Hinsicht keine klaren Vorstellungen haben, und daß dieser Umstand, wenn sie später davon erfahren, häufig der Anlaß für schwerwiegende Zweifel und Gewissensbedenken ist. Dementsprechend muß die Beratung mit darauf ausgerichtet sein, derartigen Konfliktsituationen vorzubeugen." (\abbrBTDrucks{} 7/1981 [neu] \abbrS{} 16)}
\end{sourcepara}
\begin{transpara}
\transrn{168}\transquoteinner{「這首先是指關於手術方式及其可能健康後果之諮詢。然而，如同有意選用『醫療的』一詞所表明者，此種諮詢不得限於這一純粹醫學面向。毋寧，凡在個案中可能且適當之範圍內，諮詢必須及於懷孕婦女目前及未來之整體處境，至少及於此一處境可能受終止妊娠影響之部分；同時，依醫師另一項任務，亦須將未出生生命之保護納入其中。因此，醫師也必須向懷孕婦女說明：此一手術將毀滅人類生命，以及該生命正處於何種發展階段。經驗顯示，許多婦女對此並無清楚認識；而當其日後知悉此一情況時，往往會因此產生重大疑慮與良心不安。此一經驗已在公開聽證中例如由 Pross（AP VI S. 2255, 2256）及 Rolinski（AP VI S. 2221）所確認。因此，諮詢也必須以預防此類衝突情境為目標。」（BTDrucks. 7/1981 [neu] S. 16）}
\end{transpara}

\syncleft
\begin{sourcepara}
\sourcern{169}Eine Aufklärung in der hier vorgesehenen Weise mit dem verfassungsrechtlich gebotenen Ziel, auf eine Fortsetzung der Schwangerschaft hinzuwirken, kann von dem Arzt, der von der Schwangeren gerade zu dem Zwecke aufgesucht wird, daß er den Schwangerschaftsabbruch vornehme, nicht erwartet werden. Da nach dem Ergebnis der bisherigen Umfragen und nach den Stellungnahmen repräsentativer ärztlicher Standesgremien angenommen werden muß, daß die Mehrzahl der Ärzte die Vornahme von nicht indizierten Schwangerschaftsabbrüchen ablehnt, werden sich vor allem solche Ärzte zur Verfügung stellen, die im Schwangerschaftsabbruch entweder ein gewinnbringendes Geschäft sehen oder jedem Wunsch einer Frau nach Schwangerschaftsabbruch zu entsprechen geneigt sind, weil sie darin lediglich eine Manifestation des Selbstbestimmungsrechts oder ein Mittel zur Emanzipation der Frau erblicken. In beiden Fällen ist eine Beeinflussung der Schwangeren im Sinne einer Fortsetzung der Schwangerschaft durch den Arzt sehr unwahrscheinlich.
\end{sourcepara}
\begin{transpara}
\transrn{169}懷孕婦女尋求該醫師，正是為使其實施終止妊娠；不能期待此醫師以此處所設想之方式進行說明，並以憲法所要求、促使繼續妊娠為目標。既然依既有調查結果及具代表性醫師職業團體之意見，必須認為多數醫師拒絕實施無指標事由之終止妊娠，則主要會願意提供此服務者，將是將終止妊娠視為有利可圖業務之醫師，或因僅將其視為自我決定權之表現或婦女解放手段，而傾向滿足每一名婦女終止妊娠願望之醫師。在兩種情形中，醫師均極不可能以繼續妊娠之方向影響懷孕婦女。
\end{transpara}

\syncleft
\begin{sourcepara}
\sourcern{170}Dies zeigen die Erfahrungen in England. Dort muß die (sehr weit gefaßte) Indikation von zwei beliebigen Ärzten festgestellt werden. Dies hat dazu geführt, daß bei den darauf spezialisierten Privatärzten praktisch jeder gewünschte Schwangerschaftsabbruch ausgeführt wird. Das Auftreten gewerbsmäßiger Vermittler, die Frauen diesen Privatkliniken zuführen, ist eine besonders unerfreuliche, aber schwer vermeidbare Nebenerscheinung (vgl. Lane-Report, \abbrBd{} 1 \abbrNr{} 436 und 452).
\end{sourcepara}
\begin{transpara}
\transrn{170}英國經驗顯示此點。在該處，（範圍極廣之）指標事由須由任意兩名醫師確認。其結果是，於專門從事此業務之私人醫師處，實際上任何被要求之終止妊娠均會實施。營業性中介者將婦女帶往此等私人診所之現象，是特別令人不快但難以避免之附隨現象（參見 Lane-Report, Bd. 1 Nr. 436 und 452）。
\end{transpara}

\syncleft
\begin{sourcepara}
\sourcern{171}c) Weiterhin ist den Erfolgsaussichten abträglich, daß der Unterrichtung und Beratung der Schwangerschaftsabbruch unmittelbar folgen kann. Eine ernsthafte Auseinandersetzung der Schwangeren und ihrer Angehörigen mit den in der Beratung ihr entgegengehaltenen Argumenten ist unter diesen Umständen nicht zu erwarten. Die vom Bundesjustizministerium dem Sonderausschuß für die Strafrechtsreform vorgelegte Formulierungsalternative für § 218c sah deshalb vor, daß der Schwangerschaftsabbruch erst vorgenommen werden dürfe, nachdem mindestens drei Tage seit der Unterrichtung über die zur Verfügung stehenden Hilfen (§ 218 \abbrAbs{} 1 \abbrNr{} 1) verstrichen seien (Sonderausschuß 7. \abbrWp{}, 30. Sitzung, \abbrStenBer{} \abbrS{} 1659). Indessen wurde nach dem Bericht des Sonderausschusses "auf eine strafrechtlich erzwungene Karenzzeit zwischen den Beratungen und dem Eingriff ... verzichtet. Dies könnte in Einzelfällen für die Schwangere je nach ihrem Wohnort und ihrer persönlichen Situation unzumutbare Schwierigkeiten mit der Folge mit sich bringen, daß die Schwangere auf die Beratung verzichtet" (\abbrBTDrucks{} 7/1981 (neu) \abbrS{} 17). Für die zum Schwangerschaftsabbruch entschlossene Frau kommt es somit nur darauf an, einen willfährigen Arzt zu finden; da er sowohl die soziale wie die ärztliche Beratung vornehmen und schließlich auch den Eingriff durchführen darf, ist von ihm nicht der ernsthafte Versuch zu erwarten, die Schwangere von ihrem Entschluß abzubringen.
\end{sourcepara}
\begin{transpara}
\transrn{171}c) 此外，告知與諮詢之後得立即接續終止妊娠，也不利於成功機會。在此等情況下，不能期待懷孕婦女及其親屬認真面對諮詢中向其提出之論點。因此，聯邦司法部提交刑法改革特別委員會之第218c條替代文字曾規定，終止妊娠須待關於可利用協助之告知（第218條第1項第1款）後至少三日屆滿始得實施（Sonderausschuß 7. Wp., 30. Sitzung, StenBer. S. 1659）。然而，依特別委員會報告，「已放棄以刑法強制諮詢與介入之間的等待期間……。此在個案中，依懷孕婦女居住地及個人情況，可能帶來不可期待之困難，導致懷孕婦女放棄諮詢」（BTDrucks. 7/1981 [neu] S. 17）。對決意終止妊娠之婦女而言，關鍵遂僅在找到一名順從其意願之醫師；既然該醫師既得進行社會及醫師諮詢，最後又得實施介入，即不能期待其會認真嘗試勸阻懷孕婦女放棄其決定。
\end{transpara}

\bisubsubsection{III.}{III.}
\syncleft
\begin{sourcepara}
\sourcern{172}Zusammenfassend ist zur verfassungsrechtlichen Beurteilung der im Fünften Strafrechtsreformgesetz getroffenen Fristenregelung folgendes auszuführen:
\end{sourcepara}
\begin{transpara}
\transrn{172}綜上所述，對於刑法改革第五法中期限規定的合憲性評估，必須說明以下幾點：
\end{transpara}

\syncleft
\begin{sourcepara}
\sourcern{173}Es ist mit der dem Gesetzgeber obliegenden Lebensschutzpflicht unvereinbar, daß Schwangerschaftsabbrüche auch dann rechtlich nicht mißbilligt und nicht unter Strafe gestellt werden, wenn sie aus Gründen erfolgen, die vor der Wertordnung des Grundgesetzes keinen Bestand haben. Zwar wäre die Einschränkung der Strafbarkeit verfassungsrechtlich nicht zu beanstanden, wenn sie mit anderen Maßnahmen verbunden wäre, die den Wegfall des Strafschutzes in ihrer Wirkung zumindest auszugleichen vermöchten. Das ist indes - wie dargelegt - offensichtlich nicht der Fall. Die parlamentarischen Auseinandersetzungen um die Reform des Abtreibungsrechts haben zwar die Einsicht vertieft, daß es die vornehmste Aufgabe des Staates ist, die Abtötung ungeborenen Lebens durch Aufklärung über vorbeugende Schwangerschaftsverhütung einerseits sowie durch wirksame soziale Förderungsmaßnahmen und durch eine allgemeine Veränderung der gesellschaftlichen Auffassungen andererseits zu verhindern. Jedoch vermögen weder die gegenwärtig angebotenen und gewährten Hilfen dieser Art noch die im Fünften Strafrechtsreformgesetz vorgesehene Beratung den individuellen Lebensschutz zu ersetzen, den eine Strafnorm grundsätzlich auch heute noch in den Fällen gewährt, in denen für den Abbruch einer Schwangerschaft kein nach der Wertordnung des Grundgesetzes achtenswerter Grund besteht.
\end{sourcepara}
\begin{transpara}
\transrn{173}若終止妊娠係基於不能在《基本法》價值秩序前成立之理由而實施，法律卻仍不予否定評價且不科以刑罰，則此與立法者所負生命保護義務不相容。誠然，若限制可罰性同時結合其他措施，而其效果至少足以補償刑法保護之消失，則在憲法上尚非不可接受。然而，如前所述，本案顯然並非如此。圍繞墮胎法改革之議會爭論固然加深一項認識：國家最重要之任務，在於一方面透過預防懷孕之資訊，另一方面透過有效社會促進措施及社會觀念之一般變化，防止未出生生命遭殺害。然而，無論是目前提供及給予之此類協助，或《第五刑法改革法》所規定之諮詢，均不能取代刑法規範今日仍原則上提供之個別生命保護；該保護適用於依《基本法》價值秩序並無值得尊重之終止妊娠理由的案件。
\end{transpara}

\syncleft
\begin{sourcepara}
\sourcern{174}Wenn der Gesetzgeber die bisherige undifferenzierte Strafandrohung für den Schwangerschaftsabbruch als ein fragwürdiges Mittel des Lebensschutzes ansieht, so entbindet ihn dies doch nicht von der Verpflichtung, zumindest den Versuch zu unternehmen, durch eine differenziertere strafrechtliche Regelung einen besseren Lebensschutz zu erreichen, indem er diejenigen Fälle unter Strafe stellt, in denen der Schwangerschaftsabbruch verfassungsrechtlich zu mißbilligen ist. Eine klare Abgrenzung dieser Fallgruppe gegenüber den anderen Fällen, in denen die Fortsetzung der Schwangerschaft der Frau nicht zumutbar ist, wird die rechtsbewußtseinsbildende Kraft der Strafnorm verstärken. Wer überhaupt den Vorrang des Lebensschutzes vor dem Anspruch der Frau auf freie Lebensgestaltung anerkennt, wird in diesen durch keine besondere Indikation gedeckten Fällen den Unrechtsgehalt der Tat nicht bestreiten können. Wenn der Staat diese Fälle nicht nur für strafbar erklärt, sondern sie auch in der Rechtspraxis verfolgt und bestraft, wird dies im Rechtsbewußtsein der Allgemeinheit weder als ungerecht noch als unsozial empfunden werden.
\end{sourcepara}
\begin{transpara}
\transrn{174}若立法機關認為迄今對終止妊娠不加區分之刑罰威嚇，是一項可疑之生命保護手段，這仍不免除其至少嘗試透過較具區分性之刑法規制，達成較佳生命保護之義務；亦即，將終止妊娠在憲法上應受否定評價之案件置於刑罰之下。將此類案件與其他繼續妊娠對婦女不可期待之案件清楚區分，將強化刑罰規範形塑法律意識之力量。凡承認生命保護優先於婦女自由生活形成請求者，均不能否認，在此等無特殊指標事由支撐之案件中，行為具有不法內涵。若國家不僅宣告此等案件可罰，且在法律實務中追訴並處罰，在一般法律意識中，這既不會被感受為不正義，也不會被感受為反社會。
\end{transpara}

\syncleft
\begin{sourcepara}
\sourcern{175}Die leidenschaftliche Diskussion der Abtreibungsproblematik mag Anlaß zu der Befürchtung geben, daß in einem Teil der Bevölkerung der Wert des ungeborenen Lebens nicht mehr voll erkannt wird. Das gibt jedoch dem Gesetzgeber nicht das Recht zur Resignation. Er muß vielmehr den ernsthaften Versuch unternehmen, durch eine Differenzierung der Strafandrohung einen wirksameren Lebensschutz und eine Regelung zu erreichen, die auch vom allgemeinen Rechtsbewußtsein getragen wird.
\end{sourcepara}
\begin{transpara}
\transrn{175}對終止妊娠問題的熱烈討論可能會引起人們的擔憂，即部分人口不再充分認識到未出生生命的價值。然而，這並不賦予立法機關辭職的權利。相反，他必須做出認真的嘗試，透過區分刑罰威嚇來實現更有效的生命保護和普遍法律意識支持的監管。
\end{transpara}

\bisubsubsection{IV.}{IV.}
\syncleft
\begin{sourcepara}
\sourcern{176}Die im Fünften Strafrechtsreformgesetz getroffene Regelung wird bisweilen mit dem Hinweis verteidigt, daß in anderen demokratischen Ländern der westlichen Welt in jüngster Zeit die strafrechtlichen Vorschriften über den Schwangerschaftsabbruch in ähnlicher oder noch weitergehender Weise "liberalisiert" oder "modernisiert" worden seien; dies sei ein Anzeichen dafür, daß die Neuregelung jedenfalls der allgemeinen Entwicklung der Anschauungen auf diesem Gebiet entspreche und mit fundamentalen sozialethischen und rechtlichen Prinzipien nicht unvereinbar sei.
\end{sourcepara}
\begin{transpara}
\transrn{176}有時有人為《刑法改革第五法》的規定辯護，指出在西方世界的其他民主國家，有關終止妊娠的刑法條項最近也以類似甚至更廣泛的方式「自由化」或「現代化」；這表明新規定符合該領域的普遍觀點，與基本的社會道德和法律原則並不矛盾。
\end{transpara}

\syncleft
\begin{sourcepara}
\sourcern{177}Diese Erwägungen können die hier zu treffende Entscheidung nicht beeinflussen. Abgesehen davon, daß alle diese ausländischen Regelungen in ihren eigenen Ländern stark umstritten sind, unterscheiden sich die rechtlichen Maßstäbe, die dort für das Handeln des Gesetzgebers gelten, wesentlich von denen der Bundesrepublik Deutschland.
\end{sourcepara}
\begin{transpara}
\transrn{177}這些考慮因素不能影響此處所做的決定。除了所有這些外國法規在其本國都存在很大爭議之外，適用於立法機關行為的法律標準與德意志聯邦共和國的法律標準也有很大不同。
\end{transpara}

\syncleft
\begin{sourcepara}
\sourcern{178}Dem Grundgesetz liegen Prinzipien der Staatsgestaltung zugrunde, die sich nur aus der geschichtlichen Erfahrung und der geistig-sittlichen Auseinandersetzung mit dem vorangegangenen System des Nationalsozialismus erklären lassen. Gegenüber der Allmacht des totalitären Staates, der schrankenlose Herrschaft über alle Bereiche des sozialen Lebens für sich beanspruchte und dem bei der Verfolgung seiner Staatsziele die Rücksicht auch auf das Leben des Einzelnen grundsätzlich nichts bedeutete, hat das Grundgesetz eine wertgebundene Ordnung aufgerichtet, die den einzelnen Menschen und seine Würde in den Mittelpunkt aller seiner Regelungen stellt. Dem liegt, wie das Bundesverfassungsgericht bereits früh ausgesprochen hat (\abbrBVerfGE{} 2, 1 [12]), die Vorstellung zugrunde, daß der Mensch in der Schöpfungsordnung einen eigenen selbständigen Wert besitzt, der die unbedingte Achtung vor dem Leben jedes einzelnen Menschen, auch dem scheinbar sozial "wertlosen", unabdingbar fordert und der es deshalb ausschließt, solches Leben ohne rechtfertigenden Grund zu vernichten. Diese Grundentscheidung der Verfassung bestimmt Gestaltung und Auslegung der gesamten Rechtsordnung. Auch der Gesetzgeber ist ihr gegenüber nicht frei; gesellschaftspolitische Zweckmäßigkeitserwägungen, ja staatspolitische Notwendigkeiten können diese verfassungsrechtliche Schranke nicht überwinden (\abbrBVerfGE{} 1, 14 [36]). Auch ein allgemeiner Wandel der hierüber in der Bevölkerung herrschenden Anschauungen - falls er überhaupt festzustellen wäre - würde daran nichts ändern können. Das Bundesverfassungsgericht, dem von der Verfassung aufgetragen ist, die Beachtung ihrer grundlegenden Prinzipien durch alle Staatsorgane zu überwachen und gegebenenfalls durchzusetzen, kann seine Entscheidungen nur an diesen Prinzipien orientieren, zu deren Entfaltung es selbst in seiner Rechtsprechung entscheidend beigetragen hat. Damit wird kein absprechendes Urteil über andere Rechtsordnungen gefällt, "die diese Erfahrungen mit einem Unrechtssystem nicht gemacht haben und die aufgrund einer anders verlaufenen geschichtlichen Entwicklung, anderer staatspolitischer Gegebenheiten und staatsphilosophischer Grundauffassungen eine solche Entscheidung für sich nicht getroffen haben" (\abbrBVerfGE{} 18, 112 [117]).
\end{sourcepara}
\begin{transpara}
\transrn{178}《基本法》所依據的國家形成原則只能透過歷史經驗以及與先前國家社會主義制度的知識和道德辯論來解釋。與極權國家的無所不能相比，極權國家聲稱對社會生活的所有領域進行無限的控制，在追求國家目標時考慮個人的生活根本沒有任何意義，《基本法》建立了一種以價值為基礎的秩序，將個人及其尊嚴置於其所有法規的中心。正如聯邦憲法法院早先指出的那樣(BVerfGE 2, 1 [12])，這是基於人類在創造秩序中擁有自己獨立價值的理念，這必然要求無條件尊重每個個體的生命，即使是那些在社會上看起來「毫無價值」的人，因此排除了無正當理由破壞這種生命的可能性。憲法的這一根本決定決定了整個法律體系的設計和解釋。即使是立法機關也不能隨意處理它；社會政治權宜考慮，實際上是國家政策的必要性，無法克服這一憲法障礙（BVerfGE 1, 14 [36]）。即使人們對這個問題的普遍看法發生了普遍變化——如果能夠被發現的話——也無法改變這一點。憲法授權聯邦憲法法院監督並在必要時強制所有國家機關遵守其基本原則，聯邦憲法法院只能根據這些原則做出決定，而聯邦憲法法院本身在其判例中對這些原則的發展做出了決定性的貢獻。這意味著，不會對其他「沒有經歷過不公正制度的經歷，並且由於不同的歷史發展、其他國家政治環境和基本國家哲學觀點而沒有為自己做出這樣的決定」的法律制度做出負面判斷（BVerfGE 18, 112 [117]）。
\end{transpara}

\bisubsection{E.}{E.}
\syncleft
\begin{sourcepara}
\sourcern{179}Nach alledem ist § 218a \abbrStGB{} in der Fassung des Fünften Strafrechtsreformgesetzes mit \abbrArt{} 2 \abbrAbs{} 2 Satz 1 in Verbindung mit \abbrArt{} 1 \abbrAbs{} 1 \abbrGG{}\ggfnLifeDignityDE{} insoweit unvereinbar, als er den Schwangerschaftsabbruch auch dann von der Strafbarkeit ausnimmt, wenn keine Gründe vorliegen, die nach den vorstehenden Ausführungen vor der Wertordnung des Grundgesetzes Bestand haben. In diesem Umfang war die Nichtigkeit der Vorschrift festzustellen. Es ist Sache des Gesetzgebers, die Fälle des indizierten und des nicht indizierten Schwangerschaftsabbruchs näher voneinander anzugrenzen. Im Interesse der Rechtsklarheit bis zum Inkrafttreten der gesetzlichen Regelung erschien es geboten, gemäß § 35 \abbrBVerfGG{} eine Anordnung des aus dem Urteilstenor ersichtlichen Inhalts zu erlassen.
\end{sourcepara}
\begin{transpara}
\transrn{179}綜上所述，《第五刑法改革法》版本《刑法》第218a條，在其即使於不存在依前述說明可在《基本法》價值秩序前成立之理由時，仍將終止妊娠排除於可罰性之外之範圍內，與《基本法》第2條第2項第1句結合第1條第1項\ggfnLifeDignityZH{}不相容。在此範圍內，應確認該規定無效。由立法者負責進一步界定具指標事由與未具指標事由之終止妊娠案件。為在法律規制生效前維持法律明確性，依《聯邦憲法法院法》第35條發布一項內容如判決主文所示之命令，乃屬必要。
\end{transpara}

\syncleft
\begin{sourcepara}
\sourcern{180}Es bestand kein Anlaß, weitere Vorschriften des Fünften Strafrechtsreformgesetzes für nichtig zu erklären.
\end{sourcepara}
\begin{transpara}
\transrn{180}沒有理由宣布《刑法改革第五法》的任何進一步規定無效。
\end{transpara}

\syncleft
\begin{sourcepara}
{\centering \abbrDr{} Benda, Ritterspach, \abbrDr{} Haager, Rupp-v.Brünneck, \abbrDr{} Böhmer, \abbrDr{} Faller, \abbrDr{} Brox, \abbrDr{} Simon\par}
\end{sourcepara}
\begin{transpara}
Dr. Benda、Ritterspach、Dr. Haager、Rupp-v.Brünneck、Dr. Böhmer、Dr. Faller、Dr. Brox、Dr. Simon
\end{transpara}

\bisection{Abweichende Meinung}{不同意見}
\syncleft
\begin{sourcepara}
Abweichende Meinung der Richterin Rupp-v. Brünneck und des Richters \abbrDr{} Simon zum Urteil des Ersten Senats des Bundesverfassungsgerichts vom 25.Februar 1975 - 1 \abbrBvF{} 1, 2, 3, 4, 5, 6/74 -
\end{sourcepara}
\begin{transpara}
Rupp-v. Brünneck 法官與 Simon 法官就聯邦憲法法院第一庭1975年2月25日判決（1 BvF 1, 2, 3, 4, 5, 6/74）所提出之不同意見
\end{transpara}

\syncleft
\begin{sourcepara}
\sourcern{1}Das Leben jedes einzelnen Menschen ist selbstverständlich ein zentraler Wert der Rechtsordnung. Unbestritten umfaßt die verfassungsrechtliche Pflicht zum Schutz dieses Lebens auch seine Vorstufe vor der Geburt. Die Auseinandersetzungen im Parlament und vor dem Bundesverfassungsgericht betrafen nicht das Ob, sondern allein das Wie dieses Schutzes. Die Entscheidung hierüber gehört in die Verantwortung des Gesetzgebers. Aus der Verfassung kann unter Umständen eine Pflicht des Staates hergeleitet werden, den Schwangerschaftsabbruch in jedem Stadium der Schwangerschaft unter Strafe zu stellen. Der Gesetzgeber durfte sich sowohl für die Beratungs- und Fristenregelung wie für die Indikationenlösung entscheiden.
\end{sourcepara}
\begin{transpara}
\transrn{1}每一個個別人的生命，當然是法律秩序之核心價值。毫無疑問，保護此一生命之憲法義務，亦包含其出生前之前階段。議會及聯邦憲法法院前之爭論，並非是否保護，而僅是如何保護。此項決定屬於立法者責任。憲法或許在特定情況下可導出國家義務，要求將妊娠任何階段之終止妊娠置於刑罰之下；但立法者得選擇諮詢與期限規定，亦得選擇指標方案。
\end{transpara}

\syncleft
\begin{sourcepara}
\sourcern{2}Eine entgegengesetzte Verfassungsauslegung ist mit dem freiheitlichen Charakter der Grundrechtsnormen nicht vereinbar und verlagert in folgenschwerem Ausmaß Entscheidungskompetenzen auf das Bundesverfassungsgericht (A). Bei der Beurteilung des Fünften Strafrechtsreformgesetzes vernachlässigt die Mehrheit die Singularität des Schwangerschaftsabbruchs im Verhältnis zu anderen Gefährdungen des Lebens (B I 1). Sie würdigt nicht hinreichend die vom Gesetzgeber vorgefundene soziale Problematik sowie die Ziele der dringlichen Reform (B I 2). Schon weil jede Lösung Stückwerk bleibt, ist es verfassungsrechtlich nicht zu beanstanden, daß der deutsche Gesetzgeber - im Einklang mit den Reformen in anderen westlichen Kulturstaaten (B III) - sozialpolitischen Maßnahmen den Vorrang vor weitgehend wirkungslosen Strafdrohungen gegeben hat (B I 3-5). Eine gesetzliche "Mißbilligung" sittlich nicht achtenswerten Verhaltens ohne Rücksicht auf ihre tatsächliche Schutzwirkung schreibt die Verfassung nirgends vor (B II).
\end{sourcepara}
\begin{transpara}
\transrn{2}相反之憲法解釋，與基本權規範之自由性格不相容，並以後果重大之程度，將決定權限轉移至聯邦憲法法院（A）。評價《第五刑法改革法》時，多數意見忽略終止妊娠相較於其他生命危害之獨特性（B I 1）。其亦未充分評價立法者所面對之社會問題，以及迫切改革之目標（B I 2）。正因任何解決方案均仍屬片段，德國立法者與其他西方文化國家之改革一致（B III），使社會政策措施優先於大致無效之刑罰威嚇（B I 3-5），在憲法上並無可非難之處。憲法並未在任何地方規定：不問實際保護效果如何，法律均須對倫理上不值得尊重之行為作成「否定評價」（B II）。
\end{transpara}

\bicompositeheading{A.}{A.}{I.}{I.}
\syncleft
\begin{sourcepara}
\sourcern{3}Die Befugnis des Bundesverfassungsgerichts, Entscheidungen des parlamentarischen Gesetzgebers zu annullieren, erfordert einen sparsamen Gebrauch, wenn eine Verschiebung der Gewichte zwischen den Verfassungsorganen vermieden werden soll. Das Gebot richterlicher Selbstbeschränkung (judicial self-restraint), das als das "Lebenselexier" der Rechtsprechung des Bundesverfassungsgerichts bezeichnet worden ist (Leibholz, \abbrVVDStRL{} 20 [1963], \abbrS{} 119), gilt vor allem, wenn es sich nicht um die Abwehr von Übergriffen der staatlichen Gewalt handelt, sondern wenn dem Volk unmittelbar legitimierten Gesetzgeber im Wege der verfassungsgerichtlichen Kontrolle Vorschriften für die positive Gestaltung der Sozialordnung gemacht werden sollen. Hier darf das Bundesverfassungsgericht nicht der Versuchung erliegen, selbst die Funktion des zu kontrollierenden Organs zu übernehmen, soll nicht auf lange Sicht die Stellung der Verfassungsgerichtsbarkeit gefährdet werden.
\end{sourcepara}
\begin{transpara}
\transrn{3}聯邦憲法法院撤銷議會立法者決定之權限，若要避免憲法機關間權重移動，即須節制使用。司法自我克制（judicial self-restraint）之要求，曾被稱為聯邦憲法法院判例之「生命靈藥」（Leibholz, VVDStRL 20 [1963], S. 119），尤其適用於下列情形：所涉並非抵禦國家權力之侵害，而是透過憲法法院審查（verfassungsgerichtliche Kontrolle），對由人民直接賦予民主正當性之立法者，作成積極形塑社會秩序之規定。若不欲長期危及憲法司法之地位，聯邦憲法法院在此不得屈服於誘惑，自行接管受審查機關之功能。
\end{transpara}

\syncleft
\begin{sourcepara}
\sourcern{4}1. Die in diesem Verfahren begehrte Prüfung verläßt den Boden der klassischen verfassungsgerichtlichen Kontrolle. Die im Zentrum unserer Verfassung stehenden Grundrechtsnormen gewährleisten als Abwehrrechte dem Bürger im Verhältnis zum Staat einen Bereich freier, eigenverantwortlicher Lebensgestaltung. Die klassische Funktion des Bundesverfassungsgerichts liegt insoweit darin, Verletzungen dieses Freiheitsraums durch übermäßige Eingriffe der staatlichen Gewalt abzuwehren. In der Skala der staatlichen Eingriffsmöglichkeiten stehen Strafvorschriften an der Spitze: Sie befehlen dem Bürger ein bestimmtes Verhalten und unterwerfen ihn bei Zuwiderhandlungen empfindlichen Freiheitsbeschränkungen oder finanziellen Belastungen. Verfassungsgerichtliche Kontrolle solcher Vorschriften bedeutet daher die Prüfung, ob der mit dem Erlaß oder der Anwendung der Strafvorschrift verbundene Eingriff in die grundrechtlich geschützte Freiheitssphäre zulässig ist, ob also der Staat überhaupt oder in dem vorgesehenen Umfang strafen darf.
\end{sourcepara}
\begin{transpara}
\transrn{4}1. 本程序所尋求的審查脫離了經典憲法審查的基礎。我國憲法核心的基本權利規範作為防禦性權利，保證公民在與國家的關係中享有自由、自我負責的生活。聯邦憲法法院的經典功能是防止國家權力過度干預而侵犯這一自由領域。在國家干預選項的規模中，刑事法規處於首位：它們命令公民以某種方式行事，如果公民違反這些規定，他們的自由或經濟負擔就會受到嚴格限制。因此，對此類條項的憲法審查意味著審查與刑事條項的頒布或適用相關的對受憲法保護的自由領域的干預是否允許，即是否允許國家施以刑罰或達到所設想的程度。
\end{transpara}

\syncleft
\begin{sourcepara}
\sourcern{5}Im vorliegenden Verfassungsstreit wird gerade umgekehrt erstmalig zur Prüfung gestellt, ob der Staat strafen muß, nämlich ob die Aufhebung der Strafvorschrift gegen den Schwangerschaftsabbruch in den ersten drei Monaten der Schwangerschaft mit den Grundrechten vereinbar ist. Es liegt aber auf der Hand, daß das Absehen von Strafe das Gegenteil eines staatlichen Eingriffs ist. Da die teilweise Rücknahme der Strafvorschrift nicht zur Begünstigung von Schwangerschaftsabbrüchen geschah, sondern weil sich die bisherige Strafdrohung nach der durch Erfahrung erhärteten, unwiderlegten Annahme des Gesetzgebers weitgehend als wirkungslos erwiesen hat, läßt sich nicht einmal mittelbar ein staatlicher "Eingriff" in das ungeborene Leben konstruieren. Weil es daran fehlt, hat der Österreichische Verfassungsgerichtshof eine Verletzung des nach innerösterreichischem Recht geltenden Grundrechtskatalogs durch die dortige Fristenregelung verneint (Vgl. das Erkenntnis vom 11. Oktober 1974 - G 8/74 - II 2 b der Entscheidungsgründe, \abbrEuGRZ{} 1975, \abbrS{} 74 [76].).
\end{sourcepara}
\begin{transpara}
\transrn{5}在目前的憲法爭議中，國家是否必須科以刑罰的問題，首次探討的恰恰是相反的問題，即廢除懷孕前三個月終止妊娠的刑事條項是否與基本權利相一致。但很明顯，不科以刑罰與國家干預相反。由於刑事條項的部分撤回並不是為了促進終止妊娠，而是因為根據立法者無可辯駁的假設並已被經驗所證實，先前的刑罰威嚇已被證明基本上無效，因此甚至不能間接構建國家對胎兒生命的「干預」。由於這一點缺失，奧地利憲法法院通過期限規定否認了對奧地利國內法適用的基本權利目錄的侵犯（參見 1974 年 10 月 11 日的決定 - G 8/74 - II 2 b 決定理由，EuGRZ 1975，第 74 頁 [76]）。
\end{transpara}

\syncleft
\begin{sourcepara}
\sourcern{6}2. Da die Grundrechte als Abwehrrechte von vornherein ungeeignet sind, den Gesetzgeber an der Beseitigung von Strafvorschriften zu hindern, will die Senatsmehrheit die Grundlage dafür in der weitergehenden Bedeutung der Grundrechte als objektive Wertentscheidungen finden (C I 3 und C III b). Danach normieren die Grundrechte nicht nur Abwehrrechte des Einzelnen gegen den Staat, sondern enthalten zugleich objektive Wertentscheidungen, deren Verwirklichung durch aktives Handeln ständige Aufgabe der staatlichen Gewalt ist. Diese Auffassung ist vom Bundesverfassungsgericht in dem begrüßenswerten Bemühen entwickelt worden, den Grundrechten in ihrem freiheitssichernden und auf soziale Gerechtigkeit angelegten Gehalt größere Wirksamkeit zu verleihen. Die Senatsmehrheit berücksichtigt jedoch die für die verfassungsgerichtliche Kontrolle wesentlichen Unterschiede der beiden Grundrechtsaspekte nicht hinreichend.
\end{sourcepara}
\begin{transpara}
\transrn{6}2. 由於作為防禦性權利的基本權利從一開始就不適合阻止立法機關廢除刑事條項，本庭多數意見派希望在作為客觀價值決定的基本權利的更廣泛含義中找到這一點的基礎（C I 3和C III b）。據此，基本權利不僅規範了個人對抗國家的自衛權利，而且包含客觀的價值決定，透過積極行動實現這一價值決定是國家權力的永恆任務。這一觀點是由聯邦憲法法院提出的，旨在使基本權利的內容更加有效，確保自由並以社會正義為目標。然而，本庭多數意見派並沒有充分考慮憲法審查所必需的基本權利的兩個方面之間的差異。
\end{transpara}

\syncleft
\begin{sourcepara}
\sourcern{7}Als Abwehrrechte haben die Grundrechte einen verhältnismäßig deutlich erkennbaren Inhalt; in ihrer Auslegung und Anwendung hat die Rechtsprechung praktikable, allgemein anerkannte Kriterien zur Kontrolle staatlicher Eingriffe - etwa den Grundsatz der Verhältnismäßigkeit - entwickelt. Demgegenüber ist es regelmäßig eine höchst komplexe Frage, wie eine Wertentscheidung durch aktive Maßnahmen des Gesetzgebers zu verwirklichen ist. Die notwendig allgemein gehaltenen Wertentscheidungen könnten insoweit etwa als Verfassungsaufträge charakterisiert werden, die zwar für alles staatliche Handeln richtungweisend, aber notwendig auf eine Umsetzung in verbindliche Regelungen angewiesen sind. Je nach der Beurteilung der tatsächlichen Verhältnisse, der konkreten Zielsetzungen und ihrer Priorität, der Eignung der denkbaren Mittel und Wege sind sehr verschiedene Lösungen möglich. Die Entscheidung, die häufig Kompromisse voraussetzt und sich im Verfahren des trial and error vollzieht, gehört nach dem Grundsatz der Gewaltenteilung und dem demokratischen Prinzip in die Verantwortung des vom Volk unmittelbar legitimierten Gesetzgebers (Vgl. dazu näher unsere abweichende Meinung im Hochschul-Urteil, \abbrBVerfGE{} 37, 148 [150, 153, 155 \abbrF{}]).
\end{sourcepara}
\begin{transpara}
\transrn{7}基本權利作為防禦性權利，具有相對清晰的識別內容；在其解釋和應用中，判例法制定了控制國家干預的實用的、普遍認可的標準——例如比例原則。相較之下，立法機關如何透過積極措施落實價值決策通常是一個高度複雜的問題。價值決定必然是普遍性的，可以被描述為憲法授權，雖然指導所有國家行動，但必然取決於具有約束力的法規的實施。根據對實際情況的評估、具體目標及其優先順序以及可設想的手段和方式的適用性，可能有非常不同的解決方案。這個決定往往需要妥協，並透過反覆試驗來進行，根據分權原則和民主原則，立法機關負責，由人民直接合法化（見我們在大學判決中的反對意見，BVerfGE 37, 148 [150, 153, 155 f.]）。
\end{transpara}

\syncleft
\begin{sourcepara}
\sourcern{8}Freilich kann schon wegen der wachsenden Bedeutung fördernder Sozialmaßnahmen für die Effektuierung der Grundrechte auch in diesem Bereich nicht auf jede verfassungsgerichtliche Kontrolle verzichtet werden; die Erarbeitung eines geeigneten, die Gestaltungsfreiheit des Gesetzgebers respektierenden Instrumentariums wird möglicherweise zu den Hauptaufgaben der Rechtsprechung in den nächsten Jahrzehnten gehören. Solange es aber daran noch fehlt, droht die Gefahr, daß die verfassungsgerichtliche Kontrolle sich nicht auf die Nachprüfung der vom Gesetzgeber getroffenen Entscheidung beschränkt, sondern diese durch eine andere, vom Gericht für besser gehaltene ersetzt. Diese Gefahr besteht in erhöhtem Maße, wenn - wie hier - in stark kontroversen Fragen eine nach langen Auseinandersetzungen getroffene Entscheidung der Parlamentsmehrheit von der unterlegenen Minderheit vor dem Bundesverfassungsgericht angegriffen wird. Unbeschadet der legitimen Befugnis der Antragsberechtigten, verfassungsrechtliche Zweifel auf diesem Wege klären zu lassen, gerät hier das Bundesverfassungsgericht unversehens in die Lage, als politische Schiedsinstanz für die Auswahl zwischen konkurrierenden Gesetzgebungsprojekten in Anspruch genommen zu werden.
\end{sourcepara}
\begin{transpara}
\transrn{8}誠然，鑑於促進性社會措施對實現基本權之重要性日益增加，在此領域亦不能放棄一切憲法法院審查；發展一套尊重立法者形成自由之適當工具，可能會成為未來數十年判例法之主要任務。惟只要此等工具尚不存在，即有一項危險：憲法法院審查不再限於覆核立法者所作決定，而是以法院認為較佳之另一決定取而代之。當如本案般，在高度爭議問題上，議會多數經長期爭論後作成之決定，由敗北少數向聯邦憲法法院攻擊時，此種危險更形升高。在不損及聲請權人以此途徑澄清憲法疑義之正當權限下，聯邦憲法法院在此不知不覺即被置於一種位置：被請求作為政治仲裁機關，在相互競爭之立法計畫間作成選擇。
\end{transpara}

\syncleft
\begin{sourcepara}
\sourcern{9}Der Gedanke der objektiven Wertentscheidung darf aber nicht zum Vehikel werden, um spezifisch gesetzgeberische Funktionen in der Gestaltung der Sozialordnung auf das Bundesverfassungsgericht zu verlagern. Sonst würde das Gericht in eine Rolle gedrängt, für die es weder kompetent noch ausgerüstet ist. Daher sollte das Bundesverfassungsgericht weiter die Zurückhaltung wahren, die es bis zum Hochschul-Urteil geübt hat (vgl. \abbrBVerfGE{} 4, 7 [18]; 27, 253 [283]; 33, 303 [333 \abbrF{}]; 35, 148 - abw.M. - [152 \abbrFF{}]; 36, 321 [330 \abbrFF{}]). Es darf dem Gesetzgeber nur dann entgegentreten, wenn er eine Wertentscheidung ganz außer acht gelassen hat oder die Art und Weise ihrer Realisierung offensichtlich fehlsam ist. Demgegenüber legt die Mehrheit dem Gesetzgeber trotz vermeintlicher Anerkennung seiner Gestaltungsfreiheit faktisch zur Last, er habe eine an sich anerkannte Wertentscheidung nach ihrer Auffassung nicht bestmöglich verwirklicht. Sollte das zum allgemeinen Prüfungsmaßstab werden, so wäre damit das Gebot richterlicher Selbstbeschränkung preisgegeben.
\end{sourcepara}
\begin{transpara}
\transrn{9}然而，客觀價值決定的概念絕不能成為將社會秩序設計中的具體立法功能轉移給聯邦憲法法院的工具。否則，法院將被迫承擔其既無能力也無能力承擔的角色。因此，聯邦憲法法院應繼續行使其所行使的限制，直至大學作出裁決（參見 BVerfGE 4, 7 [18]；27, 253 [283]；33, 303 [333 ff.]；35, 148 - abw.M. - [152 ff.]；只有當立法機關完全忽視價值決定或其實施方式明顯不正確時，它才可能反對立法機關。相較之下，多數意見儘管據稱承認其設計自由，但實際上指責立法機關沒有以他們認為的最佳方式實施公認的價值決策。如果這成為審查的普遍標準，那麼司法自我約束的要求就被放棄了。
\end{transpara}

\bisubsubsection{II.}{II.}
\syncleft
\begin{sourcepara}
\sourcern{10}1. Unser stärkstes Bedenken richtet sich dagegen, daß erstmals in der verfassungsgerichtlichen Rechtsprechung eine objektive Wertentscheidung dazu dienen soll, eine Pflicht des Gesetzgebers zum Erlaß von Strafnormen, also zum stärksten denkbaren Eingriff in den Freiheitsbereich des Bürgers zu postulieren. Dies verkehrt die Funktion der Grundrechte in ihr Gegenteil. Wenn die in einer Grundrechtsnorm enthaltene objektive Wertentscheidung zum Schutz eines bestimmten Rechtsgutes genügen soll, um daraus die Pflicht zum Strafen herzuleiten, so könnten die Grundrechte unter der Hand aus einem Hort der Freiheitssicherung zur Grundlage einer Fülle von freiheitsbeschränkenden Reglementierungen werden. Was für den Schutz des Lebens gilt, kann auch für andere Rechtsgüter von hohem Rang - etwa körperliche Unversehrtheit, Freiheit, Ehe und Familie - in Anspruch genommen werden.
\end{sourcepara}
\begin{transpara}
\transrn{10}1. 我們最關心的是，在憲法判例中，客觀價值決定首次旨在假定立法機關有義務頒布刑事規範，即對公民的自由施加最大程度的干擾。這將基本權利的功能逆轉為相反的功能。如果基本權利規範所包含的客觀價值決定足以保護某種法益，從而衍生出刑罰義務，那麼基本權利就可能秘密地從保障自由的避風港變成大量限制自由的法規的基礎。適用於保護生命的內容也適用於其他高度重要的法律利益，例如身體健全、自由、婚姻和家庭。
\end{transpara}

\syncleft
\begin{sourcepara}
\sourcern{11}Selbstverständlich setzt die Verfassung voraus, daß der Staat zum Schutz eines geordneten Zusammenlebens auch seine Strafgewalt gebrauchen kann; der Sinn der Grundrechte geht jedoch nicht dahin, solchen Einsatz zu fordern, sondern ihm Grenzen zu ziehen. So hat der Supreme Court der Vereinigten Staaten die Bestrafung von Schwangerschaftsabbrüchen, die mit Einwilligung der Schwangeren im ersten Drittel der Schwangerschaft durch einen Arzt vorgenommen werden, sogar als Grundrechtsverletzung angesehen (Roe v. Wade, 410 \abbrUS{} 113 [1973] = 93 \abbrS{}Ct. 705 = 41 USLaw Week 4213.). Dies ginge nach deutschem Verfassungsrecht zwar zu weit. Jedoch braucht auch nach dem freiheitlichen Charakter unserer Verfassung der Gesetzgeber eine verfassungsrechtliche Rechtfertigung dafür, daß er straft, nicht aber dafür, daß er von Strafe absieht, weil nach seiner Auffassung eine Strafdrohung keinen Erfolg verspricht oder aus anderen Gründen als unangemessene Reaktion erscheint (vgl. \abbrBVerfGE{} 22, 49 [78]; 27, 18 [28]; 32, 40 [48]).
\end{sourcepara}
\begin{transpara}
\transrn{11}當然，憲法預設國家為保護有序共同生活，亦得使用其刑罰權力；然而，基本權利之意義並非要求國家動用刑罰，而是為此種動用劃定界限。例如，美國最高法院甚至認為，對於懷孕前三分之一期間、經懷孕婦女同意並由醫師實施之終止妊娠加以處罰，構成基本權侵害（Roe v. Wade, 410 U.S. 113 [1973] = 93 S.Ct. 705 = 41 USLaw Week 4213）。依德國憲法法理，此一見解固然走得過遠。然而，即使依我國憲法之自由性格，立法機關也須為其處罰提出憲法上正當化理由；相反地，若立法機關認為刑罰威嚇無成功希望，或基於其他理由認為刑罰是不相當反應，而決定不予處罰，則無須為此提出同樣的正當化（參見 BVerfGE 22, 49 [78]; 27, 18 [28]; 32, 40 [48]）。
\end{transpara}

\syncleft
\begin{sourcepara}
\sourcern{12}Die entgegengesetzte Grundrechtsinterpretation führt zwangsläufig zu einer nicht minder bedenklichen Ausweitung der verfassungsgerichtlichen Kontrolle: Es ist nicht mehr allein zu prüfen, ob eine Strafvorschrift zu weit in den Rechtsbereich des Bürgers eingreift, sondern auch umgekehrt, ob der Staat zu wenig straft. Dabei wird das Bundesverfassungsgericht entgegen der Mehrheitsmeinung (D I) nicht bei der Frage stehenbleiben können, ob der Erlaß irgendeiner Strafnorm gleich welchen Inhalts geboten ist, sondern klären müssen, welche Strafsanktion zum Schutz des jeweiligen Rechtsguts hinreicht. In letzter Konsequenz könnte das Gericht sogar zu der Prüfung genötigt sein, ob die Anwendung einer Strafnorm im Einzelfall dem Schutzgedanken genügt.
\end{sourcepara}
\begin{transpara}
\transrn{12}對基本權利的相反解釋不可避免地導致憲法審查的同樣令人擔憂的擴張：不再需要僅僅審查刑事條項是否過多侵犯公民的法律領域，而且還需要反過來審查國家刑罰是否過少。與多數意見（D I）相反，聯邦憲法法院不能停留在製定任何刑法（無論其內容如何）是否有必要的問題上，而必須澄清哪種刑事制裁足以保護各自的法益。最終，法院甚至可能被迫審查個別案件中刑事條項的適用是否滿足保護的理念。
\end{transpara}

\syncleft
\begin{sourcepara}
\sourcern{13}Ein verfassungsrechtliches Festschreiben von Strafnormen - wie die Mehrheit es fordert - ist schließlich auch deswegen abzulehnen, weil gerade die Leitgedanken des Strafrechts nach den Erfahrungen der letzten Jahrzehnte und der zu erwartenden Entwicklung im Bereich der Sozialwissenschaften einen schnellen und starken Wandel unterliegen. Dies zeigt nicht allein ein Blick auf die grundlegenden Veränderungen etwa in der Beurteilung der Sittlichkeitsdelikte - z.B. Homosexualität, Ehegattenkuppelei, Exhibitionismus -, sondern läßt sich speziell für die Strafvorschriften gegen Abtreibung belegen. Die Straflosigkeit der ethisch (kriminologisch) indizierten Abtreibung, die heute der ganz überwiegenden Rechtsauffassung entspricht, war noch in den sechziger Jahren äußerst umstritten (Vgl. die Auseinandersetzung im Bundesrat und Bundestag [Niederschrift über die 254. Sitzung des Bundesratsrechtsausschusses vom 26. Juni 1962 \abbrS{} 30 \abbrFF{}; Verhandlungen des Bundesrats 1962 \abbrS{} 140 \abbrF{}, 153, 154 \abbrF{}; Verhandlunben des Deutschen Bundestages, 4. \abbrWp{}, SenBer. der 70. Sitzung vom 28. März 1963, \abbrS{} 3188, 3208, 3210, 3217, 3221]; s. ferner die Nachweise bei Lang-Hinrichsen, JZ 1963, \abbrS{} 725 \abbrFF{}). Die von der Bundesregierung 1960 und 1962 vorgelegten Entwürfe eines Strafgesetzbuches lehnten diese Indikation ausdrücklich ab (Vgl. §§ 140 \abbrF{}, 157 und die Begründung \abbrBTDrucks{} III/2150 \abbrS{} 262, 274 \abbrF{}; \abbrBTDrucks{} IV/650 \abbrS{} 278, 292 \abbrF{}); zur sozialen und eugenischen Indikation begnügten sie sich mit dem Hinweis, daß sich die Ablehnung "von selbst verstehe" (\abbrBTDrucks{} III/2150 \abbrS{} 262, IV/650 \abbrS{} 278.).
\end{sourcepara}
\begin{transpara}
\transrn{13}如多數意見所要求，將刑事規範以憲法方式固定下來，最終也應予拒絕；因為依過去數十年之經驗，以及社會科學領域可預期之發展，刑法之指導思想正迅速且劇烈地變動。這不僅可從對風化犯罪之評價發生根本變化看出，例如同性戀、配偶媒介性行為、暴露癖；亦可特別就反墮胎刑罰規定加以證明。今日已符合絕大多數法律見解之倫理上（犯罪學上）有指標事由之墮胎不罰，在1960年代仍極具爭議（參見聯邦參議院與聯邦議院之爭論：1962年6月26日聯邦參議院法律委員會第254次會議紀錄 S. 30 ff.; Verhandlungen des Bundesrats 1962 S. 140 f., 153, 154 f.; Verhandlungen des Deutschen Bundestages, 4. Wp., StenBer. der 70. Sitzung vom 28. März 1963, S. 3188, 3208, 3210, 3217, 3221；另見 Lang-Hinrichsen, JZ 1963, S. 725 ff. 所列證據）。聯邦政府於1960年及1962年提出之刑法典草案明文拒絕此一指標事由（參見 §§ 140 f., 157 及理由書 BTDrucks. III/2150 S. 262, 274 f.; BTDrucks. IV/650 S. 278, 292 f.）；至於社會與優生指標事由，草案僅以拒絕乃「不言自明」作為說明（BTDrucks. III/2150 S. 262; IV/650 S. 278）。
\end{transpara}

\syncleft
\begin{sourcepara}
\sourcern{14}2. Auch die Entstehungsgeschichte des Grundgesetzes spricht dagegen, aus Grundrechtsnormen eine Pflicht zum Strafen abzuleiten. Wo der Parlamentarische Rat Strafsanktionen von Verfassungs wegen für geboten erachtete, hat er dies ausdrücklich in das Grundgesetz aufgenommen: \abbrArt{} 26 \abbrAbs{} 1 für die Vorbereitung eines Angriffskrieges und \abbrArt{} 143 in der ursprünglichen Fassung für Hochverrat\ggfnArtTwoSixAndOneFourThreeDE{}.
\end{sourcepara}
\begin{transpara}
\transrn{14}2. 《基本法》的起源也反對從基本權利規範中推導出刑罰義務。如果議會理事會認為刑事制裁在憲法上是必要的，則明確將其納入《基本法》：第 26 條第1項針對準備侵略戰爭，以及原版第 143 條針對叛國罪\ggfnArtTwoSixAndOneFourThreeZH{}。
\end{transpara}

\syncleft
\begin{sourcepara}
\sourcern{15}Demgegenüber finden sich in den Materialien zu \abbrArt{} 2 \abbrAbs{} 2 \abbrGG{}\ggfnArtTwoTwoDE{}, wie auch die Mehrheit einräumt (C I 1 d), keine Anhaltspunkte für eine Verpflichtung, das ungeborene Leben strafrechtlich zu schützen. Eine nähere Analyse der Entstehungsgeschichte des Artikels spricht darüber hinaus dafür, daß die strafrechtliche Bewertung von Schwangerschaftsabbrüchen bewußt der eigenverantwortlichen Entscheidung des einfachen Gesetzgebers überlassen bleiben sollte. Die insoweit erheblichen Äußerungen der Abgeordneten Heuss und Greve und die Ablehnung des Antrages des Abgeordneten Seebohm (C I 1 d mit den dortigen Nachweisen) müssen vor ihrem zeitgeschichtlichen Hintergrund verstanden werden. In der Weimarer Zeit war die Bestrafung der Abtreibung außerordentlich umstritten; es handelte sich damals um ein ungleich schwerer wiegendes Problem, weil es die heute allgemein verbreiteten und leicht anwendbaren Mittel zur Empfängnisverhütung noch nicht gab. Dieser Zustand dauerte zur Zeit des Parlamentarischen Rates unverändert fort. Wenn unter diesen Umständen die beantragte Aufnahme einer ausdrücklichen Bestimmung über den Schutz keimenden Lebens abgelehnt worden ist, so ist das in Verbindung mit den erwähnten Äußerungen nur dahin zu verstehen, daß die Reform des umstrittenen § 218 \abbrStGB{} durch die Verfassung nicht präjudiziert werden sollte.
\end{sourcepara}
\begin{transpara}
\transrn{15}相較之下，在《基本法》第二條第二項\ggfnArtTwoTwoZH{}的資料中，正如多數意見所承認的那樣(C I 1 d)，並沒有證據表明刑法規定有保護未出生生命的義務。對條文淵源的進一步分析也表明，對終止妊娠的刑法評估應有意識地交由普通立法者獨立決定。赫斯議員和格雷夫議員的重要聲明以及對西博姆議員提案的否決（C I 1 d 及其證據）必須結合其當代歷史背景來理解。在魏瑪時期，終止妊娠的刑罰引起了極大的爭議。在當時，這是一個更嚴重的問題，因為當時廣泛使用且易於使用的避孕手段還不存在。這種情況在議會理事會召開時仍然沒有改變。如果在這種情況下，列入保護萌芽生命的明確條項的請求被拒絕，那麼結合上述陳述，只能被理解為憲法不應損害有爭議的刑法第218條的改革。
\end{transpara}

\syncleft
\begin{sourcepara}
\sourcern{16}Ein gegenteiliger Standpunkt läßt sich nicht damit begründen, daß die Aufnahme des \abbrArt{} 2 \abbrAbs{} 2 \abbrGG{}\ggfnArtTwoTwoDE{} unstreitig der Reaktion auf die unmenschliche Ideologie und Praxis der nationalsozialistischen Regimes entsprang (Vgl. aber C I 1 a und D IV des Urteils.). Diese Reaktion bezieht sich auf die Massenvernichtung menschlichen Lebens von Staats wegen in Konzentrationslagern und bei Geisteskranken, auf behördlich angeordnete Sterilisierungen und Zwangsabtreibungen, auf medizinische Versuche mit Menschen gegen deren Willen und auf die in zahllosen anderen staatlichen Maßnahmen zum Ausdruck kommende Mißachtung des individuellen Lebens und der Menschenwürde.
\end{sourcepara}
\begin{transpara}
\transrn{16}相反的立場不能以下列事實為理由，即《基本法》第2條第2項之納入毫無爭議地源自於對國家社會主義政權的不人道意識形態和實踐的反應（但請參閱判決書的 C I 1 a 和 D IV）。這種反應指的是國家在集中營和精神病患者中大規模毀滅人類生命，官方下令絕育和強迫終止妊娠，違背人們意願對人們進行醫學實驗，以及無數其他國家措施所表現出的對個人生命和人類尊嚴的漠視。
\end{transpara}

\syncleft
\begin{sourcepara}
\sourcern{17}Hieraus Schlußfolgerungen für die verfassungsrechtliche Bewertung einer nicht vom Staat, sondern von der Schwangeren selbst oder mit ihrem Willen von Dritten vorgenommenen Abtötung der Leibesfrucht zu ziehen, ist um so weniger am Platze, als das nationalsozialistische Regime entsprechend seiner biologisch-bevölkerungspolitischen Ideologie gerade dazu einen rigorosen Standpunkt eingenommen hatte. Neben neuen Vorschriften gegen die Werbung für Abtreibungen oder Abtreibungsmittel wurde durch entsprechende staatliche Maßnahmen darauf hingewirkt, im Gegensatz zur Praxis in der Weimarer Zeit eine striktere Anwendung der Strafbestimmungen durchzusetzen (Vgl. über den Anstieg der Verurteilungen im Dritten Reich: Dotzauer, Abtreibung, in: Handwörterbuch der Kriminologie [hrsg. von Sieverts], 2. \abbrAufl{}, \abbrBd{} I [1966], \abbrS{} 10 \abbrF{}). Diese an sich schon hohen Strafdrohungen wurden 1943 wesentlich verschärft. Während bisher sowohl für die Schwangere wie für den nicht gewerbsmäßig handelnden Helfer nur Gefängnisstrafe vorgesehen war, wurde nunmehr die Selbstabtreibung in besonders schweren Fällen mit Zuchthaus belegt. Die Fremdabtreibung war, abgesehen von minder schweren Fällen, stets mit Zuchthaus zu bestrafen; hatte der Täter "dadurch die Lebenskraft des deutschen Volkes fortgesetzt beeinträchtigt", sogar mit Todesstrafe. Angesichts dieser, bei der Entstehung des Grundgesetzes noch unveränderten, lediglich durch das alliierte Verbot grausamer oder übermäßig hoher Strafen in ihrer Anwendung gemilderten Bestimmungen können die Gründe, die zur Aufnahme des \abbrArt{} 2 \abbrAbs{} 2 \abbrGG{}\ggfnArtTwoTwoDE{} führten, schlechterdings nicht zugunsten einer verfassungsrechtlichen Pflicht zur Bestrafung von Abtreibungen herangezogen werden. Vielmehr gebietet die mit dem Grundgesetz vollzogene, entschiedene Abkehr vom totalitären nationalsozialistischen Staat eher umgekehrt Zurückhaltung im Umgang mit der Kriminalstrafe, deren verfehlter Gebrauch in der Geschichte der Menschheit schon unendlich viel Leid angerichtet hat.
\end{sourcepara}
\begin{transpara}
\transrn{17}由此推論出，對並非由國家、而是由懷孕婦女本人或第三人依其意願所為之殺害胎兒行為（Abtötung der Leibesfrucht），在憲法評價上應採相同結論，尤其不妥；因為國家社會主義政權正是依其生物學及人口政策意識形態，在此問題上採取嚴苛立場。除新增禁止墮胎或墮胎手段廣告之規定外，亦透過相應國家措施，與魏瑪時期實務相反，促成刑罰規定之更嚴格適用（關於第三帝國定罪數增加，參見 Dotzauer, Abtreibung, in: Handwörterbuch der Kriminologie [hrsg. von Sieverts], 2. Aufl., Bd. I [1966], S. 10 f.）。此等刑罰威嚇本已嚴厲，1943年又大幅加重。此前對懷孕婦女及非業務性協助者僅規定一般監禁（Gefängnisstrafe）；此後，自行墮胎於特別嚴重情形下改處重懲役刑（Zuchthaus）。他人墮胎（Fremdabtreibung）除較輕情形外，一律處重懲役刑；若行為人「因而持續損害德國人民生命力」，甚至處死刑。鑑於此等規定於《基本法》制定時仍未變更，僅因盟軍禁止殘酷或過度刑罰而在適用上有所緩和，促成《基本法》第2條第2項\ggfnArtTwoTwoZH{}納入之理由，絕不能被用來支持一項以刑罰制裁墮胎之憲法義務。毋寧，《基本法》所完成之對極權國家社會主義國家的堅決背離，反而要求對刑罰保持克制；刑罰之錯誤使用，在人類歷史中已造成無窮苦難。
\end{transpara}

\bisubsection{B.}{B.}
\syncleft
\begin{sourcepara}
\sourcern{18}Selbst wenn man entgegen unserer Auffassung mit der Mehrheit eine verfassungsrechtliche Pflicht zum Strafen für denkbar hält, kann dem Gesetzgeber hier kein Verfassungsverstoß zur Last gelegt werden. Die Mehrheitsbegründung begegnet - ohne daß es eines Eingehens auf jedes Detail bedarf - den folgenden Einwänden:
\end{sourcepara}
\begin{transpara}
\transrn{18}縱使有違我等見解而從多數意見，認刑罰之憲法義務在概念上可得設想，立法機關於此亦不得被歸責為違憲。多數意見之論理——無須逐一深入每一細節——面臨如下反對意見：
\end{transpara}

\bisubsubsection{I.}{I.}
\syncleft
\begin{sourcepara}
\sourcern{19}Auch nach Meinung der Mehrheit soll eine verfassungsrechtliche Pflicht zum Strafen nur als ultima ratio in Betracht kommen (C III 2 b). Macht man damit wirklich Ernst, so setzt eine solche Pflicht zunächst voraus, daß geeignete Mittel milderer Art fehlen oder ihr Einsatz sich als wirkungslos erwiesen hat; darüber hinaus muß die Strafsanktion geeignet und erforderlich sein, um das erstrebte Ziel überhaupt oder besser zu erreichen. Beides muß - folgt man der bisherigen Rechtsprechung (vgl. etwa \abbrBVerfGE{} 17, 306 [313 \abbrF{}]) - zweifelsfrei nachgewiesen werden. Denn wenn schon die Zulässigkeit einer bestehenden Strafvorschrift davon abhängt, ob sie zum Schutz des jeweiligen Rechtsgutes geeignet und erforderlich ist, dann bedarf es eines solchen Nachweises erst recht, wenn der Gesetzgeber sogar gegen seinen Willen zum Strafen gezwungen werden soll. Soweit es dafür auf die Beurteilung der Sachlage und der Effektivität beabsichtigter Maßnahmen ankommt, hat das Gericht die Auffassung des Gesetzgebers zugrunde zu legen, solange sie nicht als offensichtlich irrig widerlegt wird (vgl. \abbrBVerfGE{} 7, 377 [412]; 24, 367 [406]; 35, 148 - abw.M. - [165]).
\end{sourcepara}
\begin{transpara}
\transrn{19}多數意見也認為，憲法規定的刑罰義務只能被視為最後手段（C III 2 b）。如果一個人真的認真對待這一點，那麼這種義務首先就以缺乏合適的、較溫和的手段為先決條件，或者這些手段的使用已被證明是無效的；此外，刑罰性制裁必須是適當和必要的，以便完全或更好地實現預期目標。如果遵循先前的判例法（參見 BVerfGE 17, 306 [313 f.]），兩者都必須得到毫無疑問的證明。如果現有刑事條項的可採性取決於其是否適合和有必要保護各自的法益，那麼如果立法機關要被迫施以刑罰，甚至違背其意願，則更需要這樣的證據。如果對事實情況和預期措施的有效性的評估很重要，法院必須以立法機關的意見為基礎，只要立法機關的意見沒有被認為明顯錯誤而被駁斥（參見 BVerfGE 7, 377 [412]; 24, 367 [406]; 35, 148 - abw.M. - [165]）。
\end{transpara}

\syncleft
\begin{sourcepara}
\sourcern{20}Diesen Anforderungen genügt die Urteilsbegründung nicht. Sie verstrickt sich wiederholt in Widersprüche und kehrt am Ende die Beweislast geradezu um: Ein Verzicht auf Strafsanktionen soll dem Gesetzgeber nur noch erlaubt sein, wenn zweifelsfrei feststeht, daß die von ihm bevorzugten milderen Maßnahmen zur Erfüllung der Schutzpflicht "zumindest" ebenso wirksam oder gar wirksamer sind (D III, C III 3; vgl. auch D II 2 c).
\end{sourcepara}
\begin{transpara}
\transrn{20}判決理由不符合這些要求。它一再陷入矛盾之中，並最終顛倒了舉證責任：只有在明確無疑地表明立法機關傾向於履行保護義務的較溫和措施「至少」同樣有效或甚至更有效的情況下，才應允許立法機關放棄刑罰性制裁（D III，C III 3；另參見D II 2 c）。
\end{transpara}

\syncleft
\begin{sourcepara}
\sourcern{21}1. Die zunächst eindrucksvollen Ausführungen über den unbestrittenen hohen Rang des Lebensschutzes vernachlässigen die Singularität des Schwangerschaftsabbruchs im Verhältnis zu anderen Gefährdungen menschliches Lebens. Es handelt sich hier nicht um die akademische Frage, ob der Einsatz der staatlichen Strafgewalt zum Schutz vor Mördern und Totschlägern, auf die in keiner anderen Weise präventiv eingewirkt werden kann, unumgänglich ist. In der europäischen, kirchlich beeinflußten Rechtsgeschichte ist stets zwischen geborenem und ungeborenem Leben unterschieden worden. Auch die Wertentscheidung der Verfassung läßt bei der Auswahl der erforderlichen Schutzmaßnahmen Raum für eine solche Differenzierung, zumal da das Grundrecht des \abbrArt{} 2 \abbrAbs{} 2 \abbrGG{}\ggfnArtTwoTwoDE{} nicht - wie die Mehrheit formuliert (C II 1) - "umfassend" gewährleistet ist, sondern unter Gesetzesvorbehalt steht. Anderenfalls ließe sich weder die ethische noch die eugenische oder gar die soziale Indikation begründen.
\end{sourcepara}
\begin{transpara}
\transrn{21}1. 關於保護生命無可爭議的高度優先性的最初令人印象深刻的陳述忽視了終止妊娠相對於人類生命的其他威脅的特殊性。這不是一個學術問題，即使用國家刑事權力來防止殺人犯和過失殺人行為是否是不可避免的，而殺人犯和過失殺人行為是無法透過任何其他方式預防的。在歐洲法律史上，受教會的影響，出生的生命和未出生的生命一直是有區別的。憲法價值決定也為選擇必要的保護措施留有區分的空間，特別是《基本法》第二條第二項\ggfnArtTwoTwoZH{}的基本權利並非如多數意見所說（C II 1）那樣「全面」得到保障，而是受法律約束。否則，倫理學、優生學甚至社會學指標都無法被證明是合理的。
\end{transpara}

\syncleft
\begin{sourcepara}
\sourcern{22}Auch die Mehrheit bezweifelt nicht die Berechtigung dieser Unterscheidung (C III 2 a), trennt aber wiederum nicht zwischen den verschiedenen Aspekten der Grundrechtsnorm. Soweit es sich um die Abwehr staatlicher Eingriffe handelt, kann selbstverständlich nicht zwischen dem vor- und nachgeburtlichen Entwicklungsstadium unterschieden werden; der Embryo ist insoweit als potentieller Grundrechtsträger durchgängig in gleicher Weise zu schützen wie jedes geborene Menschenleben. Diese rechtliche Gleichbehandlung läßt sich schon auf die Verletzung ungeborenen Lebens durch Dritte gegen den Willen der Schwangeren nur begrenzt übertragen, keinesfalls aber auf die Weigerung der Frau, die Menschwerdung ihrer Leibesfrucht im eigenen Körper zuzulassen.
\end{sourcepara}
\begin{transpara}
\transrn{22}多數意見也不懷疑這種區分的合理性（C III 2 a），但同樣不區分基本權利規範的不同面向。就防禦國家干預而言，當然不可能區分出生前和出生後的發展階段；在這方面，胚胎作為基本權利的潛在承載者，必須像每個出生的人類生命一樣受到保護。這種法律上的平等待遇只能有限地適用於第三方違背懷孕婦女意願侵犯未出生生命的行為，而絕不能適用於懷孕婦女拒絕讓子宮投生到自己體內的情況。
\end{transpara}

\syncleft
\begin{sourcepara}
\sourcern{23}Die Besonderheit, daß in der Person der Schwangeren eine singuläre Einheit von "Täter" und "Opfer" vorliegt (So auch die Mehrheit unter C II 2, C III 3), fällt rechtlich bereits deswegen ins Gewicht, weil der Schwangeren - anders als dem Adressaten von Strafvorschriften gegen Tötungsdelikte - weit mehr abverlangt wird als nur ein Unterlassen: Sie soll nicht nur die mit dem Austragen der Leibesfrucht verbundenen tiefgreifenden Veränderungen ihrer Gesundheit und ihres Wohlbefindens dulden, sondern auch die Eingriffe in ihre Lebensgestaltung hinnehmen, die sich aus Schwangerschaft und Geburt ergeben, besonders die mütterliche Verantwortung für die weitere Entwicklung des Kindes nach der Geburt tragen. Anders als bei den genannten Tötungsdelikten kann und muß der Gesetzgeber ferner davon ausgehen, daß das Schutzobjekt - die Leibesfrucht - am wirksamsten durch die Mutter selbst geschützt wird und daß deren Bereitschaft zum Austragen der Leibesfrucht durch Maßnahmen verschiedenster Art gestärkt werden kann. Da es von Natur aus keiner Strafvorschrift bedarf, um die mütterliche Schutzbeziehung herzustellen und zu sichern, läßt sich schon deswegen fragen, ob einer Störung dieser Beziehung, wie sie bei Schwangerschaftsabbrüchen zutage tritt, gerade durch eine Strafsanktion in geeigneter Weise begegnet werden kann. Jedenfalls darf der Gesetzgeber darauf wegen der genannten Besonderheiten anders reagieren als auf die Tötung menschlichen Lebens durch Dritte.
\end{sourcepara}
\begin{transpara}
\transrn{23}懷孕婦女身上存在著「犯罪者」和「受害者」的單一統一體（這也是 C II 2、C III 3 中的大多數）的特殊性，這一點具有重要的法律意義，因為懷孕婦女與針對殺人罪的刑事條項的受害人不同，需要做的不僅僅是克制：她不僅應該容忍與懷胎相關的健康和福祉的深刻變化，而且還應該接受懷孕和乾預而對其生活方式造成的深刻變化，而且還應該接受懷孕的變化，而且還應該接受懷孕和懷孕而對其生活方式造成的深刻變化，而且還應該接受懷孕和乾預而對其生活方式造成的深刻變化，而且還應該接受懷孕的變化，而且還應該接受懷孕和懷孕而對其生活方式造成的深刻變化，而且還應該接受懷孕和乾預而對其生活方式造成的深刻變化，而且還應該接受懷孕和懷孕而對其生活方式的變化出生，特別是對孩子出生後的進一步發展承擔母親的責任。與上述殺人罪相反，立法機關可以而且必須假定，保護對象──胎兒──是由母親本人最有效地保護的，並且可以透過多種措施增強母親生育胎兒的意願。由於刑事法規並沒有固有的必要性來建立和保障孕產婦保護關係，因此可以提出這樣的問題：是否可以透過刑事制裁來充分解決終止妊娠中發生的這種關係的破壞。無論如何，由於所提及的特殊性，立法機關的反應可能與第三人殺害人命的反應不同。
\end{transpara}

\syncleft
\begin{sourcepara}
\sourcern{24}Nach Auffassung der unterzeichnenden Richterin ist die Weigerung der Schwangeren, die Menschwerdung ihrer Leibesfrucht in ihrem Körper zuzulassen, nicht allein nach dem natürlichen Empfindungen der Frau, sondern auch rechtlich etwas wesentlich anderes als die Vernichtung selbständig existenten Lebens. Schon deswegen verbietet es sich von vornherein, die Abtreibung im ersten Stadium der Schwangerschaft mit Mord oder vorsätzlicher Tötung prinzipiell gleichzustellen. Erst recht ist es verfehlt, wenn nicht unsachlich, die Fristenlösung in die Nähe der Euthanasie oder gar der "Tötung unwerten Lebens" zu rücken, um sie von daher zu diskriminieren - wie dies in der öffentlichen Diskussion geschehen ist. Der Umstand, daß erst in einem längeren Entwicklungsprozeß ein vom mütterlichen Organismus trennbares selbständig existentes Lebewesen entsteht, legt es vielmehr nahe oder läßt es wenigsten zu, bei der rechtlichen Beurteilung zeitliche, dieser Entwicklung entsprechende Zäsuren zu berücksichtigen (Vgl. das Erkenntnis des Österreichischen Verfassungsgerichtshofs unter II 5 b der Entscheidungsgründe, \abbrEuGRZ{} 1975, \abbrS{} 74 [80]; Lay, JZ 1970, \abbrS{} 465 \abbrFF{}; noch weiter gehend bezeichnet es Roman Herzog [JR 1969, \abbrS{} 441] als Sache des Gesetzgebers, "von welcher Entwicklungsphase an [der] staatliche Schutz für das werdende Leben wirksam werden soll".). Die biologische Kontinuität der Gesamtentwicklung bis zur Geburt (Siehe C I 1 b) - deren Beginn bei konsequenter Anwendung der Mehrheitsauffassung nicht erst bei der Einnistung, sondern bei der Empfängnis anzusetzen wäre - ändert nichts daran, daß den verschiedenen Entwicklungsstufen des Embryos eine Veränderung in der Einstellung der Schwangeren im Sinne einer wachsenden mütterlichen Bindung entspricht. Demgemäß ist es für das Rechtsbewußtsein der Schwangeren, aber auch für das allgemeine Rechtsbewußtsein nicht das gleiche, ob ein Schwangerschaftsabbruch im ersten Stadium der Schwangerschaft oder in einer späteren Phase stattfindet. Dies hat zu allen Zeiten in in- und ausländischen Rechtsordnungen seinen Niederschlag in einer an solche zeitlichen Einschnitte anknüpfenden verschiedenen strafrechtlichen Bewertung der Abtreibung gefunden, wie etwa der Supreme Court eindrucksvoll dargestellt hat (410 \abbrUS{} 113 [132 \abbrFF{}, 160 \abbrF{}]). Für den deutschen Rechtsraum verdient Hervorhebung, daß das Kirchenrecht, gestützt auf die Beseelungslehre, bis zum Ende des 19. Jahrhunderts die Abtreibung in die Zeitspanne bis zum 80. Tag nach der Empfängnis als straflos angesehen hat; auch das weltliche Strafrecht sah bis zum Erlaß des Strafgesetzbuches von 1871 zeitliche Abstufungen für die Höhe der Strafdrohung vor (Vgl. Dähn in: Das Abtreibungsverbot des § 218 \abbrStGB{} [hrsg. von Baumann], 2. \abbrAufl{}, 1972, \abbrS{} 331 \abbrF{}; Supreme Court, \abbrAAO{}, 134; Simson-Geerds, Straftaten gegen die Person und Sittlichkeitsdelikte in rechtsvergleichender Sicht, 1969, \abbrS{} 87; Sonderausschuß für die Strafrechtsreform, 7. \abbrWp{}, Anlage zur 15. Sitzung, \abbrStenBer{} \abbrS{} 690 \abbrFF{}, 697 \abbrF{}).
\end{sourcepara}
\begin{transpara}
\transrn{24}下列署名法官認為，懷孕婦女拒絕讓胎兒投生在她的體內，不僅符合婦女的自然感受，而且從法律上講，這與破壞獨立存在的生命有著本質上的不同。僅出於這個原因，從一開始就禁止將懷孕第一階段的終止妊娠等同於謀殺或故意殺人。如果不是不客觀的話，將期限規定接近安樂死甚至「殺死不值得的生命」以歧視他們——正如公眾討論中所發生的那樣，這是特別錯誤的。一個可以與母體有機體分離的獨立存在的生物體只有在較長的發育時期才能形成，這一事實表明，或至少允許，在法律評估中應考慮到與這種發育相對應的時間轉折點（參見奧地利憲法法院根據該決定理由的 II 5 b 的結論，EuGRZ 1975，第 74 頁 [80] 74 [80]； [JR 1969，p. 441]將其描述為立法機關的問題「國家對正在發展中的生命的保護應該從哪個階段開始生效。」）。直到出生為止的整體發育的生物學連續性（見 C I 1 b）——如果多數意見的觀點一致，則不應從植入開始，而是從受孕開始——並沒有改變這樣一個事實，即胚胎發育的各個階段對應於懷孕婦女在母性紐帶不斷增長的意義上的態度變化。因此，無論是懷孕初期或懷孕後期終止妊娠，無論是懷孕婦女的法律意識或一般法律意識都是不同的。正如美國最高法院令人印象深刻地表明的那樣，這一點始終反映在國內外法律體系中與此類時間變化相關的各種終止妊娠刑法評估中（410 U.S. 113 [132 ff., 160 f.]）。對於德國法律領域來說，值得強調的是，基於靈魂教義的教會法認為，直到19世紀末，受孕後第80天以內的終止妊娠不會受處罰；在 1871 年《刑法典》頒布之前，世俗刑法還規定了刑罰程度的時間分級（參見 Dähn，參見：The abortion ban of § 218 StGB [Baumann 編）頁；美國最高法院，同上，134；Simson-Geerds，《相對法律中針對個人的犯罪和道德犯罪》觀點，1969 年，第 87 頁；刑法改革特別委員會，第 7 屆會議附錄，第 690 頁，第 697 頁）。
\end{transpara}

\syncleft
\begin{sourcepara}
\sourcern{25}Der unterzeichnende Richter neigt dazu, diesen weiteren Überlegungen zum Verhältnis zwischen der Schwangeren und ihrer Leibesfrucht rechtlich eine geringere Bedeutung beizumessen. Wenn jedoch die Rücknahme der Strafdrohung in den ersten drei Monaten der Schwangerschaft aus anderen - bereits genannten oder noch zu erörternden - Gründen verfassungsrechtlich nicht zu beanstanden ist, dann handelt der Gesetzgeber jedenfalls nicht sachfremd, wenn er den genannten Umständen bei seiner Regelung Rechnung trägt.
\end{sourcepara}
\begin{transpara}
\transrn{25}署名法官傾向於認為，關於懷孕婦女與其胎兒間關係之上述進一步考量，在法律上重要性較低。然而，若基於其他已述或待述理由，撤回妊娠最初三個月內之刑罰威嚇在憲法上並無可指摘，則立法機關於其規制中考量上述情事，無論如何並非事理無關。
\end{transpara}

\syncleft
\begin{sourcepara}
\sourcern{26}2. Die Prüfung, ob trotz der genannten Besonderheiten auch zum Schutz ungeborenen Lebens eine Pflicht zum Strafen als ultima ratio zu fordern ist, muß von dem sozialen Problem ausgehen, das den Gesetzgeber zu seiner Regelung veranlaßt hat. In der Mehrheitsbegründung finden sich lediglich knappe Hinweise auf die Komplexität dieses Problems und - im Zusammenhang mit der Indikationenregelung - einige Erörterungen über die sozialen Ursachen für Abtreibungen (C, C III 2, C III 3, D II, D II 3); insgesamt unterbleibt jedoch wegen der mehr dogmatischen Betrachtungsweise eine hinreichende Würdigung der vom Gesetzgeber vorgefundenen Verhältnisse und der daraus folgenden Schwierigkeiten für die allseits als notwendig anerkannte Reform.
\end{sourcepara}
\begin{transpara}
\transrn{26}2. 儘管有上述的特殊性，是否應將刑罰義務作為保護未出生生命的最後手段進行審查，必須基於促使立法機關對其進行規範的社會問題。多數聲明僅簡要提及該問題的複雜性，並結合指標事由規定，對終止妊娠的社會原因進行了一些討論（C、C III 2、C III 3、D II、D II 3）；但總體而言，由於比較教條的做法，沒有對立法機關發現的條件以及由此產生的普遍認為必要的改革的困難進行充分的評估。
\end{transpara}

\syncleft
\begin{sourcepara}
\sourcern{27}a) Diese Verhältnisse sind in erster Linie geprägt durch die enorme Dunkelziffer, die nicht damit bagatellisiert werden kann, daß - verständlicherweise - keine sicheren Daten zu ermitteln sind. Nach den Berichten des Sonderausschusses für die Strafrechtsreform ist " nach ernst zu nehmenden Untersuchungen" von 75.000 bis 300.000 illegalen Abtreibungen jährlich auszugehen (Vgl. \abbrBTDrucks{} 7/1981 (neu) \abbrS{}6; 7/1982, \abbrS{} 5; 7/1983 \abbrS{} 5, je mit weiteren Nachweisen); innerhalb dieses Bereichs bewegen sich auch die von den Sachverständigen bei der öffentlichen Anhörung vor dem Sonderausschuß genannten Zahlen (Vgl. Sonderausschuß für die Strafrechtsreform, 6. \abbrWp{}, 74., 75. und 76. Sitzung, \abbrStenBer{} \abbrS{} 2173, 2218, 2241.). Noch bis vor kurzem, d.h. vor Beginn der Diskussion im Parlament wurden allgemein weit höhere Zahlen angenommen (Vgl. die Nachweise in der Antwort des Bundesministers der Justiz auf eine Kleine Anfrage der Fraktion der CDU/CSU, \abbrBTDrucks{} VI/2025 \abbrS{} 3; s. auch E.-W. Böckenförde, der 200.000 bis 400.000 illegale Aborte zugrunde legt (Stimmen der Zeit, \abbrBd{} 188 [1971], \abbrS{} 147 [152]).).
\end{sourcepara}
\begin{transpara}
\transrn{27}a) 此等情況首先以極高的黑數為特徵；不能因為可以理解地無法取得確切資料，就將此一情況輕描淡寫。依刑法改革特別委員會之報告，根據「值得認真看待之研究」，每年非法墮胎數量應估計為75,000至300,000件（參見 BTDrucks. 7/1981 [neu] S. 6; 7/1982, S. 5; 7/1983, S. 5，均附進一步資料）。特別委員會公開聽證中專家所提出之數字，也落在此一範圍內（參見 Sonderausschuß für die Strafrechtsreform, 6. Wp., 74., 75. und 76. Sitzung, StenBer. S. 2173, 2218, 2241）。直到不久以前，即議會討論開始前，一般甚至假定更高之數字（參見聯邦司法部長答覆 CDU/CSU 黨團小質詢之資料，BTDrucks. VI/2025 S. 3；另見 E.-W. Böckenförde，其以200,000至400,000件非法墮胎為基礎，Stimmen der Zeit, Bd. 188 [1971], S. 147 [152]）。
\end{transpara}

\syncleft
\begin{sourcepara}
\sourcern{28}Auch wenn man nur die niedrigsten Schätzungen zugrunde legt, bleibt die Zahl erschreckend hoch. Ihr steht eine verschwindend geringe Zahl behördlich bekanntgewordener Abtreibungsfälle und gerichtlicher Verurteilungen gegenüber: für 1971 bei 584 bekanntgewordenen Straftaten 184 Verurteilungen, für 1972 bei 476 Fällen 154 Verurteilungen (Statistisches Jahrbuch für die Bundesrepublik Deutschland, 1973, \abbrS{} 117, 121; 1974, \abbrS{} 116, 121.). Ganz überwiegend wird dabei nur auf Geldstrafen erkannt; ausnahmsweise verhängte kurzfristige Freiheitsstrafen werden meist zur Bewährung ausgesetzt (Vgl. Statistisches Bundesamt, Fachserie A "Bevölkerung und Kultur", Reihe 9 "Rechtspflege", 1972, \abbrS{} 100 \abbrF{}, 144 \abbrF{}, 160 \abbrF{}). Die hierin erkennbare Nichtbeachtung der Strafvorschrift bedeutet nicht allein eine Entwertung des werdenden Lebens, sondern wirkt sich auch korrumpierend für die allgemeine Rechtsgeltung aus, zumal da unter diesen Umständen die strafgerichtliche Verfolgung zum reinen Zufall wird.
\end{sourcepara}
\begin{transpara}
\transrn{28}即使僅以最低估計為基礎，數字仍高得驚人。與此相對，為主管機關所知之墮胎案件及法院定罪數量微乎其微：1971年，在584件已知犯罪中有184件定罪；1972年，在476件案件中有154件定罪（Statistisches Jahrbuch für die Bundesrepublik Deutschland, 1973, S. 117, 121; 1974, S. 116, 121）。其中絕大多數僅判處罰金；例外判處之短期自由刑通常宣告緩刑（參見 Statistisches Bundesamt, Fachserie A "Bevölkerung und Kultur", Reihe 9 "Rechtspflege", 1972, S. 100 f., 144 f., 160 f.）。此處所顯示之刑罰條項不受遵守，不僅意味著正在發展中生命之貶值，也腐蝕一般法律效力；尤其在此等情況下，刑事追訴變成純然偶然。
\end{transpara}

\syncleft
\begin{sourcepara}
\sourcern{29}Den Gesetzgeber konnte es ferner nicht gleichgültig lassen, daß illegale Schwangerschaftsabbrüche auch heute noch zu Gesundheitsschäden führen, und zwar nicht nur bei Abtreibungen durch "Kurpfuscher" und "Engelmacherinnen", sondern in größerem Umfang auch deswegen, weil bei ärztlichen Eingriffen die Illegalität den vollen Einsatz des modernen Instrumentariums und des erforderlichen Hilfspersonals beeinträchtigt oder eine notwendige Nachbehandlung verhindert. Als Übelstand erscheint weiter die kommerzielle Ausbeutung abtreibungswilliger Frauen im In- und Ausland und die damit verbundene soziale Ungleichheit; bessergestellte Frauen können, besonders durch Reisen ins benachbarte Ausland, weit leichter eine Abtreibung durch einen Arzt erreichen als ärmere oder weniger gewandte. Schließlich kommt die sogenannte Folgekriminalität hinzu; so steht die Erpressung mit der Kenntnis eines illegalen Aborts in der Reihe der Erpressungsmittel an dritter Stelle (Geerds, Erpressung, in Handwörterbuch der Kriminologie, \abbrAAO{}, \abbrS{} 182).
\end{sourcepara}
\begin{transpara}
\transrn{29}此外，立法機關亦不能對下列事實無動於衷：非法終止妊娠至今仍會造成健康損害，且不僅發生於由「庸醫」與「墮胎婆」實施之墮胎；更大範圍內，還因醫療介入之非法性妨礙現代器械及必要輔助人員之充分投入，或阻礙必要後續治療。尋求墮胎之婦女在國內外遭商業剝削，以及與此相連之社會不平等，亦屬弊端；處境較佳之婦女，尤其透過前往鄰近外國，遠比貧困或較不熟練者更容易取得醫師實施之墮胎。最後，尚有所謂後續犯罪；例如以知悉非法墮胎為手段之勒索，在勒索手段序列中位居第三（Geerds, Erpressung, in Handwörterbuch der Kriminologie, a.a.O., S. 182）。
\end{transpara}

\syncleft
\begin{sourcepara}
\sourcern{30}b) Für die Entscheidung des Gesetzgebers, wie diese Zustände am besten zu reformieren sind, war besonders bedeutsam, daß der Entschluß zum Schwangerschaftsabbruch regelmäßig aus einer Konfliktsituation erwächst, die auf mannigfaltigen, stark von den Verhältnissen des Einzelfalles geprägten Motivationen beruht. Neben wirtschaftlichen oder materiellen Gründen - etwa unzulängliche Wohnverhältnisse, nicht ausreichendes oder ungesichertes Einkommen der vielleicht schon vielköpfigen Familie, Notwendigkeit einer Erwerbstätigkeit beider Ehegatten - stehen persönliche Gründe: die immer noch nicht beseitigte gesellschaftliche Diskriminierung der unverheirateten Mutter, der Druck des Erzeugers oder der Familie, die Furcht vor einer Gefährdung der Partnerbeziehung oder einem Zerwürfnis mit den Eltern, der Wunsch oder die Notwendigkeit, eine begonnene Ausbildung fortzusetzen oder einen Beruf weiter auszuüben, Schwierigkeiten in der Ehe, das Gefühl, der Betreuung weiterer Kinder physisch oder psychisch nicht mehr gewachsen zu sein, bei Alleinstehenden auch die Ablehnung einer nicht zu verantwortenden Heimerziehung des Kindes. Die Angst der Schwangeren, daß die ungewollte Schwangerschaft zu einem irreparablen Einbruch in ihre persönliche Lebensgestaltung oder den Lebensstandard der Familie führen wird, das Empfinden, daß sie beim Austragen ihrer Leibesfrucht nicht auf eine wirksame Hilfe der Umwelt rechnen kann, sondern daß die nachteiligen Folgen eines nicht nur von ihr zu verantwortenden Verhaltens sie allein treffen, lassen ihr die Schwangerschaftsunterbrechung häufig als einzigen Ausweg erscheinen. Selbst wo nach der persönlichen Situation nicht-einsichtige Motivationen der Bequemlichkeit, des Egoismus, besonders des Konsumstrebens im Vordergrund stehen, kann dies nicht ausschließlich der Frau zur Last gelegt werden, sondern spiegelt zugleich die materialistische, weitgehend kinderfeindliche Gesamteinstellung der "Wohlstandsgesellschaft" wider. Auch haben Staat und Gesellschaft bisher noch keine hinreichenden Einrichtungen und Lebensformen entwickelt, die es der Frau ermöglichen, Mutterschaft und Familienleben mit einer chancengleichen persönlichen Entfaltung, besonders auf beruflichem Gebiet, zu verbinden (Vgl. zu alledem etwa das Memorandum des Bensberger Kreises zur Reform des § 218 \abbrStGB{} (Publik-Forum Sonderdruck) unter I 1 sowie die Ausführungen der Sachverständigen und des Regierungsvertreters vor dem Sonderausschuß für die Strafrechtsreform, 6. \abbrWp{}, 74., 75. und 76. Sitzung, \abbrStenBer{} \abbrS{} 2219 mit der tabellarischen Übersicht in Anlage 3 \abbrS{} 2368 (Rolinski); \abbrS{} 2233 (Dotzauer); \abbrS{} 2251 \abbrFF{} (Pross); 7. \abbrWp{}, 23. Sitzung, \abbrStenBer{} \abbrS{} 1390 \abbrF{}; s. ferner die Berichte des Sonderausschusses für die Strafrechtsreform, \abbrBTDrucks{} 7/1981 (neu) \abbrS{} 7; 7/1982 \abbrS{} 7; 7/1983 \abbrS{} 7; 7/1984 (neu) \abbrS{} 5.).
\end{sourcepara}
\begin{transpara}
\transrn{30}b) 對於立法機關如何最好地改革這些條件的決定來說，特別重要的是，終止妊娠的決定通常是在衝突情況下做出的，這種情況是基於各種動機，而這些動機受到個案情況的強烈影響。除了經濟或物質原因——例如生活條件不足、家庭收入不足或不穩定（可能已經有大量成員）、配偶雙方都需要工作——還有個人原因：對未婚母親的社會歧視尚未消除、來自父親或家庭的壓力、擔心伴侶關係受到威脅或與父母不和、希望或需要繼續培訓或繼續從事職業、婚姻中的困難等。心理上有能力照顧額外的孩子，對於單身人士來說，拒絕在家中撫養孩子，因為他們不負責撫養孩子。懷孕婦女擔心意外懷孕會對她的個人生活方式或家庭生活水平造成不可挽回的破壞，感覺自己不能依靠環境的有效幫助來懷上胎兒，而且她不是唯一的行為所產生的不良後果會影響到她一個人，這常常使終止妊娠在她看來是唯一的出路。即使基於個人情況，便利、利己主義、特別是消費慾望等違反直覺的動機處於顯著位置，這也不能僅歸咎於女性，同時也反映了「富裕社會」物質至上、很大程度上敵視兒童的整體態度。國家和社會尚未建立足夠的機構和生活方式，使婦女能夠將母性和家庭生活與個人發展的平等機會結合起來，特別是在職業領域（以上種種，例如參見本斯伯格圈（Bensberger Kreis）關於刑法第218條改革之備忘錄（《公共論壇》特別印刷品）第 I.1 節，以及專家和政府代表在刑法改革特別委員會之陳述，第六屆第74、75及76次會議，會議記錄第2219頁，含附件3第2368頁之表格總覽（Rolinski）；第2233頁（Dotzauer）；第2251頁及以下（Pross）；第七屆第23次會議，會議記錄第1390頁及以下；另參見刑法改革特別委員會之報告，BTDrucks. 7/1981 (neu) 第7頁；7/1982 第7頁；7/1983 第7頁；7/1984 (neu) 第5頁。）
\end{transpara}

\syncleft
\begin{sourcepara}
\sourcern{31}3. In dieser Gesamtsituation ist die "Eindämmung der Abtreibungsseuche" nicht nur ein "gesellschaftspolitisch erwünschtes Ziel" (D II 2 b), sondern gerade auch im Sinne eines besseren Lebensschutzes und zur Wiederherstellung der Glaubwürdigkeit der Rechtsordnung dringend geboten. Im Bemühen um eine Lösung dieser höchst schwierigen Aufgabe hat der Gesetzgeber alle wesentlichen Gesichtspunkte erschöpfend gewürdigt. Die Reform des § 218 \abbrStGB{} hatte bereits die in dieser Frage tief gespaltene Öffentlichkeit nachhaltig beschäftigt. Die parlamentarischen Beratungen sind auf diesem Hintergrund mit großem Ernst und ungewöhnlicher Gründlichkeit durchgeführt worden. Hierbei wurden die Wertentscheidungen der Verfassung ausdrücklich herangezogen; namentlich bestand Einigkeit über die staatliche Pflicht zum Schutz des ungeborenen Lebens. Bei der Ermittlung der maßgeblichen Faktoren und Argumente für eine sachgerechte Entscheidung entsprach das Vorgehen der Gesetzgebungsorgane ganz dem, was das KPD-Urteil als charakteristisch für eine legitime Willensbildung im freiheitlich-demokratischen Staat ansieht (\abbrBVerfGE{} 5, 85 [135, 197 \abbrF{}]).
\end{sourcepara}
\begin{transpara}
\transrn{31}3. 在這一總體情況下，「遏制終止妊娠流行」不僅是「社會政治期望目標」（D II 2 b），而且從更好地保護生命和恢復法律體系可信度的意義上來說也是迫切需要的。為了找到解決這一極其艱鉅任務的方法，立法機關對所有重要方面進行了詳盡的評估。刑法第218條的改革已經對公眾產生了持久的影響，公眾對此問題有深刻的分歧。在此背景下，議會審議工作異常嚴肅、異常徹底。這裡明確使用了憲法的價值決定；特別是，各方就國家保護未出生生命的責任達成了一致。在確定適當決定的相關因素和論點時，立法機構的方法完全符合 KPD 判決所認為的自由民主國家合法決策的特徵 (BVerfGE 5, 85 [135, 197 f.])。
\end{transpara}

\syncleft
\begin{sourcepara}
\sourcern{32}Bei der gewählten Lösung konnte der Gesetzgeber davon ausgehen, daß angesichts des Versagens der Strafsanktion die geeigneten Mittel zur Abhilfe im sozialen und gesellschaftlichen Bereich zu suchen sind, daß es darauf ankommt, einerseits durch präventive psychologische, sozial- und gesellschaftspolitische Förderungsmaßnahmen der Mutter das Austragen des Kindes zu erleichtern und ihre eigene Bereitschaft dazu zu stärken, andererseits durch bessere Information über die Möglichkeiten der Empfängnisverhütung die Zahl ungewollter Schwangerschaften zu vermindern. Auch die Mehrheit bezweifelt offenbar nicht (C III 1, D II, D III), daß solche Maßnahmen insgesamt gesehen am wirksamsten sind und am ehesten einer Effektuierung der Grundrechte im Sinne größerer Freiheit und vermehrter sozialer Gerechtigkeit entsprechen.
\end{sourcepara}
\begin{transpara}
\transrn{32}透過所選的解決方案，立法機關能夠認為，鑑於刑事制裁的失敗，應在社會和社會領域尋求適當的補救手段，重要的是，一方面，透過預防性心理、社會和社會政治支持措施，使母親更容易懷上孩子，並增強她自己這樣做的意願；另一方面，透過更好地了解懷孕的選擇，減少意外懷孕的數量。避孕。多數意見顯然並不懷疑（C III 1，D II，D III），這些措施總體上是最有效的，並且最接近實現更大自由和增強社會正義意義上的基本權利。
\end{transpara}

\syncleft
\begin{sourcepara}
\sourcern{33}Förderungsmaßnahmen dieser Art können in ein Strafgesetz schon wegen der verschiedenen staatlichen Kompetenzen begreiflicherweise nur begrenzt Eingang finden. Das Fünfte Strafrechtsreformgesetz enthält daher insoweit lediglich eine Pflicht zur Beratung. Nach der Konzeption des Gesetzgebers soll dadurch die Schwangere - ohne Furcht vor Strafe - aus ihrer Isolierung herausgeholt, die Bewältigung ihrer Schwierigkeiten soll ihr durch offene Kontakte mit ihrer Umwelt und durch eine auf ihre persönliche Konfliktsituation abgestellte individuelle Beratung erleichtert werden. Daß die vorgeschriebene Beratung dem Schutz werdenden Lebens dienen soll, indem sie die Bereitschaft zum Austragen der Leibesfrucht weckt und stärkt, wo dem nicht schwerwiegende Gründe entgegenstehen, ergibt sich bereits aus den in der Urteilsbegründung zitierten Gesetzesmaterialien und der dort erwähnten Entschließung der Bundestagsmehrheit (A I 6 d, D II, D II 3, D II 3 b; s. weiter die Ausführungen von Regierungsvertretern und Abgeordneten in der 2. und 3. Beratung [Abg. de With, Frau Funcke, Frau Eilers, \abbrDr{} Eppler, Scheu, Bundesminister Frau \abbrDr{} Focke und Bundeskanzler Brandt] 95. Sitzung, \abbrStenBer{} \abbrS{} 6383 [D]; 6384 [A]; 6391 [A]; 6402 [B]; 96. Sitzung, \abbrStenBer{} \abbrS{} 6471 [B]; 6482 [B]; 6499; 6500 [B].).
\end{sourcepara}
\begin{transpara}
\transrn{33}此類促進措施，基於國家權限之不同，固然只能有限地納入刑法。因此，《第五刑法改革法》就此僅包含諮詢義務。依立法者構想，懷孕婦女應在無須畏懼刑罰下脫離孤立；其困難之克服，應透過與其環境之開放接觸，以及針對其個人衝突處境之個別諮詢而較為容易。所規定之諮詢應服務於正在形成之生命保護，亦即在無重大相反理由時，喚起並強化將胎兒懷至足月之意願；此由判決理由所引法律材料及其中提及之聯邦議院多數決議已可得知（A I 6 d, D II, D II 3, D II 3 b；另參政府代表及議員於第二、三讀之發言：[de With 議員、Funcke 女士、Eilers 女士、Dr. Eppler、Scheu、聯邦部長 Dr. Focke 女士及聯邦總理 Brandt]，95. Sitzung, StenBer. S. 6383 [D]; 6384 [A]; 6391 [A]; 6402 [B]; 96. Sitzung, StenBer. S. 6471 [B]; 6482 [B]; 6499; 6500 [B]）。
\end{transpara}

\syncleft
\begin{sourcepara}
\sourcern{34}Wir bestreiten nicht, daß diese Beratungsregelung - wie im Urteil dargelegt (D II 3) - noch Schwächen aufweist. Soweit diese sich nicht durch eine verfassungskonforme Auslegung des Gesetzes und entsprechende Durchführungsregelungen der Länder hätten ausräumen lassen, wäre aber eine verfassungsrechtliche Beanstandung allein auf diese Mängel zu beschränken gewesen und durfte nicht die Fristen- und Beratungsregelung in ihrer Gesamtkonzeption in Frage stellen. Nicht zuletzt hängt der Erfolg einer Beratungsregelung wesentlich davon ab, ob der beratenen Frau Hilfen angeboten oder vermittelt werden können, die ihr Auswege aus ihren Schwierigkeiten eröffnen. Fehlt es daran, dann ist auch das Strafrecht nichts anderes als ein Alibi für das Defizit wirksamer Hilfen; Verantwortung und Lasten werden allein auf die schwächsten Glieder der Gesellschaft abgewälzt. Die Mehrheit erklärt sich - in Übereinstimmung mit der bisherigen Rechtsprechung - außerstande, insoweit die Gestaltungsfreiheit des Gesetzgebers zu beschränken und ihm eine Erweiterung sozial-präventiver Maßnahmen vorzuschreiben (C III 1). Wenn aber dafür richterliche Selbstbeschränkung gilt, dann darf das Verfassungsgericht den Gesetzgeber erst recht nicht zum Einsatz der Strafgewalt als des stärksten staatlichen Zwangsmittels zwingen, um soziale Pflichtversäumnis (Vgl. dazu Rudolphi, Straftaten gegen das werdende Leben, \abbrZStrW{} 83 [1971], \abbrS{} 105 [114 \abbrF{}, 128 \abbrF{}, 134].) durch Strafdrohung zu kompensieren. Dies entspricht gewiß nicht der Funktion des Strafrechts in einem freiheitlichen Sozialstaat.
\end{sourcepara}
\begin{transpara}
\transrn{34}我們並不否認，如判決所述（D II 3），此項諮詢規定仍有弱點。惟若此等弱點不能透過對法律作合憲解釋及各邦相應施行規制加以排除，憲法上之指摘亦應僅限於此等瑕疵，而不應質疑期限與諮詢規制之整體構想。尤其，諮詢規制之成功，在重要程度上取決於能否向受諮詢之婦女提供或轉介協助，使其得以擺脫困難。若欠缺此等協助，刑法即不過是有效協助不足之藉口；責任與負擔將完全轉嫁於社會最弱勢成員。多數意見表示，其與既有判例一致，無法就此限制立法者形成自由，並命令其擴張社會預防措施（C III 1）。然而，若司法自我節制適用於此，憲法法院更不得強迫立法者運用刑罰權此一最強國家強制手段，以刑罰威嚇補償社會義務之怠忽（參見 Rudolphi, Straftaten gegen das werdende Leben, ZStrW 83 [1971], S. 105 [114 f., 128 f., 134]）。這顯然不符合自由社會國中刑法之功能。
\end{transpara}

\syncleft
\begin{sourcepara}
\sourcern{35}4. Auch die Mehrheit erkennt die gesetzgeberische Absicht, durch Beratung Leben zu erhalten, als "achtenswertes Ziel" an (D II 2 b, D II 1), hält aber - in Übereinstimmung mit den Antragstellern - die Anordnung flankierender Strafsanktionen für unabdingbar, weil der durchgängige Verzicht auf Bestrafung in den Fällen, in denen der Schwangerschaftsabbruch auf keinerlei achtenswerten Gründen beruhe, eine "Schutzlücke" hinterlasse (A II 2 c, C III 2 b, D II 2).
\end{sourcepara}
\begin{transpara}
\transrn{35}4. 多數意見也承認透過諮詢保護生命的立法意圖是一個「值得尊敬的目標」（D II 2 b、D II 1），但與聲請人一致認為，下令附帶刑事制裁至關重要，因為在並非基於任何值得尊敬的理由的情況下完全免除刑罰會留下「保護缺口」（A II 2 c、C III 2 b、D II 2）。
\end{transpara}

\syncleft
\begin{sourcepara}
\sourcern{36}a) Die Eignung von Strafsanktionen für den beabsichtigten Lebensschutz erscheint jedoch von vornherein als zweifelhaft. Auch die Mehrheit räumt ein, daß die bisherige generelle Strafbarkeit des Schwangerschaftsabbruchs das werdende Leben gerade nicht ausreichend geschützt und möglicherweise sogar dazu beigetragen hat, andere wirksame Schutzmaßnahmen zu vernachlässigen (D II). Sie glaubt - wenn sie sich dessen auch nicht ganz sicher ist (Vgl. D III) -, diesem Versagen des Strafschutzes werde durch eine differenzierende Strafdrohung abgeholfen, nach der der Schwangerschaftsabbruch in den im Gesetzgebungsverfahren erörterten Indikationsfällen straflos bleiben soll. Hinsichtlich der schon bisher anerkannten oder praktizierten medizinischen, ethischen und eugenischen Indikation bringt diese Indikationenlösung freilich keine wesentliche Veränderung des bisherigen unbefriedigenden Rechtszustandes. Eine wirkliche Differenzierung liegt nur in der Anerkennung der sozialen Indikation, sofern der Gesetzgeber bei der ihm obliegenden Abgrenzung nicht übermäßig streng verfährt und wenigstens hier die erwähnte Wechselbeziehung zwischen geschuldeter sozialer Hilfe und vertretbarem Strafen beachtet: Je weniger der Staat seinerseits zur Hilfe imstande ist, desto fragwürdiger und zugleich wirkungsloser sind Strafdrohungen gegenüber Frauen, die sich ihrerseits der Pflicht zum Austragen der Leibesfrucht nicht gewachsen fühlen.
\end{sourcepara}
\begin{transpara}
\transrn{36}a) 然而，刑罰制裁對預定生命保護是否適合，從一開始即顯有疑義。多數意見亦承認，迄今對終止妊娠之一概可罰性，並未充分保護正在形成之生命，甚至可能促成忽視其他有效保護措施（D II）。其相信，雖然並非全然確定（參見 D III），刑法保護之此一失敗，可透過一項差別化刑罰威嚇加以補救；依該威嚇，在立法程序中討論之指標事由案件中，終止妊娠應保持不罰。就迄今已承認或實行之醫學、倫理及優生指標事由而言，此一指標方案固然不會對既有不令人滿意之法狀態帶來重大變化。真正的差別化僅在於承認社會指標事由；其前提是立法者於其所負界定任務中不過度嚴格，並至少在此注意前述應提供之社會協助與可正當化刑罰間之相互關係：國家本身越無能力提供協助，對那些自認無力承擔將胎兒懷至足月義務之婦女而言，刑罰威嚇就越可疑且越無效。
\end{transpara}

\syncleft
\begin{sourcepara}
\sourcern{37}Die von der Mehrheit insgesamt zugunsten der Indikationenlösung angestellten Erwägungen verdienen rechtspolitisch gewiß Berücksichtigung. Verfassungsrechtlich ist aber entscheidend, daß sich bei realistischer Betrachtung auf keinem Wege, auch nicht mit differenzierenden Strafdrohungen ein lückenloser Lebensschutz erreichen läßt und daß daher keine Lösung verfassungsrechtlich "festgeschrieben" werden kann. Die Mehrheit bleibt schon den ihr obliegenden Nachweis schuldig, daß im Zeitalter des "Abtreibungstourismus" von innerstaatlichen Strafvorschriften ein günstiger Einfluß gerade auf solche Frauen zu erwarten ist, die ohne einsichtigen Grund zur Abtreibung entschlossen sind. Wenn überhaupt, so läßt sich ein solcher Erfolg nur in einer gewissen Zahl von Fällen - besonders bei Angehörigen sozial schwächerer Gruppen - erreichen. Bei an sich beeinflußbaren Frauen zeigt sich die ambivalente Wirkung von Strafdrohungen unter anderem darin, daß diese einerseits einen gewissen Rückhalt gegen das Abtreibungsverlangen des Erzeugers oder der Familie bieten mögen, andererseits zur Erhöhung der Aborte beitragen können, indem sie die Schwangere in die Isolierung treiben, dadurch erst recht solchen Pressionen aussetzen und zu Kurzschlußhandlungen veranlassen.
\end{sourcepara}
\begin{transpara}
\transrn{37}多數意見整體上為支持指標方案所提出之考量，在法政策上當然值得納入。但憲法上決定性者在於：現實觀察下，無論循何途徑，即使藉差別化刑罰威嚇，亦不能達成毫無缺口之生命保護；因此，任何解決方案均不能在憲法上被「固定化」。多數意見已未履行其所負證明責任：在「墮胎旅遊」時代，國內刑罰規定正可期待對那些無可理解理由而決意墮胎之婦女產生有利影響。即使存在此種效果，也只能在一定數量案件中達成，尤其是在社會較弱勢群體成員中。對於本可受影響之婦女，刑罰威嚇之雙重效果亦表現在：一方面，它或許能對抗生父或家庭之墮胎要求而提供某種支撐；另一方面，它也可能增加墮胎，因其將懷孕婦女推入孤立，使其更暴露於此類壓力，並促成一時衝動行為。
\end{transpara}

\syncleft
\begin{sourcepara}
\sourcern{38}b) Wie immer aber die Schutzwirkung von Strafdrohungen zu beurteilen sein mag, jedenfalls beruht ihre teilweise Rücknahme auf Erwägungen, die gerade unter dem Gesichtspunkt des Lebensschutzes Gewicht haben und sich - mindestens bei einer verbesserten Beratungsregelung - keinesfalls als offensichtlich fehlsam widerlegen lassen.
\end{sourcepara}
\begin{transpara}
\transrn{38}b）無論如何判斷刑罰威嚇的保護作用，其部分撤回是基於從保護生命的角度來看特別重要的考慮，並且至少在改進的諮詢規則下絕不能被駁斥為明顯錯誤。
\end{transpara}

\syncleft
\begin{sourcepara}
\sourcern{39}Der Gesetzgeber hatte bei seiner Konzeption das ganze Spektrum der Abtreibungsproblematik vor Augen, besonders die Vielzahl jener Schwangeren, die einer Beeinflussung zugänglich sind. Er durfte davon ausgehen, daß sich Frauen normalerweise nicht leichten Herzens und ohne Grund einem solchen Eingriff unterziehen. In aller Regel liegt ein ernst zu nehmender, jedenfalls verständlicher Konflikt vor; die Entscheidung zum Schwangerschaftsabbruch wird "in Tiefen der Persönlichkeit getroffen, in die der Appell des Strafgesetzes nicht eindringt" (Rolinski, Sonderausschuß für die Strafrechtsreform, 6. \abbrWp{}, 74., 75. und 76. Sitzung, \abbrStenBer{} \abbrS{} 2219.). Gerade in diesen Fällen setzt nach Auffassung des Gesetzgebers die erfolgreiche Verwirklichung der Beratungsregelung zwingend voraus, daß eine gleichzeitige Strafdrohung entfällt; denn zur Abtreibung geneigte Frauen werden die Beratungsstellen nicht aufsuchen, solange sie befürchten müssen, dadurch ihre Entscheidungsfreiheit zu verlieren und sich durch das Bekanntwerden ihrer Schwangerschaft im Falle eines späteren illegalen Eingriffs strafrechtlicher Verfolgung auszusetzen. Diese auf das Urteil vieler Sachverständiger gestützte und zudem der Lebenserfahrung entsprechende Auffassung ist weder von den Antragstellern in der mündlichen Verhandlung noch von der Senatsmehrheit widerlegt worden.
\end{sourcepara}
\begin{transpara}
\transrn{39}立法者在其構想中著眼於整個墮胎問題範圍，尤其是眾多可能受影響之懷孕婦女。其得出發於如下假設：婦女通常不會輕率且無理由地接受此類介入。通常存在一項應予認真看待、至少可理解之衝突；終止妊娠之決定，是「在人格深處作成，而刑法之訴求無法深入其間」（Rolinski, Sonderausschuß für die Strafrechtsreform, 6. Wp., 74., 75. und 76. Sitzung, StenBer. S. 2219）。正是在此等案件中，依立法者見解，諮詢規制之成功實現必然要求同時存在之刑罰威嚇消失；因為傾向墮胎之婦女，只要必須擔心因此喪失決定自由，並因懷孕被知悉而於日後非法介入時暴露於刑事追訴，即不會尋求諮詢機構。此一見解建立於眾多專家判斷之上，且符合生活經驗，未經言詞辯論中之聲請人或本庭多數意見駁倒。
\end{transpara}

\syncleft
\begin{sourcepara}
\sourcern{40}Der Gesetzgeber befand sich daher in dem Dilemma, daß sich nach seiner Beurteilung präventive Beratung und repressive Strafdrohung in ihrer lebensschützenden Wirkung teilweise ausschließen. Seine Erwägung, auf eine mögliche abortverhindernde Wirkung von Strafsanktionen in einer wahrscheinlich geringen Zahl von Fällen zu verzichten, um möglicherweise in einer größeren Zahl von Fällen anderes Leben zu retten, läßt sich nicht damit abtun, dies sei eine "pauschale Abwägung von Leben gegen Leben", die mit der verfassungsrechtlichen Pflicht zum individuellen Schutz jedes einzelnen ungeborenen Lebens unvereinbar sei (So D II 2 b-c). Mit dieser Argumentation verschließt sich die Mehrheit in schwer verständlicher Weise der Einsicht, daß sie selbst tut, was sie dem Gesetzgeber vorwirft. Denn sie nötigt ihrerseits den Gesetzgeber sogar von Verfassungs wegen zu einer Verrechnung, indem sie ihn durch die Forderung nach Beibehaltung der Strafvorschrift zwingt, solches ungeborene Leben schutzlos zu lassen, das bei einer Rücknahme der Strafdrohung durch geeignete Beratung bewahrt bleiben könnte.
\end{sourcepara}
\begin{transpara}
\transrn{40}因此，立法者處於一項兩難：依其判斷，預防性諮詢與壓制性刑罰威嚇，在其生命保護效果上部分相互排斥。其考量在可能較少數案件中放棄刑罰制裁之防止墮胎效果，以便可能在更多案件中拯救其他生命，不能以「生命對生命之概括衡量」為由逕予排除，並稱其與個別保護每一未出生生命之憲法義務不相容（如 D II 2 b-c）。多數意見以此論證，以難以理解之方式拒絕承認：其本身正在做其指責立法者之事。因為多數意見甚至基於憲法迫使立法者進行計算；其要求維持刑罰規定，因而迫使立法者使某些未出生生命不受保護，而此等生命在撤回刑罰威嚇後，本可能透過適當諮詢獲得保存。
\end{transpara}

\syncleft
\begin{sourcepara}
\sourcern{41}Der Rigorismus der Mehrheit verträgt sich zudem schwer mit der ausdrücklichen Zulassung einer Abwägung nicht nur von Leben gegen Leben, sondern sogar von Leben gegen minderrangige Rechtsgüter bei indizierten Schwangerschaftsabbrüchen. Soweit bei der Indikationenregelung über diese Abwägung staatlich ermächtigte Gutachterstellen befinden müssen, durfte der Gesetzgeber es als einen spezifischen Nachteil dieser Lösung werten, daß danach das Abtöten einer Leibesfrucht ausdrücklich amtlich legitimiert wird. Die Mehrheit läßt freilich für den Fall der sozialen Indikation offen, ob die Prüfung der Voraussetzungen vorweg durch Gutachterstellen vorzunehmen oder einem späteren Strafverfahren zu überlassen ist (C III 3). Der zweite Weg würde aber ein wesentliches Anliegen der Reform verfehlen, weil er eine rechtsstaatlich höchst bedenkliche Unsicherheit für die betroffenen Frauen und die beteiligten Ärzte zur Folge hätte.
\end{sourcepara}
\begin{transpara}
\transrn{41}多數意見之嚴格立場，亦難以與其明確允許在指標事由終止妊娠中進行權衡相容；該權衡不僅可能發生於生命與生命之間，甚至可能發生於生命與較低位階法益之間。就指標規制而言，若此種權衡必須由國家授權之鑑定機構作成，立法者本可將此解決方案之一項特有缺點評價為：殺害胎兒因而受到明示之官方正當化。誠然，多數意見就社會指標事由之情形，仍未決定其要件應預先由鑑定機構審查，或留待日後刑事程序處理（C III 3）。然而，第二種途徑將錯失改革之重要目標，因其會使受影響之婦女與參與之醫師承受一種在法治國上高度可疑之不確定性。
\end{transpara}

\syncleft
\begin{sourcepara}
\sourcern{42}5. Da nach alledem jede Lösung unter dem Gesichtspunkt des Lebensschutzes Stückwerk bleibt, durfte der Gesetzgeber zugunsten der Fristenlösung weitere - von der Mehrheit außer acht gelassene - verfassungsrechtliche, gesundheits- und kriminalpolitische Gesichtspunkte berücksichtigen. Er konnte besonders davon ausgehen, daß diese Regelung die Eigenverantwortung der Frau und Mutter in einer ihr Lebensschicksal betreffenden Frage am besten respektiert und vermeidet, sie den schon mit der Verfahrensprozedur vor einer Gutachterstelle verbundenen Eingriffen in ihren Persönlichkeitsbereich auszusetzen. Er durfte auch berücksichtigen, daß der Schutz des werdenden Lebens über die physische Existenz hinausgeht und die Lebenschancen für ein nach entsprechender Beratung von der Mutter eigenverantwortlich angenommenes Kind besser sind, als wenn sie es nur aus Angst vor Strafe austrägt. Wesentlich konnte ferner sein, daß die mit illegalen Abtreibungen verbundenen Gesundheitsschäden wegfallen und das Rechtsbewußtsein nicht mehr durch eine leer laufende Strafdrohung oder deren bagatellisierende Anwendung erschüttert wird.
\end{sourcepara}
\begin{transpara}
\transrn{42}5. 既然如此，從生命保護觀點出發的任何解決方案終究都只是片段性的，立法者遂得為支持期限規定，考量其他被多數意見忽略的憲法、衛生政策與刑事政策觀點。立法者尤其可以認為，這項規定最能尊重婦女及母親在關乎其人生際遇問題上的自主負責，並避免使其人格領域暴露於專家機構程序本身即已伴隨的干預之下。立法者亦得考量，對發展中生命的保護並不止於物理存在；相較於母親僅因害怕刑罰而將孩子懷至出生，經相應諮詢後由母親自主負責地接納的孩子，其生命機會更佳。再者，與非法墮胎相伴而生的健康損害得以消除，且法律意識不再因空轉的刑罰威嚇或其輕忽化適用而受到動搖，此亦可能具有重要性。
\end{transpara}

\syncleft
\begin{sourcepara}
\sourcern{43}Nicht zuletzt war es nicht offensichtlich fehlsam, wenn der Gesetzgeber, gestützt auf Erfahrungen im Ausland, einen wesentlichen Nachteil der Indikationenlösung darin sah, daß es als schwierig, wenn nicht unmöglich erscheint, objektivierbare, einheitliche Abgrenzungsmerkmale für die - unter dem Blickpunkt der Reform allein relevante - soziale Indikation zu finden (Vgl. u.a. \abbrBTDrucks{} 7/1981 [neu] \abbrS{} 12 sowie den unter A I 5 zitierten Alternativ-Entwurf, \abbrAAO{}, \abbrS{} 27.). Die starken Kontroversen bei den Gesetzberatungen haben offenbar gemacht, daß gerade in diesem Bereich kein Konsens über die Grenze des Zulässigen besteht. Voraussichtlich wird daher die behördliche Beurteilung darüber, wann die Gefahr einer schwerwiegenden sozialen Notlage vorliegt und welche anderen Maßnahmen zur Abwendung dieser Gefahr von der Schwangeren hinzunehmen sind, regional und nach der persönlichen Einstellung der Gutachter und Richter weit auseinandergehen. Das Ergebnis wäre eine schwer erträgliche Rechtsunsicherheit und Rechtsungleichheit für die betroffenen Frauen und die beteiligten Ärzte sowie weiterhin ein Ausweichen in die Illegalität.
\end{sourcepara}
\begin{transpara}
\transrn{43}最後，立法機關基於外國經驗，認為指標方案之一項重大缺點在於：要為社會指標事由找到客觀化且一致的界限特徵，即使並非不可能，也顯得困難，而此一社會指標事由正是在改革觀點下唯一相關者。此種看法並非顯然錯誤（參見例如 BTDrucks. 7/1981 [neu] S. 12，以及上文 A I 5 所引替代草案，a.a.O., S. 27）。立法審議中激烈爭議已顯示，正是在此一領域，對於容許界限並無共識。因此，行政機關判斷何時存在嚴重社會困境之危險，以及懷孕婦女為排除此一危險須忍受哪些其他措施，預期將因地區及鑑定人、法官之個人立場而有相當差異。其結果將是使受影響婦女及參與醫師難以承受之法律不確定與法律不平等，並且仍將逃向非法領域。
\end{transpara}

\syncleft
\begin{sourcepara}
\sourcern{44}Aus allen diesen Gründen durfte der Gesetzgeber den Versuch wagen, die jetzigen unhaltbaren Zustände im Wege der Beratungs- und Fristenlösung zu reformieren, auch wenn eine sichere Voraussage der weiteren Entwicklung nicht möglich ist. Da auch die Mehrheit zutreffend davon ausgeht, daß das bekannte Zahlenmaterial einen sicheren Schluß in der einen oder anderen Richtung nicht erlaubt (D II 2 c), bedarf es keines weiteren Eingehens auf die kritischen Äußerungen gegen die Prognose des Gesetzgebers (Nach neuesten Berichten aus der \abbrDDR{}, wo seit 1972 die Fristenregelung gilt, soll die Zahl der Schwangerschaftsunterbrechungen in den beiden letzten Jahren wesentlich zurückgegangen sein. Dies wird auf die intensiven staatlichen Maßnahmen zur Förderung junger Familien und den Ausbau der Ehe- und Sexualberatung zurückgeführt [vgl. Mehlan, Das deutsche Gesundheitswesen 1974, \abbrS{} 2216 \abbrFF{}].).
\end{sourcepara}
\begin{transpara}
\transrn{44}基於所有這些理由，立法機關即使無法可靠預測後續發展，仍得嘗試以諮詢方案與期限規定改革當前不可維持之狀態。既然多數意見亦正確地認為，既有統計資料並不容許朝任一方向作成確定結論（D II 2 c），則無須再進一步處理針對立法機關預測所提出之批評意見。（依來自德意志民主共和國之最新報告，該地自1972年起施行期限規定，最近兩年終止妊娠之數量據稱已大幅下降。此一情形被歸因於國家密集採取促進年輕家庭之措施，以及婚姻與性諮詢之擴充〔參見 Mehlan, Das deutsche Gesundheitswesen 1974, S. 2216 ff.〕。）
\end{transpara}

\bisubsubsection{II.}{II.}
\syncleft
\begin{sourcepara}
\sourcern{45}Die Mehrheit begründet die Aufrechterhaltung einer - differenzierten - Strafdrohung nachdrücklich damit, die verfassungsrechtlich geforderte "Mißbilligung" nichtindizierter Schwangerschaftsabbrüche müsse klar zum Ausdruck kommen (C II 3, C III 2 b, C III 3, D II 1, D II 2, D III.). Soweit damit die general-präventive Wirkung des Strafrechts angesprochen sein soll, d.h. das Bestreben, eine Tat durch Auferlegung eines Übels zu mißbilligen und dadurch das tatsächliche Verhalten der Rechtsunterworfenen zu beeinflussen, ist - wie erörtert - nicht dargetan, daß die Indikationenlösung ihrerseits effektiven Lebensschutz gewährleistet. Es ist daher vielleicht kein Zufall, daß die Mehrheit "zweispurig" argumentiert: Sie verlangt auch unabhängig von der angestrebten tatsächlichen Wirkung eine Mißbilligung als Ausdruck eines sozialethischen Unwerturteils, das die unmotivierten Abtreibungen deutlich als Unrecht kennzeichnet.
\end{sourcepara}
\begin{transpara}
\transrn{45}多數意見認為必須明確表達憲法規定的「反對」未經指標事由的終止妊娠，從而證明維持差別化的刑罰威嚇是合理的（C II 3、C III 2 b、C III 3、D II 1、D II 2、D III.）。就其目的而言，這是為了解決刑法的一般預防效果，即通過施加邪惡來反對某種行為，從而影響受法律約束的人的實際行為，但正如所討論的那樣，並沒有表明指標事由解決方案反過來保證了對生命的有效保護。因此，多數意見以「雙管齊下」的方式爭論也許並非巧合：無論預期的實際效果如何，他們都要求反對，以此表達對無價值的社會倫理判斷，這​​種判斷明確地將無動機的終止妊娠定性為不公正。
\end{transpara}

\syncleft
\begin{sourcepara}
\sourcern{46}1. Es kann dahingestellt bleiben, wieweit die neuere Strafrechtswissenschaft einer derartigen Auffassung über die Funktion des Strafrechts und seine Zuordnung zur Ethik zustimmt (Vgl. etwa Baumann, Strafrecht, Allgemeiner Teil, 6. \abbrAufl{}, 1974, \abbrS{} 7 \abbrF{}, 27 \abbrF{}, ders., Das Verhältnis von Moral und Recht, in Moral [hrsg. von Anselm Hertz, 1972] \abbrS{} 60 \abbrFF{}; Hanack, Verhandlungen des 47. \abbrDJT{} 1968, \abbrBd{} I, \abbrS{} A 29 \abbrFF{}; Sax, Grundsätze der Strafrechtspflege, in: Bettermann-Nipperdey-Scheuner, Die Grundrechte, \abbrBd{} III/2, 1959, \abbrS{} 930 \abbrF{}, 955 \abbrF{}; Arthur Kaufmann, Recht und Sittlichkeit, 1964, \abbrS{} 42 \abbrFF{}) und ob damit nicht doch das Strafrecht zum Selbstzweck erhoben wird. Selbstverständlich sind unmotivierte Schwangerschaftsabbrüche ethisch verwerflich. Gegenüber der Mehrheitsbegründung ist aber zunächst zu bedenken, daß das Absehen von Strafe hier wie auch sonst nicht den Schluß aufdrängt, ein nicht mehr strafbares Verhalten werde gebilligt. Hierfür ist namentlich dann kein Raum, wenn der Gesetzgeber eine Strafvorschrift aufhebt, weil sie nach seiner Meinung wirkungslos oder sogar schädlich ist oder weil dem bisher strafbaren, sozialschädlichen Verhalten auf andere Weise begegnet werden soll. So wird wohl niemand aus der Aufhebung oder Einschränkung der Strafvorschriften gegen Prostitution, Drogenmißbrauch, Ehebruch oder Ehegattenkuppelei schließen, entsprechende Handlungen erfreuten sich nunmehr der offiziellen Anerkennung als rechtmäßig und sittlich. Die Auseinandersetzungen um die Reform des § 218 \abbrStGB{} bieten keinen Anhaltspunkt dafür, als sei das Abtöten ungeborenen Lebens im Ernst als "normaler sozialer Vorgang" angesehen worden.
\end{sourcepara}
\begin{transpara}
\transrn{46}1. 近代刑法學在多大程度上同意此種關於刑法功能及其與倫理關係之見解，可以不予決定（參見例如 Baumann, Strafrecht, Allgemeiner Teil, 6. Aufl., 1974, S. 7 f., 27 f.; ders., Das Verhältnis von Moral und Recht, in Moral [hrsg. von Anselm Hertz, 1972] S. 60 ff.; Hanack, Verhandlungen des 47. DJT 1968, Bd. I, S. A 29 ff.; Sax, Grundsätze der Strafrechtspflege, in: Bettermann-Nipperdey-Scheuner, Die Grundrechte, Bd. III/2, 1959, S. 930 f., 955 f.; Arthur Kaufmann, Recht und Sittlichkeit, 1964, S. 42 ff.）；也可以不問此種見解是否終究將刑法提升為自我目的。無疑地，無動機之終止妊娠在倫理上應受非難。然而，面對多數意見之理由，首先必須考慮：在此處以及在其他情形中，不予處罰並不必然導出一項結論，即一項不再受罰之行為即被認可。尤其在立法機關廢止某一刑罰條項，是因為其認為該條項無效、甚至有害，或因為要以其他方式處理過去受罰且有社會危害性之行為時，更無從作此結論。因此，恐怕沒有人會從廢止或限制關於賣淫、藥物濫用、通姦或配偶間媒介性行為之刑罰規定，推論出相應行為如今已獲官方承認為合法且合乎道德。圍繞《刑法》第218條改革之爭論，並無任何根據足以表明，殺害未出生生命曾被嚴肅地視為「正常的社會過程」。
\end{transpara}

\syncleft
\begin{sourcepara}
\sourcern{47}Soweit die Mehrheit in diesem Zusammenhang zur Beurteilung der Beratungs- und Fristenregelung im Fünften Strafrechtsreformgesetz das bisher von den gesetzgebenden Körperschaften noch nicht abschließend behandelte Strafrechtsreform-Ergänzungsgesetz heranzieht (D II 1) kommt es darauf schon deswegen nicht an, weil beide Gesetze auch nach Auffassung der Mehrheit inhaltlich voneinander unabhängig sind (B 4). Erst wenn das Strafrechtsreform-Ergänzungsgesetz erlassen würde, wäre selbständig zu prüfen, ob die geplante generelle Kostenerstattung und Lohnfortzahlung bei nicht strafbaren Schwangerschaftsabbrüchen eine unzulässige staatliche Förderung für die nicht indizierten Fälle enthielte oder ob dies aus bestimmten gewichtigen Gründen noch hinzunehmen wäre, etwa wegen der Bekämpfung der mit illegalen Abtreibungen verbundenen Gesundheitsgefahren, die den Supreme Court sogar zum verfassungsrechtlichen Verbot des Strafens veranlaßt haben (410 \abbrUS{} 113 [148 \abbrFF{}, 162 \abbrFF{}].). Im ersten Fall ließe sich dieser Mangel innerhalb des Strafrechtsreform-Ergänzungsgesetzes korrigieren, z.B. durch Beschränkung der Kostenerstattung auf indizierte Fälle, wobei die erforderliche Prüfung auch nach dem Schwangerschaftsabbruch, also ohne Zeitdruck vorgenommen werden könnte. (Auf diese Weise hätte sich übrigens zugleich die gewünschte Mißbilligung unmotivierter Abtreibungen erreichen lassen).
\end{sourcepara}
\begin{transpara}
\transrn{47}多數意見在此脈絡中，為評價《第五刑法改革法》之諮詢與期限規定，而援引迄今尚未由立法機關最終處理之《刑法改革補充法》（D II 1）；然而，僅因兩部法律即使依多數意見亦在內容上彼此獨立（B 4），此點即已無關緊要。只有在《刑法改革補充法》制定後，才須獨立審查：對不受刑罰之終止妊娠所計畫的一般費用補償及工資繼續支付，是否對非指標事由案件構成不容許之國家促進；或是否基於特定重大理由仍可容忍，例如為防治非法墮胎所伴隨之健康危險，而此等危險甚至曾促使美國最高法院作成憲法上禁止處罰之判斷（410 U.S. 113 [148 ff., 162 ff.]）。在前一種情形中，此一缺陷可於《刑法改革補充法》內加以修正，例如將費用補償限於指標事由案件；所需審查亦可於終止妊娠後進行，因而不受時間壓力。（順帶一提，藉此也可同時達成對無動機墮胎所期望之否定評價。）
\end{transpara}

\syncleft
\begin{sourcepara}
\sourcern{48}2. Unser wesentlicher Einwand richtet sich dagegen, daß die Mehrheit nicht darlegt, woraus verfassungsrechtlich das Erfordernis der Mißbilligung als einer selbständigen Pflicht hergeleitet werden soll. Nach unserer Auffassung schreibt die Verfassung nirgends vor, ethisch verwerfliches oder strafwürdiges Verhalten müsse per se Rücksicht auf den damit erzielten Effekt mit Hilfe des Gesetzesrechts mißbilligt werden. In einem pluralistischen, weltanschaulich neutralen und freiheitlichen demokratischen Gemeinwesen bleibt es den gesellschaftlichen Kräften überlassen, Gesinnungspostulate zu statuieren. Der Staat hat darin Enthaltsamkeit zu üben; seine Aufgabe ist der Schutz der von der Verfassung gewährleisteten und anerkannten Rechtsgüter. Für die verfassungsrechtliche Entscheidung kommt es allein darauf an, ob die Strafvorschrift zwingend geboten ist, um einen effektiven Schutz des werdenden Lebens unter Berücksichtigung der schutzwürdigen Interessen der Frau zu sichern.
\end{sourcepara}
\begin{transpara}
\transrn{48}2. 我們主要反對的是，多數意見並未說明，要求作為獨立義務之否定評價，如何能由憲法導出。我們認為，憲法並未規定，不問其實際保護效果如何，均須藉助法律規範本身否定評價倫理上不值得尊重之行為，或應受刑罰之行為。在一個多元、意識形態中立且自由民主之社會中，形塑輿論預設由社會力量承擔。國家在此必須自我節制；其任務是保護受憲法保障及承認之法益。就憲法裁判而言，唯一重要者是：考量婦女值得保護之利益後，刑罰規定是否為確保正在發展中生命之有效保護所不可或缺。
\end{transpara}

\bisubsubsection{III.}{III.}
\syncleft
\begin{sourcepara}
\sourcern{49}Daß die Entscheidung des deutschen Gesetzgebers für die Fristen- und Beratungsregelung weder einer sittlich noch rechtlich zu mißbilligenden Grundhaltung entspringt noch von offensichtlich falschen Prämissen in der Beurteilung der Lebensverhältnisse ausgeht, wird durch gleiche oder ähnliche Reformvorschriften in zahlreichen ausländischen Staaten bestätigt. In Österreich, Frankreich, Dänemark und Schweden ist der von einem Arzt mit Einwilligung der Schwangeren vorgenommene Schwangerschaftsabbruch in den ersten zwölf (in Frankreich zehn) Wochen der Schwangerschaft nicht strafbar; in Großbritannien und den Niederlanden gilt eine Indikationenregelung, die in der praktischen Anwendung auf das gleiche hinausläuft (Für die Vereinigten Staaten vgl. oben A II 1 der abweichenden Meinung.). Diese Staaten können sich zum Teil einer eindrucksvollen Verfassungstradition rühmen und stehen sämtlich in der unbedingten Achtung vor dem Leben jedes einzelnen Menschen der Bundesrepublik gewiß nicht nach; einige von ihnen haben ebenfalls geschichtliche Erfahrungen mit einem menschenverachtenden Unrechtssystem. Ihre Entscheidung erforderte eine Auseinandersetzung mit den gleichen rechtlichen und sozialen Problemen wie in der Bundesrepublik. In allen diesen Staaten gilt zudem rechtsverbindlich die Europäische Menschenrechtskonvention, deren Artikel 2 in Absatz 1 ("Das Recht jedes Menschen auf das Leben wird gesetzlich geschützt") der Verfassungsvorschrift des \abbrArt{} 2 \abbrAbs{} 2 Satz 1 \abbrGG{}\ggfnArtTwoTwoOneDE{} nahekommt und der im ganzen eher weitergehen könnte als die innerdeutsche Norm. Der Österreichische Verfassungsgerichtshof hat ausdrücklich festgestellt, daß die dortige Fristenlösung mit der Menschenrechtskonvention, der in Österreich Verfassungsrang zukommt, vereinbar ist (A.a.O. unter II 3 b der Entscheidungsgründe, \abbrEuGRZ{}, 1975, \abbrS{} 74 [77 \abbrF{}]. Auch in Frankreich hat die Konvention vor innerfranzösischen Gesetzen, vgl. \abbrArt{} 55 der französischen Verfassung, s. auch die Entscheidung des Conseil constitutionnel vom 15. Januar 1975, Journal officiel vom 16. Januar 1975 \abbrS{} 671 = \abbrEuGRZ{} 1975, \abbrS{} 54.).
\end{sourcepara}
\begin{transpara}
\transrn{49}德國立法機關選擇期限與諮詢規制，既非出於道德上或法律上應受非難之基本態度，也非以對生活狀況之顯然錯誤判斷為前提；此點可由許多外國相同或類似之改革規定加以確認。在奧地利、法國、丹麥及瑞典，由醫師經懷孕婦女同意，在妊娠前十二週（法國為十週）內實施之終止妊娠不受處罰；在英國及荷蘭，則適用一套在實際運用上導向相同結果之指標規制（美國部分參見本不同意見上文 A II 1）。這些國家中有些擁有令人印象深刻之憲政傳統，且其對每一個人生命之無條件尊重，當然不遜於聯邦共和國；其中若干國家同樣有過蔑視人性之不法體制的歷史經驗。其決定也要求處理與聯邦共和國相同之法律與社會問題。此外，歐洲人權公約對所有這些國家均具法律拘束力；其第2條第1項「每個人的生命權受法律保護」接近《基本法》第2條第2項第1句之憲法規定，整體而言甚至可能比德國國內規範更進一步。奧地利憲法法院已明確認定，當地期限規定與歐洲人權公約相容，而該公約在奧地利具有憲法位階（a.a.O. unter II 3 b der Entscheidungsgründe, EuGRZ 1975, S. 74 [77 f.]）。在法國，該公約亦優先於法國國內法律，參見法國憲法第55條；另見 Conseil constitutionnel 1975年1月15日之決定，Journal officiel vom 16. Januar 1975 S. 671 = EuGRZ 1975, S. 54。
\end{transpara}

\bisubsubsection{IV.}{IV.}
\syncleft
\begin{sourcepara}
\sourcern{50}Insgesamt war es daher nach unserer Ansicht dem Gesetzgeber durch die Verfassung nicht verwehrt, auf eine nach seiner unwiderlegten Auffassung weitgehend wirkungslose, inadäquate und sogar schädliche Strafdrohung zu verzichten. Sein Versuch, dem in den gegenwärtigen Zuständen offenbar werdenden Unvermögen von Staat und Gesellschaft im Dienste des Lebensschutzes durch sozial adäquatere Mittel abzuhelfen, mag unvollkommen sein; er entspricht aber dem Geist des Grundgesetzes mehr als die Forderung nach Strafe und Mißbilligung.
\end{sourcepara}
\begin{transpara}
\transrn{50}總之，我們認為，憲法並不阻止立法者放棄其認為大致無效、不充分甚至有害之刑罰威嚇。其試圖以更充分之社會手段，至少部分補償國家與社會目前顯然無力保護生命之狀態；此一嘗試或許仍不完美，但較刑罰與否定評價之命令更符合《基本法》精神。
\end{transpara}

\syncleft
\begin{sourcepara}
{\centering Rupp-v. Brünneck, \abbrDr{} Simon\par}
\end{sourcepara}
\begin{transpara}
{\centering Rupp-v. Brünneck、Dr. Simon\par}
\end{transpara}

\bilingualstop

\phantomsection
\addcontentsline{toc}{section}{Glossar / 譯語表}
\ifbilingualcolumns
  \begin{center}
  \begin{tabular}{@{}p{0.44\textwidth}@{\hspace{0.52cm}}|@{\hspace{0.52cm}}p{0.44\textwidth}@{}}
  \makebox[\linewidth][c]{\begin{minipage}{0.86\linewidth}
  \centering
  {\Large\bfseries Glossar\par}
  \vspace{0.28em}
  {\footnotesize\color{pendingink}Ausgewählte Terminologie des Urteils\par}
  \end{minipage}}
  &
  \makebox[\linewidth][c]{\begin{minipage}{0.86\linewidth}
  \centering\songfont
  {\Large\bfseries 譯語表\par}
  \vspace{0.28em}
  {\footnotesize\color{pendingink}本文件採用之主要法律術語對照\par}
  \end{minipage}}
  \end{tabular}
  \end{center}
  \vspace{0.45em}
\else
  \ifshowtranslation
    \section*{譯語表}
  \else
    \section*{Glossar}
  \fi
\fi

\begin{termbox}
\noindent{\songfont\footnotesize\color{pendingink}
譯語說明：本文正文原則上將 \textit{Schwangerschaftsabbruch} 譯為「終止妊娠」；若德文原文使用 \textit{Abtreibung}，或語境須凸顯較具評價色彩之用語時，譯為「墮胎」或於語境中另行標示。\par}
\vspace{0.45em}
\begin{termcolumns}
\begin{termitems}
\item[Schwangerschaftsabbruch] 終止妊娠
\item[Abtreibung] 墮胎；原文帶非正式或貶抑色彩時保留此譯
\item[Fristenregelung] 期限規定
\item[Fristenlösung] 期限規定
\item[Indikation] 指標事由
\item[Indikationenregelung / Indikationenlösung] 指標事由規定；指標方案
\item[medizinische Indikation] 醫學指標事由
\item[eugenische Indikation] 優生指標事由
\item[ethische Indikation] 倫理指標事由
\item[kriminologische Indikation] 犯罪學指標事由
\item[soziale Indikation / Notlagenindikation] 社會指標事由；困境指標事由
\item[werdendes Leben / sich entwickelndes Leben] 正在發展中的生命
\item[ungeborenes Leben] 未出生生命
\item[Leibesfrucht] 胎兒；依語境亦譯「子宮內胎兒」
\item[Nasciturus] \textit{nasciturus}；胎兒；未出生者
\item[Lebensschutz] 生命保護
\item[Schutzpflicht des Staates] 國家保護義務
\item[Rechtsgut] 法益
\item[Selbstbestimmungsrecht] 自決權
\item[Persönlichkeitsbereich / Persönlichkeitssphäre] 人格領域；人格私密領域
\item[Menschenwürde] 人性尊嚴
\item[Wertordnung des Grundgesetzes] 基本法的價值秩序
\item[Mißbilligung / Unwerturteil] 否定評價；負面價值判斷
\item[Strafbarkeit] 可罰性；刑事責任
\item[Straffreiheit / straflos] 不罰；不受刑罰處罰
\item[Strafdrohung] 刑罰威嚇
\item[Strafsanktion] 刑事制裁
\item[Strafgewalt] 刑罰權力
\item[Beratung / Beratungsregelung] 諮詢；諮詢規定
\item[Beratungsstelle] 諮詢機構
\item[Gutachterstelle] 鑑定機構
\item[Zumutbarkeit / unzumutbar] 可期待性；不可期待
\item[Notlage] 困境；嚴重困境
\item[Gefahr] 危險
\item[Beeinträchtigung des Gesundheitszustandes] 健康狀態之損害
\item[Empfängnis] 受孕
\item[seit / nach der Empfängnis] 受孕後（如：受孕後十二週）；5. StrRG 採此計算基準
\item[Schwangerschaftswochen] 妊娠週數（自末次月經起算，較受孕早約兩週）；如：妊娠最初十二週
\item[Nidation] 著床
\item[Austragen der Leibesfrucht] 將胎兒孕育至出生
\item[Gesetzgeber] 立法者；立法機關
\item[Bundestag] 聯邦議院
\item[Bundesrat] 聯邦參議院
\item[Bundesverfassungsgericht] 聯邦憲法法院
\item[Erster Senat] 第一庭
\item[abweichende Meinung] 不同意見
\item[Antrag（向法院）] 聲請
\item[Antrag（議會動議）] 提案；對應的 Ablehnung 譯「否決」
\item[Antrag（一般行政）] 申請
\item[Antragsteller] 聲請人
\item[einstweilige Anordnung] 暫時處分
\item[Entscheidungsformel / Tenor] 主文
\item[Entscheidungsgründe / Gründe] 理由
\end{termitems}
\end{termcolumns}

\termsectionheading{德文用語簡稱}
\begin{termcolumns}
\begin{abbritems}
\item[Abs.] \textit{Absatz}；項
\item[AE] \textit{Alternativ-Entwurf}；替代草案
\item[Art.] \textit{Artikel}；條
\item[BGBl.] \textit{Bundesgesetzblatt}；聯邦法律公報
\item[BGHSt] \textit{Entscheidungen des Bundesgerichtshofes in Strafsachen}；聯邦普通法院刑事判決彙編
\item[BRDrucks.] \textit{Bundesrats-Drucksache}；聯邦參議院文件
\item[BTDrucks.] \textit{Bundestags-Drucksache}；聯邦議院文件
\item[BVerfG] \textit{Bundesverfassungsgericht}；聯邦憲法法院
\item[BVerfGE] \textit{Entscheidungen des Bundesverfassungsgerichts}；聯邦憲法法院判決彙編
\item[BVerfGG] \textit{Bundesverfassungsgerichtsgesetz}；聯邦憲法法院法
\item[DDR] \textit{Deutsche Demokratische Republik}；德意志民主共和國
\item[DJT] \textit{Deutscher Juristentag}；德國法學家大會
\item[EuGRZ] \textit{Europäische Grundrechte-Zeitschrift}；歐洲基本權期刊
\item[f. / ff.] \textit{folgende / fortfolgende}；以下一頁；以下數頁
\item[GG] \textit{Grundgesetz}；德國基本法
\item[JR] \textit{Juristische Rundschau}；法學評論
\item[JZ] \textit{JuristenZeitung}；法學家報
\item[MdB] \textit{Mitglied des Bundestages}；聯邦議院議員
\item[Nr.] \textit{Nummer}；款、目或號，依語境
\item[RGBl.] \textit{Reichsgesetzblatt}；帝國法律公報
\item[RGSt] \textit{Entscheidungen des Reichsgerichts in Strafsachen}；帝國法院刑事判決彙編
\item[S.] \textit{Seite / Satz}；頁；於條文引用中亦可能為句
\item[StenBer.] \textit{Stenographischer Bericht}；速記紀錄
\item[StGB] \textit{Strafgesetzbuch}；德國刑法典
\item[StrRG] \textit{Strafrechtsreformgesetz}；刑法改革法
\item[U.S.] \textit{United States Reports}；美國最高法院判決彙編
\item[Wp.] \textit{Wahlperiode}；屆期
\item[ZStrW] \textit{Zeitschrift für die gesamte Strafrechtswissenschaft}；整體刑法學期刊
\end{abbritems}
\end{termcolumns}
\end{termbox}

\end{document}
